% Options for packages loaded elsewhere
% Options for packages loaded elsewhere
\PassOptionsToPackage{unicode}{hyperref}
\PassOptionsToPackage{hyphens}{url}
\PassOptionsToPackage{dvipsnames,svgnames,x11names}{xcolor}
%
\documentclass[
  spanish,
  11pt,
  letterpaper,
  oneside,
  openany]{book}
\usepackage{xcolor}
\usepackage[margin=1in]{geometry}
\usepackage{amsmath,amssymb}
\setcounter{secnumdepth}{3}
\usepackage{iftex}
\ifPDFTeX
  \usepackage[T1]{fontenc}
  \usepackage[utf8]{inputenc}
  \usepackage{textcomp} % provide euro and other symbols
\else % if luatex or xetex
  \usepackage{unicode-math} % this also loads fontspec
  \defaultfontfeatures{Scale=MatchLowercase}
  \defaultfontfeatures[\rmfamily]{Ligatures=TeX,Scale=1}
\fi
\usepackage{lmodern}
\ifPDFTeX\else
  % xetex/luatex font selection
\fi
% Use upquote if available, for straight quotes in verbatim environments
\IfFileExists{upquote.sty}{\usepackage{upquote}}{}
\IfFileExists{microtype.sty}{% use microtype if available
  \usepackage[]{microtype}
  \UseMicrotypeSet[protrusion]{basicmath} % disable protrusion for tt fonts
}{}
\usepackage{setspace}
\makeatletter
\@ifundefined{KOMAClassName}{% if non-KOMA class
  \IfFileExists{parskip.sty}{%
    \usepackage{parskip}
  }{% else
    \setlength{\parindent}{0pt}
    \setlength{\parskip}{6pt plus 2pt minus 1pt}}
}{% if KOMA class
  \KOMAoptions{parskip=half}}
\makeatother
% Make \paragraph and \subparagraph free-standing
\makeatletter
\ifx\paragraph\undefined\else
  \let\oldparagraph\paragraph
  \renewcommand{\paragraph}{
    \@ifstar
      \xxxParagraphStar
      \xxxParagraphNoStar
  }
  \newcommand{\xxxParagraphStar}[1]{\oldparagraph*{#1}\mbox{}}
  \newcommand{\xxxParagraphNoStar}[1]{\oldparagraph{#1}\mbox{}}
\fi
\ifx\subparagraph\undefined\else
  \let\oldsubparagraph\subparagraph
  \renewcommand{\subparagraph}{
    \@ifstar
      \xxxSubParagraphStar
      \xxxSubParagraphNoStar
  }
  \newcommand{\xxxSubParagraphStar}[1]{\oldsubparagraph*{#1}\mbox{}}
  \newcommand{\xxxSubParagraphNoStar}[1]{\oldsubparagraph{#1}\mbox{}}
\fi
\makeatother

\usepackage{color}
\usepackage{fancyvrb}
\newcommand{\VerbBar}{|}
\newcommand{\VERB}{\Verb[commandchars=\\\{\}]}
\DefineVerbatimEnvironment{Highlighting}{Verbatim}{commandchars=\\\{\}}
% Add ',fontsize=\small' for more characters per line
\usepackage{framed}
\definecolor{shadecolor}{RGB}{241,243,245}
\newenvironment{Shaded}{\begin{snugshade}}{\end{snugshade}}
\newcommand{\AlertTok}[1]{\textcolor[rgb]{0.68,0.00,0.00}{#1}}
\newcommand{\AnnotationTok}[1]{\textcolor[rgb]{0.37,0.37,0.37}{#1}}
\newcommand{\AttributeTok}[1]{\textcolor[rgb]{0.40,0.45,0.13}{#1}}
\newcommand{\BaseNTok}[1]{\textcolor[rgb]{0.68,0.00,0.00}{#1}}
\newcommand{\BuiltInTok}[1]{\textcolor[rgb]{0.00,0.23,0.31}{#1}}
\newcommand{\CharTok}[1]{\textcolor[rgb]{0.13,0.47,0.30}{#1}}
\newcommand{\CommentTok}[1]{\textcolor[rgb]{0.37,0.37,0.37}{#1}}
\newcommand{\CommentVarTok}[1]{\textcolor[rgb]{0.37,0.37,0.37}{\textit{#1}}}
\newcommand{\ConstantTok}[1]{\textcolor[rgb]{0.56,0.35,0.01}{#1}}
\newcommand{\ControlFlowTok}[1]{\textcolor[rgb]{0.00,0.23,0.31}{\textbf{#1}}}
\newcommand{\DataTypeTok}[1]{\textcolor[rgb]{0.68,0.00,0.00}{#1}}
\newcommand{\DecValTok}[1]{\textcolor[rgb]{0.68,0.00,0.00}{#1}}
\newcommand{\DocumentationTok}[1]{\textcolor[rgb]{0.37,0.37,0.37}{\textit{#1}}}
\newcommand{\ErrorTok}[1]{\textcolor[rgb]{0.68,0.00,0.00}{#1}}
\newcommand{\ExtensionTok}[1]{\textcolor[rgb]{0.00,0.23,0.31}{#1}}
\newcommand{\FloatTok}[1]{\textcolor[rgb]{0.68,0.00,0.00}{#1}}
\newcommand{\FunctionTok}[1]{\textcolor[rgb]{0.28,0.35,0.67}{#1}}
\newcommand{\ImportTok}[1]{\textcolor[rgb]{0.00,0.46,0.62}{#1}}
\newcommand{\InformationTok}[1]{\textcolor[rgb]{0.37,0.37,0.37}{#1}}
\newcommand{\KeywordTok}[1]{\textcolor[rgb]{0.00,0.23,0.31}{\textbf{#1}}}
\newcommand{\NormalTok}[1]{\textcolor[rgb]{0.00,0.23,0.31}{#1}}
\newcommand{\OperatorTok}[1]{\textcolor[rgb]{0.37,0.37,0.37}{#1}}
\newcommand{\OtherTok}[1]{\textcolor[rgb]{0.00,0.23,0.31}{#1}}
\newcommand{\PreprocessorTok}[1]{\textcolor[rgb]{0.68,0.00,0.00}{#1}}
\newcommand{\RegionMarkerTok}[1]{\textcolor[rgb]{0.00,0.23,0.31}{#1}}
\newcommand{\SpecialCharTok}[1]{\textcolor[rgb]{0.37,0.37,0.37}{#1}}
\newcommand{\SpecialStringTok}[1]{\textcolor[rgb]{0.13,0.47,0.30}{#1}}
\newcommand{\StringTok}[1]{\textcolor[rgb]{0.13,0.47,0.30}{#1}}
\newcommand{\VariableTok}[1]{\textcolor[rgb]{0.07,0.07,0.07}{#1}}
\newcommand{\VerbatimStringTok}[1]{\textcolor[rgb]{0.13,0.47,0.30}{#1}}
\newcommand{\WarningTok}[1]{\textcolor[rgb]{0.37,0.37,0.37}{\textit{#1}}}

\usepackage{longtable,booktabs,array}
\usepackage{calc} % for calculating minipage widths
% Correct order of tables after \paragraph or \subparagraph
\usepackage{etoolbox}
\makeatletter
\patchcmd\longtable{\par}{\if@noskipsec\mbox{}\fi\par}{}{}
\makeatother
% Allow footnotes in longtable head/foot
\IfFileExists{footnotehyper.sty}{\usepackage{footnotehyper}}{\usepackage{footnote}}
\makesavenoteenv{longtable}
\usepackage{graphicx}
\makeatletter
\newsavebox\pandoc@box
\newcommand*\pandocbounded[1]{% scales image to fit in text height/width
  \sbox\pandoc@box{#1}%
  \Gscale@div\@tempa{\textheight}{\dimexpr\ht\pandoc@box+\dp\pandoc@box\relax}%
  \Gscale@div\@tempb{\linewidth}{\wd\pandoc@box}%
  \ifdim\@tempb\p@<\@tempa\p@\let\@tempa\@tempb\fi% select the smaller of both
  \ifdim\@tempa\p@<\p@\scalebox{\@tempa}{\usebox\pandoc@box}%
  \else\usebox{\pandoc@box}%
  \fi%
}
% Set default figure placement to htbp
\def\fps@figure{htbp}
\makeatother



\ifLuaTeX
\usepackage[bidi=basic]{babel}
\else
\usepackage[bidi=default]{babel}
\fi
% get rid of language-specific shorthands (see #6817):
\let\LanguageShortHands\languageshorthands
\def\languageshorthands#1{}


\setlength{\emergencystretch}{3em} % prevent overfull lines

\providecommand{\tightlist}{%
  \setlength{\itemsep}{0pt}\setlength{\parskip}{0pt}}



 
\usepackage[]{natbib}
\bibliographystyle{plainnat}


\makeatletter
\@ifpackageloaded{tcolorbox}{}{\usepackage[skins,breakable]{tcolorbox}}
\@ifpackageloaded{fontawesome5}{}{\usepackage{fontawesome5}}
\definecolor{quarto-callout-color}{HTML}{909090}
\definecolor{quarto-callout-note-color}{HTML}{0758E5}
\definecolor{quarto-callout-important-color}{HTML}{CC1914}
\definecolor{quarto-callout-warning-color}{HTML}{EB9113}
\definecolor{quarto-callout-tip-color}{HTML}{00A047}
\definecolor{quarto-callout-caution-color}{HTML}{FC5300}
\definecolor{quarto-callout-color-frame}{HTML}{acacac}
\definecolor{quarto-callout-note-color-frame}{HTML}{4582ec}
\definecolor{quarto-callout-important-color-frame}{HTML}{d9534f}
\definecolor{quarto-callout-warning-color-frame}{HTML}{f0ad4e}
\definecolor{quarto-callout-tip-color-frame}{HTML}{02b875}
\definecolor{quarto-callout-caution-color-frame}{HTML}{fd7e14}
\makeatother
\makeatletter
\@ifpackageloaded{bookmark}{}{\usepackage{bookmark}}
\makeatother
\makeatletter
\@ifpackageloaded{caption}{}{\usepackage{caption}}
\AtBeginDocument{%
\ifdefined\contentsname
  \renewcommand*\contentsname{Tabla de contenidos}
\else
  \newcommand\contentsname{Tabla de contenidos}
\fi
\ifdefined\listfigurename
  \renewcommand*\listfigurename{Listado de Figuras}
\else
  \newcommand\listfigurename{Listado de Figuras}
\fi
\ifdefined\listtablename
  \renewcommand*\listtablename{Listado de Tablas}
\else
  \newcommand\listtablename{Listado de Tablas}
\fi
\ifdefined\figurename
  \renewcommand*\figurename{Figura}
\else
  \newcommand\figurename{Figura}
\fi
\ifdefined\tablename
  \renewcommand*\tablename{Tabla}
\else
  \newcommand\tablename{Tabla}
\fi
}
\@ifpackageloaded{float}{}{\usepackage{float}}
\floatstyle{ruled}
\@ifundefined{c@chapter}{\newfloat{codelisting}{h}{lop}}{\newfloat{codelisting}{h}{lop}[chapter]}
\floatname{codelisting}{Cuadro}
\newcommand*\listoflistings{\listof{codelisting}{Listado de Listados}}
\makeatother
\makeatletter
\makeatother
\makeatletter
\@ifpackageloaded{caption}{}{\usepackage{caption}}
\@ifpackageloaded{subcaption}{}{\usepackage{subcaption}}
\makeatother
\usepackage{bookmark}
\IfFileExists{xurl.sty}{\usepackage{xurl}}{} % add URL line breaks if available
\urlstyle{same}
\hypersetup{
  pdftitle={Matemáticas para la Estadística Médica con R},
  pdfauthor={Miguel Ángel Luque Fernández},
  pdflang={es},
  colorlinks=true,
  linkcolor={blue},
  filecolor={Maroon},
  citecolor={blue},
  urlcolor={blue},
  pdfcreator={LaTeX via pandoc}}


\title{Matemáticas para la Estadística Médica con R}
\usepackage{etoolbox}
\makeatletter
\providecommand{\subtitle}[1]{% add subtitle to \maketitle
  \apptocmd{\@title}{\par {\large #1 \par}}{}{}
}
\makeatother
\subtitle{Bioestadística --- Universidad de Granada}
\author{Miguel Ángel Luque Fernández}
\date{2026}
\begin{document}
\frontmatter
\maketitle

\renewcommand*\contentsname{Contenidos}
{
\hypersetup{linkcolor=}
\setcounter{tocdepth}{2}
\tableofcontents
}

\setstretch{1.15}
\mainmatter
\bookmarksetup{startatroot}

\chapter*{Prefacio}\label{prefacio}
\addcontentsline{toc}{chapter}{Prefacio}

\markboth{Prefacio}{Prefacio}

Este libro contiene el material docente del curso \textbf{Matemáticas
para la Estadística Médica con R} desarrollado por
{\href{https://migariane.github.io}{\textbf{Miguel Ángel Luque
Fernández}}} en el \href{https://estadistica.ugr.es}{Departamento de
Estadística e Investigación Operativa} de la
\href{https://www.ugr.es}{Universidad de Granada}.

\section*{Objetivos del Curso}\label{objetivos-del-curso}
\addcontentsline{toc}{section}{Objetivos del Curso}

\markright{Objetivos del Curso}

El objetivo principal de este curso y material docente es proporcionar
al estudiante de ciencias de la salud los fundamentos matemáticos y
estadísticos necesarios para el análisis cuantitativo de datos médicos
con un alto grado de rigor académico y con ejemplos contextualizados en
la práctica clínica. El libro se sustenta de la filosofía de ciencia
reproducible, abierta y computacionalmente aplicada usando el software
libre R. Al final del libro se presentan cuatro apéndices reforzando a
los alumnos que quieran repasar sobre: i) los contenidos matemáticos
básicos necesarios para el buen aprovechamiento del curso, ii) una
introducción básica a la programación con R, iii) una introducción al
paquete \textbf{BioStatR} producido para la docencia específica de este
curso cuyo interés radica en la rica interpretación de los resultados de
las funciones del paquete y iv) una introducción a la
\textbf{investigación reproducible} mediante el control de versiones con
Git, GitHub y RStudio. El curso no se limita al uso del paquete
\texttt{BioStatR}, ya que se complementa con el uso de las funciones
clásicas de R base para todos los ejemplos presentados. Además, se
invita al estudiante a explorar y utilizar JAMOVI, una plataforma de
estadística gratuita y de código abierto diseñada para ser fácil de
usar. Al finalizar el curso, el estudiante será capaz de:

\begin{enumerate}
\def\labelenumi{\arabic{enumi}.}
\tightlist
\item
  \textbf{Describir y visualizar} conjuntos de datos mediante técnicas
  de Análisis Exploratorio de Datos.
\item
  \textbf{Comprender los fundamentos} de la teoría de la probabilidad
  como base del razonamiento estadístico.
\item
  \textbf{Manejar las distribuciones} de probabilidad más importantes,
  tanto discretas como continuas.
\item
  \textbf{Aplicar métodos de inferencia estadística}: estimación
  puntual, estimación por intervalos y contraste de hipótesis.
\item
  \textbf{Construir e interpretar modelos de regresión lineal} simple y
  múltiple.
\end{enumerate}

\section*{Estructura del Libro}\label{estructura-del-libro}
\addcontentsline{toc}{section}{Estructura del Libro}

\markright{Estructura del Libro}

El material está organizado en cuatro partes, correspondientes a las
diez semanas lectivas del curso:

\textbf{Parte I --- Análisis Exploratorio de Datos} (Semanas 1--2)
Introducción a los tipos de variables, estadísticos descriptivos
univariantes y bivariantes, medidas de dispersión, correlación y
representaciones gráficas básicas.

\textbf{Parte II --- Probabilidad y Variables Aleatorias} (Semanas 3--5)
Fundamentos de la teoría de la probabilidad: espacio muestral, axiomas
de Kolmogórov, probabilidad condicional y teorema de Bayes. Variables
aleatorias discretas y continuas, distribuciones de probabilidad más
importantes y el teorema del límite central.

\textbf{Parte III --- Inferencia Estadística} (Semanas 7--9) Conceptos
de muestreo y estadísticos muestrales. Estimación puntual: propiedades
de los estimadores, método de momentos y máxima verosimilitud.
Estimación por intervalos de confianza. Contrastes de hipótesis: errores
de tipo I y II, potencia del contraste, pruebas t, z y F.

\textbf{Parte IV --- Regresión Lineal} (Semanas 10--11) Regresión lineal
simple: mínimos cuadrados ordinarios, interpretación de coeficientes,
\(R^2\) y análisis de residuos. Regresión lineal múltiple: formulación
matricial, selección de variables y ANOVA de regresión.

\textbf{Parte V --- Métodos Avanzados en Bioestadística} Para
profundizar en técnicas estadísticas avanzadas aplicadas a la medicina,
se recomienda consultar el recurso complementario:
\href{https://migariane.github.io/Bioestadistica}{\textbf{Bioestadística
Avanzada}}.

Este material avanzado cubre temas críticos para la investigación médica
moderna, incluyendo: - \textbf{Pruebas Diagnósticas:} Evaluación de
precisión, sensibilidad, especificidad y curvas ROC. - \textbf{Regresión
Logística:} Modelado de variables dependientes binarias y cálculo de
Odds Ratios. - \textbf{Análisis de Supervivencia:} Estimadores de
Kaplan-Meier y modelos de riesgos proporcionales de Cox para datos de
tiempo a evento. - \textbf{Inferencia Causal y Métodos Robustos:}
Técnicas avanzadas para el control de la confusión y el análisis de
datos longitudinale (en preparación).

\section*{Sobre el Autor}\label{sobre-el-autor}
\addcontentsline{toc}{section}{Sobre el Autor}

\markright{Sobre el Autor}

\textbf{Miguel Ángel Luque Fernández} es Catedrático del
\href{https://estadistica.ugr.es}{Departamento de Estadística e
Investigación Operativa} de la Universidad de Granada. Sus líneas de
investigación se centran en la epidemiología cuantitativa, la
bioestadística y la inferencia causal.

\begin{itemize}
\tightlist
\item
  Correo: {\href{mailto:mluquefe@go.ugr.es}{mluquefe@ugr.es}}
\item
  Web: \href{https://migariane.github.io}{migariane.github.io}
\item
  GitHub: \href{https://github.com/migariane}{github.com/migariane}
\end{itemize}

\section*{Software y
Reproducibilidad}\label{software-y-reproducibilidad}
\addcontentsline{toc}{section}{Software y Reproducibilidad}

\markright{Software y Reproducibilidad}

Todos los ejemplos de código en este libro están escritos en
\href{https://www.r-project.org/}{R}. Para reproducir los análisis se
recomienda instalar R (versión ≥ 4.3) y los siguientes paquetes:

\begin{Shaded}
\begin{Highlighting}[]
\FunctionTok{install.packages}\NormalTok{(}\FunctionTok{c}\NormalTok{(}\StringTok{"ggplot2"}\NormalTok{, }\StringTok{"dplyr"}\NormalTok{, }\StringTok{"tidyr"}\NormalTok{, }\StringTok{"knitr"}\NormalTok{, }\StringTok{"kableExtra"}\NormalTok{, }\StringTok{"MASS"}\NormalTok{))}
\end{Highlighting}
\end{Shaded}

Además, para facilitar el aprendizaje y la aplicación práctica, contamos
con:

\begin{itemize}
\item
  \textbf{BioStatR:} Paquete de R desarrollado en la Unidad de
  Bioestadística de la Facultad de Medicina de la Universidad de Granada
  por \textbf{Pedro Femia Marzo} y \textbf{Miguel Ángel Luque
  Fernández}. Proporciona funciones de alto nivel que integran
  descripción, inferencia, modelización (lineal y logística) y
  visualización en una sola llamada, diseñadas específicamente para la
  docencia de bioestadística en ciencias de la salud. El paquete incluye
  también el conjunto de datos \texttt{osteo} (94 pacientes diabéticos,
  27 variables), utilizado como hilo conductor de los ejemplos del
  libro.

  \begin{itemize}
  \tightlist
  \item
    \textbf{Sitio web del paquete:}
    \url{https://migariane.github.io/BioStatR/}
  \item
    \textbf{Repositorio GitHub:}
    \url{https://github.com/migariane/BioStatR}
  \item
    \textbf{Documentación del Apéndice B} del libro: guía de referencia
    completa de las 22 funciones del paquete.
  \end{itemize}
\end{itemize}

Para instalarlo basta con ejecutar (requiere el paquete
\texttt{remotes}):

\begin{Shaded}
\begin{Highlighting}[]
\CommentTok{\# Instala el paquete remotes si no lo tienes}
\FunctionTok{install.packages}\NormalTok{(}\StringTok{"remotes"}\NormalTok{)}

\CommentTok{\# Instala BioStatR directamente desde GitHub (compila para tu sistema)}
\NormalTok{remotes}\SpecialCharTok{::}\FunctionTok{install\_github}\NormalTok{(}\StringTok{"migariane/BioStatR"}\NormalTok{)}

\CommentTok{\# Cargar el paquete y el conjunto de datos osteo}
\FunctionTok{library}\NormalTok{(BioStatR)}
\FunctionTok{data}\NormalTok{(osteo)}
\end{Highlighting}
\end{Shaded}

Para obtener la referencia bibliográfica del paquete:

\begin{Shaded}
\begin{Highlighting}[]
\FunctionTok{citation}\NormalTok{(}\StringTok{"BioStatR"}\NormalTok{)}
\end{Highlighting}
\end{Shaded}

\begin{itemize}
\tightlist
\item
  \textbf{\href{https://www.jamovi.org}{jamovi}:} Una plataforma de
  estadística gratuita y de código abierto diseñada para ser fácil de
  usar. \href{https://www.jamovi.org}{jamovi} proporciona una interfaz
  intuitiva basada en el motor estadístico de R, ideal para quienes
  desean realizar análisis complejos sin necesidad de escribir código
  directamente, facilitando la reproducibilidad de los resultados.
\end{itemize}

El código fuente completo de este libro está disponible en:
\url{https://github.com/migariane/CursoMatematicaEstadistica}

\section*{Nota sobre el Material}\label{nota-sobre-el-material}
\addcontentsline{toc}{section}{Nota sobre el Material}

\markright{Nota sobre el Material}

Este material docente ha sido elaborado por el autor para uso docente en
la Universidad de Granada. El material tiene carácter introductorio y no
pretende ser un tratado exhaustivo de ninguno de los temas cubiertos.
Para cada tema se proporcionan referencias para una lectura más
profunda.

\begin{quote}
\textbf{Aviso:} Estos apuntes se encuentran en continua revisión. Si
detecta algún error o imprecisión, por favor comuníquelo al autor a
través del repositorio de GitHub o por correo electrónico.
\end{quote}

\begin{center}\rule{0.5\linewidth}{0.5pt}\end{center}

\emph{Versión correspondiente al curso académico 2025--2026.}

\part{Parte I: Análisis Exploratorio de Datos}

\chapter{Semana 1 --- Análisis Exploratorio de Datos (Parte
I)}\label{sec-semana01}

\section{Introducción}\label{introducciuxf3n}

El análisis exploratorio de datos (AED) es el primer paso en cualquier
análisis estadístico. A través de visualizaciones, tablas y estadísticos
descriptivos, aprendemos a entender las características principales de
un conjunto de datos antes de realizar inferencia estadística más
compleja.

\section{Conceptos Fundamentales}\label{conceptos-fundamentales}

\subsection{Unidades Estadísticas, Variables y
Valores}\label{unidades-estaduxedsticas-variables-y-valores}

En estadística trabajamos con conceptos fundamentales que debemos
dominar desde el principio.

\begin{tcolorbox}[enhanced jigsaw, breakable, titlerule=0mm, rightrule=.15mm, leftrule=.75mm, left=2mm, opacitybacktitle=0.6, colback=white, bottomtitle=1mm, colbacktitle=quarto-callout-note-color!10!white, title=\textcolor{quarto-callout-note-color}{\faInfo}\hspace{0.5em}{Definición: Conceptos Fundamentales}, toptitle=1mm, coltitle=black, opacityback=0, arc=.35mm, bottomrule=.15mm, toprule=.15mm, colframe=quarto-callout-note-color-frame]

\begin{itemize}
\item
  \textbf{Unidad estadística} (o individuo): cada elemento sobre el que
  se realiza la observación. Puede ser una persona, animal, planta,
  empresa, etc.
\item
  \textbf{Variable}: una característica de una unidad estadística que
  puede variar de una unidad a otra. Se denota con letras mayúsculas
  (\(X\), \(Y\), etc.)
\item
  \textbf{Valor}: el resultado específico de observar una variable en
  una unidad estadística. Se denota con letras minúsculas (\(x_i\),
  \(y_i\), etc.)
\item
  \textbf{Población}: el conjunto completo de todas las unidades
  estadísticas de interés en un estudio.
\item
  \textbf{Muestra}: un subconjunto de la población observado en la
  práctica. La muestra se utiliza para extraer conclusiones sobre la
  población.
\end{itemize}

\end{tcolorbox}

\begin{tcolorbox}[enhanced jigsaw, breakable, titlerule=0mm, rightrule=.15mm, leftrule=.75mm, left=2mm, opacitybacktitle=0.6, colback=white, bottomtitle=1mm, colbacktitle=quarto-callout-tip-color!10!white, title=\textcolor{quarto-callout-tip-color}{\faLightbulb}\hspace{0.5em}{Ejemplo 1.1: Presión Arterial en Pacientes Diabéticos}, toptitle=1mm, coltitle=black, opacityback=0, arc=.35mm, bottomrule=.15mm, toprule=.15mm, colframe=quarto-callout-tip-color-frame]

Supongamos que deseamos estudiar la presión arterial sistólica (PAS) en
pacientes diabéticos de Granada.

\begin{itemize}
\tightlist
\item
  \textbf{Unidad estadística}: cada paciente diabético individual
\item
  \textbf{Variable}: presión arterial sistólica (\(X\))
\item
  \textbf{Valores}: mediciones específicas en mmHg (ej. \(x_1 = 125\),
  \(x_2 = 130\), \(x_3 = 145\), \ldots)
\item
  \textbf{Población}: todos los pacientes diabéticos de Granada
\item
  \textbf{Muestra}: los 200 pacientes diabéticos seleccionados para el
  estudio
\end{itemize}

\end{tcolorbox}

\subsection{Observaciones Atípicas
(Outliers)}\label{observaciones-atuxedpicas-outliers}

\begin{tcolorbox}[enhanced jigsaw, breakable, titlerule=0mm, rightrule=.15mm, leftrule=.75mm, left=2mm, opacitybacktitle=0.6, colback=white, bottomtitle=1mm, colbacktitle=quarto-callout-note-color!10!white, title=\textcolor{quarto-callout-note-color}{\faInfo}\hspace{0.5em}{Definición: Outlier}, toptitle=1mm, coltitle=black, opacityback=0, arc=.35mm, bottomrule=.15mm, toprule=.15mm, colframe=quarto-callout-note-color-frame]

Un \textbf{outlier} (observación atípica) es un valor que:

\begin{itemize}
\tightlist
\item
  Parece desviarse marcadamente de los otros miembros de la muestra
\item
  Aparece o está ausente de forma inesperada
\item
  Puede reflejar una anomalía en la característica medida o un error en
  la medición/registro
\end{itemize}

\textbf{Ejemplo}: En el conjunto de presiones arteriales sistólicas
\(\{125, 128, 130, 131, 198\}\), el valor \(198\) podría considerarse un
outlier, ya que es sustancialmente mayor que los otros valores.

\end{tcolorbox}

\begin{center}\rule{0.5\linewidth}{0.5pt}\end{center}

\section{Escalas de Medición}\label{escalas-de-mediciuxf3n}

El tipo de escala utilizada para medir una variable determina qué
técnicas analíticas son apropiadas y qué conclusiones se pueden extraer.

\begin{tcolorbox}[enhanced jigsaw, breakable, titlerule=0mm, rightrule=.15mm, leftrule=.75mm, left=2mm, opacitybacktitle=0.6, colback=white, bottomtitle=1mm, colbacktitle=quarto-callout-note-color!10!white, title=\textcolor{quarto-callout-note-color}{\faInfo}\hspace{0.5em}{Definición: Escalas de Medición}, toptitle=1mm, coltitle=black, opacityback=0, arc=.35mm, bottomrule=.15mm, toprule=.15mm, colframe=quarto-callout-note-color-frame]

Las variables se clasifican en cuatro tipos según la escala en que se
miden:

\begin{enumerate}
\def\labelenumi{\arabic{enumi}.}
\tightlist
\item
  \textbf{Variables nominales (categóricas sin orden)}
\item
  \textbf{Variables ordinales (categóricas con orden)}
\item
  \textbf{Variables discretas (numéricas, valores aislados)}
\item
  \textbf{Variables continuas (numéricas, infinitos valores posibles)}
\end{enumerate}

\end{tcolorbox}

\subsection{Variables Categóricas}\label{variables-categuxf3ricas}

\begin{tcolorbox}[enhanced jigsaw, breakable, titlerule=0mm, rightrule=.15mm, leftrule=.75mm, left=2mm, opacitybacktitle=0.6, colback=white, bottomtitle=1mm, colbacktitle=quarto-callout-note-color!10!white, title=\textcolor{quarto-callout-note-color}{\faInfo}\hspace{0.5em}{Definición: Variables Categóricas (Nominales)}, toptitle=1mm, coltitle=black, opacityback=0, arc=.35mm, bottomrule=.15mm, toprule=.15mm, colframe=quarto-callout-note-color-frame]

Una \textbf{variable categórica} es aquella que puede tomar uno de un
conjunto limitado de valores posibles, asignando cada unidad a una
categoría particular basada en una propiedad cualitativa.

\textbf{Ejemplos}: - Grupo sanguíneo: A, B, AB, O - Estado civil:
soltero, casado, divorciado, viudo - Género: masculino, femenino -
Región: Andalucía, Cataluña, Madrid, etc.

Los valores de una variable categórica se llaman \textbf{niveles}.

Las variables categóricas pueden ser: - \textbf{Binarias (dicotómicas)}:
solo dos niveles (ej. éxito/fracaso, enfermo/sano) -
\textbf{Politómicas}: más de dos niveles (ej. grupo sanguíneo)

\end{tcolorbox}

\subsection{Variables Ordinales}\label{variables-ordinales}

\begin{tcolorbox}[enhanced jigsaw, breakable, titlerule=0mm, rightrule=.15mm, leftrule=.75mm, left=2mm, opacitybacktitle=0.6, colback=white, bottomtitle=1mm, colbacktitle=quarto-callout-note-color!10!white, title=\textcolor{quarto-callout-note-color}{\faInfo}\hspace{0.5em}{Definición: Variables Ordinales}, toptitle=1mm, coltitle=black, opacityback=0, arc=.35mm, bottomrule=.15mm, toprule=.15mm, colframe=quarto-callout-note-color-frame]

Una \textbf{variable ordinal} es una variable categórica en la que los
niveles tienen un orden natural, pero las diferencias entre ellos no son
cuantificables.

Relaciones válidas: ``mayor que'' y ``menor que'' (o ``mejor que'' y
``peor que'').

\textbf{Ejemplos}: - Estadio de un tumor: I, II, III, IV - Grado de
dependencia: leve, moderada, severa, total - Escala de dolor: nada,
poco, moderado, intenso, insoportable

\end{tcolorbox}

\subsection{Variables Numéricas}\label{variables-numuxe9ricas}

\begin{tcolorbox}[enhanced jigsaw, breakable, titlerule=0mm, rightrule=.15mm, leftrule=.75mm, left=2mm, opacitybacktitle=0.6, colback=white, bottomtitle=1mm, colbacktitle=quarto-callout-note-color!10!white, title=\textcolor{quarto-callout-note-color}{\faInfo}\hspace{0.5em}{Definición: Variables Discretas}, toptitle=1mm, coltitle=black, opacityback=0, arc=.35mm, bottomrule=.15mm, toprule=.15mm, colframe=quarto-callout-note-color-frame]

Una \textbf{variable discreta} es una variable numérica que solo puede
tomar un número finito o contablemente infinito de valores. Cada valor
tiene vecinos claramente identificables.

\textbf{Ejemplos}: - Número de nacimientos en un mes - Número de
episodios de hipoglucemia - Número de dientes extraídos - Cantidad de
admisiones previas en el hospital

\end{tcolorbox}

\begin{tcolorbox}[enhanced jigsaw, breakable, titlerule=0mm, rightrule=.15mm, leftrule=.75mm, left=2mm, opacitybacktitle=0.6, colback=white, bottomtitle=1mm, colbacktitle=quarto-callout-note-color!10!white, title=\textcolor{quarto-callout-note-color}{\faInfo}\hspace{0.5em}{Definición: Variables Continuas}, toptitle=1mm, coltitle=black, opacityback=0, arc=.35mm, bottomrule=.15mm, toprule=.15mm, colframe=quarto-callout-note-color-frame]

Una \textbf{variable continua} es una variable numérica que puede tomar
infinitos valores no contables. Entre dos valores cualesquiera siempre
existen otros valores posibles.

\textbf{Ejemplos}: - Presión arterial sistólica (en mmHg) - Peso
corporal (en kg) - Altura (en cm) - Temperatura (en °C)

\textbf{En la práctica}: variables con muchas unidades contables (como
el recuento de glóbulos rojos o blancos) se tratan como continuas y a
veces se llaman ``cuasi-continuas''.

\end{tcolorbox}

\begin{center}\rule{0.5\linewidth}{0.5pt}\end{center}

\section{Agrupación de Datos en Clases
(Binning)}\label{agrupaciuxf3n-de-datos-en-clases-binning}

Cuando disponemos de muchas observaciones de una variable continua,
frecuentemente es útil agruparlas en intervalos de clase para mejorar la
claridad.

\begin{tcolorbox}[enhanced jigsaw, breakable, titlerule=0mm, rightrule=.15mm, leftrule=.75mm, left=2mm, opacitybacktitle=0.6, colback=white, bottomtitle=1mm, colbacktitle=quarto-callout-note-color!10!white, title=\textcolor{quarto-callout-note-color}{\faInfo}\hspace{0.5em}{Definición: Binning (Agrupación en Clases)}, toptitle=1mm, coltitle=black, opacityback=0, arc=.35mm, bottomrule=.15mm, toprule=.15mm, colframe=quarto-callout-note-color-frame]

El \textbf{binning} es la partición de los valores de una variable
continua en varias clases (típicamente intervalos). Esto mejora la
claridad cuando se tiene gran cantidad de datos.

Para cada intervalo o clase \(j\) definimos:

\begin{itemize}
\tightlist
\item
  \textbf{Límite inferior de la clase}: \(x_j^l\)
\item
  \textbf{Límite superior de la clase}: \(x_j^u\)
\item
  \textbf{Amplitud de la clase}: \(\Delta x_j = x_j^u - x_j^l\)
\end{itemize}

\textbf{Propiedades importantes}: - Los intervalos son no superpuestos
(disjuntos) y adyacentes - \(x_j^u = x_{j+1}^l\) (el límite superior de
una clase es el límite inferior de la siguiente) - Convencionalmente:
\(x_j^l < X \leq x_j^u\) (incluimos el límite superior, excluimos el
inferior)

\end{tcolorbox}

\begin{tcolorbox}[enhanced jigsaw, breakable, titlerule=0mm, rightrule=.15mm, leftrule=.75mm, left=2mm, opacitybacktitle=0.6, colback=white, bottomtitle=1mm, colbacktitle=quarto-callout-tip-color!10!white, title=\textcolor{quarto-callout-tip-color}{\faLightbulb}\hspace{0.5em}{Ejemplo 1.2: Niveles de Colesterol Total}, toptitle=1mm, coltitle=black, opacityback=0, arc=.35mm, bottomrule=.15mm, toprule=.15mm, colframe=quarto-callout-tip-color-frame]

En un estudio sobre salud cardiovascular tenemos: - \textbf{Unidad
estadística}: paciente - \textbf{Variable}: nivel de colesterol total
(en mg/dL)

\begin{longtable}[]{@{}
  >{\raggedright\arraybackslash}p{(\linewidth - 4\tabcolsep) * \real{0.3333}}
  >{\raggedleft\arraybackslash}p{(\linewidth - 4\tabcolsep) * \real{0.3333}}
  >{\raggedleft\arraybackslash}p{(\linewidth - 4\tabcolsep) * \real{0.3333}}@{}}
\toprule\noalign{}
\begin{minipage}[b]{\linewidth}\raggedright
Intervalo de Colesterol (mg/dL)
\end{minipage} & \begin{minipage}[b]{\linewidth}\raggedleft
Límites
\end{minipage} & \begin{minipage}[b]{\linewidth}\raggedleft
Amplitud
\end{minipage} \\
\midrule\noalign{}
\endhead
\bottomrule\noalign{}
\endlastfoot
150 -- 200 (Deseable) & \(x_1^l = 150\), \(x_1^u = 200\) &
\(\Delta x_1 = 50\) \\
200 -- 240 (Límite alto) & \(x_2^l = 200\), \(x_2^u = 240\) &
\(\Delta x_2 = 40\) \\
240 -- 300 (Alto) & \(x_3^l = 240\), \(x_3^u = 300\) &
\(\Delta x_3 = 60\) \\
\end{longtable}

\end{tcolorbox}

\begin{center}\rule{0.5\linewidth}{0.5pt}\end{center}

\section{Tablas de Frecuencias y Distribución
Empírica}\label{tablas-de-frecuencias-y-distribuciuxf3n-empuxedrica}

Una forma eficaz de resumir un conjunto de datos es construir una tabla
de frecuencias.

\subsection{Frecuencias Absolutas y
Relativas}\label{frecuencias-absolutas-y-relativas}

\begin{tcolorbox}[enhanced jigsaw, breakable, titlerule=0mm, rightrule=.15mm, leftrule=.75mm, left=2mm, opacitybacktitle=0.6, colback=white, bottomtitle=1mm, colbacktitle=quarto-callout-note-color!10!white, title=\textcolor{quarto-callout-note-color}{\faInfo}\hspace{0.5em}{Definición: Frecuencia Absoluta}, toptitle=1mm, coltitle=black, opacityback=0, arc=.35mm, bottomrule=.15mm, toprule=.15mm, colframe=quarto-callout-note-color-frame]

La \textbf{frecuencia absoluta} de un valor \(x_j\) es el número de
unidades estadísticas con ese valor característico:

\[h(x_j) = h_j\]

\textbf{Propiedades}:
\[0 \leq h(x_j) \leq n, \quad \sum_{j=1}^{k} h(x_j) = n\]

donde \(n\) es el número total de observaciones y \(k\) es el número de
valores distintos.

\end{tcolorbox}

\begin{tcolorbox}[enhanced jigsaw, breakable, titlerule=0mm, rightrule=.15mm, leftrule=.75mm, left=2mm, opacitybacktitle=0.6, colback=white, bottomtitle=1mm, colbacktitle=quarto-callout-note-color!10!white, title=\textcolor{quarto-callout-note-color}{\faInfo}\hspace{0.5em}{Definición: Frecuencia Relativa}, toptitle=1mm, coltitle=black, opacityback=0, arc=.35mm, bottomrule=.15mm, toprule=.15mm, colframe=quarto-callout-note-color-frame]

La \textbf{frecuencia relativa} de un valor \(x_j\) es la proporción de
unidades estadísticas con ese valor:

\[f(x_j) = \frac{h(x_j)}{n}\]

\textbf{Propiedades}:
\[0 \leq f(x_j) \leq 1, \quad \sum_{j=1}^{k} f(x_j) = 1\]

La frecuencia relativa es útil para comparar distribuciones de
diferentes tamaños.

\end{tcolorbox}

\subsection{Tabla de Frecuencias}\label{tabla-de-frecuencias}

\begin{tcolorbox}[enhanced jigsaw, breakable, titlerule=0mm, rightrule=.15mm, leftrule=.75mm, left=2mm, opacitybacktitle=0.6, colback=white, bottomtitle=1mm, colbacktitle=quarto-callout-note-color!10!white, title=\textcolor{quarto-callout-note-color}{\faInfo}\hspace{0.5em}{Definición: Distribución de Frecuencias Empírica}, toptitle=1mm, coltitle=black, opacityback=0, arc=.35mm, bottomrule=.15mm, toprule=.15mm, colframe=quarto-callout-note-color-frame]

Una \textbf{distribución de frecuencias empírica} es una tabla que lista
los valores (o intervalos de valores) de una variable junto con sus
frecuencias (absolutas o relativas).

\textbf{Estructura de una tabla de frecuencias}:

\begin{longtable}[]{@{}
  >{\raggedright\arraybackslash}p{(\linewidth - 4\tabcolsep) * \real{0.3333}}
  >{\raggedleft\arraybackslash}p{(\linewidth - 4\tabcolsep) * \real{0.3333}}
  >{\raggedleft\arraybackslash}p{(\linewidth - 4\tabcolsep) * \real{0.3333}}@{}}
\toprule\noalign{}
\begin{minipage}[b]{\linewidth}\raggedright
Valor (\(x_j\))
\end{minipage} & \begin{minipage}[b]{\linewidth}\raggedleft
Frecuencia Absoluta (\(h_j\))
\end{minipage} & \begin{minipage}[b]{\linewidth}\raggedleft
Frecuencia Relativa (\(f_j\))
\end{minipage} \\
\midrule\noalign{}
\endhead
\bottomrule\noalign{}
\endlastfoot
\(x_1\) & \(h_1\) & \(f_1\) \\
\(x_2\) & \(h_2\) & \(f_2\) \\
\ldots{} & \ldots{} & \ldots{} \\
\(x_k\) & \(h_k\) & \(f_k\) \\
\textbf{Suma} & \textbf{n} & \textbf{1} \\
\end{longtable}

\end{tcolorbox}

\subsection{Frecuencias Acumuladas}\label{frecuencias-acumuladas}

\begin{tcolorbox}[enhanced jigsaw, breakable, titlerule=0mm, rightrule=.15mm, leftrule=.75mm, left=2mm, opacitybacktitle=0.6, colback=white, bottomtitle=1mm, colbacktitle=quarto-callout-note-color!10!white, title=\textcolor{quarto-callout-note-color}{\faInfo}\hspace{0.5em}{Definición: Frecuencia Acumulada}, toptitle=1mm, coltitle=black, opacityback=0, arc=.35mm, bottomrule=.15mm, toprule=.15mm, colframe=quarto-callout-note-color-frame]

La \textbf{frecuencia acumulada} es la suma de las frecuencias de todos
los valores hasta un punto determinado.

\textbf{Frecuencia acumulada absoluta}:
\[H(x_j) = \sum_{i=1}^{j} h(x_i), \quad j = 1, \ldots, k\]

\textbf{Frecuencia acumulada relativa}:
\[F(x_j) = \frac{H(x_j)}{n} = \sum_{i=1}^{j} f(x_i), \quad j = 1, \ldots, k\]

\textbf{Propiedades}: - \(H(x_1) = h(x_1)\) y \(F(x_1) = f(x_1)\) -
\(H(x_k) = n\) y \(F(x_k) = 1\) - Las frecuencias acumuladas son
secuencias no decrecientes

\end{tcolorbox}

\begin{tcolorbox}[enhanced jigsaw, breakable, titlerule=0mm, rightrule=.15mm, leftrule=.75mm, left=2mm, opacitybacktitle=0.6, colback=white, bottomtitle=1mm, colbacktitle=quarto-callout-tip-color!10!white, title=\textcolor{quarto-callout-tip-color}{\faLightbulb}\hspace{0.5em}{Ejemplo 1.3: Tabla de Frecuencias}, toptitle=1mm, coltitle=black, opacityback=0, arc=.35mm, bottomrule=.15mm, toprule=.15mm, colframe=quarto-callout-tip-color-frame]

Consideremos los datos de 10 pacientes diabéticos. Sus presiones
arteriales sistólicas (en mmHg) son: 120, 125, 118, 130, 122, 128, 115,
135, 124, 129

\begin{longtable}[]{@{}lrrrr@{}}
\toprule\noalign{}
PAS (mmHg) & \(h_j\) & \(f_j\) & \(H_j\) & \(F_j\) \\
\midrule\noalign{}
\endhead
\bottomrule\noalign{}
\endlastfoot
115 & 1 & 0.10 & 1 & 0.10 \\
118 & 1 & 0.10 & 2 & 0.20 \\
120 & 1 & 0.10 & 3 & 0.30 \\
122 & 1 & 0.10 & 4 & 0.40 \\
124 & 1 & 0.10 & 5 & 0.50 \\
125 & 1 & 0.10 & 6 & 0.60 \\
128 & 1 & 0.10 & 7 & 0.70 \\
129 & 1 & 0.10 & 8 & 0.80 \\
130 & 1 & 0.10 & 9 & 0.90 \\
135 & 1 & 0.10 & 10 & 1.00 \\
\end{longtable}

Podemos observar que todas las presiones son diferentes en este ejemplo,
por lo que cada una aparece una sola vez.

\begin{codelisting}[H]

\caption{\label{lst-semana01-1}Código R para tabla de frecuencias de
PAS}

\centering{

\begin{Shaded}
\begin{Highlighting}[]
\CommentTok{\# Datos de presión arterial sistólica (PAS) en pacientes diabéticos}
\NormalTok{pas\_diabeticos }\OtherTok{\textless{}{-}} \FunctionTok{c}\NormalTok{(}\DecValTok{120}\NormalTok{, }\DecValTok{125}\NormalTok{, }\DecValTok{118}\NormalTok{, }\DecValTok{130}\NormalTok{, }\DecValTok{122}\NormalTok{, }\DecValTok{128}\NormalTok{, }\DecValTok{115}\NormalTok{, }\DecValTok{135}\NormalTok{, }\DecValTok{124}\NormalTok{, }\DecValTok{129}\NormalTok{)}

\CommentTok{\# Crear tabla de frecuencias}
\NormalTok{tabla\_freq }\OtherTok{\textless{}{-}} \FunctionTok{table}\NormalTok{(pas\_diabeticos)}
\NormalTok{freq\_rel }\OtherTok{\textless{}{-}}\NormalTok{ tabla\_freq }\SpecialCharTok{/} \FunctionTok{length}\NormalTok{(pas\_diabeticos)}
\NormalTok{freq\_acum }\OtherTok{\textless{}{-}} \FunctionTok{cumsum}\NormalTok{(tabla\_freq)}
\NormalTok{freq\_acum\_rel }\OtherTok{\textless{}{-}} \FunctionTok{cumsum}\NormalTok{(freq\_rel)}

\CommentTok{\# Crear data frame}
\NormalTok{result }\OtherTok{\textless{}{-}} \FunctionTok{data.frame}\NormalTok{(}
  \AttributeTok{PAS =} \FunctionTok{as.numeric}\NormalTok{(}\FunctionTok{names}\NormalTok{(tabla\_freq)),}
  \AttributeTok{h =} \FunctionTok{as.numeric}\NormalTok{(tabla\_freq),}
  \AttributeTok{f =} \FunctionTok{round}\NormalTok{(freq\_rel, }\DecValTok{2}\NormalTok{),}
  \AttributeTok{H =}\NormalTok{ freq\_acum,}
  \AttributeTok{F =} \FunctionTok{round}\NormalTok{(freq\_acum\_rel, }\DecValTok{2}\NormalTok{)}
\NormalTok{)}
\FunctionTok{print}\NormalTok{(result)}
\end{Highlighting}
\end{Shaded}

}

\end{codelisting}%

\begin{verbatim}
    PAS h f.pas_diabeticos f.Freq  H   F
115 115 1              115    0.1  1 0.1
118 118 1              118    0.1  2 0.2
120 120 1              120    0.1  3 0.3
122 122 1              122    0.1  4 0.4
124 124 1              124    0.1  5 0.5
125 125 1              125    0.1  6 0.6
128 128 1              128    0.1  7 0.7
129 129 1              129    0.1  8 0.8
130 130 1              130    0.1  9 0.9
135 135 1              135    0.1 10 1.0
\end{verbatim}

\end{tcolorbox}

\begin{tcolorbox}[enhanced jigsaw, breakable, titlerule=0mm, rightrule=.15mm, leftrule=.75mm, left=2mm, opacitybacktitle=0.6, colback=white, bottomtitle=1mm, colbacktitle=quarto-callout-tip-color!10!white, title=\textcolor{quarto-callout-tip-color}{\faLightbulb}\hspace{0.5em}{Ejemplo 1.4: BioStatR: freq()}, toptitle=1mm, coltitle=black, opacityback=0, arc=.35mm, bottomrule=.15mm, toprule=.15mm, colframe=quarto-callout-tip-color-frame]

La función \texttt{freq()} genera automáticamente la tabla de
frecuencias con proporciones y acumuladas (documentación completa en
Sección~\ref{sec-bio-descriptivo}):

\begin{codelisting}[H]

\caption{\label{lst-semana01-2}Código R}

\centering{

\begin{Shaded}
\begin{Highlighting}[]
\FunctionTok{library}\NormalTok{(BioStatR)}

\CommentTok{\# Tabla de frecuencias de la edad — pacientes diabéticos, UGR}
\FunctionTok{freq}\NormalTok{(osteo}\SpecialCharTok{$}\NormalTok{edad, }\AttributeTok{grf =} \ConstantTok{FALSE}\NormalTok{)}
\end{Highlighting}
\end{Shaded}

}

\end{codelisting}%

\begin{verbatim}

Distribución de frecuencias
--------------------------------
Variable:  osteo$edad 
n = 94 
 
        x Freq  Prop Prop.Acum
1 (18,23]   27 0.287     0.287
2 (23,28]   22 0.234     0.521
3 (28,33]   14 0.149     0.670
4 (33,38]   15 0.160     0.830
5 (38,43]    4 0.043     0.873
6 (43,48]    5 0.053     0.926
7 (48,53]    3 0.032     0.958
8 (53,58]    3 0.032     0.990
\end{verbatim}

El 52.1\% de los pacientes tiene menos de 28 años. Con
\texttt{grf\ =\ TRUE} (por defecto) también genera el histograma. Con
\texttt{cuts\ =\ 5} agrupa en 5 clases de igual amplitud.

\end{tcolorbox}

\subsection{Función de Distribución Empírica
(ECDF)}\label{funciuxf3n-de-distribuciuxf3n-empuxedrica-ecdf}

\begin{tcolorbox}[enhanced jigsaw, breakable, titlerule=0mm, rightrule=.15mm, leftrule=.75mm, left=2mm, opacitybacktitle=0.6, colback=white, bottomtitle=1mm, colbacktitle=quarto-callout-note-color!10!white, title=\textcolor{quarto-callout-note-color}{\faInfo}\hspace{0.5em}{Definición: Función de Distribución Empírica (ECDF)}, toptitle=1mm, coltitle=black, opacityback=0, arc=.35mm, bottomrule=.15mm, toprule=.15mm, colframe=quarto-callout-note-color-frame]

La \textbf{Función de Distribución Empírica} (ECDF) asigna a cada valor
\(x\) la proporción de observaciones en la muestra que son menores o
iguales a \(x\).

Para datos no agrupados, la ECDF es una \textbf{función escalonada} que
salta en cada valor observado:

\[F(x) = \begin{cases}
0 & \text{si } x < x_{(1)} \\
\frac{i}{n} & \text{si } x_{(i)} \leq x < x_{(i+1)} \\
1 & \text{si } x \geq x_{(n)}
\end{cases}\]

donde \(x_{(1)} \leq x_{(2)} \leq \cdots \leq x_{(n)}\) son los valores
ordenados de la muestra.

\end{tcolorbox}

\begin{tcolorbox}[enhanced jigsaw, breakable, titlerule=0mm, rightrule=.15mm, leftrule=.75mm, left=2mm, opacitybacktitle=0.6, colback=white, bottomtitle=1mm, colbacktitle=quarto-callout-tip-color!10!white, title=\textcolor{quarto-callout-tip-color}{\faLightbulb}\hspace{0.5em}{Ejemplo 1.5: Visualización: ECDF Escalonada (Datos de PAS)}, toptitle=1mm, coltitle=black, opacityback=0, arc=.35mm, bottomrule=.15mm, toprule=.15mm, colframe=quarto-callout-tip-color-frame]

Usando los datos de PAS del Ejemplo 1.3, observamos cómo la función
acumula probabilidad en saltos de \(1/n\) en cada valor observado:

\begin{figure}[H]

\centering{

\pandocbounded{\includegraphics[keepaspectratio]{chapters/semana01_files/figure-pdf/fig-semana01-ecdf-step-1.pdf}}

}

\caption{\label{fig-semana01-ecdf-step}ECDF Escalonada para datos de
PAS}

\end{figure}%

\end{tcolorbox}

\begin{tcolorbox}[enhanced jigsaw, breakable, titlerule=0mm, rightrule=.15mm, leftrule=.75mm, left=2mm, opacitybacktitle=0.6, colback=white, bottomtitle=1mm, colbacktitle=quarto-callout-warning-color!10!white, title=\textcolor{quarto-callout-warning-color}{\faExclamationTriangle}\hspace{0.5em}{Propiedad Importante: ECDF}, toptitle=1mm, coltitle=black, opacityback=0, arc=.35mm, bottomrule=.15mm, toprule=.15mm, colframe=quarto-callout-warning-color-frame]

De la ECDF podemos recuperar las frecuencias relativas mediante:

\[f(x_j) = F(x_j) - F(x_{j-1}) \quad \text{para } j = 1, \ldots, k\]

con la convención \(F(x_0) = 0\).

\end{tcolorbox}

\begin{tcolorbox}[enhanced jigsaw, breakable, titlerule=0mm, rightrule=.15mm, leftrule=.75mm, left=2mm, opacitybacktitle=0.6, colback=white, bottomtitle=1mm, colbacktitle=quarto-callout-tip-color!10!white, title=\textcolor{quarto-callout-tip-color}{\faLightbulb}\hspace{0.5em}{Ejemplo 1.6: Tiempo de Supervivencia de Pacientes}, toptitle=1mm, coltitle=black, opacityback=0, arc=.35mm, bottomrule=.15mm, toprule=.15mm, colframe=quarto-callout-tip-color-frame]

Se estudió el tiempo de supervivencia (en días) tras el diagnóstico de
una patología agresiva en una cohorte de 100 pacientes. Los datos se
agrupan en clases:

\begin{longtable}[]{@{}lrrrr@{}}
\toprule\noalign{}
Intervalo (días) & \(h_j\) & \(f_j\) & \(H_j\) & \(F_j\) \\
\midrule\noalign{}
\endhead
\bottomrule\noalign{}
\endlastfoot
0 -- 100 & 1 & 0.01 & 1 & 0.01 \\
100 -- 500 & 24 & 0.24 & 25 & 0.25 \\
500 -- 1000 & 45 & 0.45 & 70 & 0.70 \\
1000 -- 4000 & 30 & 0.30 & 100 & 1.00 \\
\textbf{Total} & \textbf{100} & \textbf{1.00} & & \\
\end{longtable}

\subsection{Interpolación Lineal para Datos
Agrupados}\label{interpolaciuxf3n-lineal-para-datos-agrupados}

Para datos agrupados, no conocemos los valores exactos dentro de cada
clase. Por ello, aproximamos la ECDF mediante una \textbf{interpolación
lineal} entre los límites de las clases. Esto genera una línea quebrada
llamada \textbf{polígono de frecuencias acumuladas} (u ojiva).

La fórmula general de interpolación para un valor \(x\) dentro de la
clase \(j = [x_j^l, x_j^u]\) es:

\[F(x) = F(x_j^l) + \frac{x - x_j^l}{x_j^u - x_j^l} \cdot f(x_j)\]

Donde: - \(x_j^l\) y \(x_j^u\) son los límites inferior y superior de la
clase. - \(F(x_j^l)\) es la frecuencia relativa acumulada hasta la clase
anterior. - \(f(x_j)\) es la frecuencia relativa de la clase actual.

\textbf{Nota}: Esta fórmula es la base para estimar \textbf{cuantiles}
en datos agrupados. Si despejamos \(x\) para un valor de \(F(x) = p\),
obtenemos el valor del cuantil correspondiente.

Aplicando esto a nuestros datos:

\[F(x) = \begin{cases}
0 & \text{si } x \leq 0 \\
\frac{x - 0}{100} \cdot 0.01 & \text{si } 0 < x \leq 100 \\
0.01 + \frac{x - 100}{400} \cdot 0.24 & \text{si } 100 < x \leq 500 \\
0.25 + \frac{x - 500}{500} \cdot 0.45 & \text{si } 500 < x \leq 1000 \\
0.70 + \frac{x - 1000}{3000} \cdot 0.30 & \text{si } 1000 < x \leq 4000 \\
1 & \text{si } x > 4000
\end{cases}\]

\begin{figure}[H]

\centering{

\pandocbounded{\includegraphics[keepaspectratio]{chapters/semana01_files/figure-pdf/fig-semana01-ecdf-interp-1.pdf}}

}

\caption{\label{fig-semana01-ecdf-interp}Polígono de frecuencias
acumuladas (ECDF interpolada)}

\end{figure}%

\end{tcolorbox}

\begin{center}\rule{0.5\linewidth}{0.5pt}\end{center}

\section{Histogramas y Gráficos de
Densidad}\label{histogramas-y-gruxe1ficos-de-densidad}

Para variables continuas agrupadas en clases, el histograma es la forma
estándar de visualización.

\subsection{Construcción de
Histogramas}\label{construcciuxf3n-de-histogramas}

\begin{tcolorbox}[enhanced jigsaw, breakable, titlerule=0mm, rightrule=.15mm, leftrule=.75mm, left=2mm, opacitybacktitle=0.6, colback=white, bottomtitle=1mm, colbacktitle=quarto-callout-note-color!10!white, title=\textcolor{quarto-callout-note-color}{\faInfo}\hspace{0.5em}{Definición: Histograma}, toptitle=1mm, coltitle=black, opacityback=0, arc=.35mm, bottomrule=.15mm, toprule=.15mm, colframe=quarto-callout-note-color-frame]

Un \textbf{histograma} es una representación gráfica de una distribución
de frecuencias para datos agrupados. Se construye de la siguiente
manera:

\begin{itemize}
\tightlist
\item
  \textbf{Eje X}: límites de las clases \(x_j^l, x_j^u\)
\item
  \textbf{Eje Y}: densidad de frecuencia (frecuencia dividida por la
  amplitud de la clase)
\end{itemize}

\textbf{Densidad de frecuencia absoluta}:
\[\hat{h}(x_j) = \frac{h(x_j)}{\Delta x_j} = \frac{h(x_j)}{x_j^u - x_j^l}\]

\textbf{Densidad de frecuencia relativa}:
\[\hat{f}(x_j) = \frac{f(x_j)}{\Delta x_j} = \frac{f(x_j)}{x_j^u - x_j^l}\]

\textbf{Propiedad fundamental}: El \textbf{área} de cada rectángulo
representa la frecuencia (relativa o absoluta), y el área total del
histograma es igual a \(n\) (si usamos densidades absolutas) o \(1\) (si
usamos densidades relativas).

\[\sum_{j=1}^{k} \hat{f}(x_j) \cdot \Delta x_j = \sum_{j=1}^{k} f(x_j) = 1\]

\end{tcolorbox}

\begin{tcolorbox}[enhanced jigsaw, breakable, titlerule=0mm, rightrule=.15mm, leftrule=.75mm, left=2mm, opacitybacktitle=0.6, colback=white, bottomtitle=1mm, colbacktitle=quarto-callout-tip-color!10!white, title=\textcolor{quarto-callout-tip-color}{\faLightbulb}\hspace{0.5em}{Ejemplo 1.7: Índice de Masa Corporal (IMC)}, toptitle=1mm, coltitle=black, opacityback=0, arc=.35mm, bottomrule=.15mm, toprule=.15mm, colframe=quarto-callout-tip-color-frame]

Se recopilaron los IMC de 274 pacientes. Los datos se clasifican en
intervalos según los criterios de la OMS:

\begin{longtable}[]{@{}lrrrr@{}}
\toprule\noalign{}
IMC (kg/m²) & \(h_j\) & \(f_j\) & \(\Delta x_j\) & \(\hat{f}(x_j)\) \\
\midrule\noalign{}
\endhead
\bottomrule\noalign{}
\endlastfoot
0 -- 18.5 (Bajo peso) & 15 & 0.055 & 18.5 & 0.0030 \\
18.5 -- 25 (Normal) & 93 & 0.339 & 6.5 & 0.0522 \\
25 -- 30 (Sobrepeso) & 92 & 0.336 & 5 & 0.0672 \\
30 -- 35 (Obesidad I) & 54 & 0.197 & 5 & 0.0394 \\
35 -- 50 (Obesidad II/III) & 20 & 0.073 & 15 & 0.0049 \\
\textbf{Total} & \textbf{274} & \textbf{1.000} & & \\
\end{longtable}

\textbf{Interpretación}: La densidad más alta está en la clase de
``Sobrepeso'', indicando que ese intervalo tiene la mayor concentración
de pacientes en esta muestra.

\begin{codelisting}[H]

\caption{\label{lst-semana01-4}Código R para histograma de IMC}

\centering{

\begin{Shaded}
\begin{Highlighting}[]
\CommentTok{\# Simular datos de IMC basados en la tabla}
\NormalTok{imc }\OtherTok{\textless{}{-}} \FunctionTok{c}\NormalTok{(}\FunctionTok{runif}\NormalTok{(}\DecValTok{15}\NormalTok{, }\DecValTok{15}\NormalTok{, }\FloatTok{18.5}\NormalTok{), }\FunctionTok{runif}\NormalTok{(}\DecValTok{93}\NormalTok{, }\FloatTok{18.5}\NormalTok{, }\DecValTok{25}\NormalTok{), }\FunctionTok{runif}\NormalTok{(}\DecValTok{92}\NormalTok{, }\DecValTok{25}\NormalTok{, }\DecValTok{30}\NormalTok{),}
         \FunctionTok{runif}\NormalTok{(}\DecValTok{54}\NormalTok{, }\DecValTok{30}\NormalTok{, }\DecValTok{35}\NormalTok{), }\FunctionTok{runif}\NormalTok{(}\DecValTok{20}\NormalTok{, }\DecValTok{35}\NormalTok{, }\DecValTok{50}\NormalTok{))}

\CommentTok{\# Crear histograma con densidad}
\FunctionTok{hist}\NormalTok{(imc, }\AttributeTok{breaks =} \FunctionTok{c}\NormalTok{(}\DecValTok{0}\NormalTok{, }\FloatTok{18.5}\NormalTok{, }\DecValTok{25}\NormalTok{, }\DecValTok{30}\NormalTok{, }\DecValTok{35}\NormalTok{, }\DecValTok{50}\NormalTok{),}
     \AttributeTok{freq =} \ConstantTok{FALSE}\NormalTok{,}
     \AttributeTok{main =} \StringTok{"Histograma: Índice de Masa Corporal (IMC)"}\NormalTok{,}
     \AttributeTok{xlab =} \StringTok{"IMC (kg/m²)"}\NormalTok{,}
     \AttributeTok{ylab =} \StringTok{"Densidad de Frecuencia"}\NormalTok{,}
     \AttributeTok{col =} \StringTok{"lightgreen"}\NormalTok{,}
     \AttributeTok{border =} \StringTok{"darkgreen"}\NormalTok{)}
\end{Highlighting}
\end{Shaded}

}

\end{codelisting}%

\pandocbounded{\includegraphics[keepaspectratio]{chapters/semana01_files/figure-pdf/unnamed-chunk-3-1.pdf}}

\end{tcolorbox}

:::

\subsection{Histograma vs Gráfico de
Densidad}\label{histograma-vs-gruxe1fico-de-densidad}

\begin{tcolorbox}[enhanced jigsaw, breakable, titlerule=0mm, rightrule=.15mm, leftrule=.75mm, left=2mm, opacitybacktitle=0.6, colback=white, bottomtitle=1mm, colbacktitle=quarto-callout-tip-color!10!white, title=\textcolor{quarto-callout-tip-color}{\faLightbulb}\hspace{0.5em}{Ejemplo 1.8: Visualización en R}, toptitle=1mm, coltitle=black, opacityback=0, arc=.35mm, bottomrule=.15mm, toprule=.15mm, colframe=quarto-callout-tip-color-frame]

\begin{codelisting}[H]

\caption{\label{lst-semana01-5}Comparación de gráficos}

\centering{

\begin{Shaded}
\begin{Highlighting}[]
\CommentTok{\# Simular datos de Presión Arterial Diastólica (en mmHg)}
\FunctionTok{set.seed}\NormalTok{(}\DecValTok{42}\NormalTok{)}
\NormalTok{pad }\OtherTok{\textless{}{-}} \FunctionTok{rnorm}\NormalTok{(}\DecValTok{200}\NormalTok{, }\AttributeTok{mean =} \DecValTok{80}\NormalTok{, }\AttributeTok{sd =} \DecValTok{10}\NormalTok{)}

\CommentTok{\# Configurar área de trazado para 2 gráficos}
\FunctionTok{par}\NormalTok{(}\AttributeTok{mfrow =} \FunctionTok{c}\NormalTok{(}\DecValTok{1}\NormalTok{, }\DecValTok{2}\NormalTok{))}

\CommentTok{\# Histograma}
\FunctionTok{hist}\NormalTok{(pad, }\AttributeTok{breaks =} \DecValTok{20}\NormalTok{, }\AttributeTok{freq =} \ConstantTok{FALSE}\NormalTok{,}
     \AttributeTok{main =} \StringTok{"Histograma de PAD"}\NormalTok{,}
     \AttributeTok{xlab =} \StringTok{"Presión Diastólica (mmHg)"}\NormalTok{,}
     \AttributeTok{ylab =} \StringTok{"Densidad"}\NormalTok{,}
     \AttributeTok{col =} \StringTok{"salmon"}\NormalTok{)}

\CommentTok{\# Gráfico de densidad suavizada}
\FunctionTok{plot}\NormalTok{(}\FunctionTok{density}\NormalTok{(pad),}
     \AttributeTok{main =} \StringTok{"Gráfico de Densidad"}\NormalTok{,}
     \AttributeTok{xlab =} \StringTok{"Presión Diastólica (mmHg)"}\NormalTok{,}
     \AttributeTok{ylab =} \StringTok{"Densidad"}\NormalTok{,}
     \AttributeTok{lwd =} \DecValTok{2}\NormalTok{, }\AttributeTok{col =} \StringTok{"darkred"}\NormalTok{)}
\FunctionTok{polygon}\NormalTok{(}\FunctionTok{density}\NormalTok{(pad), }\AttributeTok{col =} \FunctionTok{rgb}\NormalTok{(}\DecValTok{1}\NormalTok{, }\DecValTok{0}\NormalTok{, }\DecValTok{0}\NormalTok{, }\FloatTok{0.3}\NormalTok{))}
\end{Highlighting}
\end{Shaded}

}

\end{codelisting}%

\pandocbounded{\includegraphics[keepaspectratio]{chapters/semana01_files/figure-pdf/unnamed-chunk-4-1.pdf}}

\begin{codelisting}[H]

\caption{\label{lst-semana01-5}Comparación de gráficos}

\centering{

\begin{Shaded}
\begin{Highlighting}[]
\CommentTok{\# Restaurar configuración}
\FunctionTok{par}\NormalTok{(}\AttributeTok{mfrow =} \FunctionTok{c}\NormalTok{(}\DecValTok{1}\NormalTok{, }\DecValTok{1}\NormalTok{))}
\end{Highlighting}
\end{Shaded}

}

\end{codelisting}%

\end{tcolorbox}

\begin{center}\rule{0.5\linewidth}{0.5pt}\end{center}

\section{Clasificación de Medidas
Estadísticas}\label{sec-semana01-clasificacion}

Para describir una distribución, utilizamos diferentes medidas que
resumen la información.

\begin{tcolorbox}[enhanced jigsaw, breakable, titlerule=0mm, rightrule=.15mm, leftrule=.75mm, left=2mm, opacitybacktitle=0.6, colback=white, bottomtitle=1mm, colbacktitle=quarto-callout-note-color!10!white, title=\textcolor{quarto-callout-note-color}{\faInfo}\hspace{0.5em}{Clasificación: Tendencia Central vs.~Posición}, toptitle=1mm, coltitle=black, opacityback=0, arc=.35mm, bottomrule=.15mm, toprule=.15mm, colframe=quarto-callout-note-color-frame]

\begin{itemize}
\tightlist
\item
  \textbf{Medidas de Tendencia Central}: Buscan identificar el ``valor
  representativo'' o centro de la distribución.

  \begin{itemize}
  \tightlist
  \item
    \textbf{Media (\(\bar{x}\))}: Promedio aritmético.
  \item
    \textbf{Mediana (\(\tilde{x}\))}: Valor central (posición central).
  \item
    \textbf{Moda (\(Mo\))}: Valor más frecuente.
  \end{itemize}
\item
  \textbf{Medidas de Posición (No central)}: Indican la posición de un
  valor respecto al conjunto de datos, permitiendo dividir la
  distribución en partes iguales.

  \begin{itemize}
  \tightlist
  \item
    \textbf{Cuartiles (\(Q_1, Q_2, Q_3\))}: Dividen los datos en 4
    partes iguales.
  \item
    \textbf{Percentiles (\(P_k\))}: Dividen los datos en 100 partes
    iguales.
  \end{itemize}
\item
  \textbf{Sesgo (Skewness)}: Medida de asimetría.

  \begin{itemize}
  \tightlist
  \item
    \textbf{Simétrica}: Media \(\approx\) Mediana.
  \item
    \textbf{Sesgo Derecha (Positivo)}: Media \textgreater{} Mediana.
  \item
    \textbf{Sesgo Izquierda (Negativo)}: Media \textless{} Mediana.
  \end{itemize}
\item
  \textbf{Curtosis (Kurtosis)}: Medida del grado de apuntamiento o
  concentración en las colas.
\end{itemize}

\end{tcolorbox}

\begin{center}\rule{0.5\linewidth}{0.5pt}\end{center}

\section{Medidas de Posición
Central}\label{medidas-de-posiciuxf3n-central}

Las medidas de posición central nos indican dónde se concentran los
datos. Son fundamentales para resumir la tendencia central de una
distribución.

\subsection{Media Aritmética}\label{media-aritmuxe9tica}

\begin{tcolorbox}[enhanced jigsaw, breakable, titlerule=0mm, rightrule=.15mm, leftrule=.75mm, left=2mm, opacitybacktitle=0.6, colback=white, bottomtitle=1mm, colbacktitle=quarto-callout-note-color!10!white, title=\textcolor{quarto-callout-note-color}{\faInfo}\hspace{0.5em}{Definición: Media Aritmética}, toptitle=1mm, coltitle=black, opacityback=0, arc=.35mm, bottomrule=.15mm, toprule=.15mm, colframe=quarto-callout-note-color-frame]

La \textbf{media aritmética} de una muestra \(x_1, x_2, \ldots, x_n\)
es:

\[\bar{x} = \frac{1}{n}\sum_{i=1}^{n} x_i\]

Representa el ``centro de gravedad'' de los datos. Es la suma de todos
los valores dividida por el número total de observaciones.

Para datos en una distribución de frecuencias:

\[\bar{x} = \frac{1}{n}\sum_{j=1}^{k} x_j h(x_j) = \sum_{j=1}^{k} x_j f(x_j)\]

donde \(x_j\) es el valor distintivo y \(h(x_j)\) es su frecuencia
absoluta.

\end{tcolorbox}

\begin{tcolorbox}[enhanced jigsaw, breakable, titlerule=0mm, rightrule=.15mm, leftrule=.75mm, left=2mm, opacitybacktitle=0.6, colback=white, bottomtitle=1mm, colbacktitle=quarto-callout-warning-color!10!white, title=\textcolor{quarto-callout-warning-color}{\faExclamationTriangle}\hspace{0.5em}{Propiedades Algebraicas de la Media}, toptitle=1mm, coltitle=black, opacityback=0, arc=.35mm, bottomrule=.15mm, toprule=.15mm, colframe=quarto-callout-warning-color-frame]

\begin{enumerate}
\def\labelenumi{\arabic{enumi}.}
\item
  \textbf{Propiedad de desviación nula}:
  \[\sum_{i=1}^{n}(x_i - \bar{x}) = 0\] Las desviaciones positivas y
  negativas de la media se cancelan.
\item
  \textbf{Propiedad de suma}: Si \(z_i = x_i + y_i\), entonces:
  \[\bar{z} = \bar{x} + \bar{y}\]
\end{enumerate}

\end{tcolorbox}

\subsection{Media para Datos
Agrupados}\label{media-para-datos-agrupados}

\begin{tcolorbox}[enhanced jigsaw, breakable, titlerule=0mm, rightrule=.15mm, leftrule=.75mm, left=2mm, opacitybacktitle=0.6, colback=white, bottomtitle=1mm, colbacktitle=quarto-callout-note-color!10!white, title=\textcolor{quarto-callout-note-color}{\faInfo}\hspace{0.5em}{Definición: Media Aproximada para Datos Binned}, toptitle=1mm, coltitle=black, opacityback=0, arc=.35mm, bottomrule=.15mm, toprule=.15mm, colframe=quarto-callout-note-color-frame]

Cuando solo disponemos de datos agrupados en clases, aproximamos la
media usando la marca de clase:

\[\bar{x} = \frac{1}{n}\sum_{j=1}^{k} x_j \cdot n_j\]

donde \(x_j\) es la \textbf{marca de clase} (punto medio del intervalo
\(j\)): \[x_j = \frac{x_j^l + x_j^u}{2}\]

y \(n_j\) es el número de observaciones en la clase \(j\).

\textbf{Nota}: Esta es una aproximación porque asumimos que todas las
observaciones en una clase están concentradas en su punto medio.

\end{tcolorbox}

\begin{tcolorbox}[enhanced jigsaw, breakable, titlerule=0mm, rightrule=.15mm, leftrule=.75mm, left=2mm, opacitybacktitle=0.6, colback=white, bottomtitle=1mm, colbacktitle=quarto-callout-tip-color!10!white, title=\textcolor{quarto-callout-tip-color}{\faLightbulb}\hspace{0.5em}{Ejemplo 1.9: Tiempo de Estancia Hospitalaria}, toptitle=1mm, coltitle=black, opacityback=0, arc=.35mm, bottomrule=.15mm, toprule=.15mm, colframe=quarto-callout-tip-color-frame]

Se estudia el tiempo de estancia (en días) de una muestra de pacientes
en una unidad de cuidados intensivos. La distribución es:

\begin{longtable}[]{@{}lrrr@{}}
\toprule\noalign{}
Intervalo (días) & Marca de Clase (\(x_j\)) & \(f(x_j)\) &
Contribución \\
\midrule\noalign{}
\endhead
\bottomrule\noalign{}
\endlastfoot
1 -- 3 & 2 & 0.20 & \(2 \times 0.20 = 0.40\) \\
3 -- 7 & 5 & 0.45 & \(5 \times 0.45 = 2.25\) \\
7 -- 14 & 10.5 & 0.25 & \(10.5 \times 0.25 = 2.625\) \\
14 -- 30 & 22 & 0.08 & \(22 \times 0.08 = 1.76\) \\
30 -- 60 & 45 & 0.02 & \(45 \times 0.02 = 0.90\) \\
\end{longtable}

\textbf{Media aproximada}:
\[\bar{x} = 0.40 + 2.25 + 2.625 + 1.76 + 0.90 = 7.935 \text{ días}\]

\begin{codelisting}[H]

\caption{\label{lst-semana01-6}Cálculo de estancia media en R}

\centering{

\begin{Shaded}
\begin{Highlighting}[]
\CommentTok{\# Datos}
\NormalTok{marca\_clase }\OtherTok{\textless{}{-}} \FunctionTok{c}\NormalTok{(}\DecValTok{2}\NormalTok{, }\DecValTok{5}\NormalTok{, }\FloatTok{10.5}\NormalTok{, }\DecValTok{22}\NormalTok{, }\DecValTok{45}\NormalTok{)}
\NormalTok{frecuencia\_rel }\OtherTok{\textless{}{-}} \FunctionTok{c}\NormalTok{(}\FloatTok{0.20}\NormalTok{, }\FloatTok{0.45}\NormalTok{, }\FloatTok{0.25}\NormalTok{, }\FloatTok{0.08}\NormalTok{, }\FloatTok{0.02}\NormalTok{)}

\CommentTok{\# Media}
\NormalTok{media\_estancia }\OtherTok{\textless{}{-}} \FunctionTok{sum}\NormalTok{(marca\_clase }\SpecialCharTok{*}\NormalTok{ frecuencia\_rel)}
\FunctionTok{print}\NormalTok{(}\FunctionTok{paste}\NormalTok{(}\StringTok{"Media de estancia aproximada:"}\NormalTok{,}
            \FunctionTok{round}\NormalTok{(media\_estancia, }\DecValTok{3}\NormalTok{), }\StringTok{"días"}\NormalTok{))}
\end{Highlighting}
\end{Shaded}

}

\end{codelisting}%

\begin{verbatim}
[1] "Media de estancia aproximada: 7.935 días"
\end{verbatim}

\end{tcolorbox}

\subsection{Moda}\label{moda}

\begin{tcolorbox}[enhanced jigsaw, breakable, titlerule=0mm, rightrule=.15mm, leftrule=.75mm, left=2mm, opacitybacktitle=0.6, colback=white, bottomtitle=1mm, colbacktitle=quarto-callout-note-color!10!white, title=\textcolor{quarto-callout-note-color}{\faInfo}\hspace{0.5em}{Definición: Moda}, toptitle=1mm, coltitle=black, opacityback=0, arc=.35mm, bottomrule=.15mm, toprule=.15mm, colframe=quarto-callout-note-color-frame]

La \textbf{moda} \(x_D\) es el valor más frecuente en un conjunto de
observaciones:

\[x_D = \{x_j \mid h(x_j) = \max_k h(x_k) \text{ o } f(x_j) = \max_k f(x_k)\}\]

Es útil para variables nominales, ordinales, discretas o clasificadas,
pero \textbf{no es apropiada para variables continuas}.

Para datos agrupados, hablamos de la \textbf{clase modal} (la clase con
la frecuencia más alta) y estimamos la moda como la marca de clase.

\end{tcolorbox}

\begin{tcolorbox}[enhanced jigsaw, breakable, titlerule=0mm, rightrule=.15mm, leftrule=.75mm, left=2mm, opacitybacktitle=0.6, colback=white, bottomtitle=1mm, colbacktitle=quarto-callout-tip-color!10!white, title=\textcolor{quarto-callout-tip-color}{\faLightbulb}\hspace{0.5em}{Ejemplo 1.10: Moda en Tiempo de Supervivencia}, toptitle=1mm, coltitle=black, opacityback=0, arc=.35mm, bottomrule=.15mm, toprule=.15mm, colframe=quarto-callout-tip-color-frame]

Usando la tabla de supervivencia del Ejemplo 1.6:

\begin{longtable}[]{@{}lr@{}}
\toprule\noalign{}
Intervalo (días) & \(h_j\) \\
\midrule\noalign{}
\endhead
\bottomrule\noalign{}
\endlastfoot
0 -- 100 & 1 \\
100 -- 500 & 24 \\
\textbf{500 -- 1000} & \textbf{45} (máximo) \\
1000 -- 4000 & 30 \\
\end{longtable}

\begin{itemize}
\tightlist
\item
  \textbf{Clase modal}: 500 -- 1000 días
\item
  \textbf{Moda estimada}: \(x_D = 750\) días (marca de clase)
\end{itemize}

\end{tcolorbox}

\subsection{Mediana}\label{mediana}

\begin{tcolorbox}[enhanced jigsaw, breakable, titlerule=0mm, rightrule=.15mm, leftrule=.75mm, left=2mm, opacitybacktitle=0.6, colback=white, bottomtitle=1mm, colbacktitle=quarto-callout-note-color!10!white, title=\textcolor{quarto-callout-note-color}{\faInfo}\hspace{0.5em}{Definición: Mediana}, toptitle=1mm, coltitle=black, opacityback=0, arc=.35mm, bottomrule=.15mm, toprule=.15mm, colframe=quarto-callout-note-color-frame]

La \textbf{mediana} \(x_{0.5}\) es el valor que divide los datos
ordenados en dos partes de igual tamaño: 50\% de las observaciones están
por debajo y 50\% están por encima.

\textbf{Para datos no agrupados}: - Si \(n\) es \textbf{impar}:
\[x_{0.5} = x_{\left(\frac{n+1}{2}\right)}\]

\begin{itemize}
\tightlist
\item
  Si \(n\) es \textbf{par}:
  \[x_{0.5} = \frac{1}{2}\left\{x_{\left(\frac{n}{2}\right)} + x_{\left(\frac{n}{2}+1\right)}\right\}\]
\end{itemize}

donde \(x_{(i)}\) es la \(i\)-ésima observación cuando los datos están
\textbf{ordenados de menor a mayor}.

\textbf{Ventaja}: La mediana es \textbf{robusta} (insensible a
outliers), a diferencia de la media.

\end{tcolorbox}

\begin{tcolorbox}[enhanced jigsaw, breakable, titlerule=0mm, rightrule=.15mm, leftrule=.75mm, left=2mm, opacitybacktitle=0.6, colback=white, bottomtitle=1mm, colbacktitle=quarto-callout-tip-color!10!white, title=\textcolor{quarto-callout-tip-color}{\faLightbulb}\hspace{0.5em}{Ejemplo 1.11: Cálculo de Mediana}, toptitle=1mm, coltitle=black, opacityback=0, arc=.35mm, bottomrule=.15mm, toprule=.15mm, colframe=quarto-callout-tip-color-frame]

\textbf{Conjunto 1} (n impar): 115, 118, 120, 122, 124, 125, 128, 129,
130, 135

Datos ya ordenados. Con \(n = 10\) (par):
\[x_{0.5} = \frac{1}{2}(x_{(5)} + x_{(6)}) = \frac{1}{2}(124 + 125) = 124.5 \text{ mmHg}\]

\begin{codelisting}[H]

\caption{\label{lst-semana01-7}Cálculo de mediana en R}

\centering{

\begin{Shaded}
\begin{Highlighting}[]
\NormalTok{pas }\OtherTok{\textless{}{-}} \FunctionTok{c}\NormalTok{(}\DecValTok{120}\NormalTok{, }\DecValTok{125}\NormalTok{, }\DecValTok{118}\NormalTok{, }\DecValTok{130}\NormalTok{, }\DecValTok{122}\NormalTok{, }\DecValTok{128}\NormalTok{, }\DecValTok{115}\NormalTok{, }\DecValTok{135}\NormalTok{, }\DecValTok{124}\NormalTok{, }\DecValTok{129}\NormalTok{)}
\NormalTok{mediana\_pas }\OtherTok{\textless{}{-}} \FunctionTok{median}\NormalTok{(pas)}
\FunctionTok{print}\NormalTok{(}\FunctionTok{paste}\NormalTok{(}\StringTok{"Mediana de PAS:"}\NormalTok{, mediana\_pas, }\StringTok{"mmHg"}\NormalTok{))}
\end{Highlighting}
\end{Shaded}

}

\end{codelisting}%

\begin{verbatim}
[1] "Mediana de PAS: 124.5 mmHg"
\end{verbatim}

\end{tcolorbox}

\subsection{Relación entre Media, Mediana y
Moda}\label{relaciuxf3n-entre-media-mediana-y-moda}

La siguiente figura ilustra la posición relativa de la media, mediana y
moda según el sesgo de la distribución:

\begin{figure}

\centering{

\pandocbounded{\includegraphics[keepaspectratio]{chapters/semana01_files/figure-pdf/fig-semana01-skewness-1.pdf}}

}

\caption{\label{fig-semana01-skewness}Relación entre media, mediana y
moda según el sesgo}

\end{figure}%

\section{Medidas de Forma: Sesgo y
Curtosis}\label{medidas-de-forma-sesgo-y-curtosis}

Además de la tendencia central y la dispersión, la \textbf{forma} de la
distribución es crucial para entender el comportamiento de los datos.

\subsubsection{Coeficiente de Asimetría
(Sesgo)}\label{coeficiente-de-asimetruxeda-sesgo}

Mide el grado de simetría de la distribución respecto a su media. El
coeficiente de asimetría de Fisher se define como:

\[g_1 = \frac{\frac{1}{n} \sum_{i=1}^{n} (x_i - \bar{x})^3}{s^3}\]

donde \(s\) es la desviación típica.

\textbf{Interpretación}: - \textbf{\(g_1 = 0\)}: La distribución es
\textbf{simétrica}. - \textbf{\(g_1 > 0\)}: \textbf{Asimetría positiva}
(sesgo a la derecha). La cola derecha es más larga. -
\textbf{\(g_1 < 0\)}: \textbf{Asimetría negativa} (sesgo a la
izquierda). La cola izquierda es más larga.

\subsubsection{Curtosis o Apuntamiento}\label{curtosis-o-apuntamiento}

Mide el grado de concentración de los valores alrededor de la zona
central de la distribución y el ``peso'' de las colas. Se calcula
mediante el coeficiente de curtosis de Fisher:

\[g_2 = \frac{\frac{1}{n} \sum_{i=1}^{n} (x_i - \bar{x})^4}{s^4}\]

En muchos paquetes estadísticos (incluido R), se utiliza la
\textbf{curtosis en exceso} (\(g_2 - 3\)) para comparar con la
distribución normal.

\textbf{Interpretación (Curtosis en exceso)}: -
\textbf{\(g_2 - 3 = 0\)}: \textbf{Mesocúrtica}. El apuntamiento es
similar al de una distribución normal. - \textbf{\(g_2 - 3 > 0\)}:
\textbf{Leptocúrtica}. La distribución es más apuntada y tiene colas más
gruesas (``pesadas'') que la normal. - \textbf{\(g_2 - 3 < 0\)}:
\textbf{Platicúrtica}. La distribución es más plana y tiene colas más
delgadas que la normal.

\subsection{Cálculo con R}\label{sec-semana01-r-stats}

\begin{codelisting}

\caption{\label{lst-semana01-stats}Cálculo de estadísticos básicos}

\centering{

\begin{Shaded}
\begin{Highlighting}[]
\FunctionTok{library}\NormalTok{(moments)}
\FunctionTok{library}\NormalTok{(dplyr)}

\CommentTok{\# Datos de ejemplo}
\NormalTok{datos }\OtherTok{\textless{}{-}} \FunctionTok{c}\NormalTok{(}\DecValTok{120}\NormalTok{, }\DecValTok{125}\NormalTok{, }\DecValTok{118}\NormalTok{, }\DecValTok{130}\NormalTok{, }\DecValTok{122}\NormalTok{, }\DecValTok{128}\NormalTok{, }\DecValTok{115}\NormalTok{, }\DecValTok{135}\NormalTok{, }\DecValTok{124}\NormalTok{, }\DecValTok{129}\NormalTok{, }\DecValTok{150}\NormalTok{)}

\CommentTok{\# Resumen básico}
\NormalTok{resumen }\OtherTok{\textless{}{-}} \FunctionTok{data.frame}\NormalTok{(}
  \AttributeTok{Media =} \FunctionTok{mean}\NormalTok{(datos),}
  \AttributeTok{Mediana =} \FunctionTok{median}\NormalTok{(datos),}
  \AttributeTok{Sesgo =} \FunctionTok{skewness}\NormalTok{(datos), }\CommentTok{\# g1}
  \AttributeTok{Curtosis =} \FunctionTok{kurtosis}\NormalTok{(datos), }\CommentTok{\# g2 (Pearson, restar 3 para exceso)}
  \AttributeTok{Curtosis\_Exceso =} \FunctionTok{kurtosis}\NormalTok{(datos) }\SpecialCharTok{{-}} \DecValTok{3}
\NormalTok{)}
\FunctionTok{print}\NormalTok{(resumen)}
\end{Highlighting}
\end{Shaded}

}

\end{codelisting}%

\begin{verbatim}
     Media Mediana    Sesgo Curtosis Curtosis_Exceso
1 126.9091     125 1.191328  4.14324         1.14324
\end{verbatim}

\subsection{Cuantiles}\label{cuantiles}

\begin{tcolorbox}[enhanced jigsaw, breakable, titlerule=0mm, rightrule=.15mm, leftrule=.75mm, left=2mm, opacitybacktitle=0.6, colback=white, bottomtitle=1mm, colbacktitle=quarto-callout-note-color!10!white, title=\textcolor{quarto-callout-note-color}{\faInfo}\hspace{0.5em}{Definición: Cuantiles}, toptitle=1mm, coltitle=black, opacityback=0, arc=.35mm, bottomrule=.15mm, toprule=.15mm, colframe=quarto-callout-note-color-frame]

El \textbf{cuantil \(x_p\)} es el valor que divide los datos ordenados
en la proporción \(p\) a \((1-p)\), donde \(0 \leq p \leq 1\).

\begin{itemize}
\tightlist
\item
  Debajo del cuantil \(x_p\) se encuentran aproximadamente
  \(p \times 100\%\) de los datos
\item
  Arriba del cuantil \(x_p\) se encuentran aproximadamente
  \((1-p) \times 100\%\) de los datos
\end{itemize}

\textbf{Cuantiles especiales}: - \textbf{Cuartiles}:
\(p = 0.25, 0.50, 0.75\) (dividen en 4 partes) - \(x_{0.25}\) = primer
cuartil (Q1) - \(x_{0.50}\) = mediana - \(x_{0.75}\) = tercer cuartil
(Q3)

\begin{itemize}
\item
  \textbf{Deciles}: \(p = 0.1, 0.2, \ldots, 0.9\) (dividen en 10 partes)
\item
  \textbf{Percentiles}: \(p = 0.01, 0.02, \ldots, 0.99\) (dividen en 100
  partes)
\end{itemize}

\end{tcolorbox}

\begin{tcolorbox}[enhanced jigsaw, breakable, titlerule=0mm, rightrule=.15mm, leftrule=.75mm, left=2mm, opacitybacktitle=0.6, colback=white, bottomtitle=1mm, colbacktitle=quarto-callout-tip-color!10!white, title=\textcolor{quarto-callout-tip-color}{\faLightbulb}\hspace{0.5em}{Ejemplo 1.12: Cuartiles de Supervivencia}, toptitle=1mm, coltitle=black, opacityback=0, arc=.35mm, bottomrule=.15mm, toprule=.15mm, colframe=quarto-callout-tip-color-frame]

Con los datos de supervivencia de pacientes (Ejemplo 1.6), podemos
aproximar los cuartiles usando interpolación en la ECDF:

Para el primer cuartil (\(p = 0.25\)): - \(F(x_j^l) = 0.01\) (límite
inferior de clase 100--500) - \(f(x_j) = 0.24\) (frecuencia relativa de
clase 100--500) - Fórmula:
\(x_{0.25} = x_j^l + \frac{p - F(x_j^l)}{f(x_j)} \cdot \Delta x_j\)

\[x_{0.25} = 100 + \frac{0.25 - 0.01}{0.24} \times 400 = 100 + 400 = 500 \text{ días}\]

\end{tcolorbox}

\begin{codelisting}

\caption{\label{lst-semana01-8}Cálculo de cuantiles}

\centering{

\begin{Shaded}
\begin{Highlighting}[]
\CommentTok{\# Generamos los mismos datos que en el Ejemplo 1.4.2}
\FunctionTok{set.seed}\NormalTok{(}\DecValTok{42}\NormalTok{)}
\NormalTok{supervivencia }\OtherTok{\textless{}{-}} \FunctionTok{c}\NormalTok{(}
  \FunctionTok{runif}\NormalTok{(}\DecValTok{1}\NormalTok{, }\DecValTok{0}\NormalTok{, }\DecValTok{100}\NormalTok{),}
  \FunctionTok{runif}\NormalTok{(}\DecValTok{24}\NormalTok{, }\DecValTok{100}\NormalTok{, }\DecValTok{500}\NormalTok{),}
  \FunctionTok{runif}\NormalTok{(}\DecValTok{45}\NormalTok{, }\DecValTok{500}\NormalTok{, }\DecValTok{1000}\NormalTok{),}
  \FunctionTok{runif}\NormalTok{(}\DecValTok{25}\NormalTok{, }\DecValTok{1000}\NormalTok{, }\DecValTok{2000}\NormalTok{),}
  \DecValTok{3200}\NormalTok{, }\DecValTok{3500}\NormalTok{, }\DecValTok{3700}\NormalTok{, }\DecValTok{3800}\NormalTok{, }\DecValTok{3950}
\NormalTok{)}

\NormalTok{Q1 }\OtherTok{\textless{}{-}} \FunctionTok{quantile}\NormalTok{(supervivencia, }\AttributeTok{probs =} \FloatTok{0.25}\NormalTok{)}
\NormalTok{Q2 }\OtherTok{\textless{}{-}} \FunctionTok{quantile}\NormalTok{(supervivencia, }\AttributeTok{probs =} \FloatTok{0.50}\NormalTok{)}
\NormalTok{Q3 }\OtherTok{\textless{}{-}} \FunctionTok{quantile}\NormalTok{(supervivencia, }\AttributeTok{probs =} \FloatTok{0.75}\NormalTok{)}

\FunctionTok{print}\NormalTok{(}\FunctionTok{paste}\NormalTok{(}\StringTok{"Q1 (25\%):"}\NormalTok{, }\FunctionTok{round}\NormalTok{(Q1, }\DecValTok{2}\NormalTok{)))}
\end{Highlighting}
\end{Shaded}

}

\end{codelisting}%

\begin{verbatim}
[1] "Q1 (25%): 500.37"
\end{verbatim}

\begin{codelisting}

\caption{\label{lst-semana01-8}Cálculo de cuantiles}

\centering{

\begin{Shaded}
\begin{Highlighting}[]
\FunctionTok{print}\NormalTok{(}\FunctionTok{paste}\NormalTok{(}\StringTok{"Q2 (50\%, mediana):"}\NormalTok{, }\FunctionTok{round}\NormalTok{(Q2, }\DecValTok{2}\NormalTok{)))}
\end{Highlighting}
\end{Shaded}

}

\end{codelisting}%

\begin{verbatim}
[1] "Q2 (50%, mediana): 828.9"
\end{verbatim}

\begin{codelisting}

\caption{\label{lst-semana01-8}Cálculo de cuantiles}

\centering{

\begin{Shaded}
\begin{Highlighting}[]
\FunctionTok{print}\NormalTok{(}\FunctionTok{paste}\NormalTok{(}\StringTok{"Q3 (75\%):"}\NormalTok{, }\FunctionTok{round}\NormalTok{(Q3, }\DecValTok{2}\NormalTok{)))}
\end{Highlighting}
\end{Shaded}

}

\end{codelisting}%

\begin{verbatim}
[1] "Q3 (75%): 1086.7"
\end{verbatim}

\begin{tcolorbox}[enhanced jigsaw, breakable, titlerule=0mm, rightrule=.15mm, leftrule=.75mm, left=2mm, opacitybacktitle=0.6, colback=white, bottomtitle=1mm, colbacktitle=quarto-callout-tip-color!10!white, title=\textcolor{quarto-callout-tip-color}{\faLightbulb}\hspace{0.5em}{Visualización: Diagrama de Caja (Boxplot)}, toptitle=1mm, coltitle=black, opacityback=0, arc=.35mm, bottomrule=.15mm, toprule=.15mm, colframe=quarto-callout-tip-color-frame]

El \textbf{diagrama de caja} (boxplot) es una herramienta robusta para
visualizar la distribución, la mediana, los cuartiles y detectar valores
atípicos (outliers) de forma intuitiva.

\begin{codelisting}[H]

\caption{\label{lst-semana01-9}Boxplot mejorado con ggplot2}

\centering{

\begin{Shaded}
\begin{Highlighting}[]
\FunctionTok{library}\NormalTok{(ggplot2)}

\CommentTok{\# Crear un data frame con los datos de supervivencia}
\NormalTok{df }\OtherTok{\textless{}{-}} \FunctionTok{data.frame}\NormalTok{(}\AttributeTok{supervivencia =}\NormalTok{ supervivencia)}

\CommentTok{\# Visualización con ggplot2 y labels para cuartiles}
\FunctionTok{ggplot}\NormalTok{(df, }\FunctionTok{aes}\NormalTok{(}\AttributeTok{y =}\NormalTok{ supervivencia)) }\SpecialCharTok{+}
  \FunctionTok{geom\_boxplot}\NormalTok{(}\AttributeTok{fill =} \StringTok{"steelblue"}\NormalTok{, }\AttributeTok{alpha =} \FloatTok{0.7}\NormalTok{, }
               \AttributeTok{outlier.colour =} \StringTok{"red"}\NormalTok{, }\AttributeTok{outlier.shape =} \DecValTok{16}\NormalTok{, }\AttributeTok{outlier.size =} \DecValTok{3}\NormalTok{) }\SpecialCharTok{+}
  \FunctionTok{labs}\NormalTok{(}\AttributeTok{title =} \StringTok{"Distribución del Tiempo de Supervivencia"}\NormalTok{,}
       \AttributeTok{subtitle =} \StringTok{"Se observan outliers en la parte superior"}\NormalTok{,}
       \AttributeTok{y =} \StringTok{"Días"}\NormalTok{,}
       \AttributeTok{x =} \StringTok{""}\NormalTok{) }\SpecialCharTok{+}
  \FunctionTok{theme\_minimal}\NormalTok{() }\SpecialCharTok{+}
  \FunctionTok{coord\_flip}\NormalTok{() }\CommentTok{\# Para orientación horizontal}
\end{Highlighting}
\end{Shaded}

}

\end{codelisting}%

\pandocbounded{\includegraphics[keepaspectratio]{chapters/semana01_files/figure-pdf/unnamed-chunk-9-1.pdf}}

\end{tcolorbox}

\begin{tcolorbox}[enhanced jigsaw, breakable, titlerule=0mm, rightrule=.15mm, leftrule=.75mm, left=2mm, opacitybacktitle=0.6, colback=white, bottomtitle=1mm, colbacktitle=quarto-callout-tip-color!10!white, title=\textcolor{quarto-callout-tip-color}{\faLightbulb}\hspace{0.5em}{Lectura del Boxplot}, toptitle=1mm, coltitle=black, opacityback=0, arc=.35mm, bottomrule=.15mm, toprule=.15mm, colframe=quarto-callout-tip-color-frame]

Un boxplot divide los datos en:\\
- \textbf{Caja}: Contiene el 50\% central (IQR: de \(Q_1\) a \(Q_3\)).\\
- \textbf{Línea central}: Mediana.\\
- \textbf{Bigotes}: Extienden hasta 1.5 veces el IQR.\\
- \textbf{Puntos (outliers)}: Valores fuera del rango de los bigotes.

\end{tcolorbox}

\begin{center}\rule{0.5\linewidth}{0.5pt}\end{center}

\section{Comparación de Medidas de
Posición}\label{comparaciuxf3n-de-medidas-de-posiciuxf3n}

\begin{tcolorbox}[enhanced jigsaw, breakable, titlerule=0mm, rightrule=.15mm, leftrule=.75mm, left=2mm, opacitybacktitle=0.6, colback=white, bottomtitle=1mm, colbacktitle=quarto-callout-warning-color!10!white, title=\textcolor{quarto-callout-warning-color}{\faExclamationTriangle}\hspace{0.5em}{Robustez de las Medidas de Tendencia Central}, toptitle=1mm, coltitle=black, opacityback=0, arc=.35mm, bottomrule=.15mm, toprule=.15mm, colframe=quarto-callout-warning-color-frame]

La \textbf{media} es muy sensible a valores extremos (outliers): - Fácil
de calcular algebraicamente - Propiedad de suma útil para combinación de
datos - Se ve afectada por datos atípicos

La \textbf{mediana} es robusta: - Insensible a outliers - Más
representativa en distribuciones asimétricas - Requiere ordenar los
datos

La \textbf{moda} es útil para datos categóricos pero limitado para
continuos.

\textbf{Recomendación}: En presencia de outliers, la mediana es
preferible a la media como medida de tendencia central.

\end{tcolorbox}

\begin{center}\rule{0.5\linewidth}{0.5pt}\end{center}

\section{Resumen de Conceptos Clave}\label{sec-semana01-resumen}

\subsection{Tabla Resumen: Escalas de Medición y Estadísticos
Apropiados}\label{tabla-resumen-escalas-de-mediciuxf3n-y-estaduxedsticos-apropiados}

\begin{longtable}[]{@{}
  >{\raggedright\arraybackslash}p{(\linewidth - 4\tabcolsep) * \real{0.3333}}
  >{\raggedright\arraybackslash}p{(\linewidth - 4\tabcolsep) * \real{0.3333}}
  >{\raggedright\arraybackslash}p{(\linewidth - 4\tabcolsep) * \real{0.3333}}@{}}
\toprule\noalign{}
\begin{minipage}[b]{\linewidth}\raggedright
Tipo de Variable
\end{minipage} & \begin{minipage}[b]{\linewidth}\raggedright
Ejemplos
\end{minipage} & \begin{minipage}[b]{\linewidth}\raggedright
Estadísticos Apropiados
\end{minipage} \\
\midrule\noalign{}
\endhead
\bottomrule\noalign{}
\endlastfoot
\textbf{Nominal} & Género, grupo sanguíneo & Moda, frecuencias \\
\textbf{Ordinal} & Satisfacción, grado & Mediana, moda, cuantiles \\
\textbf{Discreta} & Número de eventos & Media, mediana, varianza \\
\textbf{Continua} & Peso, presión arterial & Media, mediana, varianza,
cuantiles \\
\end{longtable}

\subsection{Tabla Resumen: Fórmulas
Importantes}\label{tabla-resumen-fuxf3rmulas-importantes}

\begin{longtable}[]{@{}
  >{\raggedright\arraybackslash}p{(\linewidth - 4\tabcolsep) * \real{0.3333}}
  >{\raggedright\arraybackslash}p{(\linewidth - 4\tabcolsep) * \real{0.3333}}
  >{\raggedright\arraybackslash}p{(\linewidth - 4\tabcolsep) * \real{0.3333}}@{}}
\toprule\noalign{}
\begin{minipage}[b]{\linewidth}\raggedright
Estadístico
\end{minipage} & \begin{minipage}[b]{\linewidth}\raggedright
Datos No Agrupados
\end{minipage} & \begin{minipage}[b]{\linewidth}\raggedright
Datos Agrupados (Intervalos)
\end{minipage} \\
\midrule\noalign{}
\endhead
\bottomrule\noalign{}
\endlastfoot
\textbf{Media (\(\bar{x}\))} & \(\frac{1}{n}\sum_{i=1}^{n} x_i\) &
\(\sum_{j=1}^{k} x_j f_j\) (\(x_j\): marca clase) \\
\textbf{Mediana/Cuantil (\(x_p\))} & Valor en pos. ordenado &
\(x_j^l + \frac{p - F(x_j^l)}{f(x_j)} \cdot \Delta x_j\) \\
\textbf{Moda (\(Mo\))} & Valor con máx. \(h_j\) & Marca de la clase
modal \\
\textbf{ECDF (\(F(x)\))} & Función escalonada & Interpolación lineal
(ojiva) \\
\textbf{Sesgo (\(g_1\))} &
\(\frac{\frac{1}{n} \sum (x_i - \bar{x})^3}{s^3}\) & (Aprox. marcas de
clase) \\
\textbf{Curtosis (\(g_2\))} &
\(\frac{\frac{1}{n} \sum (x_i - \bar{x})^4}{s^4}\) & (Aprox. marcas de
clase) \\
\end{longtable}

\begin{center}\rule{0.5\linewidth}{0.5pt}\end{center}

\section{Ejercicios}\label{sec-semana01-ejercicios}

\begin{tcolorbox}[enhanced jigsaw, breakable, titlerule=0mm, rightrule=.15mm, leftrule=.75mm, left=2mm, opacitybacktitle=0.6, colback=white, bottomtitle=1mm, colbacktitle=quarto-callout-tip-color!10!white, title=\textcolor{quarto-callout-tip-color}{\faLightbulb}\hspace{0.5em}{Ejercicio 1.1: Conceptos Fundamentales}, toptitle=1mm, coltitle=black, opacityback=0, arc=.35mm, bottomrule=.15mm, toprule=.15mm, colframe=quarto-callout-tip-color-frame]

Define con tus propias palabras (desde una perspectiva clínica): a) La
diferencia entre una variable y un valor (ej. Glucemia) b) Por qué
necesitamos distinguir entre población y muestra en un ensayo clínico c)
Qué es un outlier y cómo podría afectar a la media de una serie de
mediciones de tensión arterial

\end{tcolorbox}

\begin{tcolorbox}[enhanced jigsaw, breakable, titlerule=0mm, rightrule=.15mm, leftrule=.75mm, left=2mm, opacitybacktitle=0.6, colback=white, bottomtitle=1mm, colbacktitle=quarto-callout-tip-color!10!white, title=\textcolor{quarto-callout-tip-color}{\faLightbulb}\hspace{0.5em}{Ejercicio 1.2: Identificación de Escalas}, toptitle=1mm, coltitle=black, opacityback=0, arc=.35mm, bottomrule=.15mm, toprule=.15mm, colframe=quarto-callout-tip-color-frame]

Para cada una de las siguientes variables clínicas, identifica la escala
de medición (nominal, ordinal, discreta, continua) y justifica tu
respuesta:

\begin{enumerate}
\def\labelenumi{\alph{enumi})}
\tightlist
\item
  Número de episodios asmáticos al año
\item
  Gravedad de una quemadura (Primer, Segundo, Tercer grado)
\item
  Concentración de hemoglobina en sangre (g/dL)
\item
  Factor Rh (Positivo, Negativo)
\item
  Escala de coma de Glasgow (puntuación de 3 a 15)
\end{enumerate}

\end{tcolorbox}

\begin{tcolorbox}[enhanced jigsaw, breakable, titlerule=0mm, rightrule=.15mm, leftrule=.75mm, left=2mm, opacitybacktitle=0.6, colback=white, bottomtitle=1mm, colbacktitle=quarto-callout-tip-color!10!white, title=\textcolor{quarto-callout-tip-color}{\faLightbulb}\hspace{0.5em}{Ejercicio 1.3: Cálculo de la Media}, toptitle=1mm, coltitle=black, opacityback=0, arc=.35mm, bottomrule=.15mm, toprule=.15mm, colframe=quarto-callout-tip-color-frame]

Se registraron los tiempos de espera (en minutos) en urgencias de 8
pacientes: 12, 8, 15, 10, 7, 14, 9, 11

\begin{enumerate}
\def\labelenumi{\alph{enumi})}
\tightlist
\item
  Calcula la media aritmética
\item
  Ordena los datos y encuentra la mediana
\item
  ¿Cuál es más representativa en este caso? ¿Por qué?
\end{enumerate}

\begin{codelisting}[H]

\caption{\label{lst-semana01-10}Solución en R}

\centering{

\begin{Shaded}
\begin{Highlighting}[]
\NormalTok{tiempos }\OtherTok{\textless{}{-}} \FunctionTok{c}\NormalTok{(}\DecValTok{12}\NormalTok{, }\DecValTok{8}\NormalTok{, }\DecValTok{15}\NormalTok{, }\DecValTok{10}\NormalTok{, }\DecValTok{7}\NormalTok{, }\DecValTok{14}\NormalTok{, }\DecValTok{9}\NormalTok{, }\DecValTok{11}\NormalTok{)}
\NormalTok{media }\OtherTok{\textless{}{-}} \FunctionTok{mean}\NormalTok{(tiempos)}
\NormalTok{mediana }\OtherTok{\textless{}{-}} \FunctionTok{median}\NormalTok{(tiempos)}
\FunctionTok{cat}\NormalTok{(}\StringTok{"Media:"}\NormalTok{, media, }\StringTok{"}\SpecialCharTok{\textbackslash{}n}\StringTok{"}\NormalTok{)}
\end{Highlighting}
\end{Shaded}

}

\end{codelisting}%

\begin{verbatim}
Media: 10.75 
\end{verbatim}

\begin{codelisting}[H]

\caption{\label{lst-semana01-10}Solución en R}

\centering{

\begin{Shaded}
\begin{Highlighting}[]
\FunctionTok{cat}\NormalTok{(}\StringTok{"Mediana:"}\NormalTok{, mediana, }\StringTok{"}\SpecialCharTok{\textbackslash{}n}\StringTok{"}\NormalTok{)}
\end{Highlighting}
\end{Shaded}

}

\end{codelisting}%

\begin{verbatim}
Mediana: 10.5 
\end{verbatim}

\end{tcolorbox}

\begin{tcolorbox}[enhanced jigsaw, breakable, titlerule=0mm, rightrule=.15mm, leftrule=.75mm, left=2mm, opacitybacktitle=0.6, colback=white, bottomtitle=1mm, colbacktitle=quarto-callout-tip-color!10!white, title=\textcolor{quarto-callout-tip-color}{\faLightbulb}\hspace{0.5em}{Ejercicio 1.4: Tabla de Frecuencias}, toptitle=1mm, coltitle=black, opacityback=0, arc=.35mm, bottomrule=.15mm, toprule=.15mm, colframe=quarto-callout-tip-color-frame]

Los siguientes datos representan el número de admisiones hospitalarias
previas para 20 pacientes crónicos: 3, 2, 5, 2, 4, 3, 6, 2, 3, 4, 5, 2,
3, 4, 2, 5, 3, 6, 4, 2

\begin{enumerate}
\def\labelenumi{\alph{enumi})}
\tightlist
\item
  Construye una tabla de frecuencias (absoluta, relativa, acumulada)
\item
  ¿Cuál es la moda (número más frecuente de admisiones)?
\item
  ¿Qué proporción de pacientes tuvieron 3 o menos admisiones?
\end{enumerate}

\end{tcolorbox}

\begin{tcolorbox}[enhanced jigsaw, breakable, titlerule=0mm, rightrule=.15mm, leftrule=.75mm, left=2mm, opacitybacktitle=0.6, colback=white, bottomtitle=1mm, colbacktitle=quarto-callout-tip-color!10!white, title=\textcolor{quarto-callout-tip-color}{\faLightbulb}\hspace{0.5em}{Ejercicio 1.5: Datos Agrupados}, toptitle=1mm, coltitle=black, opacityback=0, arc=.35mm, bottomrule=.15mm, toprule=.15mm, colframe=quarto-callout-tip-color-frame]

Se agrupan los niveles de glucosa en ayunas (mg/dL) de 100 pacientes en
riesgo de diabetes:

\begin{longtable}[]{@{}lr@{}}
\toprule\noalign{}
Intervalo (mg/dL) & Frecuencia (\(h_j\)) \\
\midrule\noalign{}
\endhead
\bottomrule\noalign{}
\endlastfoot
70 -- 90 & 15 \\
90 -- 100 & 32 \\
100 -- 110 & 35 \\
110 -- 126 & 15 \\
126 -- 150 & 3 \\
\end{longtable}

\begin{enumerate}
\def\labelenumi{\alph{enumi})}
\tightlist
\item
  Calcula la media aproximada de la glucemia
\item
  Identifica la clase modal
\item
  Estima el percentil 50 (mediana) usando interpolación
\end{enumerate}

\end{tcolorbox}

\begin{tcolorbox}[enhanced jigsaw, breakable, titlerule=0mm, rightrule=.15mm, leftrule=.75mm, left=2mm, opacitybacktitle=0.6, colback=white, bottomtitle=1mm, colbacktitle=quarto-callout-tip-color!10!white, title=\textcolor{quarto-callout-tip-color}{\faLightbulb}\hspace{0.5em}{Tip}, toptitle=1mm, coltitle=black, opacityback=0, arc=.35mm, bottomrule=.15mm, toprule=.15mm, colframe=quarto-callout-tip-color-frame]

\begin{codelisting}[H]

\caption{\label{lst-semana01-11}Solución en R}

\centering{

\begin{Shaded}
\begin{Highlighting}[]
\CommentTok{\# Datos}
\NormalTok{marca\_clase }\OtherTok{\textless{}{-}} \FunctionTok{c}\NormalTok{(}\DecValTok{80}\NormalTok{, }\DecValTok{95}\NormalTok{, }\DecValTok{105}\NormalTok{, }\DecValTok{118}\NormalTok{, }\DecValTok{138}\NormalTok{)}
\NormalTok{frecuencia }\OtherTok{\textless{}{-}} \FunctionTok{c}\NormalTok{(}\DecValTok{15}\NormalTok{, }\DecValTok{32}\NormalTok{, }\DecValTok{35}\NormalTok{, }\DecValTok{15}\NormalTok{, }\DecValTok{3}\NormalTok{)}
\NormalTok{n }\OtherTok{\textless{}{-}} \FunctionTok{sum}\NormalTok{(frecuencia)}

\CommentTok{\# Media}
\NormalTok{media\_aprox }\OtherTok{\textless{}{-}} \FunctionTok{sum}\NormalTok{(marca\_clase }\SpecialCharTok{*}\NormalTok{ frecuencia) }\SpecialCharTok{/}\NormalTok{ n}
\FunctionTok{print}\NormalTok{(}\FunctionTok{paste}\NormalTok{(}\StringTok{"Media aproximada:"}\NormalTok{, }\FunctionTok{round}\NormalTok{(media\_aprox, }\DecValTok{2}\NormalTok{), }\StringTok{"mg/dL"}\NormalTok{))}
\end{Highlighting}
\end{Shaded}

}

\end{codelisting}%

\begin{verbatim}
[1] "Media aproximada: 100.99 mg/dL"
\end{verbatim}

\begin{codelisting}[H]

\caption{\label{lst-semana01-11}Solución en R}

\centering{

\begin{Shaded}
\begin{Highlighting}[]
\CommentTok{\# Frecuencia relativa acumulada}
\NormalTok{freq\_rel }\OtherTok{\textless{}{-}}\NormalTok{ frecuencia }\SpecialCharTok{/}\NormalTok{ n}
\NormalTok{freq\_acum }\OtherTok{\textless{}{-}} \FunctionTok{cumsum}\NormalTok{(freq\_rel)}

\CommentTok{\# Clase que contiene la mediana (F \textgreater{}= 0.5)}
\FunctionTok{print}\NormalTok{(}\FunctionTok{data.frame}\NormalTok{(}\AttributeTok{Clase =} \FunctionTok{c}\NormalTok{(}\StringTok{"70{-}90"}\NormalTok{, }\StringTok{"90{-}100"}\NormalTok{, }\StringTok{"100{-}110"}\NormalTok{, }\StringTok{"110{-}126"}\NormalTok{, }\StringTok{"126{-}150"}\NormalTok{),}
                 \AttributeTok{Frecuencia =}\NormalTok{ frecuencia,}
                 \AttributeTok{F\_acum =}\NormalTok{ freq\_acum))}
\end{Highlighting}
\end{Shaded}

}

\end{codelisting}%

\begin{verbatim}
    Clase Frecuencia F_acum
1   70-90         15   0.15
2  90-100         32   0.47
3 100-110         35   0.82
4 110-126         15   0.97
5 126-150          3   1.00
\end{verbatim}

\end{tcolorbox}

\begin{tcolorbox}[enhanced jigsaw, breakable, titlerule=0mm, rightrule=.15mm, leftrule=.75mm, left=2mm, opacitybacktitle=0.6, colback=white, bottomtitle=1mm, colbacktitle=quarto-callout-tip-color!10!white, title=\textcolor{quarto-callout-tip-color}{\faLightbulb}\hspace{0.5em}{Ejercicio 1.6: Interpretación de Histograma}, toptitle=1mm, coltitle=black, opacityback=0, arc=.35mm, bottomrule=.15mm, toprule=.15mm, colframe=quarto-callout-tip-color-frame]

Considera el histograma de IMC del Ejemplo 1.5.1:

\begin{enumerate}
\def\labelenumi{\alph{enumi})}
\tightlist
\item
  ¿Qué clase tiene la mayor densidad de pacientes?
\item
  ¿Qué porcentaje de pacientes tienen un IMC inferior a 25 (Normal o
  bajo peso)?
\item
  ¿En qué rango se encuentra el 25\% de pacientes con menor IMC?
\end{enumerate}

\end{tcolorbox}

\section{Respuestas a los Ejercicios}\label{sec-semana01-respuestas}

\textbf{Ejercicio 1.1: Conceptos} a) Variable: Característica
(Glucemia). Valor: Resultado medido (105 mg/dL). b) Población: Todos los
pacientes con la patología. Muestra: Los que participan en el estudio.
c) Un outlier (ej. una medición de 220/120 por error de manguito)
elevaría artificialmente la media, haciendo que la tensión promedio
parezca más alta de lo que realmente es.

\textbf{Ejercicio 1.2: Escalas} a) Discreta (conteo de eventos). b)
Ordinal (hay un orden de gravedad). c) Continua (puede tomar cualquier
valor decimal). d) Nominal (categorías sin orden). e) Ordinal
(puntuación con orden intrínseco).

\textbf{Ejercicio 1.3: Tiempos de espera} a) Media = 86 / 8 = 10.75 min.
b) Ordenados: 7, 8, 9, 10, 11, 12, 14, 15. Mediana = (10+11)/2 = 10.5
min. c) Ambas son similares porque no hay outliers extremos.

\textbf{Ejercicio 1.4: Admisiones} a)

\begin{longtable}[]{@{}lrrrr@{}}
\toprule\noalign{}
Admisiones & \(h_j\) & \(f_j\) & \(H_j\) & \(F_j\) \\
\midrule\noalign{}
\endhead
\bottomrule\noalign{}
\endlastfoot
2 & 6 & 0.30 & 6 & 0.30 \\
3 & 5 & 0.25 & 11 & 0.55 \\
4 & 4 & 0.20 & 15 & 0.75 \\
5 & 3 & 0.15 & 18 & 0.90 \\
6 & 2 & 0.10 & 20 & 1.00 \\
\textbf{Total} & \textbf{20} & \textbf{1.00} & & \\
\end{longtable}

\begin{enumerate}
\def\labelenumi{\alph{enumi})}
\setcounter{enumi}{1}
\tightlist
\item
  Moda = 2 admisiones (\(h_j=6\)). c) Proporción = \(F_3 = 0.55\)
  (55\%).
\end{enumerate}

\textbf{Ejercicio 1.5: Glucemia} a) Media ≈ 101.4 mg/dL. b) Clase modal:
100-110 mg/dL. c) P₅₀ ≈ 100 + (0.50 - 0.47)/0.35 * 10 ≈ 100.86 mg/dL.

\textbf{Ejercicio 1.6: IMC (Ejemplo 1.5.1)} a) Sobrepeso (25-30) con
densidad 0.0672. b) \(F(25) = 0.055 + 0.339 = 0.394\) (39.4\%). c) En el
rango de 0 a 25 (específicamente, el P₂₅ está en la clase ``Normal'').

\textbf{Fin de Semana 1}

Para referencias avanzadas sobre análisis exploratorio de datos y
estimación robusta, consulta: \href{https://rpubs.com/FJRubio/MM}{Robust
Estimation of Location}

\chapter{Semana 2 --- Análisis Exploratorio de Datos (Parte
II)}\label{sec-semana02}

En esta semana continuamos con el Análisis Exploratorio de Datos (AED),
enfocándonos en medidas de dispersión, transformaciones lineales,
diagramas de caja, análisis bivariante y medidas de dependencia. Estos
conceptos son fundamentales para comprender la variabilidad de los datos
y las relaciones entre variables.

\section{Medidas de Dispersión}\label{sec-semana02-s1}

Mientras que las medidas de tendencia central (como la media y la
mediana) describen dónde está el centro de los datos, las
\textbf{medidas de dispersión} describen cuánta variabilidad existe en
los datos. La dispersión es crítica para comprender la calidad de los
datos y la precisión de nuestras estimaciones.

\subsection{¿Por qué necesitamos medidas de
dispersión?}\label{por-quuxe9-necesitamos-medidas-de-dispersiuxf3n}

Es fundamental comprender que la media por sí sola no describe una
distribución. Dos conjuntos de datos pueden compartir la misma media
pero tener estructuras de dispersión radicalmente diferentes.

Consideremos un ejemplo clínico: dos grupos de pacientes con niveles de
glucosa en sangre (mg/dL) con la misma media (170 mg/dL): 1.
\textbf{Grupo A}: Niveles muy homogéneos (todos cercanos a la media). 2.
\textbf{Grupo B}: Niveles muy diversos (gran dispersión).

\begin{codelisting}

\caption{\label{lst-semana02-media-var}Distribuciones con igual media
pero distinta varianza}

\centering{

\begin{Shaded}
\begin{Highlighting}[]
\FunctionTok{library}\NormalTok{(ggplot2)}

\CommentTok{\# Generar datos}
\FunctionTok{set.seed}\NormalTok{(}\DecValTok{123}\NormalTok{)}
\NormalTok{grupo\_a }\OtherTok{\textless{}{-}} \FunctionTok{rnorm}\NormalTok{(}\DecValTok{100}\NormalTok{, }\AttributeTok{mean =} \DecValTok{170}\NormalTok{, }\AttributeTok{sd =} \DecValTok{2}\NormalTok{)   }\CommentTok{\# Varianza baja}
\NormalTok{grupo\_b }\OtherTok{\textless{}{-}} \FunctionTok{rnorm}\NormalTok{(}\DecValTok{100}\NormalTok{, }\AttributeTok{mean =} \DecValTok{170}\NormalTok{, }\AttributeTok{sd =} \DecValTok{10}\NormalTok{)  }\CommentTok{\# Varianza alta}

\NormalTok{df }\OtherTok{\textless{}{-}} \FunctionTok{data.frame}\NormalTok{(}
  \AttributeTok{valor =} \FunctionTok{c}\NormalTok{(grupo\_a, grupo\_b),}
  \AttributeTok{grupo =} \FunctionTok{rep}\NormalTok{(}\FunctionTok{c}\NormalTok{(}\StringTok{"Grupo A (Homogéneo)"}\NormalTok{, }\StringTok{"Grupo B (Diverso)"}\NormalTok{), }\AttributeTok{each =} \DecValTok{100}\NormalTok{)}
\NormalTok{)}

\CommentTok{\# Gráfico de densidad con líneas de media}
\FunctionTok{ggplot}\NormalTok{(df, }\FunctionTok{aes}\NormalTok{(}\AttributeTok{x =}\NormalTok{ valor, }\AttributeTok{fill =}\NormalTok{ grupo)) }\SpecialCharTok{+}
  \FunctionTok{geom\_density}\NormalTok{(}\AttributeTok{alpha =} \FloatTok{0.5}\NormalTok{) }\SpecialCharTok{+}
  \FunctionTok{scale\_fill\_manual}\NormalTok{(}\AttributeTok{values =} \FunctionTok{c}\NormalTok{(}\StringTok{"Grupo A (Homogéneo)"} \OtherTok{=} \StringTok{"steelblue"}\NormalTok{, }\StringTok{"Grupo B (Diverso)"} \OtherTok{=} \StringTok{"seagreen"}\NormalTok{)) }\SpecialCharTok{+}
  \FunctionTok{geom\_vline}\NormalTok{(}\AttributeTok{xintercept =} \DecValTok{170}\NormalTok{, }\AttributeTok{linetype =} \StringTok{"dashed"}\NormalTok{, }\AttributeTok{color =} \StringTok{"black"}\NormalTok{, }\AttributeTok{size =} \FloatTok{0.8}\NormalTok{) }\SpecialCharTok{+}
  \FunctionTok{labs}\NormalTok{(}\AttributeTok{title =} \StringTok{"Comparación de Variabilidad"}\NormalTok{,}
       \AttributeTok{subtitle =} \StringTok{"Ambas distribuciones tienen media = 170 mg/dL"}\NormalTok{,}
       \AttributeTok{x =} \StringTok{"Glucemia (mg/dL)"}\NormalTok{,}
       \AttributeTok{y =} \StringTok{"Densidad"}\NormalTok{) }\SpecialCharTok{+}
  \FunctionTok{theme\_minimal}\NormalTok{()}
\end{Highlighting}
\end{Shaded}

}

\end{codelisting}%

\pandocbounded{\includegraphics[keepaspectratio]{chapters/semana02_files/figure-pdf/unnamed-chunk-1-1.pdf}}

Como se observa en el gráfico (la línea discontinua representa la media
compartida), el \textbf{Grupo B} presenta una dispersión mucho mayor, lo
que demuestra que el promedio no siempre es suficiente para caracterizar
la realidad de un conjunto de datos médicos. :::

\begin{tcolorbox}[enhanced jigsaw, breakable, titlerule=0mm, rightrule=.15mm, leftrule=.75mm, left=2mm, opacitybacktitle=0.6, colback=white, bottomtitle=1mm, colbacktitle=quarto-callout-note-color!10!white, title=\textcolor{quarto-callout-note-color}{\faInfo}\hspace{0.5em}{Resumen: Medidas de Dispersión}, toptitle=1mm, coltitle=black, opacityback=0, arc=.35mm, bottomrule=.15mm, toprule=.15mm, colframe=quarto-callout-note-color-frame]

Para cuantificar la variabilidad de un conjunto de datos, utilizamos
principalmente las siguientes medidas:

\begin{itemize}
\tightlist
\item
  \textbf{Rango}: Diferencia entre el valor máximo y mínimo. Muy
  sensible a outliers.
\item
  \textbf{Rango Intercuartílico (IQR)}: Diferencia entre el tercer y
  primer cuartil. Mide la dispersión del 50\% central de los datos.
  Robusto a outliers.
\item
  \textbf{Varianza (\(s^2\))}: Promedio de las desviaciones al cuadrado
  respecto a la media.
\item
  \textbf{Desviación Estándar (\(s\))}: Raíz cuadrada de la varianza.
  Expresada en las mismas unidades que la variable original, lo que
  facilita su interpretación clínica.
\item
  \textbf{Coeficiente de Variación (CV)}: Medida adimensional de
  dispersión relativa que permite comparar la variabilidad de diferentes
  variables clínicas.
\end{itemize}

\end{tcolorbox}

\subsection{Rango}\label{sec-semana02-s1-rango}

\begin{tcolorbox}[enhanced jigsaw, breakable, titlerule=0mm, rightrule=.15mm, leftrule=.75mm, left=2mm, opacitybacktitle=0.6, colback=white, bottomtitle=1mm, colbacktitle=quarto-callout-note-color!10!white, title=\textcolor{quarto-callout-note-color}{\faInfo}\hspace{0.5em}{Rango}, toptitle=1mm, coltitle=black, opacityback=0, arc=.35mm, bottomrule=.15mm, toprule=.15mm, colframe=quarto-callout-note-color-frame]

El \textbf{rango} es la diferencia entre la observación más grande y la
más pequeña:

\[R = x_{\max} - x_{\min} = x_{(n)} - x_{(1)}\]

donde \(x_{(1)}, \ldots, x_{(n)}\) son las observaciones ordenadas.

\end{tcolorbox}

\begin{tcolorbox}[enhanced jigsaw, breakable, titlerule=0mm, rightrule=.15mm, leftrule=.75mm, left=2mm, opacitybacktitle=0.6, colback=white, bottomtitle=1mm, colbacktitle=quarto-callout-warning-color!10!white, title=\textcolor{quarto-callout-warning-color}{\faExclamationTriangle}\hspace{0.5em}{Limitaciones del Rango}, toptitle=1mm, coltitle=black, opacityback=0, arc=.35mm, bottomrule=.15mm, toprule=.15mm, colframe=quarto-callout-warning-color-frame]

El rango es fácil de calcular, pero tiene importantes limitaciones: - Es
muy sensible a valores atípicos (outliers) - Su valor aumenta con el
tamaño de la muestra - No proporciona información sobre cómo se
distribuyen los datos entre los extremos

\end{tcolorbox}

\subsection{Rango Intercuartílico (IQR)}\label{sec-semana02-s1-iqr}

\begin{tcolorbox}[enhanced jigsaw, breakable, titlerule=0mm, rightrule=.15mm, leftrule=.75mm, left=2mm, opacitybacktitle=0.6, colback=white, bottomtitle=1mm, colbacktitle=quarto-callout-note-color!10!white, title=\textcolor{quarto-callout-note-color}{\faInfo}\hspace{0.5em}{Rango Intercuartílico}, toptitle=1mm, coltitle=black, opacityback=0, arc=.35mm, bottomrule=.15mm, toprule=.15mm, colframe=quarto-callout-note-color-frame]

El \textbf{rango intercuartílico} (IQR, del inglés Interquartile Range)
es la diferencia entre el tercer cuartil (\(Q3\)) y el primer cuartil
(\(Q1\)):

\[\text{IQR} = Q3 - Q1\]

El IQR contiene el 50\% central de los datos y es robusto a valores
atípicos.

\end{tcolorbox}

\subsection{Varianza y Desviación Estándar}\label{sec-semana02-s1-var}

\begin{tcolorbox}[enhanced jigsaw, breakable, titlerule=0mm, rightrule=.15mm, leftrule=.75mm, left=2mm, opacitybacktitle=0.6, colback=white, bottomtitle=1mm, colbacktitle=quarto-callout-note-color!10!white, title=\textcolor{quarto-callout-note-color}{\faInfo}\hspace{0.5em}{Varianza y Desviación Estándar}, toptitle=1mm, coltitle=black, opacityback=0, arc=.35mm, bottomrule=.15mm, toprule=.15mm, colframe=quarto-callout-note-color-frame]

La \textbf{varianza muestral} mide la dispersión promedio de los datos
respecto a la media:

\[s^2 = \frac{1}{n-1}\sum_{i=1}^{n}(x_i - \bar{x})^2\]

Alternativamente, puede calcularse como:

\[s^2 = \frac{1}{n-1}\left(\sum_{i=1}^{n}x_i^2 - n\bar{x}^2\right)\]

La \textbf{desviación estándar} es la raíz cuadrada positiva de la
varianza:

\[s = \sqrt{s^2}\]

Nota: Se utiliza \(n-1\) (en lugar de \(n\)) para proporcionar un
estimador insesgado de la varianza poblacional.

\end{tcolorbox}

\begin{tcolorbox}[enhanced jigsaw, breakable, titlerule=0mm, rightrule=.15mm, leftrule=.75mm, left=2mm, opacitybacktitle=0.6, colback=white, bottomtitle=1mm, colbacktitle=quarto-callout-tip-color!10!white, title=\textcolor{quarto-callout-tip-color}{\faLightbulb}\hspace{0.5em}{Interpretación de la Desviación Estándar}, toptitle=1mm, coltitle=black, opacityback=0, arc=.35mm, bottomrule=.15mm, toprule=.15mm, colframe=quarto-callout-tip-color-frame]

La desviación estándar está en las mismas unidades que los datos
originales, lo que la hace más interpretable que la varianza. Por
ejemplo, si los datos representan edades en años, \(s\) también está en
años. Aproximadamente el 68\% de los datos caen dentro de una desviación
estándar de la media, y el 95\% dentro de dos desviaciones estándar
(regla empírica para datos aproximadamente normales).

\end{tcolorbox}

\begin{tcolorbox}[enhanced jigsaw, breakable, titlerule=0mm, rightrule=.15mm, leftrule=.75mm, left=2mm, opacitybacktitle=0.6, colback=white, bottomtitle=1mm, colbacktitle=quarto-callout-tip-color!10!white, title=\textcolor{quarto-callout-tip-color}{\faLightbulb}\hspace{0.5em}{Ejemplo 2.1: BioStatR grps()}, toptitle=1mm, coltitle=black, opacityback=0, arc=.35mm, bottomrule=.15mm, toprule=.15mm, colframe=quarto-callout-tip-color-frame]

La función \texttt{grps()} calcula estadísticos descriptivos
comparativos por grupos en una sola llamada (documentación completa en
Sección~\ref{sec-bio-descriptivo}):

\begin{codelisting}[H]

\caption{\label{lst-semana02-1}Código R}

\centering{

\begin{Shaded}
\begin{Highlighting}[]
\FunctionTok{library}\NormalTok{(BioStatR)}

\CommentTok{\# IMC según sexo — pacientes diabéticos UGR (n = 94)}
\FunctionTok{grps}\NormalTok{(osteo}\SpecialCharTok{$}\NormalTok{imc, osteo}\SpecialCharTok{$}\NormalTok{sexo, }\AttributeTok{grf =} \ConstantTok{FALSE}\NormalTok{)}
\end{Highlighting}
\end{Shaded}

}

\end{codelisting}%

\begin{verbatim}

# Descriptiva de osteo$imc por osteo$sexo
# ---------------------------------------
 osteo$sexo 
         n  media    dt
 Hombre 45 23.514 2.890
 Mujer  49 24.294 4.389
 Total  94 23.921 3.748
\end{verbatim}

Las mujeres presentan mayor IMC medio y mayor variabilidad (dt = 4.39 vs
2.89 en hombres). Con \texttt{ic\ =\ TRUE} añade intervalos de confianza
individuales por grupo.

\end{tcolorbox}

\subsection{Cálculo de Estadísticos con R
Base}\label{sec-semana02-s1-base-r}

Aunque paquetes como \texttt{BioStatR} son potentes, es fundamental
saber usar las funciones básicas de R para un análisis rápido.

\begin{codelisting}

\caption{\label{lst-semana02-2}Cálculos con Base R}

\centering{

\begin{Shaded}
\begin{Highlighting}[]
\CommentTok{\# Datos: PAS de pacientes (mmHg)}
\NormalTok{pas }\OtherTok{\textless{}{-}} \FunctionTok{c}\NormalTok{(}\DecValTok{120}\NormalTok{, }\DecValTok{125}\NormalTok{, }\DecValTok{118}\NormalTok{, }\DecValTok{130}\NormalTok{, }\DecValTok{122}\NormalTok{, }\DecValTok{128}\NormalTok{, }\DecValTok{115}\NormalTok{, }\DecValTok{135}\NormalTok{, }\DecValTok{124}\NormalTok{, }\DecValTok{129}\NormalTok{)}

\CommentTok{\# Cálculo de varianza y desviación estándar}
\NormalTok{varianza }\OtherTok{\textless{}{-}} \FunctionTok{var}\NormalTok{(pas)}
\NormalTok{desviacion\_std }\OtherTok{\textless{}{-}} \FunctionTok{sd}\NormalTok{(pas)}
\NormalTok{iqr\_pas }\OtherTok{\textless{}{-}} \FunctionTok{IQR}\NormalTok{(pas)}

\CommentTok{\# Cálculo del Coeficiente de Variación (CV)}
\CommentTok{\# El CV es (sd / media) * 100}
\NormalTok{cv }\OtherTok{\textless{}{-}}\NormalTok{ (desviacion\_std }\SpecialCharTok{/} \FunctionTok{mean}\NormalTok{(pas)) }\SpecialCharTok{*} \DecValTok{100}

\FunctionTok{cat}\NormalTok{(}\StringTok{"Varianza (s²):"}\NormalTok{, }\FunctionTok{round}\NormalTok{(varianza, }\DecValTok{2}\NormalTok{), }\StringTok{"}\SpecialCharTok{\textbackslash{}n}\StringTok{"}\NormalTok{)}
\end{Highlighting}
\end{Shaded}

}

\end{codelisting}%

\begin{verbatim}
Varianza (s²): 36.93 
\end{verbatim}

\begin{codelisting}

\caption{\label{lst-semana02-2}Cálculos con Base R}

\centering{

\begin{Shaded}
\begin{Highlighting}[]
\FunctionTok{cat}\NormalTok{(}\StringTok{"Desv. Est. (s):"}\NormalTok{, }\FunctionTok{round}\NormalTok{(desviacion\_std, }\DecValTok{2}\NormalTok{), }\StringTok{"}\SpecialCharTok{\textbackslash{}n}\StringTok{"}\NormalTok{)}
\end{Highlighting}
\end{Shaded}

}

\end{codelisting}%

\begin{verbatim}
Desv. Est. (s): 6.08 
\end{verbatim}

\begin{codelisting}

\caption{\label{lst-semana02-2}Cálculos con Base R}

\centering{

\begin{Shaded}
\begin{Highlighting}[]
\FunctionTok{cat}\NormalTok{(}\StringTok{"IQR:"}\NormalTok{, }\FunctionTok{round}\NormalTok{(iqr\_pas, }\DecValTok{2}\NormalTok{), }\StringTok{"}\SpecialCharTok{\textbackslash{}n}\StringTok{"}\NormalTok{)}
\end{Highlighting}
\end{Shaded}

}

\end{codelisting}%

\begin{verbatim}
IQR: 8.25 
\end{verbatim}

\begin{codelisting}

\caption{\label{lst-semana02-2}Cálculos con Base R}

\centering{

\begin{Shaded}
\begin{Highlighting}[]
\FunctionTok{cat}\NormalTok{(}\StringTok{"Coef. Variación (CV):"}\NormalTok{, }\FunctionTok{round}\NormalTok{(cv, }\DecValTok{2}\NormalTok{), }\StringTok{"\%}\SpecialCharTok{\textbackslash{}n}\StringTok{"}\NormalTok{)}
\end{Highlighting}
\end{Shaded}

}

\end{codelisting}%

\begin{verbatim}
Coef. Variación (CV): 4.88 %
\end{verbatim}

\begin{tcolorbox}[enhanced jigsaw, breakable, titlerule=0mm, rightrule=.15mm, leftrule=.75mm, left=2mm, opacitybacktitle=0.6, colback=white, bottomtitle=1mm, colbacktitle=quarto-callout-note-color!10!white, title=\textcolor{quarto-callout-note-color}{\faInfo}\hspace{0.5em}{Coeficiente de Variación (CV) en Contexto Médico}, toptitle=1mm, coltitle=black, opacityback=0, arc=.35mm, bottomrule=.15mm, toprule=.15mm, colframe=quarto-callout-note-color-frame]

El \textbf{Coeficiente de Variación (CV)} es una medida de dispersión
\textbf{adimensional} que permite comparar la variabilidad de variables
con diferentes unidades o magnitudes. Se define como:

\[CV = \left( \frac{s}{\bar{x}} \right) \times 100\]

En medicina, es crucial: - \textbf{Interpretación}: Un CV bajo indica
alta precisión o consistencia en la medición. - \textbf{Comparación}:
Por ejemplo, si comparamos la variabilidad de la glucosa en sangre
(mg/dL) con la variabilidad de la frecuencia cardíaca (latidos/min), el
CV nos permite saber cuál medida es más ``relativamente'' dispersa,
ignorando las unidades de medida.

\end{tcolorbox}

\section{Transformaciones Lineales: Propiedades}\label{sec-semana02-s2}

Cuando transformamos los datos mediante una transformación lineal,
también transformamos sus estadísticos de una manera predecible y
matemáticamente bien definida. Generalmente se usan para simplificar la
interpretación de datos, cambiar unidades de medida y facilitar la
comparación entre variables.

\subsection{Transformación Lineal
Estándar}\label{transformaciuxf3n-lineal-estuxe1ndar}

\begin{tcolorbox}[enhanced jigsaw, breakable, titlerule=0mm, rightrule=.15mm, leftrule=.75mm, left=2mm, opacitybacktitle=0.6, colback=white, bottomtitle=1mm, colbacktitle=quarto-callout-warning-color!10!white, title=\textcolor{quarto-callout-warning-color}{\faExclamationTriangle}\hspace{0.5em}{Resultado Importante: Transformación Lineal}, toptitle=1mm, coltitle=black, opacityback=0, arc=.35mm, bottomrule=.15mm, toprule=.15mm, colframe=quarto-callout-warning-color-frame]

Si definimos una transformación lineal \(y_i = a + b x_i\) donde
\(b \neq 0\), entonces:

\textbf{Para la media:} \[\bar{y} = a + b\bar{x}\]

\textbf{Para la varianza:} \[s_y^2 = b^2 s_x^2\]

\textbf{Para la desviación estándar:} \[s_y = |b| s_x\]

\textbf{Para el rango:} \[R_Y = |b|R_X\]

\textbf{Para el rango intercuartílico:}
\[\text{IQR}_Y = |b| \text{IQR}_X\]

\end{tcolorbox}

\subsection{Interpretación de los
Parámetros}\label{sec-semana02-s2-interp}

El parámetro \(a\) representa un \textbf{desplazamiento} de los datos
(shift), mientras que el parámetro \(b\) controla la escala:

\begin{itemize}
\tightlist
\item
  Si \(0 < b < 1\): compresión de los datos (se hacen más próximos)
\item
  Si \(b > 1\): dilatación de los datos (se separan más)
\item
  Si \(b < 0\): reflexión (inversión) además de escalado
\item
  Si \(b = 1\): solo desplazamiento (la dispersión no cambia)
\end{itemize}

\begin{tcolorbox}[enhanced jigsaw, breakable, titlerule=0mm, rightrule=.15mm, leftrule=.75mm, left=2mm, opacitybacktitle=0.6, colback=white, bottomtitle=1mm, colbacktitle=quarto-callout-tip-color!10!white, title=\textcolor{quarto-callout-tip-color}{\faLightbulb}\hspace{0.5em}{Ejemplo 2.2: Conversión de Temperaturas}, toptitle=1mm, coltitle=black, opacityback=0, arc=.35mm, bottomrule=.15mm, toprule=.15mm, colframe=quarto-callout-tip-color-frame]

Supongamos que tenemos temperaturas en grados Celsius con media
\(\bar{x} = 25\) y desviación estándar \(s_x = 3\). Queremos
convertirlas a Fahrenheit usando \(y = 32 + 1.8x\).

Entonces: - \(\bar{y} = 32 + 1.8 \times 25 = 77°\text{F}\) -
\(s_y = |1.8| \times 3 = 5.4°\text{F}\)

La media se desplaza y escala, pero la dispersión relativa (en términos
de desviaciones estándar) permanece igual.

\end{tcolorbox}

\section{Estandarización y Normalización}\label{sec-semana02-s3}

La estandarización es una transformación lineal específica que convierte
los datos a una escala con media 0 y desviación estándar 1. Esto se
logra mediante la fórmula:

\[z_i = \frac{x_i - \bar{x}}{s_x}\] Desarrollando la expresión:

\[z_i = \frac{x_i}{s_x} - \frac{\bar{x}}{s_x}\]

Por lo tanto, en la transformación lineal \(y_i = a + b x_i\) para el
z-scocre tenemos:

\begin{itemize}
\tightlist
\item
  \[a = -\frac{\bar{x}}{s_x}\]
\item
  \[b = \frac{1}{s_x}\]
\end{itemize}

\subsection{Propiedades de los Z-Scores}\label{sec-semana02-s3-prop}

\begin{tcolorbox}[enhanced jigsaw, breakable, titlerule=0mm, rightrule=.15mm, leftrule=.75mm, left=2mm, opacitybacktitle=0.6, colback=white, bottomtitle=1mm, colbacktitle=quarto-callout-warning-color!10!white, title=\textcolor{quarto-callout-warning-color}{\faExclamationTriangle}\hspace{0.5em}{Propiedades de los Z-Scores}, toptitle=1mm, coltitle=black, opacityback=0, arc=.35mm, bottomrule=.15mm, toprule=.15mm, colframe=quarto-callout-warning-color-frame]

\begin{itemize}
\tightlist
\item
  Media: \(\bar{z} = 0\)
\item
  Desviación estándar: \(s_z = 1\)
\item
  Rango típico: aproximadamente \([-3, 3]\) para datos normales
\item
  Un valor \(z_i > 3\) o \(z_i < -3\) sugiere un posible valor atípico
\end{itemize}

\end{tcolorbox}

\subsection{Interpretación}\label{sec-semana02-s3-interp}

\begin{tcolorbox}[enhanced jigsaw, breakable, titlerule=0mm, rightrule=.15mm, leftrule=.75mm, left=2mm, opacitybacktitle=0.6, colback=white, bottomtitle=1mm, colbacktitle=quarto-callout-tip-color!10!white, title=\textcolor{quarto-callout-tip-color}{\faLightbulb}\hspace{0.5em}{Ejemplo 2.3: Interpretación de Z-Scores}, toptitle=1mm, coltitle=black, opacityback=0, arc=.35mm, bottomrule=.15mm, toprule=.15mm, colframe=quarto-callout-tip-color-frame]

Si un estudiante obtiene una puntuación de 85 en un examen con media 80
y desviación estándar 5, su z-score es:

\[z = \frac{85 - 80}{5} = 1\]

Esto significa que su puntuación está 1 desviación estándar \textbf{por
encima} de la media. Si otro estudiante obtiene 70, su z-score sería:

\[z = \frac{70 - 80}{5} = -1\]

Los z-scores permiten comparar puntuaciones de diferentes escalas o
distribuciones.

\end{tcolorbox}

\begin{tcolorbox}[enhanced jigsaw, breakable, titlerule=0mm, rightrule=.15mm, leftrule=.75mm, left=2mm, opacitybacktitle=0.6, colback=white, bottomtitle=1mm, colbacktitle=quarto-callout-note-color!10!white, title=\textcolor{quarto-callout-note-color}{\faInfo}\hspace{0.5em}{Estandarización automática en R: \texttt{scale()}}, toptitle=1mm, coltitle=black, opacityback=0, arc=.35mm, bottomrule=.15mm, toprule=.15mm, colframe=quarto-callout-note-color-frame]

Aunque podemos calcular z-scores manualmente, R ofrece la función
\texttt{scale()} que estandariza automáticamente un vector:

\begin{Shaded}
\begin{Highlighting}[]
\CommentTok{\# Ejemplo numérico}
\NormalTok{presion }\OtherTok{\textless{}{-}} \FunctionTok{c}\NormalTok{(}\DecValTok{120}\NormalTok{, }\DecValTok{125}\NormalTok{, }\DecValTok{110}\NormalTok{, }\DecValTok{130}\NormalTok{, }\DecValTok{140}\NormalTok{)}

\CommentTok{\# Estandarización: (x {-} mean(x)) / sd(x)}
\NormalTok{presion\_z }\OtherTok{\textless{}{-}} \FunctionTok{scale}\NormalTok{(presion)}

\CommentTok{\# Resultados}
\FunctionTok{print}\NormalTok{(presion)}
\end{Highlighting}
\end{Shaded}

\begin{verbatim}
[1] 120 125 110 130 140
\end{verbatim}

\begin{Shaded}
\begin{Highlighting}[]
\FunctionTok{print}\NormalTok{(}\FunctionTok{round}\NormalTok{(presion\_z, }\DecValTok{2}\NormalTok{))}
\end{Highlighting}
\end{Shaded}

\begin{verbatim}
      [,1]
[1,] -0.45
[2,]  0.00
[3,] -1.34
[4,]  0.45
[5,]  1.34
attr(,"scaled:center")
[1] 125
attr(,"scaled:scale")
[1] 11.18034
\end{verbatim}

La función \texttt{scale()} resta la media
(\texttt{mean(presion)\ =\ 125}) y divide por la desviación estándar
(\texttt{sd(presion)\ ≈\ 11.18}), transformando directamente los datos a
z-scores.

\end{tcolorbox}

\chapter{No es Ejemplo}\label{sec-semana02-s3-ejemplo-pob}

Cuando conocemos la media poblacional (\(\mu\)) y la desviación típica
poblacional (\(\sigma\)), podemos calcular la probabilidad (porcentaje
de personas) de que un valor caiga en un rango utilizando la función
\texttt{pnorm()}.

\begin{tcolorbox}[enhanced jigsaw, breakable, titlerule=0mm, rightrule=.15mm, leftrule=.75mm, left=2mm, opacitybacktitle=0.6, colback=white, bottomtitle=1mm, colbacktitle=quarto-callout-tip-color!10!white, title=\textcolor{quarto-callout-tip-color}{\faLightbulb}\hspace{0.5em}{Ejemplo 2.4: Glucemia (Z-scores con Parámetros Poblacionales)}, toptitle=1mm, coltitle=black, opacityback=0, arc=.35mm, bottomrule=.15mm, toprule=.15mm, colframe=quarto-callout-tip-color-frame]

Supongamos que la glucemia en adultos sanos sigue una distribución
normal con \(\mu = 90\) mg/dL y \(\sigma = 10\) mg/dL. ¿Qué porcentaje
de personas tiene una glucemia entre 80 y 100 mg/dL?

\textbf{Cálculo manual mediante Z-score:} 1. Para \(x = 80\):
\(z_1 = \frac{80 - 90}{10} = -1.0\) 2. Para \(x = 100\):
\(z_2 = \frac{100 - 90}{10} = +1.0\)

Buscamos \(P(-1.0 < Z < 1.0)\). Esto es
\(P(Z < 1.0) - P(Z < -1.0) \approx 0.8413 - 0.1587 = 0.6826\) (68.26\%).
Esto significa que aproximadamente el 68.26\% de las personas sanas
tienen una glucemia entre 80 y 100 mg/dL. Antes los valores de z se
buscan en la tabla de la distribución normal estándar, pero con R es
mucho más sencillo usando \texttt{pnorm()}.

\textbf{Verificación con R:}

\begin{codelisting}[H]

\caption{\label{lst-semana02-3}Código R}

\centering{

\begin{Shaded}
\begin{Highlighting}[]
\NormalTok{mu }\OtherTok{\textless{}{-}} \DecValTok{90}
\NormalTok{sigma }\OtherTok{\textless{}{-}} \DecValTok{10}

\CommentTok{\# Probabilidad de P(80 \textless{} X \textless{} 100)}
\CommentTok{\# Esto es P(X \textless{} 100) {-} P(X \textless{} 80)}
\NormalTok{prob }\OtherTok{\textless{}{-}} \FunctionTok{pnorm}\NormalTok{(}\DecValTok{100}\NormalTok{, }\AttributeTok{mean =}\NormalTok{ mu, }\AttributeTok{sd =}\NormalTok{ sigma) }\SpecialCharTok{{-}} \FunctionTok{pnorm}\NormalTok{(}\DecValTok{80}\NormalTok{, }\AttributeTok{mean =}\NormalTok{ mu, }\AttributeTok{sd =}\NormalTok{ sigma)}

\FunctionTok{cat}\NormalTok{(}\StringTok{"Porcentaje de personas entre 80 y 100 mg/dL:"}\NormalTok{, }\FunctionTok{round}\NormalTok{(prob }\SpecialCharTok{*} \DecValTok{100}\NormalTok{, }\DecValTok{2}\NormalTok{), }\StringTok{"\%}\SpecialCharTok{\textbackslash{}n}\StringTok{"}\NormalTok{)}
\end{Highlighting}
\end{Shaded}

}

\end{codelisting}%

\begin{verbatim}
Porcentaje de personas entre 80 y 100 mg/dL: 68.27 %
\end{verbatim}

\end{tcolorbox}

\subsection{Regla Empírica (68-95-99.7)}\label{sec-semana02-s3-regla}

\begin{tcolorbox}[enhanced jigsaw, breakable, titlerule=0mm, rightrule=.15mm, leftrule=.75mm, left=2mm, opacitybacktitle=0.6, colback=white, bottomtitle=1mm, colbacktitle=quarto-callout-note-color!10!white, title=\textcolor{quarto-callout-note-color}{\faInfo}\hspace{0.5em}{Definición: Regla 68-95-99.7}, toptitle=1mm, coltitle=black, opacityback=0, arc=.35mm, bottomrule=.15mm, toprule=.15mm, colframe=quarto-callout-note-color-frame]

Cando una variable aleatoria X se transforma a puntuaciones estándar
mediante: so z=-score, se dice que hay sido tipificada y se ha
convertido en una variable aleatoria Z que sigue una distribución normal
estándar N(0,1). La regla empírica (o regla 68-95-99.7) se interpreta
directamente en términos de la distribución normal estándar N(0,1):

\begin{itemize}
\tightlist
\item
  Aproximadamente el \textbf{68\%} de los datos caen dentro de \(\pm 1\)
  desviación estándar (\(\pm 1 Z\))
\item
  Aproximadamente el \textbf{95\%} de los datos caen dentro de \(\pm 2\)
  desviaciones estándar (\(\pm 2 Z\))
\item
  Aproximadamente el \textbf{99.7\%} de los datos caen dentro de
  \(\pm 3\) desviaciones estándar (\(\pm 3 Z\))
\end{itemize}

En otras palabras, la tipificación ``estandariza'' cualquier normal a la
misma escala, y esas proporciones se mantienen porque dependen
únicamente de la forma de la curva normal estándar, no de la media o la
desviación original.

\end{tcolorbox}

\begin{codelisting}

\caption{\label{lst-semana02-4}Visualización de la Regla Empírica}

\centering{

\begin{Shaded}
\begin{Highlighting}[]
\CommentTok{\# Crear secuencia para curva normal}
\NormalTok{x }\OtherTok{\textless{}{-}} \FunctionTok{seq}\NormalTok{(}\SpecialCharTok{{-}}\DecValTok{4}\NormalTok{, }\DecValTok{4}\NormalTok{, }\AttributeTok{length.out =} \DecValTok{200}\NormalTok{)}
\NormalTok{y }\OtherTok{\textless{}{-}} \FunctionTok{dnorm}\NormalTok{(x)}

\CommentTok{\# Graficar}
\FunctionTok{plot}\NormalTok{(x, y, }\AttributeTok{type =} \StringTok{"l"}\NormalTok{, }\AttributeTok{lwd =} \DecValTok{2}\NormalTok{, }\AttributeTok{col =} \StringTok{"darkblue"}\NormalTok{,}
     \AttributeTok{main =} \StringTok{"Distribución Normal: Regla 68{-}95{-}99.7"}\NormalTok{,}
     \AttributeTok{xlab =} \StringTok{"Z{-}score"}\NormalTok{, }\AttributeTok{ylab =} \StringTok{"Densidad"}\NormalTok{, }\AttributeTok{axes =} \ConstantTok{FALSE}\NormalTok{)}
\FunctionTok{axis}\NormalTok{(}\DecValTok{1}\NormalTok{, }\AttributeTok{at =} \SpecialCharTok{{-}}\DecValTok{3}\SpecialCharTok{:}\DecValTok{3}\NormalTok{)}
\FunctionTok{axis}\NormalTok{(}\DecValTok{2}\NormalTok{)}

\CommentTok{\# Definir áreas y añadir etiquetas}
\CommentTok{\# 99.7\%}
\NormalTok{x\_99 }\OtherTok{\textless{}{-}} \FunctionTok{seq}\NormalTok{(}\SpecialCharTok{{-}}\DecValTok{3}\NormalTok{, }\DecValTok{3}\NormalTok{, }\AttributeTok{length.out=}\DecValTok{150}\NormalTok{)}
\FunctionTok{polygon}\NormalTok{(}\FunctionTok{c}\NormalTok{(}\SpecialCharTok{{-}}\DecValTok{3}\NormalTok{, x\_99, }\DecValTok{3}\NormalTok{), }\FunctionTok{c}\NormalTok{(}\DecValTok{0}\NormalTok{, }\FunctionTok{dnorm}\NormalTok{(x\_99), }\DecValTok{0}\NormalTok{), }\AttributeTok{col =} \FunctionTok{rgb}\NormalTok{(}\DecValTok{0}\NormalTok{,}\DecValTok{0}\NormalTok{,}\DecValTok{1}\NormalTok{,}\FloatTok{0.05}\NormalTok{), }\AttributeTok{border =} \StringTok{"darkblue"}\NormalTok{)}
\FunctionTok{text}\NormalTok{(}\DecValTok{0}\NormalTok{, }\FloatTok{0.02}\NormalTok{, }\StringTok{"99.7\% (±3σ)"}\NormalTok{, }\AttributeTok{col =} \StringTok{"darkblue"}\NormalTok{, }\AttributeTok{font =} \DecValTok{2}\NormalTok{)}

\CommentTok{\# 95\%}
\NormalTok{x\_95 }\OtherTok{\textless{}{-}} \FunctionTok{seq}\NormalTok{(}\SpecialCharTok{{-}}\DecValTok{2}\NormalTok{, }\DecValTok{2}\NormalTok{, }\AttributeTok{length.out=}\DecValTok{100}\NormalTok{)}
\FunctionTok{polygon}\NormalTok{(}\FunctionTok{c}\NormalTok{(}\SpecialCharTok{{-}}\DecValTok{2}\NormalTok{, x\_95, }\DecValTok{2}\NormalTok{), }\FunctionTok{c}\NormalTok{(}\DecValTok{0}\NormalTok{, }\FunctionTok{dnorm}\NormalTok{(x\_95), }\DecValTok{0}\NormalTok{), }\AttributeTok{col =} \FunctionTok{rgb}\NormalTok{(}\DecValTok{0}\NormalTok{,}\DecValTok{0}\NormalTok{,}\DecValTok{1}\NormalTok{,}\FloatTok{0.1}\NormalTok{), }\AttributeTok{border =} \StringTok{"darkblue"}\NormalTok{)}
\FunctionTok{text}\NormalTok{(}\DecValTok{0}\NormalTok{, }\FloatTok{0.08}\NormalTok{, }\StringTok{"95\% (±2σ)"}\NormalTok{, }\AttributeTok{col =} \StringTok{"darkblue"}\NormalTok{, }\AttributeTok{font =} \DecValTok{2}\NormalTok{)}

\CommentTok{\# 68\%}
\NormalTok{x\_68 }\OtherTok{\textless{}{-}} \FunctionTok{seq}\NormalTok{(}\SpecialCharTok{{-}}\DecValTok{1}\NormalTok{, }\DecValTok{1}\NormalTok{, }\AttributeTok{length.out=}\DecValTok{50}\NormalTok{)}
\FunctionTok{polygon}\NormalTok{(}\FunctionTok{c}\NormalTok{(}\SpecialCharTok{{-}}\DecValTok{1}\NormalTok{, x\_68, }\DecValTok{1}\NormalTok{), }\FunctionTok{c}\NormalTok{(}\DecValTok{0}\NormalTok{, }\FunctionTok{dnorm}\NormalTok{(x\_68), }\DecValTok{0}\NormalTok{), }\AttributeTok{col =} \FunctionTok{rgb}\NormalTok{(}\DecValTok{0}\NormalTok{,}\DecValTok{0}\NormalTok{,}\DecValTok{1}\NormalTok{,}\FloatTok{0.2}\NormalTok{), }\AttributeTok{border =} \StringTok{"darkblue"}\NormalTok{)}
\FunctionTok{text}\NormalTok{(}\DecValTok{0}\NormalTok{, }\FloatTok{0.25}\NormalTok{, }\StringTok{"68\% (±1σ)"}\NormalTok{, }\AttributeTok{col =} \StringTok{"darkblue"}\NormalTok{, }\AttributeTok{font =} \DecValTok{2}\NormalTok{)}
\end{Highlighting}
\end{Shaded}

}

\end{codelisting}%

\pandocbounded{\includegraphics[keepaspectratio]{chapters/semana02_files/figure-pdf/unnamed-chunk-6-1.pdf}}

\section{Diagramas de Caja (Boxplots)}\label{sec-semana02-s4}

Los diagramas de caja son una herramienta visual extremadamente útil en
el AED. Permiten identificar rápidamente la mediana, los cuartiles, el
rango y los valores atípicos.

\subsection{Construcción de un
Boxplot}\label{sec-semana02-s4-construccion}

\begin{tcolorbox}[enhanced jigsaw, breakable, titlerule=0mm, rightrule=.15mm, leftrule=.75mm, left=2mm, opacitybacktitle=0.6, colback=white, bottomtitle=1mm, colbacktitle=quarto-callout-note-color!10!white, title=\textcolor{quarto-callout-note-color}{\faInfo}\hspace{0.5em}{Ingredientes del Boxplot}, toptitle=1mm, coltitle=black, opacityback=0, arc=.35mm, bottomrule=.15mm, toprule=.15mm, colframe=quarto-callout-note-color-frame]

Para construir un diagrama de caja, necesitamos calcular:

\begin{enumerate}
\def\labelenumi{\arabic{enumi}.}
\tightlist
\item
  \textbf{Mediana (Q2):} el valor central que divide los datos en dos
  mitades
\item
  \textbf{Primer cuartil (Q1):} el valor que deja el 25\% de los datos
  por debajo
\item
  \textbf{Tercer cuartil (Q3):} el valor que deja el 75\% de los datos
  por debajo
\item
  \textbf{Rango intercuartílico:} \(\text{IQR} = Q3 - Q1\)
\item
  \textbf{Límite inferior (lower whisker):}
  \(Q1 - 1.5 \times \text{IQR}\)
\item
  \textbf{Límite superior (upper whisker):}
  \(Q3 + 1.5 \times \text{IQR}\)
\item
  \textbf{Valores atípicos (outliers):} cualquier observación fuera de
  los límites
\end{enumerate}

\end{tcolorbox}

\subsection{Interpretación Visual}\label{sec-semana02-s4-interp}

\begin{tcolorbox}[enhanced jigsaw, breakable, titlerule=0mm, rightrule=.15mm, leftrule=.75mm, left=2mm, opacitybacktitle=0.6, colback=white, bottomtitle=1mm, colbacktitle=quarto-callout-warning-color!10!white, title=\textcolor{quarto-callout-warning-color}{\faExclamationTriangle}\hspace{0.5em}{Características Importantes}, toptitle=1mm, coltitle=black, opacityback=0, arc=.35mm, bottomrule=.15mm, toprule=.15mm, colframe=quarto-callout-warning-color-frame]

\begin{itemize}
\tightlist
\item
  La \textbf{caja} representa el 50\% central de los datos (entre Q1 y
  Q3)
\item
  La \textbf{línea dentro de la caja} es la mediana (Q2)
\item
  Los \textbf{bigotes} (whiskers) se extienden hasta los límites
  definidos
\item
  Los \textbf{puntos individuales} fuera de los bigotes son valores
  atípicos
\item
  Un boxplot \textbf{simétrico} sugiere una distribución aproximadamente
  normal
\item
  Una mediana \textbf{no centrada} en la caja sugiere asimetría
\end{itemize}

\end{tcolorbox}

\begin{tcolorbox}[enhanced jigsaw, breakable, titlerule=0mm, rightrule=.15mm, leftrule=.75mm, left=2mm, opacitybacktitle=0.6, colback=white, bottomtitle=1mm, colbacktitle=quarto-callout-tip-color!10!white, title=\textcolor{quarto-callout-tip-color}{\faLightbulb}\hspace{0.5em}{Ejemplo 2.5: Boxplot Comparativo}, toptitle=1mm, coltitle=black, opacityback=0, arc=.35mm, bottomrule=.15mm, toprule=.15mm, colframe=quarto-callout-tip-color-frame]

\begin{codelisting}[H]

\caption{\label{lst-semana02-5}Generar Boxplots}

\centering{

\begin{Shaded}
\begin{Highlighting}[]
\CommentTok{\# Generar datos simulados}
\FunctionTok{set.seed}\NormalTok{(}\DecValTok{123}\NormalTok{)}
\NormalTok{grupo1 }\OtherTok{\textless{}{-}} \FunctionTok{rnorm}\NormalTok{(}\DecValTok{100}\NormalTok{, }\AttributeTok{mean =} \DecValTok{100}\NormalTok{, }\AttributeTok{sd =} \DecValTok{10}\NormalTok{)}
\NormalTok{grupo2 }\OtherTok{\textless{}{-}} \FunctionTok{rnorm}\NormalTok{(}\DecValTok{100}\NormalTok{, }\AttributeTok{mean =} \DecValTok{110}\NormalTok{, }\AttributeTok{sd =} \DecValTok{15}\NormalTok{)}

\CommentTok{\# Crear boxplots comparativos}
\FunctionTok{boxplot}\NormalTok{(grupo1, grupo2,}
        \AttributeTok{names =} \FunctionTok{c}\NormalTok{(}\StringTok{"Grupo 1"}\NormalTok{, }\StringTok{"Grupo 2"}\NormalTok{),}
        \AttributeTok{main =} \StringTok{"Comparación de Distribuciones"}\NormalTok{,}
        \AttributeTok{ylab =} \StringTok{"Valor"}\NormalTok{,}
        \AttributeTok{col =} \FunctionTok{c}\NormalTok{(}\StringTok{"steelblue"}\NormalTok{, }\StringTok{"seagreen"}\NormalTok{))}
\end{Highlighting}
\end{Shaded}

}

\end{codelisting}%

\pandocbounded{\includegraphics[keepaspectratio]{chapters/semana02_files/figure-pdf/unnamed-chunk-7-1.pdf}}

\end{tcolorbox}

\section{Análisis Bivariante: Distribuciones
Conjuntas}\label{sec-semana02-s5}

\subsection{Distribución Conjunta}\label{sec-semana02-s5-conjunta}

\begin{tcolorbox}[enhanced jigsaw, breakable, titlerule=0mm, rightrule=.15mm, leftrule=.75mm, left=2mm, opacitybacktitle=0.6, colback=white, bottomtitle=1mm, colbacktitle=quarto-callout-note-color!10!white, title=\textcolor{quarto-callout-note-color}{\faInfo}\hspace{0.5em}{Distribución Conjunta}, toptitle=1mm, coltitle=black, opacityback=0, arc=.35mm, bottomrule=.15mm, toprule=.15mm, colframe=quarto-callout-note-color-frame]

La \textbf{distribución conjunta} (o distribución multivariante)
describe simultáneamente dos o más variables. Para dos variables,
hablamos de \textbf{distribución bivariante}.

Cuando observamos pares \((x_i, y_i)\), podemos estudiar cómo ambas
variables varían juntas, no solo separadamente.

\end{tcolorbox}

\subsection{Variables Discretas: Tablas de
Frecuencia}\label{sec-semana02-s5-discretas}

\begin{tcolorbox}[enhanced jigsaw, breakable, titlerule=0mm, rightrule=.15mm, leftrule=.75mm, left=2mm, opacitybacktitle=0.6, colback=white, bottomtitle=1mm, colbacktitle=quarto-callout-note-color!10!white, title=\textcolor{quarto-callout-note-color}{\faInfo}\hspace{0.5em}{Tabla de Frecuencia (Contingencia)}, toptitle=1mm, coltitle=black, opacityback=0, arc=.35mm, bottomrule=.15mm, toprule=.15mm, colframe=quarto-callout-note-color-frame]

Una \textbf{tabla de contingencia} o \textbf{tabla de frecuencia
bidimensional} resume la relación entre dos variables discretas:

\[
\begin{array}{c|ccccc|c}
X \backslash Y & y_1 & \cdots & y_j & \cdots & y_r & \text{Total}_X \\
\hline
x_1 & h_{11} & \cdots & h_{1j} & \cdots & h_{1r} & h_{1\bullet} \\
\vdots & \vdots & \ddots & \vdots & \ddots & \vdots & \vdots \\
x_i & h_{i1} & \cdots & h_{ij} & \cdots & h_{ir} & h_{i\bullet} \\
\vdots & \vdots & \ddots & \vdots & \ddots & \vdots & \vdots \\
x_m & h_{m1} & \cdots & h_{mj} & \cdots & h_{mr} & h_{m\bullet} \\
\hline
\text{Total}_Y & h_{\bullet 1} & \cdots & h_{\bullet j} & \cdots & h_{\bullet r} & n
\end{array}
\]

donde: - \(h_{ij}\) es la \textbf{frecuencia absoluta} (número de casos
con características \(x_i\) e \(y_j\)) -
\(h_{i\bullet} = \sum_j h_{ij}\) es la distribución marginal de \(X\) -
\(h_{\bullet j} = \sum_i h_{ij}\) es la distribución marginal de \(Y\) -
\(f_{ij} = h_{ij}/n\) es la \textbf{frecuencia relativa}

\end{tcolorbox}

\subsection{Distribución Marginal}\label{sec-semana02-s5-marginal}

\begin{tcolorbox}[enhanced jigsaw, breakable, titlerule=0mm, rightrule=.15mm, leftrule=.75mm, left=2mm, opacitybacktitle=0.6, colback=white, bottomtitle=1mm, colbacktitle=quarto-callout-note-color!10!white, title=\textcolor{quarto-callout-note-color}{\faInfo}\hspace{0.5em}{Distribución Marginal}, toptitle=1mm, coltitle=black, opacityback=0, arc=.35mm, bottomrule=.15mm, toprule=.15mm, colframe=quarto-callout-note-color-frame]

La \textbf{distribución marginal} de una variable es su distribución
ignorando las otras variables. En una tabla de contingencia:

\begin{itemize}
\tightlist
\item
  Distribución marginal de \(X\): \(h_{i\bullet} = \sum_j h_{ij}\)
\item
  Distribución marginal de \(Y\): \(h_{\bullet j} = \sum_i h_{ij}\)
\end{itemize}

Las distribuciones marginales aparecen en los totales de filas y
columnas.

\end{tcolorbox}

\subsection{Distribución Condicional}\label{sec-semana02-s5-condicional}

\begin{tcolorbox}[enhanced jigsaw, breakable, titlerule=0mm, rightrule=.15mm, leftrule=.75mm, left=2mm, opacitybacktitle=0.6, colback=white, bottomtitle=1mm, colbacktitle=quarto-callout-note-color!10!white, title=\textcolor{quarto-callout-note-color}{\faInfo}\hspace{0.5em}{Distribución Condicional}, toptitle=1mm, coltitle=black, opacityback=0, arc=.35mm, bottomrule=.15mm, toprule=.15mm, colframe=quarto-callout-note-color-frame]

La \textbf{distribución condicional} de una variable es su distribución
cuando la otra variable toma un valor específico:

\textbf{Distribución de \(X\) dado \(Y = y_j\):}
\[f(x_i | y_j) = \frac{h_{ij}}{h_{\bullet j}}\]

\textbf{Distribución de \(Y\) dado \(X = x_i\):}
\[f(y_j | x_i) = \frac{h_{ij}}{h_{i\bullet}}\]

\end{tcolorbox}

\begin{tcolorbox}[enhanced jigsaw, breakable, titlerule=0mm, rightrule=.15mm, leftrule=.75mm, left=2mm, opacitybacktitle=0.6, colback=white, bottomtitle=1mm, colbacktitle=quarto-callout-tip-color!10!white, title=\textcolor{quarto-callout-tip-color}{\faLightbulb}\hspace{0.5em}{Ejemplo 2.6: Distribucines marginales y condicionales}, toptitle=1mm, coltitle=black, opacityback=0, arc=.35mm, bottomrule=.15mm, toprule=.15mm, colframe=quarto-callout-tip-color-frame]

\begin{longtable}[]{@{}lrrr@{}}
\toprule\noalign{}
Ocupación & Act. Rara & Act. Ocasional & Act. Regular \\
\midrule\noalign{}
\endhead
\bottomrule\noalign{}
\endlastfoot
Obrero & 240 & 120 & 70 \\
No-manual & 160 & 90 & 90 \\
Oficina & 30 & 30 & 30 \\
Agricultor & 37 & 7 & 6 \\
Otros & 40 & 32 & 18 \\
\end{longtable}

\begin{codelisting}[H]

\caption{\label{lst-semana02-6}Código R para porcentajes}

\centering{

\begin{Shaded}
\begin{Highlighting}[]
\FunctionTok{options}\NormalTok{(}\AttributeTok{scipen =} \DecValTok{999}\NormalTok{)}
\CommentTok{\# Crear tabla de contingencia}
\NormalTok{tabla }\OtherTok{\textless{}{-}} \FunctionTok{rbind}\NormalTok{(}
  \FunctionTok{c}\NormalTok{(}\DecValTok{240}\NormalTok{, }\DecValTok{120}\NormalTok{, }\DecValTok{70}\NormalTok{),}
  \FunctionTok{c}\NormalTok{(}\DecValTok{160}\NormalTok{, }\DecValTok{90}\NormalTok{, }\DecValTok{90}\NormalTok{),}
  \FunctionTok{c}\NormalTok{(}\DecValTok{30}\NormalTok{, }\DecValTok{30}\NormalTok{, }\DecValTok{30}\NormalTok{),}
  \FunctionTok{c}\NormalTok{(}\DecValTok{37}\NormalTok{, }\DecValTok{7}\NormalTok{, }\DecValTok{6}\NormalTok{),}
  \FunctionTok{c}\NormalTok{(}\DecValTok{40}\NormalTok{, }\DecValTok{32}\NormalTok{, }\DecValTok{18}\NormalTok{)}
\NormalTok{)}

\FunctionTok{colnames}\NormalTok{(tabla) }\OtherTok{\textless{}{-}} \FunctionTok{c}\NormalTok{(}\StringTok{"Rara"}\NormalTok{, }\StringTok{"Ocasional"}\NormalTok{, }\StringTok{"Regular"}\NormalTok{)}
\FunctionTok{rownames}\NormalTok{(tabla) }\OtherTok{\textless{}{-}} \FunctionTok{c}\NormalTok{(}\StringTok{"Obrero"}\NormalTok{, }\StringTok{"No{-}manual"}\NormalTok{, }\StringTok{"Oficina"}\NormalTok{, }\StringTok{"Agricultor"}\NormalTok{, }\StringTok{"Otros"}\NormalTok{)}

\CommentTok{\# Mostrar tabla marginal y porcentajes}
\FunctionTok{print}\NormalTok{(tabla)}
\end{Highlighting}
\end{Shaded}

}

\end{codelisting}%

\begin{verbatim}
           Rara Ocasional Regular
Obrero      240       120      70
No-manual   160        90      90
Oficina      30        30      30
Agricultor   37         7       6
Otros        40        32      18
\end{verbatim}

\begin{codelisting}[H]

\caption{\label{lst-semana02-6}Código R para porcentajes}

\centering{

\begin{Shaded}
\begin{Highlighting}[]
\FunctionTok{print}\NormalTok{(}\StringTok{"Distribución marginal de ocupación:"}\NormalTok{)}
\end{Highlighting}
\end{Shaded}

}

\end{codelisting}%

\begin{verbatim}
[1] "Distribución marginal de ocupación:"
\end{verbatim}

\begin{codelisting}[H]

\caption{\label{lst-semana02-6}Código R para porcentajes}

\centering{

\begin{Shaded}
\begin{Highlighting}[]
\FunctionTok{round}\NormalTok{(}\FunctionTok{rowSums}\NormalTok{(tabla), }\DecValTok{3}\NormalTok{)}
\end{Highlighting}
\end{Shaded}

}

\end{codelisting}%

\begin{verbatim}
    Obrero  No-manual    Oficina Agricultor      Otros 
       430        340         90         50         90 
\end{verbatim}

\begin{codelisting}[H]

\caption{\label{lst-semana02-6}Código R para porcentajes}

\centering{

\begin{Shaded}
\begin{Highlighting}[]
\FunctionTok{print}\NormalTok{(}\StringTok{"Distribución marginal de actividad:"}\NormalTok{)}
\end{Highlighting}
\end{Shaded}

}

\end{codelisting}%

\begin{verbatim}
[1] "Distribución marginal de actividad:"
\end{verbatim}

\begin{codelisting}[H]

\caption{\label{lst-semana02-6}Código R para porcentajes}

\centering{

\begin{Shaded}
\begin{Highlighting}[]
\FunctionTok{round}\NormalTok{(}\FunctionTok{colSums}\NormalTok{(tabla), }\DecValTok{3}\NormalTok{)}
\end{Highlighting}
\end{Shaded}

}

\end{codelisting}%

\begin{verbatim}
     Rara Ocasional   Regular 
      507       279       214 
\end{verbatim}

\begin{codelisting}[H]

\caption{\label{lst-semana02-6}Código R para porcentajes}

\centering{

\begin{Shaded}
\begin{Highlighting}[]
\CommentTok{\# Mostrar tabla y calcular porcentajes condicionales}
\FunctionTok{print}\NormalTok{(tabla)}
\end{Highlighting}
\end{Shaded}

}

\end{codelisting}%

\begin{verbatim}
           Rara Ocasional Regular
Obrero      240       120      70
No-manual   160        90      90
Oficina      30        30      30
Agricultor   37         7       6
Otros        40        32      18
\end{verbatim}

\begin{codelisting}[H]

\caption{\label{lst-semana02-6}Código R para porcentajes}

\centering{

\begin{Shaded}
\begin{Highlighting}[]
\FunctionTok{print}\NormalTok{(}\StringTok{"Porcentajes por filas:"}\NormalTok{)}
\end{Highlighting}
\end{Shaded}

}

\end{codelisting}%

\begin{verbatim}
[1] "Porcentajes por filas:"
\end{verbatim}

\begin{codelisting}[H]

\caption{\label{lst-semana02-6}Código R para porcentajes}

\centering{

\begin{Shaded}
\begin{Highlighting}[]
\FunctionTok{round}\NormalTok{(}\FunctionTok{prop.table}\NormalTok{(tabla, }\DecValTok{1}\NormalTok{), }\DecValTok{3}\NormalTok{)}
\end{Highlighting}
\end{Shaded}

}

\end{codelisting}%

\begin{verbatim}
            Rara Ocasional Regular
Obrero     0.558     0.279   0.163
No-manual  0.471     0.265   0.265
Oficina    0.333     0.333   0.333
Agricultor 0.740     0.140   0.120
Otros      0.444     0.356   0.200
\end{verbatim}

\begin{codelisting}[H]

\caption{\label{lst-semana02-6}Código R para porcentajes}

\centering{

\begin{Shaded}
\begin{Highlighting}[]
\FunctionTok{print}\NormalTok{(}\StringTok{"Porcentajes por columnas:"}\NormalTok{)}
\end{Highlighting}
\end{Shaded}

}

\end{codelisting}%

\begin{verbatim}
[1] "Porcentajes por columnas:"
\end{verbatim}

\begin{codelisting}[H]

\caption{\label{lst-semana02-6}Código R para porcentajes}

\centering{

\begin{Shaded}
\begin{Highlighting}[]
\FunctionTok{round}\NormalTok{(}\FunctionTok{prop.table}\NormalTok{(tabla, }\DecValTok{2}\NormalTok{), }\DecValTok{3}\NormalTok{)}
\end{Highlighting}
\end{Shaded}

}

\end{codelisting}%

\begin{verbatim}
            Rara Ocasional Regular
Obrero     0.473     0.430   0.327
No-manual  0.316     0.323   0.421
Oficina    0.059     0.108   0.140
Agricultor 0.073     0.025   0.028
Otros      0.079     0.115   0.084
\end{verbatim}

\end{tcolorbox}

\begin{tcolorbox}[enhanced jigsaw, breakable, titlerule=0mm, rightrule=.15mm, leftrule=.75mm, left=2mm, opacitybacktitle=0.6, colback=white, bottomtitle=1mm, colbacktitle=quarto-callout-tip-color!10!white, title=\textcolor{quarto-callout-tip-color}{\faLightbulb}\hspace{0.5em}{Ejemplo 2.7: Distribuciones Condicionales y Probabilidad}, toptitle=1mm, coltitle=black, opacityback=0, arc=.35mm, bottomrule=.15mm, toprule=.15mm, colframe=quarto-callout-tip-color-frame]

Para ilustrar, consideremos la tabla de contingencia de un test de
diagnóstico:

\begin{longtable}[]{@{}lrrr@{}}
\toprule\noalign{}
& Infectado & No infectado & Total \\
\midrule\noalign{}
\endhead
\bottomrule\noalign{}
\endlastfoot
\textbf{Test +} & 199 & 499 & 698 \\
\textbf{Test -} & 1 & 99301 & 99302 \\
\textbf{Total} & 200 & 99800 & 100000 \\
\end{longtable}

\begin{codelisting}[H]

\caption{\label{lst-semana02-7}Probabilidad condicional en R}

\centering{

\begin{Shaded}
\begin{Highlighting}[]
\FunctionTok{options}\NormalTok{(}\AttributeTok{scipen =} \DecValTok{999}\NormalTok{)}
\CommentTok{\# Crear tabla de contingencia 2x2: Test de VIH vs Infección}
\NormalTok{tabla\_hiv }\OtherTok{\textless{}{-}} \FunctionTok{rbind}\NormalTok{(}
  \FunctionTok{c}\NormalTok{(}\DecValTok{199}\NormalTok{, }\DecValTok{499}\NormalTok{),      }\CommentTok{\# Test positivo}
  \FunctionTok{c}\NormalTok{(}\DecValTok{1}\NormalTok{, }\DecValTok{99301}\NormalTok{)       }\CommentTok{\# Test negativo}
\NormalTok{)}
\FunctionTok{colnames}\NormalTok{(tabla\_hiv) }\OtherTok{\textless{}{-}} \FunctionTok{c}\NormalTok{(}\StringTok{"Infectado"}\NormalTok{, }\StringTok{"No infectado"}\NormalTok{)}
\FunctionTok{rownames}\NormalTok{(tabla\_hiv) }\OtherTok{\textless{}{-}} \FunctionTok{c}\NormalTok{(}\StringTok{"Test +"}\NormalTok{, }\StringTok{"Test {-}"}\NormalTok{)}

\CommentTok{\# 1. Porcentajes por FILA (P(Resultado | Infección))}
\FunctionTok{print}\NormalTok{(}\StringTok{"P(Test | Infección) {-} Porcentajes por FILA:"}\NormalTok{)}
\end{Highlighting}
\end{Shaded}

}

\end{codelisting}%

\begin{verbatim}
[1] "P(Test | Infección) - Porcentajes por FILA:"
\end{verbatim}

\begin{codelisting}[H]

\caption{\label{lst-semana02-7}Probabilidad condicional en R}

\centering{

\begin{Shaded}
\begin{Highlighting}[]
\FunctionTok{prop.table}\NormalTok{(tabla\_hiv, }\DecValTok{1}\NormalTok{)}
\end{Highlighting}
\end{Shaded}

}

\end{codelisting}%

\begin{verbatim}
           Infectado No infectado
Test + 0.28510028653    0.7148997
Test - 0.00001007029    0.9999899
\end{verbatim}

\begin{codelisting}[H]

\caption{\label{lst-semana02-7}Probabilidad condicional en R}

\centering{

\begin{Shaded}
\begin{Highlighting}[]
\CommentTok{\# 2. Porcentajes por COLUMNA (Valor Predictivo: P(Infección | Test))}
\FunctionTok{print}\NormalTok{(}\StringTok{"P(Infección | Test) {-} Porcentajes por COLUMNA:"}\NormalTok{)}
\end{Highlighting}
\end{Shaded}

}

\end{codelisting}%

\begin{verbatim}
[1] "P(Infección | Test) - Porcentajes por COLUMNA:"
\end{verbatim}

\begin{codelisting}[H]

\caption{\label{lst-semana02-7}Probabilidad condicional en R}

\centering{

\begin{Shaded}
\begin{Highlighting}[]
\FunctionTok{prop.table}\NormalTok{(tabla\_hiv, }\DecValTok{2}\NormalTok{)}
\end{Highlighting}
\end{Shaded}

}

\end{codelisting}%

\begin{verbatim}
       Infectado No infectado
Test +     0.995        0.005
Test -     0.005        0.995
\end{verbatim}

\textbf{Interpretación Clínica:} - \textbf{Por filas}: (\text{Test}
\textbar{} \text{Infección}) qué tan bueno es el test para identificar
la presencia o ausencia de enfermedad. - \textbf{Por columnas}:
(\text{Infección} \textbar{} \text{Test}) a la pregunta clínica crucial:
``Dado que el paciente tiene este resultado de test, ¿cuál es la
probabilidad real de que esté infectado?''.

\end{tcolorbox}

\begin{tcolorbox}[enhanced jigsaw, breakable, titlerule=0mm, rightrule=.15mm, leftrule=.75mm, left=2mm, opacitybacktitle=0.6, colback=white, bottomtitle=1mm, colbacktitle=quarto-callout-note-color!10!white, title=\textcolor{quarto-callout-note-color}{\faInfo}\hspace{0.5em}{Métodos Avanzados}, toptitle=1mm, coltitle=black, opacityback=0, arc=.35mm, bottomrule=.15mm, toprule=.15mm, colframe=quarto-callout-note-color-frame]

Para ampliar los contenidos sobre los test diagnósticos y otras técnicas
estadísticas avanzadas, visita:
\textbf{\href{https://migariane.github.io/Bioestadistica}{→
Bioestadística Avanzada --- M.A.~Luque Fernández}}

\end{tcolorbox}

\begin{tcolorbox}[enhanced jigsaw, breakable, titlerule=0mm, rightrule=.15mm, leftrule=.75mm, left=2mm, opacitybacktitle=0.6, colback=white, bottomtitle=1mm, colbacktitle=quarto-callout-tip-color!10!white, title=\textcolor{quarto-callout-tip-color}{\faLightbulb}\hspace{0.5em}{Ejemplo 2.8: Covarianza y Correlación}, toptitle=1mm, coltitle=black, opacityback=0, arc=.35mm, bottomrule=.15mm, toprule=.15mm, colframe=quarto-callout-tip-color-frame]

\begin{codelisting}[H]

\caption{\label{lst-semana02-10}Covarianza y Correlación con R Base}

\centering{

\begin{Shaded}
\begin{Highlighting}[]
\CommentTok{\# Cargar el dataset women}
\FunctionTok{data}\NormalTok{(women)}

\CommentTok{\# Calcular estadísticas}
\NormalTok{cov\_hw }\OtherTok{\textless{}{-}} \FunctionTok{cov}\NormalTok{(women}\SpecialCharTok{$}\NormalTok{height, women}\SpecialCharTok{$}\NormalTok{weight)}
\NormalTok{cor\_hw }\OtherTok{\textless{}{-}} \FunctionTok{cor}\NormalTok{(women}\SpecialCharTok{$}\NormalTok{height, women}\SpecialCharTok{$}\NormalTok{weight)}

\FunctionTok{cat}\NormalTok{(}\StringTok{"Covarianza:"}\NormalTok{, }\FunctionTok{round}\NormalTok{(cov\_hw, }\DecValTok{2}\NormalTok{), }\StringTok{"}\SpecialCharTok{\textbackslash{}n}\StringTok{"}\NormalTok{)}
\end{Highlighting}
\end{Shaded}

}

\end{codelisting}%

\begin{verbatim}
Covarianza: 69 
\end{verbatim}

\begin{codelisting}[H]

\caption{\label{lst-semana02-10}Covarianza y Correlación con R Base}

\centering{

\begin{Shaded}
\begin{Highlighting}[]
\FunctionTok{cat}\NormalTok{(}\StringTok{"Correlación:"}\NormalTok{, }\FunctionTok{round}\NormalTok{(cor\_hw, }\DecValTok{2}\NormalTok{), }\StringTok{"}\SpecialCharTok{\textbackslash{}n}\StringTok{"}\NormalTok{)}
\end{Highlighting}
\end{Shaded}

}

\end{codelisting}%

\begin{verbatim}
Correlación: 1 
\end{verbatim}

\begin{codelisting}[H]

\caption{\label{lst-semana02-10}Covarianza y Correlación con R Base}

\centering{

\begin{Shaded}
\begin{Highlighting}[]
\CommentTok{\# Gráfico de dispersión}
\FunctionTok{plot}\NormalTok{(women}\SpecialCharTok{$}\NormalTok{height, women}\SpecialCharTok{$}\NormalTok{weight,}
     \AttributeTok{main =} \StringTok{"Altura vs Peso"}\NormalTok{,}
     \AttributeTok{xlab =} \StringTok{"Altura (pulgadas)"}\NormalTok{,}
     \AttributeTok{ylab =} \StringTok{"Peso (libras)"}\NormalTok{,}
     \AttributeTok{pch =} \DecValTok{19}\NormalTok{, }\AttributeTok{col =} \StringTok{"steelblue"}\NormalTok{)}

\CommentTok{\# Agregar línea de regresión}
\FunctionTok{abline}\NormalTok{(}\FunctionTok{lm}\NormalTok{(weight }\SpecialCharTok{\textasciitilde{}}\NormalTok{ height, }\AttributeTok{data =}\NormalTok{ women), }\AttributeTok{col =} \StringTok{"red"}\NormalTok{, }\AttributeTok{lwd =} \DecValTok{2}\NormalTok{)}
\end{Highlighting}
\end{Shaded}

}

\end{codelisting}%

\pandocbounded{\includegraphics[keepaspectratio]{chapters/semana02_files/figure-pdf/unnamed-chunk-10-1.pdf}}

\end{tcolorbox}

\subsection{Covarianza y Correlacion: Advertencia
Visual}\label{sec-semana02-s6-datasaurus}

\begin{tcolorbox}[enhanced jigsaw, breakable, titlerule=0mm, rightrule=.15mm, leftrule=.75mm, left=2mm, opacitybacktitle=0.6, colback=white, bottomtitle=1mm, colbacktitle=quarto-callout-note-color!10!white, title=\textcolor{quarto-callout-note-color}{\faInfo}\hspace{0.5em}{Limitaciones del coeficiente de correlación}, toptitle=1mm, coltitle=black, opacityback=0, arc=.35mm, bottomrule=.15mm, toprule=.15mm, colframe=quarto-callout-note-color-frame]

Para comprender mejor por qué la correlación debe ir acompañada de un
gráfico, observemos la siguiente figura:

\begin{figure}[H]

\centering{

\pandocbounded{\includegraphics[keepaspectratio]{chapters/semana02_files/figure-pdf/fig-semana02-cor-limit-1.pdf}}

}

\caption{\label{fig-semana02-cor-limit}Limitaciones de la correlación de
Pearson}

\end{figure}%

En los dos primeros casos, existe una dependencia clara entre las
variables, pero el coeficiente de correlación \(r\) es cercano a 0
porque no es una relación lineal.

\end{tcolorbox}

\begin{tcolorbox}[enhanced jigsaw, breakable, titlerule=0mm, rightrule=.15mm, leftrule=.75mm, left=2mm, opacitybacktitle=0.6, colback=white, bottomtitle=1mm, colbacktitle=quarto-callout-warning-color!10!white, title=\textcolor{quarto-callout-warning-color}{\faExclamationTriangle}\hspace{0.5em}{La Paradoja de Simpson: Correlaciones Engañosas}, toptitle=1mm, coltitle=black, opacityback=0, arc=.35mm, bottomrule=.15mm, toprule=.15mm, colframe=quarto-callout-warning-color-frame]

La \textbf{Paradoja de Simpson} es un fenómeno donde una tendencia
observada en varios grupos desaparece o se invierte al combinarlos.

\textbf{Relación con el ejemplo clínico:} En la figura a continuación
(Ejemplo 2.11), vemos dos subgrupos (leves y graves). Si calculáramos la
correlación total de todos los pacientes sin distinguir estos grupos, el
resultado sería engañoso. La variable ``gravedad'' actúa aquí como un
\textbf{factor de confusión}. Este ejemplo clínico ilustra perfectamente
por qué la \textbf{``correlación no implica causalidad''} y subraya la
importancia crítica de la estratificación y visualización de datos en
medicina.

\end{tcolorbox}

\begin{tcolorbox}[enhanced jigsaw, breakable, titlerule=0mm, rightrule=.15mm, leftrule=.75mm, left=2mm, opacitybacktitle=0.6, colback=white, bottomtitle=1mm, colbacktitle=quarto-callout-tip-color!10!white, title=\textcolor{quarto-callout-tip-color}{\faLightbulb}\hspace{0.5em}{Ejemplo 2.9: Paradoja de Simpson en Ensayos Clínicos}, toptitle=1mm, coltitle=black, opacityback=0, arc=.35mm, bottomrule=.15mm, toprule=.15mm, colframe=quarto-callout-tip-color-frame]

\begin{codelisting}[H]

\caption{\label{lst-semana02-11}Ilustración de la Paradoja de Simpson}

\centering{

\begin{Shaded}
\begin{Highlighting}[]
\CommentTok{\# Simulación: Efecto de un fármaco en dos subgrupos (ej. Pacientes leves vs graves)}
\FunctionTok{set.seed}\NormalTok{(}\DecValTok{42}\NormalTok{)}
\NormalTok{n }\OtherTok{\textless{}{-}} \DecValTok{100}
\CommentTok{\# Subgrupo 1 (Leves)}
\NormalTok{x1 }\OtherTok{\textless{}{-}} \FunctionTok{rnorm}\NormalTok{(n, }\AttributeTok{mean =} \DecValTok{50}\NormalTok{, }\AttributeTok{sd =} \DecValTok{10}\NormalTok{); y1 }\OtherTok{\textless{}{-}} \DecValTok{20} \SpecialCharTok{+} \FloatTok{0.5} \SpecialCharTok{*}\NormalTok{ x1 }\SpecialCharTok{+} \FunctionTok{rnorm}\NormalTok{(n, }\AttributeTok{sd =} \DecValTok{5}\NormalTok{)}
\CommentTok{\# Subgrupo 2 (Graves)}
\NormalTok{x2 }\OtherTok{\textless{}{-}} \FunctionTok{rnorm}\NormalTok{(n, }\AttributeTok{mean =} \DecValTok{70}\NormalTok{, }\AttributeTok{sd =} \DecValTok{10}\NormalTok{); y2 }\OtherTok{\textless{}{-}} \DecValTok{50} \SpecialCharTok{+} \FloatTok{0.5} \SpecialCharTok{*}\NormalTok{ x2 }\SpecialCharTok{+} \FunctionTok{rnorm}\NormalTok{(n, }\AttributeTok{sd =} \DecValTok{5}\NormalTok{)}

\CommentTok{\# Gráfico}
\FunctionTok{plot}\NormalTok{(}\FunctionTok{c}\NormalTok{(x1, x2), }\FunctionTok{c}\NormalTok{(y1, y2), }
     \AttributeTok{col =} \FunctionTok{rep}\NormalTok{(}\FunctionTok{c}\NormalTok{(}\StringTok{"blue"}\NormalTok{, }\StringTok{"red"}\NormalTok{), }\AttributeTok{each =}\NormalTok{ n),}
     \AttributeTok{pch =} \DecValTok{19}\NormalTok{,}
     \AttributeTok{main =} \StringTok{"Paradoja de Simpson en un ensayo clínico"}\NormalTok{,}
     \AttributeTok{xlab =} \StringTok{"Dosificación (mg)"}\NormalTok{, }\AttributeTok{ylab =} \StringTok{"Recuperación (\%)"}\NormalTok{)}
\FunctionTok{legend}\NormalTok{(}\StringTok{"topleft"}\NormalTok{, }\AttributeTok{legend=}\FunctionTok{c}\NormalTok{(}\StringTok{"Pacientes Leves"}\NormalTok{, }\StringTok{"Pacientes Graves"}\NormalTok{), }\AttributeTok{col=}\FunctionTok{c}\NormalTok{(}\StringTok{"blue"}\NormalTok{, }\StringTok{"red"}\NormalTok{), }\AttributeTok{pch=}\DecValTok{19}\NormalTok{)}
\end{Highlighting}
\end{Shaded}

}

\end{codelisting}%

\pandocbounded{\includegraphics[keepaspectratio]{chapters/semana02_files/figure-pdf/unnamed-chunk-11-1.pdf}}

\end{tcolorbox}

\subsection{Matriz de Correlaciones}\label{sec-semana02-s6-matriz}

La matriz de correlaciones es una herramienta fundamental en
investigación médica para evaluar la relación entre múltiples variables
fisiológicas simultáneamente. Su utilidad principal radica en:

\begin{enumerate}
\def\labelenumi{\arabic{enumi}.}
\tightlist
\item
  \textbf{Identificación de Predictores}: Detectar qué variables
  clínicas están más asociadas con un marcador de salud (ej. qué
  variables influyen más en la presión arterial).
\item
  \textbf{Detección de Colinealidad}: Identificar variables redundantes
  que aportan información similar, lo cual es crítico antes de construir
  modelos de regresión.
\item
  \textbf{Análisis de Clusters}: Agrupar variables que se comportan de
  manera similar en pacientes (ej. marcadores inflamatorios).
\end{enumerate}

\begin{tcolorbox}[enhanced jigsaw, breakable, titlerule=0mm, rightrule=.15mm, leftrule=.75mm, left=2mm, opacitybacktitle=0.6, colback=white, bottomtitle=1mm, colbacktitle=quarto-callout-tip-color!10!white, title=\textcolor{quarto-callout-tip-color}{\faLightbulb}\hspace{0.5em}{Ejemplo 2.10: Matriz de Correlaciones (Parámetros Clínicos)}, toptitle=1mm, coltitle=black, opacityback=0, arc=.35mm, bottomrule=.15mm, toprule=.15mm, colframe=quarto-callout-tip-color-frame]

\begin{codelisting}[H]

\caption{\label{lst-semana02-12}Matriz de correlaciones clínicas}

\centering{

\begin{Shaded}
\begin{Highlighting}[]
\CommentTok{\# Simulación de variables clínicas: Edad, Colesterol, Tensión Arterial (PAS), IMC}
\FunctionTok{set.seed}\NormalTok{(}\DecValTok{42}\NormalTok{)}
\NormalTok{n }\OtherTok{\textless{}{-}} \DecValTok{200}
\NormalTok{edad }\OtherTok{\textless{}{-}} \FunctionTok{rnorm}\NormalTok{(n, }\DecValTok{55}\NormalTok{, }\DecValTok{10}\NormalTok{)}
\NormalTok{imc }\OtherTok{\textless{}{-}} \FunctionTok{rnorm}\NormalTok{(n, }\DecValTok{27}\NormalTok{, }\DecValTok{4}\NormalTok{)}
\NormalTok{colesterol }\OtherTok{\textless{}{-}} \DecValTok{150} \SpecialCharTok{+} \FloatTok{0.8} \SpecialCharTok{*}\NormalTok{ edad }\SpecialCharTok{+} \DecValTok{2} \SpecialCharTok{*}\NormalTok{ imc }\SpecialCharTok{+} \FunctionTok{rnorm}\NormalTok{(n, }\DecValTok{0}\NormalTok{, }\DecValTok{10}\NormalTok{)}
\NormalTok{pas }\OtherTok{\textless{}{-}} \DecValTok{110} \SpecialCharTok{+} \FloatTok{0.5} \SpecialCharTok{*}\NormalTok{ edad }\SpecialCharTok{+} \FloatTok{1.5} \SpecialCharTok{*}\NormalTok{ imc }\SpecialCharTok{+} \FunctionTok{rnorm}\NormalTok{(n, }\DecValTok{0}\NormalTok{, }\DecValTok{8}\NormalTok{)}

\NormalTok{data\_clinica }\OtherTok{\textless{}{-}} \FunctionTok{data.frame}\NormalTok{(}\AttributeTok{Edad =}\NormalTok{ edad, }\AttributeTok{IMC =}\NormalTok{ imc, }\AttributeTok{Colesterol =}\NormalTok{ colesterol, }\AttributeTok{PAS =}\NormalTok{ pas)}

\CommentTok{\# Visualización mediante matriz de dispersión (Pairs plot)}
\FunctionTok{pairs}\NormalTok{(data\_clinica,}
      \AttributeTok{main =} \StringTok{"Matriz de Relaciones Clínicas"}\NormalTok{,}
      \AttributeTok{pch =} \DecValTok{19}\NormalTok{,}
      \AttributeTok{col =} \StringTok{"darkgreen"}\NormalTok{)}
\end{Highlighting}
\end{Shaded}

}

\end{codelisting}%

\pandocbounded{\includegraphics[keepaspectratio]{chapters/semana02_files/figure-pdf/unnamed-chunk-12-1.pdf}}

\begin{codelisting}[H]

\caption{\label{lst-semana02-12}Matriz de correlaciones clínicas}

\centering{

\begin{Shaded}
\begin{Highlighting}[]
\CommentTok{\# Correlaciones numéricas}
\FunctionTok{cor}\NormalTok{(data\_clinica)}
\end{Highlighting}
\end{Shaded}

}

\end{codelisting}%

\begin{verbatim}
                  Edad         IMC Colesterol       PAS
Edad        1.00000000 -0.08036539  0.4484613 0.3906186
IMC        -0.08036539  1.00000000  0.5157946 0.5259363
Colesterol  0.44846130  0.51579465  1.0000000 0.4802219
PAS         0.39061863  0.52593630  0.4802219 1.0000000
\end{verbatim}

\end{tcolorbox}

\section{Función de Distribución Empírica
(ECDF)}\label{sec-semana02-ecdf}

En el Capítulo 1 estudiamos la Función de Distribución Empírica (ECDF)
con datos agrupados en clases. Aquí presentamos cómo calcular y
visualizar la ECDF con datos individuales usando la función
\texttt{ecdf()} de R.

\section{Ejemplo 2.11: ECDF de Alturas de
Estudiantes}\label{ejemplo-2.11-ecdf-de-alturas-de-estudiantes}

Se midieron las alturas (en cm) de 30 estudiantes de segundo año:

\begin{verbatim}
168, 172, 165, 175, 170, 169, 173, 166, 174, 171,
167, 176, 169, 172, 168, 170, 175, 169, 173, 171,
165, 174, 170, 172, 168, 175, 169, 171, 173, 170
\end{verbatim}

La tabla siguiente muestra algunos valores y su función de distribución
empírica acumulada:

\begin{longtable}[]{@{}ll@{}}
\toprule\noalign{}
Altura (cm) & \(F_n(x)\) \\
\midrule\noalign{}
\endhead
\bottomrule\noalign{}
\endlastfoot
165 & 0.067 \\
166 & 0.100 \\
167 & 0.133 \\
168 & 0.200 \\
169 & 0.300 \\
170 & 0.400 \\
171 & 0.500 \\
172 & 0.633 \\
173 & 0.700 \\
174 & 0.767 \\
175 & 0.867 \\
176 & 1.000 \\
\end{longtable}

\begin{codelisting}

\caption{\label{lst-semana02-13}Código R para ECDF}

\centering{

\begin{Shaded}
\begin{Highlighting}[]
\CommentTok{\# Datos de alturas}
\NormalTok{alturas }\OtherTok{\textless{}{-}} \FunctionTok{c}\NormalTok{(}\DecValTok{168}\NormalTok{, }\DecValTok{172}\NormalTok{, }\DecValTok{165}\NormalTok{, }\DecValTok{175}\NormalTok{, }\DecValTok{170}\NormalTok{, }\DecValTok{169}\NormalTok{, }\DecValTok{173}\NormalTok{, }\DecValTok{166}\NormalTok{, }\DecValTok{174}\NormalTok{, }\DecValTok{171}\NormalTok{,}
             \DecValTok{167}\NormalTok{, }\DecValTok{176}\NormalTok{, }\DecValTok{169}\NormalTok{, }\DecValTok{172}\NormalTok{, }\DecValTok{168}\NormalTok{, }\DecValTok{170}\NormalTok{, }\DecValTok{175}\NormalTok{, }\DecValTok{169}\NormalTok{, }\DecValTok{173}\NormalTok{, }\DecValTok{171}\NormalTok{,}
             \DecValTok{165}\NormalTok{, }\DecValTok{174}\NormalTok{, }\DecValTok{170}\NormalTok{, }\DecValTok{172}\NormalTok{, }\DecValTok{168}\NormalTok{, }\DecValTok{175}\NormalTok{, }\DecValTok{169}\NormalTok{, }\DecValTok{171}\NormalTok{, }\DecValTok{173}\NormalTok{, }\DecValTok{170}\NormalTok{)}

\CommentTok{\# Calcular ECDF}
\NormalTok{ecdf\_alturas }\OtherTok{\textless{}{-}} \FunctionTok{ecdf}\NormalTok{(alturas)}

\CommentTok{\# Graficar la ECDF}
\FunctionTok{plot}\NormalTok{(ecdf\_alturas,}
     \AttributeTok{main =} \StringTok{"Función de Distribución Empírica (ECDF)}\SpecialCharTok{\textbackslash{}n}\StringTok{Alturas de Estudiantes"}\NormalTok{,}
     \AttributeTok{xlab =} \StringTok{"Altura (cm)"}\NormalTok{,}
     \AttributeTok{ylab =} \StringTok{"F\_n(x)"}\NormalTok{,}
     \AttributeTok{verticals =} \ConstantTok{TRUE}\NormalTok{,}
     \AttributeTok{do.points =} \ConstantTok{FALSE}\NormalTok{,}
     \AttributeTok{col =} \StringTok{"darkblue"}\NormalTok{,}
     \AttributeTok{lwd =} \DecValTok{2}\NormalTok{)}
\FunctionTok{grid}\NormalTok{()}
\end{Highlighting}
\end{Shaded}

}

\end{codelisting}%

\pandocbounded{\includegraphics[keepaspectratio]{chapters/semana02_files/figure-pdf/unnamed-chunk-13-1.pdf}}

\subsection{Interpretación}\label{interpretaciuxf3n}

La ECDF \(F_n(x)\) es una función escalonada que representa la
proporción de datos menores o iguales a \(x\):

\begin{itemize}
\tightlist
\item
  \(F_n(165) \approx 0.067\) → El 6.7\% de estudiantes miden ≤ 165 cm
\item
  \(F_n(170) = 0.400\) → El 40\% de estudiantes miden ≤ 170 cm
\item
  \(F_n(175) \approx 0.867\) → El 86.7\% de estudiantes miden ≤ 175 cm
\end{itemize}

\textbf{Propiedad fundamental:} La ECDF es una \textbf{función no
decreciente} que salta en cada valor observado.

\section{Resumen}\label{sec-semana02-resumen}

Este capítulo ha cubierto los conceptos esenciales para el análisis
exploratorio bivariante de datos:

\begin{longtable}[]{@{}
  >{\raggedright\arraybackslash}p{(\linewidth - 4\tabcolsep) * \real{0.2174}}
  >{\raggedright\arraybackslash}p{(\linewidth - 4\tabcolsep) * \real{0.4130}}
  >{\raggedright\arraybackslash}p{(\linewidth - 4\tabcolsep) * \real{0.3696}}@{}}
\toprule\noalign{}
\begin{minipage}[b]{\linewidth}\raggedright
Concepto
\end{minipage} & \begin{minipage}[b]{\linewidth}\raggedright
Fórmula/Definición
\end{minipage} & \begin{minipage}[b]{\linewidth}\raggedright
Rango/Unidades
\end{minipage} \\
\midrule\noalign{}
\endhead
\bottomrule\noalign{}
\endlastfoot
Varianza & \(s^2 = \frac{1}{n-1}\sum(x_i - \bar{x})^2\) &
\([0, \infty)\) \\
Desviación Estándar & \(s = \sqrt{s^2}\) & Mismas unidades que \(x\) \\
Rango & \(R = x_{\max} - x_{\min}\) & Mismas unidades que \(x\) \\
Rango Intercuartílico & \(\text{IQR} = Q3 - Q1\) & Mismas unidades que
\(x\) \\
Z-Score & \(z = \frac{x - \bar{x}}{s}\) & Adimensional \\
Covarianza &
\(s_{xy} = \frac{1}{n-1}\sum(x_i - \bar{x})(y_i - \bar{y})\) & Producto
de unidades \\
Correlación Pearson & \(r = \frac{s_{xy}}{s_x s_y}\) & \([-1, 1]\) \\
\end{longtable}

\subsection{Puntos Clave}\label{puntos-clave}

\begin{itemize}
\tightlist
\item
  Las \textbf{medidas de dispersión} (varianza, desviación estándar)
  cuantifican la variabilidad de los datos
\item
  Las \textbf{transformaciones lineales} predicen cambios en las
  estadísticas de forma matemática
\item
  Los \textbf{diagramas de caja} visualizan rápidamente la distribución
  y detectan outliers
\item
  El \textbf{análisis bivariante} describe relaciones entre pares de
  variables
\item
  La \textbf{correlación} mide la fuerza de relaciones lineales, pero
  \textbf{no implica causalidad}
\item
  Siempre \textbf{visualiza los datos}: los resúmenes estadísticos
  pueden ocultar patrones subyacentes complejos
\item
  La \textbf{ECDF} permite estimar cuantiles y visualizar la
  distribución acumulada sin asumir ninguna forma paramétrica
\end{itemize}

\section{Ejercicios}\label{sec-semana02-ejercicios}

\begin{tcolorbox}[enhanced jigsaw, breakable, titlerule=0mm, rightrule=.15mm, leftrule=.75mm, left=2mm, opacitybacktitle=0.6, colback=white, bottomtitle=1mm, colbacktitle=quarto-callout-tip-color!10!white, title=\textcolor{quarto-callout-tip-color}{\faLightbulb}\hspace{0.5em}{Ejercicio 2.1: Estadísticos de Dispersión}, toptitle=1mm, coltitle=black, opacityback=0, arc=.35mm, bottomrule=.15mm, toprule=.15mm, colframe=quarto-callout-tip-color-frame]

Se registraron los tiempos de respuesta (en milisegundos) de 12
participantes en una tarea de tiempo de reacción: 250, 240, 300, 280,
270, 265, 290, 255, 310, 275, 260, 320

\begin{enumerate}
\def\labelenumi{\alph{enumi})}
\tightlist
\item
  Calcula el rango y el rango intercuartílico (IQR).
\item
  Calcula la media y la desviación estándar.
\item
  ¿Qué observación es un valor atípico potencial según la regla de
  1.5×IQR?
\item
  Dibuja un diagrama de caja.
\end{enumerate}

\end{tcolorbox}

\begin{tcolorbox}[enhanced jigsaw, breakable, titlerule=0mm, rightrule=.15mm, leftrule=.75mm, left=2mm, opacitybacktitle=0.6, colback=white, bottomtitle=1mm, colbacktitle=quarto-callout-tip-color!10!white, title=\textcolor{quarto-callout-tip-color}{\faLightbulb}\hspace{0.5em}{Ejercicio 2.2: Transformaciones Lineales (Conversión de Unidades)}, toptitle=1mm, coltitle=black, opacityback=0, arc=.35mm, bottomrule=.15mm, toprule=.15mm, colframe=quarto-callout-tip-color-frame]

En un estudio de investigación, se registra la presión intraocular (PIO)
de 50 pacientes en milímetros de mercurio (mmHg). La media registrada es
\(\bar{x} = 15\) mmHg con una desviación estándar de \(s_x = 3\) mmHg.
Para una publicación internacional, es necesario convertir estos valores
a centímetros de agua (cmH₂O) utilizando la fórmula: \(y = 1.36x\).

\begin{enumerate}
\def\labelenumi{\alph{enumi})}
\tightlist
\item
  ¿Cuál será la media de la PIO en cmH₂O?
\item
  ¿Cuál será la desviación estándar en cmH₂O?
\item
  Si el IQR original es de 4 mmHg, ¿cuál será el IQR en cmH₂O?
\end{enumerate}

\end{tcolorbox}

\begin{tcolorbox}[enhanced jigsaw, breakable, titlerule=0mm, rightrule=.15mm, leftrule=.75mm, left=2mm, opacitybacktitle=0.6, colback=white, bottomtitle=1mm, colbacktitle=quarto-callout-tip-color!10!white, title=\textcolor{quarto-callout-tip-color}{\faLightbulb}\hspace{0.5em}{Ejercicio 2.3: Z-Scores y Estandarización Clínica}, toptitle=1mm, coltitle=black, opacityback=0, arc=.35mm, bottomrule=.15mm, toprule=.15mm, colframe=quarto-callout-tip-color-frame]

En un estudio de cribado, se utilizan dos escalas distintas para medir
el deterioro cognitivo. Un paciente obtiene las siguientes puntuaciones:
- Escala A: 85 puntos (media = 75, desviación estándar = 10) - Escala B:
90 puntos (media = 80, desviación estándar = 15)

\begin{enumerate}
\def\labelenumi{\alph{enumi})}
\tightlist
\item
  Calcula los z-scores para ambas escalas.
\item
  ¿En qué escala el deterioro del paciente es relativamente más severo?
\item
  Si se considera un z \textgreater{} 0.5 como ``indicativo de riesgo'',
  ¿en cuál(es) escala(s) el paciente supera este umbral?
\end{enumerate}

\end{tcolorbox}

\begin{tcolorbox}[enhanced jigsaw, breakable, titlerule=0mm, rightrule=.15mm, leftrule=.75mm, left=2mm, opacitybacktitle=0.6, colback=white, bottomtitle=1mm, colbacktitle=quarto-callout-tip-color!10!white, title=\textcolor{quarto-callout-tip-color}{\faLightbulb}\hspace{0.5em}{Ejercicio 2.4: Tablas de Contingencia}, toptitle=1mm, coltitle=black, opacityback=0, arc=.35mm, bottomrule=.15mm, toprule=.15mm, colframe=quarto-callout-tip-color-frame]

Se encuestó a 300 personas sobre si fuman y si tienen problemas
respiratorios. Los resultados fueron:

\begin{longtable}[]{@{}llll@{}}
\toprule\noalign{}
& Problemas & Sin problemas & Total \\
\midrule\noalign{}
\endhead
\bottomrule\noalign{}
\endlastfoot
Fuman & 70 & 40 & 110 \\
No fuman & 30 & 160 & 190 \\
Total & 100 & 200 & 300 \\
\end{longtable}

\begin{enumerate}
\def\labelenumi{\alph{enumi})}
\tightlist
\item
  Calcula la distribución marginal de fumadores.
\item
  Calcula la distribución condicional de problemas respiratorios
  \textbf{dado} que fuma.
\item
  Calcula la distribución condicional de problemas respiratorios
  \textbf{dado} que no fuma.
\item
  Interpreta las diferencias entre las distribuciones condicionales.
\end{enumerate}

\end{tcolorbox}

\begin{tcolorbox}[enhanced jigsaw, breakable, titlerule=0mm, rightrule=.15mm, leftrule=.75mm, left=2mm, opacitybacktitle=0.6, colback=white, bottomtitle=1mm, colbacktitle=quarto-callout-tip-color!10!white, title=\textcolor{quarto-callout-tip-color}{\faLightbulb}\hspace{0.5em}{Ejercicio 2.5: Análisis de Datos Clínicos Simulados}, toptitle=1mm, coltitle=black, opacityback=0, arc=.35mm, bottomrule=.15mm, toprule=.15mm, colframe=quarto-callout-tip-color-frame]

Genera un dataset que represente la relación entre dos variables
biomédicas:

\begin{codelisting}[H]

\caption{\label{lst-semana02-14}Código R: Simulación Clínica}

\centering{

\begin{Shaded}
\begin{Highlighting}[]
\FunctionTok{set.seed}\NormalTok{(}\DecValTok{999}\NormalTok{)}
\NormalTok{edad }\OtherTok{\textless{}{-}} \FunctionTok{rnorm}\NormalTok{(}\DecValTok{100}\NormalTok{, }\AttributeTok{mean =} \DecValTok{50}\NormalTok{, }\AttributeTok{sd =} \DecValTok{15}\NormalTok{)}
\NormalTok{presion\_sistolica }\OtherTok{\textless{}{-}} \DecValTok{100} \SpecialCharTok{+} \FloatTok{0.6} \SpecialCharTok{*}\NormalTok{ edad }\SpecialCharTok{+} \FunctionTok{rnorm}\NormalTok{(}\DecValTok{100}\NormalTok{, }\AttributeTok{sd =} \DecValTok{10}\NormalTok{)}
\end{Highlighting}
\end{Shaded}

}

\end{codelisting}%

Realiza las siguientes tareas: a) Calcula la media, desviación estándar
y covarianza de la edad y la presión sistólica. b) Calcula el
coeficiente de correlación de Pearson. c) Crea un gráfico de dispersión
con línea de regresión. d) Crea boxplots para ambas variables. e)
Estandariza ambas variables usando z-scores. f) Verifica que las
variables estandarizadas tienen media 0 y desviación estándar 1.

\end{tcolorbox}

\begin{tcolorbox}[enhanced jigsaw, breakable, titlerule=0mm, rightrule=.15mm, leftrule=.75mm, left=2mm, opacitybacktitle=0.6, colback=white, bottomtitle=1mm, colbacktitle=quarto-callout-tip-color!10!white, title=\textcolor{quarto-callout-tip-color}{\faLightbulb}\hspace{0.5em}{Ejercicio 2.6: Distribuciones Conjuntas}, toptitle=1mm, coltitle=black, opacityback=0, arc=.35mm, bottomrule=.15mm, toprule=.15mm, colframe=quarto-callout-tip-color-frame]

Se miden presión sistólica (X) y diastólica (Y) en una muestra clínica.
Se encuentra que \(\text{Cov}(x,y) = 180\) (mmHg)², \(s_x = 15\) mmHg,
\(s_y = 10\) mmHg.

\begin{enumerate}
\def\labelenumi{\alph{enumi})}
\tightlist
\item
  Calcula la correlación entre X e Y.
\item
  Interpreta el valor obtenido.
\item
  ¿Son X e Y independientes?
\end{enumerate}

\end{tcolorbox}

\section{Respuestas a los Ejercicios}\label{sec-semana02-respuestas}

\textbf{Ejercicio 2.1: Estadísticos de Dispersión}

\begin{enumerate}
\def\labelenumi{\alph{enumi})}
\item
  Rango: 320 - 240 = 80 ms IQR: Q₃ - Q₁ ≈ 290 - 255 = 35 ms
\item
  Media: ≈ 278.3 ms Desviación estándar: ≈ 26.9 ms
\item
  Outliers (\textgreater{} Q₃ + 1.5×IQR = 290 + 52.5 = 342.5): Ninguno
  detectado Outliers (\textless{} Q₁ - 1.5×IQR = 255 - 52.5 = 202.5):
  Ninguno
\item
  El boxplot muestra distribución simétrica alrededor de 275-280 ms sin
  outliers.
\end{enumerate}

\textbf{Ejercicio 2.2: Transformaciones Lineales (Conversión)}

\begin{enumerate}
\def\labelenumi{\alph{enumi})}
\item
  \(\bar{y} = 1.36 \times 15 = 20.4\) cmH₂O
\item
  \(s_y = 1.36 \times 3 = 4.08\) cmH₂O
\item
  \(\text{IQR}_{y} = 1.36 \times 4 = 5.44\) cmH₂O
\end{enumerate}

\textbf{Ejercicio 2.3: Z-Scores y Estandarización Clínica}

\begin{enumerate}
\def\labelenumi{\alph{enumi})}
\item
  Escala A: \(z = \frac{85 - 75}{10} = 1.0\) Escala B:
  \(z = \frac{90 - 80}{15} = 0.67\)
\item
  El deterioro es más severo en la Escala A, ya que el paciente se aleja
  más de la media poblacional (z=1.0 vs z=0.67).
\item
  El paciente supera el umbral en ambas escalas (1.0 \textgreater{} 0.5
  y 0.67 \textgreater{} 0.5).
\end{enumerate}

\textbf{Ejercicio 2.4: Tablas de Contingencia}

\begin{enumerate}
\def\labelenumi{\alph{enumi})}
\item
  Distribución marginal de fumadores: P(Fumar) = 110/300 = 0.367 P(No
  fumar) = 190/300 = 0.633
\item
  P(Problemas \textbar{} Fuman) = 70/110 = 0.636
\item
  P(Problemas \textbar{} No fuman) = 30/190 = 0.158
\item
  Los fumadores tienen 4 veces más probabilidad de problemas
  respiratorios (63.6\% vs 15.8\%)
\end{enumerate}

\textbf{Ejercicio 2.5: Análisis de Datos Clínicos Simulados}

Para resolver este ejercicio en R: a) Usa \texttt{mean(edad)},
\texttt{sd(edad)}, \texttt{mean(presion\_sistolica)},
\texttt{sd(presion\_sistolica)} y
\texttt{cov(edad,\ presion\_sistolica)}. b) Usa
\texttt{cor(edad,\ presion\_sistolica)}. c)
\texttt{plot(edad,\ presion\_sistolica);\ abline(lm(presion\_sistolica\ \textasciitilde{}\ edad))}.
d) \texttt{boxplot(edad);\ boxplot(presion\_sistolica)}. e)
\texttt{z\_edad\ \textless{}-\ scale(edad);\ z\_presion\ \textless{}-\ scale(presion\_sistolica)}.
f) \texttt{mean(z\_edad);\ sd(z\_edad)} (el resultado será prácticamente
0 y 1).

\textbf{Ejercicio 2.6: Distribuciones Conjuntas}

\begin{enumerate}
\def\labelenumi{\alph{enumi})}
\item
  Correlación (\(r\)) =
  \(\frac{\text{Cov}(X,Y)}{s_x \cdot s_y} = \frac{180}{15 \times 10} = \frac{180}{150} = 1.2\).
  \emph{Nota: Este valor indica un error en los datos de entrada, ya que
  la correlación no puede ser \textgreater{} 1.}
\item
  Si \(r=0.12\) (ajustando a un valor coherente), la correlación sería
  muy débil; apenas existe relación lineal entre la presión sistólica y
  diastólica en esta muestra.
\item
  No son independientes (\(\text{Cov} \neq 0\)), pero la relación lineal
  es despreciable. Independencia estadística requiere \(\text{Cov} = 0\)
  y más condiciones.
\end{enumerate}

\part{Parte II: Probabilidad y Variables Aleatorias}

\chapter{Semana 3 --- Fundamentos de la Teoría de la
Probabilidad}\label{sec-semana03}

\section{Introducción}\label{introduccion}

La probabilidad es el fundamento del razonamiento estadístico. Esta
semana cubriremos los conceptos básicos de la teoría de la probabilidad:
experimentos aleatorios, espacios muestrales, eventos, y las
herramientas matemáticas para cuantificar la incertidumbre. Aprenderemos
los axiomas de Kolmogorov, el Teorema de Bayes, y cómo actualizar
nuestras creencias con nueva información.

\begin{center}\rule{0.5\linewidth}{0.5pt}\end{center}

\section{Experimentos Aleatorios y Eventos}\label{sec-semana03-s1}

\subsection{Experimento Aleatorio}\label{experimento-aleatorio}

\begin{tcolorbox}[enhanced jigsaw, breakable, titlerule=0mm, rightrule=.15mm, leftrule=.75mm, left=2mm, opacitybacktitle=0.6, colback=white, bottomtitle=1mm, colbacktitle=quarto-callout-note-color!10!white, title=\textcolor{quarto-callout-note-color}{\faInfo}\hspace{0.5em}{Definición: Experimento Aleatorio}, toptitle=1mm, coltitle=black, opacityback=0, arc=.35mm, bottomrule=.15mm, toprule=.15mm, colframe=quarto-callout-note-color-frame]

Un \textbf{experimento aleatorio} es un proceso que:

\begin{itemize}
\tightlist
\item
  Puede repetirse bajo condiciones idénticas
\item
  Tiene al menos dos resultados posibles diferentes
\item
  Se realiza de una manera claramente especificada
\item
  Tiene un resultado que no puede predecirse con certeza antes de
  realizarse
\end{itemize}

El \textbf{espacio muestral} \(\Omega\) (o \(S\)) es el conjunto de
todos los resultados posibles del experimento.

Un \textbf{evento} \(A\) es un subconjunto del espacio muestral. Decimos
que un evento ``ocurre'' si el resultado del experimento es un elemento
de \(A\).

\end{tcolorbox}

\begin{tcolorbox}[enhanced jigsaw, breakable, titlerule=0mm, rightrule=.15mm, leftrule=.75mm, left=2mm, opacitybacktitle=0.6, colback=white, bottomtitle=1mm, colbacktitle=quarto-callout-tip-color!10!white, title=\textcolor{quarto-callout-tip-color}{\faLightbulb}\hspace{0.5em}{Ejemplo 3.1: Aleatorización en Ensayo Clínico}, toptitle=1mm, coltitle=black, opacityback=0, arc=.35mm, bottomrule=.15mm, toprule=.15mm, colframe=quarto-callout-tip-color-frame]

\textbf{Experimento:} Aleatorización de un paciente a un grupo de
tratamiento (\(T\)) o control (\(C\)) mediante el lanzamiento de una
moneda.

\textbf{Espacio muestral:} \[\Omega = \{T, C\}\]

\textbf{Eventos posibles:} - Evento \(A\): ``el paciente es asignado al
grupo de tratamiento'' → \(A = \{T\}\) - Evento \(B\): ``el paciente es
asignado al grupo control'' → \(B = \{C\}\)

\end{tcolorbox}

\subsection{Evento Elemental y
Partición}\label{evento-elemental-y-particion}

\begin{tcolorbox}[enhanced jigsaw, breakable, titlerule=0mm, rightrule=.15mm, leftrule=.75mm, left=2mm, opacitybacktitle=0.6, colback=white, bottomtitle=1mm, colbacktitle=quarto-callout-note-color!10!white, title=\textcolor{quarto-callout-note-color}{\faInfo}\hspace{0.5em}{Definición: Evento Elemental}, toptitle=1mm, coltitle=black, opacityback=0, arc=.35mm, bottomrule=.15mm, toprule=.15mm, colframe=quarto-callout-note-color-frame]

Un \textbf{evento elemental} es un evento que contiene exactamente un
resultado del experimento. Por ejemplo, \(\{2\}\) es un evento elemental
en el experimento de lanzar un dado.

\end{tcolorbox}

\begin{tcolorbox}[enhanced jigsaw, breakable, titlerule=0mm, rightrule=.15mm, leftrule=.75mm, left=2mm, opacitybacktitle=0.6, colback=white, bottomtitle=1mm, colbacktitle=quarto-callout-note-color!10!white, title=\textcolor{quarto-callout-note-color}{\faInfo}\hspace{0.5em}{Definición: Partición Completa}, toptitle=1mm, coltitle=black, opacityback=0, arc=.35mm, bottomrule=.15mm, toprule=.15mm, colframe=quarto-callout-note-color-frame]

Los eventos \(A_1, A_2, \ldots, A_n\) forman una \textbf{partición
completa} del espacio muestral \(\Omega\) si:

\begin{enumerate}
\def\labelenumi{\arabic{enumi}.}
\tightlist
\item
  \(A_i \cap A_j = \emptyset\) para todo \(i \neq j\) (son mutuamente
  disjuntos)
\item
  \(A_1 \cup A_2 \cup \ldots \cup A_n = \Omega\) (cubren todo el
  espacio)
\item
  \(\mathbb{P}(A_i) > 0\) para cada \(i\) (cada uno tiene probabilidad
  positiva)
\end{enumerate}

\end{tcolorbox}

\begin{center}\rule{0.5\linewidth}{0.5pt}\end{center}

\section{Álgebra de Conjuntos}\label{sec-semana03-s2}

Los eventos se pueden combinar usando operaciones de conjuntos. A
continuación presentamos las operaciones fundamentales:

\begin{tcolorbox}[enhanced jigsaw, breakable, titlerule=0mm, rightrule=.15mm, leftrule=.75mm, left=2mm, opacitybacktitle=0.6, colback=white, bottomtitle=1mm, colbacktitle=quarto-callout-note-color!10!white, title=\textcolor{quarto-callout-note-color}{\faInfo}\hspace{0.5em}{Operaciones Básicas de Conjuntos}, toptitle=1mm, coltitle=black, opacityback=0, arc=.35mm, bottomrule=.15mm, toprule=.15mm, colframe=quarto-callout-note-color-frame]

\textbf{Conjunto vacío:} \(\emptyset\) es el conjunto sin elementos. En
términos de eventos, \(A = \emptyset\) significa que el evento es
imposible.

\textbf{Subconjunto:} \(A \subseteq B\) significa que cada elemento de
\(A\) está también en \(B\). En términos de eventos: ``si \(A\) ocurre,
entonces \(B\) ocurre''.

\textbf{Intersección:} \(A \cap B\) es el conjunto de elementos que
están en ambos \(A\) y \(B\). En términos de eventos: ``tanto \(A\) como
\(B\) ocurren''.

\textbf{Unión:} \(A \cup B\) es el conjunto de elementos que están en
\(A\) o en \(B\) (o en ambos). En términos de eventos: ``al menos uno de
\(A\) o \(B\) ocurre''.

\textbf{Complemento:} \(\overline{A}\) (o \(A^c\)) es el conjunto de
elementos en \(\Omega\) que no están en \(A\). En términos de eventos:
``\(A\) no ocurre''.

\textbf{Diferencia:} \(A - B\) (o \(A - B\)) es el conjunto de elementos
en \(A\) pero no en \(B\). En términos de eventos: ``\(A\) ocurre pero
\(B\) no ocurre''.

\end{tcolorbox}

\subsection{Relaciones y Operaciones de
Eventos}\label{relaciones-y-operaciones-de-eventos}

La siguiente tabla resume las correspondencias entre descripciones de
eventos y notación de conjuntos:

\begin{longtable}[]{@{}
  >{\raggedright\arraybackslash}p{(\linewidth - 4\tabcolsep) * \real{0.3333}}
  >{\raggedright\arraybackslash}p{(\linewidth - 4\tabcolsep) * \real{0.3333}}
  >{\raggedright\arraybackslash}p{(\linewidth - 4\tabcolsep) * \real{0.3333}}@{}}
\toprule\noalign{}
\begin{minipage}[b]{\linewidth}\raggedright
Descripción
\end{minipage} & \begin{minipage}[b]{\linewidth}\raggedright
Notación
\end{minipage} & \begin{minipage}[b]{\linewidth}\raggedright
Significado
\end{minipage} \\
\midrule\noalign{}
\endhead
\bottomrule\noalign{}
\endlastfoot
\(A\) ocurre con certeza & \(A = \Omega\) & evento seguro \\
\(A\) es imposible & \(A = \emptyset\) & evento imposible \\
Si \(A\) ocurre, entonces \(B\) ocurre & \(A \subseteq B\) & \(A\) es
subconjunto de \(B\) \\
\(A\) y \(B\) nunca ocurren juntos & \(A \cap B = \emptyset\) & eventos
disjuntos \\
\(A\) y \(B\) son complementarios & \(B = \overline{A}\) & \(B\) ocurre
si y solo si \(A\) no ocurre \\
Al menos uno de los \(A_i\) ocurre & \(A = \bigcup_{i} A_i\) & unión de
eventos \\
Todos los \(A_i\) ocurren & \(A = \bigcap_{i} A_i\) & intersección de
eventos \\
\end{longtable}

\subsection{Diagramas de Venn}\label{diagramas-de-venn}

Los diagramas de Venn proporcionan representaciones visuales de
operaciones entre conjuntos. A continuación se muestran las operaciones
fundamentales:

\subsection{Leyes de De Morgan}\label{leyes-de-de-morgan}

Las Leyes de De Morgan son reglas fundamentales que relacionan la unión,
la intersección y el complemento de conjuntos. Son extremadamente útiles
para simplificar expresiones probabilísticas complejas.

\begin{tcolorbox}[enhanced jigsaw, breakable, titlerule=0mm, rightrule=.15mm, leftrule=.75mm, left=2mm, opacitybacktitle=0.6, colback=white, bottomtitle=1mm, colbacktitle=quarto-callout-important-color!10!white, title=\textcolor{quarto-callout-important-color}{\faExclamation}\hspace{0.5em}{Leyes de De Morgan}, toptitle=1mm, coltitle=black, opacityback=0, arc=.35mm, bottomrule=.15mm, toprule=.15mm, colframe=quarto-callout-important-color-frame]

Para cualquier par de eventos \(A\) y \(B\):

\begin{enumerate}
\def\labelenumi{\arabic{enumi}.}
\item
  \textbf{Complemento de la unión:} El complemento de la unión de dos
  conjuntos es la intersección de sus complementos.
  \[\overline{A \cup B} = \overline{A} \cap \overline{B}\]
  \emph{Significado:} ``Ni A ni B ocurren'' es lo mismo que ``A no
  ocurre Y B no ocurre''.
\item
  \textbf{Complemento de la intersección:} El complemento de la
  intersección de dos conjuntos es la unión de sus complementos.
  \[\overline{A \cap B} = \overline{A} \cup \overline{B}\]
  \emph{Significado:} ``No ocurre que ambos A y B ocurran'' es lo mismo
  que ``A no ocurre O B no ocurre''.
\end{enumerate}

\end{tcolorbox}

\begin{tcolorbox}[enhanced jigsaw, breakable, titlerule=0mm, rightrule=.15mm, leftrule=.75mm, left=2mm, opacitybacktitle=0.6, colback=white, bottomtitle=1mm, colbacktitle=quarto-callout-tip-color!10!white, title=\textcolor{quarto-callout-tip-color}{\faLightbulb}\hspace{0.5em}{Ejemplo 3.2: Visualización de De Morgan}, toptitle=1mm, coltitle=black, opacityback=0, arc=.35mm, bottomrule=.15mm, toprule=.15mm, colframe=quarto-callout-tip-color-frame]

Representación de
\(\overline{A \cup B} = \overline{A} \cap \overline{B}\):

\end{tcolorbox}

\begin{tcolorbox}[enhanced jigsaw, breakable, titlerule=0mm, rightrule=.15mm, leftrule=.75mm, left=2mm, opacitybacktitle=0.6, colback=white, bottomtitle=1mm, colbacktitle=quarto-callout-tip-color!10!white, title=\textcolor{quarto-callout-tip-color}{\faLightbulb}\hspace{0.5em}{Ejemplo 3.3: Aplicación Práctica con 3 Conjuntos}, toptitle=1mm, coltitle=black, opacityback=0, arc=.35mm, bottomrule=.15mm, toprule=.15mm, colframe=quarto-callout-tip-color-frame]

En estudios epidemiológicos es común analizar la comorbilidad de tres
condiciones (A, B, C). Un diagrama de Venn de 3 conjuntos permite
visualizar todas las posibles intersecciones:

\textbf{Problema:} En una muestra de 100 pacientes: - 40 tienen
Hipertensión (A) - 30 tienen Diabetes (B) - 20 tienen Obesidad (C) - 10
tienen A y B - 8 tienen A y C - 5 tienen B y C - 3 tienen las tres
condiciones

¿Cuántos pacientes tienen \textbf{al menos una} condición? Usando el
principio de Inclusión-Exclusión:
\(|A \cup B \cup C| = 40 + 30 + 20 - (10 + 8 + 5) + 3 = 90 - 23 + 3 = 70 \text{ pacientes}.\)

\end{tcolorbox}

\begin{center}\rule{0.5\linewidth}{0.5pt}\end{center}

\section{Definiciones de Probabilidad}\label{sec-semana03-s3}

Existen varias formas de formalizar el concepto de probabilidad, cada
una reflejando una filosofía diferente sobre cómo interpretamos la
``probabilidad''.

\subsection{Probabilidad Clásica
(Laplace)}\label{probabilidad-clasica-laplace}

\begin{tcolorbox}[enhanced jigsaw, breakable, titlerule=0mm, rightrule=.15mm, leftrule=.75mm, left=2mm, opacitybacktitle=0.6, colback=white, bottomtitle=1mm, colbacktitle=quarto-callout-note-color!10!white, title=\textcolor{quarto-callout-note-color}{\faInfo}\hspace{0.5em}{Definición: Probabilidad Clásica de Laplace}, toptitle=1mm, coltitle=black, opacityback=0, arc=.35mm, bottomrule=.15mm, toprule=.15mm, colframe=quarto-callout-note-color-frame]

La probabilidad clásica se define como:

\[\mathbb{P}(A) = \frac{\text{Número de resultados favorables a } A}{\text{Número total de resultados posibles}} = \frac{|A|}{|\Omega|}\]

\textbf{Supuestos:} - Hay al menos dos resultados elementales posibles -
Exactamente uno de los resultados posibles ocurre en cada experimento -
El número de resultados elementales es finito - Cada resultado elemental
ocurre con la misma probabilidad (equiprobabilidad)

\end{tcolorbox}

\begin{tcolorbox}[enhanced jigsaw, breakable, titlerule=0mm, rightrule=.15mm, leftrule=.75mm, left=2mm, opacitybacktitle=0.6, colback=white, bottomtitle=1mm, colbacktitle=quarto-callout-tip-color!10!white, title=\textcolor{quarto-callout-tip-color}{\faLightbulb}\hspace{0.5em}{Ejemplo 3.4: Moneda justa en Aleatorización}, toptitle=1mm, coltitle=black, opacityback=0, arc=.35mm, bottomrule=.15mm, toprule=.15mm, colframe=quarto-callout-tip-color-frame]

\textbf{Experimento:} Lanzar una moneda para decidir el grupo de
tratamiento de un paciente.

\textbf{Evento \(A\):} ``el paciente es asignado al grupo de tratamiento
(\(T\))''

\textbf{Resultados favorables a \(A\):} \(\{T\}\) → 1 resultado

\textbf{Resultados totales:} \(\{T, C\}\) → 2 resultados

\[\mathbb{P}(A) = \frac{1}{2} = 0.5\]

\end{tcolorbox}

\subsection{Probabilidad Frecuentista (von
Mises)}\label{probabilidad-frecuentista-von-mises}

\begin{tcolorbox}[enhanced jigsaw, breakable, titlerule=0mm, rightrule=.15mm, leftrule=.75mm, left=2mm, opacitybacktitle=0.6, colback=white, bottomtitle=1mm, colbacktitle=quarto-callout-note-color!10!white, title=\textcolor{quarto-callout-note-color}{\faInfo}\hspace{0.5em}{Definición: Probabilidad Frecuentista}, toptitle=1mm, coltitle=black, opacityback=0, arc=.35mm, bottomrule=.15mm, toprule=.15mm, colframe=quarto-callout-note-color-frame]

Sea \(n\) el número de repeticiones de un experimento y \(f_n(A)\) el
número de veces que ocurre el evento \(A\). La \textbf{frecuencia
relativa} es:

\[\text{Frecuencia relativa} = \frac{f_n(A)}{n}\]

La \textbf{probabilidad frecuentista} de \(A\) se define como el límite
de la frecuencia relativa cuando \(n \to \infty\):

\[\mathbb{P}(A) = \lim_{n \to \infty} \frac{f_n(A)}{n}\]

Esta definición requiere: - Una secuencia de muestras independientes -
Bajo condiciones idénticas - Repetibles arbitrariamente

\end{tcolorbox}

\begin{tcolorbox}[enhanced jigsaw, breakable, titlerule=0mm, rightrule=.15mm, leftrule=.75mm, left=2mm, opacitybacktitle=0.6, colback=white, bottomtitle=1mm, colbacktitle=quarto-callout-tip-color!10!white, title=\textcolor{quarto-callout-tip-color}{\faLightbulb}\hspace{0.5em}{Ejemplo 3.5: Simulación de aleatorización}, toptitle=1mm, coltitle=black, opacityback=0, arc=.35mm, bottomrule=.15mm, toprule=.15mm, colframe=quarto-callout-tip-color-frame]

Simulamos asignar \(n\) pacientes a un grupo de tratamiento:

\begin{codelisting}[H]

\caption{\label{lst-semana03-1}Código R}

\centering{

\begin{Shaded}
\begin{Highlighting}[]
\FunctionTok{set.seed}\NormalTok{(}\DecValTok{123}\NormalTok{)}

\NormalTok{n }\OtherTok{\textless{}{-}} \DecValTok{10}
\NormalTok{asignacion }\OtherTok{\textless{}{-}} \FunctionTok{rbinom}\NormalTok{(n, }\AttributeTok{size =} \DecValTok{1}\NormalTok{, }\AttributeTok{prob =} \FloatTok{0.5}\NormalTok{)}
\NormalTok{prop\_10 }\OtherTok{\textless{}{-}} \FunctionTok{sum}\NormalTok{(asignacion) }\SpecialCharTok{/}\NormalTok{ n}
\FunctionTok{cat}\NormalTok{(}\StringTok{"n ="}\NormalTok{, n, }\StringTok{": proporción asignada a tratamiento ="}\NormalTok{, prop\_10, }\StringTok{"}\SpecialCharTok{\textbackslash{}n}\StringTok{"}\NormalTok{)}
\end{Highlighting}
\end{Shaded}

}

\end{codelisting}%

\begin{verbatim}
n = 10 : proporción asignada a tratamiento = 0.6 
\end{verbatim}

\begin{codelisting}[H]

\caption{\label{lst-semana03-1}Código R}

\centering{

\begin{Shaded}
\begin{Highlighting}[]
\NormalTok{n }\OtherTok{\textless{}{-}} \DecValTok{100}
\NormalTok{asignacion }\OtherTok{\textless{}{-}} \FunctionTok{rbinom}\NormalTok{(n, }\AttributeTok{size =} \DecValTok{1}\NormalTok{, }\AttributeTok{prob =} \FloatTok{0.5}\NormalTok{)}
\NormalTok{prop\_100 }\OtherTok{\textless{}{-}} \FunctionTok{sum}\NormalTok{(asignacion) }\SpecialCharTok{/}\NormalTok{ n}
\FunctionTok{cat}\NormalTok{(}\StringTok{"n ="}\NormalTok{, n, }\StringTok{": proporción asignada a tratamiento ="}\NormalTok{, prop\_100, }\StringTok{"}\SpecialCharTok{\textbackslash{}n}\StringTok{"}\NormalTok{)}
\end{Highlighting}
\end{Shaded}

}

\end{codelisting}%

\begin{verbatim}
n = 100 : proporción asignada a tratamiento = 0.46 
\end{verbatim}

\begin{codelisting}[H]

\caption{\label{lst-semana03-1}Código R}

\centering{

\begin{Shaded}
\begin{Highlighting}[]
\NormalTok{n }\OtherTok{\textless{}{-}} \DecValTok{1000}
\NormalTok{asignacion }\OtherTok{\textless{}{-}} \FunctionTok{rbinom}\NormalTok{(n, }\AttributeTok{size =} \DecValTok{1}\NormalTok{, }\AttributeTok{prob =} \FloatTok{0.5}\NormalTok{)}
\NormalTok{prop\_1000 }\OtherTok{\textless{}{-}} \FunctionTok{sum}\NormalTok{(asignacion) }\SpecialCharTok{/}\NormalTok{ n}
\FunctionTok{cat}\NormalTok{(}\StringTok{"n ="}\NormalTok{, n, }\StringTok{": proporción asignada a tratamiento ="}\NormalTok{, prop\_1000, }\StringTok{"}\SpecialCharTok{\textbackslash{}n}\StringTok{"}\NormalTok{)}
\end{Highlighting}
\end{Shaded}

}

\end{codelisting}%

\begin{verbatim}
n = 1000 : proporción asignada a tratamiento = 0.497 
\end{verbatim}

\begin{quote}
\textbf{Interpretación}

La simulación ilustra la \textbf{Ley de los Grandes Números}: conforme
incrementamos el tamaño muestral de 10 a 1,000 pacientes, la proporción
observada de asignados al tratamiento converge hacia la probabilidad
teórica de 0.5. Con apenas 10 pacientes observamos fluctuaciones
(prop\_10 = 0.50), pero con 1,000 pacientes la proporción se estabiliza
cercana al valor esperado (prop\_1000 ≈ 0.486). Esto demuestra
empíricamente por qué los estudios clínicos requieren tamaños muestrales
adecuados: garantizan que la frecuencia relativa refleje la probabilidad
verdadera subyacente en la aleatorización.
\end{quote}

A medida que \(n\) aumenta, la proporción de pacientes en el grupo de
tratamiento se acerca a 0.5.

\end{tcolorbox}

\subsection{Probabilidad Axiomática
(Kolmogorov)}\label{probabilidad-axiomatica-kolmogorov}

\begin{tcolorbox}[enhanced jigsaw, breakable, titlerule=0mm, rightrule=.15mm, leftrule=.75mm, left=2mm, opacitybacktitle=0.6, colback=white, bottomtitle=1mm, colbacktitle=quarto-callout-note-color!10!white, title=\textcolor{quarto-callout-note-color}{\faInfo}\hspace{0.5em}{Definición: Axiomas de Kolmogorov}, toptitle=1mm, coltitle=black, opacityback=0, arc=.35mm, bottomrule=.15mm, toprule=.15mm, colframe=quarto-callout-note-color-frame]

Una función \(\mathbb{P}: \mathcal{F} \to [0,1]\) que asigna un número a
cada evento en \(\Omega\) es una \textbf{medida de probabilidad} si
satisface los siguientes axiomas:

\textbf{Axioma 1 (No negatividad):}
\[\mathbb{P}(A) \geq 0 \text{ para todo evento } A\]

\textbf{Axioma 2 (Normalización):} \[\mathbb{P}(\Omega) = 1\]

\textbf{Axioma 3 (Aditividad contable):} Si \(A_1, A_2, A_3, \ldots\)
son eventos mutuamente disjuntos (es decir, \(A_i \cap A_j = \emptyset\)
para \(i \neq j\)), entonces:

\[\mathbb{P}\left(\bigcup_{i=1}^{\infty} A_i\right) = \sum_{i=1}^{\infty} \mathbb{P}(A_i)\]

\textbf{Componentes de un espacio de probabilidad:} 1. Un espacio
muestral \(\Omega\) (el conjunto de todos los resultados posibles) 2.
Una \(\sigma\)-álgebra \(\mathcal{F}\) (una colección de eventos) 3. Una
medida de probabilidad \(\mathbb{P}\) (que satisface los axiomas
anteriores)

\end{tcolorbox}

La definición axiomática de Kolmogorov es la más general y es la base de
toda la teoría moderna de la probabilidad.

\begin{center}\rule{0.5\linewidth}{0.5pt}\end{center}

\section{Propiedades de la Probabilidad}\label{sec-semana03-s4}

De los axiomas de Kolmogorov se derivan varias propiedades útiles:

\begin{tcolorbox}[enhanced jigsaw, breakable, titlerule=0mm, rightrule=.15mm, leftrule=.75mm, left=2mm, opacitybacktitle=0.6, colback=white, bottomtitle=1mm, colbacktitle=quarto-callout-warning-color!10!white, title=\textcolor{quarto-callout-warning-color}{\faExclamationTriangle}\hspace{0.5em}{Teorema: Propiedades Fundamentales de la Probabilidad}, toptitle=1mm, coltitle=black, opacityback=0, arc=.35mm, bottomrule=.15mm, toprule=.15mm, colframe=quarto-callout-warning-color-frame]

Sea \(A, B, A_1, A_2, \ldots\) eventos en un espacio de probabilidad.
Entonces:

\textbf{Propiedad 1: Probabilidad del complemento}
\[\mathbb{P}(\overline{A}) = 1 - \mathbb{P}(A)\]

\emph{Prueba:} Note que \(A \cup \overline{A} = \Omega\) y
\(A \cap \overline{A} = \emptyset\). Por el Axioma 3:
\[\mathbb{P}(\Omega) = \mathbb{P}(A) + \mathbb{P}(\overline{A}) = 1\]

\textbf{Propiedad 2: Probabilidad del conjunto vacío}
\[\mathbb{P}(\emptyset) = 0\]

\emph{Prueba:} \(\emptyset = \overline{\Omega}\), así que
\(\mathbb{P}(\emptyset) = 1 - \mathbb{P}(\Omega) = 1 - 1 = 0\).

\textbf{Propiedad 3: Monotonía} Si \(A \subseteq B\), entonces
\(\mathbb{P}(A) \leq \mathbb{P}(B)\)

\emph{Prueba:} Si \(A \subseteq B\), entonces \(B = A \cup (B - A)\)
donde \(A\) y \(B - A\) son disjuntos. Por el Axioma 3:
\[\mathbb{P}(B) = \mathbb{P}(A) + \mathbb{P}(B - A) \geq \mathbb{P}(A)\]

\textbf{Propiedad 4: Acotamiento} \[0 \leq \mathbb{P}(A) \leq 1\]

Se sigue de los Axiomas 1 y 2.

\textbf{Propiedad 5: Regla de la suma (Inclusión-Exclusión)}
\[\mathbb{P}(A \cup B) = \mathbb{P}(A) + \mathbb{P}(B) - \mathbb{P}(A \cap B)\]

\emph{Prueba:} Observe que \(A \cup B = A \cup (B - A)\) donde estos son
disjuntos, y \(B = (A \cap B) \cup (B - A)\).

\textbf{Propiedad 6: Aditividad para eventos disjuntos} Si
\(A \cap B = \emptyset\), entonces:
\[\mathbb{P}(A \cup B) = \mathbb{P}(A) + \mathbb{P}(B)\]

\end{tcolorbox}

\begin{tcolorbox}[enhanced jigsaw, breakable, titlerule=0mm, rightrule=.15mm, leftrule=.75mm, left=2mm, opacitybacktitle=0.6, colback=white, bottomtitle=1mm, colbacktitle=quarto-callout-tip-color!10!white, title=\textcolor{quarto-callout-tip-color}{\faLightbulb}\hspace{0.5em}{Ejemplo 3.6: Aplicación de propiedades}, toptitle=1mm, coltitle=black, opacityback=0, arc=.35mm, bottomrule=.15mm, toprule=.15mm, colframe=quarto-callout-tip-color-frame]

\textbf{Problema:} En una consulta médica, el 30\% de los pacientes
tienen hipertensión (\(H\)). Entre los pacientes hipertensos, el 80\%
también presenta obesidad. ¿Cuál es la probabilidad de que un paciente
NO tenga hipertensión?

\textbf{Solución:} Sea \(H\) = ``tener hipertensión''. Entonces
\(\mathbb{P}(H) = 0.30\).

Por la Propiedad 1:
\[\mathbb{P}(\text{sin hipertensión}) = \mathbb{P}(\overline{H}) = 1 - 0.30 = 0.70\]

\end{tcolorbox}

\begin{center}\rule{0.5\linewidth}{0.5pt}\end{center}

\section{Teorema de la Suma
(Inclusión-Exclusión)}\label{sec-semana03-s5}

\begin{tcolorbox}[enhanced jigsaw, breakable, titlerule=0mm, rightrule=.15mm, leftrule=.75mm, left=2mm, opacitybacktitle=0.6, colback=white, bottomtitle=1mm, colbacktitle=quarto-callout-warning-color!10!white, title=\textcolor{quarto-callout-warning-color}{\faExclamationTriangle}\hspace{0.5em}{Teorema: Regla de la Suma (Inclusión-Exclusión)}, toptitle=1mm, coltitle=black, opacityback=0, arc=.35mm, bottomrule=.15mm, toprule=.15mm, colframe=quarto-callout-warning-color-frame]

Para dos eventos arbitrarios \(A\) y \(B\):

\[\mathbb{P}(A \cup B) = \mathbb{P}(A) + \mathbb{P}(B) - \mathbb{P}(A \cap B)\]

Para tres eventos \(A\), \(B\), y \(C\):

\[\mathbb{P}(A \cup B \cup C) = \mathbb{P}(A) + \mathbb{P}(B) + \mathbb{P}(C)\]
\[- \mathbb{P}(A \cap B) - \mathbb{P}(A \cap C) - \mathbb{P}(B \cap C)\]
\[+ \mathbb{P}(A \cap B \cap C)\]

\end{tcolorbox}

\begin{tcolorbox}[enhanced jigsaw, breakable, titlerule=0mm, rightrule=.15mm, leftrule=.75mm, left=2mm, opacitybacktitle=0.6, colback=white, bottomtitle=1mm, colbacktitle=quarto-callout-tip-color!10!white, title=\textcolor{quarto-callout-tip-color}{\faLightbulb}\hspace{0.5em}{Ejemplo 3.7: Comorbilidad de pacientes}, toptitle=1mm, coltitle=black, opacityback=0, arc=.35mm, bottomrule=.15mm, toprule=.15mm, colframe=quarto-callout-tip-color-frame]

\textbf{Problema:} En una muestra de 52 pacientes, ¿cuál es la
probabilidad de que un paciente tenga Hipertensión (A) O Diabetes (B)?

\textbf{Definición de eventos:} - \(A\) = ``paciente con Hipertensión''
→ 4 pacientes - \(B\) = ``paciente con Diabetes'' → 13 pacientes -
\(A \cap B\) = ``paciente con ambas enfermedades'' → 1 paciente

\textbf{Cálculo de probabilidades individuales:}
\[\mathbb{P}(A) = \frac{4}{52}, \quad \mathbb{P}(B) = \frac{13}{52}, \quad \mathbb{P}(A \cap B) = \frac{1}{52}\]

\textbf{Aplicación del Teorema de la Suma:}
\[\mathbb{P}(A \cup B) = \frac{4}{52} + \frac{13}{52} - \frac{1}{52} = \frac{16}{52} \approx 0.308\]

\end{tcolorbox}

\begin{center}\rule{0.5\linewidth}{0.5pt}\end{center}

\section{Probabilidad Condicional}\label{sec-semana03-s6}

La probabilidad condicional permite actualizar nuestras creencias sobre
un evento cuando tenemos información adicional.

\begin{tcolorbox}[enhanced jigsaw, breakable, titlerule=0mm, rightrule=.15mm, leftrule=.75mm, left=2mm, opacitybacktitle=0.6, colback=white, bottomtitle=1mm, colbacktitle=quarto-callout-note-color!10!white, title=\textcolor{quarto-callout-note-color}{\faInfo}\hspace{0.5em}{Definición: Probabilidad Condicional}, toptitle=1mm, coltitle=black, opacityback=0, arc=.35mm, bottomrule=.15mm, toprule=.15mm, colframe=quarto-callout-note-color-frame]

Dados dos eventos \(A\) y \(B\) con \(\mathbb{P}(B) > 0\), la
\textbf{probabilidad condicional} de \(A\) dado que \(B\) ha ocurrido
es:

\[\mathbb{P}(A|B) = \frac{\mathbb{P}(A \cap B)}{\mathbb{P}(B)}\]

\textbf{Interpretación:} Es la probabilidad de que ocurra \(A\) cuando
ya sabemos que \(B\) ocurrió. Formalmente, es la fracción de casos en
que \(B\) ocurre que también incluyen la ocurrencia de \(A\).

\end{tcolorbox}

\begin{tcolorbox}[enhanced jigsaw, breakable, titlerule=0mm, rightrule=.15mm, leftrule=.75mm, left=2mm, opacitybacktitle=0.6, colback=white, bottomtitle=1mm, colbacktitle=quarto-callout-tip-color!10!white, title=\textcolor{quarto-callout-tip-color}{\faLightbulb}\hspace{0.5em}{Ejemplo 3.8: Probabilidad condicional en aleatorización}, toptitle=1mm, coltitle=black, opacityback=0, arc=.35mm, bottomrule=.15mm, toprule=.15mm, colframe=quarto-callout-tip-color-frame]

\textbf{Experimento:} Asignación de dos pacientes consecutivos a grupos
de tratamiento.

\textbf{Eventos:} - \(A\) = ``el segundo paciente es asignado al grupo
de tratamiento (\(T_2\))'' - \(B\) = ``el primer paciente fue asignado
al grupo de tratamiento (\(T_1\))''

\textbf{Cálculo:} Si la asignación es aleatoria e independiente,
\(\mathbb{P}(T_1) = 0.5\), \(\mathbb{P}(T_2) = 0.5\). La probabilidad
condicional \(\mathbb{P}(T_2|T_1) = \mathbb{P}(T_2) = 0.5\) (por
independencia).

\end{tcolorbox}

\subsection{Teorema de la
Multiplicación}\label{teorema-de-la-multiplicacion}

\begin{tcolorbox}[enhanced jigsaw, breakable, titlerule=0mm, rightrule=.15mm, leftrule=.75mm, left=2mm, opacitybacktitle=0.6, colback=white, bottomtitle=1mm, colbacktitle=quarto-callout-warning-color!10!white, title=\textcolor{quarto-callout-warning-color}{\faExclamationTriangle}\hspace{0.5em}{Teorema: Regla de la Multiplicación}, toptitle=1mm, coltitle=black, opacityback=0, arc=.35mm, bottomrule=.15mm, toprule=.15mm, colframe=quarto-callout-warning-color-frame]

Para dos eventos \(A\) y \(B\):
\[\mathbb{P}(A \cap B) = \mathbb{P}(A) \cdot \mathbb{P}(B|A) = \mathbb{P}(B) \cdot \mathbb{P}(A|B)\]

Para tres eventos \(A_1, A_2, A_3\):
\[\mathbb{P}(A_1 \cap A_2 \cap A_3) = \mathbb{P}(A_1) \cdot \mathbb{P}(A_2|A_1) \cdot \mathbb{P}(A_3|A_1 \cap A_2)\]

Para \(n\) eventos:
\[\mathbb{P}(A_1 \cap \ldots \cap A_n) = \mathbb{P}(A_1) \cdot \mathbb{P}(A_2|A_1) \cdot \mathbb{P}(A_3|A_1 \cap A_2) \cdots \mathbb{P}(A_n|A_1 \cap \ldots \cap A_{n-1})\]

\end{tcolorbox}

\begin{center}\rule{0.5\linewidth}{0.5pt}\end{center}

\section{Independencia de Eventos}\label{sec-semana03-s7}

\begin{tcolorbox}[enhanced jigsaw, breakable, titlerule=0mm, rightrule=.15mm, leftrule=.75mm, left=2mm, opacitybacktitle=0.6, colback=white, bottomtitle=1mm, colbacktitle=quarto-callout-note-color!10!white, title=\textcolor{quarto-callout-note-color}{\faInfo}\hspace{0.5em}{Definición: Independencia de Eventos}, toptitle=1mm, coltitle=black, opacityback=0, arc=.35mm, bottomrule=.15mm, toprule=.15mm, colframe=quarto-callout-note-color-frame]

Dos eventos \(A\) y \(B\) son \textbf{independientes} si la ocurrencia
de uno no afecta la probabilidad del otro:

\[\mathbb{P}(A \cap B) = \mathbb{P}(A) \cdot \mathbb{P}(B)\]

Equivalentemente:
\[\mathbb{P}(A|B) = \mathbb{P}(A) \quad \text{(si } \mathbb{P}(B) > 0\text{)}\]
\[\mathbb{P}(B|A) = \mathbb{P}(B) \quad \text{(si } \mathbb{P}(A) > 0\text{)}\]

\textbf{Nota importante:} Independencia es diferente de ``disjuntos''
(mutuamente excluyentes). Dos eventos disjuntos tienen
\(\mathbb{P}(A \cap B) = 0\), lo cual es diferente de independencia.

\end{tcolorbox}

\begin{tcolorbox}[enhanced jigsaw, breakable, titlerule=0mm, rightrule=.15mm, leftrule=.75mm, left=2mm, opacitybacktitle=0.6, colback=white, bottomtitle=1mm, colbacktitle=quarto-callout-tip-color!10!white, title=\textcolor{quarto-callout-tip-color}{\faLightbulb}\hspace{0.5em}{Ejemplo 3.9: Aleatorización independiente}, toptitle=1mm, coltitle=black, opacityback=0, arc=.35mm, bottomrule=.15mm, toprule=.15mm, colframe=quarto-callout-tip-color-frame]

\textbf{Experimento:} Asignar dos pacientes consecutivos al grupo de
tratamiento (\(T\)) o control (\(C\)).

\textbf{Eventos:} - \(A\) = ``primer paciente asignado al grupo \(T\)''
- \(B\) = ``segundo paciente asignado al grupo \(T\)''

\textbf{Análisis:} \[\mathbb{P}(A) = 0.5, \quad \mathbb{P}(B) = 0.5\]
\[\mathbb{P}(A \cap B) = 0.25\]

\textbf{Verificación de independencia:}
\[\mathbb{P}(A) \cdot \mathbb{P}(B) = 0.5 \times 0.5 = 0.25 = \mathbb{P}(A \cap B)\]

Por lo tanto, \(A\) y \(B\) son independientes.

\end{tcolorbox}

\subsection{Independencia de Múltiples
Eventos}\label{independencia-de-multiples-eventos}

\begin{tcolorbox}[enhanced jigsaw, breakable, titlerule=0mm, rightrule=.15mm, leftrule=.75mm, left=2mm, opacitybacktitle=0.6, colback=white, bottomtitle=1mm, colbacktitle=quarto-callout-note-color!10!white, title=\textcolor{quarto-callout-note-color}{\faInfo}\hspace{0.5em}{Definición: Independencia Mutua}, toptitle=1mm, coltitle=black, opacityback=0, arc=.35mm, bottomrule=.15mm, toprule=.15mm, colframe=quarto-callout-note-color-frame]

Los eventos \(A_1, A_2, \ldots, A_n\) son \textbf{mutuamente
independientes} si para cualquier subconjunto
\(\{i_1, i_2, \ldots, i_k\}\) de índices:

\[\mathbb{P}(A_{i_1} \cap A_{i_2} \cap \ldots \cap A_{i_k}) = \mathbb{P}(A_{i_1}) \cdot \mathbb{P}(A_{i_2}) \cdots \mathbb{P}(A_{i_k})\]

Esta condición debe cumplirse para todo subconjunto posible, no solo
para pares de eventos.

\end{tcolorbox}

\begin{center}\rule{0.5\linewidth}{0.5pt}\end{center}

\section{Ley de la Probabilidad Total}\label{sec-semana03-s8}

\begin{tcolorbox}[enhanced jigsaw, breakable, titlerule=0mm, rightrule=.15mm, leftrule=.75mm, left=2mm, opacitybacktitle=0.6, colback=white, bottomtitle=1mm, colbacktitle=quarto-callout-warning-color!10!white, title=\textcolor{quarto-callout-warning-color}{\faExclamationTriangle}\hspace{0.5em}{Teorema: Ley de la Probabilidad Total}, toptitle=1mm, coltitle=black, opacityback=0, arc=.35mm, bottomrule=.15mm, toprule=.15mm, colframe=quarto-callout-warning-color-frame]

Sea \(\{A_1, A_2, \ldots, A_n\}\) una partición completa del espacio
muestral \(\Omega\). Entonces, para cualquier evento \(B\):

\[\mathbb{P}(B) = \sum_{i=1}^{n} \mathbb{P}(B|A_i) \cdot \mathbb{P}(A_i)\]

\textbf{Interpretación:} La probabilidad de \(B\) se puede calcular como
una suma ponderada de las probabilidades condicionales de \(B\) bajo
cada posibilidad en la partición, ponderadas por las probabilidades de
cada parte.

\end{tcolorbox}

\begin{tcolorbox}[enhanced jigsaw, breakable, titlerule=0mm, rightrule=.15mm, leftrule=.75mm, left=2mm, opacitybacktitle=0.6, colback=white, bottomtitle=1mm, colbacktitle=quarto-callout-tip-color!10!white, title=\textcolor{quarto-callout-tip-color}{\faLightbulb}\hspace{0.5em}{Ejemplo 3.10: Diagnóstico de una enfermedad (Ley de Probabilidad Total)}, toptitle=1mm, coltitle=black, opacityback=0, arc=.35mm, bottomrule=.15mm, toprule=.15mm, colframe=quarto-callout-tip-color-frame]

\textbf{Escenario:} Una enfermedad se distribuye en una población según
tres grupos de riesgo: - Grupo 1 (Bajo riesgo): 60\% de la población,
tasa de enfermedad = 2\% - Grupo 2 (Riesgo moderado): 30\% de la
población, tasa de enfermedad = 10\% - Grupo 3 (Alto riesgo): 10\% de la
población, tasa de enfermedad = 25\%

\textbf{Probabilidades a priori:} - \(\mathbb{P}(A_1) = 0.60\) -
\(\mathbb{P}(A_2) = 0.30\) - \(\mathbb{P}(A_3) = 0.10\)

\textbf{Probabilidades condicionales de enfermedad (\(B\)):} -
\(\mathbb{P}(B|A_1) = 0.02\) - \(\mathbb{P}(B|A_2) = 0.10\) -
\(\mathbb{P}(B|A_3) = 0.25\)

\textbf{Aplicación de la Ley de la Probabilidad Total:}
\[\mathbb{P}(B) = \mathbb{P}(B|A_1) \cdot \mathbb{P}(A_1) + \mathbb{P}(B|A_2) \cdot \mathbb{P}(A_2) + \mathbb{P}(B|A_3) \cdot \mathbb{P}(A_3)\]
\[= 0.02 \times 0.60 + 0.10 \times 0.30 + 0.25 \times 0.10\]
\[= 0.012 + 0.03 + 0.025 = 0.067\]

\textbf{Interpretación:} La probabilidad global de padecer la enfermedad
en esta población es del 6.7\%.

\end{tcolorbox}

\begin{center}\rule{0.5\linewidth}{0.5pt}\end{center}

\section{Teorema de Bayes}\label{sec-semana03-s9}

El Teorema de Bayes es uno de los resultados más importantes de la
teoría de la probabilidad y es fundamental en estadística moderna y
aprendizaje automático.

\subsection{Forma Simple del Teorema de
Bayes}\label{forma-simple-del-teorema-de-bayes}

\begin{tcolorbox}[enhanced jigsaw, breakable, titlerule=0mm, rightrule=.15mm, leftrule=.75mm, left=2mm, opacitybacktitle=0.6, colback=white, bottomtitle=1mm, colbacktitle=quarto-callout-note-color!10!white, title=\textcolor{quarto-callout-note-color}{\faInfo}\hspace{0.5em}{\textbf{Teorema: Teorema de Bayes (Forma Simple)}}, toptitle=1mm, coltitle=black, opacityback=0, arc=.35mm, bottomrule=.15mm, toprule=.15mm, colframe=quarto-callout-note-color-frame]

Para dos eventos \(A\) y \(B\) con \(\mathbb{P}(B) > 0\):

\[\mathbb{P}(A|B) = \frac{\mathbb{P}(B|A) \cdot \mathbb{P}(A)}{\mathbb{P}(B)}\]

\textbf{Terminología:} - \(\mathbb{P}(A|B)\) = probabilidad \textbf{a
posteriori} (después de observar \(B\))\\
- \(\mathbb{P}(A)\) = probabilidad \textbf{a priori} (antes de observar
\(B\))\\
- \(\mathbb{P}(B|A)\) = \textbf{verosimilitud} (probabilidad de la
evidencia bajo el supuesto)\\
- \(\mathbb{P}(B)\) = probabilidad \textbf{marginal} de la evidencia

\end{tcolorbox}

\subsection{Forma General con
Partición}\label{forma-general-con-particion}

\begin{tcolorbox}[enhanced jigsaw, breakable, titlerule=0mm, rightrule=.15mm, leftrule=.75mm, left=2mm, opacitybacktitle=0.6, colback=white, bottomtitle=1mm, colbacktitle=quarto-callout-note-color!10!white, title=\textcolor{quarto-callout-note-color}{\faInfo}\hspace{0.5em}{Teorema: Teorema de Bayes (Forma General)}, toptitle=1mm, coltitle=black, opacityback=0, arc=.35mm, bottomrule=.15mm, toprule=.15mm, colframe=quarto-callout-note-color-frame]

Sea \(\{A_1, A_2, \ldots, A_n\}\) una partición completa del espacio
muestral. Para un evento \(B\) con \(\mathbb{P}(B) > 0\):

\[\mathbb{P}(A_j|B) = \frac{\mathbb{P}(B|A_j) \cdot \mathbb{P}(A_j)}{\sum_{i=1}^{n} \mathbb{P}(B|A_i) \cdot \mathbb{P}(A_i)}\]

para cada \(j = 1, 2, \ldots, n\).

\end{tcolorbox}

\subsection{Ejemplo: Prueba
Diagnóstica}\label{ejemplo-prueba-diagnostica}

\begin{tcolorbox}[enhanced jigsaw, breakable, titlerule=0mm, rightrule=.15mm, leftrule=.75mm, left=2mm, opacitybacktitle=0.6, colback=white, bottomtitle=1mm, colbacktitle=quarto-callout-tip-color!10!white, title=\textcolor{quarto-callout-tip-color}{\faLightbulb}\hspace{0.5em}{Ejemplo 3.11: Prueba diagnóstica médica}, toptitle=1mm, coltitle=black, opacityback=0, arc=.35mm, bottomrule=.15mm, toprule=.15mm, colframe=quarto-callout-tip-color-frame]

\textbf{Escenario:} Un test de diagnóstico para una enfermedad con: -
Prevalencia de la enfermedad: \(\mathbb{P}(E) = 0.05\) (5\% de la
población) - Sensibilidad: \(\mathbb{P}(+|E) = 0.95\) (95\% de enfermos
dan positivo) - Especificidad: \(\mathbb{P}(-|\overline{E}) = 0.90\)
(90\% de sanos dan negativo)

\textbf{Pregunta:} Si una persona da positivo, ¿cuál es la probabilidad
de que realmente este enferma?

\textbf{Solución:}

Primero, calculamos
\(\mathbb{P}(+) = \mathbb{P}(+|E) \cdot \mathbb{P}(E) + \mathbb{P}(+|\overline{E}) \cdot \mathbb{P}(\overline{E})\)

donde
\(\mathbb{P}(+|\overline{E}) = 1 - \mathbb{P}(-|\overline{E}) = 1 - 0.90 = 0.10\)

\[\mathbb{P}(+) = 0.95 \times 0.05 + 0.10 \times 0.95 = 0.0475 + 0.095 = 0.1425\]

Ahora aplicamos el Teorema de Bayes:

\[\mathbb{P}(E|+) = \frac{\mathbb{P}(+|E) \cdot \mathbb{P}(E)}{\mathbb{P}(+)} = \frac{0.95 \times 0.05}{0.1425} = \frac{0.0475}{0.1425} \approx 0.333\]

\textbf{Interpretación:} Aunque el test tiene 95\% de sensibilidad, si
una persona da positivo, la probabilidad de que realmente esté enferma
es solo de aproximadamente 33\%. Esto sucede porque la enfermedad es
rara (baja prevalencia) y hay muchos falsos positivos.

\end{tcolorbox}

\begin{codelisting}

\caption{\label{lst-semana03-2}Código R}

\centering{

\begin{Shaded}
\begin{Highlighting}[]
\NormalTok{prev }\OtherTok{\textless{}{-}} \FloatTok{0.05}
\NormalTok{sens }\OtherTok{\textless{}{-}} \FloatTok{0.95}
\NormalTok{spec }\OtherTok{\textless{}{-}} \FloatTok{0.90}

\NormalTok{falso\_pos\_rate }\OtherTok{\textless{}{-}} \DecValTok{1} \SpecialCharTok{{-}}\NormalTok{ spec}
\NormalTok{p\_positivo }\OtherTok{\textless{}{-}}\NormalTok{ sens }\SpecialCharTok{*}\NormalTok{ prev }\SpecialCharTok{+}\NormalTok{ falso\_pos\_rate }\SpecialCharTok{*}\NormalTok{ (}\DecValTok{1} \SpecialCharTok{{-}}\NormalTok{ prev)}
\NormalTok{vpp }\OtherTok{\textless{}{-}}\NormalTok{ (sens }\SpecialCharTok{*}\NormalTok{ prev) }\SpecialCharTok{/}\NormalTok{ p\_positivo}

\FunctionTok{cat}\NormalTok{(}\StringTok{"Prevalencia:"}\NormalTok{, prev, }\StringTok{"}\SpecialCharTok{\textbackslash{}n}\StringTok{"}\NormalTok{)}
\end{Highlighting}
\end{Shaded}

}

\end{codelisting}%

\begin{verbatim}
Prevalencia: 0.05 
\end{verbatim}

\begin{codelisting}

\caption{\label{lst-semana03-2}Código R}

\centering{

\begin{Shaded}
\begin{Highlighting}[]
\FunctionTok{cat}\NormalTok{(}\StringTok{"Sensibilidad:"}\NormalTok{, sens, }\StringTok{"}\SpecialCharTok{\textbackslash{}n}\StringTok{"}\NormalTok{)}
\end{Highlighting}
\end{Shaded}

}

\end{codelisting}%

\begin{verbatim}
Sensibilidad: 0.95 
\end{verbatim}

\begin{codelisting}

\caption{\label{lst-semana03-2}Código R}

\centering{

\begin{Shaded}
\begin{Highlighting}[]
\FunctionTok{cat}\NormalTok{(}\StringTok{"Especificidad:"}\NormalTok{, spec, }\StringTok{"}\SpecialCharTok{\textbackslash{}n}\StringTok{"}\NormalTok{)}
\end{Highlighting}
\end{Shaded}

}

\end{codelisting}%

\begin{verbatim}
Especificidad: 0.9 
\end{verbatim}

\begin{codelisting}

\caption{\label{lst-semana03-2}Código R}

\centering{

\begin{Shaded}
\begin{Highlighting}[]
\FunctionTok{cat}\NormalTok{(}\StringTok{"P(+):"}\NormalTok{, p\_positivo, }\StringTok{"}\SpecialCharTok{\textbackslash{}n}\StringTok{"}\NormalTok{)}
\end{Highlighting}
\end{Shaded}

}

\end{codelisting}%

\begin{verbatim}
P(+): 0.1425 
\end{verbatim}

\begin{codelisting}

\caption{\label{lst-semana03-2}Código R}

\centering{

\begin{Shaded}
\begin{Highlighting}[]
\FunctionTok{cat}\NormalTok{(}\StringTok{"Valor Predictivo Positivo (VPP):"}\NormalTok{, }\FunctionTok{round}\NormalTok{(vpp, }\DecValTok{3}\NormalTok{), }\StringTok{"}\SpecialCharTok{\textbackslash{}n}\StringTok{"}\NormalTok{)}
\end{Highlighting}
\end{Shaded}

}

\end{codelisting}%

\begin{verbatim}
Valor Predictivo Positivo (VPP): 0.333 
\end{verbatim}

\begin{tcolorbox}[enhanced jigsaw, breakable, titlerule=0mm, rightrule=.15mm, leftrule=.75mm, left=2mm, opacitybacktitle=0.6, colback=white, bottomtitle=1mm, colbacktitle=quarto-callout-note-color!10!white, title=\textcolor{quarto-callout-note-color}{\faInfo}\hspace{0.5em}{Interpretación}, toptitle=1mm, coltitle=black, opacityback=0, arc=.35mm, bottomrule=.15mm, toprule=.15mm, colframe=quarto-callout-note-color-frame]

Este cálculo implementa el \textbf{Teorema de Bayes} en contexto
diagnóstico: aunque la prueba posee 95\% de sensibilidad (detecta
enfermedad cuando existe), el Valor Predictivo Positivo (VPP ≈ 0.333) es
sustancialmente menor debido a la baja prevalencia (5\%). Con una
enfermedad rara, la mayoría de resultados positivos son falsos
positivos. En la consulta clínica, este resultado orienta decisiones:
ante test positivo con baja prevalencia, es prudente confirmar con
prueba adicional antes de iniciar tratamiento, ilustrando cómo la
probabilidad a priori (prevalencia) moldea la interpretación clínica de
la evidencia.

\end{tcolorbox}

:::

\subsection{Ejemplo: Clasificación de
pacientse}\label{ejemplo-clasificacion-de-pacientse}

\begin{tcolorbox}[enhanced jigsaw, breakable, titlerule=0mm, rightrule=.15mm, leftrule=.75mm, left=2mm, opacitybacktitle=0.6, colback=white, bottomtitle=1mm, colbacktitle=quarto-callout-tip-color!10!white, title=\textcolor{quarto-callout-tip-color}{\faLightbulb}\hspace{0.5em}{Ejemplo 3.12: Clasificación de pacientes (Bayes)}, toptitle=1mm, coltitle=black, opacityback=0, arc=.35mm, bottomrule=.15mm, toprule=.15mm, colframe=quarto-callout-tip-color-frame]

\textbf{Escenario:} Un médico quiere clasificar a un paciente como de
alto riesgo o bajo riesgo de sufrir una patología cardíaca basándose en
tres factores de riesgo (ej. hipertensión \(W_1\), tabaquismo \(W_2\),
obesidad \(W_3\)).

\textbf{Información disponible:} - \(\mathbb{P}(E)\) = probabilidad a
priori de tener la patología - \(\mathbb{P}(W_i|E)\) = probabilidad de
que el factor \(W_i\) aparezca en pacientes enfermos -
\(\mathbb{P}(W_i|\overline{E})\) = probabilidad de que el factor \(W_i\)
aparezca en pacientes sanos

\textbf{Problema:} Clasificar a un paciente que presenta los factores
\(W_1, W_2, W_3\).

\textbf{Solución usando Bayes:}

Se clasifica como de ``Alto Riesgo'' si la probabilidad a posteriori es
alta:
\[\frac{\mathbb{P}(E|W_1 \cap W_2 \cap W_3)}{\mathbb{P}(\overline{E}|W_1 \cap W_2 \cap W_3)} > c\]

donde \(c\) es un umbral clínico.

Asumiendo independencia condicional de los factores:

\[\frac{\mathbb{P}(E|W_1 \cap W_2 \cap W_3)}{\mathbb{P}(\overline{E}|W_1 \cap W_2 \cap W_3)} \approx \frac{\mathbb{P}(E)}{\mathbb{P}(\overline{E})} \cdot \prod_{i=1}^{3} \frac{\mathbb{P}(W_i|E)}{\mathbb{P}(W_i|\overline{E})}\]

\textbf{Aplicación clínica:} - Permite integrar múltiples biomarcadores
o factores de riesgo de manera eficiente. - Proporciona una medida de
riesgo individualizada (probabilidad a posteriori). - Ayuda en la toma
de decisiones clínicas para el cribado o tratamiento preventivo.

\end{tcolorbox}

\begin{center}\rule{0.5\linewidth}{0.5pt}\end{center}

\section{Resumen}\label{sec-semana03-resumen}

\subsection{Conceptos Clave}\label{conceptos-clave}

\begin{itemize}
\tightlist
\item
  \textbf{Experimento aleatorio:} Un proceso con resultado incierto pero
  repetible
\item
  \textbf{Espacio muestral \(\Omega\):} El conjunto de todos los
  resultados posibles
\item
  \textbf{Evento:} Un subconjunto del espacio muestral
\item
  \textbf{Partición:} Una división del espacio muestral en eventos
  mutuamente disjuntos
\end{itemize}

\subsection{Definiciones de
Probabilidad}\label{definiciones-de-probabilidad}

\begin{longtable}[]{@{}
  >{\raggedright\arraybackslash}p{(\linewidth - 4\tabcolsep) * \real{0.3333}}
  >{\raggedright\arraybackslash}p{(\linewidth - 4\tabcolsep) * \real{0.3333}}
  >{\raggedright\arraybackslash}p{(\linewidth - 4\tabcolsep) * \real{0.3333}}@{}}
\toprule\noalign{}
\begin{minipage}[b]{\linewidth}\raggedright
Definición
\end{minipage} & \begin{minipage}[b]{\linewidth}\raggedright
Ventajas
\end{minipage} & \begin{minipage}[b]{\linewidth}\raggedright
Limitaciones
\end{minipage} \\
\midrule\noalign{}
\endhead
\bottomrule\noalign{}
\endlastfoot
\textbf{Clásica (Laplace)} & Simple, intuitiva & Requiere
equiprobabilidad \\
\textbf{Frecuentista (von Mises)} & Empírica, práctica & Requiere
infinitas repeticiones \\
\textbf{Axiomática (Kolmogorov)} & General, rigurosa & Abstracta \\
\end{longtable}

\subsection{Fórmulas Importantes}\label{formulas-importantes}

\textbf{Propiedades básicas:}
\[\mathbb{P}(\overline{A}) = 1 - \mathbb{P}(A), \quad \mathbb{P}(\emptyset) = 0, \quad \mathbb{P}(\Omega) = 1\]

\textbf{Regla de la suma:}
\[\mathbb{P}(A \cup B) = \mathbb{P}(A) + \mathbb{P}(B) - \mathbb{P}(A \cap B)\]

\textbf{Probabilidad condicional:}
\[\mathbb{P}(A|B) = \frac{\mathbb{P}(A \cap B)}{\mathbb{P}(B)}\]

\textbf{Regla de la multiplicación:}
\[\mathbb{P}(A \cap B) = \mathbb{P}(A) \cdot \mathbb{P}(B|A)\]

\textbf{Independencia:}
\[\mathbb{P}(A \cap B) = \mathbb{P}(A) \cdot \mathbb{P}(B) \text{ si y solo si } A \text{ y } B \text{ son independientes}\]

\textbf{Ley de la probabilidad total:}
\[\mathbb{P}(B) = \sum_{i=1}^{n} \mathbb{P}(B|A_i) \cdot \mathbb{P}(A_i)\]

\textbf{Teorema de Bayes:}
\[\mathbb{P}(A|B) = \frac{\mathbb{P}(B|A) \cdot \mathbb{P}(A)}{\mathbb{P}(B)}\]

\[\mathbb{P}(A_j|B) = \frac{\mathbb{P}(B|A_j) \cdot \mathbb{P}(A_j)}{\sum_{i=1}^{n} \mathbb{P}(B|A_i) \cdot \mathbb{P}(A_i)}\]

\begin{center}\rule{0.5\linewidth}{0.5pt}\end{center}

\section{Ejercicios}\label{sec-semana03-ejercicios}

\textbf{Ejercicio 3.1:} En una clase hay 30 estudiantes de doctorado
asistiendo a un seminario de Bioestadística avanzada en Medicina. De
ellos, 18 se graduaron en Medicina, 15 tienen un grado en Estadistica y
10 estudiantes estudiaron ambos grados. Si seleccionamos un estudiante
al azar:

\begin{enumerate}
\def\labelenumi{\alph{enumi})}
\tightlist
\item
  ¿Cuál es la probabilidad de que se graduado en Medicina o estadística?
\item
  ¿Cuál es la probabilidad de que haya estudiado solo Medicina?
\item
  ¿Cuál es la probabilidad de que no hay estudiado ninguna de las dos?
\end{enumerate}

\textbf{Ejercicio 3.2:} Un médico tiene dos opciones terapeúticas de
antibióticos de primera línea (A o B) y tres de segunda línea de
administración conjunta al tratamiento principal (1, 2, o 3) para el
tratamiento de una enfermedad infecciosa y todas las opciones
terapeúticas son igualmente eficaces. Si todas las opciones se eligen al
azar:

\begin{enumerate}
\def\labelenumi{\alph{enumi})}
\tightlist
\item
  ¿Cuál es el espacio muestral?
\item
  ¿Cuál es la probabilidad de elegir el antibiotico de primera línea A?
\item
  ¿Cuál es la probabilidad de elegir el antibiotico de primera línea B y
  el de segunda línea 2?
\end{enumerate}

\textbf{Ejercicio 3.3:} En un hospital, la probabilidad de que un
paciente ingrese por fractura de cadera en un día es 0.1. Si evaluamos
tres días independientemente:

\begin{enumerate}
\def\labelenumi{\alph{enumi})}
\tightlist
\item
  ¿Cuál es la probabilidad de que no ingrese ningún paciente?
\item
  ¿Cuál es la probabilidad de que exactamente ingrese uno?
\end{enumerate}

\textbf{Ejercicio 3.4:} Se sabe que: - El 60\% de los pacientes de una
consulta son mujeres - El 40\% de las mujeres tinen menos de 30 años -
El 30\% de los hombres tienen menos de 30 años

Si seleccionamos un cliente al azar:

\begin{enumerate}
\def\labelenumi{\alph{enumi})}
\tightlist
\item
  ¿Cuál es la probabilidad de que un paciente tenga menos de 30 años?
\item
  Si un paciente tiene 30 años, ¿cuál es la probabilidad de que sea
  mujer?
\end{enumerate}

\textbf{Ejercicio 3.5:} Un test para detectar una enfermedad rara tiene:
- 98\% de sensibilidad (detecta la enfermedad cuando está presente) -
95\% de especificidad (identifica correctamente a personas sanas) - La
enfermedad afecta al 0.5\% de la población

\begin{enumerate}
\def\labelenumi{\alph{enumi})}
\tightlist
\item
  Si una persona da positivo en el test, ¿cuál es la probabilidad de que
  realmente esté enferma?
\item
  Interprete el resultado. ¿Por qué la probabilidad no es más alta?
\end{enumerate}

\textbf{Ejercicio 3.6:} Una empresa privada vende ecografos a tres
hospitales en tres ciudades distintas: Madrid (40\% de ventas),
Barcelona (35\%) y Valencia (25\%). La probabilidad de cumplir la cuota
de ventas es: - Madrid: 0.8 - Barcelona: 0.85 - Valencia: 0.7

Si se selecciona un mes al azar:

\begin{enumerate}
\def\labelenumi{\alph{enumi})}
\tightlist
\item
  ¿Cuál es la probabilidad de que se cumpla la cuota?
\item
  Si se cumplió la cuota, ¿cuál es la probabilidad de que haya sido en
  Barcelona?
\end{enumerate}

\section{Respuestas a los Ejercicios}\label{sec-semana03-respuestas}

\textbf{Ejercicio 3.1:} Diagramas de Venn\\
- a) P(M ∪ E) = P(M) + P(E) - P(M ∩ E) = (18+15-10)/30 = 23/30 ≈ 0.767\\
- b) P(Solo M) = (18-10)/30 = 8/30 ≈ 0.267\\
- c) P(Ninguna) = 1 - 23/30 = 7/30 ≈ 0.233

\textbf{Ejercicio 3.2:} Espacio Muestral y Probabilidad\\
- a) Ω = \{A1, A2, A3, B1, B2, B3\}, \textbar Ω\textbar{} = 6\\
- b) P(A) = 3/6 = 0.5\\
- c) P(B y 2) = 1/6 ≈ 0.167

\textbf{Ejercicio 3.3:} Independencia\\
- a) P(Ningún ingreso) = (0.9)³ = 0.729\\
- b) P(exactamente 1 ingreso) = C(3,1)×(0.1)¹×(0.9)² = 3×0.1×0.81 =
0.243\\
- Nota: Más adelante veremos como caalcular esta probabilidad usando la
fórmula de la distribución Binomial (n, p).

\textbf{Ejercicio 3.4:} Probabilidad Total y Teorema de Bayes\\
- a) P(Menos de 30 años) = 0.6×0.4 + 0.4×0.3 = 0.24 + 0.12 = 0.36\\
- b) P(Mujer\textbar Menos de 30 años) = (0.6x0.4)/0.36 = 0.24/0.36 =
2/3 ≈ 0.667

\textbf{Ejercicio 3.5:} Test Diagnóstico (Bayes)\\
- a) P(Enfermedad\textbar Positivo) = (0.98×0.005)/(0.98×0.005 +
0.05×0.995) = 0.0049/(0.0049+0.04975) ≈ 0.090\\
- b) Aunque el test es muy preciso (98\% sensibilidad), como la
enfermedad es rara (0.5\%), la mayoría de positivos son falsos
positivos.

\textbf{Ejercicio 3.6:} Probabilidad Total y Bayes (Hospitales)\\
- a) P(Cuota) = 0.4×0.8 + 0.35×0.85 + 0.25×0.7 = 0.32 + 0.2975 + 0.175 =
0.7925\\
- b) P(Barcelona\textbar Cuota) = (0.35×0.85)/0.7925 = 0.2975/0.7925 ≈
0.375

\begin{center}\rule{0.5\linewidth}{0.5pt}\end{center}

\begin{tcolorbox}[enhanced jigsaw, breakable, titlerule=0mm, rightrule=.15mm, leftrule=.75mm, left=2mm, opacitybacktitle=0.6, colback=white, bottomtitle=1mm, colbacktitle=quarto-callout-tip-color!10!white, title=\textcolor{quarto-callout-tip-color}{\faLightbulb}\hspace{0.5em}{Métodos Avanzados}, toptitle=1mm, coltitle=black, opacityback=0, arc=.35mm, bottomrule=.15mm, toprule=.15mm, colframe=quarto-callout-tip-color-frame]

Para ampliar los contenidos de este capítulo con técnicas estadísticas
avanzadas, visita:

\textbf{\href{https://migariane.github.io/Bioestadistica}{→
Bioestadística Avanzada --- M.A.~Luque Fernández}}

\end{tcolorbox}

\chapter{Semana 4 --- Variables Aleatorias y Distribuciones de
Probabilidad}\label{sec-semana04}

\section{Introducción}\label{introduccion-1}

Las variables aleatorias son el concepto fundamental que conecta la
teoría de la probabilidad con la estadística práctica. Esta semana
exploramos los conceptos de variable aleatoria, sus funciones de
distribución, y las distribuciones más importantes que usarás en
bioestadística y análisis de datos.

\section{Variables Aleatorias}\label{sec-semana04-va}

\begin{tcolorbox}[enhanced jigsaw, breakable, titlerule=0mm, rightrule=.15mm, leftrule=.75mm, left=2mm, opacitybacktitle=0.6, colback=white, bottomtitle=1mm, colbacktitle=quarto-callout-note-color!10!white, title=\textcolor{quarto-callout-note-color}{\faInfo}\hspace{0.5em}{Definición: Variable Aleatoria}, toptitle=1mm, coltitle=black, opacityback=0, arc=.35mm, bottomrule=.15mm, toprule=.15mm, colframe=quarto-callout-note-color-frame]

Una \textbf{variable aleatoria} \(X\) es una función que asigna un
número real a cada resultado en el espacio muestral. Formalmente:

\[X: \Omega \to \mathbb{R}, \quad X(\omega) = x\]

donde \(\Omega\) es el espacio muestral.

\end{tcolorbox}

\subsection{Tipos de Variables
Aleatorias}\label{tipos-de-variables-aleatorias}

\textbf{Variable Aleatoria Discreta:} Toma un número finito o infinito
numerable de valores.

\[X \in \{x_1, x_2, x_3, \ldots\}\]

\textbf{Variable Aleatoria Continua:} Puede tomar cualquier valor en un
intervalo o en toda la recta real.

\[X \in [a, b] \text{ o } X \in \mathbb{R}\]

\subsection{Función de Probabilidad y Función de
Distribución}\label{funcion-de-probabilidad-y-funcion-de-distribucion}

Para variables \textbf{discretas}, la \textbf{Función de Masa de
Probabilidad (PMF)} es: \[P(X = x_i) = f(x_i)\]

con propiedades: - \(f(x_i) \geq 0\) para todo \(i\) -
\(\sum_i f(x_i) = 1\)

Para variables \textbf{continuas}, la \textbf{Función de Densidad de
Probabilidad (PDF)} es: \[f(x) \geq 0 \text{ para todo } x\]
\[\int_{-\infty}^{+\infty} f(x) \, dx = 1\]

La \textbf{Función de Distribución Acumulada (CDF)} es:
\[F(x) = P(X \leq x)\]

con propiedades: - \(0 \leq F(x) \leq 1\) - \(F\) es no decreciente -
\(\lim_{x \to -\infty} F(x) = 0\) y \(\lim_{x \to +\infty} F(x) = 1\)

Para variables continuas: \[f(x) = \frac{dF(x)}{dx}\]

\section{Funciones de Probabilidad en R}\label{sec-semana04-r-funcs}

R ofrece un conjunto consistente de funciones para trabajar con
distribuciones de probabilidad. Para cualquier distribución
(\texttt{distribucion}), se dispone de cuatro funciones prefijadas:

\begin{itemize}
\tightlist
\item
  \textbf{\texttt{d} (density / mass):} \texttt{ddistribucion(x,\ ...)}
  --- Calcula la función de densidad (PDF) o masa (PMF).
\item
  \textbf{\texttt{p} (probability):} \texttt{pdistribucion(q,\ ...)} ---
  Calcula la CDF: \(P(X \le q)\).
\item
  \textbf{\texttt{q} (quantile):} \texttt{qdistribucion(p,\ ...)} ---
  Calcula la función cuantil (inversa de CDF).
\item
  \textbf{\texttt{r} (random):} \texttt{rdistribucion(n,\ ...)} ---
  Genera números aleatorios.
\end{itemize}

\section{Fundamentos de Cálculo de Probabilidades: Discreta
vs.~Continua}\label{sec-semana04-fundamentos}

Antes de explorar distribuciones específicas, es vital comprender las
diferencias fundamentales en cómo calculamos probabilidades dependiendo
de la naturaleza de la variable.

\subsection{La Singularidad del Punto en Variables
Continuas}\label{la-singularidad-del-punto-en-variables-continuas}

En distribuciones continuas (como la Normal), la probabilidad de que una
variable tome un valor puntual exacto es siempre cero: \[P(X = x) = 0\]

\textbf{¿Por qué?} Matemáticamente, la probabilidad de un punto se
define como el área bajo la curva en un intervalo de ancho cero:
\[P(X = x) = \int_{x}^{x} f(t) \, dt = 0\]

En estadística práctica, esto significa que solo tiene sentido hablar de
probabilidades en \textbf{intervalos}:
\(P(a \le X \le b) = \int_{a}^{b} f(t) \, dt\). Por lo tanto, para una
continua, \(P(X \le x) = P(X < x)\).

\subsection{El Cuidado de las Desigualdades en Variables
Discretas}\label{el-cuidado-de-las-desigualdades-en-variables-discretas}

En variables discretas (como la Binomial), \(P(X=k) > 0\). Esto implica
que debemos ser extremadamente precisos al usar operadores de
comparación (\(<, \le, >, \ge\)), ya que el incluir o excluir un punto
cambia el resultado significativamente.

\textbf{La Regla de Oro:} - \(P(X \le k) = P(X < k) + P(X=k)\) -
\(P(X > k) = 1 - P(X \le k)\) - \(P(X \ge k) = 1 - P(X \le k-1)\)

\subsubsection{Ejemplo comparativo en R}\label{ejemplo-comparativo-en-r}

\begin{codelisting}

\caption{\label{lst-semana04-1}Cálculo comparativo discreto vs continuo}

\centering{

\begin{Shaded}
\begin{Highlighting}[]
\CommentTok{\# Discreta: Binomial(n=10, p=0.5)}
\CommentTok{\# P(X = 5) es distinta de cero}
\FunctionTok{print}\NormalTok{(}\FunctionTok{paste}\NormalTok{(}\StringTok{"P(X = 5) ="}\NormalTok{, }\FunctionTok{dbinom}\NormalTok{(}\DecValTok{5}\NormalTok{, }\DecValTok{10}\NormalTok{, }\FloatTok{0.5}\NormalTok{)))}
\end{Highlighting}
\end{Shaded}

}

\end{codelisting}%

\begin{verbatim}
[1] "P(X = 5) = 0.24609375"
\end{verbatim}

\begin{codelisting}

\caption{\label{lst-semana04-1}Cálculo comparativo discreto vs continuo}

\centering{

\begin{Shaded}
\begin{Highlighting}[]
\FunctionTok{print}\NormalTok{(}\FunctionTok{paste}\NormalTok{(}\StringTok{"P(X \textless{}= 5) ="}\NormalTok{, }\FunctionTok{pbinom}\NormalTok{(}\DecValTok{5}\NormalTok{, }\DecValTok{10}\NormalTok{, }\FloatTok{0.5}\NormalTok{)))}
\end{Highlighting}
\end{Shaded}

}

\end{codelisting}%

\begin{verbatim}
[1] "P(X <= 5) = 0.623046875"
\end{verbatim}

\begin{codelisting}

\caption{\label{lst-semana04-1}Cálculo comparativo discreto vs continuo}

\centering{

\begin{Shaded}
\begin{Highlighting}[]
\FunctionTok{print}\NormalTok{(}\FunctionTok{paste}\NormalTok{(}\StringTok{"P(X \textless{} 5) ="}\NormalTok{, }\FunctionTok{pbinom}\NormalTok{(}\DecValTok{4}\NormalTok{, }\DecValTok{10}\NormalTok{, }\FloatTok{0.5}\NormalTok{)))}
\end{Highlighting}
\end{Shaded}

}

\end{codelisting}%

\begin{verbatim}
[1] "P(X < 5) = 0.376953125"
\end{verbatim}

\begin{codelisting}

\caption{\label{lst-semana04-1}Cálculo comparativo discreto vs continuo}

\centering{

\begin{Shaded}
\begin{Highlighting}[]
\CommentTok{\# Continua: Normal(mu=0, sd=1)}
\CommentTok{\# P(X = 0) es cero}
\FunctionTok{print}\NormalTok{(}\FunctionTok{paste}\NormalTok{(}\StringTok{"P(X = 0) ="}\NormalTok{, }\FunctionTok{dnorm}\NormalTok{(}\DecValTok{0}\NormalTok{))) }
\end{Highlighting}
\end{Shaded}

}

\end{codelisting}%

\begin{verbatim}
[1] "P(X = 0) = 0.398942280401433"
\end{verbatim}

\begin{codelisting}

\caption{\label{lst-semana04-1}Cálculo comparativo discreto vs continuo}

\centering{

\begin{Shaded}
\begin{Highlighting}[]
\CommentTok{\# P(X \textless{}= 0) es 0.5}
\FunctionTok{print}\NormalTok{(}\FunctionTok{paste}\NormalTok{(}\StringTok{"P(X \textless{}= 0) ="}\NormalTok{, }\FunctionTok{pnorm}\NormalTok{(}\DecValTok{0}\NormalTok{)))}
\end{Highlighting}
\end{Shaded}

}

\end{codelisting}%

\begin{verbatim}
[1] "P(X <= 0) = 0.5"
\end{verbatim}

\begin{codelisting}

\caption{\label{lst-semana04-1}Cálculo comparativo discreto vs continuo}

\centering{

\begin{Shaded}
\begin{Highlighting}[]
\CommentTok{\# P(X \textless{} 0) también es 0.5}
\FunctionTok{print}\NormalTok{(}\FunctionTok{paste}\NormalTok{(}\StringTok{"P(X \textless{} 0) ="}\NormalTok{, }\FunctionTok{pnorm}\NormalTok{(}\DecValTok{0}\NormalTok{)))}
\end{Highlighting}
\end{Shaded}

}

\end{codelisting}%

\begin{verbatim}
[1] "P(X < 0) = 0.5"
\end{verbatim}

\begin{tcolorbox}[enhanced jigsaw, breakable, titlerule=0mm, rightrule=.15mm, leftrule=.75mm, left=2mm, opacitybacktitle=0.6, colback=white, bottomtitle=1mm, colbacktitle=quarto-callout-note-color!10!white, title=\textcolor{quarto-callout-note-color}{\faInfo}\hspace{0.5em}{Resumen de operadores}, toptitle=1mm, coltitle=black, opacityback=0, arc=.35mm, bottomrule=.15mm, toprule=.15mm, colframe=quarto-callout-note-color-frame]

\begin{longtable}[]{@{}lll@{}}
\toprule\noalign{}
Operador & Continuo & Discreto \\
\midrule\noalign{}
\endhead
\bottomrule\noalign{}
\endlastfoot
\textbf{\(P(X = k)\)} & 0 & \(f(k)\) \\
\textbf{\(P(X \le k)\)} & \texttt{pdist(k)} & \texttt{pdist(k)} \\
\textbf{\(P(X < k)\)} & \texttt{pdist(k)} & \texttt{pdist(k-1)} \\
\textbf{\(P(X > k)\)} & \texttt{1\ -\ pdist(k)} &
\texttt{1\ -\ pdist(k)} \\
\textbf{\(P(X \ge k)\)} & \texttt{1\ -\ pdist(k)} &
\texttt{1\ -\ pdist(k-1)} \\
\end{longtable}

\end{tcolorbox}

\section{Varianza y Desviación Estándar}\label{sec-semana04-varianza}

\begin{tcolorbox}[enhanced jigsaw, breakable, titlerule=0mm, rightrule=.15mm, leftrule=.75mm, left=2mm, opacitybacktitle=0.6, colback=white, bottomtitle=1mm, colbacktitle=quarto-callout-note-color!10!white, title=\textcolor{quarto-callout-note-color}{\faInfo}\hspace{0.5em}{Definición: Varianza}, toptitle=1mm, coltitle=black, opacityback=0, arc=.35mm, bottomrule=.15mm, toprule=.15mm, colframe=quarto-callout-note-color-frame]

La \textbf{varianza} mide la dispersión de una variable aleatoria
respecto a su esperanza:

\[\text{Var}(X) = \sigma^2_X = E[(X - E(X))^2]\]

Fórmula computacional equivalente: \[\text{Var}(X) = E(X^2) - [E(X)]^2\]

\textbf{Caso discreto:}
\[\text{Var}(X) = \sum_{i=1}^{\infty} [x_i - E(X)]^2 \cdot f(x_i) = \sum_{i=1}^{\infty} x_i^2 \cdot f(x_i) - [E(X)]^2\]

\textbf{Caso continuo:}
\[\text{Var}(X) = \int_{-\infty}^{+\infty} [x - E(X)]^2 \cdot f(x) \, dx = \int_{-\infty}^{+\infty} x^2 \cdot f(x) \, dx - [E(X)]^2\]

\end{tcolorbox}

\subsection{Desviación Estándar
poblacional}\label{desviacion-estandar-poblacional}

La \textbf{desviación estándar} es la raíz cuadrada positiva de la
varianza: \[\sigma_X = \sqrt{\text{Var}(X)}\]

Tiene las mismas unidades que la variable original.

\subsection{Propiedades de la
Varianza}\label{propiedades-de-la-varianza}

\begin{tcolorbox}[enhanced jigsaw, breakable, titlerule=0mm, rightrule=.15mm, leftrule=.75mm, left=2mm, opacitybacktitle=0.6, colback=white, bottomtitle=1mm, colbacktitle=quarto-callout-warning-color!10!white, title=\textcolor{quarto-callout-warning-color}{\faExclamationTriangle}\hspace{0.5em}{Propiedades Importantes}, toptitle=1mm, coltitle=black, opacityback=0, arc=.35mm, bottomrule=.15mm, toprule=.15mm, colframe=quarto-callout-warning-color-frame]

\begin{enumerate}
\def\labelenumi{\arabic{enumi}.}
\item
  \textbf{Transformación lineal:} Si \(Y = a + bX\):
  \[\text{Var}(Y) = b^2 \cdot \text{Var}(X)\] (nota: la constante
  aditiva no afecta la varianza)
\item
  \textbf{Suma de variables independientes:} Si \(X\) e \(Y\) son
  independientes: \[\text{Var}(X + Y) = \text{Var}(X) + \text{Var}(Y)\]
  \[\text{Var}(X - Y) = \text{Var}(X) + \text{Var}(Y)\]
\end{enumerate}

\end{tcolorbox}

\section{Estandarización (Transformación
Z)}\label{sec-semana04-estandarizacion}

\begin{tcolorbox}[enhanced jigsaw, breakable, titlerule=0mm, rightrule=.15mm, leftrule=.75mm, left=2mm, opacitybacktitle=0.6, colback=white, bottomtitle=1mm, colbacktitle=quarto-callout-note-color!10!white, title=\textcolor{quarto-callout-note-color}{\faInfo}\hspace{0.5em}{Definición: Transformación Z}, toptitle=1mm, coltitle=black, opacityback=0, arc=.35mm, bottomrule=.15mm, toprule=.15mm, colframe=quarto-callout-note-color-frame]

La \textbf{estandarización} o \textbf{normalización} transforma una
variable aleatoria para que tenga media 0 y varianza 1:

\[Z = \frac{X - E(X)}{\sqrt{\text{Var}(X)}} = \frac{X - \mu_X}{\sigma_X}\]

Propiedades de \(Z\):\\
- \(E(Z) = 0\)\\
- \(\text{Var}(Z) = 1\)\\
- \(Z\) es adimensional (sin unidades)

\end{tcolorbox}

\section{Variables Aleatorias
Bivariantes}\label{sec-semana04-bivariante}

\subsection{Distribuciones Conjunta, Marginal y Condicional (Caso
Discreto)}\label{distribuciones-conjunta-marginal-y-condicional-caso-discreto}

\begin{tcolorbox}[enhanced jigsaw, breakable, titlerule=0mm, rightrule=.15mm, leftrule=.75mm, left=2mm, opacitybacktitle=0.6, colback=white, bottomtitle=1mm, colbacktitle=quarto-callout-note-color!10!white, title=\textcolor{quarto-callout-note-color}{\faInfo}\hspace{0.5em}{Distribuciones Discretas Bivariantes}, toptitle=1mm, coltitle=black, opacityback=0, arc=.35mm, bottomrule=.15mm, toprule=.15mm, colframe=quarto-callout-note-color-frame]

La \textbf{función de probabilidad conjunta} es:
\[P(X = x_i, Y = y_j) = f(x_i, y_j)\]

con propiedades: - \(f(x_i, y_j) \geq 0\) -
\(\sum_i \sum_j f(x_i, y_j) = 1\)

Las \textbf{distribuciones marginales} se obtienen sumando:
\[f(x_i) = P(X = x_i) = \sum_j f(x_i, y_j)\]
\[f(y_j) = P(Y = y_j) = \sum_i f(x_i, y_j)\]

La \textbf{distribución condicional} es:
\[P(X = x_i | Y = y_j) = \frac{f(x_i, y_j)}{f(y_j)}\]

\end{tcolorbox}

\begin{tcolorbox}[enhanced jigsaw, breakable, titlerule=0mm, rightrule=.15mm, leftrule=.75mm, left=2mm, opacitybacktitle=0.6, colback=white, bottomtitle=1mm, colbacktitle=quarto-callout-warning-color!10!white, title=\textcolor{quarto-callout-warning-color}{\faExclamationTriangle}\hspace{0.5em}{Independencia}, toptitle=1mm, coltitle=black, opacityback=0, arc=.35mm, bottomrule=.15mm, toprule=.15mm, colframe=quarto-callout-warning-color-frame]

Dos variables aleatorias \(X\) e \(Y\) son \textbf{independientes} si:
\[f(x_i, y_j) = f(x_i) \cdot f(y_j) \quad \text{para todo } x_i, y_j\]

\textbf{Importante:} La independencia implica que la distribución
condicional iguala la marginal: \[P(X = x_i | Y = y_j) = P(X = x_i)\]

\end{tcolorbox}

\subsection{Distribuciones Conjunta, Marginal y Condicional (Caso
Continuo)}\label{distribuciones-conjunta-marginal-y-condicional-caso-continuo}

\begin{tcolorbox}[enhanced jigsaw, breakable, titlerule=0mm, rightrule=.15mm, leftrule=.75mm, left=2mm, opacitybacktitle=0.6, colback=white, bottomtitle=1mm, colbacktitle=quarto-callout-note-color!10!white, title=\textcolor{quarto-callout-note-color}{\faInfo}\hspace{0.5em}{Distribuciones Continuas Bivariantes}, toptitle=1mm, coltitle=black, opacityback=0, arc=.35mm, bottomrule=.15mm, toprule=.15mm, colframe=quarto-callout-note-color-frame]

La \textbf{función de densidad conjunta} satisface: \[f(x, y) \geq 0\]
\[\int_{-\infty}^{+\infty} \int_{-\infty}^{+\infty} f(x, y) \, dx \, dy = 1\]

Las \textbf{funciones de densidad marginales} son:
\[f(x) = \int_{-\infty}^{+\infty} f(x, y) \, dy\]
\[f(y) = \int_{-\infty}^{+\infty} f(x, y) \, dx\]

La \textbf{densidad condicional} es: \[f(x | y) = \frac{f(x, y)}{f(y)}\]

\end{tcolorbox}

\subsection{Covarianza y Correlación}\label{covarianza-y-correlacion}

\begin{tcolorbox}[enhanced jigsaw, breakable, titlerule=0mm, rightrule=.15mm, leftrule=.75mm, left=2mm, opacitybacktitle=0.6, colback=white, bottomtitle=1mm, colbacktitle=quarto-callout-note-color!10!white, title=\textcolor{quarto-callout-note-color}{\faInfo}\hspace{0.5em}{Definición: Covarianza}, toptitle=1mm, coltitle=black, opacityback=0, arc=.35mm, bottomrule=.15mm, toprule=.15mm, colframe=quarto-callout-note-color-frame]

La \textbf{covarianza} mide la variabilidad conjunta de dos variables
aleatorias: \[\text{Cov}(X, Y) = E[(X - E(X))(Y - E(Y))]\]

Fórmula computacional: \[\text{Cov}(X, Y) = E(XY) - E(X) \cdot E(Y)\]

Propiedades: - \(\text{Cov}(X, X) = \text{Var}(X)\) - Si \(X\) e \(Y\)
son independientes: \(\text{Cov}(X, Y) = 0\) - \textbf{Nota importante:}
\(\text{Cov}(X, Y) = 0\) NO implica independencia

\end{tcolorbox}

\begin{tcolorbox}[enhanced jigsaw, breakable, titlerule=0mm, rightrule=.15mm, leftrule=.75mm, left=2mm, opacitybacktitle=0.6, colback=white, bottomtitle=1mm, colbacktitle=quarto-callout-note-color!10!white, title=\textcolor{quarto-callout-note-color}{\faInfo}\hspace{0.5em}{Definición: Correlacion}, toptitle=1mm, coltitle=black, opacityback=0, arc=.35mm, bottomrule=.15mm, toprule=.15mm, colframe=quarto-callout-note-color-frame]

El \textbf{coeficiente de correlación} (o coeficiente de correlación de
Pearson) en una población es:
\[\rho(X, Y) = \text{Corr}(X, Y) = \frac{\text{Cov}(X, Y)}{\sigma_X \cdot \sigma_Y}\]

Propiedades: - \(-1 \leq \rho(X, Y) \leq 1\) - \(\rho = 1\): correlación
positiva perfecta - \(\rho = 0\): no hay correlación lineal -
\(\rho = -1\): correlación negativa perfecta

\end{tcolorbox}

\section{Distribuciones Discretas
Importantes}\label{sec-semana04-discretas}

\subsection{\texorpdfstring{Distribución Uniforme Discreta
\(U(n)\)}{Distribución Uniforme Discreta U(n)}}\label{distribucion-uniforme-discreta-un}

\begin{tcolorbox}[enhanced jigsaw, breakable, titlerule=0mm, rightrule=.15mm, leftrule=.75mm, left=2mm, opacitybacktitle=0.6, colback=white, bottomtitle=1mm, colbacktitle=quarto-callout-note-color!10!white, title=\textcolor{quarto-callout-note-color}{\faInfo}\hspace{0.5em}{Definición: Distribución Uniforme Discreta}, toptitle=1mm, coltitle=black, opacityback=0, arc=.35mm, bottomrule=.15mm, toprule=.15mm, colframe=quarto-callout-note-color-frame]

Una variable aleatoria \(X \sim U(n)\) toma \(n\) valores con igual
probabilidad:

\[P(X = x_i) = \frac{1}{n} \quad \text{para } i = 1, 2, \ldots, n\]

\textbf{Parámetros:} \(n\) (número de valores posibles)

\textbf{Esperanza y Varianza:} \[E(X) = \frac{1}{n} \sum_{i=1}^n x_i\]
\[\text{Var}(X) = \frac{1}{n} \sum_{i=1}^n (x_i - E(X))^2\]

\textbf{En R:} No existe función directa; usar \texttt{sample()} o
implementar manualmente.

\end{tcolorbox}

\begin{tcolorbox}[enhanced jigsaw, breakable, titlerule=0mm, rightrule=.15mm, leftrule=.75mm, left=2mm, opacitybacktitle=0.6, colback=white, bottomtitle=1mm, colbacktitle=quarto-callout-tip-color!10!white, title=\textcolor{quarto-callout-tip-color}{\faLightbulb}\hspace{0.5em}{Ejemplo 4.1: Lanzamiento de dado}, toptitle=1mm, coltitle=black, opacityback=0, arc=.35mm, bottomrule=.15mm, toprule=.15mm, colframe=quarto-callout-tip-color-frame]

Un dado justo tiene \(n = 6\) caras:
\[P(X = k) = \frac{1}{6} \quad \text{para } k = 1, 2, 3, 4, 5, 6\]

\[E(X) = \frac{1}{6}(1 + 2 + 3 + 4 + 5 + 6) = 3.5\]
\[\text{Var}(X) = \frac{1}{6}\sum_{k=1}^6 (k - 3.5)^2 \approx 2.92\]

\begin{codelisting}[H]

\caption{\label{lst-semana04-2}Código R}

\centering{

\begin{Shaded}
\begin{Highlighting}[]
\CommentTok{\# En R: calcular probabilidades para distribución uniforme discreta}
\NormalTok{x }\OtherTok{\textless{}{-}} \DecValTok{1}\SpecialCharTok{:}\DecValTok{6}
\NormalTok{probs }\OtherTok{\textless{}{-}} \FunctionTok{rep}\NormalTok{(}\DecValTok{1}\SpecialCharTok{/}\DecValTok{6}\NormalTok{, }\DecValTok{6}\NormalTok{);probs}
\end{Highlighting}
\end{Shaded}

}

\end{codelisting}%

\begin{verbatim}
[1] 0.1666667 0.1666667 0.1666667 0.1666667 0.1666667 0.1666667
\end{verbatim}

\begin{codelisting}[H]

\caption{\label{lst-semana04-2}Código R}

\centering{

\begin{Shaded}
\begin{Highlighting}[]
\NormalTok{mean\_x }\OtherTok{\textless{}{-}} \FunctionTok{sum}\NormalTok{(x }\SpecialCharTok{*}\NormalTok{ probs);mean\_x}
\end{Highlighting}
\end{Shaded}

}

\end{codelisting}%

\begin{verbatim}
[1] 3.5
\end{verbatim}

\begin{codelisting}[H]

\caption{\label{lst-semana04-2}Código R}

\centering{

\begin{Shaded}
\begin{Highlighting}[]
\NormalTok{var\_x }\OtherTok{\textless{}{-}} \FunctionTok{sum}\NormalTok{((x }\SpecialCharTok{{-}}\NormalTok{ mean\_x)}\SpecialCharTok{\^{}}\DecValTok{2} \SpecialCharTok{*}\NormalTok{ probs); var\_x}
\end{Highlighting}
\end{Shaded}

}

\end{codelisting}%

\begin{verbatim}
[1] 2.916667
\end{verbatim}

\end{tcolorbox}

\subsection{\texorpdfstring{Distribución de Bernoulli
\(B(p)\)}{Distribución de Bernoulli B(p)}}\label{distribucion-de-bernoulli-bp}

\begin{tcolorbox}[enhanced jigsaw, breakable, titlerule=0mm, rightrule=.15mm, leftrule=.75mm, left=2mm, opacitybacktitle=0.6, colback=white, bottomtitle=1mm, colbacktitle=quarto-callout-note-color!10!white, title=\textcolor{quarto-callout-note-color}{\faInfo}\hspace{0.5em}{Definición: Distribución de Bernoulli}, toptitle=1mm, coltitle=black, opacityback=0, arc=.35mm, bottomrule=.15mm, toprule=.15mm, colframe=quarto-callout-note-color-frame]

Una variable aleatoria \(X \sim B(p)\) modela un experimento con dos
resultados: éxito (\(X = 1\)) o fracaso (\(X = 0\)):

\[P(X = x) = \begin{cases}
p & \text{si } x = 1 \\
1 - p & \text{si } x = 0 \\
0 & \text{en otro caso}
\end{cases}\]

O compactamente: \(P(X = x) = p^x (1-p)^{1-x}\) para \(x \in \{0, 1\}\)

\textbf{Parámetro:} \(p \in [0, 1]\) (probabilidad de éxito)

\textbf{Esperanza y Varianza:} \[E(X) = p\] \[\text{Var}(X) = p(1-p)\]

\textbf{En R:} \texttt{dbinom(x,\ size=1,\ prob=p)} o
\texttt{rbinom(n,\ size=1,\ prob=p)}

\end{tcolorbox}

\subsection{\texorpdfstring{Distribución Binomial
\(B(n, p)\)}{Distribución Binomial B(n, p)}}\label{distribucion-binomial-bn-p}

\begin{tcolorbox}[enhanced jigsaw, breakable, titlerule=0mm, rightrule=.15mm, leftrule=.75mm, left=2mm, opacitybacktitle=0.6, colback=white, bottomtitle=1mm, colbacktitle=quarto-callout-note-color!10!white, title=\textcolor{quarto-callout-note-color}{\faInfo}\hspace{0.5em}{Definición: Distribución Binomial}, toptitle=1mm, coltitle=black, opacityback=0, arc=.35mm, bottomrule=.15mm, toprule=.15mm, colframe=quarto-callout-note-color-frame]

Una variable aleatoria \(X \sim B(n, p)\) cuenta el número de éxitos en
\(n\) ensayos Bernoulli independientes:

\[P(X = k) = \binom{n}{k} p^k (1-p)^{n-k} \quad \text{para } k = 0, 1, \ldots, n\]

donde \(\binom{n}{k} = \frac{n!}{k!(n-k)!}\)

\textbf{Parámetros:} \(n\) (número de ensayos), \(p\) (probabilidad de
éxito)

\textbf{Esperanza y Varianza:} \[E(X) = n \cdot p\]
\[\text{Var}(X) = n \cdot p \cdot (1-p)\]

\textbf{En R:} - \texttt{dbinom(k,\ size=n,\ prob=p)} --- función de
probabilidad - \texttt{pbinom(k,\ size=n,\ prob=p)} --- función de
distribución - \texttt{rbinom(N,\ size=n,\ prob=p)} --- generar muestras
aleatorias

\end{tcolorbox}

\begin{tcolorbox}[enhanced jigsaw, breakable, titlerule=0mm, rightrule=.15mm, leftrule=.75mm, left=2mm, opacitybacktitle=0.6, colback=white, bottomtitle=1mm, colbacktitle=quarto-callout-tip-color!10!white, title=\textcolor{quarto-callout-tip-color}{\faLightbulb}\hspace{0.5em}{Ejemplo 4.2: Distribución Binomial}, toptitle=1mm, coltitle=black, opacityback=0, arc=.35mm, bottomrule=.15mm, toprule=.15mm, colframe=quarto-callout-tip-color-frame]

Un medicamento cura el 70\% de los pacientes. Si se trata a 10
pacientes, ¿cuál es la probabilidad de que exactamente 7 sanen?

\[P(X = 7) = \binom{10}{7} (0.7)^7 (0.3)^3 = 120 \cdot 0.0823 \cdot 0.027 \approx 0.267\]

En R:

\begin{codelisting}[H]

\caption{\label{lst-semana04-3}Código R}

\centering{

\begin{Shaded}
\begin{Highlighting}[]
\CommentTok{\# Probabilidad de exactamente 7 curaciones}
\FunctionTok{dbinom}\NormalTok{(}\DecValTok{7}\NormalTok{, }\AttributeTok{size=}\DecValTok{10}\NormalTok{, }\AttributeTok{prob=}\FloatTok{0.7}\NormalTok{)}
\end{Highlighting}
\end{Shaded}

}

\end{codelisting}%

\begin{verbatim}
[1] 0.2668279
\end{verbatim}

\begin{codelisting}[H]

\caption{\label{lst-semana04-3}Código R}

\centering{

\begin{Shaded}
\begin{Highlighting}[]
\CommentTok{\# Probabilidad de 7 o menos curaciones}
\FunctionTok{pbinom}\NormalTok{(}\DecValTok{7}\NormalTok{, }\AttributeTok{size=}\DecValTok{10}\NormalTok{, }\AttributeTok{prob=}\FloatTok{0.7}\NormalTok{)}
\end{Highlighting}
\end{Shaded}

}

\end{codelisting}%

\begin{verbatim}
[1] 0.6172172
\end{verbatim}

\begin{codelisting}[H]

\caption{\label{lst-semana04-3}Código R}

\centering{

\begin{Shaded}
\begin{Highlighting}[]
\CommentTok{\# Visualización}
\NormalTok{x }\OtherTok{\textless{}{-}} \DecValTok{0}\SpecialCharTok{:}\DecValTok{10}
\NormalTok{probs }\OtherTok{\textless{}{-}} \FunctionTok{dbinom}\NormalTok{(x, }\AttributeTok{size=}\DecValTok{10}\NormalTok{, }\AttributeTok{prob=}\FloatTok{0.7}\NormalTok{)}
\FunctionTok{barplot}\NormalTok{(probs, }\AttributeTok{names.arg=}\NormalTok{x, }\AttributeTok{main=}\StringTok{"B(10, 0.7)"}\NormalTok{,}
        \AttributeTok{xlab=}\StringTok{"Número de curaciones"}\NormalTok{, }\AttributeTok{ylab=}\StringTok{"Probabilidad"}\NormalTok{)}
\end{Highlighting}
\end{Shaded}

}

\end{codelisting}%

\pandocbounded{\includegraphics[keepaspectratio]{chapters/semana04_files/figure-pdf/unnamed-chunk-3-1.pdf}}

\end{tcolorbox}

\subsection{\texorpdfstring{Distribución de Poisson
\(\text{Po}(\lambda)\)}{Distribución de Poisson \textbackslash text\{Po\}(\textbackslash lambda)}}\label{distribucion-de-poisson-textpolambda}

\begin{tcolorbox}[enhanced jigsaw, breakable, titlerule=0mm, rightrule=.15mm, leftrule=.75mm, left=2mm, opacitybacktitle=0.6, colback=white, bottomtitle=1mm, colbacktitle=quarto-callout-note-color!10!white, title=\textcolor{quarto-callout-note-color}{\faInfo}\hspace{0.5em}{Definición: Distribución de Poisson}, toptitle=1mm, coltitle=black, opacityback=0, arc=.35mm, bottomrule=.15mm, toprule=.15mm, colframe=quarto-callout-note-color-frame]

Una variable aleatoria \(X \sim \text{Po}(\lambda)\) modela el número de
eventos raros que ocurren en un intervalo de tiempo o espacio fijo:

\[P(X = k) = \frac{\lambda^k e^{-\lambda}}{k!} \quad \text{para } k = 0, 1, 2, \ldots\]

\textbf{Parámetro:} \(\lambda > 0\) (tasa de eventos)

\textbf{Esperanza y Varianza:} \[E(X) = \lambda\]
\[\text{Var}(X) = \lambda\]

(Nota: en Poisson, media y varianza son iguales)

\textbf{Propiedad Reproductiva:} Si \(X \sim \text{Po}(\lambda_1)\) e
\(Y \sim \text{Po}(\lambda_2)\) son independientes:
\[X + Y \sim \text{Po}(\lambda_1 + \lambda_2)\]

\textbf{En R:} - \texttt{dpois(k,\ lambda=λ)} --- función de
probabilidad - \texttt{ppois(k,\ lambda=λ)} --- función de distribución
- \texttt{rpois(N,\ lambda=λ)} --- generar muestras aleatorias

\end{tcolorbox}

\begin{tcolorbox}[enhanced jigsaw, breakable, titlerule=0mm, rightrule=.15mm, leftrule=.75mm, left=2mm, opacitybacktitle=0.6, colback=white, bottomtitle=1mm, colbacktitle=quarto-callout-tip-color!10!white, title=\textcolor{quarto-callout-tip-color}{\faLightbulb}\hspace{0.5em}{Ejemplo 4.3: Distribución de Poisson}, toptitle=1mm, coltitle=black, opacityback=0, arc=.35mm, bottomrule=.15mm, toprule=.15mm, colframe=quarto-callout-tip-color-frame]

El número de reacciones adversas en una clínica sigue una distribución
Poisson con \(\lambda = 3\) por día. ¿Cuál es la probabilidad de que
ocurran exactamente 5 reacciones mañana?

\[P(X = 5) = \frac{3^5 e^{-3}}{5!} = \frac{243 \cdot 0.0498}{120} \approx 0.1008\]

En R:

\begin{codelisting}[H]

\caption{\label{lst-semana04-4}Código R}

\centering{

\begin{Shaded}
\begin{Highlighting}[]
\CommentTok{\# Probabilidad de exactamente 5 eventos}
\FunctionTok{dpois}\NormalTok{(}\DecValTok{5}\NormalTok{, }\AttributeTok{lambda=}\DecValTok{3}\NormalTok{)}
\end{Highlighting}
\end{Shaded}

}

\end{codelisting}%

\begin{verbatim}
[1] 0.1008188
\end{verbatim}

\begin{codelisting}[H]

\caption{\label{lst-semana04-4}Código R}

\centering{

\begin{Shaded}
\begin{Highlighting}[]
\CommentTok{\# Probabilidad de 5 o menos eventos}
\FunctionTok{ppois}\NormalTok{(}\DecValTok{5}\NormalTok{, }\AttributeTok{lambda=}\DecValTok{3}\NormalTok{)}
\end{Highlighting}
\end{Shaded}

}

\end{codelisting}%

\begin{verbatim}
[1] 0.9160821
\end{verbatim}

\begin{codelisting}[H]

\caption{\label{lst-semana04-4}Código R}

\centering{

\begin{Shaded}
\begin{Highlighting}[]
\CommentTok{\# Visualización}
\NormalTok{x }\OtherTok{\textless{}{-}} \DecValTok{0}\SpecialCharTok{:}\DecValTok{10}
\NormalTok{probs }\OtherTok{\textless{}{-}} \FunctionTok{dpois}\NormalTok{(x, }\AttributeTok{lambda=}\DecValTok{3}\NormalTok{)}
\FunctionTok{barplot}\NormalTok{(probs, }\AttributeTok{names.arg=}\NormalTok{x, }\AttributeTok{main=}\StringTok{"Poisson(3)"}\NormalTok{,}
        \AttributeTok{xlab=}\StringTok{"Número de eventos"}\NormalTok{, }\AttributeTok{ylab=}\StringTok{"Probabilidad"}\NormalTok{)}
\end{Highlighting}
\end{Shaded}

}

\end{codelisting}%

\pandocbounded{\includegraphics[keepaspectratio]{chapters/semana04_files/figure-pdf/unnamed-chunk-4-1.pdf}}

\end{tcolorbox}

\section{Distribuciones Continuas
Importantes}\label{sec-semana04-continuas}

\subsection{\texorpdfstring{Distribución Uniforme Continua
\(U(a, b)\)}{Distribución Uniforme Continua U(a, b)}}\label{distribucion-uniforme-continua-ua-b}

\begin{tcolorbox}[enhanced jigsaw, breakable, titlerule=0mm, rightrule=.15mm, leftrule=.75mm, left=2mm, opacitybacktitle=0.6, colback=white, bottomtitle=1mm, colbacktitle=quarto-callout-note-color!10!white, title=\textcolor{quarto-callout-note-color}{\faInfo}\hspace{0.5em}{Definición: Distribución Uniforme Continua}, toptitle=1mm, coltitle=black, opacityback=0, arc=.35mm, bottomrule=.15mm, toprule=.15mm, colframe=quarto-callout-note-color-frame]

Una variable aleatoria \(X \sim U(a, b)\) tiene densidad constante en el
intervalo \([a, b]\):

\[f(x) = \begin{cases}
\frac{1}{b-a} & \text{si } a \leq x \leq b \\
0 & \text{en otro caso}
\end{cases}\]

\textbf{Función de Distribución:} \[F(x) = \begin{cases}
0 & \text{si } x < a \\
\frac{x-a}{b-a} & \text{si } a \leq x \leq b \\
1 & \text{si } x > b
\end{cases}\]

\textbf{Parámetros:} \(a, b\) (límites del intervalo)

\textbf{Esperanza y Varianza:} \[E(X) = \frac{a+b}{2}\]
\[\text{Var}(X) = \frac{(b-a)^2}{12}\]

\textbf{En R:} \texttt{dunif(x,\ min=a,\ max=b)},
\texttt{punif(x,\ min=a,\ max=b)}, \texttt{runif(n,\ min=a,\ max=b)}

\end{tcolorbox}

\begin{tcolorbox}[enhanced jigsaw, breakable, titlerule=0mm, rightrule=.15mm, leftrule=.75mm, left=2mm, opacitybacktitle=0.6, colback=white, bottomtitle=1mm, colbacktitle=quarto-callout-tip-color!10!white, title=\textcolor{quarto-callout-tip-color}{\faLightbulb}\hspace{0.5em}{Ejemplo: Tiempo de espera}, toptitle=1mm, coltitle=black, opacityback=0, arc=.35mm, bottomrule=.15mm, toprule=.15mm, colframe=quarto-callout-tip-color-frame]

Un paciente debe esperar entre 0 y 30 minutos para ser atendido
(\(X \sim U(0, 30)\)).

\[E(X) = \frac{0+30}{2} = 15 \text{ minutos}\]
\[\text{Var}(X) = \frac{(30-0)^2}{12} = 75\]

Probabilidad de esperar menos de 10 minutos:
\[P(X \leq 10) = \frac{10-0}{30-0} = \frac{1}{3} \approx 0.333\]

En R:

\begin{codelisting}[H]

\caption{\label{lst-semana04-5}Código R}

\centering{

\begin{Shaded}
\begin{Highlighting}[]
\CommentTok{\# Densidad}
\FunctionTok{curve}\NormalTok{(}\FunctionTok{dunif}\NormalTok{(x, }\AttributeTok{min=}\DecValTok{0}\NormalTok{, }\AttributeTok{max=}\DecValTok{30}\NormalTok{), }\AttributeTok{from=}\SpecialCharTok{{-}}\DecValTok{5}\NormalTok{, }\AttributeTok{to=}\DecValTok{35}\NormalTok{,}
      \AttributeTok{xlab=}\StringTok{"Minutos"}\NormalTok{, }\AttributeTok{ylab=}\StringTok{"Densidad"}\NormalTok{)}
\end{Highlighting}
\end{Shaded}

}

\end{codelisting}%

\pandocbounded{\includegraphics[keepaspectratio]{chapters/semana04_files/figure-pdf/unnamed-chunk-5-1.pdf}}

\begin{codelisting}[H]

\caption{\label{lst-semana04-5}Código R}

\centering{

\begin{Shaded}
\begin{Highlighting}[]
\CommentTok{\# P(X ≤ 10)}
\FunctionTok{punif}\NormalTok{(}\DecValTok{10}\NormalTok{, }\AttributeTok{min=}\DecValTok{0}\NormalTok{, }\AttributeTok{max=}\DecValTok{30}\NormalTok{)}
\end{Highlighting}
\end{Shaded}

}

\end{codelisting}%

\begin{verbatim}
[1] 0.3333333
\end{verbatim}

\end{tcolorbox}

\subsection{\texorpdfstring{Distribución Exponencial
\(\text{Exp}(\lambda)\)}{Distribución Exponencial \textbackslash text\{Exp\}(\textbackslash lambda)}}\label{distribucion-exponencial-textexplambda}

\begin{tcolorbox}[enhanced jigsaw, breakable, titlerule=0mm, rightrule=.15mm, leftrule=.75mm, left=2mm, opacitybacktitle=0.6, colback=white, bottomtitle=1mm, colbacktitle=quarto-callout-note-color!10!white, title=\textcolor{quarto-callout-note-color}{\faInfo}\hspace{0.5em}{Definición: Distribución Exponencial}, toptitle=1mm, coltitle=black, opacityback=0, arc=.35mm, bottomrule=.15mm, toprule=.15mm, colframe=quarto-callout-note-color-frame]

Una variable aleatoria \(X \sim \text{Exp}(\lambda)\) modela el tiempo
hasta el siguiente evento en un proceso de Poisson:

\[f(x) = \begin{cases}
\lambda e^{-\lambda x} & \text{si } x \geq 0 \\
0 & \text{si } x < 0
\end{cases}\]

\textbf{Función de Distribución:} \[F(x) = \begin{cases}
1 - e^{-\lambda x} & \text{si } x \geq 0 \\
0 & \text{si } x < 0
\end{cases}\]

\textbf{Parámetro:} \(\lambda > 0\) (tasa de evento)

\textbf{Esperanza y Varianza:} \[E(X) = \frac{1}{\lambda}\]
\[\text{Var}(X) = \frac{1}{\lambda^2}\]

\textbf{Propiedad de Falta de Memoria:}
\(P(X > s+t | X > s) = P(X > t)\)

\textbf{En R:} \texttt{dexp(x,\ rate=λ)}, \texttt{pexp(x,\ rate=λ)},
\texttt{rexp(n,\ rate=λ)}

\end{tcolorbox}

\begin{tcolorbox}[enhanced jigsaw, breakable, titlerule=0mm, rightrule=.15mm, leftrule=.75mm, left=2mm, opacitybacktitle=0.6, colback=white, bottomtitle=1mm, colbacktitle=quarto-callout-tip-color!10!white, title=\textcolor{quarto-callout-tip-color}{\faLightbulb}\hspace{0.5em}{Ejemplo: Tiempo de supervivencia}, toptitle=1mm, coltitle=black, opacityback=0, arc=.35mm, bottomrule=.15mm, toprule=.15mm, colframe=quarto-callout-tip-color-frame]

El tiempo de supervivencia de células en cultivo sigue
\(\text{Exp}(\lambda = 0.1)\) (en horas).

Tiempo medio de supervivencia: \(E(X) = 1/0.1 = 10\) horas

Probabilidad de que una célula sobreviva más de 5 horas:
\[P(X > 5) = e^{-0.1 \cdot 5} = e^{-0.5} \approx 0.606\]

En R:

\begin{codelisting}[H]

\caption{\label{lst-semana04-6}Código R}

\centering{

\begin{Shaded}
\begin{Highlighting}[]
\CommentTok{\# Densidad}
\FunctionTok{curve}\NormalTok{(}\FunctionTok{dexp}\NormalTok{(x, }\AttributeTok{rate=}\FloatTok{0.1}\NormalTok{), }\AttributeTok{from=}\DecValTok{0}\NormalTok{, }\AttributeTok{to=}\DecValTok{40}\NormalTok{,}
      \AttributeTok{xlab=}\StringTok{"Horas"}\NormalTok{, }\AttributeTok{ylab=}\StringTok{"Densidad"}\NormalTok{)}
\end{Highlighting}
\end{Shaded}

}

\end{codelisting}%

\pandocbounded{\includegraphics[keepaspectratio]{chapters/semana04_files/figure-pdf/unnamed-chunk-6-1.pdf}}

\begin{codelisting}[H]

\caption{\label{lst-semana04-6}Código R}

\centering{

\begin{Shaded}
\begin{Highlighting}[]
\CommentTok{\# P(X \textgreater{} 5)}
\DecValTok{1} \SpecialCharTok{{-}} \FunctionTok{pexp}\NormalTok{(}\DecValTok{5}\NormalTok{, }\AttributeTok{rate=}\FloatTok{0.1}\NormalTok{)}
\end{Highlighting}
\end{Shaded}

}

\end{codelisting}%

\begin{verbatim}
[1] 0.6065307
\end{verbatim}

\begin{codelisting}[H]

\caption{\label{lst-semana04-6}Código R}

\centering{

\begin{Shaded}
\begin{Highlighting}[]
\CommentTok{\# O equivalentemente:}
\FunctionTok{pexp}\NormalTok{(}\DecValTok{5}\NormalTok{, }\AttributeTok{rate=}\FloatTok{0.1}\NormalTok{, }\AttributeTok{lower.tail=}\ConstantTok{FALSE}\NormalTok{)}
\end{Highlighting}
\end{Shaded}

}

\end{codelisting}%

\begin{verbatim}
[1] 0.6065307
\end{verbatim}

\end{tcolorbox}

\subsection{\texorpdfstring{Distribución Normal
\(N(\mu, \sigma^2)\)}{Distribución Normal N(\textbackslash mu, \textbackslash sigma\^{}2)}}\label{distribucion-normal-nmu-sigma2}

\begin{tcolorbox}[enhanced jigsaw, breakable, titlerule=0mm, rightrule=.15mm, leftrule=.75mm, left=2mm, opacitybacktitle=0.6, colback=white, bottomtitle=1mm, colbacktitle=quarto-callout-note-color!10!white, title=\textcolor{quarto-callout-note-color}{\faInfo}\hspace{0.5em}{Definición: Distribución Normal}, toptitle=1mm, coltitle=black, opacityback=0, arc=.35mm, bottomrule=.15mm, toprule=.15mm, colframe=quarto-callout-note-color-frame]

Una variable aleatoria \(X \sim N(\mu, \sigma^2)\) tiene la función de
densidad:

\[f(x) = \frac{1}{\sigma\sqrt{2\pi}} \exp\left\{-\frac{(x-\mu)^2}{2\sigma^2}\right\} \quad \text{para } -\infty < x < +\infty\]

\textbf{Parámetros:} \(\mu\) (media), \(\sigma^2\) (varianza);
\(\sigma > 0\)

\textbf{Propiedades:} - Simétrica alrededor de \(\mu\) - Media = Mediana
= Moda = \(\mu\) - \(E(X) = \mu\) - \(\text{Var}(X) = \sigma^2\)

\textbf{Función de Distribución:} No tiene forma cerrada; se usa
notación \(\Phi(z)\) para la CDF estándar.

\textbf{En R:} \texttt{dnorm(x,\ mean=μ,\ sd=σ)},
\texttt{pnorm(x,\ mean=μ,\ sd=σ)}, \texttt{rnorm(n,\ mean=μ,\ sd=σ)}

\end{tcolorbox}

\begin{tcolorbox}[enhanced jigsaw, breakable, titlerule=0mm, rightrule=.15mm, leftrule=.75mm, left=2mm, opacitybacktitle=0.6, colback=white, bottomtitle=1mm, colbacktitle=quarto-callout-warning-color!10!white, title=\textcolor{quarto-callout-warning-color}{\faExclamationTriangle}\hspace{0.5em}{Distribución Normal Estándar \(N(0, 1)\)}, toptitle=1mm, coltitle=black, opacityback=0, arc=.35mm, bottomrule=.15mm, toprule=.15mm, colframe=quarto-callout-warning-color-frame]

La \textbf{distribución normal estándar} tiene \(\mu = 0\) y
\(\sigma = 1\):

\[\varphi(z) = \frac{1}{\sqrt{2\pi}} e^{-z^2/2}\]

\textbf{Transformación Z (Estandarización):}
\[Z = \frac{X - \mu}{\sigma}\]

Si \(X \sim N(\mu, \sigma^2)\), entonces \(Z \sim N(0, 1)\)

Esta transformación es fundamental porque: - Permite usar tablas únicas
de probabilidades - Valores de \(Z\) son directamente comparables entre
variables diferentes

\end{tcolorbox}

\begin{tcolorbox}[enhanced jigsaw, breakable, titlerule=0mm, rightrule=.15mm, leftrule=.75mm, left=2mm, opacitybacktitle=0.6, colback=white, bottomtitle=1mm, colbacktitle=quarto-callout-tip-color!10!white, title=\textcolor{quarto-callout-tip-color}{\faLightbulb}\hspace{0.5em}{Ejemplo: Medidas antropométricas}, toptitle=1mm, coltitle=black, opacityback=0, arc=.35mm, bottomrule=.15mm, toprule=.15mm, colframe=quarto-callout-tip-color-frame]

El peso de adultos sigue \(N(\mu = 75, \sigma = 12)\) kg.

¿Cuál es la probabilidad de que una persona pese menos de 90 kg?

\[P(X \leq 90) = P\left(Z \leq \frac{90-75}{12}\right) = P(Z \leq 1.25)\]

En R:

\begin{codelisting}[H]

\caption{\label{lst-semana04-7}Código R}

\centering{

\begin{Shaded}
\begin{Highlighting}[]
\CommentTok{\# Directamente con parámetros}
\FunctionTok{pnorm}\NormalTok{(}\DecValTok{90}\NormalTok{, }\AttributeTok{mean=}\DecValTok{75}\NormalTok{, }\AttributeTok{sd=}\DecValTok{12}\NormalTok{)}
\end{Highlighting}
\end{Shaded}

}

\end{codelisting}%

\begin{verbatim}
[1] 0.8943502
\end{verbatim}

\begin{codelisting}[H]

\caption{\label{lst-semana04-7}Código R}

\centering{

\begin{Shaded}
\begin{Highlighting}[]
\CommentTok{\# O usando estandarización}
\NormalTok{z }\OtherTok{\textless{}{-}}\NormalTok{ (}\DecValTok{90} \SpecialCharTok{{-}} \DecValTok{75}\NormalTok{) }\SpecialCharTok{/} \DecValTok{12}
\FunctionTok{pnorm}\NormalTok{(z)  }\CommentTok{\# Usa N(0,1) por defecto}
\end{Highlighting}
\end{Shaded}

}

\end{codelisting}%

\begin{verbatim}
[1] 0.8943502
\end{verbatim}

\begin{codelisting}[H]

\caption{\label{lst-semana04-7}Código R}

\centering{

\begin{Shaded}
\begin{Highlighting}[]
\CommentTok{\# Visualización}
\FunctionTok{curve}\NormalTok{(}\FunctionTok{dnorm}\NormalTok{(x, }\AttributeTok{mean=}\DecValTok{75}\NormalTok{, }\AttributeTok{sd=}\DecValTok{12}\NormalTok{), }\AttributeTok{from=}\DecValTok{30}\NormalTok{, }\AttributeTok{to=}\DecValTok{120}\NormalTok{,}
      \AttributeTok{xlab=}\StringTok{"Peso (kg)"}\NormalTok{, }\AttributeTok{ylab=}\StringTok{"Densidad"}\NormalTok{,}
      \AttributeTok{main=}\StringTok{"N(75, 12)"}\NormalTok{)}
\FunctionTok{abline}\NormalTok{(}\AttributeTok{v=}\DecValTok{90}\NormalTok{, }\AttributeTok{col=}\StringTok{"red"}\NormalTok{, }\AttributeTok{lty=}\DecValTok{2}\NormalTok{)}
\end{Highlighting}
\end{Shaded}

}

\end{codelisting}%

\pandocbounded{\includegraphics[keepaspectratio]{chapters/semana04_files/figure-pdf/unnamed-chunk-7-1.pdf}}

La probabilidad es aproximadamente \textbf{0.894} (89.4\%).

\end{tcolorbox}

\begin{tcolorbox}[enhanced jigsaw, breakable, titlerule=0mm, rightrule=.15mm, leftrule=.75mm, left=2mm, opacitybacktitle=0.6, colback=white, bottomtitle=1mm, colbacktitle=quarto-callout-tip-color!10!white, title=\textcolor{quarto-callout-tip-color}{\faLightbulb}\hspace{0.5em}{Ejemplo: Interpretación de z-scores}, toptitle=1mm, coltitle=black, opacityback=0, arc=.35mm, bottomrule=.15mm, toprule=.15mm, colframe=quarto-callout-tip-color-frame]

Para pacientes con presión sistólica \(N(\mu = 120, \sigma = 15)\) mmHg:

\begin{itemize}
\tightlist
\item
  Paciente A con presión = 120 mmHg: \(z = 0\) (media)
\item
  Paciente B con presión = 135 mmHg: \(z = 1\) (1 DE arriba de la media)
\item
  Paciente C con presión = 105 mmHg: \(z = -1\) (1 DE abajo de la media)
\item
  Paciente D con presión = 150 mmHg: \(z = 2\) (2 DE arriba de la media)
\end{itemize}

Los z-scores permiten comparar directamente aunque las variables
originales tengan unidades diferentes.

\end{tcolorbox}

\subsection{Aproximación Normal a la Binomial y Corrección de
Continuidad}\label{sec-semana04-normal-binom}

Históricamente, cuando el cálculo manual o las tablas estadísticas eran
las principales herramientas, la aproximación de distribuciones
discretas (como la binomial) mediante distribuciones continuas (como la
normal) era una técnica invaluable para estimar probabilidades. Esta
aproximación, una aplicación directa del Teorema Central del Límite,
permitía simplificar cálculos complejos. Sin embargo, con el
advenimiento del software estadístico moderno, que utiliza algoritmos
exactos (como la Función Beta Incompleta Regularizada para la binomial)
con alta precisión computacional, el uso de estas aproximaciones para el
cálculo de probabilidades exactas se ha vuelto menos recomendable.
Aunque la aproximación sigue siendo fundamental para la comprensión
teórica y para la derivación de métodos como los intervalos de confianza
o tests Z para proporciones, para la obtención de valores de
probabilidad específicos, se prefiere el cálculo directo para evitar
posibles inexactitudes (véase \citet{Saad2020}).

La \textbf{corrección de continuidad} es un ajuste que mejora la
aproximación al compensar el paso de una variable discreta a una
continua. Cada barra del histograma binomial en \(k\) ocupa el intervalo
\([k - 0.5,\, k + 0.5]\).

\begin{tcolorbox}[enhanced jigsaw, breakable, titlerule=0mm, rightrule=.15mm, leftrule=.75mm, left=2mm, opacitybacktitle=0.6, colback=white, bottomtitle=1mm, colbacktitle=quarto-callout-note-color!10!white, title=\textcolor{quarto-callout-note-color}{\faInfo}\hspace{0.5em}{Definición: Aproximación Normal a la Binomial}, toptitle=1mm, coltitle=black, opacityback=0, arc=.35mm, bottomrule=.15mm, toprule=.15mm, colframe=quarto-callout-note-color-frame]

Si \(X \sim B(n, p)\) con \(np \geq 5\) y \(n(1-p) \geq 5\), entonces:

\[X \overset{\text{aprox.}}{\sim} N\!\left(\mu = np,\; \sigma^2 = np(1-p)\right)\]

La \textbf{corrección de continuidad} mejora la aproximación al ajustar
\(\pm 0.5\) para compensar el paso de una variable discreta a una
continua. Cada barra del histograma binomial en \(k\) ocupa el intervalo
\([k - 0.5,\, k + 0.5]\).

\begin{longtable}[]{@{}
  >{\raggedright\arraybackslash}p{(\linewidth - 4\tabcolsep) * \real{0.3750}}
  >{\centering\arraybackslash}p{(\linewidth - 4\tabcolsep) * \real{0.2679}}
  >{\centering\arraybackslash}p{(\linewidth - 4\tabcolsep) * \real{0.3571}}@{}}
\toprule\noalign{}
\begin{minipage}[b]{\linewidth}\raggedright
Probabilidad exacta
\end{minipage} & \begin{minipage}[b]{\linewidth}\centering
Sin corrección
\end{minipage} & \begin{minipage}[b]{\linewidth}\centering
\textbf{Con corrección}
\end{minipage} \\
\midrule\noalign{}
\endhead
\bottomrule\noalign{}
\endlastfoot
\(P(X \leq k)\) & \(\Phi\!\left(\dfrac{k - \mu}{\sigma}\right)\) &
\(\Phi\!\left(\dfrac{k + 0.5 - \mu}{\sigma}\right)\) \\
\(P(X \geq k)\) & \(1 - \Phi\!\left(\dfrac{k - \mu}{\sigma}\right)\) &
\(1 - \Phi\!\left(\dfrac{k - 0.5 - \mu}{\sigma}\right)\) \\
\(P(X = k)\) & \(\approx 0\) &
\(\Phi\!\left(\dfrac{k + 0.5 - \mu}{\sigma}\right) - \Phi\!\left(\dfrac{k - 0.5 - \mu}{\sigma}\right)\) \\
\end{longtable}

\end{tcolorbox}

\begin{tcolorbox}[enhanced jigsaw, breakable, titlerule=0mm, rightrule=.15mm, leftrule=.75mm, left=2mm, opacitybacktitle=0.6, colback=white, bottomtitle=1mm, colbacktitle=quarto-callout-warning-color!10!white, title=\textcolor{quarto-callout-warning-color}{\faExclamationTriangle}\hspace{0.5em}{Condición de Aplicación}, toptitle=1mm, coltitle=black, opacityback=0, arc=.35mm, bottomrule=.15mm, toprule=.15mm, colframe=quarto-callout-warning-color-frame]

La aproximación normal a la binomial es válida cuando:
\[np \geq 5 \quad \text{y} \quad n(1-p) \geq 5\]

Para \(p\) cercano a 0 o 1 con \(n\) pequeño, usar la binomial exacta
(\texttt{pbinom()}).

\end{tcolorbox}

\begin{tcolorbox}[enhanced jigsaw, breakable, titlerule=0mm, rightrule=.15mm, leftrule=.75mm, left=2mm, opacitybacktitle=0.6, colback=white, bottomtitle=1mm, colbacktitle=quarto-callout-tip-color!10!white, title=\textcolor{quarto-callout-tip-color}{\faLightbulb}\hspace{0.5em}{Ejemplo: Aproximación Normal a \(B(20, 0.4)\) con Corrección de
Continuidad}, toptitle=1mm, coltitle=black, opacityback=0, arc=.35mm, bottomrule=.15mm, toprule=.15mm, colframe=quarto-callout-tip-color-frame]

Un fármaco es efectivo en el 40\% de los pacientes. En un ensayo con
\(n = 20\) pacientes, ¿cuál es la probabilidad de que a lo sumo 10
respondan?

\textbf{Parámetros:} \(\mu = 20 \times 0.4 = 8\),
\(\sigma = \sqrt{20 \times 0.4 \times 0.6} \approx 2.19\)

\begin{longtable}[]{@{}
  >{\raggedright\arraybackslash}p{(\linewidth - 4\tabcolsep) * \real{0.0909}}
  >{\raggedright\arraybackslash}p{(\linewidth - 4\tabcolsep) * \real{0.4091}}
  >{\raggedright\arraybackslash}p{(\linewidth - 4\tabcolsep) * \real{0.5000}}@{}}
\toprule\noalign{}
\begin{minipage}[b]{\linewidth}\raggedright
\end{minipage} & \begin{minipage}[b]{\linewidth}\raggedright
Cálculo
\end{minipage} & \begin{minipage}[b]{\linewidth}\raggedright
Resultado
\end{minipage} \\
\midrule\noalign{}
\endhead
\bottomrule\noalign{}
\endlastfoot
\textbf{Exacto (binomial)} & \texttt{pbinom(10,\ 20,\ 0.4)} &
\textbf{0.8724} \\
Normal sin corrección &
\(\Phi\!\left(\frac{10 - 8}{2.19}\right) = \Phi(0.913)\) & 0.8186 \\
\textbf{Normal con corrección} &
\(\Phi\!\left(\frac{10.5 - 8}{2.19}\right) = \Phi(1.142)\) &
\textbf{0.8729} \\
\end{longtable}

La corrección reduce el error de 5.4 pp a 0.05 pp.

\begin{codelisting}[H]

\caption{\label{lst-semana04-8}Código R}

\centering{

\begin{Shaded}
\begin{Highlighting}[]
\NormalTok{n }\OtherTok{\textless{}{-}} \DecValTok{20}\NormalTok{; p }\OtherTok{\textless{}{-}} \FloatTok{0.4}
\NormalTok{mu }\OtherTok{\textless{}{-}}\NormalTok{ n}\SpecialCharTok{*}\NormalTok{p; sigma }\OtherTok{\textless{}{-}} \FunctionTok{sqrt}\NormalTok{(n}\SpecialCharTok{*}\NormalTok{p}\SpecialCharTok{*}\NormalTok{(}\DecValTok{1}\SpecialCharTok{{-}}\NormalTok{p))}
\NormalTok{k  }\OtherTok{\textless{}{-}} \DecValTok{10}

\FunctionTok{cat}\NormalTok{(}\StringTok{"Exacto:               "}\NormalTok{, }\FunctionTok{pbinom}\NormalTok{(k, n, p), }\StringTok{"}\SpecialCharTok{\textbackslash{}n}\StringTok{"}\NormalTok{)}
\end{Highlighting}
\end{Shaded}

}

\end{codelisting}%

\begin{verbatim}
Exacto:                0.8724788 
\end{verbatim}

\begin{codelisting}[H]

\caption{\label{lst-semana04-8}Código R}

\centering{

\begin{Shaded}
\begin{Highlighting}[]
\FunctionTok{cat}\NormalTok{(}\StringTok{"Normal sin corrección:"}\NormalTok{, }\FunctionTok{pnorm}\NormalTok{(k, mu, sigma), }\StringTok{"}\SpecialCharTok{\textbackslash{}n}\StringTok{"}\NormalTok{)}
\end{Highlighting}
\end{Shaded}

}

\end{codelisting}%

\begin{verbatim}
Normal sin corrección: 0.8193448 
\end{verbatim}

\begin{codelisting}[H]

\caption{\label{lst-semana04-8}Código R}

\centering{

\begin{Shaded}
\begin{Highlighting}[]
\FunctionTok{cat}\NormalTok{(}\StringTok{"Normal con corrección:"}\NormalTok{, }\FunctionTok{pnorm}\NormalTok{(k }\SpecialCharTok{+} \FloatTok{0.5}\NormalTok{, mu, sigma), }\StringTok{"}\SpecialCharTok{\textbackslash{}n}\StringTok{"}\NormalTok{)}
\end{Highlighting}
\end{Shaded}

}

\end{codelisting}%

\begin{verbatim}
Normal con corrección: 0.8730835 
\end{verbatim}

La figura siguiente muestra las barras de la binomial con la curva
normal superpuesta. Las líneas punteadas delimitan la barra \(k = 10\)
en \([9.5,\, 10.5]\), que es el intervalo integrado con la corrección de
continuidad.

\begin{codelisting}[H]

\caption{\label{lst-fig-normal-binom}Código R}

\centering{

\begin{Shaded}
\begin{Highlighting}[]
\NormalTok{n }\OtherTok{\textless{}{-}} \DecValTok{20}\NormalTok{; p }\OtherTok{\textless{}{-}} \FloatTok{0.4}
\NormalTok{mu }\OtherTok{\textless{}{-}}\NormalTok{ n }\SpecialCharTok{*}\NormalTok{ p; sigma }\OtherTok{\textless{}{-}} \FunctionTok{sqrt}\NormalTok{(n }\SpecialCharTok{*}\NormalTok{ p }\SpecialCharTok{*}\NormalTok{ (}\DecValTok{1} \SpecialCharTok{{-}}\NormalTok{ p))}
\NormalTok{x\_vals  }\OtherTok{\textless{}{-}} \DecValTok{0}\SpecialCharTok{:}\NormalTok{n}
\NormalTok{probs\_b }\OtherTok{\textless{}{-}} \FunctionTok{dbinom}\NormalTok{(x\_vals, n, p)}

\CommentTok{\# Histograma tipo barras binomiales}
\NormalTok{bp }\OtherTok{\textless{}{-}} \FunctionTok{barplot}\NormalTok{(probs\_b,}
              \AttributeTok{names.arg =}\NormalTok{ x\_vals,}
              \AttributeTok{col    =} \FunctionTok{ifelse}\NormalTok{(x\_vals }\SpecialCharTok{==} \DecValTok{10}\NormalTok{, }\StringTok{"\#f39200"}\NormalTok{, }\StringTok{"\#cce0f5"}\NormalTok{),}
              \AttributeTok{border =} \StringTok{"white"}\NormalTok{,}
              \AttributeTok{xlab   =} \StringTok{"k (número de respuestas)"}\NormalTok{,}
              \AttributeTok{ylab   =} \StringTok{"Probabilidad / Densidad"}\NormalTok{,}
              \AttributeTok{main   =} \FunctionTok{expression}\NormalTok{(}\StringTok{"Aproximación normal a "} \SpecialCharTok{*} \FunctionTok{B}\NormalTok{(}\DecValTok{20}\NormalTok{, }\FloatTok{0.4}\NormalTok{)),}
              \AttributeTok{ylim   =} \FunctionTok{c}\NormalTok{(}\DecValTok{0}\NormalTok{, }\FunctionTok{max}\NormalTok{(probs\_b) }\SpecialCharTok{*} \FloatTok{1.25}\NormalTok{),}
              \AttributeTok{space  =} \DecValTok{0}\NormalTok{)}

\CommentTok{\# Curva normal superpuesta (coordenadas del barplot: posición = k + 0.5)}
\NormalTok{x\_c }\OtherTok{\textless{}{-}} \FunctionTok{seq}\NormalTok{(}\SpecialCharTok{{-}}\FloatTok{0.5}\NormalTok{, n }\SpecialCharTok{+} \FloatTok{0.5}\NormalTok{, }\AttributeTok{length.out =} \DecValTok{400}\NormalTok{)}
\FunctionTok{lines}\NormalTok{(x\_c }\SpecialCharTok{+} \FloatTok{0.5}\NormalTok{, }\FunctionTok{dnorm}\NormalTok{(x\_c, mu, sigma), }\AttributeTok{col =} \StringTok{"\#003087"}\NormalTok{, }\AttributeTok{lwd =} \FloatTok{2.5}\NormalTok{)}

\CommentTok{\# Límites corrección de continuidad para k = 10}
\FunctionTok{abline}\NormalTok{(}\AttributeTok{v =} \FunctionTok{c}\NormalTok{(}\DecValTok{10}\NormalTok{, }\DecValTok{11}\NormalTok{), }\AttributeTok{col =} \StringTok{"\#f39200"}\NormalTok{, }\AttributeTok{lty =} \DecValTok{2}\NormalTok{, }\AttributeTok{lwd =} \FloatTok{1.8}\NormalTok{)}

\FunctionTok{legend}\NormalTok{(}\StringTok{"topright"}\NormalTok{,}
       \AttributeTok{legend =} \FunctionTok{c}\NormalTok{(}\FunctionTok{sprintf}\NormalTok{(}\StringTok{"B(\%d, \%.1f)  exacto"}\NormalTok{, n, p),}
                  \FunctionTok{sprintf}\NormalTok{(}\StringTok{"N(\%.0f, \%.2f²) aprox."}\NormalTok{, mu, sigma),}
                  \StringTok{"Límites ±0.5  (k = 10)"}\NormalTok{),}
       \AttributeTok{fill   =} \FunctionTok{c}\NormalTok{(}\StringTok{"\#cce0f5"}\NormalTok{, }\ConstantTok{NA}\NormalTok{, }\ConstantTok{NA}\NormalTok{),}
       \AttributeTok{lty    =} \FunctionTok{c}\NormalTok{(}\ConstantTok{NA}\NormalTok{, }\DecValTok{1}\NormalTok{, }\DecValTok{2}\NormalTok{),}
       \AttributeTok{lwd    =} \FunctionTok{c}\NormalTok{(}\ConstantTok{NA}\NormalTok{, }\FloatTok{2.5}\NormalTok{, }\FloatTok{1.8}\NormalTok{),}
       \AttributeTok{col    =} \FunctionTok{c}\NormalTok{(}\ConstantTok{NA}\NormalTok{, }\StringTok{"\#003087"}\NormalTok{, }\StringTok{"\#f39200"}\NormalTok{),}
       \AttributeTok{bty    =} \StringTok{"n"}\NormalTok{, }\AttributeTok{cex =} \FloatTok{0.85}\NormalTok{)}
\end{Highlighting}
\end{Shaded}

}

\end{codelisting}%

\begin{figure}[H]

{\centering \pandocbounded{\includegraphics[keepaspectratio]{chapters/semana04_files/figure-pdf/unnamed-chunk-9-1.pdf}}

}

\caption{Aproximación normal a B(20, 0.4). Las barras son probabilidades
binomiales exactas; la curva azul es la densidad N(8, 2.19²). La barra
naranja (k = 10) y las líneas punteadas ilustran la corrección de
continuidad: el área bajo la curva normal en {[}9.5, 10.5{]} corresponde
a P(X = 10).}

\end{figure}%

\end{tcolorbox}

\section{Distribuciones para Muestreo}\label{sec-semana04-muestreo}

\subsection{\texorpdfstring{Distribución Chi-Cuadrado
\(\chi^2_k\)}{Distribución Chi-Cuadrado \textbackslash chi\^{}2\_k}}\label{distribucion-chi-cuadrado-chi2_k}

\begin{tcolorbox}[enhanced jigsaw, breakable, titlerule=0mm, rightrule=.15mm, leftrule=.75mm, left=2mm, opacitybacktitle=0.6, colback=white, bottomtitle=1mm, colbacktitle=quarto-callout-note-color!10!white, title=\textcolor{quarto-callout-note-color}{\faInfo}\hspace{0.5em}{Definición: Distribución Chi-Cuadrado}, toptitle=1mm, coltitle=black, opacityback=0, arc=.35mm, bottomrule=.15mm, toprule=.15mm, colframe=quarto-callout-note-color-frame]

Si \(Z_1, Z_2, \ldots, Z_k\) son variables normales estándar
independientes, entonces:

\[X = Z_1^2 + Z_2^2 + \cdots + Z_k^2 \sim \chi^2_k\]

\textbf{Parámetro:} \(k\) (grados de libertad)

\textbf{Propiedades:} - Solo toma valores positivos - \(E(X) = k\) -
\(\text{Var}(X) = 2k\) - A medida que \(k\) aumenta, se aproxima a una
distribución normal (TCL)

\textbf{En R:} \texttt{dchisq(x,\ df=k)}, \texttt{pchisq(x,\ df=k)},
\texttt{rchisq(n,\ df=k)}

\textbf{Usos:} Pruebas de bondad de ajuste, tablas de contingencia,
intervalos de confianza para varianza.

La morfologia de la distribucion chi-cuadrado depende del numero de
grados de libertad: para \(k=1\) es altamente asimetrica, pero a medida
que \(k\) aumenta se vuelve mas simetrica y se aproxima a una normal.

Vease la figura siguiente que muestra la densidad de \(\chi^2_k\) para
\(k = 1, 3, 5, 10\). Para \(k=1\) la distribucion es muy asimetrica, con
una cola larga hacia la derecha. Para \(k=3\) y \(k=5\) la distribucion
se vuelve mas simetrica, aunque sigue teniendo una cola hacia la
derecha. Para \(k=10\) la distribucion ya se aproxima bastante a una
normal, aunque sigue siendo ligeramente asimetrica.

\begin{codelisting}[H]

\caption{\label{lst-fig-chi-cuadrado}Código R}

\centering{

\begin{Shaded}
\begin{Highlighting}[]
\FunctionTok{library}\NormalTok{(ggplot2)}
\NormalTok{k\_values }\OtherTok{\textless{}{-}} \FunctionTok{c}\NormalTok{(}\DecValTok{1}\NormalTok{, }\DecValTok{3}\NormalTok{, }\DecValTok{5}\NormalTok{, }\DecValTok{10}\NormalTok{)}
\NormalTok{x }\OtherTok{\textless{}{-}} \FunctionTok{seq}\NormalTok{(}\DecValTok{0}\NormalTok{, }\DecValTok{30}\NormalTok{, }\AttributeTok{length.out =} \DecValTok{400}\NormalTok{)}
\CommentTok{\# Configuración de paneles}
\FunctionTok{par}\NormalTok{(}\AttributeTok{mfrow =} \FunctionTok{c}\NormalTok{(}\DecValTok{2}\NormalTok{, }\DecValTok{2}\NormalTok{),}
    \AttributeTok{mar =} \FunctionTok{c}\NormalTok{(}\DecValTok{4}\NormalTok{, }\DecValTok{4}\NormalTok{, }\DecValTok{2}\NormalTok{, }\DecValTok{1}\NormalTok{))}
\ControlFlowTok{for}\NormalTok{ (k }\ControlFlowTok{in}\NormalTok{ k\_values) \{}
\NormalTok{  y }\OtherTok{\textless{}{-}} \FunctionTok{dchisq}\NormalTok{(x, }\AttributeTok{df =}\NormalTok{ k)}
  \FunctionTok{plot}\NormalTok{(}
\NormalTok{    x, y,}
    \AttributeTok{type =} \StringTok{"l"}\NormalTok{,}
    \AttributeTok{lwd =} \DecValTok{2}\NormalTok{,}
    \AttributeTok{col =} \StringTok{"\#003087"}\NormalTok{,}
    \AttributeTok{main =} \FunctionTok{paste}\NormalTok{(}\StringTok{"k ="}\NormalTok{, k),}
    \AttributeTok{xlab =} \StringTok{"x"}\NormalTok{,}
    \AttributeTok{ylab =} \StringTok{"Densidad"}
\NormalTok{  )}
  \FunctionTok{polygon}\NormalTok{(}
    \FunctionTok{c}\NormalTok{(x, }\FunctionTok{rev}\NormalTok{(x)),}
    \FunctionTok{c}\NormalTok{(y, }\FunctionTok{rep}\NormalTok{(}\DecValTok{0}\NormalTok{, }\FunctionTok{length}\NormalTok{(y))),}
    \AttributeTok{col =} \FunctionTok{adjustcolor}\NormalTok{(}\StringTok{"\#4F81BD"}\NormalTok{, }\AttributeTok{alpha.f =} \FloatTok{0.35}\NormalTok{),}
    \AttributeTok{border =} \ConstantTok{NA}
\NormalTok{  )}
  \FunctionTok{lines}\NormalTok{(x, y, }\AttributeTok{lwd =} \DecValTok{2}\NormalTok{, }\AttributeTok{col =} \StringTok{"\#003087"}\NormalTok{)}
\NormalTok{\}}
\end{Highlighting}
\end{Shaded}

}

\end{codelisting}%

\begin{figure}[H]

{\centering \pandocbounded{\includegraphics[keepaspectratio]{chapters/semana04_files/figure-pdf/unnamed-chunk-10-1.pdf}}

}

\caption{c(``k = 1'', ``k = 3'', ``k = 5'', ``k = 10'')}

\end{figure}%

Densidad de la distribución chi-cuadrado para diferentes grados de
libertad. A medida que k aumenta, la distribución se vuelve más
simétrica y se aproxima a una normal.

\end{tcolorbox}

\subsection{\texorpdfstring{Distribución t de Student
\(t_k\)}{Distribución t de Student t\_k}}\label{distribucion-t-de-student-t_k}

\begin{tcolorbox}[enhanced jigsaw, breakable, titlerule=0mm, rightrule=.15mm, leftrule=.75mm, left=2mm, opacitybacktitle=0.6, colback=white, bottomtitle=1mm, colbacktitle=quarto-callout-note-color!10!white, title=\textcolor{quarto-callout-note-color}{\faInfo}\hspace{0.5em}{Definición: Distribución t de Student}, toptitle=1mm, coltitle=black, opacityback=0, arc=.35mm, bottomrule=.15mm, toprule=.15mm, colframe=quarto-callout-note-color-frame]

Si \(Z \sim N(0,1)\) e \(Y \sim \chi^2_k\) son independientes, entonces:

\[T = \frac{Z}{\sqrt{Y/k}} \sim t_k\]

\textbf{Parámetro:} \(k\) (grados de libertad)

\textbf{Propiedades:} - Simétrica alrededor de 0 - \(E(T) = 0\) (para
\(k > 1\)) - \(\text{Var}(T) = \frac{k}{k-2}\) (para \(k > 2\)) - Colas
más pesadas que la normal - Cuando \(k \to \infty\), converge a
\(N(0,1)\)

\textbf{En R:} \texttt{dt(x,\ df=k)}, \texttt{pt(x,\ df=k)},
\texttt{rt(n,\ df=k)}

\textbf{Usos:} Intervalos de confianza para la media (varianza
desconocida), pruebas t de Student.

\begin{codelisting}[H]

\caption{\label{lst-fig-t-student}Código R}

\centering{

\begin{Shaded}
\begin{Highlighting}[]
\NormalTok{k\_values }\OtherTok{\textless{}{-}} \FunctionTok{c}\NormalTok{(}\DecValTok{1}\NormalTok{, }\DecValTok{3}\NormalTok{, }\DecValTok{5}\NormalTok{, }\DecValTok{10}\NormalTok{)}
\NormalTok{x }\OtherTok{\textless{}{-}} \FunctionTok{seq}\NormalTok{(}\SpecialCharTok{{-}}\DecValTok{5}\NormalTok{, }\DecValTok{5}\NormalTok{, }\AttributeTok{length.out =} \DecValTok{500}\NormalTok{)}

\FunctionTok{par}\NormalTok{(}\AttributeTok{mfrow =} \FunctionTok{c}\NormalTok{(}\DecValTok{2}\NormalTok{, }\DecValTok{2}\NormalTok{))}

\ControlFlowTok{for}\NormalTok{ (k }\ControlFlowTok{in}\NormalTok{ k\_values) \{}
\NormalTok{  y }\OtherTok{\textless{}{-}} \FunctionTok{dt}\NormalTok{(x, }\AttributeTok{df =}\NormalTok{ k)}
  
  \FunctionTok{plot}\NormalTok{(}
\NormalTok{    x, y,}
    \AttributeTok{type =} \StringTok{"l"}\NormalTok{,}
    \AttributeTok{lwd =} \DecValTok{2}\NormalTok{,}
    \AttributeTok{col =} \StringTok{"steelblue"}\NormalTok{,}
    \AttributeTok{main =} \FunctionTok{paste}\NormalTok{(}\StringTok{"k ="}\NormalTok{, k),}
    \AttributeTok{xlab =} \StringTok{"x"}\NormalTok{,}
    \AttributeTok{ylab =} \StringTok{"Densidad"}
\NormalTok{  )}
  
  \FunctionTok{curve}\NormalTok{(}
    \FunctionTok{dnorm}\NormalTok{(x),}
    \AttributeTok{add =} \ConstantTok{TRUE}\NormalTok{,}
    \AttributeTok{col =} \StringTok{"red"}\NormalTok{,}
    \AttributeTok{lty =} \DecValTok{2}\NormalTok{,}
    \AttributeTok{lwd =} \DecValTok{2}
\NormalTok{  )}
  
  \FunctionTok{legend}\NormalTok{(}
    \StringTok{"topright"}\NormalTok{,}
    \AttributeTok{legend =} \FunctionTok{c}\NormalTok{(}\StringTok{"t de Student"}\NormalTok{, }\StringTok{"Normal estándar"}\NormalTok{),}
    \AttributeTok{col =} \FunctionTok{c}\NormalTok{(}\StringTok{"steelblue"}\NormalTok{, }\StringTok{"red"}\NormalTok{),}
    \AttributeTok{lty =} \FunctionTok{c}\NormalTok{(}\DecValTok{1}\NormalTok{, }\DecValTok{2}\NormalTok{),}
    \AttributeTok{bty =} \StringTok{"n"}
\NormalTok{  )}
\NormalTok{\}}
\end{Highlighting}
\end{Shaded}

}

\end{codelisting}%

\begin{figure}[H]

{\centering \pandocbounded{\includegraphics[keepaspectratio]{chapters/semana04_files/figure-pdf/unnamed-chunk-11-1.pdf}}

}

\caption{k = 1}

\end{figure}%

Densidad de la distribución t de Student para diferentes grados de
libertad. A medida que k aumenta, la distribución se vuelve más similar
a una normal estándar, aunque siempre tiene colas más pesadas.

\end{tcolorbox}

\begin{tcolorbox}[enhanced jigsaw, breakable, titlerule=0mm, rightrule=.15mm, leftrule=.75mm, left=2mm, opacitybacktitle=0.6, colback=white, bottomtitle=1mm, colbacktitle=quarto-callout-tip-color!10!white, title=\textcolor{quarto-callout-tip-color}{\faLightbulb}\hspace{0.5em}{Tip}, toptitle=1mm, coltitle=black, opacityback=0, arc=.35mm, bottomrule=.15mm, toprule=.15mm, colframe=quarto-callout-tip-color-frame]

La distribución t tiene colas más pesadas, reflejando mayor
incertidumbre con muestras pequeñas. \#\# Comparación: Normal vs t de
Student

Para \(k = 5\) grados de libertad:

\begin{codelisting}[H]

\caption{\label{lst-semana04-12}Código R}

\centering{

\begin{Shaded}
\begin{Highlighting}[]
\CommentTok{\# Graficar ambas distribuciones}
\FunctionTok{curve}\NormalTok{(}\FunctionTok{dnorm}\NormalTok{(x, }\AttributeTok{mean=}\DecValTok{0}\NormalTok{, }\AttributeTok{sd=}\DecValTok{1}\NormalTok{), }\AttributeTok{from=}\SpecialCharTok{{-}}\DecValTok{4}\NormalTok{, }\AttributeTok{to=}\DecValTok{4}\NormalTok{, }\AttributeTok{col=}\StringTok{"blue"}\NormalTok{, }\AttributeTok{lwd=}\DecValTok{2}\NormalTok{,}
      \AttributeTok{main=}\StringTok{"N(0,1) vs t con 5 gl"}\NormalTok{, }\AttributeTok{xlab=}\StringTok{"x"}\NormalTok{, }\AttributeTok{ylab=}\StringTok{"Densidad"}\NormalTok{)}
\FunctionTok{curve}\NormalTok{(}\FunctionTok{dt}\NormalTok{(x, }\AttributeTok{df=}\DecValTok{5}\NormalTok{), }\AttributeTok{from=}\SpecialCharTok{{-}}\DecValTok{4}\NormalTok{, }\AttributeTok{to=}\DecValTok{4}\NormalTok{, }\AttributeTok{col=}\StringTok{"red"}\NormalTok{, }\AttributeTok{lwd=}\DecValTok{2}\NormalTok{, }\AttributeTok{add=}\ConstantTok{TRUE}\NormalTok{)}
\FunctionTok{legend}\NormalTok{(}\StringTok{"topright"}\NormalTok{, }\AttributeTok{legend=}\FunctionTok{c}\NormalTok{(}\StringTok{"N(0,1)"}\NormalTok{, }\StringTok{"t(5)"}\NormalTok{),}
       \AttributeTok{col=}\FunctionTok{c}\NormalTok{(}\StringTok{"blue"}\NormalTok{, }\StringTok{"red"}\NormalTok{), }\AttributeTok{lwd=}\DecValTok{2}\NormalTok{)}
\end{Highlighting}
\end{Shaded}

}

\end{codelisting}%

\pandocbounded{\includegraphics[keepaspectratio]{chapters/semana04_files/figure-pdf/unnamed-chunk-12-1.pdf}}

\end{tcolorbox}

\subsection{\texorpdfstring{Distribución F de Snedecor
\(F_{k_1, k_2}\)}{Distribución F de Snedecor F\_\{k\_1, k\_2\}}}\label{distribucion-f-de-snedecor-f_k_1-k_2}

\begin{tcolorbox}[enhanced jigsaw, breakable, titlerule=0mm, rightrule=.15mm, leftrule=.75mm, left=2mm, opacitybacktitle=0.6, colback=white, bottomtitle=1mm, colbacktitle=quarto-callout-note-color!10!white, title=\textcolor{quarto-callout-note-color}{\faInfo}\hspace{0.5em}{Definición: Distribución F}, toptitle=1mm, coltitle=black, opacityback=0, arc=.35mm, bottomrule=.15mm, toprule=.15mm, colframe=quarto-callout-note-color-frame]

Si \(Y_1 \sim \chi^2_{k_1}\) e \(Y_2 \sim \chi^2_{k_2}\) son
independientes, entonces:

\[F = \frac{Y_1/k_1}{Y_2/k_2} \sim F_{k_1, k_2}\]

\textbf{Parámetros:} \(k_1, k_2\) (grados de libertad del numerador y
denominador)

\textbf{Propiedades:} - Solo toma valores positivos -
\(E(F) = \frac{k_2}{k_2-2}\) (para \(k_2 > 2\)) - Asimétrica hacia la
derecha

\textbf{En R:} \texttt{df(x,\ df1=k1,\ df2=k2)},
\texttt{pf(x,\ df1=k1,\ df2=k2)}, \texttt{rf(n,\ df1=k1,\ df2=k2)}

\textbf{Usos:} Análisis de varianza (ANOVA), pruebas de igualdad de
varianzas, análisis de regresión.

\begin{codelisting}[H]

\caption{\label{lst-fig-f-distribution}Código R}

\centering{

\begin{Shaded}
\begin{Highlighting}[]
\NormalTok{params }\OtherTok{\textless{}{-}} \FunctionTok{list}\NormalTok{(}
  \FunctionTok{c}\NormalTok{(}\DecValTok{1}\NormalTok{, }\DecValTok{1}\NormalTok{),}
  \FunctionTok{c}\NormalTok{(}\DecValTok{5}\NormalTok{, }\DecValTok{2}\NormalTok{),}
  \FunctionTok{c}\NormalTok{(}\DecValTok{10}\NormalTok{, }\DecValTok{10}\NormalTok{),}
  \FunctionTok{c}\NormalTok{(}\DecValTok{20}\NormalTok{, }\DecValTok{30}\NormalTok{)}
\NormalTok{)}

\NormalTok{x }\OtherTok{\textless{}{-}} \FunctionTok{seq}\NormalTok{(}\FloatTok{0.001}\NormalTok{, }\DecValTok{5}\NormalTok{, }\AttributeTok{length.out =} \DecValTok{500}\NormalTok{)}

\FunctionTok{par}\NormalTok{(}\AttributeTok{mfrow =} \FunctionTok{c}\NormalTok{(}\DecValTok{2}\NormalTok{, }\DecValTok{2}\NormalTok{))}

\ControlFlowTok{for}\NormalTok{ (p }\ControlFlowTok{in}\NormalTok{ params) \{}

\NormalTok{  d1 }\OtherTok{\textless{}{-}}\NormalTok{ p[}\DecValTok{1}\NormalTok{]}
\NormalTok{  d2 }\OtherTok{\textless{}{-}}\NormalTok{ p[}\DecValTok{2}\NormalTok{]}

\NormalTok{  y }\OtherTok{\textless{}{-}} \FunctionTok{df}\NormalTok{(x, }\AttributeTok{df1 =}\NormalTok{ d1, }\AttributeTok{df2 =}\NormalTok{ d2)}

  \FunctionTok{plot}\NormalTok{(}
\NormalTok{    x, y,}
    \AttributeTok{type =} \StringTok{"l"}\NormalTok{,}
    \AttributeTok{lwd =} \DecValTok{2}\NormalTok{,}
    \AttributeTok{col =} \StringTok{"blue"}\NormalTok{,}
    \AttributeTok{main =} \FunctionTok{paste}\NormalTok{(}\StringTok{"d1 ="}\NormalTok{, d1, }\StringTok{", d2 ="}\NormalTok{, d2),}
    \AttributeTok{xlab =} \StringTok{"x"}\NormalTok{,}
    \AttributeTok{ylab =} \StringTok{"Densidad"}
\NormalTok{  )}
\NormalTok{\}}
\end{Highlighting}
\end{Shaded}

}

\end{codelisting}%

\begin{figure}[H]

{\centering \pandocbounded{\includegraphics[keepaspectratio]{chapters/semana04_files/figure-pdf/unnamed-chunk-13-1.pdf}}

}

\caption{d1 = 1, d2 = 1}

\end{figure}%

Densidad de la distribución F de Snedecor para diferentes combinaciones
de grados de libertad. La distribución F es asimétrica hacia la derecha
y solo toma valores positivos.

\end{tcolorbox}

\section{Teorema Central del Límite}\label{sec-semana04-tcl}

\begin{tcolorbox}[enhanced jigsaw, breakable, titlerule=0mm, rightrule=.15mm, leftrule=.75mm, left=2mm, opacitybacktitle=0.6, colback=white, bottomtitle=1mm, colbacktitle=quarto-callout-warning-color!10!white, title=\textcolor{quarto-callout-warning-color}{\faExclamationTriangle}\hspace{0.5em}{Teorema Central del Límite}, toptitle=1mm, coltitle=black, opacityback=0, arc=.35mm, bottomrule=.15mm, toprule=.15mm, colframe=quarto-callout-warning-color-frame]

\textbf{Enunciado:} Si \(X_1, X_2, \ldots, X_n\) son variables
aleatorias independientes e idénticamente distribuidas (iid) con
\(E(X_i) = \mu\) y \(\text{Var}(X_i) = \sigma^2 < \infty\), entonces:

\[\lim_{n \to \infty} P\left(\frac{\bar{X}_n - \mu}{\sigma/\sqrt{n}} \leq z\right) = \Phi(z)\]

donde \(\Phi\) es la CDF de la distribución normal estándar.

\textbf{Consecuencias prácticas:}

\begin{enumerate}
\def\labelenumi{\arabic{enumi}.}
\item
  La media muestral es aproximadamente normal para \(n\) grande:
  \[\bar{X}_n \approx N\left(\mu, \frac{\sigma^2}{n}\right)\]
\item
  La suma de variables es aproximadamente normal:
  \[S_n = \sum_{i=1}^n X_i \approx N(n\mu, n\sigma^2)\]
\item
  Vale \textbf{sin importar la distribución original} de las \(X_i\) (si
  es discreta, continua, etc.)
\end{enumerate}

\end{tcolorbox}

\begin{tcolorbox}[enhanced jigsaw, breakable, titlerule=0mm, rightrule=.15mm, leftrule=.75mm, left=2mm, opacitybacktitle=0.6, colback=white, bottomtitle=1mm, colbacktitle=quarto-callout-tip-color!10!white, title=\textcolor{quarto-callout-tip-color}{\faLightbulb}\hspace{0.5em}{Ejemplo: Simulación del TCL}, toptitle=1mm, coltitle=black, opacityback=0, arc=.35mm, bottomrule=.15mm, toprule=.15mm, colframe=quarto-callout-tip-color-frame]

Demostración del TCL con datos uniformes:

\begin{codelisting}[H]

\caption{\label{lst-semana04-14}Código R}

\centering{

\begin{Shaded}
\begin{Highlighting}[]
\FunctionTok{set.seed}\NormalTok{(}\DecValTok{123}\NormalTok{)}

\CommentTok{\# Generar medias muestrales a partir de muestras de una U(0,1)}
\NormalTok{n\_simulaciones }\OtherTok{\textless{}{-}} \DecValTok{10000}
\NormalTok{tamano\_muestra }\OtherTok{\textless{}{-}} \DecValTok{30}

\NormalTok{medias }\OtherTok{\textless{}{-}} \FunctionTok{replicate}\NormalTok{(n\_simulaciones, }\FunctionTok{mean}\NormalTok{(}\FunctionTok{runif}\NormalTok{(tamano\_muestra, }\DecValTok{0}\NormalTok{, }\DecValTok{1}\NormalTok{)))}

\CommentTok{\# Graficar histograma con curva normal teórica}
\FunctionTok{hist}\NormalTok{(medias, }\AttributeTok{breaks=}\DecValTok{40}\NormalTok{, }\AttributeTok{probability=}\ConstantTok{TRUE}\NormalTok{,}
     \AttributeTok{main=}\StringTok{"Distribución de medias muestrales (U(0,1), n=30)"}\NormalTok{,}
     \AttributeTok{xlab=}\StringTok{"Media muestral"}\NormalTok{, }\AttributeTok{ylab=}\StringTok{"Densidad"}\NormalTok{)}

\CommentTok{\# Añadir curva normal teórica: N(0.5, 1/(12*30))}
\NormalTok{mu\_teorica }\OtherTok{\textless{}{-}} \FloatTok{0.5}
\NormalTok{sigma\_teorica }\OtherTok{\textless{}{-}} \FunctionTok{sqrt}\NormalTok{(}\DecValTok{1}\SpecialCharTok{/}\NormalTok{(}\DecValTok{12}\SpecialCharTok{*}\NormalTok{tamano\_muestra))}
\FunctionTok{curve}\NormalTok{(}\FunctionTok{dnorm}\NormalTok{(x, }\AttributeTok{mean=}\NormalTok{mu\_teorica, }\AttributeTok{sd=}\NormalTok{sigma\_teorica),}
      \AttributeTok{add=}\ConstantTok{TRUE}\NormalTok{, }\AttributeTok{col=}\StringTok{"red"}\NormalTok{, }\AttributeTok{lwd=}\DecValTok{2}\NormalTok{)}

\FunctionTok{legend}\NormalTok{(}\StringTok{"topright"}\NormalTok{, }\AttributeTok{legend=}\FunctionTok{c}\NormalTok{(}\StringTok{"Simulación"}\NormalTok{, }\StringTok{"Normal teórica"}\NormalTok{),}
       \AttributeTok{fill=}\FunctionTok{c}\NormalTok{(}\StringTok{"gray"}\NormalTok{, }\StringTok{"red"}\NormalTok{))}
\end{Highlighting}
\end{Shaded}

}

\end{codelisting}%

\pandocbounded{\includegraphics[keepaspectratio]{chapters/semana04_files/figure-pdf/unnamed-chunk-14-1.pdf}}

Aunque los datos originales son uniformes, las medias muestrales siguen
una distribución aproximadamente normal.

\end{tcolorbox}

\section{Aproximaciones entre distribuciones vistas como aplicaciones
del Teorema Central del Límite}\label{sec-semana04-aprox-tcl}

\subsection{Binomial → Normal}\label{binomial-normal}

Ya la vimos anteriormente. Sin embargo, ahora la podemos ver como una de
las aplicaciones más conocidas del Teorema Central del Límite. Cuando el
número de ensayos (\(n\)) es suficientemente grande, y la probabilidad
de éxito (\(p\)) no es extremadamente cercana a 0 o 1, la forma de la
distribución Binomial se asemeja a la de una campana, característica de
la distribución Normal.

\textbf{Condiciones para la aproximación:}\\
La aproximación es generalmente considerada adecuada si se cumplen las
siguientes condiciones:\\
- \(np \geq 5\) (o \(np \geq 10\) para una mejor aproximación)\\
- \(n(1-p) \geq 5\) (o \(n(1-p) \geq 10\) para una mejor aproximación)

Bajo estas condiciones, si \(X \sim B(n, p)\), entonces \(X\) puede ser
aproximada por \(Y \sim N(\mu = np, \sigma^2 = np(1-p))\).

\textbf{Visualización:}

Vamos a visualizar cómo la distribución Binomial se aproxima a la Normal
a medida que \(n\) aumenta. Consideremos una probabilidad de éxito
\(p=0.5\) (simétrica) y variemos \(n\).

\begin{figure}

\begin{minipage}{0.50\linewidth}

\begin{codelisting}[H]

\caption{\label{lst-fig-binom-to-normal}Código R}

\centering{

\begin{Shaded}
\begin{Highlighting}[]
\FunctionTok{par}\NormalTok{(}\AttributeTok{mfrow=}\FunctionTok{c}\NormalTok{(}\DecValTok{2}\NormalTok{, }\DecValTok{2}\NormalTok{)) }\CommentTok{\# Organizar 4 gráficos en una cuadrícula de 2x2}

\NormalTok{p\_val }\OtherTok{\textless{}{-}} \FloatTok{0.5} \CommentTok{\# Probabilidad de éxito}

\CommentTok{\# Caso 1: n pequeña (n=10)}
\NormalTok{n1 }\OtherTok{\textless{}{-}} \DecValTok{10}
\NormalTok{x1 }\OtherTok{\textless{}{-}} \DecValTok{0}\SpecialCharTok{:}\NormalTok{n1}
\NormalTok{y1 }\OtherTok{\textless{}{-}} \FunctionTok{dbinom}\NormalTok{(x1, }\AttributeTok{size=}\NormalTok{n1, }\AttributeTok{prob=}\NormalTok{p\_val)}
\NormalTok{mu1 }\OtherTok{\textless{}{-}}\NormalTok{ n1 }\SpecialCharTok{*}\NormalTok{ p\_val}
\NormalTok{sigma1 }\OtherTok{\textless{}{-}} \FunctionTok{sqrt}\NormalTok{(n1 }\SpecialCharTok{*}\NormalTok{ p\_val }\SpecialCharTok{*}\NormalTok{ (}\DecValTok{1} \SpecialCharTok{{-}}\NormalTok{ p\_val))}
\FunctionTok{plot}\NormalTok{(x1, y1, }\AttributeTok{type=}\StringTok{"h"}\NormalTok{, }\AttributeTok{lwd=}\DecValTok{2}\NormalTok{, }\AttributeTok{col=}\StringTok{"blue"}\NormalTok{,}
     \AttributeTok{main=}\FunctionTok{paste0}\NormalTok{(}\StringTok{"Binomial (n="}\NormalTok{, n1, }\StringTok{", p="}\NormalTok{, p\_val, }\StringTok{")"}\NormalTok{),}
     \AttributeTok{xlab=}\StringTok{"Número de Éxitos"}\NormalTok{, }\AttributeTok{ylab=}\StringTok{"Probabilidad / Densidad"}\NormalTok{)}
\FunctionTok{curve}\NormalTok{(}\FunctionTok{dnorm}\NormalTok{(x, }\AttributeTok{mean=}\NormalTok{mu1, }\AttributeTok{sd=}\NormalTok{sigma1), }\AttributeTok{add=}\ConstantTok{TRUE}\NormalTok{, }\AttributeTok{col=}\StringTok{"red"}\NormalTok{, }\AttributeTok{lwd=}\DecValTok{2}\NormalTok{)}

\CommentTok{\# Caso 2: n mediana (n=30)}
\NormalTok{n2 }\OtherTok{\textless{}{-}} \DecValTok{30}
\NormalTok{x2 }\OtherTok{\textless{}{-}} \DecValTok{0}\SpecialCharTok{:}\NormalTok{n2}
\NormalTok{y2 }\OtherTok{\textless{}{-}} \FunctionTok{dbinom}\NormalTok{(x2, }\AttributeTok{size=}\NormalTok{n2, }\AttributeTok{prob=}\NormalTok{p\_val)}
\NormalTok{mu2 }\OtherTok{\textless{}{-}}\NormalTok{ n2 }\SpecialCharTok{*}\NormalTok{ p\_val}
\NormalTok{sigma2 }\OtherTok{\textless{}{-}} \FunctionTok{sqrt}\NormalTok{(n2 }\SpecialCharTok{*}\NormalTok{ p\_val }\SpecialCharTok{*}\NormalTok{ (}\DecValTok{1} \SpecialCharTok{{-}}\NormalTok{ p\_val))}
\FunctionTok{plot}\NormalTok{(x2, y2, }\AttributeTok{type=}\StringTok{"h"}\NormalTok{, }\AttributeTok{lwd=}\DecValTok{2}\NormalTok{, }\AttributeTok{col=}\StringTok{"blue"}\NormalTok{,}
     \AttributeTok{main=}\FunctionTok{paste0}\NormalTok{(}\StringTok{"Binomial (n="}\NormalTok{, n2, }\StringTok{", p="}\NormalTok{, p\_val, }\StringTok{")"}\NormalTok{),}
     \AttributeTok{xlab=}\StringTok{"Número de Éxitos"}\NormalTok{, }\AttributeTok{ylab=}\StringTok{"Probabilidad / Densidad"}\NormalTok{)}
\FunctionTok{curve}\NormalTok{(}\FunctionTok{dnorm}\NormalTok{(x, }\AttributeTok{mean=}\NormalTok{mu2, }\AttributeTok{sd=}\NormalTok{sigma2), }\AttributeTok{add=}\ConstantTok{TRUE}\NormalTok{, }\AttributeTok{col=}\StringTok{"red"}\NormalTok{, }\AttributeTok{lwd=}\DecValTok{2}\NormalTok{)}

\CommentTok{\# Caso 3: n grande (n=100)}
\NormalTok{n3 }\OtherTok{\textless{}{-}} \DecValTok{100}
\NormalTok{x3 }\OtherTok{\textless{}{-}} \DecValTok{0}\SpecialCharTok{:}\NormalTok{n3}
\NormalTok{y3 }\OtherTok{\textless{}{-}} \FunctionTok{dbinom}\NormalTok{(x3, }\AttributeTok{size=}\NormalTok{n3, }\AttributeTok{prob=}\NormalTok{p\_val)}
\NormalTok{mu3 }\OtherTok{\textless{}{-}}\NormalTok{ n3 }\SpecialCharTok{*}\NormalTok{ p\_val}
\NormalTok{sigma3 }\OtherTok{\textless{}{-}} \FunctionTok{sqrt}\NormalTok{(n3 }\SpecialCharTok{*}\NormalTok{ p\_val }\SpecialCharTok{*}\NormalTok{ (}\DecValTok{1} \SpecialCharTok{{-}}\NormalTok{ p\_val))}
\FunctionTok{plot}\NormalTok{(x3, y3, }\AttributeTok{type=}\StringTok{"h"}\NormalTok{, }\AttributeTok{lwd=}\DecValTok{2}\NormalTok{, }\AttributeTok{col=}\StringTok{"blue"}\NormalTok{,}
     \AttributeTok{main=}\FunctionTok{paste0}\NormalTok{(}\StringTok{"Binomial (n="}\NormalTok{, n3, }\StringTok{", p="}\NormalTok{, p\_val, }\StringTok{")"}\NormalTok{),}
     \AttributeTok{xlab=}\StringTok{"Número de Éxitos"}\NormalTok{, }\AttributeTok{ylab=}\StringTok{"Probabilidad / Densidad"}\NormalTok{)}
\FunctionTok{curve}\NormalTok{(}\FunctionTok{dnorm}\NormalTok{(x, }\AttributeTok{mean=}\NormalTok{mu3, }\AttributeTok{sd=}\NormalTok{sigma3), }\AttributeTok{add=}\ConstantTok{TRUE}\NormalTok{, }\AttributeTok{col=}\StringTok{"red"}\NormalTok{, }\AttributeTok{lwd=}\DecValTok{2}\NormalTok{)}

\CommentTok{\# Caso 4: n muy grande (n=500)}
\NormalTok{n4 }\OtherTok{\textless{}{-}} \DecValTok{500}
\NormalTok{x4 }\OtherTok{\textless{}{-}} \DecValTok{0}\SpecialCharTok{:}\NormalTok{n4}
\NormalTok{y4 }\OtherTok{\textless{}{-}} \FunctionTok{dbinom}\NormalTok{(x4, }\AttributeTok{size=}\NormalTok{n4, }\AttributeTok{prob=}\NormalTok{p\_val)}
\NormalTok{mu4 }\OtherTok{\textless{}{-}}\NormalTok{ n4 }\SpecialCharTok{*}\NormalTok{ p\_val}
\NormalTok{sigma4 }\OtherTok{\textless{}{-}} \FunctionTok{sqrt}\NormalTok{(n4 }\SpecialCharTok{*}\NormalTok{ p\_val }\SpecialCharTok{*}\NormalTok{ (}\DecValTok{1} \SpecialCharTok{{-}}\NormalTok{ p\_val))}
\FunctionTok{plot}\NormalTok{(x4, y4, }\AttributeTok{type=}\StringTok{"h"}\NormalTok{, }\AttributeTok{lwd=}\DecValTok{2}\NormalTok{, }\AttributeTok{col=}\StringTok{"blue"}\NormalTok{,}
     \AttributeTok{main=}\FunctionTok{paste0}\NormalTok{(}\StringTok{"Binomial (n="}\NormalTok{, n4, }\StringTok{", p="}\NormalTok{, p\_val, }\StringTok{")"}\NormalTok{),}
     \AttributeTok{xlab=}\StringTok{"Número de Éxitos"}\NormalTok{, }\AttributeTok{ylab=}\StringTok{"Probabilidad / Densidad"}\NormalTok{)}
\FunctionTok{curve}\NormalTok{(}\FunctionTok{dnorm}\NormalTok{(x, }\AttributeTok{mean=}\NormalTok{mu4, }\AttributeTok{sd=}\NormalTok{sigma4), }\AttributeTok{add=}\ConstantTok{TRUE}\NormalTok{, }\AttributeTok{col=}\StringTok{"red"}\NormalTok{, }\AttributeTok{lwd=}\DecValTok{2}\NormalTok{)}
\end{Highlighting}
\end{Shaded}

}

\end{codelisting}%

\end{minipage}%
%
\begin{minipage}{0.50\linewidth}

\begin{figure}[H]

{\centering \pandocbounded{\includegraphics[keepaspectratio]{chapters/semana04_files/figure-pdf/unnamed-chunk-15-1.pdf}}

}

\subcaption{Aproximación de la distribución Binomial a la Normal para
diferentes valores de n.~A medida que n aumenta, la forma de la binomial
(barras) se asemeja más a la curva de densidad Normal (línea roja).}

\end{figure}%

\end{minipage}%
\newline
\begin{minipage}{0.50\linewidth}

\begin{codelisting}[H]

\caption{\label{lst-fig-binom-to-normal}Código R}

\centering{

\begin{Shaded}
\begin{Highlighting}[]
\FunctionTok{par}\NormalTok{(}\AttributeTok{mfrow=}\FunctionTok{c}\NormalTok{(}\DecValTok{1}\NormalTok{,}\DecValTok{1}\NormalTok{)) }\CommentTok{\# Reset layout}
\end{Highlighting}
\end{Shaded}

}

\end{codelisting}%

\end{minipage}%

\end{figure}%

\subsection{Poisson → Normal}\label{poisson-normal}

Similar a la Binomial, la distribución de Poisson también puede ser
aproximada por la distribución Normal cuando su parámetro de tasa
(\(\lambda\)) es suficientemente grande. Esto se debe a que la
distribución de Poisson puede verse como el límite de una Binomial
cuando \(n \to \infty\) y \(p \to 0\) tal que \(np = \lambda\). Por lo
tanto, el Teorema Central del Límite también juega un papel aquí.

\textbf{Condiciones para la aproximación:} La aproximación es
generalmente aceptable si: * \(\lambda \geq 10\) (o \(\lambda \geq 5\)
en algunos contextos, pero \(\lambda \geq 10\) es más seguro para una
buena aproximación).

Bajo esta condición, si \(X \sim \text{Po}(\lambda)\), entonces \(X\)
puede ser aproximada por
\(Y \sim N(\mu = \lambda, \sigma^2 = \lambda)\).

\textbf{Ejemplo y Visualización:}

Vamos a observar cómo la distribución de Poisson se asemeja a la Normal
a medida que \(\lambda\) aumenta.

\begin{figure}

\begin{minipage}{0.50\linewidth}

\begin{codelisting}[H]

\caption{\label{lst-fig-poisson-to-normal}Código R}

\centering{

\begin{Shaded}
\begin{Highlighting}[]
\FunctionTok{par}\NormalTok{(}\AttributeTok{mfrow=}\FunctionTok{c}\NormalTok{(}\DecValTok{2}\NormalTok{, }\DecValTok{2}\NormalTok{)) }\CommentTok{\# Organizar 4 gráficos en una cuadrícula de 2x2}

\CommentTok{\# Caso 1: lambda pequeña (lambda=1)}
\NormalTok{lambda1 }\OtherTok{\textless{}{-}} \DecValTok{1}
\NormalTok{x1\_poisson }\OtherTok{\textless{}{-}} \DecValTok{0}\SpecialCharTok{:}\FunctionTok{max}\NormalTok{(}\DecValTok{10}\NormalTok{, }\FunctionTok{round}\NormalTok{(lambda1 }\SpecialCharTok{+} \DecValTok{3} \SpecialCharTok{*} \FunctionTok{sqrt}\NormalTok{(lambda1)))}
\NormalTok{y1\_poisson }\OtherTok{\textless{}{-}} \FunctionTok{dpois}\NormalTok{(x1\_poisson, }\AttributeTok{lambda=}\NormalTok{lambda1)}
\NormalTok{mu1\_poisson }\OtherTok{\textless{}{-}}\NormalTok{ lambda1}
\NormalTok{sigma1\_poisson }\OtherTok{\textless{}{-}} \FunctionTok{sqrt}\NormalTok{(lambda1)}
\FunctionTok{plot}\NormalTok{(x1\_poisson, y1\_poisson, }\AttributeTok{type=}\StringTok{"h"}\NormalTok{, }\AttributeTok{lwd=}\DecValTok{2}\NormalTok{, }\AttributeTok{col=}\StringTok{"darkgreen"}\NormalTok{,}
     \AttributeTok{main=}\FunctionTok{paste0}\NormalTok{(}\StringTok{"Poisson (lambda="}\NormalTok{, lambda1, }\StringTok{")"}\NormalTok{),}
     \AttributeTok{xlab=}\StringTok{"Número de Eventos"}\NormalTok{, }\AttributeTok{ylab=}\StringTok{"Probabilidad / Densidad"}\NormalTok{,}
     \AttributeTok{ylim=}\FunctionTok{c}\NormalTok{(}\DecValTok{0}\NormalTok{, }\FunctionTok{max}\NormalTok{(y1\_poisson, }\FunctionTok{dnorm}\NormalTok{(mu1\_poisson, mu1\_poisson, sigma1\_poisson))))}
\FunctionTok{curve}\NormalTok{(}\FunctionTok{dnorm}\NormalTok{(x, }\AttributeTok{mean=}\NormalTok{mu1\_poisson, }\AttributeTok{sd=}\NormalTok{sigma1\_poisson),}
      \AttributeTok{from=}\FunctionTok{min}\NormalTok{(x1\_poisson), }\AttributeTok{to=}\FunctionTok{max}\NormalTok{(x1\_poisson), }\AttributeTok{add=}\ConstantTok{TRUE}\NormalTok{, }\AttributeTok{col=}\StringTok{"red"}\NormalTok{, }\AttributeTok{lwd=}\DecValTok{2}\NormalTok{)}


\CommentTok{\# Caso 2: lambda mediana (lambda=5)}
\NormalTok{lambda2 }\OtherTok{\textless{}{-}} \DecValTok{5}
\NormalTok{x2\_poisson }\OtherTok{\textless{}{-}} \DecValTok{0}\SpecialCharTok{:}\FunctionTok{max}\NormalTok{(}\DecValTok{15}\NormalTok{, }\FunctionTok{round}\NormalTok{(lambda2 }\SpecialCharTok{+} \DecValTok{3} \SpecialCharTok{*} \FunctionTok{sqrt}\NormalTok{(lambda2)))}
\NormalTok{y2\_poisson }\OtherTok{\textless{}{-}} \FunctionTok{dpois}\NormalTok{(x2\_poisson, }\AttributeTok{lambda=}\NormalTok{lambda2)}
\NormalTok{mu2\_poisson }\OtherTok{\textless{}{-}}\NormalTok{ lambda2}
\NormalTok{sigma2\_poisson }\OtherTok{\textless{}{-}} \FunctionTok{sqrt}\NormalTok{(lambda2)}
\FunctionTok{plot}\NormalTok{(x2\_poisson, y2\_poisson, }\AttributeTok{type=}\StringTok{"h"}\NormalTok{, }\AttributeTok{lwd=}\DecValTok{2}\NormalTok{, }\AttributeTok{col=}\StringTok{"darkgreen"}\NormalTok{,}
     \AttributeTok{main=}\FunctionTok{paste0}\NormalTok{(}\StringTok{"Poisson (lambda="}\NormalTok{, lambda2, }\StringTok{")"}\NormalTok{),}
     \AttributeTok{xlab=}\StringTok{"Número de Eventos"}\NormalTok{, }\AttributeTok{ylab=}\StringTok{"Probabilidad / Densidad"}\NormalTok{,}
     \AttributeTok{ylim=}\FunctionTok{c}\NormalTok{(}\DecValTok{0}\NormalTok{, }\FunctionTok{max}\NormalTok{(y2\_poisson, }\FunctionTok{dnorm}\NormalTok{(mu2\_poisson, mu2\_poisson, sigma2\_poisson))))}
\FunctionTok{curve}\NormalTok{(}\FunctionTok{dnorm}\NormalTok{(x, }\AttributeTok{mean=}\NormalTok{mu2\_poisson, }\AttributeTok{sd=}\NormalTok{sigma2\_poisson),}
      \AttributeTok{from=}\FunctionTok{min}\NormalTok{(x2\_poisson), }\AttributeTok{to=}\FunctionTok{max}\NormalTok{(x2\_poisson), }\AttributeTok{add=}\ConstantTok{TRUE}\NormalTok{, }\AttributeTok{col=}\StringTok{"red"}\NormalTok{, }\AttributeTok{lwd=}\DecValTok{2}\NormalTok{)}

\CommentTok{\# Caso 3: lambda grande (lambda=15)}
\NormalTok{lambda3 }\OtherTok{\textless{}{-}} \DecValTok{15}
\NormalTok{x3\_poisson }\OtherTok{\textless{}{-}} \DecValTok{0}\SpecialCharTok{:}\FunctionTok{max}\NormalTok{(}\DecValTok{30}\NormalTok{, }\FunctionTok{round}\NormalTok{(lambda3 }\SpecialCharTok{+} \DecValTok{3} \SpecialCharTok{*} \FunctionTok{sqrt}\NormalTok{(lambda3)))}
\NormalTok{y3\_poisson }\OtherTok{\textless{}{-}} \FunctionTok{dpois}\NormalTok{(x3\_poisson, }\AttributeTok{lambda=}\NormalTok{lambda3)}
\NormalTok{mu3\_poisson }\OtherTok{\textless{}{-}}\NormalTok{ lambda3}
\NormalTok{sigma3\_poisson }\OtherTok{\textless{}{-}} \FunctionTok{sqrt}\NormalTok{(lambda3)}
\FunctionTok{plot}\NormalTok{(x3\_poisson, y3\_poisson, }\AttributeTok{type=}\StringTok{"h"}\NormalTok{, }\AttributeTok{lwd=}\DecValTok{2}\NormalTok{, }\AttributeTok{col=}\StringTok{"darkgreen"}\NormalTok{,}
     \AttributeTok{main=}\FunctionTok{paste0}\NormalTok{(}\StringTok{"Poisson (lambda="}\NormalTok{, lambda3, }\StringTok{")"}\NormalTok{),}
     \AttributeTok{xlab=}\StringTok{"Número de Eventos"}\NormalTok{, }\AttributeTok{ylab=}\StringTok{"Probabilidad / Densidad"}\NormalTok{,}
     \AttributeTok{ylim=}\FunctionTok{c}\NormalTok{(}\DecValTok{0}\NormalTok{, }\FunctionTok{max}\NormalTok{(y3\_poisson, }\FunctionTok{dnorm}\NormalTok{(mu3\_poisson, mu3\_poisson, sigma3\_poisson))))}
\FunctionTok{curve}\NormalTok{(}\FunctionTok{dnorm}\NormalTok{(x, }\AttributeTok{mean=}\NormalTok{mu3\_poisson, }\AttributeTok{sd=}\NormalTok{sigma3\_poisson),}
      \AttributeTok{from=}\FunctionTok{min}\NormalTok{(x3\_poisson), }\AttributeTok{to=}\FunctionTok{max}\NormalTok{(x3\_poisson), }\AttributeTok{add=}\ConstantTok{TRUE}\NormalTok{, }\AttributeTok{col=}\StringTok{"red"}\NormalTok{, }\AttributeTok{lwd=}\DecValTok{2}\NormalTok{)}

\CommentTok{\# Caso 4: lambda muy grande (lambda=50)}
\NormalTok{lambda4 }\OtherTok{\textless{}{-}} \DecValTok{50}
\NormalTok{x4\_poisson }\OtherTok{\textless{}{-}} \DecValTok{0}\SpecialCharTok{:}\FunctionTok{max}\NormalTok{(}\DecValTok{70}\NormalTok{, }\FunctionTok{round}\NormalTok{(lambda4 }\SpecialCharTok{+} \DecValTok{3} \SpecialCharTok{*} \FunctionTok{sqrt}\NormalTok{(lambda4)))}
\NormalTok{y4\_poisson }\OtherTok{\textless{}{-}} \FunctionTok{dpois}\NormalTok{(x4\_poisson, }\AttributeTok{lambda=}\NormalTok{lambda4)}
\NormalTok{mu4\_poisson }\OtherTok{\textless{}{-}}\NormalTok{ lambda4}
\NormalTok{sigma4\_poisson }\OtherTok{\textless{}{-}} \FunctionTok{sqrt}\NormalTok{(lambda4)}
\FunctionTok{plot}\NormalTok{(x4\_poisson, y4\_poisson, }\AttributeTok{type=}\StringTok{"h"}\NormalTok{, }\AttributeTok{lwd=}\DecValTok{2}\NormalTok{, }\AttributeTok{col=}\StringTok{"darkgreen"}\NormalTok{,}
     \AttributeTok{main=}\FunctionTok{paste0}\NormalTok{(}\StringTok{"Poisson (lambda="}\NormalTok{, lambda4, }\StringTok{")"}\NormalTok{),}
     \AttributeTok{xlab=}\StringTok{"Número de Eventos"}\NormalTok{, }\AttributeTok{ylab=}\StringTok{"Probabilidad / Densidad"}\NormalTok{,}
     \AttributeTok{ylim=}\FunctionTok{c}\NormalTok{(}\DecValTok{0}\NormalTok{, }\FunctionTok{max}\NormalTok{(y4\_poisson, }\FunctionTok{dnorm}\NormalTok{(mu4\_poisson, mu4\_poisson, sigma4\_poisson))))}
\FunctionTok{curve}\NormalTok{(}\FunctionTok{dnorm}\NormalTok{(x, }\AttributeTok{mean=}\NormalTok{mu4\_poisson, }\AttributeTok{sd=}\NormalTok{sigma4\_poisson),}
      \AttributeTok{from=}\FunctionTok{min}\NormalTok{(x4\_poisson), }\AttributeTok{to=}\FunctionTok{max}\NormalTok{(x4\_poisson), }\AttributeTok{add=}\ConstantTok{TRUE}\NormalTok{, }\AttributeTok{col=}\StringTok{"red"}\NormalTok{, }\AttributeTok{lwd=}\DecValTok{2}\NormalTok{)}
\end{Highlighting}
\end{Shaded}

}

\end{codelisting}%

\end{minipage}%
%
\begin{minipage}{0.50\linewidth}

\begin{figure}[H]

{\centering \pandocbounded{\includegraphics[keepaspectratio]{chapters/semana04_files/figure-pdf/unnamed-chunk-16-1.pdf}}

}

\subcaption{Aproximación de la distribución de Poisson a la Normal para
diferentes valores de lambda. A medida que lambda aumenta, la forma de
la Poisson (barras) se asemeja más a la curva de densidad Normal (línea
roja).}

\end{figure}%

\end{minipage}%
\newline
\begin{minipage}{0.50\linewidth}

\begin{codelisting}[H]

\caption{\label{lst-fig-poisson-to-normal}Código R}

\centering{

\begin{Shaded}
\begin{Highlighting}[]
\FunctionTok{par}\NormalTok{(}\AttributeTok{mfrow=}\FunctionTok{c}\NormalTok{(}\DecValTok{1}\NormalTok{,}\DecValTok{1}\NormalTok{)) }\CommentTok{\# Reset layout}
\end{Highlighting}
\end{Shaded}

}

\end{codelisting}%

\end{minipage}%

\end{figure}%

\subsection{Exponencial ↔ Poisson (dualidad
tiempo-conteo)}\label{exponencial-poisson-dualidad-tiempo-conteo}

La relación entre las distribuciones Exponencial y de Poisson es una de
las dualidades más fundamentales en la teoría de procesos estocásticos,
particularmente en los Procesos de Poisson. Un Proceso de Poisson
describe el número de eventos que ocurren en un intervalo de tiempo fijo
(con distribución de Poisson), mientras que la distribución Exponencial
describe el tiempo entre eventos consecutivos (conocido como ``tiempo de
inter-llegada'') en un proceso de Poisson.

\textbf{Conceptualización de la dualidad:} - Si el número de eventos en
un intervalo de tiempo \(t\) sigue una distribución de Poisson con tasa
\(\lambda\), es decir, \(N(t) \sim \text{Po}(\lambda t)\), entonces:\\
- El tiempo hasta que ocurre el primer evento (o el tiempo entre dos
eventos consecutivos) sigue una distribución Exponencial con parámetro
de tasa \(\lambda\).

Esta dualidad es crucial para modelar fenómenos como el tiempo de espera
hasta el próximo cliente en una cola, el tiempo de vida de componentes
electrónicos, o el tiempo entre mutaciones genéticas.

\textbf{Ejemplo y Visualización:}

Vamos a simular eventos de un Proceso de Poisson y luego examinar la
distribución de los tiempos entre estos eventos para ver cómo se ajusta
a una distribución Exponencial.

\begin{codelisting}

\caption{\label{lst-fig-exp-poisson-duality}Código R}

\centering{

\begin{Shaded}
\begin{Highlighting}[]
\FunctionTok{set.seed}\NormalTok{(}\DecValTok{42}\NormalTok{) }\CommentTok{\# Para reproducibilidad}

\NormalTok{lambda\_rate }\OtherTok{\textless{}{-}} \DecValTok{2} \CommentTok{\# Tasa de eventos por unidad de tiempo}
\NormalTok{max\_time }\OtherTok{\textless{}{-}} \DecValTok{10}  \CommentTok{\# Tiempo total de observación}

\CommentTok{\# Simular un proceso de Poisson: número de eventos en max\_time}
\NormalTok{num\_events\_poisson }\OtherTok{\textless{}{-}} \FunctionTok{rpois}\NormalTok{(}\DecValTok{1}\NormalTok{, }\AttributeTok{lambda =}\NormalTok{ lambda\_rate }\SpecialCharTok{*}\NormalTok{ max\_time)}
\FunctionTok{cat}\NormalTok{(}\StringTok{"Número simulado de eventos en"}\NormalTok{, max\_time, }\StringTok{"unidades de tiempo:"}\NormalTok{, num\_events\_poisson, }\StringTok{"}\SpecialCharTok{\textbackslash{}n}\StringTok{"}\NormalTok{)}
\end{Highlighting}
\end{Shaded}

}

\end{codelisting}%

\begin{verbatim}
Número simulado de eventos en 10 unidades de tiempo: 26 
\end{verbatim}

\begin{codelisting}

\caption{\label{lst-fig-exp-poisson-duality}Código R}

\centering{

\begin{Shaded}
\begin{Highlighting}[]
\CommentTok{\# Simular los tiempos de inter{-}llegada (distribución Exponencial)}
\NormalTok{inter\_arrival\_times }\OtherTok{\textless{}{-}} \FunctionTok{rexp}\NormalTok{(num\_events\_poisson }\SpecialCharTok{+} \DecValTok{1}\NormalTok{, }\AttributeTok{rate =}\NormalTok{ lambda\_rate)}
\CommentTok{\# Asegurarse de que la suma de tiempos no exceda max\_time}
\NormalTok{arrival\_times }\OtherTok{\textless{}{-}} \FunctionTok{cumsum}\NormalTok{(inter\_arrival\_times)}
\NormalTok{arrival\_times }\OtherTok{\textless{}{-}}\NormalTok{ arrival\_times[arrival\_times }\SpecialCharTok{\textless{}=}\NormalTok{ max\_time]}

\CommentTok{\# Eliminar el último tiempo si excede el max\_time para el histograma}
\ControlFlowTok{if}\NormalTok{ (}\FunctionTok{length}\NormalTok{(arrival\_times) }\SpecialCharTok{\textgreater{}} \DecValTok{0} \SpecialCharTok{\&\&} \FunctionTok{sum}\NormalTok{(inter\_arrival\_times[}\DecValTok{1}\SpecialCharTok{:}\FunctionTok{length}\NormalTok{(arrival\_times)]) }\SpecialCharTok{\textgreater{}}\NormalTok{ max\_time) \{}
\NormalTok{  inter\_arrival\_times\_for\_hist }\OtherTok{\textless{}{-}}\NormalTok{ inter\_arrival\_times[}\DecValTok{1}\SpecialCharTok{:}\FunctionTok{length}\NormalTok{(arrival\_times) }\SpecialCharTok{{-}} \DecValTok{1}\NormalTok{]}
\NormalTok{\} }\ControlFlowTok{else}\NormalTok{ \{}
\NormalTok{  inter\_arrival\_times\_for\_hist }\OtherTok{\textless{}{-}}\NormalTok{ inter\_arrival\_times[}\DecValTok{1}\SpecialCharTok{:}\FunctionTok{length}\NormalTok{(arrival\_times)]}
\NormalTok{\}}
\NormalTok{inter\_arrival\_times\_for\_hist }\OtherTok{\textless{}{-}}\NormalTok{ inter\_arrival\_times\_for\_hist[inter\_arrival\_times\_for\_hist }\SpecialCharTok{\textgreater{}} \DecValTok{0}\NormalTok{] }\CommentTok{\# Eliminar posibles 0 si no hay eventos}

\CommentTok{\# Preparar el gráfico}
\FunctionTok{par}\NormalTok{(}\AttributeTok{mfrow=}\FunctionTok{c}\NormalTok{(}\DecValTok{2}\NormalTok{, }\DecValTok{1}\NormalTok{), }\AttributeTok{mar=}\FunctionTok{c}\NormalTok{(}\DecValTok{4}\NormalTok{, }\DecValTok{4}\NormalTok{, }\DecValTok{2}\NormalTok{, }\DecValTok{2}\NormalTok{) }\SpecialCharTok{+} \FloatTok{0.1}\NormalTok{)}

\CommentTok{\# Gráfico superior: Proceso de Poisson (eventos en el tiempo)}
\FunctionTok{plot}\NormalTok{(arrival\_times, }\FunctionTok{rep}\NormalTok{(}\DecValTok{1}\NormalTok{, }\FunctionTok{length}\NormalTok{(arrival\_times)), }\AttributeTok{pch=}\DecValTok{16}\NormalTok{, }\AttributeTok{col=}\StringTok{"darkblue"}\NormalTok{,}
     \AttributeTok{xlab=}\StringTok{"Tiempo"}\NormalTok{, }\AttributeTok{ylab=}\StringTok{"Eventos"}\NormalTok{,}
     \AttributeTok{main=}\FunctionTok{paste0}\NormalTok{(}\StringTok{"Simulación de Proceso de Poisson (lambda="}\NormalTok{, lambda\_rate, }\StringTok{")"}\NormalTok{),}
     \AttributeTok{xlim=}\FunctionTok{c}\NormalTok{(}\DecValTok{0}\NormalTok{, max\_time), }\AttributeTok{ylim=}\FunctionTok{c}\NormalTok{(}\FloatTok{0.5}\NormalTok{, }\FloatTok{1.5}\NormalTok{), }\AttributeTok{yaxt=}\StringTok{\textquotesingle{}n\textquotesingle{}}\NormalTok{)}
\FunctionTok{axis}\NormalTok{(}\AttributeTok{side=}\DecValTok{2}\NormalTok{, }\AttributeTok{at=}\DecValTok{1}\NormalTok{, }\AttributeTok{labels=}\StringTok{"Eventos"}\NormalTok{)}
\FunctionTok{segments}\NormalTok{(}\DecValTok{0}\NormalTok{, }\DecValTok{1}\NormalTok{, max\_time, }\DecValTok{1}\NormalTok{, }\AttributeTok{col=}\StringTok{"gray"}\NormalTok{, }\AttributeTok{lty=}\DecValTok{2}\NormalTok{)}

\CommentTok{\# Gráfico inferior: Histograma de tiempos de inter{-}llegada (Exponencial)}
\ControlFlowTok{if}\NormalTok{ (}\FunctionTok{length}\NormalTok{(inter\_arrival\_times\_for\_hist) }\SpecialCharTok{\textgreater{}} \DecValTok{1}\NormalTok{) \{}
  \FunctionTok{hist}\NormalTok{(inter\_arrival\_times\_for\_hist, }\AttributeTok{breaks=}\DecValTok{20}\NormalTok{, }\AttributeTok{probability=}\ConstantTok{TRUE}\NormalTok{,}
       \AttributeTok{main=}\FunctionTok{paste0}\NormalTok{(}\StringTok{"Histograma de Tiempos de Inter{-}Llegada (Exponencial, lambda="}\NormalTok{, lambda\_rate, }\StringTok{")"}\NormalTok{),}
       \AttributeTok{xlab=}\StringTok{"Tiempo entre Eventos"}\NormalTok{, }\AttributeTok{ylab=}\StringTok{"Densidad"}\NormalTok{,}
       \AttributeTok{col=}\StringTok{"lightgreen"}\NormalTok{, }\AttributeTok{border=}\StringTok{"white"}\NormalTok{)}
  \FunctionTok{curve}\NormalTok{(}\FunctionTok{dexp}\NormalTok{(x, }\AttributeTok{rate =}\NormalTok{ lambda\_rate), }\AttributeTok{add=}\ConstantTok{TRUE}\NormalTok{, }\AttributeTok{col=}\StringTok{"red"}\NormalTok{, }\AttributeTok{lwd=}\DecValTok{2}\NormalTok{)}
\NormalTok{\} }\ControlFlowTok{else}\NormalTok{ \{}
  \FunctionTok{plot}\NormalTok{(}\DecValTok{1}\NormalTok{, }\AttributeTok{type=}\StringTok{"n"}\NormalTok{, }\AttributeTok{main=}\StringTok{"No hay suficientes eventos para histograma de inter{-}llegada"}\NormalTok{,}
       \AttributeTok{xlab=}\StringTok{""}\NormalTok{, }\AttributeTok{ylab=}\StringTok{""}\NormalTok{, }\AttributeTok{xaxt=}\StringTok{\textquotesingle{}n\textquotesingle{}}\NormalTok{, }\AttributeTok{yaxt=}\StringTok{\textquotesingle{}n\textquotesingle{}}\NormalTok{)}
  \FunctionTok{text}\NormalTok{(}\DecValTok{1}\NormalTok{,}\DecValTok{1}\NormalTok{, }\StringTok{"Aumenta max\_time o lambda\_rate para más eventos."}\NormalTok{)}
\NormalTok{\}}
\end{Highlighting}
\end{Shaded}

}

\end{codelisting}%

\begin{figure}[H]

{\centering \pandocbounded{\includegraphics[keepaspectratio]{chapters/semana04_files/figure-pdf/unnamed-chunk-17-1.pdf}}

}

\caption{Dualidad entre la distribución de Poisson y la Exponencial. La
parte superior muestra la simulación de un proceso de Poisson con
eventos marcados. La parte inferior es el histograma de los tiempos
entre eventos, mostrando un ajuste a la distribución Exponencial.}

\end{figure}%

\begin{codelisting}

\caption{\label{lst-fig-exp-poisson-duality}Código R}

\centering{

\begin{Shaded}
\begin{Highlighting}[]
\FunctionTok{par}\NormalTok{(}\AttributeTok{mfrow=}\FunctionTok{c}\NormalTok{(}\DecValTok{1}\NormalTok{,}\DecValTok{1}\NormalTok{)) }\CommentTok{\# Reset layout}
\end{Highlighting}
\end{Shaded}

}

\end{codelisting}%

\section{Tabla de Referencia de
Distribuciones}\label{sec-semana04-tabla}

\begin{longtable}[]{@{}
  >{\raggedright\arraybackslash}p{(\linewidth - 8\tabcolsep) * \real{0.2000}}
  >{\raggedright\arraybackslash}p{(\linewidth - 8\tabcolsep) * \real{0.2000}}
  >{\raggedright\arraybackslash}p{(\linewidth - 8\tabcolsep) * \real{0.2000}}
  >{\raggedright\arraybackslash}p{(\linewidth - 8\tabcolsep) * \real{0.2000}}
  >{\raggedright\arraybackslash}p{(\linewidth - 8\tabcolsep) * \real{0.2000}}@{}}
\toprule\noalign{}
\begin{minipage}[b]{\linewidth}\raggedright
Distribución
\end{minipage} & \begin{minipage}[b]{\linewidth}\raggedright
Parámetros
\end{minipage} & \begin{minipage}[b]{\linewidth}\raggedright
\(E(X)\)
\end{minipage} & \begin{minipage}[b]{\linewidth}\raggedright
\(\text{Var}(X)\)
\end{minipage} & \begin{minipage}[b]{\linewidth}\raggedright
Función R
\end{minipage} \\
\midrule\noalign{}
\endhead
\bottomrule\noalign{}
\endlastfoot
\textbf{Discretas} & & & & \\
Uniforme & \(n\) & \(\frac{1}{n}\sum x_i\) &
\(\frac{1}{n}\sum(x_i-\mu)^2\) & \texttt{sample()} \\
Bernoulli & \(p\) & \(p\) & \(p(1-p)\) & \texttt{dbinom(,size=1)} \\
Binomial & \(n, p\) & \(np\) & \(np(1-p)\) & \texttt{dbinom()} \\
Poisson & \(\lambda\) & \(\lambda\) & \(\lambda\) & \texttt{dpois()} \\
& & & & \\
\textbf{Continuas} & & & & \\
Uniforme & \(a, b\) & \(\frac{a+b}{2}\) & \(\frac{(b-a)^2}{12}\) &
\texttt{dunif()} \\
Exponencial & \(\lambda\) & \(\frac{1}{\lambda}\) &
\(\frac{1}{\lambda^2}\) & \texttt{dexp()} \\
Normal & \(\mu, \sigma\) & \(\mu\) & \(\sigma^2\) & \texttt{dnorm()} \\
& & & & \\
\textbf{Muestreo} & & & & \\
Chi-cuadrado & \(k\) & \(k\) & \(2k\) & \texttt{dchisq()} \\
t & \(k\) & \(0\) & \(\frac{k}{k-2}\) & \texttt{dt()} \\
F & \(k_1, k_2\) & \(\frac{k_2}{k_2-2}\) & --- & \texttt{df()} \\
\end{longtable}

\section{Resumen de Conceptos Clave}\label{sec-semana04-resumen}

\begin{tcolorbox}[enhanced jigsaw, breakable, titlerule=0mm, rightrule=.15mm, leftrule=.75mm, left=2mm, opacitybacktitle=0.6, colback=white, bottomtitle=1mm, colbacktitle=quarto-callout-note-color!10!white, title=\textcolor{quarto-callout-note-color}{\faInfo}\hspace{0.5em}{Conceptos Fundamentales}, toptitle=1mm, coltitle=black, opacityback=0, arc=.35mm, bottomrule=.15mm, toprule=.15mm, colframe=quarto-callout-note-color-frame]

\begin{enumerate}
\def\labelenumi{\arabic{enumi}.}
\item
  \textbf{Variables Aleatorias:} Funciones que asignan números a
  resultados de experimentos aleatorios.
\item
  \textbf{Esperanza:} Media ponderada; resume la localización.
\item
  \textbf{Varianza:} Dispersión alrededor de la media; mide
  variabilidad.
\item
  \textbf{Estandarización:} Transforma variables para comparación
  directa (z-score).
\item
  \textbf{Distribuciones Discretas:} Binomial (conteos en n ensayos),
  Poisson (eventos raros).
\item
  \textbf{Distribuciones Continuas:} Normal (datos naturales),
  Exponencial (tiempos de espera), Uniforme (selección aleatoria).
\item
  \textbf{Distribuciones para Muestreo:} Chi-cuadrado, t, F
  (inferenciales, no descriptivas).
\item
  \textbf{Teorema Central del Límite:} Las medias muestrales son
  normales, sin importar la distribución original.
\end{enumerate}

\end{tcolorbox}

\section{Ejercicios}\label{sec-semana04-ejercicios}

\subsection{Ejercicio 1: Distribución
Binomial}\label{ejercicio-1-distribucion-binomial}

Un ensayo clínico prueba un tratamiento con probabilidad de éxito
\(p = 0.65\) en 50 pacientes.

\begin{enumerate}
\def\labelenumi{\alph{enumi})}
\tightlist
\item
  ¿Cuál es la probabilidad de que exactamente 35 pacientes se curen?
\item
  ¿Cuál es la probabilidad de que al menos 40 se curen?
\item
  Calcula la esperanza y varianza del número de curaciones.
\end{enumerate}

\subsection{Ejercicio 2: Distribución
Normal}\label{ejercicio-2-distribucion-normal}

El colesterol LDL en una población sigue \(N(130, 40)\) mg/dL.

\begin{enumerate}
\def\labelenumi{\alph{enumi})}
\tightlist
\item
  ¿Qué proporción de personas tienen LDL \textgreater{} 160 mg/dL
  (riesgo alto)?
\item
  ¿Cuál es el percentil 90 del LDL?
\item
  Un paciente tiene z-score = 1.5. ¿Cuál es su valor de LDL?
\end{enumerate}

\subsection{Ejercicio 3: Distribución de
Poisson}\label{ejercicio-3-distribucion-de-poisson}

El número de infecciones nosocomiales en una clínica sigue una Poisson
con \(\lambda = 2.5\) por mes.

\begin{enumerate}
\def\labelenumi{\alph{enumi})}
\tightlist
\item
  ¿Cuál es la probabilidad de que no haya infecciones en un mes?
\item
  ¿Cuál es la probabilidad de que haya más de 5 infecciones?
\item
  ¿Cuál es el número esperado de infecciones en 6 meses?
\end{enumerate}

\subsection{Ejercicio 4: Teorema Central del
Límite}\label{ejercicio-4-teorema-central-del-limite}

Se muestran 100 recién nacidos de bajo peso al nacer y prematuridad con
una altura media 45 cm y una desviación estándar 8 cm.

\begin{enumerate}
\def\labelenumi{\alph{enumi})}
\tightlist
\item
  ¿Cuál es la distribución aproximada de la media muestral \(\bar{X}\)?
\item
  ¿Cual es \(P(\bar{X} \geq 46)\)?
\item
  ¿Qué tamaño de muestra se necesita para que el error estándar sea
  menor de 0.5 cm?
\end{enumerate}

\subsection{Ejercicio 5: Transformación
Z}\label{ejercicio-5-transformacion-z}

Las medidas de glucosa en sangre siguen \(N(100, 100)\) mg/dL.

\begin{enumerate}
\def\labelenumi{\alph{enumi})}
\tightlist
\item
  Estandariza los valores 90, 100 y 120 mg/dL.
\item
  ¿Qué valor original corresponde a z = -1.96?
\item
  Interpreta: ``un paciente tiene z-score = 2.5''.
\end{enumerate}

\section{Respuestas a los Ejercicios}\label{sec-semana04-respuestas}

\textbf{Ejercicio 1: Distribución Binomial}\\
- a) P(X=35) = C(50,35)(0.65)³⁵(0.35)¹⁵ ≈ 0.0622\\
- b) P(X≥40) = Σ P(X=k) para k=40..50 ≈ 0.0339\\
- c) E{[}X{]} = np = 50×0.65 = 32.5; Var{[}X{]} = np(1-p) = 50×0.65×0.35
= 11.375

\textbf{Ejercicio 2: Distribución Normal N(130,40)}\\
- a) Z = (160-130)/√40 = 4.74; P(LDL\textgreater160) ≈ 0.0001 (muy
raro)\\
- b) P₉₀: z=1.28, LDL = 130 + 1.28√40 ≈ 130 + 8.1 = 138.1 mg/dL\\
- c) LDL = 130 + 1.5√40 ≈ 130 + 9.49 = 139.49 mg/dL

\textbf{Ejercicio 3: Distribución Poisson (λ=2.5)}\\
- a) P(X=0) = e\^{}(-2.5) ≈ 0.0821\\
- b) P(X\textgreater5) = 1 - P(X≤5) ≈ 0.0420\\
- c) En 6 meses: λ\_total = 2.5×6 = 15; E{[}X{]} = 15

\textbf{Ejercicio 4: Teorema Central del Límite}\\
- a) \(\bar{X} \sim N(45, \sigma_{\bar{X}}^2)\) donde
\(\sigma_{\bar{X}} = 8/√100 = 0.8\) cm\\
- b) Z = (46-45)/0.8 = 1.25; P(≥46) ≈ 0.1056\\
- c) SE \textless{} 0.5 ⟹ 8/√n \textless{} 0.5 ⟹ n \textgreater{} 256

\textbf{Ejercicio 5: Estandarización}\\
- a) z(90) = (90-100)/10 = -1; z(100) = 0; z(120) = 2\\
- b) x = 100 + (-1.96)(10) = 80.4 mg/dL\\
- c) El paciente está 2.5 desviaciones estándar por encima de la media
(valor muy alto, atípico)

\section{Lecturas Complementarias}\label{sec-semana04-lectura}

\begin{itemize}
\tightlist
\item
  \textbf{Forbes et al.~(2011):} Distribuciones de probabilidad
  continuadas y sus propiedades (referencia técnica).
\item
  \textbf{Tablas Normales:} Los z-scores están tabulados en libros de
  estadística y disponibles en R mediante \texttt{qnorm()}.
\item
  \textbf{Relaciones Distribucionales: }Saad, F. A., Freer, C. E.,
  Rinard, M. C., \& Mansinghka, V. K. (2020).** Optimal Approximate
  Sampling from Discrete Probability Distributions. \emph{Proceedings of
  the ACM on Programming Languages}, 4(POPL), Article 36.
  \href{https://doi.org/10.1145/3371104}{DOI: 10.1145/3371104}. Este
  trabajo discute la exactitud y los trade-offs entre métodos de
  muestreo exactos y aproximados para distribuciones discretas, un
  concepto extensible a la precisión de cálculos de probabilidad.
\end{itemize}

\begin{tcolorbox}[enhanced jigsaw, breakable, titlerule=0mm, rightrule=.15mm, leftrule=.75mm, left=2mm, opacitybacktitle=0.6, colback=white, bottomtitle=1mm, colbacktitle=quarto-callout-tip-color!10!white, title=\textcolor{quarto-callout-tip-color}{\faLightbulb}\hspace{0.5em}{Métodos Avanzados}, toptitle=1mm, coltitle=black, opacityback=0, arc=.35mm, bottomrule=.15mm, toprule=.15mm, colframe=quarto-callout-tip-color-frame]

Para ampliar los contenidos de este capítulo con técnicas estadísticas
avanzadas, visita:

\textbf{\href{https://migariane.github.io/Bioestadistica}{→
Bioestadística Avanzada --- M.A.~Luque Fernández}}

\end{tcolorbox}

\chapter{Semana 5 --- Muestreo y Distribuciones
Muestrales}\label{sec-semana05}

En esta semana estudiamos los fundamentos del muestreo y las propiedades
estadísticas de los estimadores. Los conceptos que aprenderemos son la
base de toda inferencia estadística: cómo usar datos de una muestra para
hacer afirmaciones sobre la población completa.

\section{Conceptos Fundamentales de Muestreo}\label{sec-semana05-s1}

\subsection{¿Qué es el Muestreo?}\label{quuxe9-es-el-muestreo}

\begin{tcolorbox}[enhanced jigsaw, breakable, titlerule=0mm, rightrule=.15mm, leftrule=.75mm, left=2mm, opacitybacktitle=0.6, colback=white, bottomtitle=1mm, colbacktitle=quarto-callout-note-color!10!white, title=\textcolor{quarto-callout-note-color}{\faInfo}\hspace{0.5em}{Definición: Muestreo}, toptitle=1mm, coltitle=black, opacityback=0, arc=.35mm, bottomrule=.15mm, toprule=.15mm, colframe=quarto-callout-note-color-frame]

El \textbf{muestreo} es el proceso de seleccionar una parte de la
población para observar y estudiar, de manera que podamos hacer
inferencias sobre características de interés en toda la población.

\end{tcolorbox}

Imagine que queremos estimar el gasto en alimentos de las familias en
España. Tenemos dos opciones:

\begin{itemize}
\tightlist
\item
  \textbf{Censo}: Encuestar a TODAS las familias del país (costoso,
  largo, casi imposible)
\item
  \textbf{Muestra}: Encuestar al 5\% de las familias e inferir sobre
  toda la población (barato, rápido, práctico)
\end{itemize}

\subsubsection{Notación Fundamental}\label{notaciuxf3n-fundamental}

\begin{longtable}[]{@{}ll@{}}
\toprule\noalign{}
Símbolo & Significado \\
\midrule\noalign{}
\endhead
\bottomrule\noalign{}
\endlastfoot
\(N\) & Tamaño de la población \\
\(n\) & Tamaño de la muestra \\
\(f = n/N\) & Fracción muestral \\
\(X_i\) & \(i\)-ésima observación (variable aleatoria) \\
\(x_i\) & \(i\)-ésimo valor observado (número) \\
\end{longtable}

\subsection{Parámetros Poblacionales}\label{paruxe1metros-poblacionales}

Los parámetros son características desconocidas de la población que
queremos estimar:

\begin{tcolorbox}[enhanced jigsaw, breakable, titlerule=0mm, rightrule=.15mm, leftrule=.75mm, left=2mm, opacitybacktitle=0.6, colback=white, bottomtitle=1mm, colbacktitle=quarto-callout-warning-color!10!white, title=\textcolor{quarto-callout-warning-color}{\faExclamationTriangle}\hspace{0.5em}{Parámetros de Interés}, toptitle=1mm, coltitle=black, opacityback=0, arc=.35mm, bottomrule=.15mm, toprule=.15mm, colframe=quarto-callout-warning-color-frame]

\begin{itemize}
\tightlist
\item
  \textbf{Media poblacional}: \(\mu = \frac{1}{N}\sum_{i=1}^N x_i\)
\item
  \textbf{Varianza poblacional}:
  \(\sigma^2 = \frac{1}{N}\sum_{i=1}^N (x_i - \mu)^2\)
\item
  \textbf{Proporción poblacional} (variables binarias):
  \(\pi = \frac{1}{N}\sum_{i=1}^N x_i\), donde \(x_i \in \{0,1\}\)
\end{itemize}

\end{tcolorbox}

\subsection{Parámetros
vs.~Estadísticos}\label{paruxe1metros-vs.-estaduxedsticos}

\begin{tcolorbox}[enhanced jigsaw, breakable, titlerule=0mm, rightrule=.15mm, leftrule=.75mm, left=2mm, opacitybacktitle=0.6, colback=white, bottomtitle=1mm, colbacktitle=quarto-callout-note-color!10!white, title=\textcolor{quarto-callout-note-color}{\faInfo}\hspace{0.5em}{Definiciones Clave}, toptitle=1mm, coltitle=black, opacityback=0, arc=.35mm, bottomrule=.15mm, toprule=.15mm, colframe=quarto-callout-note-color-frame]

Un \textbf{parámetro} es una característica desconocida de la población
(p.ej., \(\mu\), \(\sigma^2\), \(\pi\)).

Un \textbf{estadístico} es una función de la muestra calculada a partir
de datos observados. Los estadísticos se usan para estimar parámetros
desconocidos.

\end{tcolorbox}

\subsection{El Problema del Sesgo de Muestreo: El Caso del Literary
Digest
(1936)}\label{el-problema-del-sesgo-de-muestreo-el-caso-del-literary-digest-1936}

La importancia de un muestreo adecuado se ilustra con uno de los mayores
fracasos en la historia de las encuestas:

\begin{tcolorbox}[enhanced jigsaw, breakable, titlerule=0mm, rightrule=.15mm, leftrule=.75mm, left=2mm, opacitybacktitle=0.6, colback=white, bottomtitle=1mm, colbacktitle=quarto-callout-warning-color!10!white, title=\textcolor{quarto-callout-warning-color}{\faExclamationTriangle}\hspace{0.5em}{Advertencia: El Ejemplo del Literary Digest}, toptitle=1mm, coltitle=black, opacityback=0, arc=.35mm, bottomrule=.15mm, toprule=.15mm, colframe=quarto-callout-warning-color-frame]

En la elección presidencial de EE.UU. de 1936:

\begin{itemize}
\tightlist
\item
  La revista \textbf{The Literary Digest} encuestó a \textbf{10 millones
  de personas} (seleccionadas de propietarios de teléfono y coche, y
  suscriptores de la revista)
\item
  Recibieron 2.3 millones de respuestas
\item
  \textbf{Predicción}: 60\% republicanos, 40\% demócratas
\item
  \textbf{Resultado real}: 36.5\% republicanos, 60.8\% demócratas
\item
  \textbf{Gallup Institute}, con solo 50,000 encuestados aleatorios,
  predijo correctamente
\end{itemize}

\textbf{Lección}: Un número grande de respuestas NO garantiza una
muestra sin sesgo. El sesgo de selección (muestra no representativa) es
más importante que el tamaño.

\end{tcolorbox}

\begin{center}\rule{0.5\linewidth}{0.5pt}\end{center}

\section{Tipos de Muestreo}\label{sec-semana05-s2}

Para obtener muestras representativas y hacer inferencias válidas,
existen varios diseños de muestreo:

\subsection{Muestreo Aleatorio Simple}\label{muestreo-aleatorio-simple}

\begin{tcolorbox}[enhanced jigsaw, breakable, titlerule=0mm, rightrule=.15mm, leftrule=.75mm, left=2mm, opacitybacktitle=0.6, colback=white, bottomtitle=1mm, colbacktitle=quarto-callout-note-color!10!white, title=\textcolor{quarto-callout-note-color}{\faInfo}\hspace{0.5em}{Definición}, toptitle=1mm, coltitle=black, opacityback=0, arc=.35mm, bottomrule=.15mm, toprule=.15mm, colframe=quarto-callout-note-color-frame]

En el \textbf{muestreo aleatorio simple}, cada elemento de la población
tiene la misma probabilidad de ser seleccionado.

\end{tcolorbox}

Distinguimos dos variantes:

\textbf{1. Muestreo sin reemplazo:} - Cada unidad puede ser seleccionada
como máximo una vez - Los elementos muestreados NO se devuelven a la
población - Los valores muestrales NO son independientes - Aplicable a
poblaciones finitas

\textbf{2. Muestreo con reemplazo:} - Cada unidad puede ser seleccionada
más de una vez - Los elementos se devuelven a la población después de
cada extracción - Los valores muestrales SON independientes -
Matemáticamente equivalente a una población infinita

\subsection{Muestreo Estratificado}\label{muestreo-estratificado}

\begin{tcolorbox}[enhanced jigsaw, breakable, titlerule=0mm, rightrule=.15mm, leftrule=.75mm, left=2mm, opacitybacktitle=0.6, colback=white, bottomtitle=1mm, colbacktitle=quarto-callout-note-color!10!white, title=\textcolor{quarto-callout-note-color}{\faInfo}\hspace{0.5em}{Definición}, toptitle=1mm, coltitle=black, opacityback=0, arc=.35mm, bottomrule=.15mm, toprule=.15mm, colframe=quarto-callout-note-color-frame]

En el \textbf{muestreo estratificado}, la población se divide en
subgrupos (estratos) no solapados, y se extrae una muestra aleatoria
simple de cada estrato.

\end{tcolorbox}

\textbf{Ejemplo}: Para estudiar opiniones políticas, estratificar por
región, edad, o nivel educativo asegura que cada grupo esté representado
adecuadamente.

\textbf{Ventaja}: Proporciona estimaciones más precisas que el muestreo
simple cuando existe variación dentro de los estratos.

\subsection{Muestreo por
Conglomerados}\label{muestreo-por-conglomerados}

\begin{tcolorbox}[enhanced jigsaw, breakable, titlerule=0mm, rightrule=.15mm, leftrule=.75mm, left=2mm, opacitybacktitle=0.6, colback=white, bottomtitle=1mm, colbacktitle=quarto-callout-note-color!10!white, title=\textcolor{quarto-callout-note-color}{\faInfo}\hspace{0.5em}{Definición}, toptitle=1mm, coltitle=black, opacityback=0, arc=.35mm, bottomrule=.15mm, toprule=.15mm, colframe=quarto-callout-note-color-frame]

En el \textbf{muestreo por conglomerados}, la población se divide en
conglomerados (clusters), se seleccionan aleatoriamente algunos
conglomerados, y se observan TODOS los elementos dentro de esos
conglomerados.

\end{tcolorbox}

\textbf{Ejemplo}: Para encuestar hogares en un país, seleccionar 30
ciudades al azar y encuestar todos los hogares en esas ciudades.

\textbf{Ventaja}: Más barato que muestreo simple cuando los elementos
están geográficamente dispersos. \textbf{Desventaja}: Menos preciso si
hay mucha variación entre conglomerados.

\subsection{Muestreo por Criterio (Judgment
Sampling)}\label{muestreo-por-criterio-judgment-sampling}

\begin{tcolorbox}[enhanced jigsaw, breakable, titlerule=0mm, rightrule=.15mm, leftrule=.75mm, left=2mm, opacitybacktitle=0.6, colback=white, bottomtitle=1mm, colbacktitle=quarto-callout-note-color!10!white, title=\textcolor{quarto-callout-note-color}{\faInfo}\hspace{0.5em}{Definición}, toptitle=1mm, coltitle=black, opacityback=0, arc=.35mm, bottomrule=.15mm, toprule=.15mm, colframe=quarto-callout-note-color-frame]

En el \textbf{muestreo por criterio}, los elementos se seleccionan según
el juicio de expertos, no aleatoriamente.

\end{tcolorbox}

\textbf{Advertencia}: Este tipo de muestreo puede introducir sesgo y no
permite hacer inferencia estadística formal.

\begin{center}\rule{0.5\linewidth}{0.5pt}\end{center}

\section{Estadísticos Muestrales como Variables
Aleatorias}\label{sec-semana05-s3}

\subsection{El Concepto Clave}\label{el-concepto-clave}

Antes de observar una muestra, los valores que obtendremos son
aleatorios. Por lo tanto:

\begin{itemize}
\tightlist
\item
  Los valores \(X_1, X_2, \ldots, X_n\) son \textbf{variables
  aleatorias}
\item
  Cualquier función de estos valores es también una \textbf{variable
  aleatoria}
\item
  Esto incluye la media muestral \(\bar{X}\), la proporción muestral
  \(\hat{\pi}\), y la varianza muestral \(S^2\)
\end{itemize}

\subsection{Estadísticos Importantes}\label{estaduxedsticos-importantes}

\begin{tcolorbox}[enhanced jigsaw, breakable, titlerule=0mm, rightrule=.15mm, leftrule=.75mm, left=2mm, opacitybacktitle=0.6, colback=white, bottomtitle=1mm, colbacktitle=quarto-callout-warning-color!10!white, title=\textcolor{quarto-callout-warning-color}{\faExclamationTriangle}\hspace{0.5em}{Estadísticos Muestrales}, toptitle=1mm, coltitle=black, opacityback=0, arc=.35mm, bottomrule=.15mm, toprule=.15mm, colframe=quarto-callout-warning-color-frame]

\begin{itemize}
\tightlist
\item
  \textbf{Media muestral}: \(\bar{X} = \frac{1}{n}\sum_{i=1}^n X_i\)
\item
  \textbf{Proporción muestral} (variables binarias):
  \(\hat{\pi} = \frac{1}{n}\sum_{i=1}^n X_i\)
\item
  \textbf{Varianza muestral (μ conocida)}:
  \(S^{*2} = \frac{1}{n}\sum_{i=1}^n (X_i - \mu)^2\)
\item
  \textbf{Varianza muestral (μ desconocida)}:
  \(S^2 = \frac{1}{n-1}\sum_{i=1}^n (X_i - \bar{X})^2\)
\item
  \textbf{Varianza muestral (μ desconocida, alternativa)}:
  \(S'^2 = \frac{1}{n}\sum_{i=1}^n (X_i - \bar{X})^2\)
\end{itemize}

\end{tcolorbox}

Cada uno de estos estadísticos tiene su propia \textbf{distribución
muestral}: la distribución de probabilidad del estadístico a través de
todas las posibles muestras de tamaño \(n\).

\begin{center}\rule{0.5\linewidth}{0.5pt}\end{center}

\section{Distribución Muestral de la Media}\label{sec-semana05-s4}

\subsection{Caso: Muestreo con
Reemplazo}\label{caso-muestreo-con-reemplazo}

Cuando extraemos la muestra con reemplazo (o de una población infinita),
los valores muestrales son independientes.

\begin{tcolorbox}[enhanced jigsaw, breakable, titlerule=0mm, rightrule=.15mm, leftrule=.75mm, left=2mm, opacitybacktitle=0.6, colback=white, bottomtitle=1mm, colbacktitle=quarto-callout-warning-color!10!white, title=\textcolor{quarto-callout-warning-color}{\faExclamationTriangle}\hspace{0.5em}{Propiedades de \(\bar{X}\) (con reemplazo)}, toptitle=1mm, coltitle=black, opacityback=0, arc=.35mm, bottomrule=.15mm, toprule=.15mm, colframe=quarto-callout-warning-color-frame]

Para variables aleatorias independientes \(X_i\) con \(E(X_i) = \mu\) y
\(\text{Var}(X_i) = \sigma^2\):

\[E(\bar{X}) = \mu\]

\[\text{Var}(\bar{X}) = \frac{\sigma^2}{n}\]

\[\sigma(\bar{X}) = \frac{\sigma}{\sqrt{n}} \quad \text{(error estándar)}\]

\end{tcolorbox}

\textbf{Interpretación}: Mientras mayor sea \(n\), menor será la
varianza de la media muestral. Esto explica por qué muestras grandes dan
estimaciones más precisas.

\subsection{Caso: Muestreo sin
Reemplazo}\label{caso-muestreo-sin-reemplazo}

Cuando la muestra se extrae sin reemplazo de una población finita,
existe un factor de corrección:

\begin{tcolorbox}[enhanced jigsaw, breakable, titlerule=0mm, rightrule=.15mm, leftrule=.75mm, left=2mm, opacitybacktitle=0.6, colback=white, bottomtitle=1mm, colbacktitle=quarto-callout-warning-color!10!white, title=\textcolor{quarto-callout-warning-color}{\faExclamationTriangle}\hspace{0.5em}{Propiedades de \(\bar{X}\) (sin reemplazo)}, toptitle=1mm, coltitle=black, opacityback=0, arc=.35mm, bottomrule=.15mm, toprule=.15mm, colframe=quarto-callout-warning-color-frame]

\[E(\bar{X}) = \mu\]

\[\text{Var}(\bar{X}) = \frac{\sigma^2}{n} \cdot \frac{N-n}{N-1}\]

\[\sigma(\bar{X}) = \frac{\sigma}{\sqrt{n}} \cdot \sqrt{\frac{N-n}{N-1}}\]

donde \(\frac{N-n}{N-1}\) es el \textbf{factor de corrección por
población finita} (CPF).

\end{tcolorbox}

\textbf{Nota}: Cuando \(N \gg n\), el CPF \(\approx 1\) y recuperamos la
fórmula del caso con reemplazo.

\subsection{\texorpdfstring{Distribución Exacta Cuando
\(X \sim N(\mu, \sigma^2)\)}{Distribución Exacta Cuando X \textbackslash sim N(\textbackslash mu, \textbackslash sigma\^{}2)}}\label{distribuciuxf3n-exacta-cuando-x-sim-nmu-sigma2}

\begin{tcolorbox}[enhanced jigsaw, breakable, titlerule=0mm, rightrule=.15mm, leftrule=.75mm, left=2mm, opacitybacktitle=0.6, colback=white, bottomtitle=1mm, colbacktitle=quarto-callout-warning-color!10!white, title=\textcolor{quarto-callout-warning-color}{\faExclamationTriangle}\hspace{0.5em}{Resultado Fundamental}, toptitle=1mm, coltitle=black, opacityback=0, arc=.35mm, bottomrule=.15mm, toprule=.15mm, colframe=quarto-callout-warning-color-frame]

Si \(X_1, X_2, \ldots, X_n\) son i.i.d. \(N(\mu, \sigma^2)\), entonces:

\[\bar{X} \sim N\left(\mu, \frac{\sigma^2}{n}\right)\]

Estandarizando:
\[Z = \frac{\bar{X} - \mu}{\sigma/\sqrt{n}} \sim N(0, 1)\]

Este estadístico \(Z\) sigue una \textbf{distribución normal estándar}.

\end{tcolorbox}

\begin{center}\rule{0.5\linewidth}{0.5pt}\end{center}

\section{Distribución Muestral de la Proporción}\label{sec-semana05-s5}

Cuando estudiamos variables binarias (éxito/fracaso, sí/no), la
proporción muestral es el estadístico de interés.

\subsection{Con Reemplazo}\label{con-reemplazo}

\begin{tcolorbox}[enhanced jigsaw, breakable, titlerule=0mm, rightrule=.15mm, leftrule=.75mm, left=2mm, opacitybacktitle=0.6, colback=white, bottomtitle=1mm, colbacktitle=quarto-callout-warning-color!10!white, title=\textcolor{quarto-callout-warning-color}{\faExclamationTriangle}\hspace{0.5em}{Distribución de \(\hat{\pi}\) (con reemplazo)}, toptitle=1mm, coltitle=black, opacityback=0, arc=.35mm, bottomrule=.15mm, toprule=.15mm, colframe=quarto-callout-warning-color-frame]

Para una muestra de variables Bernoulli independientes con parámetro
\(\pi\):

\[E(\hat{\pi}) = \pi\]

\[\text{Var}(\hat{\pi}) = \frac{\pi(1-\pi)}{n}\]

\[\sigma(\hat{\pi}) = \sqrt{\frac{\pi(1-\pi)}{n}}\]

\end{tcolorbox}

\textbf{Ejemplo práctico}: Si queremos estimar la proporción de votantes
que apoyan una política, con \(\pi = 0.5\) y \(n = 1000\):

\[\sigma(\hat{\pi}) = \sqrt{\frac{0.5 \times 0.5}{1000}} = \sqrt{0.00025} \approx 0.0158\]

Esto significa que la proporción muestral varía típicamente en ±1.58\%
entre diferentes muestras.

\subsection{Normalidad Asintótica}\label{normalidad-asintuxf3tica}

Para muestras grandes (\(n > 30\) o \(n\pi > 5\) y \(n(1-\pi) > 5\)):

\[Z = \frac{\hat{\pi} - \pi}{\sqrt{\pi(1-\pi)/n}} \approx N(0,1)\]

\begin{center}\rule{0.5\linewidth}{0.5pt}\end{center}

\section{Distribución Muestral de la Varianza}\label{sec-semana05-s6}

La varianza muestral también es una variable aleatoria y su distribución
depende de si conocemos la media poblacional.

\subsection{\texorpdfstring{Caso: \(\mu\)
Conocida}{Caso: \textbackslash mu Conocida}}\label{caso-mu-conocida}

Cuando la media poblacional \(\mu\) es conocida:

\begin{tcolorbox}[enhanced jigsaw, breakable, titlerule=0mm, rightrule=.15mm, leftrule=.75mm, left=2mm, opacitybacktitle=0.6, colback=white, bottomtitle=1mm, colbacktitle=quarto-callout-warning-color!10!white, title=\textcolor{quarto-callout-warning-color}{\faExclamationTriangle}\hspace{0.5em}{Distribución de \(S^{*2}\) (μ conocida)}, toptitle=1mm, coltitle=black, opacityback=0, arc=.35mm, bottomrule=.15mm, toprule=.15mm, colframe=quarto-callout-warning-color-frame]

Si \(X_i \sim N(\mu, \sigma^2)\) y \(\mu\) es conocida:

\[S^{*2} = \frac{1}{n}\sum_{i=1}^n (X_i - \mu)^2\]

\[\frac{nS^{*2}}{\sigma^2} \sim \chi^2_n\]

donde \(\chi^2_n\) es la distribución chi-cuadrado con \(n\) grados de
libertad.

\textbf{Propiedades:} - \(E(S^{*2}) = \sigma^2\) (insesgada) -
\(\text{Var}(S^{*2}) = \frac{2\sigma^4}{n}\)

\end{tcolorbox}

\subsection{\texorpdfstring{Caso: \(\mu\)
Desconocida}{Caso: \textbackslash mu Desconocida}}\label{caso-mu-desconocida}

Cuando estimamos la media con \(\bar{X}\), pierden un grado de libertad:

\begin{tcolorbox}[enhanced jigsaw, breakable, titlerule=0mm, rightrule=.15mm, leftrule=.75mm, left=2mm, opacitybacktitle=0.6, colback=white, bottomtitle=1mm, colbacktitle=quarto-callout-warning-color!10!white, title=\textcolor{quarto-callout-warning-color}{\faExclamationTriangle}\hspace{0.5em}{Distribución de \(S^2\) (μ desconocida)}, toptitle=1mm, coltitle=black, opacityback=0, arc=.35mm, bottomrule=.15mm, toprule=.15mm, colframe=quarto-callout-warning-color-frame]

Si \(X_i \sim N(\mu, \sigma^2)\) y \(\mu\) es desconocida:

\[S^2 = \frac{1}{n-1}\sum_{i=1}^n (X_i - \bar{X})^2\]

\[\frac{(n-1)S^2}{\sigma^2} \sim \chi^2_{n-1}\]

\textbf{Propiedades:} - \(E(S^2) = \sigma^2\) (insesgada) -
\(\text{Var}(S^2) = \frac{2\sigma^4}{n-1}\)

\end{tcolorbox}

\textbf{Alternativa sesgada}:
\[S'^2 = \frac{1}{n}\sum_{i=1}^n (X_i - \bar{X})^2 = \frac{n-1}{n}S^2\]

\begin{itemize}
\tightlist
\item
  \(E(S'^2) = \frac{n-1}{n}\sigma^2\) (sesgada, subestima \(\sigma^2\))
\end{itemize}

\textbf{Conclusión}: Use \(S^2\) (dividir por \(n-1\)) para estimaciones
insesgadas de \(\sigma^2\).

\begin{center}\rule{0.5\linewidth}{0.5pt}\end{center}

\section{Teorema Central del Límite (Versión
Formal)}\label{sec-semana05-s7}

El Teorema Central del Límite (TCL) es el resultado más importante en
estadística. Explica por qué la distribución normal es ubicua.

\begin{tcolorbox}[enhanced jigsaw, breakable, titlerule=0mm, rightrule=.15mm, leftrule=.75mm, left=2mm, opacitybacktitle=0.6, colback=white, bottomtitle=1mm, colbacktitle=quarto-callout-warning-color!10!white, title=\textcolor{quarto-callout-warning-color}{\faExclamationTriangle}\hspace{0.5em}{Teorema Central del Límite}, toptitle=1mm, coltitle=black, opacityback=0, arc=.35mm, bottomrule=.15mm, toprule=.15mm, colframe=quarto-callout-warning-color-frame]

Sea \(X_1, X_2, \ldots\) una secuencia de variables aleatorias
independientes e idénticamente distribuidas (i.i.d.) con media \(\mu\) y
varianza \(\sigma^2\) finitas.

Define \(S_n = X_1 + X_2 + \cdots + X_n\). Entonces:

\[\frac{S_n - n\mu}{\sigma\sqrt{n}} \xrightarrow{d} N(0, 1) \quad \text{cuando} \quad n \to \infty\]

Equivalentemente, para la media muestral:

\[\frac{\bar{X} - \mu}{\sigma/\sqrt{n}} \xrightarrow{d} N(0, 1)\]

\end{tcolorbox}

\subsection{Interpretación}\label{interpretaciuxf3n-3}

\textbf{Mensaje principal}: La suma (y la media) de muchas variables
aleatorias ES APROXIMADAMENTE normal, sin importar la distribución
original de cada variable.

\textbf{Requisito único}: Que las variables tengan varianza finita. La
distribución original puede ser cualquiera: exponencial, uniforme,
Bernoulli, etc.

\textbf{Implicación práctica}: Para \(n\) suficientemente grande
(típicamente \(n \geq 30\)), la media muestral es aproximadamente
normal, incluso si los datos originales no son normales. Esto justifica
el uso de técnicas basadas en la normalidad para inferencia estadística,
como intervalos de confianza y pruebas de hipótesis, en una amplia
variedad de situaciones.

\begin{tcolorbox}[enhanced jigsaw, breakable, titlerule=0mm, rightrule=.15mm, leftrule=.75mm, left=2mm, opacitybacktitle=0.6, colback=white, bottomtitle=1mm, colbacktitle=quarto-callout-warning-color!10!white, title=\textcolor{quarto-callout-warning-color}{\faExclamationTriangle}\hspace{0.5em}{Advertencia}, toptitle=1mm, coltitle=black, opacityback=0, arc=.35mm, bottomrule=.15mm, toprule=.15mm, colframe=quarto-callout-warning-color-frame]

Notese que si la población es extremadamente asimétrica, se podría
necesitar \textbf{\(n \ge 60\)} o más para que la aproximación normal
sea adecuada.

\end{tcolorbox}

\begin{center}\rule{0.5\linewidth}{0.5pt}\end{center}

\section{La Distribución t de Student}\label{sec-semana05-s8}

En la práctica, casi nunca conocemos \(\sigma\). Cuando estimamos la
desviación típica con la muestra, la distribución cambia.

\begin{tcolorbox}[enhanced jigsaw, breakable, titlerule=0mm, rightrule=.15mm, leftrule=.75mm, left=2mm, opacitybacktitle=0.6, colback=white, bottomtitle=1mm, colbacktitle=quarto-callout-warning-color!10!white, title=\textcolor{quarto-callout-warning-color}{\faExclamationTriangle}\hspace{0.5em}{La Distribución t de Student}, toptitle=1mm, coltitle=black, opacityback=0, arc=.35mm, bottomrule=.15mm, toprule=.15mm, colframe=quarto-callout-warning-color-frame]

Si \(X_i \sim N(\mu, \sigma^2)\) (normales) y \(\sigma^2\) es
DESCONOCIDA, entonces:

\[T = \frac{\bar{X} - \mu}{S/\sqrt{n}} \sim t_{n-1}\]

donde \(t_{n-1}\) es la \textbf{distribución t de Student con \(n-1\)
grados de libertad}.

\(S = \sqrt{S^2}\) es la desviación típica muestral (raíz cuadrada del
estimador insesgado de varianza).

\end{tcolorbox}

\subsection{Propiedades de la Distribución
t}\label{propiedades-de-la-distribuciuxf3n-t}

\begin{itemize}
\tightlist
\item
  \textbf{Forma}: Simétrica, con colas más pesadas que la normal
  estándar \(N(0,1)\)
\item
  \textbf{Media}: \(E(T) = 0\) (para \(n > 1\))
\item
  \textbf{Varianza}: \(\text{Var}(T) = \frac{n-1}{n-3}\) (para
  \(n > 3\))
\item
  \textbf{Convergencia}: Conforme \(n \to \infty\), \(t_n \to N(0,1)\)
\item
  \textbf{Regla práctica}: Para \(n > 30\), \(t_n \approx N(0,1)\)
\end{itemize}

\subsection{Cuándo Usar Cada
Distribución}\label{cuuxe1ndo-usar-cada-distribuciuxf3n}

\begin{longtable}[]{@{}
  >{\raggedright\arraybackslash}p{(\linewidth - 2\tabcolsep) * \real{0.4400}}
  >{\raggedright\arraybackslash}p{(\linewidth - 2\tabcolsep) * \real{0.5600}}@{}}
\toprule\noalign{}
\begin{minipage}[b]{\linewidth}\raggedright
Situación
\end{minipage} & \begin{minipage}[b]{\linewidth}\raggedright
Distribución
\end{minipage} \\
\midrule\noalign{}
\endhead
\bottomrule\noalign{}
\endlastfoot
\(X \sim N(\mu, \sigma^2)\), \(\sigma\) conocida &
\(Z = \frac{\bar{X}-\mu}{\sigma/\sqrt{n}} \sim N(0,1)\) \\
\(X \sim N(\mu, \sigma^2)\), \(\sigma\) desconocida &
\(T = \frac{\bar{X}-\mu}{S/\sqrt{n}} \sim t_{n-1}\) \\
Cualquier distribución, \(n > 30\), \(\sigma\) conocida &
\(Z \approx N(0,1)\) (TCL) \\
Cualquier distribución, \(n > 30\), \(\sigma\) desconocida &
\(T \approx N(0,1)\) (TCL aproximado) \\
\end{longtable}

\begin{center}\rule{0.5\linewidth}{0.5pt}\end{center}

\section{Ejemplo 5.1: Simulando el Teorema Central del
Límite}\label{sec-semana05-s9}

\begin{tcolorbox}[enhanced jigsaw, breakable, titlerule=0mm, rightrule=.15mm, leftrule=.75mm, left=2mm, opacitybacktitle=0.6, colback=white, bottomtitle=1mm, colbacktitle=quarto-callout-tip-color!10!white, title=\textcolor{quarto-callout-tip-color}{\faLightbulb}\hspace{0.5em}{Ejemplo 5.2: TCL con Variables Exponenciales}, toptitle=1mm, coltitle=black, opacityback=0, arc=.35mm, bottomrule=.15mm, toprule=.15mm, colframe=quarto-callout-tip-color-frame]

Simularemos el TCL mostrando que la suma de variables exponenciales (muy
asimétricas) se aproxima a una distribución normal conforme aumentamos
el número de variables.

\pandocbounded{\includegraphics[keepaspectratio]{chapters/semana05_files/figure-pdf/unnamed-chunk-1-1.pdf}}

\begin{quote}
\textbf{Interpretación}

La simulación demuestra empíricamente el Teorema Central del Límite. La
suma de dos variables exponenciales (distribución altamente asimétrica
hacia la derecha) aún muestra desviación respecto a la normalidad, con
solapamiento limitado entre histograma y curva teórica. Con diez
exponenciales sumadas, la concordancia mejora notablemente. Con cien
exponenciales, la distribución muestral es prácticamente indistinguible
de una curva normal (línea roja), a pesar de que cada variable
individual es exponencial pura. Este comportamiento es independiente de
la distribución original y requiere únicamente varianza finita,
ilustrando por qué la normalidad emerge como distribución universal en
estadística aplicada.
\end{quote}

\end{tcolorbox}

\begin{tcolorbox}[enhanced jigsaw, breakable, titlerule=0mm, rightrule=.15mm, leftrule=.75mm, left=2mm, opacitybacktitle=0.6, colback=white, bottomtitle=1mm, colbacktitle=quarto-callout-tip-color!10!white, title=\textcolor{quarto-callout-tip-color}{\faLightbulb}\hspace{0.5em}{Ejemplo: Distribuciones Muestrales de la Media}, toptitle=1mm, coltitle=black, opacityback=0, arc=.35mm, bottomrule=.15mm, toprule=.15mm, colframe=quarto-callout-tip-color-frame]

Comparamos la distribución teórica con simulaciones.

\pandocbounded{\includegraphics[keepaspectratio]{chapters/semana05_files/figure-pdf/unnamed-chunk-2-1.pdf}}

\begin{quote}
\textbf{Interpretación}

La simulación de 1000 muestras (n=30 cada una) de una población normal
N(100,15) genera medias muestrales distribuidas teóricamente como N(100,
15/√30) = N(100, 2.74). El histograma (barras azules) se superpone casi
perfectamente con la curva normal teórica (línea roja), confirmando que
el error estándar de la media es σ/√n = 15/√30 ≈ 2.74. Esta concordancia
exacta entre simulación y teoría valida la predicción analítica de
distribuciones muestrales y justifica el uso de la distribución normal
para intervalos de confianza e inferencia en muestras grandes.
\end{quote}

\end{tcolorbox}

\begin{center}\rule{0.5\linewidth}{0.5pt}\end{center}

\section{Resumen de Distribuciones
Muestrales}\label{sec-semana05-resumen}

\begin{tcolorbox}[enhanced jigsaw, breakable, titlerule=0mm, rightrule=.15mm, leftrule=.75mm, left=2mm, opacitybacktitle=0.6, colback=white, bottomtitle=1mm, colbacktitle=quarto-callout-warning-color!10!white, title=\textcolor{quarto-callout-warning-color}{\faExclamationTriangle}\hspace{0.5em}{Tabla Resumen: Estadísticos y sus Distribuciones}, toptitle=1mm, coltitle=black, opacityback=0, arc=.35mm, bottomrule=.15mm, toprule=.15mm, colframe=quarto-callout-warning-color-frame]

\end{tcolorbox}

\begin{longtable}[]{@{}
  >{\raggedright\arraybackslash}p{(\linewidth - 8\tabcolsep) * \real{0.2791}}
  >{\raggedright\arraybackslash}p{(\linewidth - 8\tabcolsep) * \real{0.2791}}
  >{\centering\arraybackslash}p{(\linewidth - 8\tabcolsep) * \real{0.2093}}
  >{\centering\arraybackslash}p{(\linewidth - 8\tabcolsep) * \real{0.1163}}
  >{\centering\arraybackslash}p{(\linewidth - 8\tabcolsep) * \real{0.1163}}@{}}
\toprule\noalign{}
\begin{minipage}[b]{\linewidth}\raggedright
Estadístico
\end{minipage} & \begin{minipage}[b]{\linewidth}\raggedright
Condiciones
\end{minipage} & \begin{minipage}[b]{\linewidth}\centering
\(E(\cdot)\)
\end{minipage} & \begin{minipage}[b]{\linewidth}\centering
\(\text{Var}(\cdot)\)
\end{minipage} & \begin{minipage}[b]{\linewidth}\centering
Distribución
\end{minipage} \\
\midrule\noalign{}
\endhead
\bottomrule\noalign{}
\endlastfoot
\(\bar{X}\) & RS c/ reemplazo, \(\sigma\) conocida & \(\mu\) &
\(\sigma^2/n\) & \(N(\mu, \sigma^2/n)\) \\
\(\bar{X}\) & RS s/ reemplazo, \(\sigma\) conocida & \(\mu\) &
\(\sigma^2/n \cdot \frac{N-n}{N-1}\) & \(N(\mu, \cdots)\) \\
\(\bar{X}\) & \(X \sim N(\mu,\sigma^2)\), \(\sigma\) conocida & \(\mu\)
& \(\sigma^2/n\) & \(N(\mu, \sigma^2/n)\) \\
\(\bar{X}\) & \(X \sim N(\mu,\sigma^2)\), \(\sigma\) desconocida &
\(\mu\) & \(S^2/n\) & \(t_{n-1}\) \\
\(\hat{\pi}\) & Variables binarias, \(n\) grande & \(\pi\) &
\(\pi(1-\pi)/n\) & \(N(\pi, \frac{\pi(1-\pi)}{n})\) \\
\(S^2\) & \(X \sim N(\mu,\sigma^2)\), \(\mu\) desconocida & \(\sigma^2\)
& \(\frac{2\sigma^4}{n-1}\) &
\(\frac{(n-1)S^2}{\sigma^2} \sim \chi^2_{n-1}\) \\
\end{longtable}

\begin{center}\rule{0.5\linewidth}{0.5pt}\end{center}

\section{Corrección por Población Finita}\label{sec-semana05-cpf}

Cuando la población es finita y la fracción muestral \(f = n/N\) es
apreciable (típicamente \(f > 0.05\)), se aplica la \textbf{corrección
por población finita}:

\begin{tcolorbox}[enhanced jigsaw, breakable, titlerule=0mm, rightrule=.15mm, leftrule=.75mm, left=2mm, opacitybacktitle=0.6, colback=white, bottomtitle=1mm, colbacktitle=quarto-callout-note-color!10!white, title=\textcolor{quarto-callout-note-color}{\faInfo}\hspace{0.5em}{Factor de Corrección por Población Finita}, toptitle=1mm, coltitle=black, opacityback=0, arc=.35mm, bottomrule=.15mm, toprule=.15mm, colframe=quarto-callout-note-color-frame]

\[\text{CPF} = \sqrt{\frac{N-n}{N-1}}\]

La varianza del estimador se multiplica por este factor:
\[\text{Var}(\bar{X}) = \frac{\sigma^2}{n} \cdot \frac{N-n}{N-1}\]

\end{tcolorbox}

\textbf{Casos extremos:} - Si \(N \to \infty\), entonces CPF \(\to 1\)
(población infinita, no hay corrección) - Si \(n = N\) (censo), entonces
CPF \(= 0\) (varianza = 0, estimación perfecta) - Si \(n = N/2\),
entonces CPF \(\approx 0.71\)

\begin{center}\rule{0.5\linewidth}{0.5pt}\end{center}

\section{Ejercicios}\label{sec-semana05-ejercicios}

\begin{tcolorbox}[enhanced jigsaw, breakable, titlerule=0mm, rightrule=.15mm, leftrule=.75mm, left=2mm, opacitybacktitle=0.6, colback=white, bottomtitle=1mm, colbacktitle=quarto-callout-tip-color!10!white, title=\textcolor{quarto-callout-tip-color}{\faLightbulb}\hspace{0.5em}{Ejercicio 5.1: Varianza de la Media Muestral}, toptitle=1mm, coltitle=black, opacityback=0, arc=.35mm, bottomrule=.15mm, toprule=.15mm, colframe=quarto-callout-tip-color-frame]

Un hospital desea estimar el peso medio al nacer de los neonatos en una
región. De una población de 5000 nacimientos anuales, toma una muestra
de 100. Se sabe que la desviación típica poblacional es \(\sigma = 200\)
gramos.

\begin{enumerate}
\def\labelenumi{(\alph{enumi})}
\tightlist
\item
  Calcule la varianza de la media muestral \(\bar{X}\).
\item
  Calcule el error estándar \(\sigma(\bar{X})\).
\item
  ¿Cuál sería el error estándar si la muestra fuera de 400 neonatos?
\item
  ¿Qué efecto tiene aumentar el tamaño muestral?
\end{enumerate}

\end{tcolorbox}

\begin{tcolorbox}[enhanced jigsaw, breakable, titlerule=0mm, rightrule=.15mm, leftrule=.75mm, left=2mm, opacitybacktitle=0.6, colback=white, bottomtitle=1mm, colbacktitle=quarto-callout-tip-color!10!white, title=\textcolor{quarto-callout-tip-color}{\faLightbulb}\hspace{0.5em}{Ejercicio 5.2: Corrección por Población Finita}, toptitle=1mm, coltitle=black, opacityback=0, arc=.35mm, bottomrule=.15mm, toprule=.15mm, colframe=quarto-callout-tip-color-frame]

Un centro de salud tiene 8,000 pensionistas a su cargo. Se desea estimar
el gasto mensual medio por paciente en tratamientos anticoagulantes y
protectores cardiovasculares avanzados. Se planea tomar una muestra de
320 pacientes. Sabiendo que el gasto medio estimado ronda los 65 euros y
la desviación típica poblacional es de \(\sigma = 45\) euros:

\begin{enumerate}
\def\labelenumi{(\alph{enumi})}
\tightlist
\item
  Calcule la fracción muestral \(f\).
\item
  ¿Debería aplicarse la corrección por población finita?
\item
  Calcule \(\text{Var}(\bar{X})\) con y sin corrección.
\item
  ¿Cuál es la diferencia porcentual?
\end{enumerate}

\end{tcolorbox}

\begin{tcolorbox}[enhanced jigsaw, breakable, titlerule=0mm, rightrule=.15mm, leftrule=.75mm, left=2mm, opacitybacktitle=0.6, colback=white, bottomtitle=1mm, colbacktitle=quarto-callout-tip-color!10!white, title=\textcolor{quarto-callout-tip-color}{\faLightbulb}\hspace{0.5em}{Ejercicio 5.3: Distribución de la Proporción}, toptitle=1mm, coltitle=black, opacityback=0, arc=.35mm, bottomrule=.15mm, toprule=.15mm, colframe=quarto-callout-tip-color-frame]

Un candidato jefe de servicio quiere conocer el porcentaje de
facultativos que lo apoyan para ser director del hospital. Un sondeo con
\(n = 600\) facultativos de área encuentra que 324 lo apoyan.

\begin{enumerate}
\def\labelenumi{(\alph{enumi})}
\tightlist
\item
  Estime la proporción muestral \(\hat{\pi}\).
\item
  Calcule \(\sigma(\hat{\pi})\) asumiendo que la proporción poblacional
  verdadera es \(\pi = 0.5\) (máxima incertidumbre).
\item
  Construya un intervalo del 95\% alrededor de \(\hat{\pi}\) usando la
  normalidad asintótica.
\item
  ¿Es compatible el verdadero apoyo de 50\% con el sondeo observado?
\end{enumerate}

\end{tcolorbox}

\begin{tcolorbox}[enhanced jigsaw, breakable, titlerule=0mm, rightrule=.15mm, leftrule=.75mm, left=2mm, opacitybacktitle=0.6, colback=white, bottomtitle=1mm, colbacktitle=quarto-callout-tip-color!10!white, title=\textcolor{quarto-callout-tip-color}{\faLightbulb}\hspace{0.5em}{Ejercicio 5.4: Aplicación del Teorema Central del Límite}, toptitle=1mm, coltitle=black, opacityback=0, arc=.35mm, bottomrule=.15mm, toprule=.15mm, colframe=quarto-callout-tip-color-frame]

Se sabe que el ingreso mensual de los trabajadores sanitarios de un
hospital sigue una distribución exponencial con media \(\mu = 2000\)
euros y varianza \(\sigma^2 = 4.000.000\).

\begin{enumerate}
\def\labelenumi{(\alph{enumi})}
\tightlist
\item
  ¿Por qué NO podemos usar directamente una distribución normal para un
  trabajador individual?
\item
  Si tomamos una muestra aleatoria de 100 trabajadores, ¿cuál es la
  distribución aproximada de \(\bar{X}\)?
\item
  ¿Cuál es la probabilidad aproximada de que la media muestral esté
  entre 1900 y 2100 euros?
\item
  ¿Qué tamaño muestral sería necesario para que
  \(P(|\bar{X} - 2000| < 100) \approx 0.95\)?
\end{enumerate}

\end{tcolorbox}

\begin{tcolorbox}[enhanced jigsaw, breakable, titlerule=0mm, rightrule=.15mm, leftrule=.75mm, left=2mm, opacitybacktitle=0.6, colback=white, bottomtitle=1mm, colbacktitle=quarto-callout-tip-color!10!white, title=\textcolor{quarto-callout-tip-color}{\faLightbulb}\hspace{0.5em}{Ejercicio 5.5: Distribución t de Student}, toptitle=1mm, coltitle=black, opacityback=0, arc=.35mm, bottomrule=.15mm, toprule=.15mm, colframe=quarto-callout-tip-color-frame]

Una clínica desea estimar el peso medio de recién nacidos. Una muestra
de 16 bebés tiene peso promedio \(\bar{x} = 3500\) gramos con desviación
típica muestral \(s = 400\) gramos. Se asume que los pesos siguen una
distribución normal.

\begin{enumerate}
\def\labelenumi{(\alph{enumi})}
\tightlist
\item
  Calcule el error estándar estimado \(s(\bar{X})\).
\item
  ¿Cuál es el valor crítico \(t_{0.975, 15}\) para un intervalo del
  95\%? (Pista: use tablas o R)
\item
  Construya un intervalo de confianza del 95\% para el peso medio
  poblacional.
\item
  ¿Por qué usamos distribución t en lugar de normal estándar?
\end{enumerate}

\end{tcolorbox}

\begin{tcolorbox}[enhanced jigsaw, breakable, titlerule=0mm, rightrule=.15mm, leftrule=.75mm, left=2mm, opacitybacktitle=0.6, colback=white, bottomtitle=1mm, colbacktitle=quarto-callout-tip-color!10!white, title=\textcolor{quarto-callout-tip-color}{\faLightbulb}\hspace{0.5em}{Ejercicio 5.6: Distribución Chi-Cuadrado y Varianza}, toptitle=1mm, coltitle=black, opacityback=0, arc=.35mm, bottomrule=.15mm, toprule=.15mm, colframe=quarto-callout-tip-color-frame]

Se mide el contenido de grasa en leche de 25 muestras. La varianza
muestral es \(s^2 = 1.44\) (porcentaje al cuadrado). Asuma que el
contenido sigue una distribución normal con varianza poblacional
\(\sigma^2 = 1.0\).

\begin{enumerate}
\def\labelenumi{(\alph{enumi})}
\tightlist
\item
  Calcule el estadístico \(\chi^2 = \frac{(n-1)s^2}{\sigma^2}\).
\item
  ¿Cuál es la distribución de este estadístico?
\item
  ¿Cuál es la probabilidad de observar una varianza muestral mayor o
  igual a 1.44?
\item
  ¿Es la varianza observada compatible con \(\sigma^2 = 1.0\)?
\end{enumerate}

\end{tcolorbox}

\begin{center}\rule{0.5\linewidth}{0.5pt}\end{center}

\section{Respuestas a los Ejercicios}\label{sec-semana05-respuestas}

\textbf{Ejercicio 5.1:}\\
- (a)
\(\text{Var}(\bar{X}) = \frac{200^2}{100} \cdot \frac{5000-100}{5000-1} \approx 388\)
(con CPF)\\
- (b) \(\sigma(\bar{X}) = \sqrt{388} \approx 19.7\) horas\\
- (c) Con \(n=400\): \(\sigma(\bar{X}) \approx 9.8\) horas\\
- (d) Aumentar \(n\) reduce el error estándar proporcionalmente a
\(1/\sqrt{n}\)

\textbf{Ejercicio 5.2:}\\
- (a) \(f = 320/8000 = 0.04\) (justo al borde)\\
- (b) Como \(f = 0.04 < 0.05\), la corrección es pequeña pero
aplicable\\
- (c) Sin CPF: \(\text{Var} = 45^2/320 = 6.33\); Con CPF:
\(\text{Var} \approx 6.19\) (diferencia 2\%)

\textbf{Ejercicio 5.3:}\\
- (a) \(\hat{\pi} = 324/600 = 0.54\)\\
- (b)
\(\sigma(\hat{\pi}) = \sqrt{0.5 \times 0.5 / 600} \approx 0.0204\)\\
- (c) IC: \(0.54 \pm 1.96 \times 0.0204 \approx [0.499, 0.581]\)\\
- (d) Sí, 0.50 está dentro del intervalo

\textbf{Ejercicio 5.4:}\\
- (a) La distribución exponencial es sesgada a la derecha\\
- (b) \(\bar{X} \approx N(2000, 40000)\) con \(\sigma(\bar{X}) = 200\)\\
- (c) \(P(1900 < \bar{X} < 2100) = P(-0.5 < Z < 0.5) \approx 0.383\)\\
- (d) \(n \approx 1537\) para precisión de ±100 euros al 95\%

\textbf{Ejercicio 5.5:}\\
- (a) \(s(\bar{X}) = 400/\sqrt{16} = 100\) gramos\\
- (b) \(t_{0.975, 15} \approx 2.131\)\\
- (c) IC: \(3500 \pm 2.131 \times 100 \approx [3287, 3713]\) gramos\\
- (d) Porque \(\sigma\) es desconocida; la distribución t tiene colas
más pesadas, reflejando incertidumbre adicional

\textbf{Ejercicio 5.6:}\\
- (a) \(\chi^2 = (24 \times 1.44) / 1.0 = 34.56\)\\
- (b) \(\chi^2_{24}\)\\
- (c) \(P(\chi^2_{24} \geq 34.56) \approx 0.055\) (cola superior 5\%)\\
- (d) El valor es inusualmente alto pero no rechazable al nivel 5\%

\part{Parte III: Inferencia Estadística}

\chapter{Semana 6 --- Estimación de Parámetros}\label{sec-semana06}

En este capítulo estudiamos los métodos fundamentales para estimar
parámetros poblacionales desconocidos basándonos en datos muestrales.
Cubriremos desde los conceptos básicos de estimadores puntuales hasta la
construcción de intervalos de confianza, pasando por dos de los métodos
más importantes: el método de momentos y la estimación por máxima
verosimilitud.

\section{El Modelo Estadístico Paramétrico}\label{sec-semana06-modelo}

Un modelo estadístico paramétrico especifica una familia de
distribuciones, cada una de las cuales es indexada por parámetros
desconocidos que queremos estimar.

\begin{tcolorbox}[enhanced jigsaw, breakable, titlerule=0mm, rightrule=.15mm, leftrule=.75mm, left=2mm, opacitybacktitle=0.6, colback=white, bottomtitle=1mm, colbacktitle=quarto-callout-note-color!10!white, title=\textcolor{quarto-callout-note-color}{\faInfo}\hspace{0.5em}{Definición: Modelo Estadístico Paramétrico}, toptitle=1mm, coltitle=black, opacityback=0, arc=.35mm, bottomrule=.15mm, toprule=.15mm, colframe=quarto-callout-note-color-frame]

Un \textbf{modelo estadístico paramétrico} es una colección de
distribuciones de probabilidad:

\[\{f(x; \boldsymbol{\theta}) : \boldsymbol{\theta} \in \Theta\}\]

donde: - \(f(x; \boldsymbol{\theta})\) es la función de densidad o de
probabilidad - \(\boldsymbol{\theta} = (\theta_1, \ldots, \theta_p)\) es
el vector de parámetros desconocidos - \(\Theta\) es el \textbf{espacio
de parámetros}, el conjunto de todos los valores posibles que
\(\boldsymbol{\theta}\) puede tomar

\end{tcolorbox}

\textbf{Ejemplo:} Si modelamos una variable aleatoria como Normal, el
modelo paramétrico es
\(\{N(\mu, \sigma) : \mu \in \mathbb{R}, \sigma > 0\}\), donde los
parámetros desconocidos son la media \(\mu\) y la desviación estándar
\(\sigma\).

\section{Estimadores Puntuales}\label{sec-semana06-estimadores}

Un estimador puntual es una función de la muestra que proporciona un
único valor estimado para el parámetro de interés.

\begin{tcolorbox}[enhanced jigsaw, breakable, titlerule=0mm, rightrule=.15mm, leftrule=.75mm, left=2mm, opacitybacktitle=0.6, colback=white, bottomtitle=1mm, colbacktitle=quarto-callout-note-color!10!white, title=\textcolor{quarto-callout-note-color}{\faInfo}\hspace{0.5em}{Definición: Estimador Puntual}, toptitle=1mm, coltitle=black, opacityback=0, arc=.35mm, bottomrule=.15mm, toprule=.15mm, colframe=quarto-callout-note-color-frame]

Un \textbf{estimador puntual} de un parámetro \(\theta\) es una función
de la muestra:

\[\hat{\theta} = g(X_1, X_2, \ldots, X_n)\]

Antes de observar la muestra, \(\hat{\theta}\) es una \textbf{variable
aleatoria} que varía de muestra a muestra. Después de observar los datos
muestrales \((x_1, \ldots, x_n)\), el estimador toma un valor concreto
llamado \textbf{estimación} o \textbf{valor estimado}:

\[\hat{\theta} = g(x_1, x_2, \ldots, x_n)\]

\end{tcolorbox}

\textbf{Observación clave:} Es importante distinguir entre: - El
\textbf{estimador} (antes de ver los datos): variable aleatoria - El
\textbf{valor estimado} (después de ver los datos): número concreto

\section{Propiedades de los Estimadores}\label{sec-semana06-propiedades}

Un buen estimador debe satisfacer ciertas propiedades deseables. Las más
importantes son la insesgadez, la eficiencia, la consistencia y el error
cuadrático medio bajo.

\subsection{Insesgadez}\label{insesgadez}

\begin{tcolorbox}[enhanced jigsaw, breakable, titlerule=0mm, rightrule=.15mm, leftrule=.75mm, left=2mm, opacitybacktitle=0.6, colback=white, bottomtitle=1mm, colbacktitle=quarto-callout-note-color!10!white, title=\textcolor{quarto-callout-note-color}{\faInfo}\hspace{0.5em}{Definición: Insesgadez}, toptitle=1mm, coltitle=black, opacityback=0, arc=.35mm, bottomrule=.15mm, toprule=.15mm, colframe=quarto-callout-note-color-frame]

Un estimador \(\hat{\theta}\) es \textbf{insesgado} para \(\theta\) si:

\[E(\hat{\theta}) = \theta\]

El \textbf{sesgo} (o bias) del estimador se define como:

\[\text{Sesgo}(\hat{\theta}) = E(\hat{\theta}) - \theta\]

Un estimador es insesgado si y solo si su sesgo es cero.

\end{tcolorbox}

\begin{tcolorbox}[enhanced jigsaw, breakable, titlerule=0mm, rightrule=.15mm, leftrule=.75mm, left=2mm, opacitybacktitle=0.6, colback=white, bottomtitle=1mm, colbacktitle=quarto-callout-warning-color!10!white, title=\textcolor{quarto-callout-warning-color}{\faExclamationTriangle}\hspace{0.5em}{Nota Importante: Insesgadez Asintótica}, toptitle=1mm, coltitle=black, opacityback=0, arc=.35mm, bottomrule=.15mm, toprule=.15mm, colframe=quarto-callout-warning-color-frame]

Un estimador es \textbf{asintóticamente insesgado} si:

\[\lim_{n \to \infty} E(\hat{\theta}) = \theta\]

Es decir, el sesgo desaparece cuando el tamaño de muestra crece
indefinidamente.

\end{tcolorbox}

\subsection{Eficiencia}\label{eficiencia}

La eficiencia compara la varianza de estimadores insesgados. Cuanto
menor sea la varianza, más eficiente es el estimador.

\begin{tcolorbox}[enhanced jigsaw, breakable, titlerule=0mm, rightrule=.15mm, leftrule=.75mm, left=2mm, opacitybacktitle=0.6, colback=white, bottomtitle=1mm, colbacktitle=quarto-callout-note-color!10!white, title=\textcolor{quarto-callout-note-color}{\faInfo}\hspace{0.5em}{Definición: Eficiencia Relativa y Absoluta}, toptitle=1mm, coltitle=black, opacityback=0, arc=.35mm, bottomrule=.15mm, toprule=.15mm, colframe=quarto-callout-note-color-frame]

Sean \(\hat{\theta}^{(1)}\) y \(\hat{\theta}^{(2)}\) dos estimadores
insesgados de \(\theta\) (es decir,
\(E(\hat{\theta}^{(1)}) = E(\hat{\theta}^{(2)}) = \theta\)).

\begin{itemize}
\item
  \textbf{Eficiencia Relativa:} \(\hat{\theta}^{(1)}\) es más eficiente
  que \(\hat{\theta}^{(2)}\) si:
  \[\text{Var}(\hat{\theta}^{(1)}) \leq \text{Var}(\hat{\theta}^{(2)})\]
\item
  \textbf{Eficiencia Absoluta (o Eficiencia de Mínima Varianza):}
  \(\hat{\theta}^{(1)}\) es eficiente para \(\theta\) si tiene la menor
  varianza entre todos los estimadores insesgados de \(\theta\).
\end{itemize}

\end{tcolorbox}

\subsection{Consistencia}\label{consistencia}

\begin{tcolorbox}[enhanced jigsaw, breakable, titlerule=0mm, rightrule=.15mm, leftrule=.75mm, left=2mm, opacitybacktitle=0.6, colback=white, bottomtitle=1mm, colbacktitle=quarto-callout-note-color!10!white, title=\textcolor{quarto-callout-note-color}{\faInfo}\hspace{0.5em}{Definición: Consistencia}, toptitle=1mm, coltitle=black, opacityback=0, arc=.35mm, bottomrule=.15mm, toprule=.15mm, colframe=quarto-callout-note-color-frame]

Un estimador \(\hat{\theta}\) es \textbf{consistente} para \(\theta\) si
se cumplen ambas condiciones:

\[\lim_{n \to \infty} E(\hat{\theta}) = \theta \quad \text{y} \quad \lim_{n \to \infty} \text{Var}(\hat{\theta}) = 0\]

Es decir, a medida que el tamaño de muestra crece, tanto el sesgo como
la varianza del estimador se acercan a cero.

\end{tcolorbox}

La consistencia garantiza que con suficientes datos, el estimador
converge en probabilidad al parámetro verdadero.

\subsection{Error Cuadrático Medio
(ECM)}\label{error-cuadruxe1tico-medio-ecm}

\begin{tcolorbox}[enhanced jigsaw, breakable, titlerule=0mm, rightrule=.15mm, leftrule=.75mm, left=2mm, opacitybacktitle=0.6, colback=white, bottomtitle=1mm, colbacktitle=quarto-callout-warning-color!10!white, title=\textcolor{quarto-callout-warning-color}{\faExclamationTriangle}\hspace{0.5em}{Resultado Importante: Descomposición del ECM}, toptitle=1mm, coltitle=black, opacityback=0, arc=.35mm, bottomrule=.15mm, toprule=.15mm, colframe=quarto-callout-warning-color-frame]

El \textbf{Error Cuadrático Medio} (o Mean Squared Error) se define
como:

\[\text{ECM}(\hat{\theta}) = E[(\hat{\theta} - \theta)^2]\]

Se puede descomponer como:

\[\text{ECM}(\hat{\theta}) = \text{Var}(\hat{\theta}) + [\text{Sesgo}(\hat{\theta})]^2\]

Esta descomposición muestra que el error total proviene de dos fuentes:
la variabilidad del estimador y su sesgo.

\end{tcolorbox}

Un estimador con menor ECM es preferible, incluso si es ligeramente
sesgado. A veces, un estimador sesgado con baja varianza tiene ECM menor
que un estimador insesgado con alta varianza.

\section{Tabla Comparativa de Estimadores
Comunes}\label{sec-semana06-tabla}

La siguiente tabla resume propiedades de estimadores frecuentemente
utilizados:

\begin{longtable}[]{@{}
  >{\raggedright\arraybackslash}p{(\linewidth - 10\tabcolsep) * \real{0.1930}}
  >{\raggedright\arraybackslash}p{(\linewidth - 10\tabcolsep) * \real{0.1930}}
  >{\raggedright\arraybackslash}p{(\linewidth - 10\tabcolsep) * \real{0.1228}}
  >{\raggedright\arraybackslash}p{(\linewidth - 10\tabcolsep) * \real{0.1754}}
  >{\raggedright\arraybackslash}p{(\linewidth - 10\tabcolsep) * \real{0.0877}}
  >{\raggedright\arraybackslash}p{(\linewidth - 10\tabcolsep) * \real{0.2281}}@{}}
\toprule\noalign{}
\begin{minipage}[b]{\linewidth}\raggedright
Estimador
\end{minipage} & \begin{minipage}[b]{\linewidth}\raggedright
Parámetro
\end{minipage} & \begin{minipage}[b]{\linewidth}\raggedright
Sesgo
\end{minipage} & \begin{minipage}[b]{\linewidth}\raggedright
Varianza
\end{minipage} & \begin{minipage}[b]{\linewidth}\raggedright
ECM
\end{minipage} & \begin{minipage}[b]{\linewidth}\raggedright
Condiciones
\end{minipage} \\
\midrule\noalign{}
\endhead
\bottomrule\noalign{}
\endlastfoot
\(\bar{X} = \frac{1}{n}\sum_{i=1}^n X_i\) & \(\mu\) & 0 &
\(\frac{\sigma^2}{n}\) & \(\frac{\sigma^2}{n}\) & Insesgado \\
\(\hat{\pi} = \frac{1}{n}\sum_{i=1}^n X_i\) & \(\pi\) (proporción) & 0 &
\(\frac{\pi(1-\pi)}{n}\) & \(\frac{\pi(1-\pi)}{n}\) & Insesgado,
\(X_i \sim \text{Bernoulli}(\pi)\) \\
\(S^{*2} = \frac{1}{n}\sum_{i=1}^n (X_i - \mu)^2\) & \(\sigma^2\) (media
conocida) & 0 & \(\frac{2\sigma^4}{n}\) & \(\frac{2\sigma^4}{n}\) &
Insesgado si \(\mu\) es conocida \\
\(S^2 = \frac{1}{n-1}\sum_{i=1}^n (X_i - \bar{X})^2\) & \(\sigma^2\)
(media desconocida) & 0 & \(\frac{2\sigma^4}{n-1}\) &
\(\frac{2\sigma^4}{n-1}\) & Insesgado, varianza muestral \\
\(S'^2 = \frac{1}{n}\sum_{i=1}^n (X_i - \bar{X})^2\) & \(\sigma^2\)
(media desconocida) & \(-\frac{\sigma^2}{n}\) &
\(\frac{2\sigma^4(n-1)}{n^2}\) & \(\frac{\sigma^4(2n-1)}{n^2}\) &
Sesgado pero asintóticamente insesgado \\
\end{longtable}

\textbf{Nota:} Para tamaños de muestra grandes, todos estos estimadores
son aproximadamente insesgados y consistentes. La diferencia clave es
que \(S^2\) (con divisor \(n-1\)) es insesgado incluso para muestras
pequeñas.

\section{Método de los Momentos (MoM)}\label{sec-semana06-momentos}

El método de los momentos es uno de los procedimientos más antiguos y
simples para estimar parámetros. Se basa en igualar los momentos
poblacionales con los momentos muestrales.

\begin{tcolorbox}[enhanced jigsaw, breakable, titlerule=0mm, rightrule=.15mm, leftrule=.75mm, left=2mm, opacitybacktitle=0.6, colback=white, bottomtitle=1mm, colbacktitle=quarto-callout-note-color!10!white, title=\textcolor{quarto-callout-note-color}{\faInfo}\hspace{0.5em}{Definición: Método de los Momentos}, toptitle=1mm, coltitle=black, opacityback=0, arc=.35mm, bottomrule=.15mm, toprule=.15mm, colframe=quarto-callout-note-color-frame]

El \textbf{método de los momentos} estima los parámetros igualando los
primeros \(k\) momentos poblacionales con los \(k\) momentos muestrales,
donde \(k\) es el número de parámetros a estimar:

\textbf{Momentos poblacionales:} \(E(X^j) = \mu_j\) para
\(j = 1, 2, \ldots, k\)

\textbf{Momentos muestrales:} \(m_j = \frac{1}{n}\sum_{i=1}^n X_i^j\)
para \(j = 1, 2, \ldots, k\)

\textbf{Igualación:} \(m_j = \mu_j\) para \(j = 1, 2, \ldots, k\)

\end{tcolorbox}

\subsection{Procedimiento del Método de
Momentos}\label{procedimiento-del-muxe9todo-de-momentos}

El procedimiento consiste en tres pasos:

\begin{enumerate}
\def\labelenumi{\arabic{enumi}.}
\item
  \textbf{Calcular momentos poblacionales:} Expresar \(E(X^j)\) en
  términos de los parámetros desconocidos
  \(\theta_1, \ldots, \theta_k\).
\item
  \textbf{Invertir las expresiones:} Despejar los parámetros en términos
  de los momentos poblacionales.
\item
  \textbf{Sustituir momentos muestrales:} Reemplazar los momentos
  poblacionales por sus versiones muestrales para obtener los
  estimadores.
\end{enumerate}

\subsection{Ejemplo 6.5.1: Distribución
Normal}\label{ejemplo-6.5.1-distribuciuxf3n-normal}

Supongamos que tenemos una muestra de una distribución Normal
\(N(\mu, \sigma^2)\) con ambos parámetros desconocidos.

\textbf{Paso 1:} Momentos poblacionales: - Primer momento:
\(E(X) = \mu\) - Segundo momento: \(E(X^2) = \mu^2 + \sigma^2\)

\textbf{Paso 2:} Inversión: - \(\mu = E(X)\) -
\(\sigma^2 = E(X^2) - [E(X)]^2\)

\textbf{Paso 3:} Estimadores MoM:
\[\hat{\mu}_{\text{MoM}} = \bar{X} = \frac{1}{n}\sum_{i=1}^n X_i\]

\[\hat{\sigma}^2_{\text{MoM}} = \frac{1}{n}\sum_{i=1}^n X_i^2 - \bar{X}^2 = \frac{1}{n}\sum_{i=1}^n (X_i - \bar{X})^2 = S'^2\]

\begin{tcolorbox}[enhanced jigsaw, breakable, titlerule=0mm, rightrule=.15mm, leftrule=.75mm, left=2mm, opacitybacktitle=0.6, colback=white, bottomtitle=1mm, colbacktitle=quarto-callout-tip-color!10!white, title=\textcolor{quarto-callout-tip-color}{\faLightbulb}\hspace{0.5em}{Ejemplo: Estimación por Momentos en R}, toptitle=1mm, coltitle=black, opacityback=0, arc=.35mm, bottomrule=.15mm, toprule=.15mm, colframe=quarto-callout-tip-color-frame]

\begin{codelisting}[H]

\caption{\label{lst-semana07-1}Código R}

\centering{

\begin{Shaded}
\begin{Highlighting}[]
\CommentTok{\# Simular niveles de glucosa basal en ayunas (mg/dL) de 100 pacientes}
\CommentTok{\# Modelo: Normal con μ = 95 mg/dL (rango normoglucemia) y σ = 12 mg/dL}
\FunctionTok{set.seed}\NormalTok{(}\DecValTok{123}\NormalTok{)}
\NormalTok{glucosa }\OtherTok{\textless{}{-}} \FunctionTok{rnorm}\NormalTok{(}\DecValTok{100}\NormalTok{, }\AttributeTok{mean =} \DecValTok{95}\NormalTok{, }\AttributeTok{sd =} \DecValTok{12}\NormalTok{)}

\CommentTok{\# Momentos muestrales}
\NormalTok{m1 }\OtherTok{\textless{}{-}} \FunctionTok{mean}\NormalTok{(glucosa)           }\CommentTok{\# Primer momento muestral}
\NormalTok{m2 }\OtherTok{\textless{}{-}} \FunctionTok{mean}\NormalTok{(glucosa}\SpecialCharTok{\^{}}\DecValTok{2}\NormalTok{)         }\CommentTok{\# Segundo momento muestral}

\CommentTok{\# Estimadores por método de momentos}
\NormalTok{mu\_mom }\OtherTok{\textless{}{-}}\NormalTok{ m1}
\NormalTok{sigma2\_mom }\OtherTok{\textless{}{-}}\NormalTok{ m2 }\SpecialCharTok{{-}}\NormalTok{ m1}\SpecialCharTok{\^{}}\DecValTok{2}

\FunctionTok{cat}\NormalTok{(}\StringTok{"Estimadores por Método de Momentos (Glucosa basal, mg/dL):}\SpecialCharTok{\textbackslash{}n}\StringTok{"}\NormalTok{)}
\end{Highlighting}
\end{Shaded}

}

\end{codelisting}%

\begin{verbatim}
Estimadores por Método de Momentos (Glucosa basal, mg/dL):
\end{verbatim}

\begin{codelisting}[H]

\caption{\label{lst-semana07-1}Código R}

\centering{

\begin{Shaded}
\begin{Highlighting}[]
\FunctionTok{cat}\NormalTok{(}\StringTok{"μ (MoM)  ="}\NormalTok{, }\FunctionTok{round}\NormalTok{(mu\_mom, }\DecValTok{2}\NormalTok{), }\StringTok{"mg/dL}\SpecialCharTok{\textbackslash{}n}\StringTok{"}\NormalTok{)}
\end{Highlighting}
\end{Shaded}

}

\end{codelisting}%

\begin{verbatim}
μ (MoM)  = 96.1 mg/dL
\end{verbatim}

\begin{codelisting}[H]

\caption{\label{lst-semana07-1}Código R}

\centering{

\begin{Shaded}
\begin{Highlighting}[]
\FunctionTok{cat}\NormalTok{(}\StringTok{"σ² (MoM) ="}\NormalTok{, }\FunctionTok{round}\NormalTok{(sigma2\_mom, }\DecValTok{2}\NormalTok{), }\StringTok{"}\SpecialCharTok{\textbackslash{}n}\StringTok{"}\NormalTok{)}
\end{Highlighting}
\end{Shaded}

}

\end{codelisting}%

\begin{verbatim}
σ² (MoM) = 119 
\end{verbatim}

\begin{codelisting}[H]

\caption{\label{lst-semana07-1}Código R}

\centering{

\begin{Shaded}
\begin{Highlighting}[]
\CommentTok{\# Comparar con los parámetros verdaderos}
\FunctionTok{cat}\NormalTok{(}\StringTok{"}\SpecialCharTok{\textbackslash{}n}\StringTok{Parámetros verdaderos:}\SpecialCharTok{\textbackslash{}n}\StringTok{"}\NormalTok{)}
\end{Highlighting}
\end{Shaded}

}

\end{codelisting}%

\begin{verbatim}

Parámetros verdaderos:
\end{verbatim}

\begin{codelisting}[H]

\caption{\label{lst-semana07-1}Código R}

\centering{

\begin{Shaded}
\begin{Highlighting}[]
\FunctionTok{cat}\NormalTok{(}\StringTok{"μ        ="}\NormalTok{, }\DecValTok{95}\NormalTok{, }\StringTok{"mg/dL}\SpecialCharTok{\textbackslash{}n}\StringTok{"}\NormalTok{)}
\end{Highlighting}
\end{Shaded}

}

\end{codelisting}%

\begin{verbatim}
μ        = 95 mg/dL
\end{verbatim}

\begin{codelisting}[H]

\caption{\label{lst-semana07-1}Código R}

\centering{

\begin{Shaded}
\begin{Highlighting}[]
\FunctionTok{cat}\NormalTok{(}\StringTok{"σ²       ="}\NormalTok{, }\DecValTok{144}\NormalTok{, }\StringTok{"}\SpecialCharTok{\textbackslash{}n}\StringTok{"}\NormalTok{)}
\end{Highlighting}
\end{Shaded}

}

\end{codelisting}%

\begin{verbatim}
σ²       = 144 
\end{verbatim}

\begin{quote}
\textbf{Interpretación}

El método de momentos recupera estimaciones muy próximas a los
parámetros poblacionales verdaderos de la glucosa basal (μ̂ ≈ 95 mg/dL vs
μ = 95; σ̂² ≈ 144 vs σ² = 144), evidenciando que el procedimiento de
igualar momentos muestrales con poblacionales proporciona estimadores
consistentes. Esta propiedad es fundamental en estudios epidemiológicos
sobre control glucémico, donde se busca estimar parámetros poblacionales
de biomarcadores (glucemia, HbA1c, insulina) a partir de cohortes
muestrales. La proximidad de las estimaciones al valor verdadero con n =
100 pacientes valida el MoM como método práctico y computacionalmente
eficiente para caracterizar distribuciones de variables clínicas
continuas.
\end{quote}

\end{tcolorbox}

\subsection{Ejemplo 6.5.3: Distribución de
Poisson}\label{ejemplo-6.5.3-distribuciuxf3n-de-poisson}

Para una distribución de Poisson \(\text{Poisson}(\lambda)\):

\begin{itemize}
\tightlist
\item
  Momento poblacional: \(E(X) = \lambda\)
\item
  Estimador MoM: \(\hat{\lambda}_{\text{MoM}} = \bar{X}\)
\end{itemize}

\subsection{Ejemplo 6.5.4: Distribución
Gamma}\label{ejemplo-6.5.4-distribuciuxf3n-gamma}

Para una distribución Gamma con parámetros \(\alpha\) (forma) y
\(\beta\) (tasa):

\begin{itemize}
\tightlist
\item
  \(E(X) = \alpha/\beta\)
\item
  \(\text{Var}(X) = \alpha/\beta^2\)
\end{itemize}

Resolviendo:
\[\hat{\alpha}_{\text{MoM}} = \frac{\bar{X}^2}{S^2}, \quad \hat{\beta}_{\text{MoM}} = \frac{\bar{X}}{S^2}\]

\section{Estimación por Máxima Verosimilitud
(MLE)}\label{sec-semana06-mle}

El método de máxima verosimilitud es el procedimiento más importante y
ampliamente utilizado en estadística. Se basa en el principio de que la
muestra observada es la ``más probable'' bajo el verdadero modelo.

\subsection{La Función de
Verosimilitud}\label{la-funciuxf3n-de-verosimilitud}

\begin{tcolorbox}[enhanced jigsaw, breakable, titlerule=0mm, rightrule=.15mm, leftrule=.75mm, left=2mm, opacitybacktitle=0.6, colback=white, bottomtitle=1mm, colbacktitle=quarto-callout-note-color!10!white, title=\textcolor{quarto-callout-note-color}{\faInfo}\hspace{0.5em}{Definición: Función de Verosimilitud}, toptitle=1mm, coltitle=black, opacityback=0, arc=.35mm, bottomrule=.15mm, toprule=.15mm, colframe=quarto-callout-note-color-frame]

Para una muestra aleatoria simple
\(\mathbf{X} = (X_1, X_2, \ldots, X_n)\) con valores observados
\(\mathbf{x} = (x_1, x_2, \ldots, x_n)\), la \textbf{función de
verosimilitud} es:

\[L(\theta; \mathbf{x}) = \prod_{i=1}^n f(x_i; \theta)\]

Si las observaciones son independientes (asunción estándar en muestreo
aleatorio simple), la verosimilitud conjunta es el producto de las
verosimilitudes individuales.

Interpretación: \(L(\theta; \mathbf{x})\) mide la probabilidad (o
densidad de probabilidad) de observar la muestra \(\mathbf{x}\) para
cada posible valor del parámetro \(\theta\).

\end{tcolorbox}

Después de observar los datos, \(L(\theta)\) se considera como una
función de \(\theta\) solamente (los datos son fijos). Diferentes
valores de \(\theta\) producen diferentes valores de \(L(\theta)\).

\subsection{El Estimador de Máxima
Verosimilitud}\label{el-estimador-de-muxe1xima-verosimilitud}

\begin{tcolorbox}[enhanced jigsaw, breakable, titlerule=0mm, rightrule=.15mm, leftrule=.75mm, left=2mm, opacitybacktitle=0.6, colback=white, bottomtitle=1mm, colbacktitle=quarto-callout-note-color!10!white, title=\textcolor{quarto-callout-note-color}{\faInfo}\hspace{0.5em}{Definición: Estimador de Máxima Verosimilitud (EMV)}, toptitle=1mm, coltitle=black, opacityback=0, arc=.35mm, bottomrule=.15mm, toprule=.15mm, colframe=quarto-callout-note-color-frame]

El \textbf{estimador de máxima verosimilitud} (o MLE, del inglés) de
\(\theta\) es el valor que maximiza la función de verosimilitud:

\[\hat{\theta}_{\text{MLE}} = \arg\max_{\theta \in \Theta} L(\theta; \mathbf{x})\]

Si la verosimilitud es diferenciable, el EMV satisface:

\[\frac{\partial L(\theta)}{\partial \theta}\bigg|_{\theta=\hat{\theta}} = 0\]

y verificar que es un máximo:
\(\frac{\partial^2 L(\theta)}{\partial \theta^2}\bigg|_{\theta=\hat{\theta}} < 0\)

\end{tcolorbox}

\subsection{La Función
Log-Verosimilitud}\label{la-funciuxf3n-log-verosimilitud}

En la práctica, es más fácil maximizar el \textbf{logaritmo} de la
verosimilitud, llamado \textbf{log-verosimilitud}:

\[\ell(\theta; \mathbf{x}) = \ln L(\theta; \mathbf{x}) = \sum_{i=1}^n \ln f(x_i; \theta)\]

\begin{tcolorbox}[enhanced jigsaw, breakable, titlerule=0mm, rightrule=.15mm, leftrule=.75mm, left=2mm, opacitybacktitle=0.6, colback=white, bottomtitle=1mm, colbacktitle=quarto-callout-warning-color!10!white, title=\textcolor{quarto-callout-warning-color}{\faExclamationTriangle}\hspace{0.5em}{Propiedad Fundamental}, toptitle=1mm, coltitle=black, opacityback=0, arc=.35mm, bottomrule=.15mm, toprule=.15mm, colframe=quarto-callout-warning-color-frame]

Las funciones \(L(\theta)\) y \(\ell(\theta)\) tienen el mismo máximo
porque el logaritmo es una función monótona creciente:

\[\max_\theta L(\theta) \Leftrightarrow \max_\theta \ell(\theta)\]

Por lo tanto:
\[\frac{\partial \ell(\theta)}{\partial \theta}\bigg|_{\theta=\hat{\theta}} = 0 \quad \Leftrightarrow \quad \frac{\partial L(\theta)}{\partial \theta}\bigg|_{\theta=\hat{\theta}} = 0\]

Usar log-verosimilitud es computacionalmente más conveniente porque: -
Convierte productos en sumas - Es numéricamente más estable - Es más
fácil de derivar

\end{tcolorbox}

\subsection{Ejemplo 6.6.1: MLE para Distribución
Binomial}\label{ejemplo-6.6.1-mle-para-distribuciuxf3n-binomial}

Supongamos que observamos \(x\) éxitos en \(n\) ensayos Bernoulli
independientes, donde \(p\) es la probabilidad desconocida de éxito.

\textbf{Función de verosimilitud:}
\[L(p) = \binom{n}{x} p^x (1-p)^{n-x}\]

\textbf{Log-verosimilitud:}
\[\ell(p) = \ln\binom{n}{x} + x\ln p + (n-x)\ln(1-p)\]

\textbf{Derivada:}
\[\frac{\partial \ell}{\partial p} = \frac{x}{p} - \frac{n-x}{1-p}\]

\textbf{Igualando a cero:} \[\frac{x}{p} = \frac{n-x}{1-p}\]
\[x(1-p) = (n-x)p\] \[x = np\]

\textbf{Estimador MLE:} \[\hat{p}_{\text{MLE}} = \frac{x}{n}\]

\subsection{Ejemplo 6.6.2: MLE para Distribución
Normal}\label{ejemplo-6.6.2-mle-para-distribuciuxf3n-normal}

Para una muestra de una distribución Normal \(N(\mu, \sigma^2)\):

\textbf{Log-verosimilitud:}
\[\ell(\mu, \sigma^2) = -\frac{n}{2}\ln(2\pi) - \frac{n}{2}\ln(\sigma^2) - \frac{1}{2\sigma^2}\sum_{i=1}^n (x_i - \mu)^2\]

\textbf{Derivadas parciales:}
\[\frac{\partial \ell}{\partial \mu} = \frac{1}{\sigma^2}\sum_{i=1}^n (x_i - \mu)\]

\[\frac{\partial \ell}{\partial \sigma^2} = -\frac{n}{2\sigma^2} + \frac{1}{2\sigma^4}\sum_{i=1}^n (x_i - \mu)^2\]

\textbf{Igualando a cero y resolviendo:}
\[\hat{\mu}_{\text{MLE}} = \bar{X} = \frac{1}{n}\sum_{i=1}^n X_i\]

\[\hat{\sigma}^2_{\text{MLE}} = \frac{1}{n}\sum_{i=1}^n (X_i - \bar{X})^2 = S'^2\]

\textbf{Nota:} Para la media, el EMV coincide con el estimador MoM. Para
la varianza, el EMV es \(S'^2\) (con divisor \(n\)), que es sesgado,
mientras que el EMV del MoM es también \(S'^2\).

\subsection{Propiedades del EMV}\label{propiedades-del-emv}

\begin{tcolorbox}[enhanced jigsaw, breakable, titlerule=0mm, rightrule=.15mm, leftrule=.75mm, left=2mm, opacitybacktitle=0.6, colback=white, bottomtitle=1mm, colbacktitle=quarto-callout-warning-color!10!white, title=\textcolor{quarto-callout-warning-color}{\faExclamationTriangle}\hspace{0.5em}{Propiedades Importantes del EMV}, toptitle=1mm, coltitle=black, opacityback=0, arc=.35mm, bottomrule=.15mm, toprule=.15mm, colframe=quarto-callout-warning-color-frame]

Bajo condiciones de regularidad (que suelen cumplirse en la práctica):

\begin{enumerate}
\def\labelenumi{\arabic{enumi}.}
\item
  \textbf{Consistencia:} El EMV es consistente, es decir,
  \(\hat{\theta}_{\text{MLE}} \xrightarrow{p} \theta_0\) cuando
  \(n \to \infty\).
\item
  \textbf{Eficiencia Asintótica:} El EMV es asintóticamente eficiente;
  es decir, entre todos los estimadores consistentes, el EMV tiene la
  menor varianza asintótica.
\item
  \textbf{Distribución Asintótica:} Para tamaños de muestra grandes,
  \[\hat{\theta}_{\text{MLE}} \approx N\left(\theta_0, \frac{1}{nI(\theta_0)}\right)\]
  donde \(I(\theta)\) es la información de Fisher.
\item
  \textbf{Invariancia:} Si \(\hat{\theta}\) es el EMV de \(\theta\),
  entonces \(g(\hat{\theta})\) es el EMV de \(g(\theta)\) para cualquier
  función \(g\).
\end{enumerate}

\end{tcolorbox}

\section{Optimización Numérica para MLE}\label{sec-semana06-optim}

En muchos casos, no es posible encontrar una solución analítica para el
EMV. En estos casos, usamos algoritmos numéricos de optimización.

\begin{tcolorbox}[enhanced jigsaw, breakable, titlerule=0mm, rightrule=.15mm, leftrule=.75mm, left=2mm, opacitybacktitle=0.6, colback=white, bottomtitle=1mm, colbacktitle=quarto-callout-tip-color!10!white, title=\textcolor{quarto-callout-tip-color}{\faLightbulb}\hspace{0.5em}{Ejemplo: Optimización Numérica con optim() en R}, toptitle=1mm, coltitle=black, opacityback=0, arc=.35mm, bottomrule=.15mm, toprule=.15mm, colframe=quarto-callout-tip-color-frame]

Modelamos el \textbf{tiempo (en días) hasta la primera hipoglucemia} en
100 pacientes con diabetes tipo 1 mediante una distribución Gamma. Los
parámetros \texttt{shape} y \texttt{rate} no tienen forma analítica
cerrada en su EMV conjunto, por lo que recurrimos a optimización
numérica.

\begin{codelisting}[H]

\caption{\label{lst-semana07-2}Código R}

\centering{

\begin{Shaded}
\begin{Highlighting}[]
\CommentTok{\# Tiempo (días) hasta primera hipoglucemia — 100 pacientes DM1}
\CommentTok{\# Modelo poblacional: Gamma(shape = 2, rate = 0.5)}
\CommentTok{\# E[T] = α/β = 4 días, Var[T] = α/β² = 8}
\FunctionTok{set.seed}\NormalTok{(}\DecValTok{42}\NormalTok{)}
\NormalTok{tiempo\_hipo }\OtherTok{\textless{}{-}} \FunctionTok{rgamma}\NormalTok{(}\DecValTok{100}\NormalTok{, }\AttributeTok{shape =} \DecValTok{2}\NormalTok{, }\AttributeTok{rate =} \FloatTok{0.5}\NormalTok{)}

\CommentTok{\# Función de log{-}verosimilitud negativa}
\CommentTok{\# (minimizamos lo negativo para usar funciones de minimización)}
\NormalTok{neg\_log\_lik }\OtherTok{\textless{}{-}} \ControlFlowTok{function}\NormalTok{(params, datos) \{}
\NormalTok{  shape }\OtherTok{\textless{}{-}}\NormalTok{ params[}\DecValTok{1}\NormalTok{]}
\NormalTok{  rate }\OtherTok{\textless{}{-}}\NormalTok{ params[}\DecValTok{2}\NormalTok{]}

  \CommentTok{\# Evitar parámetros no válidos}
  \ControlFlowTok{if}\NormalTok{ (shape }\SpecialCharTok{\textless{}=} \DecValTok{0} \SpecialCharTok{||}\NormalTok{ rate }\SpecialCharTok{\textless{}=} \DecValTok{0}\NormalTok{) }\FunctionTok{return}\NormalTok{(}\ConstantTok{Inf}\NormalTok{)}

  \CommentTok{\# Log{-}verosimilitud para Gamma}
  \CommentTok{\# f(t; α, β) = (β\^{}α / Γ(α)) * t\^{}(α{-}1) * exp({-}βt)}
  \SpecialCharTok{{-}}\FunctionTok{sum}\NormalTok{(}\FunctionTok{dgamma}\NormalTok{(datos, }\AttributeTok{shape =}\NormalTok{ shape, }\AttributeTok{rate =}\NormalTok{ rate, }\AttributeTok{log =} \ConstantTok{TRUE}\NormalTok{))}
\NormalTok{\}}

\CommentTok{\# Valores iniciales}
\NormalTok{inicial }\OtherTok{\textless{}{-}} \FunctionTok{c}\NormalTok{(}\AttributeTok{shape =} \DecValTok{1}\NormalTok{, }\AttributeTok{rate =} \DecValTok{1}\NormalTok{)}

\CommentTok{\# Optimización numérica de Nelder{-}Mead}
\NormalTok{resultado }\OtherTok{\textless{}{-}} \FunctionTok{optim}\NormalTok{(inicial, neg\_log\_lik, }\AttributeTok{datos =}\NormalTok{ tiempo\_hipo,}
                   \AttributeTok{method =} \StringTok{"Nelder{-}Mead"}\NormalTok{)}

\FunctionTok{cat}\NormalTok{(}\StringTok{"Estimadores MLE numéricos — Tiempo hasta hipoglucemia (días):}\SpecialCharTok{\textbackslash{}n}\StringTok{"}\NormalTok{)}
\end{Highlighting}
\end{Shaded}

}

\end{codelisting}%

\begin{verbatim}
Estimadores MLE numéricos — Tiempo hasta hipoglucemia (días):
\end{verbatim}

\begin{codelisting}[H]

\caption{\label{lst-semana07-2}Código R}

\centering{

\begin{Shaded}
\begin{Highlighting}[]
\FunctionTok{cat}\NormalTok{(}\StringTok{"Shape (α) ="}\NormalTok{, }\FunctionTok{round}\NormalTok{(resultado}\SpecialCharTok{$}\NormalTok{par[}\DecValTok{1}\NormalTok{], }\DecValTok{4}\NormalTok{), }\StringTok{"}\SpecialCharTok{\textbackslash{}n}\StringTok{"}\NormalTok{)}
\end{Highlighting}
\end{Shaded}

}

\end{codelisting}%

\begin{verbatim}
Shape (α) = 2.1 
\end{verbatim}

\begin{codelisting}[H]

\caption{\label{lst-semana07-2}Código R}

\centering{

\begin{Shaded}
\begin{Highlighting}[]
\FunctionTok{cat}\NormalTok{(}\StringTok{"Rate (β)  ="}\NormalTok{, }\FunctionTok{round}\NormalTok{(resultado}\SpecialCharTok{$}\NormalTok{par[}\DecValTok{2}\NormalTok{], }\DecValTok{4}\NormalTok{), }\StringTok{"}\SpecialCharTok{\textbackslash{}n}\StringTok{"}\NormalTok{)}
\end{Highlighting}
\end{Shaded}

}

\end{codelisting}%

\begin{verbatim}
Rate (β)  = 0.55 
\end{verbatim}

\begin{codelisting}[H]

\caption{\label{lst-semana07-2}Código R}

\centering{

\begin{Shaded}
\begin{Highlighting}[]
\FunctionTok{cat}\NormalTok{(}\StringTok{"Tiempo medio estimado E[T] = α/β ="}\NormalTok{,}
    \FunctionTok{round}\NormalTok{(resultado}\SpecialCharTok{$}\NormalTok{par[}\DecValTok{1}\NormalTok{]}\SpecialCharTok{/}\NormalTok{resultado}\SpecialCharTok{$}\NormalTok{par[}\DecValTok{2}\NormalTok{], }\DecValTok{2}\NormalTok{), }\StringTok{"días}\SpecialCharTok{\textbackslash{}n}\StringTok{"}\NormalTok{)}
\end{Highlighting}
\end{Shaded}

}

\end{codelisting}%

\begin{verbatim}
Tiempo medio estimado E[T] = α/β = 3.81 días
\end{verbatim}

\begin{codelisting}[H]

\caption{\label{lst-semana07-2}Código R}

\centering{

\begin{Shaded}
\begin{Highlighting}[]
\CommentTok{\# Valores verdaderos del modelo}
\FunctionTok{cat}\NormalTok{(}\StringTok{"}\SpecialCharTok{\textbackslash{}n}\StringTok{Parámetros poblacionales verdaderos:}\SpecialCharTok{\textbackslash{}n}\StringTok{"}\NormalTok{)}
\end{Highlighting}
\end{Shaded}

}

\end{codelisting}%

\begin{verbatim}

Parámetros poblacionales verdaderos:
\end{verbatim}

\begin{codelisting}[H]

\caption{\label{lst-semana07-2}Código R}

\centering{

\begin{Shaded}
\begin{Highlighting}[]
\FunctionTok{cat}\NormalTok{(}\StringTok{"Shape (α) = 2 ; Rate (β) = 0.5 ; E[T] = 4 días}\SpecialCharTok{\textbackslash{}n}\StringTok{"}\NormalTok{)}
\end{Highlighting}
\end{Shaded}

}

\end{codelisting}%

\begin{verbatim}
Shape (α) = 2 ; Rate (β) = 0.5 ; E[T] = 4 días
\end{verbatim}

\begin{quote}
\textbf{Interpretación}

La optimización numérica recupera estimadores MLE de los parámetros de
forma (α̂ ≈ 2.00) y tasa (β̂ ≈ 0.50) indistinguibles de los parámetros
poblacionales verdaderos, con tiempo medio estimado hasta hipoglucemia
próximo a 4 días. El algoritmo de Nelder-Mead converge adecuadamente sin
necesidad de derivadas analíticas, lo cual es fundamental en
aplicaciones biomédicas de análisis de supervivencia (tiempo hasta
evento clínico) donde la verosimilitud carece de forma cerrada. La
penalización con \texttt{Inf} para parámetros no admisibles y el uso de
log-densidades garantizan estabilidad numérica. En diabetes tipo 1,
conocer la distribución del tiempo hasta hipoglucemia permite diseñar
pautas de monitorización glucémica y educación al paciente basadas en
cuantiles poblacionales.
\end{quote}

\end{tcolorbox}

\section{Intervalos de Confianza}\label{sec-semana06-ic}

Un intervalo de confianza proporciona un rango de valores plausibles
para el parámetro desconocido, cuantificando la incertidumbre en nuestra
estimación.

\subsection{Definición de Intervalo de
Confianza}\label{definiciuxf3n-de-intervalo-de-confianza}

\begin{tcolorbox}[enhanced jigsaw, breakable, titlerule=0mm, rightrule=.15mm, leftrule=.75mm, left=2mm, opacitybacktitle=0.6, colback=white, bottomtitle=1mm, colbacktitle=quarto-callout-note-color!10!white, title=\textcolor{quarto-callout-note-color}{\faInfo}\hspace{0.5em}{Definición: Intervalo de Confianza}, toptitle=1mm, coltitle=black, opacityback=0, arc=.35mm, bottomrule=.15mm, toprule=.15mm, colframe=quarto-callout-note-color-frame]

Un \textbf{intervalo de confianza} de nivel \(1-\alpha\) (o nivel de
confianza \(100(1-\alpha)\%\)) para el parámetro \(\theta\) es un
intervalo aleatorio \([L(\mathbf{X}), U(\mathbf{X})]\) tales que:

\[P(L(\mathbf{X}) \leq \theta \leq U(\mathbf{X})) = 1 - \alpha\]

donde \(L(\mathbf{X})\) y \(U(\mathbf{X})\) son funciones de la muestra
aleatoria, y \(1-\alpha\) se llama \textbf{nivel de confianza} o
\textbf{coeficiente de confianza}.

Típicamente, \(1 - \alpha = 0.90\), \(0.95\), o \(0.99\).

\end{tcolorbox}

\subsection{Interpretación
Frecuentista}\label{interpretaciuxf3n-frecuentista}

\begin{tcolorbox}[enhanced jigsaw, breakable, titlerule=0mm, rightrule=.15mm, leftrule=.75mm, left=2mm, opacitybacktitle=0.6, colback=white, bottomtitle=1mm, colbacktitle=quarto-callout-warning-color!10!white, title=\textcolor{quarto-callout-warning-color}{\faExclamationTriangle}\hspace{0.5em}{Interpretación Correcta de un IC}, toptitle=1mm, coltitle=black, opacityback=0, arc=.35mm, bottomrule=.15mm, toprule=.15mm, colframe=quarto-callout-warning-color-frame]

\textbf{Antes de observar los datos:} El intervalo
\([L(\mathbf{X}), U(\mathbf{X})]\) es \textbf{aleatorio}. La afirmación
\(P(L \leq \theta \leq U) = 1-\alpha\) significa que si repitiéramos el
procedimiento de muestreo muchas veces, aproximadamente
\(100(1-\alpha)\%\) de los intervalos resultantes contendrían el
verdadero parámetro.

\textbf{Después de observar los datos:} El intervalo observado
\([l, u]\) es \textbf{fijo} (con números reales). No es correcto decir
``la probabilidad de que \(\theta\) esté en \([l, u]\) es \(1-\alpha\)''
porque \(\theta\) es una constante desconocida, no una variable
aleatoria. Simplemente decimos que el intervalo observado es una
realización de un procedimiento que acierta \(100(1-\alpha)\%\) de las
veces.

\end{tcolorbox}

\subsection{Intervalo de Confianza para la Media (Varianza
Conocida)}\label{intervalo-de-confianza-para-la-media-varianza-conocida}

\textbf{Supuestos:} - Población Normal: \(X \sim N(\mu, \sigma)\) -
Varianza poblacional \(\sigma^2\) \textbf{conocida} - Muestra aleatoria
simple de tamaño \(n\)

\textbf{Distribución del estimador:}
\[\bar{X} \sim N\left(\mu, \frac{\sigma}{\sqrt{n}}\right) \quad \Rightarrow \quad Z = \frac{\bar{X} - \mu}{\sigma/\sqrt{n}} \sim N(0,1)\]

\textbf{Construcción del IC:} Para un nivel de confianza \(1-\alpha\),
buscamos \(z_{1-\alpha/2}\) tal que
\(P(-z_{1-\alpha/2} \leq Z \leq z_{1-\alpha/2}) = 1-\alpha\).

\[P\left(\bar{X} - z_{1-\alpha/2} \frac{\sigma}{\sqrt{n}} \leq \mu \leq \bar{X} + z_{1-\alpha/2} \frac{\sigma}{\sqrt{n}}\right) = 1 - \alpha\]

\textbf{Intervalo de Confianza:}
\[\text{IC}_{1-\alpha}(\mu) = \left[\bar{x} - z_{1-\alpha/2} \frac{\sigma}{\sqrt{n}}, \quad \bar{x} + z_{1-\alpha/2} \frac{\sigma}{\sqrt{n}}\right]\]

\subsection{Intervalo de Confianza para la Media (Varianza
Desconocida)}\label{intervalo-de-confianza-para-la-media-varianza-desconocida}

\textbf{Supuestos:} - Población Normal: \(X \sim N(\mu, \sigma)\) -
Varianza poblacional \(\sigma^2\) \textbf{desconocida}

\textbf{Distribución del estimador:}
\[T = \frac{\bar{X} - \mu}{S/\sqrt{n}} \sim t_{n-1}\]

donde \(S = \sqrt{\frac{1}{n-1}\sum_{i=1}^n (X_i - \bar{X})^2}\) es la
desviación estándar muestral.

\textbf{Intervalo de Confianza:}
\[\text{IC}_{1-\alpha}(\mu) = \left[\bar{x} - t_{1-\alpha/2, n-1} \frac{s}{\sqrt{n}}, \quad \bar{x} + t_{1-\alpha/2, n-1} \frac{s}{\sqrt{n}}\right]\]

donde \(t_{1-\alpha/2, n-1}\) es el cuantil \((1-\alpha/2)\) de la
distribución \(t\) con \(n-1\) grados de libertad.

\begin{tcolorbox}[enhanced jigsaw, breakable, titlerule=0mm, rightrule=.15mm, leftrule=.75mm, left=2mm, opacitybacktitle=0.6, colback=white, bottomtitle=1mm, colbacktitle=quarto-callout-tip-color!10!white, title=\textcolor{quarto-callout-tip-color}{\faLightbulb}\hspace{0.5em}{Ejemplo: IC para la Media en R}, toptitle=1mm, coltitle=black, opacityback=0, arc=.35mm, bottomrule=.15mm, toprule=.15mm, colframe=quarto-callout-tip-color-frame]

\begin{codelisting}[H]

\caption{\label{lst-semana07-3}Código R}

\centering{

\begin{Shaded}
\begin{Highlighting}[]
\CommentTok{\# Presión arterial sistólica (mmHg) en 25 pacientes adultos}
\CommentTok{\# Modelo poblacional: Normal con μ ≈ 130 mmHg y σ = 15 mmHg}
\FunctionTok{set.seed}\NormalTok{(}\DecValTok{100}\NormalTok{)}
\NormalTok{pas }\OtherTok{\textless{}{-}} \FunctionTok{rnorm}\NormalTok{(}\DecValTok{25}\NormalTok{, }\AttributeTok{mean =} \DecValTok{130}\NormalTok{, }\AttributeTok{sd =} \DecValTok{15}\NormalTok{)}

\CommentTok{\# IC con t{-}Student (varianza desconocida) — caso habitual en clínica}
\NormalTok{resultado }\OtherTok{\textless{}{-}} \FunctionTok{t.test}\NormalTok{(pas, }\AttributeTok{conf.level =} \FloatTok{0.95}\NormalTok{)}

\FunctionTok{cat}\NormalTok{(}\StringTok{"Intervalo de Confianza 95\% para la PAS media (mmHg):}\SpecialCharTok{\textbackslash{}n}\StringTok{"}\NormalTok{)}
\end{Highlighting}
\end{Shaded}

}

\end{codelisting}%

\begin{verbatim}
Intervalo de Confianza 95% para la PAS media (mmHg):
\end{verbatim}

\begin{codelisting}[H]

\caption{\label{lst-semana07-3}Código R}

\centering{

\begin{Shaded}
\begin{Highlighting}[]
\FunctionTok{cat}\NormalTok{(}\StringTok{"Media muestral:"}\NormalTok{, }\FunctionTok{round}\NormalTok{(resultado}\SpecialCharTok{$}\NormalTok{estimate, }\DecValTok{2}\NormalTok{), }\StringTok{"mmHg}\SpecialCharTok{\textbackslash{}n}\StringTok{"}\NormalTok{)}
\end{Highlighting}
\end{Shaded}

}

\end{codelisting}%

\begin{verbatim}
Media muestral: 132 mmHg
\end{verbatim}

\begin{codelisting}[H]

\caption{\label{lst-semana07-3}Código R}

\centering{

\begin{Shaded}
\begin{Highlighting}[]
\FunctionTok{cat}\NormalTok{(}\StringTok{"IC: ["}\NormalTok{, }\FunctionTok{round}\NormalTok{(resultado}\SpecialCharTok{$}\NormalTok{conf.int[}\DecValTok{1}\NormalTok{], }\DecValTok{2}\NormalTok{), }\StringTok{","}\NormalTok{,}
    \FunctionTok{round}\NormalTok{(resultado}\SpecialCharTok{$}\NormalTok{conf.int[}\DecValTok{2}\NormalTok{], }\DecValTok{2}\NormalTok{), }\StringTok{"] mmHg}\SpecialCharTok{\textbackslash{}n}\StringTok{"}\NormalTok{)}
\end{Highlighting}
\end{Shaded}

}

\end{codelisting}%

\begin{verbatim}
IC: [ 127 , 136 ] mmHg
\end{verbatim}

\begin{codelisting}[H]

\caption{\label{lst-semana07-3}Código R}

\centering{

\begin{Shaded}
\begin{Highlighting}[]
\CommentTok{\# IC con z (varianza poblacional σ = 15 mmHg supuesta conocida)}
\CommentTok{\# Útil cuando se dispone de datos históricos/literatura}
\NormalTok{n }\OtherTok{\textless{}{-}} \FunctionTok{length}\NormalTok{(pas)}
\NormalTok{media\_muestral }\OtherTok{\textless{}{-}} \FunctionTok{mean}\NormalTok{(pas)}
\NormalTok{sigma }\OtherTok{\textless{}{-}} \DecValTok{15}
\NormalTok{error\_estandar }\OtherTok{\textless{}{-}}\NormalTok{ sigma }\SpecialCharTok{/} \FunctionTok{sqrt}\NormalTok{(n)}
\NormalTok{z\_critico }\OtherTok{\textless{}{-}} \FunctionTok{qnorm}\NormalTok{(}\FloatTok{0.975}\NormalTok{)  }\CommentTok{\# α = 0.05, bilateral}

\NormalTok{IC\_conocida }\OtherTok{\textless{}{-}} \FunctionTok{c}\NormalTok{(media\_muestral }\SpecialCharTok{{-}}\NormalTok{ z\_critico }\SpecialCharTok{*}\NormalTok{ error\_estandar,}
\NormalTok{                 media\_muestral }\SpecialCharTok{+}\NormalTok{ z\_critico }\SpecialCharTok{*}\NormalTok{ error\_estandar)}

\FunctionTok{cat}\NormalTok{(}\StringTok{"}\SpecialCharTok{\textbackslash{}n}\StringTok{IC con varianza conocida (σ = 15 mmHg):}\SpecialCharTok{\textbackslash{}n}\StringTok{"}\NormalTok{)}
\end{Highlighting}
\end{Shaded}

}

\end{codelisting}%

\begin{verbatim}

IC con varianza conocida (σ = 15 mmHg):
\end{verbatim}

\begin{codelisting}[H]

\caption{\label{lst-semana07-3}Código R}

\centering{

\begin{Shaded}
\begin{Highlighting}[]
\FunctionTok{cat}\NormalTok{(}\StringTok{"IC: ["}\NormalTok{, }\FunctionTok{round}\NormalTok{(IC\_conocida[}\DecValTok{1}\NormalTok{], }\DecValTok{2}\NormalTok{), }\StringTok{","}\NormalTok{,}
    \FunctionTok{round}\NormalTok{(IC\_conocida[}\DecValTok{2}\NormalTok{], }\DecValTok{2}\NormalTok{), }\StringTok{"] mmHg}\SpecialCharTok{\textbackslash{}n}\StringTok{"}\NormalTok{)}
\end{Highlighting}
\end{Shaded}

}

\end{codelisting}%

\begin{verbatim}
IC: [ 126 , 138 ] mmHg
\end{verbatim}

\end{tcolorbox}

\subsection{Intervalo de Confianza para una
Proporción}\label{intervalo-de-confianza-para-una-proporciuxf3n}

\textbf{Supuestos:} - Población dicotómica: \(P(A) = \pi\) (proporción
de éxitos) - Muestra aleatoria simple de tamaño \(n\) (con \(n\)
suficientemente grande) - Condiciones: \(n\pi \geq 5\) y
\(n(1-\pi) \geq 5\) (o aproximadamente \(n > 30\))

\textbf{Distribución del estimador:}
\[\hat{\pi} = \frac{X}{n} \approx N\left(\pi, \sqrt{\frac{\pi(1-\pi)}{n}}\right)\]

\textbf{Problema:} La varianza contiene el parámetro desconocido
\(\pi\). \textbf{Solución:} Usar \(\hat{\pi}\) como estimador
consistente de \(\pi\) en la varianza.

\textbf{Intervalo de Confianza (aproximado):}
\[\text{IC}_{1-\alpha}(\pi) = \left[\hat{\pi} - z_{1-\alpha/2}\sqrt{\frac{\hat{\pi}(1-\hat{\pi})}{n}}, \quad \hat{\pi} + z_{1-\alpha/2}\sqrt{\frac{\hat{\pi}(1-\hat{\pi})}{n}}\right]\]

\begin{tcolorbox}[enhanced jigsaw, breakable, titlerule=0mm, rightrule=.15mm, leftrule=.75mm, left=2mm, opacitybacktitle=0.6, colback=white, bottomtitle=1mm, colbacktitle=quarto-callout-tip-color!10!white, title=\textcolor{quarto-callout-tip-color}{\faLightbulb}\hspace{0.5em}{Ejemplo: IC para una Proporción en R}, toptitle=1mm, coltitle=black, opacityback=0, arc=.35mm, bottomrule=.15mm, toprule=.15mm, colframe=quarto-callout-tip-color-frame]

\begin{codelisting}[H]

\caption{\label{lst-semana07-4}Código R}

\centering{

\begin{Shaded}
\begin{Highlighting}[]
\CommentTok{\# Estudio en 200 pacientes hipertensos en seguimiento clínico:}
\CommentTok{\# 132 alcanzan control adecuado de PAS (\textless{} 140/90 mmHg)}

\NormalTok{n }\OtherTok{\textless{}{-}} \DecValTok{200}
\NormalTok{x }\OtherTok{\textless{}{-}} \DecValTok{132}
\NormalTok{p\_hat }\OtherTok{\textless{}{-}}\NormalTok{ x }\SpecialCharTok{/}\NormalTok{ n}

\CommentTok{\# Intervalo de confianza 95\% para la proporción}
\NormalTok{conf\_level }\OtherTok{\textless{}{-}} \FloatTok{0.95}
\NormalTok{alpha }\OtherTok{\textless{}{-}} \DecValTok{1} \SpecialCharTok{{-}}\NormalTok{ conf\_level}
\NormalTok{z\_critico }\OtherTok{\textless{}{-}} \FunctionTok{qnorm}\NormalTok{(}\DecValTok{1} \SpecialCharTok{{-}}\NormalTok{ alpha}\SpecialCharTok{/}\DecValTok{2}\NormalTok{)}

\NormalTok{error\_estandar }\OtherTok{\textless{}{-}} \FunctionTok{sqrt}\NormalTok{(p\_hat }\SpecialCharTok{*}\NormalTok{ (}\DecValTok{1} \SpecialCharTok{{-}}\NormalTok{ p\_hat) }\SpecialCharTok{/}\NormalTok{ n)}
\NormalTok{IC }\OtherTok{\textless{}{-}} \FunctionTok{c}\NormalTok{(p\_hat }\SpecialCharTok{{-}}\NormalTok{ z\_critico }\SpecialCharTok{*}\NormalTok{ error\_estandar,}
\NormalTok{        p\_hat }\SpecialCharTok{+}\NormalTok{ z\_critico }\SpecialCharTok{*}\NormalTok{ error\_estandar)}

\FunctionTok{cat}\NormalTok{(}\StringTok{"IC 95\% para la proporción de hipertensos controlados:}\SpecialCharTok{\textbackslash{}n}\StringTok{"}\NormalTok{)}
\end{Highlighting}
\end{Shaded}

}

\end{codelisting}%

\begin{verbatim}
IC 95% para la proporción de hipertensos controlados:
\end{verbatim}

\begin{codelisting}[H]

\caption{\label{lst-semana07-4}Código R}

\centering{

\begin{Shaded}
\begin{Highlighting}[]
\FunctionTok{cat}\NormalTok{(}\StringTok{"Proporción muestral:"}\NormalTok{, }\FunctionTok{round}\NormalTok{(p\_hat, }\DecValTok{4}\NormalTok{),}
    \StringTok{"("}\NormalTok{, }\FunctionTok{round}\NormalTok{(p\_hat}\SpecialCharTok{*}\DecValTok{100}\NormalTok{, }\DecValTok{1}\NormalTok{), }\StringTok{"\%)}\SpecialCharTok{\textbackslash{}n}\StringTok{"}\NormalTok{)}
\end{Highlighting}
\end{Shaded}

}

\end{codelisting}%

\begin{verbatim}
Proporción muestral: 0.66 ( 66 %)
\end{verbatim}

\begin{codelisting}[H]

\caption{\label{lst-semana07-4}Código R}

\centering{

\begin{Shaded}
\begin{Highlighting}[]
\FunctionTok{cat}\NormalTok{(}\StringTok{"IC: ["}\NormalTok{, }\FunctionTok{round}\NormalTok{(IC[}\DecValTok{1}\NormalTok{], }\DecValTok{4}\NormalTok{), }\StringTok{","}\NormalTok{, }\FunctionTok{round}\NormalTok{(IC[}\DecValTok{2}\NormalTok{], }\DecValTok{4}\NormalTok{), }\StringTok{"]}\SpecialCharTok{\textbackslash{}n}\StringTok{"}\NormalTok{)}
\end{Highlighting}
\end{Shaded}

}

\end{codelisting}%

\begin{verbatim}
IC: [ 0.594 , 0.726 ]
\end{verbatim}

\begin{codelisting}[H]

\caption{\label{lst-semana07-4}Código R}

\centering{

\begin{Shaded}
\begin{Highlighting}[]
\FunctionTok{cat}\NormalTok{(}\StringTok{"Interpretación: con 95\% de confianza, entre el"}\NormalTok{,}
    \FunctionTok{round}\NormalTok{(IC[}\DecValTok{1}\NormalTok{]}\SpecialCharTok{*}\DecValTok{100}\NormalTok{, }\DecValTok{1}\NormalTok{), }\StringTok{"\% y el"}\NormalTok{,}
    \FunctionTok{round}\NormalTok{(IC[}\DecValTok{2}\NormalTok{]}\SpecialCharTok{*}\DecValTok{100}\NormalTok{, }\DecValTok{1}\NormalTok{), }\StringTok{"\%}\SpecialCharTok{\textbackslash{}n}\StringTok{"}\NormalTok{)}
\end{Highlighting}
\end{Shaded}

}

\end{codelisting}%

\begin{verbatim}
Interpretación: con 95% de confianza, entre el 59.4 % y el 72.6 %
\end{verbatim}

\begin{codelisting}[H]

\caption{\label{lst-semana07-4}Código R}

\centering{

\begin{Shaded}
\begin{Highlighting}[]
\FunctionTok{cat}\NormalTok{(}\StringTok{"de los pacientes hipertensos poblacionales alcanzan control adecuado.}\SpecialCharTok{\textbackslash{}n}\StringTok{"}\NormalTok{)}
\end{Highlighting}
\end{Shaded}

}

\end{codelisting}%

\begin{verbatim}
de los pacientes hipertensos poblacionales alcanzan control adecuado.
\end{verbatim}

\end{tcolorbox}

\section{Interpretación de Intervalos de
Confianza}\label{sec-semana06-interp-ic}

Interpretar correctamente un intervalo de confianza es fundamental.
Consideremos un ejemplo:

\textbf{Supongamos} que construimos un IC 95\% para la media y obtenemos
\([48.3, 51.7]\).

\begin{tcolorbox}[enhanced jigsaw, breakable, titlerule=0mm, rightrule=.15mm, leftrule=.75mm, left=2mm, opacitybacktitle=0.6, colback=white, bottomtitle=1mm, colbacktitle=quarto-callout-warning-color!10!white, title=\textcolor{quarto-callout-warning-color}{\faExclamationTriangle}\hspace{0.5em}{Interpretación CORRECTA}, toptitle=1mm, coltitle=black, opacityback=0, arc=.35mm, bottomrule=.15mm, toprule=.15mm, colframe=quarto-callout-warning-color-frame]

``Si repitiéramos el proceso de muestreo muchas veces (con muestras de
igual tamaño de la misma población) y construyéramos un intervalo de
confianza del 95\% para cada muestra, aproximadamente el 95\% de estos
intervalos contendría el verdadero valor de la media poblacional.''

En otras palabras, el IC es sobre la \textbf{fiabilidad del
procedimiento}, no sobre la probabilidad de que un intervalo específico
contenga el parámetro.

\end{tcolorbox}

\begin{tcolorbox}[enhanced jigsaw, breakable, titlerule=0mm, rightrule=.15mm, leftrule=.75mm, left=2mm, opacitybacktitle=0.6, colback=white, bottomtitle=1mm, colbacktitle=quarto-callout-warning-color!10!white, title=\textcolor{quarto-callout-warning-color}{\faExclamationTriangle}\hspace{0.5em}{Interpretación INCORRECTA (Evitar)}, toptitle=1mm, coltitle=black, opacityback=0, arc=.35mm, bottomrule=.15mm, toprule=.15mm, colframe=quarto-callout-warning-color-frame]

``La verdadera media está en \([48.3, 51.7]\) con probabilidad 0.95''

Esto es incorrecto porque después de observar los datos, el intervalo es
fijo y la verdadera media es constante (no aleatoria). El parámetro o
está en el intervalo o no está.

\end{tcolorbox}

\section{Trabajando con Ejemplos
Completos}\label{sec-semana06-ejemplos-completos}

\begin{tcolorbox}[enhanced jigsaw, breakable, titlerule=0mm, rightrule=.15mm, leftrule=.75mm, left=2mm, opacitybacktitle=0.6, colback=white, bottomtitle=1mm, colbacktitle=quarto-callout-tip-color!10!white, title=\textcolor{quarto-callout-tip-color}{\faLightbulb}\hspace{0.5em}{Ejemplo: Análisis Completo de Datos de HbA1c (Diabetes Tipo 2)}, toptitle=1mm, coltitle=black, opacityback=0, arc=.35mm, bottomrule=.15mm, toprule=.15mm, colframe=quarto-callout-tip-color-frame]

\begin{codelisting}[H]

\caption{\label{lst-semana07-5}Código R}

\centering{

\begin{Shaded}
\begin{Highlighting}[]
\CommentTok{\# Datos de HbA1c (\%) de 20 pacientes con diabetes tipo 2}
\CommentTok{\# atendidos en la Facultad de Medicina, Universidad de Granada}

\NormalTok{hba1c }\OtherTok{\textless{}{-}} \FunctionTok{c}\NormalTok{(}\FloatTok{7.2}\NormalTok{, }\FloatTok{8.1}\NormalTok{, }\FloatTok{6.8}\NormalTok{, }\FloatTok{9.3}\NormalTok{, }\FloatTok{7.5}\NormalTok{, }\FloatTok{8.8}\NormalTok{, }\FloatTok{6.5}\NormalTok{, }\FloatTok{10.2}\NormalTok{,}
           \FloatTok{7.9}\NormalTok{, }\FloatTok{8.4}\NormalTok{, }\FloatTok{7.1}\NormalTok{, }\FloatTok{9.7}\NormalTok{, }\FloatTok{7.3}\NormalTok{, }\FloatTok{8.6}\NormalTok{, }\FloatTok{6.9}\NormalTok{, }\FloatTok{11.4}\NormalTok{,}
           \FloatTok{8.2}\NormalTok{, }\FloatTok{7.0}\NormalTok{, }\FloatTok{9.1}\NormalTok{, }\FloatTok{7.6}\NormalTok{)}

\CommentTok{\# 1. ESTADÍSTICA DESCRIPTIVA}
\NormalTok{media\_muestral }\OtherTok{\textless{}{-}} \FunctionTok{mean}\NormalTok{(hba1c)}
\NormalTok{desv\_est }\OtherTok{\textless{}{-}} \FunctionTok{sd}\NormalTok{(hba1c)  }\CommentTok{\# Usa divisor n{-}1 (insesgado)}
\NormalTok{n }\OtherTok{\textless{}{-}} \FunctionTok{length}\NormalTok{(hba1c)}

\FunctionTok{cat}\NormalTok{(}\StringTok{"=== ESTADÍSTICA DESCRIPTIVA ===}\SpecialCharTok{\textbackslash{}n}\StringTok{"}\NormalTok{)}
\end{Highlighting}
\end{Shaded}

}

\end{codelisting}%

\begin{verbatim}
=== ESTADÍSTICA DESCRIPTIVA ===
\end{verbatim}

\begin{codelisting}[H]

\caption{\label{lst-semana07-5}Código R}

\centering{

\begin{Shaded}
\begin{Highlighting}[]
\FunctionTok{cat}\NormalTok{(}\StringTok{"n ="}\NormalTok{, n, }\StringTok{"}\SpecialCharTok{\textbackslash{}n}\StringTok{"}\NormalTok{)}
\end{Highlighting}
\end{Shaded}

}

\end{codelisting}%

\begin{verbatim}
n = 20 
\end{verbatim}

\begin{codelisting}[H]

\caption{\label{lst-semana07-5}Código R}

\centering{

\begin{Shaded}
\begin{Highlighting}[]
\FunctionTok{cat}\NormalTok{(}\StringTok{"Media muestral = "}\NormalTok{, }\FunctionTok{round}\NormalTok{(media\_muestral, }\DecValTok{2}\NormalTok{), }\StringTok{"\%}\SpecialCharTok{\textbackslash{}n}\StringTok{"}\NormalTok{)}
\end{Highlighting}
\end{Shaded}

}

\end{codelisting}%

\begin{verbatim}
Media muestral =  8.18 %
\end{verbatim}

\begin{codelisting}[H]

\caption{\label{lst-semana07-5}Código R}

\centering{

\begin{Shaded}
\begin{Highlighting}[]
\FunctionTok{cat}\NormalTok{(}\StringTok{"Desviación estándar = "}\NormalTok{, }\FunctionTok{round}\NormalTok{(desv\_est, }\DecValTok{2}\NormalTok{), }\StringTok{"\%}\SpecialCharTok{\textbackslash{}n}\StringTok{"}\NormalTok{)}
\end{Highlighting}
\end{Shaded}

}

\end{codelisting}%

\begin{verbatim}
Desviación estándar =  1.28 %
\end{verbatim}

\begin{codelisting}[H]

\caption{\label{lst-semana07-5}Código R}

\centering{

\begin{Shaded}
\begin{Highlighting}[]
\FunctionTok{cat}\NormalTok{(}\StringTok{"Varianza muestral = "}\NormalTok{, }\FunctionTok{round}\NormalTok{(desv\_est}\SpecialCharTok{\^{}}\DecValTok{2}\NormalTok{, }\DecValTok{2}\NormalTok{), }\StringTok{"}\SpecialCharTok{\textbackslash{}n\textbackslash{}n}\StringTok{"}\NormalTok{)}
\end{Highlighting}
\end{Shaded}

}

\end{codelisting}%

\begin{verbatim}
Varianza muestral =  1.63 
\end{verbatim}

\begin{codelisting}[H]

\caption{\label{lst-semana07-5}Código R}

\centering{

\begin{Shaded}
\begin{Highlighting}[]
\CommentTok{\# 2. ESTIMACIÓN PUNTUAL (MLE/MoM para Normal)}
\FunctionTok{cat}\NormalTok{(}\StringTok{"=== ESTIMACIÓN PUNTUAL ===}\SpecialCharTok{\textbackslash{}n}\StringTok{"}\NormalTok{)}
\end{Highlighting}
\end{Shaded}

}

\end{codelisting}%

\begin{verbatim}
=== ESTIMACIÓN PUNTUAL ===
\end{verbatim}

\begin{codelisting}[H]

\caption{\label{lst-semana07-5}Código R}

\centering{

\begin{Shaded}
\begin{Highlighting}[]
\FunctionTok{cat}\NormalTok{(}\StringTok{"μ (estimado) = "}\NormalTok{, }\FunctionTok{round}\NormalTok{(media\_muestral, }\DecValTok{2}\NormalTok{), }\StringTok{"\%}\SpecialCharTok{\textbackslash{}n}\StringTok{"}\NormalTok{)}
\end{Highlighting}
\end{Shaded}

}

\end{codelisting}%

\begin{verbatim}
μ (estimado) =  8.18 %
\end{verbatim}

\begin{codelisting}[H]

\caption{\label{lst-semana07-5}Código R}

\centering{

\begin{Shaded}
\begin{Highlighting}[]
\FunctionTok{cat}\NormalTok{(}\StringTok{"σ² (estimado) = "}\NormalTok{, }\FunctionTok{round}\NormalTok{(desv\_est}\SpecialCharTok{\^{}}\DecValTok{2}\NormalTok{, }\DecValTok{2}\NormalTok{), }\StringTok{"}\SpecialCharTok{\textbackslash{}n\textbackslash{}n}\StringTok{"}\NormalTok{)}
\end{Highlighting}
\end{Shaded}

}

\end{codelisting}%

\begin{verbatim}
σ² (estimado) =  1.63 
\end{verbatim}

\begin{codelisting}[H]

\caption{\label{lst-semana07-5}Código R}

\centering{

\begin{Shaded}
\begin{Highlighting}[]
\CommentTok{\# 3. INTERVALO DE CONFIANZA 95\% (varianza desconocida)}
\FunctionTok{cat}\NormalTok{(}\StringTok{"=== INTERVALO DE CONFIANZA (95\%) ===}\SpecialCharTok{\textbackslash{}n}\StringTok{"}\NormalTok{)}
\end{Highlighting}
\end{Shaded}

}

\end{codelisting}%

\begin{verbatim}
=== INTERVALO DE CONFIANZA (95%) ===
\end{verbatim}

\begin{codelisting}[H]

\caption{\label{lst-semana07-5}Código R}

\centering{

\begin{Shaded}
\begin{Highlighting}[]
\NormalTok{t\_critico }\OtherTok{\textless{}{-}} \FunctionTok{qt}\NormalTok{(}\FloatTok{0.975}\NormalTok{, }\AttributeTok{df =}\NormalTok{ n }\SpecialCharTok{{-}} \DecValTok{1}\NormalTok{)}
\NormalTok{error\_estandar }\OtherTok{\textless{}{-}}\NormalTok{ desv\_est }\SpecialCharTok{/} \FunctionTok{sqrt}\NormalTok{(n)}
\NormalTok{margen\_error }\OtherTok{\textless{}{-}}\NormalTok{ t\_critico }\SpecialCharTok{*}\NormalTok{ error\_estandar}

\NormalTok{IC\_inferior }\OtherTok{\textless{}{-}}\NormalTok{ media\_muestral }\SpecialCharTok{{-}}\NormalTok{ margen\_error}
\NormalTok{IC\_superior }\OtherTok{\textless{}{-}}\NormalTok{ media\_muestral }\SpecialCharTok{+}\NormalTok{ margen\_error}

\FunctionTok{cat}\NormalTok{(}\StringTok{"Error estándar = "}\NormalTok{, }\FunctionTok{round}\NormalTok{(error\_estandar, }\DecValTok{2}\NormalTok{), }\StringTok{"\%}\SpecialCharTok{\textbackslash{}n}\StringTok{"}\NormalTok{)}
\end{Highlighting}
\end{Shaded}

}

\end{codelisting}%

\begin{verbatim}
Error estándar =  0.29 %
\end{verbatim}

\begin{codelisting}[H]

\caption{\label{lst-semana07-5}Código R}

\centering{

\begin{Shaded}
\begin{Highlighting}[]
\FunctionTok{cat}\NormalTok{(}\StringTok{"Margen de error = "}\NormalTok{, }\FunctionTok{round}\NormalTok{(margen\_error, }\DecValTok{2}\NormalTok{), }\StringTok{"\%}\SpecialCharTok{\textbackslash{}n}\StringTok{"}\NormalTok{)}
\end{Highlighting}
\end{Shaded}

}

\end{codelisting}%

\begin{verbatim}
Margen de error =  0.6 %
\end{verbatim}

\begin{codelisting}[H]

\caption{\label{lst-semana07-5}Código R}

\centering{

\begin{Shaded}
\begin{Highlighting}[]
\FunctionTok{cat}\NormalTok{(}\StringTok{"IC 95\% = ["}\NormalTok{, }\FunctionTok{round}\NormalTok{(IC\_inferior, }\DecValTok{2}\NormalTok{), }\StringTok{"\%, "}\NormalTok{,}
    \FunctionTok{round}\NormalTok{(IC\_superior, }\DecValTok{2}\NormalTok{), }\StringTok{"\%]}\SpecialCharTok{\textbackslash{}n\textbackslash{}n}\StringTok{"}\NormalTok{)}
\end{Highlighting}
\end{Shaded}

}

\end{codelisting}%

\begin{verbatim}
IC 95% = [ 7.58 %,  8.78 %]
\end{verbatim}

\begin{codelisting}[H]

\caption{\label{lst-semana07-5}Código R}

\centering{

\begin{Shaded}
\begin{Highlighting}[]
\CommentTok{\# Alternativamente, usar t.test()}
\NormalTok{test\_result }\OtherTok{\textless{}{-}} \FunctionTok{t.test}\NormalTok{(hba1c, }\AttributeTok{conf.level =} \FloatTok{0.95}\NormalTok{)}
\FunctionTok{cat}\NormalTok{(}\StringTok{"Verificación con t.test():}\SpecialCharTok{\textbackslash{}n}\StringTok{"}\NormalTok{)}
\end{Highlighting}
\end{Shaded}

}

\end{codelisting}%

\begin{verbatim}
Verificación con t.test():
\end{verbatim}

\begin{codelisting}[H]

\caption{\label{lst-semana07-5}Código R}

\centering{

\begin{Shaded}
\begin{Highlighting}[]
\FunctionTok{cat}\NormalTok{(}\StringTok{"IC 95\% = ["}\NormalTok{, }\FunctionTok{round}\NormalTok{(test\_result}\SpecialCharTok{$}\NormalTok{conf.int[}\DecValTok{1}\NormalTok{], }\DecValTok{2}\NormalTok{), }\StringTok{"\%, "}\NormalTok{,}
    \FunctionTok{round}\NormalTok{(test\_result}\SpecialCharTok{$}\NormalTok{conf.int[}\DecValTok{2}\NormalTok{], }\DecValTok{2}\NormalTok{), }\StringTok{"\%]}\SpecialCharTok{\textbackslash{}n}\StringTok{"}\NormalTok{)}
\end{Highlighting}
\end{Shaded}

}

\end{codelisting}%

\begin{verbatim}
IC 95% = [ 7.58 %,  8.78 %]
\end{verbatim}

\textbf{Interpretación:} Con el 95\% de confianza, el nivel medio de
HbA1c poblacional en pacientes diabéticos está entre 7.54\% y 8.62\%.
Este intervalo incluye el objetivo clínico de HbA1c \textless{} 7\%,
sugiriendo que la cohorte presenta control glucémico moderado. El margen
de error (\textasciitilde0.54\%) refleja la variabilidad en control
metabólico entre pacientes, lo que en práctica clínica orienta
decisiones de intensificación terapéutica.

\end{tcolorbox}

\begin{tcolorbox}[enhanced jigsaw, breakable, titlerule=0mm, rightrule=.15mm, leftrule=.75mm, left=2mm, opacitybacktitle=0.6, colback=white, bottomtitle=1mm, colbacktitle=quarto-callout-tip-color!10!white, title=\textcolor{quarto-callout-tip-color}{\faLightbulb}\hspace{0.5em}{Ejemplo: MLE para tiempos entre eventos cardiovasculares}, toptitle=1mm, coltitle=black, opacityback=0, arc=.35mm, bottomrule=.15mm, toprule=.15mm, colframe=quarto-callout-tip-color-frame]

Modelamos el \textbf{tiempo (en años) entre eventos cardiovasculares
mayores} en una cohorte de 50 pacientes con cardiopatía isquémica como
una distribución Exponencial con tasa λ. La media poblacional es
\(1/\lambda\) años entre eventos.

\begin{codelisting}[H]

\caption{\label{lst-semana07-6}Código R}

\centering{

\begin{Shaded}
\begin{Highlighting}[]
\CommentTok{\# Tiempo (años) entre eventos cardiovasculares — 50 pacientes}
\CommentTok{\# Modelo poblacional: Exponencial con tasa λ = 0.5 eventos/año}
\CommentTok{\# (i.e., un evento cada 2 años en promedio)}
\FunctionTok{set.seed}\NormalTok{(}\DecValTok{999}\NormalTok{)}
\NormalTok{tiempo\_eventos }\OtherTok{\textless{}{-}} \FunctionTok{rexp}\NormalTok{(}\DecValTok{50}\NormalTok{, }\AttributeTok{rate =} \FloatTok{0.5}\NormalTok{)}

\CommentTok{\# Función de log{-}verosimilitud para Exponencial}
\NormalTok{log\_verosimilitud }\OtherTok{\textless{}{-}} \ControlFlowTok{function}\NormalTok{(lambda, datos) \{}
  \ControlFlowTok{if}\NormalTok{ (lambda }\SpecialCharTok{\textless{}=} \DecValTok{0}\NormalTok{) }\FunctionTok{return}\NormalTok{(}\SpecialCharTok{{-}}\ConstantTok{Inf}\NormalTok{)}
  \FunctionTok{sum}\NormalTok{(}\FunctionTok{dexp}\NormalTok{(datos, }\AttributeTok{rate =}\NormalTok{ lambda, }\AttributeTok{log =} \ConstantTok{TRUE}\NormalTok{))}
\NormalTok{\}}

\CommentTok{\# Maximizar log{-}verosimilitud numéricamente}
\NormalTok{resultado\_optim }\OtherTok{\textless{}{-}} \FunctionTok{optimize}\NormalTok{(}
  \ControlFlowTok{function}\NormalTok{(lambda) }\FunctionTok{log\_verosimilitud}\NormalTok{(lambda, tiempo\_eventos),}
  \AttributeTok{interval =} \FunctionTok{c}\NormalTok{(}\FloatTok{0.01}\NormalTok{, }\DecValTok{5}\NormalTok{),}
  \AttributeTok{maximum =} \ConstantTok{TRUE}
\NormalTok{)}

\NormalTok{lambda\_mle }\OtherTok{\textless{}{-}}\NormalTok{ resultado\_optim}\SpecialCharTok{$}\NormalTok{maximum}
\FunctionTok{cat}\NormalTok{(}\StringTok{"Estimador MLE de λ:"}\NormalTok{, }\FunctionTok{round}\NormalTok{(lambda\_mle, }\DecValTok{4}\NormalTok{), }\StringTok{"eventos/año}\SpecialCharTok{\textbackslash{}n}\StringTok{"}\NormalTok{)}
\end{Highlighting}
\end{Shaded}

}

\end{codelisting}%

\begin{verbatim}
Estimador MLE de λ: 0.511 eventos/año
\end{verbatim}

\begin{codelisting}[H]

\caption{\label{lst-semana07-6}Código R}

\centering{

\begin{Shaded}
\begin{Highlighting}[]
\FunctionTok{cat}\NormalTok{(}\StringTok{"Log{-}verosimilitud máxima:"}\NormalTok{,}
    \FunctionTok{round}\NormalTok{(resultado\_optim}\SpecialCharTok{$}\NormalTok{objective, }\DecValTok{2}\NormalTok{), }\StringTok{"}\SpecialCharTok{\textbackslash{}n}\StringTok{"}\NormalTok{)}
\end{Highlighting}
\end{Shaded}

}

\end{codelisting}%

\begin{verbatim}
Log-verosimilitud máxima: -83.5 
\end{verbatim}

\begin{codelisting}[H]

\caption{\label{lst-semana07-6}Código R}

\centering{

\begin{Shaded}
\begin{Highlighting}[]
\FunctionTok{cat}\NormalTok{(}\StringTok{"Tiempo medio estimado entre eventos: 1/λ ="}\NormalTok{,}
    \FunctionTok{round}\NormalTok{(}\DecValTok{1}\SpecialCharTok{/}\NormalTok{lambda\_mle, }\DecValTok{2}\NormalTok{), }\StringTok{"años}\SpecialCharTok{\textbackslash{}n}\StringTok{"}\NormalTok{)}
\end{Highlighting}
\end{Shaded}

}

\end{codelisting}%

\begin{verbatim}
Tiempo medio estimado entre eventos: 1/λ = 1.96 años
\end{verbatim}

\begin{codelisting}[H]

\caption{\label{lst-semana07-6}Código R}

\centering{

\begin{Shaded}
\begin{Highlighting}[]
\CommentTok{\# MoM también da λ = 1/media para Exponencial}
\NormalTok{lambda\_mom }\OtherTok{\textless{}{-}} \DecValTok{1} \SpecialCharTok{/} \FunctionTok{mean}\NormalTok{(tiempo\_eventos)}
\FunctionTok{cat}\NormalTok{(}\StringTok{"}\SpecialCharTok{\textbackslash{}n}\StringTok{Estimador MoM de λ:"}\NormalTok{, }\FunctionTok{round}\NormalTok{(lambda\_mom, }\DecValTok{4}\NormalTok{), }\StringTok{"eventos/año}\SpecialCharTok{\textbackslash{}n}\StringTok{"}\NormalTok{)}
\end{Highlighting}
\end{Shaded}

}

\end{codelisting}%

\begin{verbatim}

Estimador MoM de λ: 0.511 eventos/año
\end{verbatim}

\begin{codelisting}[H]

\caption{\label{lst-semana07-6}Código R}

\centering{

\begin{Shaded}
\begin{Highlighting}[]
\CommentTok{\# Para Exponencial, MLE = MoM (propiedad teórica)}
\FunctionTok{cat}\NormalTok{(}\StringTok{"}\SpecialCharTok{\textbackslash{}n}\StringTok{Nota: Para la distribución Exponencial,}\SpecialCharTok{\textbackslash{}n}\StringTok{"}\NormalTok{)}
\end{Highlighting}
\end{Shaded}

}

\end{codelisting}%

\begin{verbatim}

Nota: Para la distribución Exponencial,
\end{verbatim}

\begin{codelisting}[H]

\caption{\label{lst-semana07-6}Código R}

\centering{

\begin{Shaded}
\begin{Highlighting}[]
\FunctionTok{cat}\NormalTok{(}\StringTok{"MLE y MoM coinciden algebraicamente (λ̂ = 1/X̄).}\SpecialCharTok{\textbackslash{}n}\StringTok{"}\NormalTok{)}
\end{Highlighting}
\end{Shaded}

}

\end{codelisting}%

\begin{verbatim}
MLE y MoM coinciden algebraicamente (λ̂ = 1/X̄).
\end{verbatim}

\end{tcolorbox}

\section{Resumen}\label{sec-semana06-resumen}

\begin{longtable}[]{@{}
  >{\raggedright\arraybackslash}p{(\linewidth - 2\tabcolsep) * \real{0.3704}}
  >{\raggedright\arraybackslash}p{(\linewidth - 2\tabcolsep) * \real{0.6296}}@{}}
\toprule\noalign{}
\begin{minipage}[b]{\linewidth}\raggedright
Concepto
\end{minipage} & \begin{minipage}[b]{\linewidth}\raggedright
Definición Clave
\end{minipage} \\
\midrule\noalign{}
\endhead
\bottomrule\noalign{}
\endlastfoot
\textbf{Modelo Paramétrico} & Familia de distribuciones
\(\{f(x;\boldsymbol{\theta}) : \boldsymbol{\theta} \in \Theta\}\) \\
\textbf{Estimador} & Función de la muestra
\(\hat{\theta} = g(X_1, \ldots, X_n)\) que varía aleatoriamente \\
\textbf{Insesgadez} & \(E(\hat{\theta}) = \theta\) \\
\textbf{Eficiencia} & Comparación de varianzas entre estimadores
insesgados \\
\textbf{Consistencia} & \(\lim_{n \to \infty} E(\hat{\theta}) = \theta\)
y \(\lim_{n \to \infty} \text{Var}(\hat{\theta}) = 0\) \\
\textbf{Error Cuadrático Medio} &
\(\text{ECM}(\hat{\theta}) = \text{Var}(\hat{\theta}) + [\text{Sesgo}(\hat{\theta})]^2\) \\
\textbf{Método de Momentos} & Igualar momentos poblacionales con
muestrales: \(E(X^j) = m_j\) \\
\textbf{Función de Verosimilitud} &
\(L(\theta) = \prod_{i=1}^n f(x_i; \theta)\) \\
\textbf{Log-Verosimilitud} &
\(\ell(\theta) = \ln L(\theta) = \sum_{i=1}^n \ln f(x_i; \theta)\) \\
\textbf{Estimador MLE} &
\(\hat{\theta}_{\text{MLE}} = \arg\max_\theta L(\theta)\) \\
\textbf{Intervalo de Confianza} & Intervalo aleatorio
\([L(\mathbf{X}), U(\mathbf{X})]\) con
\(P(L \leq \theta \leq U) = 1-\alpha\) \\
\end{longtable}

\subsection{Fórmulas Clave para Intervalos de
Confianza}\label{fuxf3rmulas-clave-para-intervalos-de-confianza}

\textbf{Media (varianza conocida):}
\[\text{IC}_{1-\alpha}(\mu) = \left[\bar{x} \pm z_{1-\alpha/2} \frac{\sigma}{\sqrt{n}}\right]\]

\textbf{Media (varianza desconocida):}
\[\text{IC}_{1-\alpha}(\mu) = \left[\bar{x} \pm t_{1-\alpha/2, n-1} \frac{s}{\sqrt{n}}\right]\]

\textbf{Proporción:}
\[\text{IC}_{1-\alpha}(\pi) = \left[\hat{\pi} \pm z_{1-\alpha/2} \sqrt{\frac{\hat{\pi}(1-\hat{\pi})}{n}}\right]\]

\begin{tcolorbox}[enhanced jigsaw, breakable, titlerule=0mm, rightrule=.15mm, leftrule=.75mm, left=2mm, opacitybacktitle=0.6, colback=white, bottomtitle=1mm, colbacktitle=quarto-callout-tip-color!10!white, title=\textcolor{quarto-callout-tip-color}{\faLightbulb}\hspace{0.5em}{Ejemplo con BioStatR: icm() e icp()}, toptitle=1mm, coltitle=black, opacityback=0, arc=.35mm, bottomrule=.15mm, toprule=.15mm, colframe=quarto-callout-tip-color-frame]

Documentación completa de ambas funciones en Sección~\ref{sec-bio-ic}.

\begin{codelisting}[H]

\caption{\label{lst-semana07-7}Código R}

\centering{

\begin{Shaded}
\begin{Highlighting}[]
\FunctionTok{library}\NormalTok{(BioStatR)}

\CommentTok{\# IC para la media de HbA1c — 94 pacientes diabéticos, Fac. Medicina UGR}
\FunctionTok{icm}\NormalTok{(}\AttributeTok{m =} \FunctionTok{mean}\NormalTok{(osteo}\SpecialCharTok{$}\NormalTok{hba1c), }\AttributeTok{s =} \FunctionTok{sd}\NormalTok{(osteo}\SpecialCharTok{$}\NormalTok{hba1c), }\AttributeTok{n =} \FunctionTok{nrow}\NormalTok{(osteo))}
\end{Highlighting}
\end{Shaded}

}

\end{codelisting}%

\begin{verbatim}

Intervalo de confianza bilateral para la media de una VA normal 
----------------------------------------------------------------
Información muestral: 
  Tamaño muestral: n =  94 
  Media: m =  8.565 
  Desviación típica: s =  1.799 
  Error estándar de la media: sem =  0.186 

Estimación: 
 95%-IC(µ):  (8.2, 8.93)
  Precisión obtenida: 0.368 
\end{verbatim}

\begin{codelisting}[H]

\caption{\label{lst-semana07-7}Código R}

\centering{

\begin{Shaded}
\begin{Highlighting}[]
\CommentTok{\# IC para la proporción de fumadores (41 de 94 pacientes)}
\FunctionTok{icp}\NormalTok{(}\AttributeTok{x =} \FunctionTok{sum}\NormalTok{(osteo}\SpecialCharTok{$}\NormalTok{tabaco }\SpecialCharTok{==} \StringTok{"Sí"}\NormalTok{), }\AttributeTok{n =} \FunctionTok{nrow}\NormalTok{(osteo))}
\end{Highlighting}
\end{Shaded}

}

\end{codelisting}%

\begin{verbatim}

Intervalo de confianza para una proporción binomial 
--------------------------------------------------- 

Información muestral: 
  Tamaño de muestra: n = 94
  Estimación puntual clásica: p=x/n = 0.436, q=(1-p)=0.564
  Casos observados: x = 41

# Método exacto (Clooper-Pearson):
  Pseudo-estimación puntual: p' = 0.44, q'=(1-p')=0.56
  95%-IC(π): (0.339, 0.542) 
  Semiamplitud: 0.102

# Método de Wilson (con cpc):
  Pseudo-estimación puntual: p' = 0.439, q'=(1-p')=0.561
  95%-IC(π): (0.335, 0.542) 
  Semiamplitud: 0.103

# Método de Wald (con cpc):
  Estimación puntual (clásica): p=x/n = 0.436, q=(1-p)=0.564
  95%-IC(π): (0.331, 0.542) 
  Precisión: 0.106

# Método de Wald ajustado (Agresti-Coull):
  Estimación puntual: p=(x+2)/(n+4) = 0.439, q=(1-p)=0.561
  95%-IC(π): (0.34, 0.537) 
  Precisión: 0.0982
\end{verbatim}

\texttt{icp()} calcula cuatro métodos simultáneamente. Para muestras
pequeñas se recomienda Clopper-Pearson (exacto); para muestras grandes
los cuatro coinciden.

\end{tcolorbox}

\section{\texorpdfstring{Cálculo del Tamaño Muestral: \texttt{nm()} y
\texttt{np()}}{Cálculo del Tamaño Muestral: nm() y np()}}\label{sec-semana06-tamano}

Antes de iniciar un estudio clínico es esencial determinar el
\textbf{tamaño muestral} necesario para alcanzar una precisión deseada
en los intervalos de confianza. Las funciones \texttt{nm()} y
\texttt{np()} del paquete \texttt{BioStatR} (ver
Sección~\ref{sec-bio-n}) calculan el tamaño muestral requerido para una
\textbf{media} o una \textbf{proporción}, respectivamente, dada una
semiamplitud objetivo \(d\) del IC.

\subsection{Fórmulas del Tamaño
Muestral}\label{fuxf3rmulas-del-tamauxf1o-muestral}

\textbf{Para la media} (con estimación piloto \(s\) de la desviación
típica): \[n = \left(\frac{z_{1-\alpha/2}\, s}{d}\right)^2\]

\textbf{Para la proporción} (con estimación piloto \(\hat{p} = x/n\)):
\[n = \left(\frac{z_{1-\alpha/2}}{d}\right)^2 \hat{p}\,(1-\hat{p})\]

donde \(d\) es la \textbf{semiamplitud} (precisión) deseada del IC y
\(z_{1-\alpha/2}\) es el cuantil normal estándar. La dependencia
cuadrática \(n \propto 1/d^2\) implica que reducir la semiamplitud a la
mitad cuadruplica el tamaño muestral.

\begin{tcolorbox}[enhanced jigsaw, breakable, titlerule=0mm, rightrule=.15mm, leftrule=.75mm, left=2mm, opacitybacktitle=0.6, colback=white, bottomtitle=1mm, colbacktitle=quarto-callout-tip-color!10!white, title=\textcolor{quarto-callout-tip-color}{\faLightbulb}\hspace{0.5em}{Ejemplo: Tamaño muestral para HbA1c y prevalencia de tabaquismo}, toptitle=1mm, coltitle=black, opacityback=0, arc=.35mm, bottomrule=.15mm, toprule=.15mm, colframe=quarto-callout-tip-color-frame]

Partiendo del estudio piloto \texttt{osteo} (94 pacientes diabéticos,
Fac. Medicina UGR):

\begin{itemize}
\tightlist
\item
  \textbf{HbA1c}: \(\bar{x} \approx 8.57\%\), \(s \approx 1.80\).
  ¿Cuántos pacientes hacen falta para reducir la semiamplitud del IC a
  \(d = 0.3\%\)?
\item
  \textbf{Tabaquismo}: 41/94 fumadores (\(\hat{p} \approx 0.436\)).
  ¿Cuántos pacientes para alcanzar precisión \(d = 0.05\) al 95\% de
  confianza?
\end{itemize}

\begin{codelisting}[H]

\caption{\label{lst-semana07-9}Código R}

\centering{

\begin{Shaded}
\begin{Highlighting}[]
\FunctionTok{library}\NormalTok{(BioStatR)}

\CommentTok{\# nm(): tamaño muestral para la media de HbA1c}
\CommentTok{\# d = 0.3\% (semiamplitud deseada), datos piloto del estudio osteo}
\FunctionTok{nm}\NormalTok{(}\AttributeTok{d =} \FloatTok{0.3}\NormalTok{, }\AttributeTok{n =} \FunctionTok{nrow}\NormalTok{(osteo),}
   \AttributeTok{m =} \FunctionTok{mean}\NormalTok{(osteo}\SpecialCharTok{$}\NormalTok{hba1c), }\AttributeTok{s =} \FunctionTok{sd}\NormalTok{(osteo}\SpecialCharTok{$}\NormalTok{hba1c), }\AttributeTok{alfa =} \FloatTok{0.05}\NormalTok{)}
\end{Highlighting}
\end{Shaded}

}

\end{codelisting}%

\begin{verbatim}


# Tamaño de muestra para la estimación de la media de una VA normal o su aproximación
# -----------------------------------------------------------------------------------

# Muestra piloto:
  Tamaño muestral:  n = 94 
  Media: m = 8.5649 
  Desviación típica: s = 1.7987 
  Error estandar de la media: sem = 0.1855 
  Precisión observada: d = 0.3684 

# Estimación del tamaño muestral:
  Precisión deseada: δ = 0.3000 
  Tamaño muestral necesario: n ≥ 142 
\end{verbatim}

\begin{codelisting}[H]

\caption{\label{lst-semana07-9}Código R}

\centering{

\begin{Shaded}
\begin{Highlighting}[]
\CommentTok{\# np(): tamaño muestral para la proporción de fumadores}
\CommentTok{\# d = 0.05 (semiamplitud deseada), conf = 0.95}
\FunctionTok{np}\NormalTok{(}\AttributeTok{x =} \FunctionTok{sum}\NormalTok{(osteo}\SpecialCharTok{$}\NormalTok{tabaco }\SpecialCharTok{==} \StringTok{"Sí"}\NormalTok{), }\AttributeTok{n =} \FunctionTok{nrow}\NormalTok{(osteo),}
   \AttributeTok{d =} \FloatTok{0.05}\NormalTok{, }\AttributeTok{conf =} \FloatTok{0.95}\NormalTok{)}
\end{Highlighting}
\end{Shaded}

}

\end{codelisting}%

\begin{verbatim}

Tamaño de muestra para estimar una proporción binomial 
-------------------------------------------------------

Información muestral 
  Tamaño de la muestra: n = 94
  Casos: x = 41
  Inferencia para la proporción basada en el método de Wald ajustado: 
  95%-IC(π): (0.3405, 0.5370) 
  precisión observada: d = 0.0982 (9.82%) 

Tamaño muestral requerido para δ = 0.05 (5.00%), conf.= 95% 
  No se distinguen casos con y sin información (la muestra actual es compatible con p =0.5): n= 385
\end{verbatim}

\begin{quote}
\textbf{Interpretación}

El cálculo del tamaño muestral cuantifica el coste estadístico de
incrementar la precisión. Para estimar la HbA1c media con una
semiamplitud del IC \(\leq 0.3\%\) (frente a los \(\approx 0.37\%\)
obtenidos con \(n = 94\)), se necesitarían aproximadamente \textbf{n ≈
138 pacientes diabéticos} --- un incremento del 47\%. Análogamente, para
estimar la prevalencia de tabaquismo con precisión \(\pm 5\%\) (IC 95\%)
se requerirían cerca de \textbf{n ≈ 380 pacientes}, cuatro veces el
tamaño del estudio piloto. En la fase de diseño de un protocolo clínico,
estos cálculos orientan el balance entre \textbf{precisión científica}
(estimaciones más estrechas guían mejor las decisiones terapéuticas) y
\textbf{viabilidad logística-financiera} (reclutamiento, seguimiento,
presupuesto). La dependencia cuadrática \(n \propto 1/d^2\) explica por
qué ganar un factor 2 en precisión cuadruplica el coste muestral --- un
argumento decisivo al justificar el tamaño muestral ante comités éticos
y agencias financiadoras.
\end{quote}

\end{tcolorbox}

\section{Ejercicios}\label{sec-semana06-ejercicios}

\begin{enumerate}
\def\labelenumi{\arabic{enumi}.}
\item
  \textbf{Concepto de Estimador:} Define con tus palabras la diferencia
  entre un parámetro poblacional, un estimador y un valor estimado. ¿Por
  qué es importante considerar un estimador como una variable aleatoria?
\item
  \textbf{Insesgadez vs.~Eficiencia:} Considera dos estimadores
  insesgados del mismo parámetro. ¿Cuál prefieres si uno tiene varianza
  dos veces mayor que el otro? Razona tu respuesta.
\item
  \textbf{Error Cuadrático Medio:} Explica por qué a veces puede ser
  preferible un estimador sesgado con baja varianza sobre uno insesgado
  con alta varianza. Proporciona un ejemplo numérico simple.
\item
  \textbf{Método de Momentos:} Para una muestra de una distribución
  Uniform\([0, \theta]\):

  \begin{itemize}
  \tightlist
  \item
    Encuentra el estimador por método de momentos de \(\theta\)
  \item
    ¿Es insesgado?
  \end{itemize}
\item
  \textbf{Máxima Verosimilitud:} Para una muestra de \(n\) observaciones
  de una distribución Poisson\((\lambda)\):

  \begin{itemize}
  \tightlist
  \item
    Escribe la función de verosimilitud
  \item
    Encuentra el EMV de \(\lambda\)
  \item
    ¿Coincide con el estimador MoM?
  \end{itemize}
\item
  \textbf{Intervalo de Confianza:} Se observa una muestra de 36 datos
  con \(\bar{x} = 50\) y \(s = 12\). Construye un IC 95\% para la media
  poblacional. ¿Cómo cambiaría el intervalo si usáramos un nivel de
  confianza del 99\%?
\item
  \textbf{Intervalo de Confianza para Proporción:} En una encuesta de
  500 personas, 325 responden afirmativamente a una pregunta. Construye
  un IC 95\% para la proporción poblacional e interpreta el resultado.
\item
  \textbf{Programación en R:} Escribe un script en R que:

  \begin{itemize}
  \tightlist
  \item
    Genere una muestra de una distribución
    Normal\((\mu=100, \sigma=15)\) de tamaño \(n=50\)
  \item
    Calcule estimadores puntuales para \(\mu\) y \(\sigma^2\)
  \item
    Construya un IC 95\% para \(\mu\)
  \item
    Verifique si el verdadero valor de \(\mu\) cae dentro del intervalo
  \end{itemize}
\end{enumerate}

\section{Respuestas a los Ejercicios}\label{sec-semana06-respuestas}

\textbf{Ejercicio 1:} Un parámetro es el valor verdadero desconocido de
la población (θ). Un estimador es una función de los datos muestrales
(T(X₁,\ldots,Xₙ)) que genera diferentes valores según la muestra. Es una
variable aleatoria porque depende de la muestra. Es importante reconocer
esto para entender su distribución y propiedades.

\textbf{Ejercicio 2:} Prefieres el estimador con menor varianza porque
tendrá una distribución más concentrada alrededor del parámetro. Si
ambos son insesgados, el de menor varianza es más eficiente.

\textbf{Ejercicio 3:} Ejemplo: θ=0, Estimador A: insesgado, Var=100;
Estimador B: sesgo=2, Var=10. ECM\_A = 100; ECM\_B = 4 + 10 = 14. B es
mejor a pesar del sesgo.

\textbf{Ejercicio 4:} Para Uniform{[}0,θ{]}, el método de momentos da θ̂
= 2X̄. No es insesgado (subestima).

\textbf{Ejercicio 5:} L(λ) = λⁿe\^{}(-nλ)∏xᵢ!/∏xᵢ!. EMV: λ̂ = X̄. Sí
coincide con MoM.

\textbf{Ejercicio 6:} n=36, X̄=50, s=12, t₀.₀₂₅,₃₅ ≈ 2.03 IC95\% = {[}50
± 2.03×(12/6){]} = {[}50 ± 4.06{]} = {[}45.94, 54.06{]} IC99\%: más
amplio ≈ {[}44.35, 55.65{]} (usa t₀.₀₀₅,₃₅ ≈ 2.72)

\textbf{Ejercicio 7:} p̂ = 325/500 = 0.65, SE = √(0.65×0.35/500) ≈ 0.0214
IC95\% = {[}0.65 ± 1.96×0.0214{]} ≈ {[}0.608, 0.692{]}

\textbf{Ejercicio 8:} Código R básico:

\begin{Shaded}
\begin{Highlighting}[]
\FunctionTok{set.seed}\NormalTok{(}\DecValTok{123}\NormalTok{)}
\NormalTok{muestra }\OtherTok{\textless{}{-}} \FunctionTok{rnorm}\NormalTok{(}\DecValTok{50}\NormalTok{, }\AttributeTok{mean=}\DecValTok{100}\NormalTok{, }\AttributeTok{sd=}\DecValTok{15}\NormalTok{)}
\NormalTok{media }\OtherTok{\textless{}{-}} \FunctionTok{mean}\NormalTok{(muestra)}
\NormalTok{var\_est }\OtherTok{\textless{}{-}} \FunctionTok{var}\NormalTok{(muestra)}
\NormalTok{se }\OtherTok{\textless{}{-}} \FunctionTok{sd}\NormalTok{(muestra)}\SpecialCharTok{/}\FunctionTok{sqrt}\NormalTok{(}\DecValTok{50}\NormalTok{)}
\NormalTok{ic }\OtherTok{\textless{}{-}} \FunctionTok{c}\NormalTok{(media }\SpecialCharTok{{-}} \FloatTok{1.96}\SpecialCharTok{*}\NormalTok{se, media }\SpecialCharTok{+} \FloatTok{1.96}\SpecialCharTok{*}\NormalTok{se)}
\CommentTok{\# Verificar: 100 ∈ ic?}
\end{Highlighting}
\end{Shaded}

\chapter{Semana 7 --- Contrastes de Hipótesis (Parte
I)}\label{sec-semana07}

\section{Introducción}\label{introducciuxf3n-1}

Un \textbf{contraste de hipótesis} (también llamado \emph{prueba de
hipótesis}) es un procedimiento estadístico para tomar una decisión
entre dos afirmaciones contrapuestas acerca de un parámetro desconocido
de la población. Utilizamos información de una muestra aleatoria para
determinar si la evidencia respalda rechazar o no una afirmación inicial
(hipótesis nula).

\textbf{Ejemplo motivador:} Un productor de monedas afirma que sus
monedas son justas (probabilidad de cara = 0.5). ¿Cómo comprobamos esta
afirmación? Si lanzamos la moneda 100 veces y obtenemos 95 caras, esto
sugiere que la moneda no es justa. Pero ¿y si obtenemos 52 caras? El
contraste de hipótesis nos proporciona un procedimiento riguroso para
tomar esta decisión considerando la aleatoriedad de los datos.

\begin{center}\rule{0.5\linewidth}{0.5pt}\end{center}

\section{Hipótesis Nula y Alternativa}\label{sec-semana07-s1}

En un contraste de hipótesis distinguimos dos hipótesis complementarias:

\begin{tcolorbox}[enhanced jigsaw, breakable, titlerule=0mm, rightrule=.15mm, leftrule=.75mm, left=2mm, opacitybacktitle=0.6, colback=white, bottomtitle=1mm, colbacktitle=quarto-callout-note-color!10!white, title=\textcolor{quarto-callout-note-color}{\faInfo}\hspace{0.5em}{Definición: Hipótesis Nula (\(H_0\)) e Hipótesis Alternativa (\(H_1\))}, toptitle=1mm, coltitle=black, opacityback=0, arc=.35mm, bottomrule=.15mm, toprule=.15mm, colframe=quarto-callout-note-color-frame]

\textbf{Hipótesis Nula (\(H_0\)):} Es la afirmación que asumimos como
cierta por defecto. Representa el status quo o la ausencia de efecto.
Siempre contiene una igualdad.

\textbf{Hipótesis Alternativa (\(H_1\)):} Es la afirmación que queremos
demostrar o investigar. Representa la desviación del valor hipotético en
\(H_0\).

Dos hipótesis son \textbf{complementarias} y \textbf{mutuamente
excluyentes}: si una es cierta, la otra es falsa. Capturan todos los
valores posibles del parámetro.

\end{tcolorbox}

\subsection{Tipos de Hipótesis}\label{tipos-de-hipuxf3tesis}

\textbf{Hipótesis Simple vs.~Compuesta:} - Una hipótesis es
\textbf{simple} si especifica un único valor del parámetro (p.ej.,
\(H_0: \mu = 100\)) - Una hipótesis es \textbf{compuesta} si especifica
múltiples valores (p.ej., \(H_1: \mu > 100\))

\textbf{Contrastes Bilaterales vs.~Unilaterales:}

\begin{longtable}[]{@{}
  >{\raggedright\arraybackslash}p{(\linewidth - 6\tabcolsep) * \real{0.0968}}
  >{\raggedright\arraybackslash}p{(\linewidth - 6\tabcolsep) * \real{0.2903}}
  >{\raggedright\arraybackslash}p{(\linewidth - 6\tabcolsep) * \real{0.4032}}
  >{\raggedright\arraybackslash}p{(\linewidth - 6\tabcolsep) * \real{0.2097}}@{}}
\toprule\noalign{}
\begin{minipage}[b]{\linewidth}\raggedright
Tipo
\end{minipage} & \begin{minipage}[b]{\linewidth}\raggedright
Hipótesis Nula
\end{minipage} & \begin{minipage}[b]{\linewidth}\raggedright
Hipótesis Alternativa
\end{minipage} & \begin{minipage}[b]{\linewidth}\raggedright
Descripción
\end{minipage} \\
\midrule\noalign{}
\endhead
\bottomrule\noalign{}
\endlastfoot
\textbf{Bilateral} & \(H_0: \theta = \theta_0\) &
\(H_1: \theta \neq \theta_0\) & Investigamos desviaciones en ambas
direcciones \\
\textbf{Unilateral derecho} & \(H_0: \theta \leq \theta_0\) &
\(H_1: \theta > \theta_0\) & Investigamos si el parámetro es mayor \\
\textbf{Unilateral izquierdo} & \(H_0: \theta \geq \theta_0\) &
\(H_1: \theta < \theta_0\) & Investigamos si el parámetro es menor \\
\end{longtable}

\subsection{Ejemplos}\label{ejemplos}

\begin{enumerate}
\def\labelenumi{\arabic{enumi}.}
\tightlist
\item
  \textbf{Prueba de moneda justa:} \(H_0: \pi = 0.5\)
  vs.~\(H_1: \pi \neq 0.5\) (bilateral)
\item
  \textbf{Control de calidad:} \(H_0: \mu \leq 1000\)
  vs.~\(H_1: \mu > 1000\) (unilateral derecho)
\item
  \textbf{Seguridad de medicamento:} \(H_0: \pi \geq 0.05\)
  vs.~\(H_1: \pi < 0.05\) (unilateral izquierdo, donde \(\pi\) es tasa
  de efectos adversos)
\end{enumerate}

\begin{center}\rule{0.5\linewidth}{0.5pt}\end{center}

\section{Errores Tipo I y Tipo II}\label{sec-semana07-s2}

Cuando tomamos una decisión basada en datos muestrales, podemos cometer
errores. La verdad sobre \(H_0\) es desconocida, pero existen dos
posibilidades:

\begin{tcolorbox}[enhanced jigsaw, breakable, titlerule=0mm, rightrule=.15mm, leftrule=.75mm, left=2mm, opacitybacktitle=0.6, colback=white, bottomtitle=1mm, colbacktitle=quarto-callout-note-color!10!white, title=\textcolor{quarto-callout-note-color}{\faInfo}\hspace{0.5em}{Definición: Errores Tipo I y Tipo II}, toptitle=1mm, coltitle=black, opacityback=0, arc=.35mm, bottomrule=.15mm, toprule=.15mm, colframe=quarto-callout-note-color-frame]

\textbf{Error Tipo I (Falso Positivo):} Rechazar \(H_0\) cuando en
realidad es verdadera. - Probabilidad:
\(P(\text{Rechazar } H_0 | H_0 \text{ es verdadera}) = \alpha\) - Nivel
de significación

\textbf{Error Tipo II (Falso Negativo):} No rechazar \(H_0\) cuando en
realidad es falsa. - Probabilidad:
\(P(\text{No rechazar } H_0 | H_1 \text{ es verdadera}) = \beta\)

\textbf{Potencia de la prueba:} Probabilidad de rechazar \(H_0\) cuando
es falsa. - Potencia = \(1 - \beta\) - Representa nuestra capacidad de
detectar un efecto verdadero

\end{tcolorbox}

\subsection{Tabla de Decisión}\label{tabla-de-decisiuxf3n}

\begin{longtable}[]{@{}
  >{\raggedright\arraybackslash}p{(\linewidth - 4\tabcolsep) * \real{0.2500}}
  >{\raggedright\arraybackslash}p{(\linewidth - 4\tabcolsep) * \real{0.3750}}
  >{\raggedright\arraybackslash}p{(\linewidth - 4\tabcolsep) * \real{0.3750}}@{}}
\toprule\noalign{}
\begin{minipage}[b]{\linewidth}\raggedright
\end{minipage} & \begin{minipage}[b]{\linewidth}\raggedright
\(H_0\) es Verdadera
\end{minipage} & \begin{minipage}[b]{\linewidth}\raggedright
\(H_0\) es Falsa
\end{minipage} \\
\midrule\noalign{}
\endhead
\bottomrule\noalign{}
\endlastfoot
\textbf{Rechazamos \(H_0\)} & Error Tipo I (probabilidad \(\alpha\)) & ✓
Decisión Correcta (probabilidad \(1-\beta\)) \\
\textbf{No Rechazamos \(H_0\)} & ✓ Decisión Correcta (probabilidad
\(1-\alpha\)) & Error Tipo II (probabilidad \(\beta\)) \\
\end{longtable}

\subsection{Analogía: El Pronóstico
Médico}\label{analoguxeda-el-pronuxf3stico-muxe9dico}

La estructura de un contraste de hipótesis puede entenderse con una
analogía clínica sobre el pronóstico de un paciente:

\begin{itemize}
\tightlist
\item
  \textbf{\(H_0\):} El paciente va a morir (pronóstico fatal asumido
  como hipótesis nula)
\item
  \textbf{\(H_1\):} El paciente va a vivir
\item
  \textbf{Error Tipo I:} Le dices que \textbf{vive}, pero \textbf{muere}
  (rechazas \(H_0\) siendo cierta)
\item
  \textbf{Error Tipo II:} Le dices que \textbf{muere}, pero
  \textbf{vive} (no rechazas \(H_0\) siendo falsa)
\end{itemize}

Esta analogía ilustra cómo las consecuencias de los errores definen la
gravedad del contraste. En este escenario, el Error Tipo I (falsa
esperanza) tiene consecuencias emocionales y logísticas muy distintas al
Error Tipo II (falsa alarma).

\subsection{\texorpdfstring{Relación entre \(\alpha\) y
\(\beta\)}{Relación entre \textbackslash alpha y \textbackslash beta}}\label{relaciuxf3n-entre-alpha-y-beta}

\textbf{Observación importante:} Para una muestra de tamaño \(n\) fijo,
\textbf{existe una relación inversa entre \(\alpha\) y \(\beta\)}:

\begin{itemize}
\tightlist
\item
  Si bajamos \(\alpha\) (menos riesgo de rechazar \(H_0\) falsamente),
  entonces \(\beta\) aumenta (mayor riesgo de no rechazar \(H_0\)
  falsamente)
\item
  Si bajamos \(\beta\) (mejor capacidad de detectar un efecto), entonces
  \(\alpha\) aumenta
\item
  La única forma de reducir ambos simultáneamente es \textbf{aumentar el
  tamaño de la muestra \(n\)}
\end{itemize}

\begin{center}\rule{0.5\linewidth}{0.5pt}\end{center}

\section{Estadístico de Prueba y Región Crítica}\label{sec-semana07-s3}

Para ejecutar un contraste de hipótesis necesitamos tres ingredientes:

\begin{tcolorbox}[enhanced jigsaw, breakable, titlerule=0mm, rightrule=.15mm, leftrule=.75mm, left=2mm, opacitybacktitle=0.6, colback=white, bottomtitle=1mm, colbacktitle=quarto-callout-note-color!10!white, title=\textcolor{quarto-callout-note-color}{\faInfo}\hspace{0.5em}{Definición: Componentes de un Contraste}, toptitle=1mm, coltitle=black, opacityback=0, arc=.35mm, bottomrule=.15mm, toprule=.15mm, colframe=quarto-callout-note-color-frame]

\textbf{1. Estadístico de Prueba (\(T_n\)):} Una función de los datos
muestrales cuya distribución es conocida bajo \(H_0\). Se calcula de la
muestra y resume la evidencia contra \(H_0\).

\textbf{2. Distribución bajo \(H_0\):} La distribución de probabilidad
del estadístico de prueba asumiendo que \(H_0\) es verdadera.

\textbf{3. Región de Rechazo (o Región Crítica):} El conjunto de valores
del estadístico de prueba para los cuales rechazamos \(H_0\). Se define
antes de recopilar datos.

\textbf{Valor Crítico:} El punto que separa la región de rechazo de la
región de no rechazo.

\end{tcolorbox}

\subsection{Procedimiento de
Decisión}\label{procedimiento-de-decisiuxf3n}

\begin{enumerate}
\def\labelenumi{\arabic{enumi}.}
\tightlist
\item
  Calcular el valor observado del estadístico de prueba:
  \(t_{obs} = T(x_1, \ldots, x_n)\)
\item
  Comparar \(t_{obs}\) con la región crítica
\item
  Si \(t_{obs}\) está en la región crítica \(\Rightarrow\) Rechazamos
  \(H_0\)
\item
  Si \(t_{obs}\) NO está en la región crítica \(\Rightarrow\) No
  rechazamos \(H_0\)
\end{enumerate}

\begin{center}\rule{0.5\linewidth}{0.5pt}\end{center}

\section{Función de Potencia}\label{sec-semana07-s4}

\begin{tcolorbox}[enhanced jigsaw, breakable, titlerule=0mm, rightrule=.15mm, leftrule=.75mm, left=2mm, opacitybacktitle=0.6, colback=white, bottomtitle=1mm, colbacktitle=quarto-callout-note-color!10!white, title=\textcolor{quarto-callout-note-color}{\faInfo}\hspace{0.5em}{Definición: Función de Potencia}, toptitle=1mm, coltitle=black, opacityback=0, arc=.35mm, bottomrule=.15mm, toprule=.15mm, colframe=quarto-callout-note-color-frame]

La \textbf{función de potencia} es una función de \(\theta\) que mide la
probabilidad de rechazar \(H_0\) para cada valor posible del parámetro:

\[\beta(\theta) = P_\theta(T_n \in R)\]

donde \(R\) es la región de rechazo y \(P_\theta\) significa
``probabilidad cuando el verdadero valor del parámetro es \(\theta\)''.

\textbf{Interpretación:} - Cuando \(\theta \in \Theta_0\) (bajo \(H_0\)
verdadera): \(\beta(\theta)\) es la probabilidad de Error Tipo I -
Cuando \(\theta \in \Theta_1\) (bajo \(H_1\) verdadera):
\(\beta(\theta) = 1 - P(\text{Error Tipo II})\) es la \textbf{potencia}

Un buen contraste tiene: - Potencia cercana a 0 para
\(\theta \in \Theta_0\) (pocos falsos positivos) - Potencia cercana a 1
para \(\theta \in \Theta_1\) (muchos verdaderos positivos)

\end{tcolorbox}

\begin{center}\rule{0.5\linewidth}{0.5pt}\end{center}

\section{Tamaño y Nivel de una Prueba}\label{sec-semana07-s5}

\begin{tcolorbox}[enhanced jigsaw, breakable, titlerule=0mm, rightrule=.15mm, leftrule=.75mm, left=2mm, opacitybacktitle=0.6, colback=white, bottomtitle=1mm, colbacktitle=quarto-callout-note-color!10!white, title=\textcolor{quarto-callout-note-color}{\faInfo}\hspace{0.5em}{Definición: Tamaño y Nivel}, toptitle=1mm, coltitle=black, opacityback=0, arc=.35mm, bottomrule=.15mm, toprule=.15mm, colframe=quarto-callout-note-color-frame]

\textbf{Tamaño de la prueba (\(\alpha\)):} Es el supremo de la función
de potencia sobre \(\Theta_0\):

\[\text{Tamaño} = \sup_{\theta \in \Theta_0} \beta(\theta)\]

En otras palabras, es la máxima probabilidad de Error Tipo I sobre todos
los valores en \(H_0\).

\textbf{Nivel de la prueba (\(\alpha\)):} Es un número entre 0 y 1
(típicamente 0.05, 0.01) tal que:

\[\sup_{\theta \in \Theta_0} \beta(\theta) \leq \alpha\]

La prueba se dice que es de ``nivel \(\alpha\)'' si controla el Error
Tipo I en este nivel.

\end{tcolorbox}

\subsection{Relación con la
Significación}\label{relaciuxf3n-con-la-significaciuxf3n}

\begin{itemize}
\tightlist
\item
  El \textbf{nivel de significación} \(\alpha\) es el umbral máximo de
  riesgo de Error Tipo I que estamos dispuestos a aceptar
\item
  Típicamente: \(\alpha = 0.05\) (5\%) o \(\alpha = 0.01\) (1\%)
\item
  Se elige \textbf{antes} de realizar el análisis
\item
  Una vez elegido, el valor crítico se determina de la distribución bajo
  \(H_0\) de forma que se cumpla el nivel
\end{itemize}

\begin{center}\rule{0.5\linewidth}{0.5pt}\end{center}

\section{Contraste Z para la Media (Varianza
Conocida)}\label{sec-semana07-s6}

Cuando la varianza poblacional \(\sigma^2\) es \textbf{conocida}
(situación rara en la práctica), usamos el contraste Z basado en la
distribución normal estándar.

\subsection{Condiciones}\label{condiciones}

\begin{itemize}
\tightlist
\item
  Muestra aleatoria de tamaño \(n\)
\item
  Datos \(X_i \sim N(\mu, \sigma^2)\) con \(\sigma^2\)
  \textbf{conocida}, O
\item
  Datos arbitrarios con \(n\) grande (Central Limit Theorem)
\item
  Queremos probar hipótesis sobre \(\mu\)
\end{itemize}

\subsection{Estadístico de Prueba}\label{estaduxedstico-de-prueba}

\begin{tcolorbox}[enhanced jigsaw, breakable, titlerule=0mm, rightrule=.15mm, leftrule=.75mm, left=2mm, opacitybacktitle=0.6, colback=white, bottomtitle=1mm, colbacktitle=quarto-callout-warning-color!10!white, title=\textcolor{quarto-callout-warning-color}{\faExclamationTriangle}\hspace{0.5em}{Resultado Importante: Contraste Z}, toptitle=1mm, coltitle=black, opacityback=0, arc=.35mm, bottomrule=.15mm, toprule=.15mm, colframe=quarto-callout-warning-color-frame]

Bajo \(H_0: \mu = \mu_0\) con \(\sigma\) conocida:

\[Z = \frac{\bar{X} - \mu_0}{\sigma/\sqrt{n}} \sim N(0, 1)\]

El estadístico \(Z\) sigue una distribución normal estándar cuando
\(H_0\) es verdadera.

\end{tcolorbox}

\subsection{Regiones de Rechazo y Valores
Críticos}\label{regiones-de-rechazo-y-valores-cruxedticos}

Para un nivel de significación \(\alpha\):

\begin{longtable}[]{@{}
  >{\raggedright\arraybackslash}p{(\linewidth - 6\tabcolsep) * \real{0.1250}}
  >{\raggedright\arraybackslash}p{(\linewidth - 6\tabcolsep) * \real{0.3958}}
  >{\raggedright\arraybackslash}p{(\linewidth - 6\tabcolsep) * \real{0.2917}}
  >{\raggedright\arraybackslash}p{(\linewidth - 6\tabcolsep) * \real{0.1875}}@{}}
\toprule\noalign{}
\begin{minipage}[b]{\linewidth}\raggedright
Tipo
\end{minipage} & \begin{minipage}[b]{\linewidth}\raggedright
Región de Rechazo
\end{minipage} & \begin{minipage}[b]{\linewidth}\raggedright
Valor Crítico
\end{minipage} & \begin{minipage}[b]{\linewidth}\raggedright
p-valor
\end{minipage} \\
\midrule\noalign{}
\endhead
\bottomrule\noalign{}
\endlastfoot
\textbf{Bilateral} & \(\|Z\| > z_{1-\alpha/2}\) & \(z_{1-\alpha/2}\) &
\(2P(Z > \|z_{obs}\|)\) \\
\textbf{Unilateral derecha} & \(Z > z_{1-\alpha}\) & \(z_{1-\alpha}\) &
\(P(Z > z_{obs})\) \\
\textbf{Unilateral izquierda} & \(Z < -z_{1-\alpha}\) &
\(-z_{1-\alpha}\) & \(P(Z < z_{obs})\) \\
\end{longtable}

\subsection{Valores Críticos Comunes (Normal
Estándar)}\label{valores-cruxedticos-comunes-normal-estuxe1ndar}

\begin{longtable}[]{@{}llllll@{}}
\toprule\noalign{}
Cuantil & 0.900 & 0.950 & 0.975 & 0.990 & 0.995 \\
\midrule\noalign{}
\endhead
\bottomrule\noalign{}
\endlastfoot
\(z_q\) & 1.28 & 1.64 & 1.96 & 2.33 & 2.58 \\
\end{longtable}

\subsection{Ejemplo: Productor de
Harina}\label{ejemplo-productor-de-harina}

\textbf{Problema:} Un productor de harina quiere verificar si sus
paquetes contienen exactamente 1 kg de harina. - Si hay mucho menos: los
clientes se quejan - Si hay mucho más: pérdida de ganancias - Se sabe
que: \(X_i \sim N(\mu, \sigma^2 = 25 \text{ g}^2)\)

\textbf{Hipótesis:} - \(H_0: \mu = 1000\) g - \(H_1: \mu \neq 1000\) g
(bilateral)

\textbf{Muestra:} \(n = 25\), \(\bar{x} = 1002\) g

\textbf{Cálculo:}
\[z_{obs} = \frac{1002 - 1000}{5/\sqrt{25}} = \frac{2}{1} = 2\]

\textbf{Decisión (con \(\alpha = 0.05\)):} - Valor crítico:
\(z_{1-0.05/2} = z_{0.975} = 1.96\) - Como \(|z_{obs}| = 2 > 1.96\),
rechazamos \(H_0\) - p-valor: \(2P(Z > 2) \approx 0.0455\)

\textbf{Conclusión:} A nivel de significación 5\%, la evidencia sugiere
que el peso medio es diferente de 1 kg.

\begin{center}\rule{0.5\linewidth}{0.5pt}\end{center}

\section{Contraste t para la Media (Varianza
Desconocida)}\label{sec-semana07-s7}

En la práctica, la varianza \(\sigma^2\) es \textbf{desconocida}. En
este caso, estimamos \(\sigma^2\) con \(S^2\) y usamos la distribución t
de Student.

\subsection{Condiciones}\label{condiciones-1}

\begin{itemize}
\tightlist
\item
  Muestra aleatoria de tamaño \(n\)
\item
  Datos \(X_i \sim N(\mu, \sigma^2)\) con \(\sigma^2\)
  \textbf{desconocida}, O
\item
  Datos aproximadamente normales
\item
  Queremos probar hipótesis sobre \(\mu\)
\end{itemize}

\subsection{Estimador de la Varianza}\label{estimador-de-la-varianza}

La varianza muestral insesgada es:

\[S^2 = \frac{\sum_{i=1}^n (X_i - \bar{X})^2}{n-1}\]

El divisor es \(n-1\) (no \(n\)) para obtener un estimador insesgado.

\subsection{Estadístico de Prueba y
Distribución}\label{estaduxedstico-de-prueba-y-distribuciuxf3n}

\begin{tcolorbox}[enhanced jigsaw, breakable, titlerule=0mm, rightrule=.15mm, leftrule=.75mm, left=2mm, opacitybacktitle=0.6, colback=white, bottomtitle=1mm, colbacktitle=quarto-callout-warning-color!10!white, title=\textcolor{quarto-callout-warning-color}{\faExclamationTriangle}\hspace{0.5em}{Resultado Importante: Contraste t}, toptitle=1mm, coltitle=black, opacityback=0, arc=.35mm, bottomrule=.15mm, toprule=.15mm, colframe=quarto-callout-warning-color-frame]

Bajo \(H_0: \mu = \mu_0\) con \(\sigma\) desconocida:

\[T = \frac{\bar{X} - \mu_0}{S/\sqrt{n}} \sim t_{n-1}\]

El estadístico \(T\) sigue una distribución t de Student con \(n-1\)
grados de libertad cuando \(H_0\) es verdadera.

\end{tcolorbox}

\subsection{Derivación de la
Distribución}\label{derivaciuxf3n-de-la-distribuciuxf3n}

Empezamos con:
\[T = \frac{\bar{X} - \mu_0}{S/\sqrt{n}} = \frac{\frac{\bar{X} - \mu_0}{\sigma/\sqrt{n}}}{\frac{S}{\sigma}} = \frac{Z}{\sqrt{\frac{C}{n-1}}}\]

donde: - \(Z = \frac{\bar{X} - \mu_0}{\sigma/\sqrt{n}} \sim N(0,1)\)
(normal estándar) - \(C = \frac{(n-1)S^2}{\sigma^2} \sim \chi^2_{n-1}\)
(chi-cuadrado con \(n-1\) gl)

Por definición de la distribución t: Si \(Z \sim N(0,1)\) y
\(C \sim \chi^2_k\) son independientes, entonces
\(\frac{Z}{\sqrt{C/k}} \sim t_k\).

Por lo tanto: \(T \sim t_{n-1}\)

\subsection{Regiones de Rechazo y Valores
Críticos}\label{regiones-de-rechazo-y-valores-cruxedticos-1}

Para un nivel de significación \(\alpha\):

\begin{longtable}[]{@{}
  >{\raggedright\arraybackslash}p{(\linewidth - 6\tabcolsep) * \real{0.1250}}
  >{\raggedright\arraybackslash}p{(\linewidth - 6\tabcolsep) * \real{0.3958}}
  >{\raggedright\arraybackslash}p{(\linewidth - 6\tabcolsep) * \real{0.2917}}
  >{\raggedright\arraybackslash}p{(\linewidth - 6\tabcolsep) * \real{0.1875}}@{}}
\toprule\noalign{}
\begin{minipage}[b]{\linewidth}\raggedright
Tipo
\end{minipage} & \begin{minipage}[b]{\linewidth}\raggedright
Región de Rechazo
\end{minipage} & \begin{minipage}[b]{\linewidth}\raggedright
Valor Crítico
\end{minipage} & \begin{minipage}[b]{\linewidth}\raggedright
p-valor
\end{minipage} \\
\midrule\noalign{}
\endhead
\bottomrule\noalign{}
\endlastfoot
\textbf{Bilateral} & \(\|T\| > t_{n-1, 1-\alpha/2}\) &
\(t_{n-1, 1-\alpha/2}\) & \(2P(T_{n-1} > \|t_{obs}\|)\) \\
\textbf{Unilateral derecha} & \(T > t_{n-1, 1-\alpha}\) &
\(t_{n-1, 1-\alpha}\) & \(P(T_{n-1} > t_{obs})\) \\
\textbf{Unilateral izquierda} & \(T < -t_{n-1, 1-\alpha}\) &
\(-t_{n-1, 1-\alpha}\) & \(P(T_{n-1} < t_{obs})\) \\
\end{longtable}

\subsection{Tabla de Cuantiles t de
Student}\label{tabla-de-cuantiles-t-de-student}

\begin{longtable}[]{@{}llllll@{}}
\toprule\noalign{}
Grados de Libertad & 0.900 & 0.950 & 0.975 & 0.990 & 0.995 \\
\midrule\noalign{}
\endhead
\bottomrule\noalign{}
\endlastfoot
\(df = 10\) & 1.37 & 1.81 & 2.23 & 2.76 & 3.17 \\
\(df = 30\) & 1.31 & 1.70 & 2.04 & 2.46 & 2.75 \\
\(df = \infty\) & 1.28 & 1.64 & 1.96 & 2.33 & 2.58 \\
\end{longtable}

\textbf{Observación:} Conforme aumenta \(n\), la distribución t converge
a la normal estándar. Para \(n > 30\), son prácticamente equivalentes.

\subsection{Ejemplo: Una Muestra t (Productor de Harina
Revisado)}\label{ejemplo-una-muestra-t-productor-de-harina-revisado}

\textbf{Hipótesis:} Mismas que antes - \(H_0: \mu = 1000\) g -
\(H_1: \mu \neq 1000\) g

\textbf{Muestra:} \(n = 25\), \(\bar{x} = 1002\) g, \(s = 4.9\) g

\textbf{Cálculo:}
\[t_{obs} = \frac{1002 - 1000}{4.9/\sqrt{25}} = \frac{2}{0.98} = 2.04\]

\textbf{Decisión (con \(\alpha = 0.05\)):} - Grados de libertad:
\(df = 24\) - Valor crítico: \(t_{24, 0.975} = 2.064\) - Como
\(|t_{obs}| = 2.04 < 2.064\), \textbf{no rechazamos} \(H_0\) - p-valor:
\(2P(T_{24} > 2.04) \approx 0.0525\)

\textbf{Conclusión:} Con estos datos, no tenemos evidencia suficiente
(al 5\%) de que el peso medio sea diferente de 1 kg.

\begin{tcolorbox}[enhanced jigsaw, breakable, titlerule=0mm, rightrule=.15mm, leftrule=.75mm, left=2mm, opacitybacktitle=0.6, colback=white, bottomtitle=1mm, colbacktitle=quarto-callout-tip-color!10!white, title=\textcolor{quarto-callout-tip-color}{\faLightbulb}\hspace{0.5em}{Código R: Contraste t de una muestra}, toptitle=1mm, coltitle=black, opacityback=0, arc=.35mm, bottomrule=.15mm, toprule=.15mm, colframe=quarto-callout-tip-color-frame]

\begin{codelisting}[H]

\caption{\label{lst-semana08-1}Código R}

\centering{

\begin{Shaded}
\begin{Highlighting}[]
\CommentTok{\# Datos de pesos de paquetes de harina (en gramos)}
\NormalTok{datos }\OtherTok{\textless{}{-}} \FunctionTok{c}\NormalTok{(}\DecValTok{1002}\NormalTok{, }\DecValTok{998}\NormalTok{, }\DecValTok{1001}\NormalTok{, }\DecValTok{1003}\NormalTok{, }\DecValTok{1000}\NormalTok{, }\DecValTok{1004}\NormalTok{, }\DecValTok{999}\NormalTok{, }\DecValTok{1005}\NormalTok{, }\DecValTok{1002}\NormalTok{, }\DecValTok{997}\NormalTok{,}
           \DecValTok{1001}\NormalTok{, }\DecValTok{1003}\NormalTok{, }\DecValTok{1000}\NormalTok{, }\DecValTok{1002}\NormalTok{, }\DecValTok{998}\NormalTok{, }\DecValTok{1004}\NormalTok{, }\DecValTok{1001}\NormalTok{, }\DecValTok{999}\NormalTok{, }\DecValTok{1003}\NormalTok{, }\DecValTok{1000}\NormalTok{,}
           \DecValTok{1002}\NormalTok{, }\DecValTok{1001}\NormalTok{, }\DecValTok{1000}\NormalTok{, }\DecValTok{1002}\NormalTok{, }\DecValTok{998}\NormalTok{)}

\CommentTok{\# Contraste: H0: mu = 1000 vs H1: mu != 1000}
\NormalTok{resultado }\OtherTok{\textless{}{-}} \FunctionTok{t.test}\NormalTok{(datos, }\AttributeTok{mu =} \DecValTok{1000}\NormalTok{, }\AttributeTok{alternative =} \StringTok{"two.sided"}\NormalTok{)}

\CommentTok{\# Ver resultados}
\FunctionTok{print}\NormalTok{(resultado)}
\end{Highlighting}
\end{Shaded}

}

\end{codelisting}%

\begin{verbatim}

    One Sample t-test

data:  datos
t = 2, df = 24, p-value = 0.03
alternative hypothesis: true mean is not equal to 1000
95 percent confidence interval:
 1000 1002
sample estimates:
mean of x 
     1001 
\end{verbatim}

\begin{codelisting}[H]

\caption{\label{lst-semana08-1}Código R}

\centering{

\begin{Shaded}
\begin{Highlighting}[]
\CommentTok{\# Extraer componentes}
\FunctionTok{cat}\NormalTok{(}\StringTok{"Estadístico t ="}\NormalTok{, }\FunctionTok{round}\NormalTok{(resultado}\SpecialCharTok{$}\NormalTok{statistic, }\DecValTok{3}\NormalTok{), }\StringTok{"}\SpecialCharTok{\textbackslash{}n}\StringTok{"}\NormalTok{)}
\end{Highlighting}
\end{Shaded}

}

\end{codelisting}%

\begin{verbatim}
Estadístico t = 2.38 
\end{verbatim}

\begin{codelisting}[H]

\caption{\label{lst-semana08-1}Código R}

\centering{

\begin{Shaded}
\begin{Highlighting}[]
\FunctionTok{cat}\NormalTok{(}\StringTok{"Grados de libertad ="}\NormalTok{, resultado}\SpecialCharTok{$}\NormalTok{parameter, }\StringTok{"}\SpecialCharTok{\textbackslash{}n}\StringTok{"}\NormalTok{)}
\end{Highlighting}
\end{Shaded}

}

\end{codelisting}%

\begin{verbatim}
Grados de libertad = 24 
\end{verbatim}

\begin{codelisting}[H]

\caption{\label{lst-semana08-1}Código R}

\centering{

\begin{Shaded}
\begin{Highlighting}[]
\FunctionTok{cat}\NormalTok{(}\StringTok{"p{-}valor ="}\NormalTok{, }\FunctionTok{round}\NormalTok{(resultado}\SpecialCharTok{$}\NormalTok{p.value, }\DecValTok{4}\NormalTok{), }\StringTok{"}\SpecialCharTok{\textbackslash{}n}\StringTok{"}\NormalTok{)}
\end{Highlighting}
\end{Shaded}

}

\end{codelisting}%

\begin{verbatim}
p-valor = 0.0257 
\end{verbatim}

\begin{codelisting}[H]

\caption{\label{lst-semana08-1}Código R}

\centering{

\begin{Shaded}
\begin{Highlighting}[]
\FunctionTok{cat}\NormalTok{(}\StringTok{"IC 95\% para mu:"}\NormalTok{, }\FunctionTok{round}\NormalTok{(resultado}\SpecialCharTok{$}\NormalTok{conf.int, }\DecValTok{2}\NormalTok{), }\StringTok{"}\SpecialCharTok{\textbackslash{}n}\StringTok{"}\NormalTok{)}
\end{Highlighting}
\end{Shaded}

}

\end{codelisting}%

\begin{verbatim}
IC 95% para mu: 1000 1002 
\end{verbatim}

\begin{codelisting}[H]

\caption{\label{lst-semana08-1}Código R}

\centering{

\begin{Shaded}
\begin{Highlighting}[]
\CommentTok{\# Unilateral derecha: H1: mu \textgreater{} 1000}
\FunctionTok{t.test}\NormalTok{(datos, }\AttributeTok{mu =} \DecValTok{1000}\NormalTok{, }\AttributeTok{alternative =} \StringTok{"greater"}\NormalTok{)}
\end{Highlighting}
\end{Shaded}

}

\end{codelisting}%

\begin{verbatim}

    One Sample t-test

data:  datos
t = 2, df = 24, p-value = 0.01
alternative hypothesis: true mean is greater than 1000
95 percent confidence interval:
 1000  Inf
sample estimates:
mean of x 
     1001 
\end{verbatim}

\begin{codelisting}[H]

\caption{\label{lst-semana08-1}Código R}

\centering{

\begin{Shaded}
\begin{Highlighting}[]
\CommentTok{\# Unilateral izquierda: H1: mu \textless{} 1000}
\FunctionTok{t.test}\NormalTok{(datos, }\AttributeTok{mu =} \DecValTok{1000}\NormalTok{, }\AttributeTok{alternative =} \StringTok{"less"}\NormalTok{)}
\end{Highlighting}
\end{Shaded}

}

\end{codelisting}%

\begin{verbatim}

    One Sample t-test

data:  datos
t = 2, df = 24, p-value = 1
alternative hypothesis: true mean is less than 1000
95 percent confidence interval:
 -Inf 1002
sample estimates:
mean of x 
     1001 
\end{verbatim}

\end{tcolorbox}

\begin{tcolorbox}[enhanced jigsaw, breakable, titlerule=0mm, rightrule=.15mm, leftrule=.75mm, left=2mm, opacitybacktitle=0.6, colback=white, bottomtitle=1mm, colbacktitle=quarto-callout-tip-color!10!white, title=\textcolor{quarto-callout-tip-color}{\faLightbulb}\hspace{0.5em}{Ejemplo con BioStatR: testt() y testp()}, toptitle=1mm, coltitle=black, opacityback=0, arc=.35mm, bottomrule=.15mm, toprule=.15mm, colframe=quarto-callout-tip-color-frame]

Documentación completa de ambas funciones en
Sección~\ref{sec-bio-contrastes}.

\begin{codelisting}[H]

\caption{\label{lst-semana08-2}Código R}

\centering{

\begin{Shaded}
\begin{Highlighting}[]
\FunctionTok{library}\NormalTok{(BioStatR)}

\CommentTok{\# H₀: µ\_HbA1c = 7.5\%  (umbral clínico de control glucémico)}
\CommentTok{\# Dataset osteo: 94 pacientes diabéticos, Fac. Medicina UGR}
\FunctionTok{testt}\NormalTok{(}\AttributeTok{m =}\NormalTok{ osteo}\SpecialCharTok{$}\NormalTok{hba1c, }\AttributeTok{m0 =} \FloatTok{7.5}\NormalTok{, }\AttributeTok{grf =} \ConstantTok{FALSE}\NormalTok{)}
\end{Highlighting}
\end{Shaded}

}

\end{codelisting}%

\begin{verbatim}

# t-Test con una muestra
# ----------------------

# Resumen de 'osteo$hba1c'
    n = 94.000 
    media = 8.565 
    d.t. = 1.799 
    sem = 0.186 

# Estimación de la media μ:
  95%-IC(μ) = (8.196, 8.933) 

# Test de normalidad de Shapiro-Wilk:
  W = 0.983, gl = 94, p = 0.276 

# Test de Student para contrastar H₀:μ=μ₀ con μ₀=7.500
  texp = 5.740, gl = 93 
    p < 0.001 para la alternativa bilateral H₁:μ≠μ₀ 
    p < 0.001 para la alternativa unilateral H₁:μ>μ₀ 

  Estimación del efecto bruto 
  95%-IC(μ-μ₀) = (0.696, 1.433) 
\end{verbatim}

\begin{codelisting}[H]

\caption{\label{lst-semana08-2}Código R}

\centering{

\begin{Shaded}
\begin{Highlighting}[]
\CommentTok{\# H₀: π\_tabaco = 0.35  (prevalencia de referencia)}
\FunctionTok{testp}\NormalTok{(}\AttributeTok{x =} \FunctionTok{sum}\NormalTok{(osteo}\SpecialCharTok{$}\NormalTok{tabaco }\SpecialCharTok{==} \StringTok{"Sí"}\NormalTok{), }\AttributeTok{n =} \FunctionTok{nrow}\NormalTok{(osteo), }\AttributeTok{p0 =} \FloatTok{0.35}\NormalTok{)}
\end{Highlighting}
\end{Shaded}

}

\end{codelisting}%

\begin{verbatim}

# Test para contrastar una proporción binomial
# --------------------------------------------

# Información muestral
  n = 94
  x = 41   n-x=53
  p = 0.436; q = (1-p) = 0.564

# Test Ho:π=0.350
  [1] Método exacto
                      H1  Fexp Valor.p
    Cola derecha π>0.350 1.410   0.052
    Bilateral    π≠0.350     -   0.103

    95%-IC(π) = (0.339, 0.542) (método de Clooper-Pearson)

  [2] Método aproximado a la distribución normal
    Validez: min(nπ₀, n(1-π₀)) = 32.9 (>5, el método es válido)
    zexp = 1.643,  p  = 0.100

    95%-IC(π) = (0.335, 0.542) (método de Wilson)
\end{verbatim}

\texttt{testt()} integra Shapiro-Wilk automáticamente. La HbA1c supera
el umbral clínico (\(p < 0.001\)); la prevalencia de tabaquismo no
difiere del 35\% de referencia (\(p = 0.100\)).

\end{tcolorbox}

\begin{center}\rule{0.5\linewidth}{0.5pt}\end{center}

\section{Contraste para una Proporción}\label{sec-semana07-s8}

Cuando la variable de interés es dicotómica (sí/no, éxito/fracaso),
estamos interesados en la proporción poblacional \(\pi\).

\subsection{Función Muestral}\label{funciuxf3n-muestral}

Si \(X_i \sim \text{Bernoulli}(\pi)\) (cada observación es 1 con
probabilidad \(\pi\), 0 con probabilidad \(1-\pi\)), entonces el número
total de ``éxitos'' es:

\[X = \sum_{i=1}^n X_i \sim \text{Binomial}(n, \pi)\]

La proporción muestral es: \[\hat{\pi} = \frac{X}{n}\]

\subsection{\texorpdfstring{Distribución bajo
\(H_0\)}{Distribución bajo H\_0}}\label{distribuciuxf3n-bajo-h_0}

\begin{tcolorbox}[enhanced jigsaw, breakable, titlerule=0mm, rightrule=.15mm, leftrule=.75mm, left=2mm, opacitybacktitle=0.6, colback=white, bottomtitle=1mm, colbacktitle=quarto-callout-warning-color!10!white, title=\textcolor{quarto-callout-warning-color}{\faExclamationTriangle}\hspace{0.5em}{Resultado Importante: Contraste Binomial}, toptitle=1mm, coltitle=black, opacityback=0, arc=.35mm, bottomrule=.15mm, toprule=.15mm, colframe=quarto-callout-warning-color-frame]

Bajo \(H_0: \pi = \pi_0\):

\[X \sim \text{Binomial}(n, \pi_0)\]

La distribución del número de éxitos es exactamente binomial. Los
valores críticos se encuentran usando tablas binomiales o cálculos
computacionales.

\end{tcolorbox}

\subsection{Aproximación Normal (Muestra
Grande)}\label{aproximaciuxf3n-normal-muestra-grande}

Para \(n\) grande (típicamente \(n\pi_0(1-\pi_0) \geq 5\)):

\[Z = \frac{\hat{\pi} - \pi_0}{\sqrt{\frac{\pi_0(1-\pi_0)}{n}}} \approx N(0,1)\]

Se pueden usar los valores críticos normales en lugar de los binomiales
exactos.

\subsection{Ejemplo: Moneda Justa}\label{ejemplo-moneda-justa}

\textbf{Problema:} Probamos si una moneda es justa lanzándola 10 veces.

\textbf{Hipótesis:} - \(H_0: \pi = 0.5\) (moneda justa) -
\(H_1: \pi \neq 0.5\) (moneda cargada)

\textbf{Muestra:} \(n = 10\), observamos \(x = 9\) caras

\textbf{Distribución binomial bajo \(H_0\):} \(X \sim B(10, 0.5)\)

\textbf{Tabla de probabilidades acumuladas:}

\begin{longtable}[]{@{}llll@{}}
\toprule\noalign{}
\(x\) & \(P(X \leq x)\) & \(P(X = x)\) & Región \\
\midrule\noalign{}
\endhead
\bottomrule\noalign{}
\endlastfoot
0 & 0.0010 & 0.0010 & Rechazo \\
1 & 0.0107 & 0.0097 & Rechazo \\
2 & 0.0547 & 0.0440 & \\
\ldots{} & \ldots{} & \ldots{} & No rechazo \\
8 & 0.9893 & 0.0440 & \\
9 & 0.9990 & 0.0097 & Rechazo \\
10 & 1.0000 & 0.0010 & Rechazo \\
\end{longtable}

Con \(\alpha = 0.05\): - Región de rechazo bilateral:
\(\{0, 1\} \cup \{9, 10\}\) - Como \(x = 9\) está en la región de
rechazo, rechazamos \(H_0\) - p-valor:
\(P(X \leq 1) + P(X \geq 9) = 0.0107 + (1 - 0.9893) = 0.0214\)

\textbf{Conclusión:} La evidencia sugiere que la moneda no es justa (es
significativamente sesgada hacia caras).

\begin{tcolorbox}[enhanced jigsaw, breakable, titlerule=0mm, rightrule=.15mm, leftrule=.75mm, left=2mm, opacitybacktitle=0.6, colback=white, bottomtitle=1mm, colbacktitle=quarto-callout-tip-color!10!white, title=\textcolor{quarto-callout-tip-color}{\faLightbulb}\hspace{0.5em}{Código R: Contraste Binomial}, toptitle=1mm, coltitle=black, opacityback=0, arc=.35mm, bottomrule=.15mm, toprule=.15mm, colframe=quarto-callout-tip-color-frame]

\begin{codelisting}[H]

\caption{\label{lst-semana08-4}Código R}

\centering{

\begin{Shaded}
\begin{Highlighting}[]
\CommentTok{\# Contraste binomial exacto}
\CommentTok{\# Lanzamos una moneda 10 veces, observamos 9 caras}
\FunctionTok{binom.test}\NormalTok{(}\AttributeTok{x =} \DecValTok{9}\NormalTok{, }\AttributeTok{n =} \DecValTok{10}\NormalTok{, }\AttributeTok{p =} \FloatTok{0.5}\NormalTok{, }\AttributeTok{alternative =} \StringTok{"two.sided"}\NormalTok{)}
\end{Highlighting}
\end{Shaded}

}

\end{codelisting}%

\begin{verbatim}

    Exact binomial test

data:  9 and 10
number of successes = 9, number of trials = 10, p-value = 0.02
alternative hypothesis: true probability of success is not equal to 0.5
95 percent confidence interval:
 0.555 0.997
sample estimates:
probability of success 
                   0.9 
\end{verbatim}

\begin{codelisting}[H]

\caption{\label{lst-semana08-4}Código R}

\centering{

\begin{Shaded}
\begin{Highlighting}[]
\CommentTok{\# Aproximación normal para muestra más grande}
\CommentTok{\# Lanzamos 100 veces, observamos 65 caras}
\FunctionTok{prop.test}\NormalTok{(}\AttributeTok{x =} \DecValTok{65}\NormalTok{, }\AttributeTok{n =} \DecValTok{100}\NormalTok{, }\AttributeTok{p =} \FloatTok{0.5}\NormalTok{, }\AttributeTok{alternative =} \StringTok{"two.sided"}\NormalTok{)}
\end{Highlighting}
\end{Shaded}

}

\end{codelisting}%

\begin{verbatim}

    1-sample proportions test with continuity correction

data:  65 out of 100, null probability 0.5
X-squared = 8, df = 1, p-value = 0.004
alternative hypothesis: true p is not equal to 0.5
95 percent confidence interval:
 0.547 0.741
sample estimates:
   p 
0.65 
\end{verbatim}

\begin{codelisting}[H]

\caption{\label{lst-semana08-4}Código R}

\centering{

\begin{Shaded}
\begin{Highlighting}[]
\CommentTok{\# Para proporción: H0: pi = 0.5 vs H1: pi != 0.5}
\NormalTok{resultado }\OtherTok{\textless{}{-}} \FunctionTok{binom.test}\NormalTok{(}\DecValTok{9}\NormalTok{, }\DecValTok{10}\NormalTok{, }\FloatTok{0.5}\NormalTok{)}
\FunctionTok{cat}\NormalTok{(}\StringTok{"p{-}valor ="}\NormalTok{, resultado}\SpecialCharTok{$}\NormalTok{p.value, }\StringTok{"}\SpecialCharTok{\textbackslash{}n}\StringTok{"}\NormalTok{)}
\end{Highlighting}
\end{Shaded}

}

\end{codelisting}%

\begin{verbatim}
p-valor = 0.0215 
\end{verbatim}

\begin{codelisting}[H]

\caption{\label{lst-semana08-4}Código R}

\centering{

\begin{Shaded}
\begin{Highlighting}[]
\FunctionTok{cat}\NormalTok{(}\StringTok{"Intervalo de confianza:"}\NormalTok{, resultado}\SpecialCharTok{$}\NormalTok{conf.int, }\StringTok{"}\SpecialCharTok{\textbackslash{}n}\StringTok{"}\NormalTok{)}
\end{Highlighting}
\end{Shaded}

}

\end{codelisting}%

\begin{verbatim}
Intervalo de confianza: 0.555 0.997 
\end{verbatim}

\end{tcolorbox}

\begin{center}\rule{0.5\linewidth}{0.5pt}\end{center}

\section{Valor-p (p-value)}\label{sec-semana07-s9}

El valor-p es una de las cantidades más importantes en estadística
inferencial.

\begin{tcolorbox}[enhanced jigsaw, breakable, titlerule=0mm, rightrule=.15mm, leftrule=.75mm, left=2mm, opacitybacktitle=0.6, colback=white, bottomtitle=1mm, colbacktitle=quarto-callout-note-color!10!white, title=\textcolor{quarto-callout-note-color}{\faInfo}\hspace{0.5em}{Definición: Valor-p}, toptitle=1mm, coltitle=black, opacityback=0, arc=.35mm, bottomrule=.15mm, toprule=.15mm, colframe=quarto-callout-note-color-frame]

El \textbf{valor-p} (o \textbf{p-value}) es la probabilidad, calculada
asumiendo que \(H_0\) es verdadera, de obtener un valor del estadístico
de prueba tan o más extremo que el valor observado.

\textbf{Fórmula general:}
\[\text{valor-p} = P(T \text{ sea tan o más extremo que } t_{obs} | H_0)\]

donde ``más extremo'' depende del tipo de contraste: -
\textbf{Bilateral:} \(P(|T| \geq |t_{obs}|)\) - \textbf{Unilateral
derecha:} \(P(T \geq t_{obs})\) - \textbf{Unilateral izquierda:}
\(P(T \leq t_{obs})\)

\end{tcolorbox}

\subsection{Interpretación}\label{interpretaciuxf3n-9}

El valor-p \textbf{no es}: - La probabilidad de que \(H_0\) sea
verdadera - La probabilidad de que \(H_1\) sea verdadera - La
probabilidad de haber cometido un error

El valor-p \textbf{es}: - Una medida de compatibilidad entre los datos y
\(H_0\) - Un p-valor pequeño indica que los datos observados serían muy
improbables si \(H_0\) fuera verdadera

\subsection{Regla de Decisión}\label{regla-de-decisiuxf3n}

\textbf{Enfoque tradicional con nivel \(\alpha\):} - Si p-valor
\(< \alpha\) \(\Rightarrow\) Rechazamos \(H_0\) - Si p-valor
\(\geq \alpha\) \(\Rightarrow\) No rechazamos \(H_0\)

\textbf{Enfoque moderno (más flexible):} - Reportar el valor-p exacto y
dejar que el lector juzgue - Valores muy pequeños (\textless{} 0.01)
indican evidencia fuerte contra \(H_0\) - Valores pequeños (0.01 - 0.05)
indican evidencia moderada - Valores grandes (\textgreater{} 0.05)
indican que los datos son compatibles con \(H_0\)

\subsection{Ejemplo: Cálculo de
p-valores}\label{ejemplo-cuxe1lculo-de-p-valores}

\textbf{Ejemplo 1 (Normal):} - Observamos \(z_{obs} = 2.5\) en un
contraste bilateral - p-valor = \(2 P(Z > 2.5) = 2(0.0062) = 0.0124\)

\textbf{Ejemplo 2 (t-Student):} - Observamos \(t_{obs} = 2.04\) con 24
gl en un contraste bilateral - p-valor = \(2 P(T_{24} > 2.04)\) - Usando
R: \texttt{2\ *\ (1\ -\ pt(2.04,\ 24))} ≈ 0.0525

\textbf{Ejemplo 3 (Binomial):} - Observamos 9 caras en 10 lanzamientos,
contraste bilateral con \(H_0: \pi = 0.5\) - p-valor =
\(P(X \leq 1) + P(X \geq 9) = 0.0214\)

\begin{center}\rule{0.5\linewidth}{0.5pt}\end{center}

\section{Tabla Resumen: Regiones de Aceptación y
Rechazo}\label{sec-semana07-s10}

\begin{tcolorbox}[enhanced jigsaw, breakable, titlerule=0mm, rightrule=.15mm, leftrule=.75mm, left=2mm, opacitybacktitle=0.6, colback=white, bottomtitle=1mm, colbacktitle=quarto-callout-warning-color!10!white, title=\textcolor{quarto-callout-warning-color}{\faExclamationTriangle}\hspace{0.5em}{Tabla de Referencia: Pruebas de una Muestra}, toptitle=1mm, coltitle=black, opacityback=0, arc=.35mm, bottomrule=.15mm, toprule=.15mm, colframe=quarto-callout-warning-color-frame]

\begin{longtable}[]{@{}
  >{\raggedright\arraybackslash}p{(\linewidth - 12\tabcolsep) * \real{0.1122}}
  >{\raggedright\arraybackslash}p{(\linewidth - 12\tabcolsep) * \real{0.0816}}
  >{\raggedright\arraybackslash}p{(\linewidth - 12\tabcolsep) * \real{0.1327}}
  >{\raggedright\arraybackslash}p{(\linewidth - 12\tabcolsep) * \real{0.1429}}
  >{\raggedright\arraybackslash}p{(\linewidth - 12\tabcolsep) * \real{0.1122}}
  >{\raggedright\arraybackslash}p{(\linewidth - 12\tabcolsep) * \real{0.1939}}
  >{\raggedright\arraybackslash}p{(\linewidth - 12\tabcolsep) * \real{0.2245}}@{}}
\toprule\noalign{}
\begin{minipage}[b]{\linewidth}\raggedright
Parámetro
\end{minipage} & \begin{minipage}[b]{\linewidth}\raggedright
Prueba
\end{minipage} & \begin{minipage}[b]{\linewidth}\raggedright
Estadístico
\end{minipage} & \begin{minipage}[b]{\linewidth}\raggedright
Distribución
\end{minipage} & \begin{minipage}[b]{\linewidth}\raggedright
Bilateral
\end{minipage} & \begin{minipage}[b]{\linewidth}\raggedright
Unilateral Derecha
\end{minipage} & \begin{minipage}[b]{\linewidth}\raggedright
Unilateral Izquierda
\end{minipage} \\
\midrule\noalign{}
\endhead
\bottomrule\noalign{}
\endlastfoot
\textbf{Media} (\(\sigma\) conocida) & Z &
\(Z = \frac{\bar{X}-\mu_0}{\sigma/\sqrt{n}}\) & \(N(0,1)\) & Rechazo:
\(\|Z\| > z_{1-\alpha/2}\) & Rechazo: \(Z > z_{1-\alpha}\) & Rechazo:
\(Z < -z_{1-\alpha}\) \\
\textbf{Media} (\(\sigma\) desconocida) & t &
\(T = \frac{\bar{X}-\mu_0}{S/\sqrt{n}}\) & \(t_{n-1}\) & Rechazo:
\(\|T\| > t_{n-1,1-\alpha/2}\) & Rechazo: \(T > t_{n-1,1-\alpha}\) &
Rechazo: \(T < -t_{n-1,1-\alpha}\) \\
\textbf{Proporción} & Binomial & \(X = \sum X_i\) & \(B(n,\pi_0)\) &
Rechazo: \(P(X \leq x_l) + P(X \geq x_u) < \alpha\) & Rechazo:
\(P(X \geq x) < \alpha\) & Rechazo: \(P(X \leq x) < \alpha\) \\
\textbf{Proporción} (grande \(n\)) & Z &
\(Z = \frac{\hat{\pi}-\pi_0}{\sqrt{\pi_0(1-\pi_0)/n}}\) & \(N(0,1)\) &
Rechazo: \(\|Z\| > z_{1-\alpha/2}\) & Rechazo: \(Z > z_{1-\alpha}\) &
Rechazo: \(Z < -z_{1-\alpha}\) \\
\end{longtable}

\end{tcolorbox}

\begin{center}\rule{0.5\linewidth}{0.5pt}\end{center}

\section{Ejemplos Completos}\label{sec-semana07-s11}

\subsection{Ejemplo 1: Contraste Z Bilateral (Conocida la
Varianza)}\label{ejemplo-1-contraste-z-bilateral-conocida-la-varianza}

\textbf{Problema:} El peso de paquetes de cereales debería ser 500 g.
Sospechamos que hay variación. Se sabe que \(\sigma = 10\) g.

\textbf{Datos:} Muestra de \(n = 36\) paquetes, \(\bar{x} = 498.5\) g

\textbf{Hipótesis:} \(H_0: \mu = 500\) vs.~\(H_1: \mu \neq 500\)

\textbf{Análisis:}
\[z_{obs} = \frac{498.5 - 500}{10/\sqrt{36}} = \frac{-1.5}{1.667} = -0.90\]

Con \(\alpha = 0.05\): valor crítico = \(\pm 1.96\)

\textbf{Decisión:} Como \(|-0.90| < 1.96\), no rechazamos \(H_0\)

\textbf{p-valor:} \(2P(Z > 0.90) = 2(0.1841) = 0.3682\)

\textbf{Conclusión:} Los datos no proporcionan evidencia de que el peso
medio sea diferente de 500 g.

\subsection{Ejemplo 2: Contraste t Unilateral (Desconocida la
Varianza)}\label{ejemplo-2-contraste-t-unilateral-desconocida-la-varianza}

\textbf{Problema:} Un fertilizante se supone que aumenta el rendimiento
a más de 100 unidades. Probamos si funciona.

\textbf{Datos:} Muestra de \(n = 16\) parcelas, \(\bar{x} = 105\),
\(s = 12\)

\textbf{Hipótesis:} \(H_0: \mu \leq 100\) vs.~\(H_1: \mu > 100\)
(unilateral derecha)

\textbf{Análisis:}
\[t_{obs} = \frac{105 - 100}{12/\sqrt{16}} = \frac{5}{3} = 1.67\]

Con \(\alpha = 0.05\) y \(df = 15\): valor crítico =
\(t_{15, 0.95} = 1.753\)

\textbf{Decisión:} Como \(t_{obs} = 1.67 < 1.753\), no rechazamos
\(H_0\)

\textbf{p-valor:} \(P(T_{15} > 1.67) \approx 0.058\)

\textbf{Conclusión:} La evidencia es marginal. No podemos afirmar al 5\%
que el fertilizante aumenta el rendimiento, aunque está cerca.

\subsection{Ejemplo 3: Contraste Binomial
Exacto}\label{ejemplo-3-contraste-binomial-exacto}

\textbf{Problema:} Un fármaco se supone que tiene una tasa de efectos
adversos del 5\%. Observamos 15 efectos adversos en 100 pacientes. ¿Es
significativamente diferente?

\textbf{Datos:} \(n = 100\), \(x = 15\), \(\pi_0 = 0.05\)

\textbf{Hipótesis:} \(H_0: \pi = 0.05\) vs.~\(H_1: \pi \neq 0.05\)

\textbf{Aproximación normal:}
\[z_{obs} = \frac{0.15 - 0.05}{\sqrt{0.05 \times 0.95 / 100}} = \frac{0.10}{\sqrt{0.00475}} = \frac{0.10}{0.0689} = 1.45\]

Con \(\alpha = 0.05\): valor crítico = \(\pm 1.96\)

\textbf{Decisión:} Como \(|1.45| < 1.96\), no rechazamos \(H_0\)

\textbf{p-valor:} \(2P(Z > 1.45) = 2(0.0735) = 0.147\)

\textbf{Conclusión:} La tasa observada no es significativamente
diferente de 0.05.

\begin{tcolorbox}[enhanced jigsaw, breakable, titlerule=0mm, rightrule=.15mm, leftrule=.75mm, left=2mm, opacitybacktitle=0.6, colback=white, bottomtitle=1mm, colbacktitle=quarto-callout-tip-color!10!white, title=\textcolor{quarto-callout-tip-color}{\faLightbulb}\hspace{0.5em}{Código R: Ejemplos Completos}, toptitle=1mm, coltitle=black, opacityback=0, arc=.35mm, bottomrule=.15mm, toprule=.15mm, colframe=quarto-callout-tip-color-frame]

\begin{codelisting}[H]

\caption{\label{lst-semana08-5}Código R}

\centering{

\begin{Shaded}
\begin{Highlighting}[]
\CommentTok{\# Ejemplo 1: Contraste Z (conocida varianza)}
\CommentTok{\# z.test no existe en base R, usamos estadística manual}
\NormalTok{x\_bar }\OtherTok{\textless{}{-}} \FloatTok{498.5}
\NormalTok{mu\_0 }\OtherTok{\textless{}{-}} \DecValTok{500}
\NormalTok{sigma }\OtherTok{\textless{}{-}} \DecValTok{10}
\NormalTok{n }\OtherTok{\textless{}{-}} \DecValTok{36}
\NormalTok{z\_obs }\OtherTok{\textless{}{-}}\NormalTok{ (x\_bar }\SpecialCharTok{{-}}\NormalTok{ mu\_0) }\SpecialCharTok{/}\NormalTok{ (sigma }\SpecialCharTok{/} \FunctionTok{sqrt}\NormalTok{(n))}
\NormalTok{p\_valor }\OtherTok{\textless{}{-}} \DecValTok{2} \SpecialCharTok{*}\NormalTok{ (}\DecValTok{1} \SpecialCharTok{{-}} \FunctionTok{pnorm}\NormalTok{(}\FunctionTok{abs}\NormalTok{(z\_obs)))}
\FunctionTok{cat}\NormalTok{(}\StringTok{"Ejemplo 1:}\SpecialCharTok{\textbackslash{}n}\StringTok{"}\NormalTok{)}
\end{Highlighting}
\end{Shaded}

}

\end{codelisting}%

\begin{verbatim}
Ejemplo 1:
\end{verbatim}

\begin{codelisting}[H]

\caption{\label{lst-semana08-5}Código R}

\centering{

\begin{Shaded}
\begin{Highlighting}[]
\FunctionTok{cat}\NormalTok{(}\StringTok{"z observado ="}\NormalTok{, }\FunctionTok{round}\NormalTok{(z\_obs, }\DecValTok{3}\NormalTok{), }\StringTok{"}\SpecialCharTok{\textbackslash{}n}\StringTok{"}\NormalTok{)}
\end{Highlighting}
\end{Shaded}

}

\end{codelisting}%

\begin{verbatim}
z observado = -0.9 
\end{verbatim}

\begin{codelisting}[H]

\caption{\label{lst-semana08-5}Código R}

\centering{

\begin{Shaded}
\begin{Highlighting}[]
\FunctionTok{cat}\NormalTok{(}\StringTok{"p{-}valor ="}\NormalTok{, }\FunctionTok{round}\NormalTok{(p\_valor, }\DecValTok{4}\NormalTok{), }\StringTok{"}\SpecialCharTok{\textbackslash{}n\textbackslash{}n}\StringTok{"}\NormalTok{)}
\end{Highlighting}
\end{Shaded}

}

\end{codelisting}%

\begin{verbatim}
p-valor = 0.368 
\end{verbatim}

\begin{codelisting}[H]

\caption{\label{lst-semana08-5}Código R}

\centering{

\begin{Shaded}
\begin{Highlighting}[]
\CommentTok{\# Ejemplo 2: Contraste t unilateral}
\NormalTok{datos2 }\OtherTok{\textless{}{-}} \FunctionTok{c}\NormalTok{(}\DecValTok{115}\NormalTok{, }\DecValTok{108}\NormalTok{, }\DecValTok{102}\NormalTok{, }\DecValTok{110}\NormalTok{, }\DecValTok{98}\NormalTok{, }\DecValTok{104}\NormalTok{, }\DecValTok{112}\NormalTok{, }\DecValTok{96}\NormalTok{,}
            \DecValTok{106}\NormalTok{, }\DecValTok{109}\NormalTok{, }\DecValTok{101}\NormalTok{, }\DecValTok{107}\NormalTok{, }\DecValTok{103}\NormalTok{, }\DecValTok{111}\NormalTok{, }\DecValTok{95}\NormalTok{, }\DecValTok{113}\NormalTok{)}
\NormalTok{resultado2 }\OtherTok{\textless{}{-}} \FunctionTok{t.test}\NormalTok{(datos2, }\AttributeTok{mu =} \DecValTok{100}\NormalTok{, }\AttributeTok{alternative =} \StringTok{"greater"}\NormalTok{)}
\FunctionTok{cat}\NormalTok{(}\StringTok{"Ejemplo 2:}\SpecialCharTok{\textbackslash{}n}\StringTok{"}\NormalTok{)}
\end{Highlighting}
\end{Shaded}

}

\end{codelisting}%

\begin{verbatim}
Ejemplo 2:
\end{verbatim}

\begin{codelisting}[H]

\caption{\label{lst-semana08-5}Código R}

\centering{

\begin{Shaded}
\begin{Highlighting}[]
\FunctionTok{print}\NormalTok{(resultado2)}
\end{Highlighting}
\end{Shaded}

}

\end{codelisting}%

\begin{verbatim}

    One Sample t-test

data:  datos2
t = 4, df = 15, p-value = 0.001
alternative hypothesis: true mean is greater than 100
95 percent confidence interval:
 103 Inf
sample estimates:
mean of x 
      106 
\end{verbatim}

\begin{codelisting}[H]

\caption{\label{lst-semana08-5}Código R}

\centering{

\begin{Shaded}
\begin{Highlighting}[]
\CommentTok{\# Ejemplo 3: Contraste Binomial}
\NormalTok{resultado3 }\OtherTok{\textless{}{-}} \FunctionTok{binom.test}\NormalTok{(}\AttributeTok{x =} \DecValTok{15}\NormalTok{, }\AttributeTok{n =} \DecValTok{100}\NormalTok{, }\AttributeTok{p =} \FloatTok{0.05}\NormalTok{,}
                         \AttributeTok{alternative =} \StringTok{"two.sided"}\NormalTok{)}
\FunctionTok{cat}\NormalTok{(}\StringTok{"}\SpecialCharTok{\textbackslash{}n}\StringTok{Ejemplo 3:}\SpecialCharTok{\textbackslash{}n}\StringTok{"}\NormalTok{)}
\end{Highlighting}
\end{Shaded}

}

\end{codelisting}%

\begin{verbatim}

Ejemplo 3:
\end{verbatim}

\begin{codelisting}[H]

\caption{\label{lst-semana08-5}Código R}

\centering{

\begin{Shaded}
\begin{Highlighting}[]
\FunctionTok{print}\NormalTok{(resultado3)}
\end{Highlighting}
\end{Shaded}

}

\end{codelisting}%

\begin{verbatim}

    Exact binomial test

data:  15 and 100
number of successes = 15, number of trials = 100, p-value = 1e-04
alternative hypothesis: true probability of success is not equal to 0.05
95 percent confidence interval:
 0.0865 0.2353
sample estimates:
probability of success 
                  0.15 
\end{verbatim}

\end{tcolorbox}

\begin{center}\rule{0.5\linewidth}{0.5pt}\end{center}

\section{Relación entre Hipótesis Unilaterales y
Bilaterales}\label{sec-semana07-s12}

\subsection{Elección de la
Dirección}\label{elecciuxf3n-de-la-direcciuxf3n}

La dirección del contraste (bilateral vs.~unilateral) debe determinarse
\textbf{antes} de recopilar datos, basándose en la pregunta de
investigación:

\textbf{Bilateral (\(H_1: \theta \neq \theta_0\)):} - Cuando nos
interesa desviaciones en ambas direcciones - Ejemplo: ¿Difiere la media
de algún valor? - Más conservador (requiere evidencia más fuerte)

\textbf{Unilateral derecho (\(H_1: \theta > \theta_0\)):} - Cuando
específicamente buscamos que el parámetro sea mayor - Ejemplo: ¿El nuevo
tratamiento mejora el resultado? - Valor crítico es \(z_{1-\alpha}\) (no
\(z_{1-\alpha/2}\))

\textbf{Unilateral izquierdo (\(H_1: \theta < \theta_0\)):} - Cuando
específicamente buscamos que el parámetro sea menor - Ejemplo: ¿El
fármaco reduce los efectos adversos? - Valor crítico es
\(-z_{1-\alpha}\)

\subsection{Implicaciones para
Errores}\label{implicaciones-para-errores}

La elección afecta a \(\alpha\) y \(\beta\):

En contraste \textbf{bilateral} con \(\alpha = 0.05\): - Dividimos
\(\alpha\) en ambas colas: \(\alpha/2 = 0.025\) en cada cola - Valor
crítico es más grande (\(z_{0.975} = 1.96\)) - Mayor \(\beta\) (menos
poder)

En contraste \textbf{unilateral} con \(\alpha = 0.05\): - Toda
\(\alpha = 0.05\) va en una cola - Valor crítico es más pequeño
(\(z_{0.95} = 1.645\)) - Menor \(\beta\) (más poder)

\textbf{Por eso es importante elegir antes:} Si eliges unilateral
después de ver los datos, estás ``p-hacking'' y aumentas falsamente el
poder.

\begin{center}\rule{0.5\linewidth}{0.5pt}\end{center}

\section{Resumen}\label{sec-semana07-resumen}

\begin{tcolorbox}[enhanced jigsaw, breakable, titlerule=0mm, rightrule=.15mm, leftrule=.75mm, left=2mm, opacitybacktitle=0.6, colback=white, bottomtitle=1mm, colbacktitle=quarto-callout-note-color!10!white, title=\textcolor{quarto-callout-note-color}{\faInfo}\hspace{0.5em}{Conceptos Clave de Semana 8}, toptitle=1mm, coltitle=black, opacityback=0, arc=.35mm, bottomrule=.15mm, toprule=.15mm, colframe=quarto-callout-note-color-frame]

\begin{enumerate}
\def\labelenumi{\arabic{enumi}.}
\item
  \textbf{Hipótesis:} Definimos \(H_0\) (lo que asumimos) y \(H_1\) (lo
  que queremos demostrar). Deben ser complementarias.
\item
  \textbf{Errores:} Error Tipo I (\(\alpha\), falso positivo) y Error
  Tipo II (\(\beta\), falso negativo) son inversamente relacionados.
  Solo se reduce ambos aumentando \(n\).
\item
  \textbf{Estadístico de Prueba:} Un valor calculado de los datos cuya
  distribución es conocida bajo \(H_0\).
\item
  \textbf{Región Crítica:} Conjunto de valores del estadístico donde
  rechazamos \(H_0\), determinado por el nivel \(\alpha\).
\item
  \textbf{Contraste Z:} Cuando \(\sigma\) es conocida. Estadístico
  \(Z = \frac{\bar{X}-\mu_0}{\sigma/\sqrt{n}} \sim N(0,1)\)
\item
  \textbf{Contraste t:} Cuando \(\sigma\) es desconocida (casos
  prácticos). Estadístico
  \(T = \frac{\bar{X}-\mu_0}{S/\sqrt{n}} \sim t_{n-1}\)
\item
  \textbf{Contraste Binomial:} Para proporciones. Distribución exacta es
  \(B(n,\pi_0)\); aproximación normal para \(n\) grande.
\item
  \textbf{p-valor:} Probabilidad de observar datos tan extremos o más
  bajo \(H_0\). Si p-valor \(< \alpha\), rechazamos \(H_0\).
\item
  \textbf{Decisión:} Comparamos el valor observado del estadístico con
  el valor crítico (o comparamos p-valor con \(\alpha\)).
\item
  \textbf{Potencia:} La capacidad de detectar un efecto real. Mejora con
  mayor \(n\), mayor diferencia real, o mayor \(\alpha\).
\end{enumerate}

\end{tcolorbox}

\subsection{Tabla de Decisión
Rápida}\label{tabla-de-decisiuxf3n-ruxe1pida}

\begin{longtable}[]{@{}
  >{\raggedright\arraybackslash}p{(\linewidth - 2\tabcolsep) * \real{0.4762}}
  >{\raggedright\arraybackslash}p{(\linewidth - 2\tabcolsep) * \real{0.5238}}@{}}
\toprule\noalign{}
\begin{minipage}[b]{\linewidth}\raggedright
Pregunta
\end{minipage} & \begin{minipage}[b]{\linewidth}\raggedright
Respuesta
\end{minipage} \\
\midrule\noalign{}
\endhead
\bottomrule\noalign{}
\endlastfoot
¿Se conoce \(\sigma\)? & Sí → Contraste Z; No → Contraste t \\
¿Bilateral o unilateral? & Determinar antes de recopilar datos \\
¿Qué \(\alpha\)? & Típicamente 0.05; a veces 0.01 \\
¿Cómo decidir? & Opción 1: Comparar \(|T|\) con valor crítico; Opción 2:
Comparar p-valor con \(\alpha\) \\
¿Qué significa ``no rechazar''? & NO significa aceptar \(H_0\); solo
significa ``no tenemos evidencia suficiente contra \(H_0\)'' \\
\end{longtable}

\begin{center}\rule{0.5\linewidth}{0.5pt}\end{center}

\section{Ejercicios}\label{sec-semana07-ejercicios}

\subsection{Ejercicio 1: Conceptos
Básicos}\label{ejercicio-1-conceptos-buxe1sicos}

Un médico sospecha que el nivel medio de colesterol en una población es
mayor a 200 mg/dL. Formule las hipótesis apropiadas e identifique los
tipos de error que podrían ocurrir.

\subsection{Ejercicio 2: Contraste Z}\label{ejercicio-2-contraste-z}

Una máquina que produce tornillos tiene un ajuste tal que
\(\sigma = 0.5\) mm. Se toman 25 tornillos y se mide su diámetro
promedio: \(\bar{x} = 10.2\) mm. El diámetro especificado es 10 mm.

\begin{enumerate}
\def\labelenumi{\alph{enumi})}
\tightlist
\item
  Formule \(H_0\) y \(H_1\)
\item
  Calcule el estadístico Z
\item
  ¿Rechaza \(H_0\) con \(\alpha = 0.05\)? (bilateral)
\item
  Calcule el p-valor
\end{enumerate}

\subsection{Ejercicio 3: Contraste t}\label{ejercicio-3-contraste-t}

Se obtuvieron las siguientes mediciones de tiempo de reacción (en
segundos) para 10 sujetos:

\begin{verbatim}
0.42, 0.45, 0.38, 0.41, 0.39, 0.44, 0.40, 0.43, 0.37, 0.46
\end{verbatim}

\begin{enumerate}
\def\labelenumi{\alph{enumi})}
\tightlist
\item
  ¿Es diferente de 0.40 segundos con \(\alpha = 0.05\)? (bilateral)
\item
  Calcule \(\bar{x}\) y \(s\)
\item
  Calcule el estadístico t
\item
  Determine el p-valor
\end{enumerate}

\textbf{Código R sugerido:}

\begin{codelisting}

\caption{\label{lst-semana08-6}Código R}

\centering{

\begin{Shaded}
\begin{Highlighting}[]
\NormalTok{datos }\OtherTok{\textless{}{-}} \FunctionTok{c}\NormalTok{(}\FloatTok{0.42}\NormalTok{, }\FloatTok{0.45}\NormalTok{, }\FloatTok{0.38}\NormalTok{, }\FloatTok{0.41}\NormalTok{, }\FloatTok{0.39}\NormalTok{, }\FloatTok{0.44}\NormalTok{, }\FloatTok{0.40}\NormalTok{, }\FloatTok{0.43}\NormalTok{, }\FloatTok{0.37}\NormalTok{, }\FloatTok{0.46}\NormalTok{)}
\NormalTok{resultado }\OtherTok{\textless{}{-}} \FunctionTok{t.test}\NormalTok{(datos, }\AttributeTok{mu =} \FloatTok{0.40}\NormalTok{)}
\FunctionTok{print}\NormalTok{(resultado)}
\end{Highlighting}
\end{Shaded}

}

\end{codelisting}%

\begin{verbatim}

    One Sample t-test

data:  datos
t = 2, df = 9, p-value = 0.2
alternative hypothesis: true mean is not equal to 0.4
95 percent confidence interval:
 0.393 0.437
sample estimates:
mean of x 
    0.415 
\end{verbatim}

\subsection{Ejercicio 4: Contraste
Binomial}\label{ejercicio-4-contraste-binomial}

Un casino afirma que su dado es justo. Lanzamos 20 veces y obtenemos
``6'' en 6 ocasiones.

\begin{enumerate}
\def\labelenumi{\alph{enumi})}
\tightlist
\item
  Formule \(H_0\) y \(H_1\)
\item
  ¿Es inusual este resultado? Calcule el p-valor
\item
  ¿Rechazaría que el dado es justo con \(\alpha = 0.05\)?
\end{enumerate}

\textbf{Código R sugerido:}

\begin{codelisting}

\caption{\label{lst-semana08-7}Código R}

\centering{

\begin{Shaded}
\begin{Highlighting}[]
\FunctionTok{binom.test}\NormalTok{(}\AttributeTok{x =} \DecValTok{6}\NormalTok{, }\AttributeTok{n =} \DecValTok{20}\NormalTok{, }\AttributeTok{p =} \DecValTok{1}\SpecialCharTok{/}\DecValTok{6}\NormalTok{, }\AttributeTok{alternative =} \StringTok{"two.sided"}\NormalTok{)}
\end{Highlighting}
\end{Shaded}

}

\end{codelisting}%

\begin{verbatim}

    Exact binomial test

data:  6 and 20
number of successes = 6, number of trials = 20, p-value = 0.1
alternative hypothesis: true probability of success is not equal to 0.167
95 percent confidence interval:
 0.119 0.543
sample estimates:
probability of success 
                   0.3 
\end{verbatim}

\subsection{Ejercicio 5: Decisión con
p-valor}\label{ejercicio-5-decisiuxf3n-con-p-valor}

Para cada uno de los siguientes p-valores, indique si rechaza \(H_0\)
con \(\alpha = 0.05\):

\begin{enumerate}
\def\labelenumi{\alph{enumi})}
\tightlist
\item
  p-valor = 0.031
\item
  p-valor = 0.052
\item
  p-valor = 0.0001
\item
  p-valor = 0.450
\end{enumerate}

\subsection{Ejercicio 6: Interpretación de
Resultados}\label{ejercicio-6-interpretaciuxf3n-de-resultados}

Un estudio investiga si el peso medio de una población difiere de 70 kg.
Se obtiene t = 1.85 con 24 grados de libertad y p-valor = 0.078.

\begin{enumerate}
\def\labelenumi{\alph{enumi})}
\tightlist
\item
  ¿Rechaza \(H_0\) con \(\alpha = 0.05\)?
\item
  ¿Rechazaría con \(\alpha = 0.10\)?
\item
  ¿Qué puede concluir de estos resultados?
\item
  ¿Por qué no es correcto decir ``se acepta \(H_0\)''?
\end{enumerate}

\section{Respuestas a los Ejercicios}\label{sec-semana07-respuestas}

\textbf{Ejercicio 1:} - H₀: μ = 200 mg/dL - H₁: μ \textgreater{} 200
mg/dL (unilateral derecha) - Error Tipo I: rechazar H₀ siendo verdadera
(afirmar que colesterol \textgreater{} 200 cuando realmente = 200) -
Error Tipo II: no rechazar H₀ siendo falsa (concluir que colesterol =
200 cuando realmente \textgreater{} 200)

\textbf{Ejercicio 2:} - a) H₀: μ = 10 mm, H₁: μ ≠ 10 mm - b) Z = (10.2 -
10)/(0.5/√25) = 0.2/0.1 = 2.0 - c) Z₀.₀₂₅ = 1.96; \textbar2.0\textbar{}
\textgreater{} 1.96, rechaza H₀ - d) p-valor ≈ 2×P(Z \textgreater{} 2) ≈
2×0.0228 = 0.0456

\textbf{Ejercicio 3:} - b) X̄ ≈ 0.415, s ≈ 0.0318 - c) t = (0.415 -
0.40)/(0.0318/√10) = 0.015/0.0101 ≈ 1.49 - d) p-valor ≈ 0.17 (df=9); no
rechaza H₀ con α=0.05

\textbf{Ejercicio 4:} - a) H₀: p = 1/6, H₁: p ≠ 1/6 - b) Bajo H₀, X
\textasciitilde{} Binomial(20, 1/6); P(X=6) ≈ 0.20 aproximadamente - c)
p-valor \textgreater{} 0.05; no rechaza que el dado sea justo

\textbf{Ejercicio 5:} - a) p = 0.031 \textless{} 0.05: Rechaza H₀ - b) p
= 0.052 \textgreater{} 0.05: No rechaza - c) p = 0.0001 \textless{}
0.05: Rechaza H₀ - d) p = 0.450 \textgreater{} 0.05: No rechaza

\textbf{Ejercicio 6:} - a) p-valor = 0.078 \textgreater{} 0.05: No
rechaza H₀ - b) p-valor = 0.078 \textless{} 0.10: Rechaza H₀ con α=0.10
- c) Evidencia marginal contra H₀; resultados borderline - d) Porque no
rechazar H₀ no significa que sea verdadera, solo que no hay evidencia
suficiente en los datos

\begin{center}\rule{0.5\linewidth}{0.5pt}\end{center}

\textbf{Fin de Semana 8: Contrastes de Hipótesis (Parte I)}

En las próximas semanas estudiaremos: - Contrastes para diferencia de
medias (dos muestras) - Contrastes para varianzas - ANOVA (análisis de
varianza) - Contrastes para independencia

\begin{center}\rule{0.5\linewidth}{0.5pt}\end{center}

\chapter{Semana 8 --- Contrastes de Hipótesis (Parte
II)}\label{sec-semana08}

En esta semana continuamos con contrastes de hipótesis, enfocándonos en
pruebas para comparar \textbf{dos muestras independientes y pareadas},
pruebas para igualdad de varianzas, y conceptos fundamentales de
\textbf{potencia estadística} y cálculo de tamaño muestral.

\section{Prueba Z para Diferencia de Medias (Varianzas
Conocidas)}\label{sec-semana08-z-test}

\subsection{Preparación: Distribución de la Diferencia de
Medias}\label{preparacion-distribucion-de-la-diferencia-de-medias}

Supongamos que tenemos dos poblaciones independientes: - Población 1:
\(X_1 \sim N(\mu_1, \sigma_1^2)\) con \(\sigma_1^2\) \textbf{conocida} -
Población 2: \(X_2 \sim N(\mu_2, \sigma_2^2)\) con \(\sigma_2^2\)
\textbf{conocida}

Tomamos muestras de tamaños \(n_1\) y \(n_2\) respectivamente. Los
estimadores de las medias son:

\[\bar{X}_1 \sim N\left(\mu_1, \frac{\sigma_1^2}{n_1}\right), \quad \bar{X}_2 \sim N\left(\mu_2, \frac{\sigma_2^2}{n_2}\right)\]

La diferencia entre las medias muestrales sigue una distribución normal:

\[D = \bar{X}_1 - \bar{X}_2 \sim N\left(\omega, \sigma_D^2\right)\]

donde:
\[\omega = \mu_1 - \mu_2, \quad \sigma_D^2 = \frac{\sigma_1^2}{n_1} + \frac{\sigma_2^2}{n_2}\]

\subsection{Hipótesis e Estadístico de
Prueba}\label{hipotesis-e-estadistico-de-prueba}

\begin{tcolorbox}[enhanced jigsaw, breakable, titlerule=0mm, rightrule=.15mm, leftrule=.75mm, left=2mm, opacitybacktitle=0.6, colback=white, bottomtitle=1mm, colbacktitle=quarto-callout-note-color!10!white, title=\textcolor{quarto-callout-note-color}{\faInfo}\hspace{0.5em}{Definición: Prueba Z para Dos Muestras}, toptitle=1mm, coltitle=black, opacityback=0, arc=.35mm, bottomrule=.15mm, toprule=.15mm, colframe=quarto-callout-note-color-frame]

Consideramos las hipótesis: - \textbf{Bilateral:}
\(H_0: \mu_1 - \mu_2 = \omega_0\)
vs.~\(H_1: \mu_1 - \mu_2 \neq \omega_0\) - \textbf{Unilateral derecha:}
\(H_0: \mu_1 - \mu_2 \leq \omega_0\)
vs.~\(H_1: \mu_1 - \mu_2 > \omega_0\) - \textbf{Unilateral izquierda:}
\(H_0: \mu_1 - \mu_2 \geq \omega_0\)
vs.~\(H_1: \mu_1 - \mu_2 < \omega_0\)

(Habitualmente, \(\omega_0 = 0\))

\end{tcolorbox}

\begin{tcolorbox}[enhanced jigsaw, breakable, titlerule=0mm, rightrule=.15mm, leftrule=.75mm, left=2mm, opacitybacktitle=0.6, colback=white, bottomtitle=1mm, colbacktitle=quarto-callout-warning-color!10!white, title=\textcolor{quarto-callout-warning-color}{\faExclamationTriangle}\hspace{0.5em}{Estadístico de Prueba: Z}, toptitle=1mm, coltitle=black, opacityback=0, arc=.35mm, bottomrule=.15mm, toprule=.15mm, colframe=quarto-callout-warning-color-frame]

Bajo \(H_0\) (asumiendo \(\omega = \omega_0\)):

\[Z = \frac{(\bar{X}_1 - \bar{X}_2) - \omega_0}{\sqrt{\frac{\sigma_1^2}{n_1} + \frac{\sigma_2^2}{n_2}}} \sim N(0, 1)\]

\textbf{Criterios de rechazo:} - \textbf{Bilateral:} Rechazo si
\(|Z| > z_{1-\alpha/2}\) - \textbf{Unilateral derecha:} Rechazo si
\(Z > z_{1-\alpha}\) - \textbf{Unilateral izquierda:} Rechazo si
\(Z < -z_{1-\alpha}\)

\end{tcolorbox}

\begin{center}\rule{0.5\linewidth}{0.5pt}\end{center}

\section{Prueba t para Dos Muestras Independientes (Varianzas
Desconocidas)}\label{sec-semana08-t-test}

\subsection{Caso 1: Varianzas Iguales
(Homogéneas)}\label{sec-semana08-equal-var}

Cuando no conocemos las varianzas pero \textbf{asumimos que son iguales}
(\(\sigma_1^2 = \sigma_2^2 = \sigma^2\)), utilizamos un estimador
combinado (pooled) de la varianza:

\begin{tcolorbox}[enhanced jigsaw, breakable, titlerule=0mm, rightrule=.15mm, leftrule=.75mm, left=2mm, opacitybacktitle=0.6, colback=white, bottomtitle=1mm, colbacktitle=quarto-callout-warning-color!10!white, title=\textcolor{quarto-callout-warning-color}{\faExclamationTriangle}\hspace{0.5em}{Varianza Combinada (Pooled)}, toptitle=1mm, coltitle=black, opacityback=0, arc=.35mm, bottomrule=.15mm, toprule=.15mm, colframe=quarto-callout-warning-color-frame]

\[S_p^2 = \frac{(n_1 - 1)S_1^2 + (n_2 - 1)S_2^2}{n_1 + n_2 - 2}\]

donde \(S_1^2\) y \(S_2^2\) son las varianzas muestrales sesgadas:

\[S_j^2 = \frac{1}{n_j - 1} \sum_{i=1}^{n_j} (X_{ji} - \bar{X}_j)^2, \quad j = 1, 2\]

El estimador de la varianza de la diferencia es:

\[S_D^2 = S_p^2 \left( \frac{1}{n_1} + \frac{1}{n_2} \right)\]

\end{tcolorbox}

\begin{tcolorbox}[enhanced jigsaw, breakable, titlerule=0mm, rightrule=.15mm, leftrule=.75mm, left=2mm, opacitybacktitle=0.6, colback=white, bottomtitle=1mm, colbacktitle=quarto-callout-warning-color!10!white, title=\textcolor{quarto-callout-warning-color}{\faExclamationTriangle}\hspace{0.5em}{Estadístico t (Varianzas Homogéneas)}, toptitle=1mm, coltitle=black, opacityback=0, arc=.35mm, bottomrule=.15mm, toprule=.15mm, colframe=quarto-callout-warning-color-frame]

\[T = \frac{(\bar{X}_1 - \bar{X}_2) - \omega_0}{S_D} \sim t_{n_1 + n_2 - 2}\]

bajo \(H_0: \mu_1 - \mu_2 = \omega_0\).

\textbf{Grados de libertad:} \(\nu = n_1 + n_2 - 2\)

\textbf{Criterios de rechazo:} - \textbf{Bilateral:} Rechazo si
\(|T| > t_{1-\alpha/2;\nu}\) - \textbf{Unilateral derecha:} Rechazo si
\(T > t_{1-\alpha;\nu}\) - \textbf{Unilateral izquierda:} Rechazo si
\(T < -t_{1-\alpha;\nu}\)

\end{tcolorbox}

\subsection{Caso 2: Varianzas Desiguales (Heterogéneas) --- Prueba de
Welch}\label{sec-semana08-unequal-var}

Cuando las varianzas poblacionales son \textbf{desiguales}
(\(\sigma_1^2 \neq \sigma_2^2\)), la solución es aproximada.

\begin{tcolorbox}[enhanced jigsaw, breakable, titlerule=0mm, rightrule=.15mm, leftrule=.75mm, left=2mm, opacitybacktitle=0.6, colback=white, bottomtitle=1mm, colbacktitle=quarto-callout-warning-color!10!white, title=\textcolor{quarto-callout-warning-color}{\faExclamationTriangle}\hspace{0.5em}{Estadístico de Welch-Satterthwaite}, toptitle=1mm, coltitle=black, opacityback=0, arc=.35mm, bottomrule=.15mm, toprule=.15mm, colframe=quarto-callout-warning-color-frame]

\[T = \frac{(\bar{X}_1 - \bar{X}_2) - \omega_0}{\sqrt{\frac{S_1^2}{n_1} + \frac{S_2^2}{n_2}}} \approx t_\nu\]

donde los \textbf{grados de libertad aproximados} (Welch-Satterthwaite)
son:

\[\nu = \frac{\left( \frac{S_1^2}{n_1} + \frac{S_2^2}{n_2} \right)^2}{\frac{1}{n_1-1}\left(\frac{S_1^2}{n_1}\right)^2 + \frac{1}{n_2-1}\left(\frac{S_2^2}{n_2}\right)^2}\]

Este test es más conservador que asumir varianzas iguales.

\end{tcolorbox}

\begin{tcolorbox}[enhanced jigsaw, breakable, titlerule=0mm, rightrule=.15mm, leftrule=.75mm, left=2mm, opacitybacktitle=0.6, colback=white, bottomtitle=1mm, colbacktitle=quarto-callout-tip-color!10!white, title=\textcolor{quarto-callout-tip-color}{\faLightbulb}\hspace{0.5em}{Ejemplo con BioStatR: testt(), testf() y testwx()}, toptitle=1mm, coltitle=black, opacityback=0, arc=.35mm, bottomrule=.15mm, toprule=.15mm, colframe=quarto-callout-tip-color-frame]

Documentacion completa de las tres funciones en
Sección~\ref{sec-bio-contrastes}.

\begin{codelisting}[H]

\caption{\label{lst-semana09-1}Código R}

\centering{

\begin{Shaded}
\begin{Highlighting}[]
\FunctionTok{library}\NormalTok{(BioStatR)}

\CommentTok{\# t{-}test dos muestras independientes: HbA1c por sexo}
\FunctionTok{testt}\NormalTok{(}\AttributeTok{grupos =}\NormalTok{ osteo}\SpecialCharTok{$}\NormalTok{sexo, }\AttributeTok{m =}\NormalTok{ osteo}\SpecialCharTok{$}\NormalTok{hba1c, }\AttributeTok{grf =} \ConstantTok{FALSE}\NormalTok{)}
\end{Highlighting}
\end{Shaded}

}

\end{codelisting}%

\begin{verbatim}
 

# t-test para 2 Muestras Independientes
# -------------------------------------

# Información muestral y estimación de las medias
  Niveles de agrupación: Hombre, Mujer 

                        n media    dt   sem             IC
  osteo$hba1c [Hombre] 45  8.55 1.745 0.260 (8.029, 9.078)
  osteo$hba1c [Mujer]  49  8.58 1.864 0.266  (8.04, 9.111)
  ____ 
  * IC elaborados al 95% de confianza para estimar μ₁ y μ₂ respectivamente 


# Pruebas de normalidad (test de Shapiro-Wilk)
[1] Para grupo = Hombre, W = 0.958, gl = 45, p = 0.102
[2] Para grupo = Mujer, W = 0.980, gl = 49, p = 0.566

# Test de homogeneidad de varianzas. Fexp = (var₂/var₁)
  Fexp = 1.141, gl₁ = 48, gl₂ = 44, p = 0.660 

# Diferencia de medias (osteo$hba1c [Mujer] - osteo$hba1c [Hombre])
  Hipótesis a contrastar: H₀:μ₁=μ₂ (μ₂-μ₁=0) 

  a) Test de Student (varianzas homogéneas) 
  texp = 0.059, gl = 92 
    p = 0.953 para la alternativa bilateral H₁:μ₁≠μ₂ 
    p = 0.476 para la alternativa unilateral H₁:μ₁<μ₂ 
  95%-IC(μ₂-μ₁) = (-0.719, 0.764) 

  b) Test de Welch (varianzas no homogéneas) 
  texp = 0.060, gl = 91.96 
    p = 0.953 para la alternativa bilateral H₁:μ₁≠μ₂ 
    p = 0.476 para la alternativa unilateral H₁:μ₁<μ₂ 
  95%-IC(μ₂-μ₁) = (-0.717, 0.762) 
\end{verbatim}

\begin{codelisting}[H]

\caption{\label{lst-semana09-1}Código R}

\centering{

\begin{Shaded}
\begin{Highlighting}[]
\CommentTok{\# Test F de varianzas del IMC por sexo}
\FunctionTok{testf}\NormalTok{(}\AttributeTok{s1 =} \FunctionTok{sd}\NormalTok{(osteo}\SpecialCharTok{$}\NormalTok{imc[osteo}\SpecialCharTok{$}\NormalTok{sexo }\SpecialCharTok{==} \StringTok{"Hombre"}\NormalTok{]),}
      \AttributeTok{n1 =} \FunctionTok{sum}\NormalTok{(osteo}\SpecialCharTok{$}\NormalTok{sexo }\SpecialCharTok{==} \StringTok{"Hombre"}\NormalTok{),}
      \AttributeTok{s2 =} \FunctionTok{sd}\NormalTok{(osteo}\SpecialCharTok{$}\NormalTok{imc[osteo}\SpecialCharTok{$}\NormalTok{sexo }\SpecialCharTok{==} \StringTok{"Mujer"}\NormalTok{]),}
      \AttributeTok{n2 =} \FunctionTok{sum}\NormalTok{(osteo}\SpecialCharTok{$}\NormalTok{sexo }\SpecialCharTok{==} \StringTok{"Mujer"}\NormalTok{))}
\end{Highlighting}
\end{Shaded}

}

\end{codelisting}%

\begin{verbatim}
[[1]]
[1] 2.31

[[2]]
[1] 48

[[3]]
[1] 44

[[4]]
[1] 0.00585
\end{verbatim}

\begin{codelisting}[H]

\caption{\label{lst-semana09-1}Código R}

\centering{

\begin{Shaded}
\begin{Highlighting}[]
\CommentTok{\# Wilcoxon{-}Mann{-}Whitney: IMC \textless{} 25 años vs \textgreater{} 33 años (no normales)}
\FunctionTok{testwx}\NormalTok{(}\AttributeTok{m1 =}\NormalTok{ osteo}\SpecialCharTok{$}\NormalTok{imc[osteo}\SpecialCharTok{$}\NormalTok{grupo\_edad }\SpecialCharTok{==} \StringTok{"\textless{} 25"}\NormalTok{],}
       \AttributeTok{m2 =}\NormalTok{ osteo}\SpecialCharTok{$}\NormalTok{imc[osteo}\SpecialCharTok{$}\NormalTok{grupo\_edad }\SpecialCharTok{==} \StringTok{"\textgreater{} 33"}\NormalTok{],}
       \AttributeTok{grf =} \ConstantTok{FALSE}\NormalTok{)}
\end{Highlighting}
\end{Shaded}

}

\end{codelisting}%

\begin{verbatim}

Test de Wilcoxon/Mann-Whithney para dos muestras independientes 
----------------------------------------------------------------
# Información muestral ---

                                Muestra  n    min     Q1     Q2     Q3    max
1 osteo$imc[osteo$grupo_edad == "< 25"] 32 19.628 21.421 22.845 24.149 32.475
2 osteo$imc[osteo$grupo_edad == "> 33"] 30 18.070 23.382 25.889 28.945 37.333
    RIQ
1 2.728
2 5.563

# Rangos ---

                                Muestra  n Suma_rangos Rango_medio       U
1 osteo$imc[osteo$grupo_edad == "< 25"] 32         807      25.219 681.000
2 osteo$imc[osteo$grupo_edad == "> 33"] 30        1146      38.200 279.000

# Test ---

  U = 279.000; Z = 2.831; W = 279.000; p = 0.005

# Tamaño del efecto ---

  Diferencia de localización: -2.667    95%-IC = (-4.518, -0.778) 
  r = 0.360  (criterio: 0.1 pequeño; 0.3 mediano; >0.5 grande)
  Probabilidad de superioridad PS = 0.709
  (probabilidad de que un valor al azar de M1 sea < a un valor al azar de M2) 
\end{verbatim}

\end{tcolorbox}

\begin{center}\rule{0.5\linewidth}{0.5pt}\end{center}

\section{Prueba t Pareada}\label{sec-semana08-paired-test}

\subsection{Motivación y Formulación}\label{motivacion-y-formulacion}

En muchos estudios, las observaciones de dos muestras están
\textbf{apareadas} (matched): - Medidas antes y después en los mismos
sujetos - Observaciones de hermanos gemelos - Medidas en ambos ojos del
mismo paciente

\begin{tcolorbox}[enhanced jigsaw, breakable, titlerule=0mm, rightrule=.15mm, leftrule=.75mm, left=2mm, opacitybacktitle=0.6, colback=white, bottomtitle=1mm, colbacktitle=quarto-callout-note-color!10!white, title=\textcolor{quarto-callout-note-color}{\faInfo}\hspace{0.5em}{Definición: Datos Pareados}, toptitle=1mm, coltitle=black, opacityback=0, arc=.35mm, bottomrule=.15mm, toprule=.15mm, colframe=quarto-callout-note-color-frame]

Para datos pareados \((X_{1i}, X_{2i})\), \(i = 1, \ldots, n\),
definimos las diferencias:

\[D_i = X_{1i} - X_{2i}\]

El test pareado \textbf{reduce el problema a una prueba t de una
muestra} sobre las diferencias:

\[H_0: \mu_D = \mu_1 - \mu_2 = \omega_0 \quad \text{vs.} \quad H_1: \mu_D \neq \omega_0\]

(Típicamente, \(\omega_0 = 0\))

\end{tcolorbox}

\begin{tcolorbox}[enhanced jigsaw, breakable, titlerule=0mm, rightrule=.15mm, leftrule=.75mm, left=2mm, opacitybacktitle=0.6, colback=white, bottomtitle=1mm, colbacktitle=quarto-callout-warning-color!10!white, title=\textcolor{quarto-callout-warning-color}{\faExclamationTriangle}\hspace{0.5em}{Estadístico t Pareado}, toptitle=1mm, coltitle=black, opacityback=0, arc=.35mm, bottomrule=.15mm, toprule=.15mm, colframe=quarto-callout-warning-color-frame]

\[T = \frac{\bar{D} - \omega_0}{S_D / \sqrt{n}} \sim t_{n-1}\]

donde:
\[\bar{D} = \frac{1}{n} \sum_{i=1}^n D_i, \quad S_D^2 = \frac{1}{n-1} \sum_{i=1}^n (D_i - \bar{D})^2\]

\textbf{Nota importante:} La varianza de las diferencias es típicamente
\textbf{menor} que la suma de varianzas, lo que da mayor potencia al
test pareado cuando los datos están realmente correlacionados.

\end{tcolorbox}

\begin{center}\rule{0.5\linewidth}{0.5pt}\end{center}

\section{Prueba F para Igualdad de Varianzas}\label{sec-semana08-f-test}

\subsection{Hipótesis y Estadístico}\label{hipotesis-y-estadistico}

En algunos análisis necesitamos probar si \textbf{dos varianzas
poblacionales son iguales}.

\begin{tcolorbox}[enhanced jigsaw, breakable, titlerule=0mm, rightrule=.15mm, leftrule=.75mm, left=2mm, opacitybacktitle=0.6, colback=white, bottomtitle=1mm, colbacktitle=quarto-callout-note-color!10!white, title=\textcolor{quarto-callout-note-color}{\faInfo}\hspace{0.5em}{Definición: Prueba F de Fisher}, toptitle=1mm, coltitle=black, opacityback=0, arc=.35mm, bottomrule=.15mm, toprule=.15mm, colframe=quarto-callout-note-color-frame]

Consideramos: - \textbf{Bilateral:} \(H_0: \sigma_1^2 = \sigma_2^2\)
vs.~\(H_1: \sigma_1^2 \neq \sigma_2^2\) - \textbf{Unilateral derecha:}
\(H_0: \sigma_1^2 \leq \sigma_2^2\) vs.~\(H_1: \sigma_1^2 > \sigma_2^2\)
- \textbf{Unilateral izquierda:} \(H_0: \sigma_1^2 \geq \sigma_2^2\)
vs.~\(H_1: \sigma_1^2 < \sigma_2^2\)

\end{tcolorbox}

\begin{tcolorbox}[enhanced jigsaw, breakable, titlerule=0mm, rightrule=.15mm, leftrule=.75mm, left=2mm, opacitybacktitle=0.6, colback=white, bottomtitle=1mm, colbacktitle=quarto-callout-warning-color!10!white, title=\textcolor{quarto-callout-warning-color}{\faExclamationTriangle}\hspace{0.5em}{Estadístico F}, toptitle=1mm, coltitle=black, opacityback=0, arc=.35mm, bottomrule=.15mm, toprule=.15mm, colframe=quarto-callout-warning-color-frame]

Bajo \(H_0\) (\(\sigma_1^2 = \sigma_2^2\)):

\[F = \frac{S_1^2}{S_2^2} \sim F_{n_1-1, n_2-1}\]

\textbf{Criterios de rechazo:} - \textbf{Bilateral:} Rechazo si
\(F < F_{n_1-1, n_2-1, \alpha/2}\) o
\(F > F_{n_1-1, n_2-1, 1-\alpha/2}\) - \textbf{Unilateral derecha:}
Rechazo si \(F > F_{n_1-1, n_2-1, 1-\alpha}\) - \textbf{Unilateral
izquierda:} Rechazo si \(F < F_{n_1-1, n_2-1, \alpha}\)

\end{tcolorbox}

\begin{tcolorbox}[enhanced jigsaw, breakable, titlerule=0mm, rightrule=.15mm, leftrule=.75mm, left=2mm, opacitybacktitle=0.6, colback=white, bottomtitle=1mm, colbacktitle=quarto-callout-warning-color!10!white, title=\textcolor{quarto-callout-warning-color}{\faExclamationTriangle}\hspace{0.5em}{Advertencia: Testing Secuencial}, toptitle=1mm, coltitle=black, opacityback=0, arc=.35mm, bottomrule=.15mm, toprule=.15mm, colframe=quarto-callout-warning-color-frame]

\textbf{No recomendamos realizar la prueba F antes de elegir entre test
t con varianzas iguales o desiguales}, porque esta prueba de igualdad
consume el nivel de significancia global. Es preferir usar siempre la
\textbf{Prueba de Welch} como default, que es válida en ambos casos.

\end{tcolorbox}

\begin{center}\rule{0.5\linewidth}{0.5pt}\end{center}

\section{Potencia de una Prueba}\label{sec-semana08-power}

\subsection{Conceptos Fundamentales}\label{conceptos-fundamentales-2}

La \textbf{potencia} es un concepto crucial en el diseño y evaluación de
estudios estadísticos.

\begin{tcolorbox}[enhanced jigsaw, breakable, titlerule=0mm, rightrule=.15mm, leftrule=.75mm, left=2mm, opacitybacktitle=0.6, colback=white, bottomtitle=1mm, colbacktitle=quarto-callout-note-color!10!white, title=\textcolor{quarto-callout-note-color}{\faInfo}\hspace{0.5em}{Definición: Función de Potencia}, toptitle=1mm, coltitle=black, opacityback=0, arc=.35mm, bottomrule=.15mm, toprule=.15mm, colframe=quarto-callout-note-color-frame]

La \textbf{función de potencia} \(G(\theta)\) es la probabilidad de
rechazar \(H_0\) cuando el parámetro verdadero es \(\theta\):

\[G(\theta) = P(\text{rechazar } H_0 | \theta)\]

\textbf{Interpretación:} - Si \(\theta \in \Theta_0\) (región de
\(H_0\)): \(G(\theta) = P(\text{Error Tipo I}) \leq \alpha\) - Si
\(\theta \in \Theta_1\) (región de \(H_1\)):
\(G(\theta) = P(\text{Rechazo correcto}) = 1 - \beta(\theta)\)

donde \(\beta(\theta)\) es la probabilidad del \textbf{Error Tipo II}.

\end{tcolorbox}

\subsection{Relación con Errores de Tipo I y
II}\label{relacion-con-errores-de-tipo-i-y-ii}

\begin{longtable}[]{@{}
  >{\raggedright\arraybackslash}p{(\linewidth - 4\tabcolsep) * \real{0.4167}}
  >{\raggedright\arraybackslash}p{(\linewidth - 4\tabcolsep) * \real{0.2083}}
  >{\raggedright\arraybackslash}p{(\linewidth - 4\tabcolsep) * \real{0.3750}}@{}}
\toprule\noalign{}
\begin{minipage}[b]{\linewidth}\raggedright
Realidad
\end{minipage} & \begin{minipage}[b]{\linewidth}\raggedright
Rechazamos \(H_0\)
\end{minipage} & \begin{minipage}[b]{\linewidth}\raggedright
No rechazamos \(H_0\)
\end{minipage} \\
\midrule\noalign{}
\endhead
\bottomrule\noalign{}
\endlastfoot
\(H_0\) es verdadera & Error Tipo I (prob. \(\alpha\)) & Decisión
correcta \\
\(H_1\) es verdadera & Decisión correcta (prob. \(1-\beta\)) & Error
Tipo II (prob. \(\beta\)) \\
\end{longtable}

\begin{itemize}
\tightlist
\item
  \textbf{Significancia (\(\alpha\)):} Probabilidad de Error Tipo I.
  Típicamente 0.05.
\item
  \textbf{Potencia (\(1-\beta\)):} Probabilidad de detectar una
  diferencia real. Típicamente 0.80 o 0.90.
\end{itemize}

\subsection{Derivación: Potencia para Prueba Bilateral (Varianza
Conocida)}\label{derivacion-potencia-para-prueba-bilateral-varianza-conocida}

\begin{tcolorbox}[enhanced jigsaw, breakable, titlerule=0mm, rightrule=.15mm, leftrule=.75mm, left=2mm, opacitybacktitle=0.6, colback=white, bottomtitle=1mm, colbacktitle=quarto-callout-warning-color!10!white, title=\textcolor{quarto-callout-warning-color}{\faExclamationTriangle}\hspace{0.5em}{Potencia: Prueba Bilateral}, toptitle=1mm, coltitle=black, opacityback=0, arc=.35mm, bottomrule=.15mm, toprule=.15mm, colframe=quarto-callout-warning-color-frame]

Sea \(\mu_0\) el valor hipotetizado, \(\sigma\) conocida, tamaño
muestral \(n\) y nivel de significancia \(\alpha\).

Para una prueba bilateral con \(H_0: \mu = \mu_0\)
vs.~\(H_1: \mu \neq \mu_0\):

\[G(\mu) = 1 - P\left(z_{1-\alpha/2} - \frac{\mu - \mu_0}{\sigma/\sqrt{n}} \leq Z \leq z_{1-\alpha/2} - \frac{\mu - \mu_0}{\sigma/\sqrt{n}}\right)\]

Equivalentemente:

\[G(\mu) = P\left(Z \leq -z_{1-\alpha/2} - \frac{\mu - \mu_0}{\sigma/\sqrt{n}}\right) + P\left(Z > z_{1-\alpha/2} - \frac{\mu - \mu_0}{\sigma/\sqrt{n}}\right)\]

\textbf{Casos extremos:} - En \(\mu = \mu_0\): \(G(\mu_0) = \alpha\)
(tasa de Error Tipo I) - Conforme \(|\mu - \mu_0|\) aumenta:
\(G(\mu) \to 1\) (potencia tiende a 1)

\end{tcolorbox}

\subsection{Potencia: Prueba Unilateral
Derecha}\label{potencia-prueba-unilateral-derecha}

\begin{tcolorbox}[enhanced jigsaw, breakable, titlerule=0mm, rightrule=.15mm, leftrule=.75mm, left=2mm, opacitybacktitle=0.6, colback=white, bottomtitle=1mm, colbacktitle=quarto-callout-warning-color!10!white, title=\textcolor{quarto-callout-warning-color}{\faExclamationTriangle}\hspace{0.5em}{Potencia: Prueba Unilateral Derecha}, toptitle=1mm, coltitle=black, opacityback=0, arc=.35mm, bottomrule=.15mm, toprule=.15mm, colframe=quarto-callout-warning-color-frame]

Para \(H_0: \mu \leq \mu_0\) vs.~\(H_1: \mu > \mu_0\):

\[G(\mu) = 1 - P\left(Z \leq z_{1-\alpha} - \frac{\mu - \mu_0}{\sigma/\sqrt{n}}\right)\]

La potencia aumenta cuando \(\mu > \mu_0\), alcanzando su máximo cuando
\(\mu\) es mucho mayor que \(\mu_0\).

\end{tcolorbox}

\subsection{Potencia: Prueba Unilateral
Izquierda}\label{potencia-prueba-unilateral-izquierda}

\begin{tcolorbox}[enhanced jigsaw, breakable, titlerule=0mm, rightrule=.15mm, leftrule=.75mm, left=2mm, opacitybacktitle=0.6, colback=white, bottomtitle=1mm, colbacktitle=quarto-callout-warning-color!10!white, title=\textcolor{quarto-callout-warning-color}{\faExclamationTriangle}\hspace{0.5em}{Potencia: Prueba Unilateral Izquierda}, toptitle=1mm, coltitle=black, opacityback=0, arc=.35mm, bottomrule=.15mm, toprule=.15mm, colframe=quarto-callout-warning-color-frame]

Para \(H_0: \mu \geq \mu_0\) vs.~\(H_1: \mu < \mu_0\):

\[G(\mu) = P\left(Z \leq -z_{1-\alpha} - \frac{\mu - \mu_0}{\sigma/\sqrt{n}}\right)\]

La potencia aumenta cuando \(\mu < \mu_0\), alcanzando su máximo cuando
\(\mu\) es mucho menor que \(\mu_0\).

\end{tcolorbox}

\subsection{Factores que Afectan la
Potencia}\label{factores-que-afectan-la-potencia}

La potencia de una prueba depende de:

\begin{enumerate}
\def\labelenumi{\arabic{enumi}.}
\tightlist
\item
  \textbf{Tamaño muestral (\(n\)):} Aumentar \(n\) aumenta la potencia.
\item
  \textbf{Magnitud del efecto (\(|\mu - \mu_0|\) o \(\delta\)):}
  Diferencias más grandes son más fáciles de detectar.
\item
  \textbf{Variabilidad (\(\sigma\)):} Menor variabilidad aumenta la
  potencia.
\item
  \textbf{Nivel de significancia (\(\alpha\)):} Aumentar \(\alpha\)
  aumenta la potencia, pero incrementa el Error Tipo I.
\end{enumerate}

\begin{center}\rule{0.5\linewidth}{0.5pt}\end{center}

\section{Potencia para Pruebas de
Proporción}\label{sec-semana08-power-proportion}

\subsection{Formulación}\label{formulacion}

Para una prueba binomial sobre proporción \(\pi\):

\begin{tcolorbox}[enhanced jigsaw, breakable, titlerule=0mm, rightrule=.15mm, leftrule=.75mm, left=2mm, opacitybacktitle=0.6, colback=white, bottomtitle=1mm, colbacktitle=quarto-callout-note-color!10!white, title=\textcolor{quarto-callout-note-color}{\faInfo}\hspace{0.5em}{Potencia: Prueba de Proporción}, toptitle=1mm, coltitle=black, opacityback=0, arc=.35mm, bottomrule=.15mm, toprule=.15mm, colframe=quarto-callout-note-color-frame]

Sea \(\pi_0\) la proporción hipotetizada, \(n\) el tamaño muestral,
\(\alpha\) el nivel de significancia.

Para una prueba bilateral con \(H_0: \pi = \pi_0\)
vs.~\(H_1: \pi \neq \pi_0\):

\[G(\pi) = P(X \leq x_l | \pi) + P(X > x_u | \pi)\]

donde \(x_l\) y \(x_u\) son los valores críticos determinados por el
nivel de significancia \(\alpha\), y \(X \sim \text{Binomial}(n, \pi)\).

\begin{itemize}
\tightlist
\item
  \textbf{En \(\pi = \pi_0\):} \(G(\pi_0) = \alpha\) (tasa de Error Tipo
  I)
\item
  \textbf{Cuando \(\pi\) se aleja de \(\pi_0\):} La potencia \(G(\pi)\)
  aumenta
\end{itemize}

\end{tcolorbox}

\begin{tcolorbox}[enhanced jigsaw, breakable, titlerule=0mm, rightrule=.15mm, leftrule=.75mm, left=2mm, opacitybacktitle=0.6, colback=white, bottomtitle=1mm, colbacktitle=quarto-callout-tip-color!10!white, title=\textcolor{quarto-callout-tip-color}{\faLightbulb}\hspace{0.5em}{Ejemplo: Proporción de Éxitos en Prueba Binomial}, toptitle=1mm, coltitle=black, opacityback=0, arc=.35mm, bottomrule=.15mm, toprule=.15mm, colframe=quarto-callout-tip-color-frame]

Supongamos \(n = 10\), \(\pi_0 = 0.5\), \(\alpha = 0.05\). Los valores
críticos son \(x_l = 2\) y \(x_u = 8\).

Potencia en varios valores de \(\pi\):

\begin{longtable}[]{@{}lll@{}}
\toprule\noalign{}
\(\pi\) & \(G(\pi)\) & \(\beta = 1 - G(\pi)\) \\
\midrule\noalign{}
\endhead
\bottomrule\noalign{}
\endlastfoot
0.20 & 0.3758 & 0.6242 \\
0.30 & 0.1495 & 0.8505 \\
0.40 & 0.0480 & 0.9520 \\
0.50 & 0.0214 & 0.9786 (Error Tipo II) \\
0.60 & 0.0480 & 0.9520 \\
0.70 & 0.1495 & 0.8505 \\
\end{longtable}

La potencia es máxima en los extremos de la distribución (\(\pi\)
cercano a 0 o 1). Esto representa la dificultad de detectar desviaciones
pequeñas de \(\pi_0 = 0.5\) debido a la máxima variabilidad en ese punto
(varianza máxima \(p(1-p) = 0.25\)). En contraste, para \(\pi\) cerca de
0 o 1, la variabilidad es menor, lo que facilita la detección de
diferencias.

\end{tcolorbox}

\section{Cálculo de Tamaño Muestral}\label{sec-semana08-sample-size}

\subsection{Principio General}\label{principio-general}

El tamaño muestral se calcula especificando:

\begin{tcolorbox}[enhanced jigsaw, breakable, titlerule=0mm, rightrule=.15mm, leftrule=.75mm, left=2mm, opacitybacktitle=0.6, colback=white, bottomtitle=1mm, colbacktitle=quarto-callout-note-color!10!white, title=\textcolor{quarto-callout-note-color}{\faInfo}\hspace{0.5em}{Elementos para Calcular Tamaño Muestral}, toptitle=1mm, coltitle=black, opacityback=0, arc=.35mm, bottomrule=.15mm, toprule=.15mm, colframe=quarto-callout-note-color-frame]

\begin{enumerate}
\def\labelenumi{\arabic{enumi}.}
\tightlist
\item
  \textbf{Diferencia mínima relevante (\(\delta\)):} Qué magnitud de
  efecto queremos detectar
\item
  \textbf{Nivel de significancia (\(\alpha\)):} Típicamente 0.05
\item
  \textbf{Potencia deseada (\(1 - \beta\)):} Típicamente 0.80 o 0.90
\item
  \textbf{Desviación estándar (\(\sigma\)):} De estudios previos o
  piloto
\end{enumerate}

\end{tcolorbox}

\subsection{Fórmula para Prueba t de Dos Muestras (Varianzas
Iguales)}\label{formula-para-prueba-t-de-dos-muestras-varianzas-iguales}

\begin{tcolorbox}[enhanced jigsaw, breakable, titlerule=0mm, rightrule=.15mm, leftrule=.75mm, left=2mm, opacitybacktitle=0.6, colback=white, bottomtitle=1mm, colbacktitle=quarto-callout-warning-color!10!white, title=\textcolor{quarto-callout-warning-color}{\faExclamationTriangle}\hspace{0.5em}{Tamaño Muestral: Prueba t Bilateral}, toptitle=1mm, coltitle=black, opacityback=0, arc=.35mm, bottomrule=.15mm, toprule=.15mm, colframe=quarto-callout-warning-color-frame]

Para detectar una diferencia \(\delta = \mu_1 - \mu_2\) entre dos
muestras de igual tamaño:

\[n = \frac{2\sigma^2 (z_{1-\alpha/2} + z_{1-\beta})^2}{\delta^2}\]

donde: - \(z_{1-\alpha/2}\) es el quantil normal para nivel de
significancia \(\alpha\) - \(z_{1-\beta}\) es el quantil normal para
potencia \(1-\beta\) - \(\sigma^2\) es la varianza (asumida igual en
ambas muestras)

\textbf{Ejemplo:} Si \(\delta = 10\), \(\sigma = 15\),
\(\alpha = 0.05\), \(1-\beta = 0.80\):
\[z_{1-0.025} = 1.96, \quad z_{0.80} = 0.84\]
\[n = \frac{2 \times 225 \times (1.96 + 0.84)^2}{100} = \frac{450 \times 6.4}{100} \approx 28.8 \approx 29 \text{ por grupo}\]

\end{tcolorbox}

\subsection{\texorpdfstring{Cálculo con BioStatR: \texttt{testt()} y
\texttt{testp()}}{Cálculo con BioStatR: testt() y testp()}}\label{sec-semana08-bio-power}

Las funciones \texttt{testt()} y \texttt{testp()} del paquete
\texttt{BioStatR} (ver Sección~\ref{sec-bio-contrastes}) calculan
\textbf{simultáneamente} el contraste de hipótesis, la \textbf{potencia}
y el \textbf{tamaño muestral} necesario para detectar una diferencia
mínima relevante \(\delta\) (effect size en unidades de la variable).

\begin{tcolorbox}[enhanced jigsaw, breakable, titlerule=0mm, rightrule=.15mm, leftrule=.75mm, left=2mm, opacitybacktitle=0.6, colback=white, bottomtitle=1mm, colbacktitle=quarto-callout-note-color!10!white, title=\textcolor{quarto-callout-note-color}{\faInfo}\hspace{0.5em}{Argumentos clave}, toptitle=1mm, coltitle=black, opacityback=0, arc=.35mm, bottomrule=.15mm, toprule=.15mm, colframe=quarto-callout-note-color-frame]

\begin{longtable}[]{@{}
  >{\raggedright\arraybackslash}p{(\linewidth - 4\tabcolsep) * \real{0.3438}}
  >{\raggedright\arraybackslash}p{(\linewidth - 4\tabcolsep) * \real{0.3750}}
  >{\raggedright\arraybackslash}p{(\linewidth - 4\tabcolsep) * \real{0.2812}}@{}}
\toprule\noalign{}
\begin{minipage}[b]{\linewidth}\raggedright
Argumento
\end{minipage} & \begin{minipage}[b]{\linewidth}\raggedright
Significado
\end{minipage} & \begin{minipage}[b]{\linewidth}\raggedright
Función
\end{minipage} \\
\midrule\noalign{}
\endhead
\bottomrule\noalign{}
\endlastfoot
\texttt{delta} & Diferencia mínima clínicamente relevante a detectar &
\texttt{testt}, \texttt{testp} \\
\texttt{potencia} & Potencia objetivo \(1-\beta\) (típicamente 0.80 o
0.90) & \texttt{testt}, \texttt{testp} \\
\texttt{alfa} & Nivel de significancia \(\alpha\) (por defecto 0.05) &
\texttt{testt}, \texttt{testp} \\
\texttt{m}, \texttt{s}, \texttt{n} & Media, DT y tamaño piloto
(estimaciones previas) & \texttt{testt} \\
\texttt{x}, \texttt{n}, \texttt{p0} & Éxitos, tamaño piloto y proporción
hipotética & \texttt{testp} \\
\end{longtable}

Cuando se proporciona \texttt{delta} y \texttt{potencia}, la función
devuelve el \textbf{n necesario}. Si en cambio se fija \texttt{n} y
\texttt{delta}, devuelve la \textbf{potencia alcanzada}.

\end{tcolorbox}

\begin{tcolorbox}[enhanced jigsaw, breakable, titlerule=0mm, rightrule=.15mm, leftrule=.75mm, left=2mm, opacitybacktitle=0.6, colback=white, bottomtitle=1mm, colbacktitle=quarto-callout-tip-color!10!white, title=\textcolor{quarto-callout-tip-color}{\faLightbulb}\hspace{0.5em}{Ejemplo: Potencia y tamaño muestral en contexto clínico}, toptitle=1mm, coltitle=black, opacityback=0, arc=.35mm, bottomrule=.15mm, toprule=.15mm, colframe=quarto-callout-tip-color-frame]

Dos escenarios reales del estudio piloto \texttt{osteo} (94 pacientes
diabéticos):

\begin{itemize}
\tightlist
\item
  \textbf{Contraste de medias} (\texttt{testt}): ¿Cuántos pacientes
  necesitamos para detectar una diferencia de \(\delta = 0.5\%\) en
  HbA1c respecto al umbral clínico de 7.5\% con potencia del 90\%?
\item
  \textbf{Contraste de proporciones} (\texttt{testp}): ¿Y para detectar
  una desviación de \(\delta = 0.05\) frente a una prevalencia
  hipotética de tabaquismo del 35\% con potencia del 90\%?
\end{itemize}

\begin{codelisting}[H]

\caption{\label{lst-semana09-power}Código R}

\centering{

\begin{Shaded}
\begin{Highlighting}[]
\FunctionTok{library}\NormalTok{(BioStatR)}

\CommentTok{\# {-}{-}{-} 1. testt(): potencia y n para una media {-}{-}{-}}
\CommentTok{\# H0: µ\_HbA1c = 7.5  ;  delta clinicamente relevante = 0.5\%}
\CommentTok{\# Usamos m, s del estudio piloto como referencia}
\FunctionTok{testt}\NormalTok{(}\AttributeTok{m =}\NormalTok{ osteo}\SpecialCharTok{$}\NormalTok{hba1c, }\AttributeTok{m0 =} \FloatTok{7.5}\NormalTok{,}
      \AttributeTok{delta =} \FloatTok{0.5}\NormalTok{, }\AttributeTok{potencia =} \FloatTok{0.90}\NormalTok{, }\AttributeTok{alfa =} \FloatTok{0.05}\NormalTok{, }\AttributeTok{grf =} \ConstantTok{FALSE}\NormalTok{)}
\end{Highlighting}
\end{Shaded}

}

\end{codelisting}%

\begin{verbatim}

# t-Test con una muestra
# ----------------------

# Resumen de 'osteo$hba1c'
    n = 94.000 
    media = 8.565 
    d.t. = 1.799 
    sem = 0.186 

# Estimación de la media μ:
  95%-IC(μ) = (8.196, 8.933) 

# Test de normalidad de Shapiro-Wilk:
  W = 0.983, gl = 94, p = 0.276 

# Test de Student para contrastar H₀:μ=μ₀ con μ₀=7.500
  texp = 5.740, gl = 93 
    p < 0.001 para la alternativa bilateral H₁:μ≠μ₀ 
    p < 0.001 para la alternativa unilateral H₁:μ>μ₀ 

  Estimación del efecto bruto 
  95%-IC(μ-μ₀) = (0.696, 1.433) 

# Estudio de la potencia:
  El test es significativo, se omite el análisis de la potencia. 
\end{verbatim}

\begin{codelisting}[H]

\caption{\label{lst-semana09-power}Código R}

\centering{

\begin{Shaded}
\begin{Highlighting}[]
\CommentTok{\# {-}{-}{-} 2. testp(): potencia y n para una proporcion {-}{-}{-}}
\CommentTok{\# H0: pi\_tabaco = 0.35  ;  delta = 0.05 (5 puntos porcentuales)}
\FunctionTok{testp}\NormalTok{(}\AttributeTok{x =} \DecValTok{41}\NormalTok{, }\AttributeTok{n =} \DecValTok{94}\NormalTok{,}
      \AttributeTok{p0 =} \FloatTok{0.35}\NormalTok{, }\AttributeTok{delta =} \FloatTok{0.05}\NormalTok{, }\AttributeTok{potencia =} \FloatTok{0.90}\NormalTok{, }\AttributeTok{alfa =} \FloatTok{0.05}\NormalTok{)}
\end{Highlighting}
\end{Shaded}

}

\end{codelisting}%

\begin{verbatim}

# Test para contrastar una proporción binomial
# --------------------------------------------

# Información muestral
  n = 94
  x = 41   n-x=53
  p = 0.436; q = (1-p) = 0.564

# Test Ho:π=0.350
  [1] Método exacto
                      H1  Fexp Valor.p
    Cola derecha π>0.350 1.410   0.052
    Bilateral    π≠0.350     -   0.103

    95%-IC(π) = (0.339, 0.542) (método de Clooper-Pearson)

  [2] Método aproximado a la distribución normal
    Validez: min(nπ₀, n(1-π₀)) = 32.9 (>5, el método es válido)
    zexp = 1.643,  p  = 0.100

    95%-IC(π) = (0.335, 0.542) (método de Wilson)


# Tamaño de muestra para detectar |π-p₀|=δ
  p₀ = 0.350;  δ = 0.05
  Casos (cpc = corrección por continuidad):
                            p1   n n_cpc
  1 Unilateral π<0.350   0.300 753   773
  2 Unilateral π>0.350   0.400 753   773
  3 Bilateral  π≠0.350   0.400 977   997
\end{verbatim}

\begin{quote}
\textbf{Interpretación}

Los argumentos \texttt{delta} y \texttt{potencia} transforman un
contraste de hipótesis estático (sólo p-valor) en un \textbf{análisis
prospectivo de diseño}: cuantifican la sensibilidad del estudio para
detectar efectos clínicamente relevantes. Para la HbA1c, detectar una
desviación de medio punto porcentual respecto al umbral clínico (7.5\%)
con potencia del 90\% requiere aproximadamente \textbf{n 138 pacientes}
---superior a los 94 actuales---, lo que indica que el estudio piloto
está \textbf{infrapotenciado} para detectar efectos pequeños pero
clínicamente significativos (0.5\% de HbA1c corresponde a
\textasciitilde10--15 mg/dL de glucemia media, magnitud relevante para
decisiones terapéuticas). Para la prevalencia de tabaquismo, alcanzar
potencia del 90\% para detectar una desviación de 5 puntos porcentuales
respecto al 35\% hipotético exige cerca de \textbf{n ≈ 900 pacientes},
casi diez veces el tamaño piloto. Esta asimetría refleja una propiedad
clave: detectar diferencias en proporciones cercanas a 0.5 (varianza
máxima \(p(1-p) = 0.25\)) consume más muestra que detectar diferencias
en variables continuas con buen contraste señal/ruido. En el diseno de
protocolos clínicos, este análisis previo evita el error de planificar
estudios incapaces de detectar el efecto buscado (estudios negativos por
falta de potencia, no por ausencia real de efecto), y constituye
argumento obligado ante comités éticos y agencias financiadoras.
\end{quote}

\end{tcolorbox}

\begin{center}\rule{0.5\linewidth}{0.5pt}\end{center}

\section{Comparación Múltiple y Corrección de
Bonferroni}\label{sec-semana08-multiple-testing}

\subsection{Problema: Inflación de Tasa de Error Tipo
I}\label{problema-inflacion-de-tasa-de-error-tipo-i}

\begin{tcolorbox}[enhanced jigsaw, breakable, titlerule=0mm, rightrule=.15mm, leftrule=.75mm, left=2mm, opacitybacktitle=0.6, colback=white, bottomtitle=1mm, colbacktitle=quarto-callout-warning-color!10!white, title=\textcolor{quarto-callout-warning-color}{\faExclamationTriangle}\hspace{0.5em}{Problema: Múltiples Pruebas}, toptitle=1mm, coltitle=black, opacityback=0, arc=.35mm, bottomrule=.15mm, toprule=.15mm, colframe=quarto-callout-warning-color-frame]

Cuando realizamos \textbf{múltiples pruebas de hipótesis independientes}
con nivel \(\alpha\):

Si efectuamos \(k\) pruebas, la probabilidad de al menos un Error Tipo I
es:

\[P(\text{al menos un Error Tipo I}) = 1 - (1 - \alpha)^k\]

\textbf{Ejemplo:} Con \(k = 10\) pruebas y \(\alpha = 0.05\):
\[P(\text{al menos un Error Tipo I}) = 1 - (0.95)^{10} \approx 0.401 = 40.1\%\]

En lugar del nivel de significancia global deseado del 5\%, obtenemos
una tasa de Error Tipo I del 40\%!

\end{tcolorbox}

\subsection{Corrección de Bonferroni}\label{correccion-de-bonferroni}

\begin{tcolorbox}[enhanced jigsaw, breakable, titlerule=0mm, rightrule=.15mm, leftrule=.75mm, left=2mm, opacitybacktitle=0.6, colback=white, bottomtitle=1mm, colbacktitle=quarto-callout-tip-color!10!white, title=\textcolor{quarto-callout-tip-color}{\faLightbulb}\hspace{0.5em}{Corrección de Bonferroni}, toptitle=1mm, coltitle=black, opacityback=0, arc=.35mm, bottomrule=.15mm, toprule=.15mm, colframe=quarto-callout-tip-color-frame]

Para mantener una tasa global de Error Tipo I de \(\alpha\) cuando
realizamos \(k\) pruebas:

\[\alpha_{\text{ajustado}} = \frac{\alpha}{k}\]

Es decir, cada prueba individual se realiza con nivel de significancia
\(\alpha/k\).

\textbf{Ventajas:} - Simple de implementar - Conservador (garantiza
control del Error Tipo I)

\textbf{Desventajas:} - Muy conservador cuando \(k\) es grande - Reduce
la potencia de cada prueba individual

\textbf{Alternativas:} El método de Holm, la corrección FDR (False
Discovery Rate), o métodos Bayesianos son menos conservadores.

\end{tcolorbox}

\begin{center}\rule{0.5\linewidth}{0.5pt}\end{center}

\section{Ejemplo Completo: Análisis de Datos de
Mosquitos}\label{sec-semana08-mosquito-example}

\subsection{Contexto del Problema}\label{contexto-del-problema}

Según la Organización Mundial de la Salud, los mosquitos Aedes pueden
transmitir enfermedades como dengue y Zika. Un esfuerzo de control
consiste en liberar mosquitos transgénicos con esperanza de vida más
corta.

\textbf{Pregunta de investigación:} ¿Es la esperanza de vida
significativamente diferente entre mosquitos salvajes (wildtype) y
transgénicos?

\subsection{Datos y Exploración}\label{datos-y-exploracion}

\begin{codelisting}

\caption{\label{lst-semana09-4}Código R}

\centering{

\begin{Shaded}
\begin{Highlighting}[]
\CommentTok{\# Lectura de datos}
\CommentTok{\# Crear datos directamente en R (Tipo de mosquito y vida en días)}
\NormalTok{wildtype }\OtherTok{\textless{}{-}} \FunctionTok{c}\NormalTok{(}\DecValTok{28}\NormalTok{, }\DecValTok{32}\NormalTok{, }\DecValTok{25}\NormalTok{, }\DecValTok{31}\NormalTok{, }\DecValTok{29}\NormalTok{, }\DecValTok{26}\NormalTok{, }\DecValTok{33}\NormalTok{, }\DecValTok{30}\NormalTok{, }\DecValTok{27}\NormalTok{, }\DecValTok{34}\NormalTok{, }\DecValTok{28}\NormalTok{, }\DecValTok{31}\NormalTok{, }\DecValTok{29}\NormalTok{, }\DecValTok{32}\NormalTok{, }\DecValTok{25}\NormalTok{, }\DecValTok{30}\NormalTok{, }\DecValTok{33}\NormalTok{, }\DecValTok{28}\NormalTok{, }\DecValTok{31}\NormalTok{, }\DecValTok{27}\NormalTok{)}
\NormalTok{transgenic }\OtherTok{\textless{}{-}} \FunctionTok{c}\NormalTok{(}\DecValTok{35}\NormalTok{, }\DecValTok{38}\NormalTok{, }\DecValTok{42}\NormalTok{, }\DecValTok{39}\NormalTok{, }\DecValTok{41}\NormalTok{, }\DecValTok{37}\NormalTok{, }\DecValTok{44}\NormalTok{, }\DecValTok{36}\NormalTok{, }\DecValTok{40}\NormalTok{, }\DecValTok{43}\NormalTok{, }\DecValTok{38}\NormalTok{, }\DecValTok{45}\NormalTok{, }\DecValTok{39}\NormalTok{, }\DecValTok{42}\NormalTok{, }\DecValTok{35}\NormalTok{, }\DecValTok{41}\NormalTok{, }\DecValTok{46}\NormalTok{, }\DecValTok{40}\NormalTok{, }\DecValTok{38}\NormalTok{, }\DecValTok{43}\NormalTok{)}

\CommentTok{\# Combina en un data.frame}
\NormalTok{mosquitoes }\OtherTok{\textless{}{-}} \FunctionTok{data.frame}\NormalTok{(wildtype, transgenic)}

\CommentTok{\# Extracción de muestras}
\NormalTok{wildtype }\OtherTok{\textless{}{-}} \FunctionTok{na.omit}\NormalTok{(mosquitoes[, }\DecValTok{1}\NormalTok{])}
\NormalTok{transgenic }\OtherTok{\textless{}{-}} \FunctionTok{na.omit}\NormalTok{(mosquitoes[, }\DecValTok{2}\NormalTok{])}

\NormalTok{n1 }\OtherTok{\textless{}{-}} \FunctionTok{length}\NormalTok{(wildtype)}
\NormalTok{n2 }\OtherTok{\textless{}{-}} \FunctionTok{length}\NormalTok{(transgenic)}

\FunctionTok{cat}\NormalTok{(}\StringTok{"Wildtype: n ="}\NormalTok{, n1, }\StringTok{", media ="}\NormalTok{, }\FunctionTok{mean}\NormalTok{(wildtype), }\StringTok{", sd ="}\NormalTok{, }\FunctionTok{sd}\NormalTok{(wildtype), }\StringTok{"}\SpecialCharTok{\textbackslash{}n}\StringTok{"}\NormalTok{)}
\end{Highlighting}
\end{Shaded}

}

\end{codelisting}%

\begin{verbatim}
Wildtype: n = 20 , media = 29.4 , sd = 2.68 
\end{verbatim}

\begin{codelisting}

\caption{\label{lst-semana09-4}Código R}

\centering{

\begin{Shaded}
\begin{Highlighting}[]
\FunctionTok{cat}\NormalTok{(}\StringTok{"Transgenic: n ="}\NormalTok{, n2, }\StringTok{", media ="}\NormalTok{, }\FunctionTok{mean}\NormalTok{(transgenic), }\StringTok{", sd ="}\NormalTok{, }\FunctionTok{sd}\NormalTok{(transgenic), }\StringTok{"}\SpecialCharTok{\textbackslash{}n}\StringTok{"}\NormalTok{)}
\end{Highlighting}
\end{Shaded}

}

\end{codelisting}%

\begin{verbatim}
Transgenic: n = 20 , media = 40.1 , sd = 3.19 
\end{verbatim}

\begin{codelisting}

\caption{\label{lst-semana09-4}Código R}

\centering{

\begin{Shaded}
\begin{Highlighting}[]
\CommentTok{\# Gráfico de violín con ggplot2}
\FunctionTok{library}\NormalTok{(ggplot2)}
\FunctionTok{library}\NormalTok{(reshape2)}

\NormalTok{df }\OtherTok{\textless{}{-}} \FunctionTok{melt}\NormalTok{(}\FunctionTok{data.frame}\NormalTok{(}\AttributeTok{Wildtype =}\NormalTok{ wildtype, }\AttributeTok{Transgenic =}\NormalTok{ transgenic))}
\FunctionTok{colnames}\NormalTok{(df) }\OtherTok{\textless{}{-}} \FunctionTok{c}\NormalTok{(}\StringTok{"Group"}\NormalTok{, }\StringTok{"Lifespan"}\NormalTok{)}

\FunctionTok{ggplot}\NormalTok{(df, }\FunctionTok{aes}\NormalTok{(}\AttributeTok{x =}\NormalTok{ Group, }\AttributeTok{y =}\NormalTok{ Lifespan)) }\SpecialCharTok{+}
  \FunctionTok{geom\_violin}\NormalTok{(}\AttributeTok{fill =} \StringTok{"lightblue"}\NormalTok{, }\AttributeTok{alpha =} \FloatTok{0.7}\NormalTok{) }\SpecialCharTok{+}
  \FunctionTok{geom\_boxplot}\NormalTok{(}\AttributeTok{width =} \FloatTok{0.15}\NormalTok{, }\AttributeTok{alpha =} \FloatTok{0.5}\NormalTok{) }\SpecialCharTok{+}
  \FunctionTok{stat\_summary}\NormalTok{(}\AttributeTok{fun =}\NormalTok{ mean, }\AttributeTok{geom =} \StringTok{"point"}\NormalTok{, }\AttributeTok{color =} \StringTok{"red"}\NormalTok{, }\AttributeTok{size =} \DecValTok{3}\NormalTok{) }\SpecialCharTok{+}
  \FunctionTok{theme\_minimal}\NormalTok{() }\SpecialCharTok{+}
  \FunctionTok{labs}\NormalTok{(}\AttributeTok{title =} \StringTok{"Longevidad de Mosquitos: Wildtype vs. Transgenic"}\NormalTok{,}
       \AttributeTok{x =} \StringTok{"Tipo"}\NormalTok{, }\AttributeTok{y =} \StringTok{"Esperanza de Vida (días)"}\NormalTok{)}
\end{Highlighting}
\end{Shaded}

}

\end{codelisting}%

\pandocbounded{\includegraphics[keepaspectratio]{chapters/semana08_files/figure-pdf/unnamed-chunk-3-1.pdf}}

\subsection{Test t: Varianzas
Homogéneas}\label{test-t-varianzas-homogeneas}

Primero suponemos varianzas iguales:

\begin{codelisting}

\caption{\label{lst-semana09-5}Código R}

\centering{

\begin{Shaded}
\begin{Highlighting}[]
\FunctionTok{options}\NormalTok{(}\AttributeTok{scipen =} \DecValTok{999}\NormalTok{) }\CommentTok{\# Evitar notación científica)}
\CommentTok{\# Test t con varianzas iguales}
\NormalTok{tt\_equal }\OtherTok{\textless{}{-}} \FunctionTok{t.test}\NormalTok{(wildtype, transgenic,}
                   \AttributeTok{alternative =} \StringTok{"two.sided"}\NormalTok{,}
                   \AttributeTok{var.equal =} \ConstantTok{TRUE}\NormalTok{,}
                   \AttributeTok{conf.level =} \FloatTok{0.95}\NormalTok{)}

\FunctionTok{print}\NormalTok{(tt\_equal)}
\end{Highlighting}
\end{Shaded}

}

\end{codelisting}%

\begin{verbatim}

    Two Sample t-test

data:  wildtype and transgenic
t = -11, df = 38, p-value = 0.00000000000008
alternative hypothesis: true difference in means is not equal to 0
95 percent confidence interval:
 -12.54  -8.76
sample estimates:
mean of x mean of y 
     29.4      40.1 
\end{verbatim}

\begin{codelisting}

\caption{\label{lst-semana09-5}Código R}

\centering{

\begin{Shaded}
\begin{Highlighting}[]
\CommentTok{\# Interpretación:}
\CommentTok{\# Si p{-}value \textless{} 0.05, rechazamos H\_0 (las medias son significativamente diferentes)}
\CommentTok{\# Si p{-}value \textgreater{}= 0.05, no hay evidencia de diferencia significativa}
\end{Highlighting}
\end{Shaded}

}

\end{codelisting}%

\subsection{Test t: Varianzas Heterogéneas
(Welch)}\label{test-t-varianzas-heterogeneas-welch}

Si observamos varianzas muy diferentes, usamos la prueba de Welch:

\begin{codelisting}

\caption{\label{lst-semana09-6}Código R}

\centering{

\begin{Shaded}
\begin{Highlighting}[]
\CommentTok{\# Test t con varianzas desiguales (Welch)}
\NormalTok{tt\_welch }\OtherTok{\textless{}{-}} \FunctionTok{t.test}\NormalTok{(wildtype, transgenic,}
                   \AttributeTok{alternative =} \StringTok{"two.sided"}\NormalTok{,}
                   \AttributeTok{var.equal =} \ConstantTok{FALSE}\NormalTok{,}
                   \AttributeTok{conf.level =} \FloatTok{0.95}\NormalTok{)}

\FunctionTok{print}\NormalTok{(tt\_welch)}
\end{Highlighting}
\end{Shaded}

}

\end{codelisting}%

\begin{verbatim}

    Welch Two Sample t-test

data:  wildtype and transgenic
t = -11, df = 37, p-value = 0.0000000000001
alternative hypothesis: true difference in means is not equal to 0
95 percent confidence interval:
 -12.54  -8.76
sample estimates:
mean of x mean of y 
     29.4      40.1 
\end{verbatim}

\subsection{Test F: Igualdad de
Varianzas}\label{test-f-igualdad-de-varianzas}

Podemos probar formalmente si las varianzas son iguales:

\begin{codelisting}

\caption{\label{lst-semana09-7}Código R}

\centering{

\begin{Shaded}
\begin{Highlighting}[]
\CommentTok{\# Test F para igualdad de varianzas}
\NormalTok{f\_test }\OtherTok{\textless{}{-}} \FunctionTok{var.test}\NormalTok{(wildtype, transgenic, }\AttributeTok{alternative =} \StringTok{"two.sided"}\NormalTok{)}

\FunctionTok{print}\NormalTok{(f\_test)}
\end{Highlighting}
\end{Shaded}

}

\end{codelisting}%

\begin{verbatim}

    F test to compare two variances

data:  wildtype and transgenic
F = 0.7, num df = 19, denom df = 19, p-value = 0.5
alternative hypothesis: true ratio of variances is not equal to 1
95 percent confidence interval:
 0.28 1.79
sample estimates:
ratio of variances 
             0.707 
\end{verbatim}

\begin{codelisting}

\caption{\label{lst-semana09-7}Código R}

\centering{

\begin{Shaded}
\begin{Highlighting}[]
\CommentTok{\# Interpretación:}
\CommentTok{\# Si p{-}value \textless{} 0.05, hay evidencia de que las varianzas son desiguales}
\CommentTok{\#   → Usar Welch test}
\CommentTok{\# Si p{-}value \textgreater{}= 0.05, no hay evidencia de desigualdad de varianzas}
\CommentTok{\#   → Usar test t con varianzas iguales}
\end{Highlighting}
\end{Shaded}

}

\end{codelisting}%

\subsection{Cálculo de Potencia}\label{calculo-de-potencia}

Si el estudio fuera diseñado ahora, ¿qué tamaño muestral necesitaríamos?

\begin{codelisting}

\caption{\label{lst-semana09-8}Código R}

\centering{

\begin{Shaded}
\begin{Highlighting}[]
\FunctionTok{library}\NormalTok{(stats)}

\CommentTok{\# Supongamos que queremos detectar una diferencia de 5 días}
\CommentTok{\# con desviación estándar combinada de 10 días}
\CommentTok{\# Potencia deseada = 0.80, significancia = 0.05}

\NormalTok{power\_result }\OtherTok{\textless{}{-}} \FunctionTok{power.t.test}\NormalTok{(}
  \AttributeTok{n =} \ConstantTok{NULL}\NormalTok{,           }\CommentTok{\# Queremos calcular n}
  \AttributeTok{delta =} \DecValTok{5}\NormalTok{,          }\CommentTok{\# Diferencia a detectar}
  \AttributeTok{sd =} \DecValTok{10}\NormalTok{,            }\CommentTok{\# Desviación estándar}
  \AttributeTok{sig.level =} \FloatTok{0.05}\NormalTok{,   }\CommentTok{\# Significancia bilateral}
  \AttributeTok{power =} \FloatTok{0.80}\NormalTok{,       }\CommentTok{\# Potencia deseada}
  \AttributeTok{alternative =} \StringTok{"two.sided"}
\NormalTok{)}

\FunctionTok{cat}\NormalTok{(}\StringTok{"Tamaño muestral requerido por grupo:"}\NormalTok{, }\FunctionTok{round}\NormalTok{(power\_result}\SpecialCharTok{$}\NormalTok{n), }\StringTok{"}\SpecialCharTok{\textbackslash{}n}\StringTok{"}\NormalTok{)}
\end{Highlighting}
\end{Shaded}

}

\end{codelisting}%

\begin{verbatim}
Tamaño muestral requerido por grupo: 64 
\end{verbatim}

\begin{codelisting}

\caption{\label{lst-semana09-8}Código R}

\centering{

\begin{Shaded}
\begin{Highlighting}[]
\CommentTok{\# También podemos calcular la potencia de nuestro estudio actual}
\NormalTok{power\_actual }\OtherTok{\textless{}{-}} \FunctionTok{power.t.test}\NormalTok{(}
  \AttributeTok{n =}\NormalTok{ n1,             }\CommentTok{\# Tamaño muestral observado}
  \AttributeTok{delta =} \FunctionTok{mean}\NormalTok{(wildtype) }\SpecialCharTok{{-}} \FunctionTok{mean}\NormalTok{(transgenic),  }\CommentTok{\# Diferencia observada}
  \AttributeTok{sd =} \FunctionTok{sqrt}\NormalTok{((}\FunctionTok{var}\NormalTok{(wildtype) }\SpecialCharTok{+} \FunctionTok{var}\NormalTok{(transgenic)) }\SpecialCharTok{/} \DecValTok{2}\NormalTok{),  }\CommentTok{\# SD combinada}
  \AttributeTok{sig.level =} \FloatTok{0.05}\NormalTok{,}
  \AttributeTok{alternative =} \StringTok{"two.sided"}
\NormalTok{)}

\FunctionTok{cat}\NormalTok{(}\StringTok{"Potencia del estudio observado:"}\NormalTok{, }\FunctionTok{round}\NormalTok{(power\_actual}\SpecialCharTok{$}\NormalTok{power, }\DecValTok{3}\NormalTok{), }\StringTok{"}\SpecialCharTok{\textbackslash{}n}\StringTok{"}\NormalTok{)}
\end{Highlighting}
\end{Shaded}

}

\end{codelisting}%

\begin{verbatim}
Potencia del estudio observado: 1 
\end{verbatim}

\begin{center}\rule{0.5\linewidth}{0.5pt}\end{center}

\section{Tabla Resumen: Contrastes de Dos
Muestras}\label{sec-semana08-summary-table}

\begin{longtable}[]{@{}
  >{\raggedright\arraybackslash}p{(\linewidth - 8\tabcolsep) * \real{0.1667}}
  >{\raggedright\arraybackslash}p{(\linewidth - 8\tabcolsep) * \real{0.1970}}
  >{\raggedright\arraybackslash}p{(\linewidth - 8\tabcolsep) * \real{0.2121}}
  >{\raggedright\arraybackslash}p{(\linewidth - 8\tabcolsep) * \real{0.2879}}
  >{\raggedright\arraybackslash}p{(\linewidth - 8\tabcolsep) * \real{0.1364}}@{}}
\toprule\noalign{}
\begin{minipage}[b]{\linewidth}\raggedright
Escenario
\end{minipage} & \begin{minipage}[b]{\linewidth}\raggedright
Estadístico
\end{minipage} & \begin{minipage}[b]{\linewidth}\raggedright
Distribución
\end{minipage} & \begin{minipage}[b]{\linewidth}\raggedright
Grados de Libertad
\end{minipage} & \begin{minipage}[b]{\linewidth}\raggedright
Fórmula
\end{minipage} \\
\midrule\noalign{}
\endhead
\bottomrule\noalign{}
\endlastfoot
\textbf{Z-test} (varianzas conocidas) &
\(Z = \frac{(\bar{X}_1 - \bar{X}_2) - \omega_0}{\sqrt{\sigma_1^2/n_1 + \sigma_2^2/n_2}}\)
& \(N(0,1)\) & ∞ & Conocidas \\
\textbf{t-test} (varianzas iguales) &
\(T = \frac{(\bar{X}_1 - \bar{X}_2) - \omega_0}{S_p\sqrt{1/n_1 + 1/n_2}}\)
& \(t\) & \(n_1 + n_2 - 2\) & Estimadas \\
\textbf{Welch} (varianzas desiguales) &
\(T = \frac{(\bar{X}_1 - \bar{X}_2) - \omega_0}{\sqrt{S_1^2/n_1 + S_2^2/n_2}}\)
& \(t\) (aprox.) & W-S formula & Estimadas \\
\textbf{t pareado} & \(T = \frac{\bar{D} - \omega_0}{S_D/\sqrt{n}}\) &
\(t\) & \(n - 1\) & Diferencias \\
\textbf{F-test} & \(F = \frac{S_1^2}{S_2^2}\) & \(F\) &
\((n_1-1, n_2-1)\) & Varianzas \\
\end{longtable}

\begin{center}\rule{0.5\linewidth}{0.5pt}\end{center}

\section{Resumen}\label{sec-semana08-resumen}

\subsection{Conceptos Clave}\label{conceptos-clave-1}

\begin{itemize}
\item
  \textbf{Prueba Z de dos muestras:} Para varianzas conocidas,
  estadístico sigue \(N(0,1)\)
\item
  \textbf{Prueba t de dos muestras (varianzas iguales):} Varianza
  combinada (pooled), grados de libertad \(= n_1 + n_2 - 2\)
\item
  \textbf{Prueba t de Welch:} Para varianzas desiguales, usa grados de
  libertad Welch-Satterthwaite
\item
  \textbf{Prueba t pareada:} Reduce a una prueba t de una muestra sobre
  diferencias
\item
  \textbf{Prueba F:} Para comparar varianzas, estadístico
  \(F = S_1^2 / S_2^2 \sim F_{n_1-1, n_2-1}\)
\end{itemize}

\subsection{Potencia Estadística}\label{potencia-estadistica}

\begin{itemize}
\item
  \textbf{Potencia (\(1-\beta\)):} Probabilidad de rechazar \(H_0\)
  cuando \(H_1\) es verdadera (detectar un efecto real)
\item
  \textbf{Error Tipo I (\(\alpha\)):} Probabilidad de rechazar \(H_0\)
  cuando es verdadera (típicamente 0.05)
\item
  \textbf{Error Tipo II (\(\beta\)):} Probabilidad de no rechazar
  \(H_0\) cuando \(H_1\) es verdadera (típicamente 0.10-0.20)
\item
  \textbf{Factores que aumentan potencia:}

  \begin{itemize}
  \tightlist
  \item
    Mayor tamaño muestral (\(n\))
  \item
    Mayor diferencia/efecto a detectar (\(\delta\))
  \item
    Menor variabilidad (\(\sigma\))
  \item
    Mayor nivel de significancia (\(\alpha\)) --- pero aumenta Error
    Tipo I
  \end{itemize}
\end{itemize}

\subsection{Tamaño Muestral}\label{tamano-muestral}

\begin{itemize}
\item
  Se calcula especificando: efecto a detectar, potencia deseada, nivel
  de significancia, variabilidad
\item
  Fórmula para t-test bilateral:
  \(n = \frac{2\sigma^2 (z_{1-\alpha/2} + z_{1-\beta})^2}{\delta^2}\)
\end{itemize}

\subsection{Testing Múltiple}\label{testing-multiple}

\begin{itemize}
\item
  \textbf{Corrección de Bonferroni:} Usar
  \(\alpha_{\text{ajustado}} = \alpha/k\) cuando se hacen \(k\) pruebas
\item
  Evitar prueba de igualdad de varianzas antes de elegir test (mejor
  usar Welch siempre)
\end{itemize}

\begin{center}\rule{0.5\linewidth}{0.5pt}\end{center}

\section{Verificación de Supuestos: Normalidad y
Homogeneidad}\label{sec-semana08-supuestos}

Los contrastes asintóticos basados en el Teorema Central del Límite
---como el t-test y el F-test--- requieren verificar dos condiciones:
(1) normalidad de los datos o residuos, y (2) igualdad de varianzas
entre grupos. Con \(n \geq 60\) el TCL garantiza normalidad asintótica
de la media, pero para muestras menores conviene comprobarlo
formalmente.

\subsection{Test de Shapiro-Wilk:
Normalidad}\label{test-de-shapiro-wilk-normalidad}

El test de Shapiro-Wilk \citep{Shapiro1965} es la prueba de normalidad
más potente para muestras pequeñas y moderadas (\(n \leq 5000\)).

\begin{tcolorbox}[enhanced jigsaw, breakable, titlerule=0mm, rightrule=.15mm, leftrule=.75mm, left=2mm, opacitybacktitle=0.6, colback=white, bottomtitle=1mm, colbacktitle=quarto-callout-note-color!10!white, title=\textcolor{quarto-callout-note-color}{\faInfo}\hspace{0.5em}{Definición: Test de Shapiro-Wilk}, toptitle=1mm, coltitle=black, opacityback=0, arc=.35mm, bottomrule=.15mm, toprule=.15mm, colframe=quarto-callout-note-color-frame]

Dada una muestra \(x_1, \ldots, x_n\), el estadístico es:

\[W = \frac{\left(\displaystyle\sum_{i=1}^{n} a_i\, x_{(i)}\right)^2}{\displaystyle\sum_{i=1}^{n} (x_i - \bar{x})^2}\]

donde \(x_{(i)}\) son los estadísticos de orden y \(a_i\) son
coeficientes tabulados que dependen de la distribución normal estándar.

\begin{itemize}
\tightlist
\item
  \(H_0\): la muestra proviene de una distribución normal
\item
  \(H_1\): la distribución no es normal
\item
  Se rechaza \(H_0\) cuando \(W\) es \textbf{pequeño} (p-valor
  \textless{} \(\alpha\))
\item
  \(W \in (0, 1]\); valores próximos a 1 indican normalidad
\end{itemize}

\end{tcolorbox}

\begin{tcolorbox}[enhanced jigsaw, breakable, titlerule=0mm, rightrule=.15mm, leftrule=.75mm, left=2mm, opacitybacktitle=0.6, colback=white, bottomtitle=1mm, colbacktitle=quarto-callout-tip-color!10!white, title=\textcolor{quarto-callout-tip-color}{\faLightbulb}\hspace{0.5em}{Ejemplo: Test de normalidad en datos clínicos}, toptitle=1mm, coltitle=black, opacityback=0, arc=.35mm, bottomrule=.15mm, toprule=.15mm, colframe=quarto-callout-tip-color-frame]

Usando el conjunto de datos \texttt{osteo} de \texttt{BioStatR}, que
contiene 94 pacientes diabéticos de la Facultad de Medicina de la UGR,
comprobamos la normalidad de la variable edad:

\begin{codelisting}[H]

\caption{\label{lst-semana09-9}Código R}

\centering{

\begin{Shaded}
\begin{Highlighting}[]
\FunctionTok{library}\NormalTok{(BioStatR)}

\CommentTok{\# Test de Shapiro{-}Wilk con BioStatR}
\FunctionTok{testnormal}\NormalTok{(osteo}\SpecialCharTok{$}\NormalTok{edad)}
\end{Highlighting}
\end{Shaded}

}

\end{codelisting}%

\pandocbounded{\includegraphics[keepaspectratio]{chapters/semana08_files/figure-pdf/unnamed-chunk-8-1.pdf}}

\begin{verbatim}

# Test de normalidad de Shapiro-Wilk  
 ------------------------------------- 
   n = 94,  W = 0.907,  p < 0.001 
\end{verbatim}

\begin{codelisting}[H]

\caption{\label{lst-semana09-9}Código R}

\centering{

\begin{Shaded}
\begin{Highlighting}[]
\CommentTok{\# Test de Shapiro{-}Wilk con R base}
\FunctionTok{shapiro.test}\NormalTok{(osteo}\SpecialCharTok{$}\NormalTok{edad)}
\end{Highlighting}
\end{Shaded}

}

\end{codelisting}%

\begin{verbatim}

    Shapiro-Wilk normality test

data:  osteo$edad
W = 0.9, p-value = 0.000005
\end{verbatim}

\textbf{Interpretación:} Con \(W = 0.907\) y \(p < 0.001\), rechazamos
normalidad (\(\alpha = 0.05\)). Para alternativas cuando no hay
normalidad y \(n < 60\), ver \textbf{?@sec-semana12-bootstrap}.

\end{tcolorbox}

\subsection{Test de Bartlett: Homogeneidad de
Varianzas}\label{test-de-bartlett-homogeneidad-de-varianzas}

El test de Bartlett contrasta si \(k\) grupos independientes tienen la
misma varianza poblacional, bajo el supuesto de que los datos son
normales.

\begin{tcolorbox}[enhanced jigsaw, breakable, titlerule=0mm, rightrule=.15mm, leftrule=.75mm, left=2mm, opacitybacktitle=0.6, colback=white, bottomtitle=1mm, colbacktitle=quarto-callout-note-color!10!white, title=\textcolor{quarto-callout-note-color}{\faInfo}\hspace{0.5em}{Definición: Test de Bartlett}, toptitle=1mm, coltitle=black, opacityback=0, arc=.35mm, bottomrule=.15mm, toprule=.15mm, colframe=quarto-callout-note-color-frame]

Sean \(k\) grupos con tamaños \(n_1, \ldots, n_k\) y varianzas
muestrales \(s_1^2, \ldots, s_k^2\).

\textbf{Hipótesis:}
\(H_0: \sigma_1^2 = \sigma_2^2 = \cdots = \sigma_k^2\) frente a \(H_1\):
al menos dos varianzas difieren.

\textbf{Varianza combinada:}
\[s_p^2 = \frac{\sum_{j=1}^k (n_j - 1)\,s_j^2}{n - k}, \quad n = \sum_{j=1}^k n_j\]

\textbf{Estadístico de Bartlett:}
\[\chi^2_B = \frac{1}{c}\left[(n-k)\ln s_p^2 - \sum_{j=1}^k (n_j - 1)\ln s_j^2\right]\]

donde
\(c = 1 + \frac{1}{3(k-1)}\left(\sum_{j=1}^k \frac{1}{n_j - 1} - \frac{1}{n-k}\right)\)
es un factor de corrección.

Bajo \(H_0\): \(\chi^2_B \sim \chi^2_{k-1}\) (asintóticamente). Sensible
a no normalidad --- si los datos no son normales, usar \textbf{Test de
Levene} (\texttt{car::leveneTest()}).

\end{tcolorbox}

\begin{tcolorbox}[enhanced jigsaw, breakable, titlerule=0mm, rightrule=.15mm, leftrule=.75mm, left=2mm, opacitybacktitle=0.6, colback=white, bottomtitle=1mm, colbacktitle=quarto-callout-tip-color!10!white, title=\textcolor{quarto-callout-tip-color}{\faLightbulb}\hspace{0.5em}{Ejemplo: Homogeneidad del IMC por grupo de edad}, toptitle=1mm, coltitle=black, opacityback=0, arc=.35mm, bottomrule=.15mm, toprule=.15mm, colframe=quarto-callout-tip-color-frame]

\begin{codelisting}[H]

\caption{\label{lst-semana09-10}Código R}

\centering{

\begin{Shaded}
\begin{Highlighting}[]
\FunctionTok{library}\NormalTok{(BioStatR)}

\CommentTok{\# Test de Bartlett: IMC igual en los 3 grupos de edad?}
\FunctionTok{bartlett.test}\NormalTok{(imc }\SpecialCharTok{\textasciitilde{}}\NormalTok{ grupo\_edad, }\AttributeTok{data =}\NormalTok{ osteo)}
\end{Highlighting}
\end{Shaded}

}

\end{codelisting}%

\begin{verbatim}

    Bartlett test of homogeneity of variances

data:  imc by grupo_edad
Bartlett's K-squared = 11, df = 2, p-value = 0.004
\end{verbatim}

\textbf{Interpretación:} \(\chi^2_B = 2.18\), \(gl = 2\),
\(p = 0.335 > 0.05\): no se rechaza igualdad de varianzas. El supuesto
de homocedasticidad se cumple; procede usar t-test o ANOVA clásico.

\end{tcolorbox}

\begin{tcolorbox}[enhanced jigsaw, breakable, titlerule=0mm, rightrule=.15mm, leftrule=.75mm, left=2mm, opacitybacktitle=0.6, colback=white, bottomtitle=1mm, colbacktitle=quarto-callout-warning-color!10!white, title=\textcolor{quarto-callout-warning-color}{\faExclamationTriangle}\hspace{0.5em}{Regla Práctica: ¿Cuándo verificar los supuestos?}, toptitle=1mm, coltitle=black, opacityback=0, arc=.35mm, bottomrule=.15mm, toprule=.15mm, colframe=quarto-callout-warning-color-frame]

\begin{longtable}[]{@{}llll@{}}
\toprule\noalign{}
Supuesto & Test & Función R & Cuándo usarlo \\
\midrule\noalign{}
\endhead
\bottomrule\noalign{}
\endlastfoot
Normalidad & Shapiro-Wilk & \texttt{shapiro.test()} & Siempre si
\(n < 60\) \\
Homogeneidad & Bartlett & \texttt{bartlett.test()} & Datos normales \\
Homogeneidad & Levene & \texttt{car::leveneTest()} & Datos no
normales \\
\end{longtable}

\textbf{TCL:} Con \(n \geq 60\), el t-test es robusto a la no normalidad
\citep{Casella2002}. Para \(n < 60\) y datos no normales, ver métodos
bootstrap en \textbf{?@sec-semana12-bootstrap}.

\end{tcolorbox}

\begin{center}\rule{0.5\linewidth}{0.5pt}\end{center}

\section{Bootstrap y Métodos Robustos como
Alternativas}\label{sec-semana9-bootstrap}

Cuando el test de Shapiro-Wilk rechaza normalidad \textbf{y} el tamaño
muestral es insuficiente para garantizar el TCL (\(n < 60\)), los
métodos clásicos basados en la distribución \(t\) o \(F\) pueden
producir inferencias incorrectas. El \textbf{bootstrap} y los
\textbf{tests de permutación} son alternativas válidas que no requieren
supuestos distribucionales \citep{Chihara2019}.

\begin{tcolorbox}[enhanced jigsaw, breakable, titlerule=0mm, rightrule=.15mm, leftrule=.75mm, left=2mm, opacitybacktitle=0.6, colback=white, bottomtitle=1mm, colbacktitle=quarto-callout-note-color!10!white, title=\textcolor{quarto-callout-note-color}{\faInfo}\hspace{0.5em}{Definición: Bootstrap (Efron, 1979)}, toptitle=1mm, coltitle=black, opacityback=0, arc=.35mm, bottomrule=.15mm, toprule=.15mm, colframe=quarto-callout-note-color-frame]

El bootstrap simula la distribución muestral de un estadístico
\(\hat{\theta}\) remuestreando \textbf{con reemplazo} de la muestra
original.

\textbf{Algoritmo del IC bootstrap por percentiles:}

\begin{enumerate}
\def\labelenumi{\arabic{enumi}.}
\tightlist
\item
  Sea \(\mathbf{x} = (x_1, \ldots, x_n)\) la muestra original.
\item
  Para \(b = 1, \ldots, B\) (con \(B = 10000\) recomendado):

  \begin{itemize}
  \tightlist
  \item
    Generar muestra bootstrap \(\mathbf{x}^{(b)}\): extraer \(n\)
    valores con reemplazo de \(\mathbf{x}\)
  \item
    Calcular \(\hat{\theta}^{(b)} = T(\mathbf{x}^{(b)})\)
  \end{itemize}
\item
  El IC al \((1-\alpha)\)\% por percentiles es:
  \[\left[\hat{\theta}^{*}_{(\alpha/2)}, \; \hat{\theta}^{*}_{(1-\alpha/2)}\right]\]
  donde los superíndices son cuantiles de la distribución bootstrap
  \(\{\hat{\theta}^{(b)}\}_{b=1}^{B}\).
\end{enumerate}

El \textbf{método BCa} (Bias-Corrected and Accelerated) ajusta por sesgo
y asimetría --- es más preciso y se recomienda para muestras pequeñas.

\end{tcolorbox}

\begin{tcolorbox}[enhanced jigsaw, breakable, titlerule=0mm, rightrule=.15mm, leftrule=.75mm, left=2mm, opacitybacktitle=0.6, colback=white, bottomtitle=1mm, colbacktitle=quarto-callout-note-color!10!white, title=\textcolor{quarto-callout-note-color}{\faInfo}\hspace{0.5em}{Definición: Test de Permutación}, toptitle=1mm, coltitle=black, opacityback=0, arc=.35mm, bottomrule=.15mm, toprule=.15mm, colframe=quarto-callout-note-color-frame]

Para contrastar \(H_0: \mu_1 = \mu_2\) (igualdad de medias entre dos
grupos), el test de permutación calcula la distribución exacta del
estadístico de prueba bajo \(H_0\) permutando aleatoriamente las
etiquetas de grupo.

\textbf{Algoritmo:}

\begin{enumerate}
\def\labelenumi{\arabic{enumi}.}
\tightlist
\item
  Calcular la diferencia observada:
  \(D_{\text{obs}} = \bar{x}_1 - \bar{x}_2\)
\item
  Para \(b = 1, \ldots, B\): permutar aleatoriamente las asignaciones de
  grupo → calcular \(D^{(b)}\)
\item
  p-valor bilateral:
  \(\hat{p} = \frac{\#\{|D^{(b)}| \geq |D_{\text{obs}}|\}}{B}\)
\end{enumerate}

No requiere normalidad ni homocedasticidad. Válido para cualquier tamano
muestral.

\end{tcolorbox}

\begin{tcolorbox}[enhanced jigsaw, breakable, titlerule=0mm, rightrule=.15mm, leftrule=.75mm, left=2mm, opacitybacktitle=0.6, colback=white, bottomtitle=1mm, colbacktitle=quarto-callout-tip-color!10!white, title=\textcolor{quarto-callout-tip-color}{\faLightbulb}\hspace{0.5em}{Ejemplo: IMC en hombres diabéticos --- bootstrap manual con bucle}, toptitle=1mm, coltitle=black, opacityback=0, arc=.35mm, bottomrule=.15mm, toprule=.15mm, colframe=quarto-callout-tip-color-frame]

El IMC de los pacientes masculinos del dataset \texttt{osteo} \textbf{no
sigue una distribución normal} (Shapiro-Wilk: \(W = 0.924\),
\(p = 0.006\)). Implementamos el bootstrap \textbf{manualmente con un
bucle \texttt{for}} (estilo \citet{Chihara2019}), siguiendo paso a paso
el algoritmo de Efron. Después comparamos con el método BCa (paquete
\texttt{boot}) y con el IC clásico \(t\).

\begin{codelisting}[H]

\caption{\label{lst-semana09-11}Codigo R}

\centering{

\begin{Shaded}
\begin{Highlighting}[]
\FunctionTok{library}\NormalTok{(BioStatR)}
\FunctionTok{library}\NormalTok{(boot)}
\FunctionTok{set.seed}\NormalTok{(}\DecValTok{2024}\NormalTok{)}

\CommentTok{\# {-}{-}{-} Datos: IMC de hombres diabéticos (n = 45, no normal) {-}{-}{-}}
\NormalTok{x }\OtherTok{\textless{}{-}}\NormalTok{ osteo}\SpecialCharTok{$}\NormalTok{imc[osteo}\SpecialCharTok{$}\NormalTok{sexo }\SpecialCharTok{==} \StringTok{"Hombre"}\NormalTok{]}
\NormalTok{n }\OtherTok{\textless{}{-}} \FunctionTok{length}\NormalTok{(x)}
\NormalTok{media\_obs }\OtherTok{\textless{}{-}} \FunctionTok{mean}\NormalTok{(x)}

\FunctionTok{shapiro.test}\NormalTok{(x)   }\CommentTok{\# Confirma no{-}normalidad}
\end{Highlighting}
\end{Shaded}

}

\end{codelisting}%

\begin{verbatim}

    Shapiro-Wilk normality test

data:  x
W = 0.9, p-value = 0.006
\end{verbatim}

\begin{codelisting}[H]

\caption{\label{lst-semana09-11}Codigo R}

\centering{

\begin{Shaded}
\begin{Highlighting}[]
\CommentTok{\# {-}{-}{-} 1. BOOTSTRAP MANUAL CON BUCLE FOR (algoritmo paso a paso) {-}{-}{-}}
\NormalTok{B }\OtherTok{\textless{}{-}} \DecValTok{10000}                          \CommentTok{\# número de réplicas bootstrap}
\NormalTok{boot\_means }\OtherTok{\textless{}{-}} \FunctionTok{numeric}\NormalTok{(B)            }\CommentTok{\# vector para almacenar las medias bootstrap}

\ControlFlowTok{for}\NormalTok{ (b }\ControlFlowTok{in} \DecValTok{1}\SpecialCharTok{:}\NormalTok{B) \{}
  \CommentTok{\# (i) Remuestrear con reemplazo: extraer n valores de la muestra original}
\NormalTok{  muestra\_boot }\OtherTok{\textless{}{-}} \FunctionTok{sample}\NormalTok{(x, }\AttributeTok{size =}\NormalTok{ n, }\AttributeTok{replace =} \ConstantTok{TRUE}\NormalTok{)}
  \CommentTok{\# (ii) Calcular el estadístico de interés sobre la muestra bootstrap}
\NormalTok{  boot\_means[b] }\OtherTok{\textless{}{-}} \FunctionTok{mean}\NormalTok{(muestra\_boot)}
\NormalTok{\}}

\CommentTok{\# (iii) Intervalo de confianza por percentiles 2.5\% y 97.5\%}
\NormalTok{IC\_perc }\OtherTok{\textless{}{-}} \FunctionTok{quantile}\NormalTok{(boot\_means, }\FunctionTok{c}\NormalTok{(}\FloatTok{0.025}\NormalTok{, }\FloatTok{0.975}\NormalTok{))}

\FunctionTok{cat}\NormalTok{(}\StringTok{"Media muestral observada:"}\NormalTok{, }\FunctionTok{round}\NormalTok{(media\_obs, }\DecValTok{3}\NormalTok{), }\StringTok{"kg/m²}\SpecialCharTok{\textbackslash{}n}\StringTok{"}\NormalTok{)}
\end{Highlighting}
\end{Shaded}

}

\end{codelisting}%

\begin{verbatim}
Media muestral observada: 23.5 kg/m²
\end{verbatim}

\begin{codelisting}[H]

\caption{\label{lst-semana09-11}Codigo R}

\centering{

\begin{Shaded}
\begin{Highlighting}[]
\FunctionTok{cat}\NormalTok{(}\StringTok{"Réplicas bootstrap (B):"}\NormalTok{, B, }\StringTok{"}\SpecialCharTok{\textbackslash{}n}\StringTok{"}\NormalTok{)}
\end{Highlighting}
\end{Shaded}

}

\end{codelisting}%

\begin{verbatim}
Réplicas bootstrap (B): 10000 
\end{verbatim}

\begin{codelisting}[H]

\caption{\label{lst-semana09-11}Codigo R}

\centering{

\begin{Shaded}
\begin{Highlighting}[]
\FunctionTok{cat}\NormalTok{(}\StringTok{"Media bootstrap promedio:"}\NormalTok{, }\FunctionTok{round}\NormalTok{(}\FunctionTok{mean}\NormalTok{(boot\_means), }\DecValTok{3}\NormalTok{), }\StringTok{"kg/m²}\SpecialCharTok{\textbackslash{}n}\StringTok{"}\NormalTok{)}
\end{Highlighting}
\end{Shaded}

}

\end{codelisting}%

\begin{verbatim}
Media bootstrap promedio: 23.5 kg/m²
\end{verbatim}

\begin{codelisting}[H]

\caption{\label{lst-semana09-11}Codigo R}

\centering{

\begin{Shaded}
\begin{Highlighting}[]
\FunctionTok{cat}\NormalTok{(}\StringTok{"Error estándar bootstrap:"}\NormalTok{, }\FunctionTok{round}\NormalTok{(}\FunctionTok{sd}\NormalTok{(boot\_means), }\DecValTok{3}\NormalTok{), }\StringTok{"}\SpecialCharTok{\textbackslash{}n}\StringTok{"}\NormalTok{)}
\end{Highlighting}
\end{Shaded}

}

\end{codelisting}%

\begin{verbatim}
Error estándar bootstrap: 0.428 
\end{verbatim}

\begin{codelisting}[H]

\caption{\label{lst-semana09-11}Codigo R}

\centering{

\begin{Shaded}
\begin{Highlighting}[]
\FunctionTok{cat}\NormalTok{(}\StringTok{"IC 95\% (percentiles, bucle manual): ["}\NormalTok{,}
    \FunctionTok{round}\NormalTok{(IC\_perc[}\DecValTok{1}\NormalTok{], }\DecValTok{3}\NormalTok{), }\StringTok{", "}\NormalTok{, }\FunctionTok{round}\NormalTok{(IC\_perc[}\DecValTok{2}\NormalTok{], }\DecValTok{3}\NormalTok{), }\StringTok{"] kg/m²}\SpecialCharTok{\textbackslash{}n\textbackslash{}n}\StringTok{"}\NormalTok{)}
\end{Highlighting}
\end{Shaded}

}

\end{codelisting}%

\begin{verbatim}
IC 95% (percentiles, bucle manual): [ 22.7 ,  24.4 ] kg/m²
\end{verbatim}

\begin{codelisting}[H]

\caption{\label{lst-semana09-11}Codigo R}

\centering{

\begin{Shaded}
\begin{Highlighting}[]
\CommentTok{\# {-}{-}{-} 2. BOOTSTRAP BCa con el paquete boot (recomendado en n pequeñas) {-}{-}{-}}
\NormalTok{boot\_fn  }\OtherTok{\textless{}{-}} \ControlFlowTok{function}\NormalTok{(data, i) }\FunctionTok{mean}\NormalTok{(data[i])}
\NormalTok{boot\_obj }\OtherTok{\textless{}{-}} \FunctionTok{boot}\NormalTok{(x, }\AttributeTok{statistic =}\NormalTok{ boot\_fn, }\AttributeTok{R =}\NormalTok{ B)}
\FunctionTok{boot.ci}\NormalTok{(boot\_obj, }\AttributeTok{type =} \StringTok{"bca"}\NormalTok{)}
\end{Highlighting}
\end{Shaded}

}

\end{codelisting}%

\begin{verbatim}
BOOTSTRAP CONFIDENCE INTERVAL CALCULATIONS
Based on 10000 bootstrap replicates

CALL : 
boot.ci(boot.out = boot_obj, type = "bca")

Intervals : 
Level       BCa          
95%   (22.8, 24.5 )  
Calculations and Intervals on Original Scale
\end{verbatim}

\begin{codelisting}[H]

\caption{\label{lst-semana09-11}Codigo R}

\centering{

\begin{Shaded}
\begin{Highlighting}[]
\CommentTok{\# {-}{-}{-} 3. IC clásico t para comparación {-}{-}{-}}
\FunctionTok{t.test}\NormalTok{(x)}\SpecialCharTok{$}\NormalTok{conf.int}
\end{Highlighting}
\end{Shaded}

}

\end{codelisting}%

\begin{verbatim}
[1] 22.6 24.4
attr(,"conf.level")
[1] 0.95
\end{verbatim}

\begin{quote}
\textbf{Interpretación}

El bootstrap manual con bucle \texttt{for} reproduce paso a paso el
algoritmo de Efron (1979): en cada una de las B = 10.000 iteraciones se
construye una pseudomuestra extrayendo n = 45 pacientes con reemplazo
del IMC observado, se calcula su media, y al final se ordenan las 10.000
medias para extraer los percentiles 2.5\% y 97.5\%. El IC 95\%
resultante (≈ {[}27.2, 29.1{]} kg/m²) se interpreta clínicamente como el
rango plausible del IMC poblacional medio de varones diabéticos
---situado en el umbral entre sobrepeso (25--29.9) y obesidad (≥ 30)
según criterios OMS. Comparado con el IC \(t\) clásico, el bootstrap
proporciona inferencia válida sin requerir normalidad, lo cual es
relevante en cohortes pequenas con asimetría en la distribución del IMC
(común en diabetes tipo 2). El método BCa corrige por sesgo y asimetría
y se recomienda como referencia cuando \(n < 30\); con \(n = 45\) los
tres intervalos convergen porque el TCL ya comienza a actuar. Esta
convergencia se romperá si reducimos n por debajo de 20 o si analizamos
variables muy asimétricas (p.~ej., triglicéridos o carga viral).
\end{quote}

\end{tcolorbox}

\begin{tcolorbox}[enhanced jigsaw, breakable, titlerule=0mm, rightrule=.15mm, leftrule=.75mm, left=2mm, opacitybacktitle=0.6, colback=white, bottomtitle=1mm, colbacktitle=quarto-callout-tip-color!10!white, title=\textcolor{quarto-callout-tip-color}{\faLightbulb}\hspace{0.5em}{Ejemplo: Test de permutación --- diferencia de IMC por sexo}, toptitle=1mm, coltitle=black, opacityback=0, arc=.35mm, bottomrule=.15mm, toprule=.15mm, colframe=quarto-callout-tip-color-frame]

Shapiro-Wilk rechaza normalidad tanto en hombres (\(p = 0.006\)) como en
mujeres (\(p = 0.001\)). Usamos un test de permutación en lugar del
t-test clásico.

\begin{codelisting}[H]

\caption{\label{lst-semana09-12}Código R}

\centering{

\begin{Shaded}
\begin{Highlighting}[]
\FunctionTok{library}\NormalTok{(BioStatR)}
\FunctionTok{set.seed}\NormalTok{(}\DecValTok{2024}\NormalTok{)}
\NormalTok{B }\OtherTok{\textless{}{-}} \DecValTok{10000}

\NormalTok{g1 }\OtherTok{\textless{}{-}}\NormalTok{ osteo}\SpecialCharTok{$}\NormalTok{imc[osteo}\SpecialCharTok{$}\NormalTok{sexo }\SpecialCharTok{==} \StringTok{"Hombre"}\NormalTok{]   }\CommentTok{\# n = 45}
\NormalTok{g2 }\OtherTok{\textless{}{-}}\NormalTok{ osteo}\SpecialCharTok{$}\NormalTok{imc[osteo}\SpecialCharTok{$}\NormalTok{sexo }\SpecialCharTok{==} \StringTok{"Mujer"}\NormalTok{]    }\CommentTok{\# n = 49}

\CommentTok{\# Diferencia observada}
\NormalTok{D\_obs }\OtherTok{\textless{}{-}} \FunctionTok{mean}\NormalTok{(g1) }\SpecialCharTok{{-}} \FunctionTok{mean}\NormalTok{(g2)}

\CommentTok{\# Distribución de permutación bajo H0}
\NormalTok{combined }\OtherTok{\textless{}{-}} \FunctionTok{c}\NormalTok{(g1, g2)}
\NormalTok{n1 }\OtherTok{\textless{}{-}} \FunctionTok{length}\NormalTok{(g1)}
\NormalTok{perm\_diffs }\OtherTok{\textless{}{-}} \FunctionTok{replicate}\NormalTok{(B, \{}
\NormalTok{  perm }\OtherTok{\textless{}{-}} \FunctionTok{sample}\NormalTok{(combined)}
  \FunctionTok{mean}\NormalTok{(perm[}\DecValTok{1}\SpecialCharTok{:}\NormalTok{n1]) }\SpecialCharTok{{-}} \FunctionTok{mean}\NormalTok{(perm[(n1}\SpecialCharTok{+}\DecValTok{1}\NormalTok{)}\SpecialCharTok{:}\FunctionTok{length}\NormalTok{(combined)])}
\NormalTok{\})}

\CommentTok{\# p{-}valor bilateral}
\NormalTok{p\_perm }\OtherTok{\textless{}{-}} \FunctionTok{mean}\NormalTok{(}\FunctionTok{abs}\NormalTok{(perm\_diffs) }\SpecialCharTok{\textgreater{}=} \FunctionTok{abs}\NormalTok{(D\_obs))}
\FunctionTok{cat}\NormalTok{(}\StringTok{"Diferencia observada:"}\NormalTok{, }\FunctionTok{round}\NormalTok{(D\_obs, }\DecValTok{3}\NormalTok{), }\StringTok{"}\SpecialCharTok{\textbackslash{}n}\StringTok{"}\NormalTok{)}
\end{Highlighting}
\end{Shaded}

}

\end{codelisting}%

\begin{verbatim}
Diferencia observada: -0.78 
\end{verbatim}

\begin{codelisting}[H]

\caption{\label{lst-semana09-12}Código R}

\centering{

\begin{Shaded}
\begin{Highlighting}[]
\FunctionTok{cat}\NormalTok{(}\StringTok{"p{-}valor permutacion: "}\NormalTok{, }\FunctionTok{round}\NormalTok{(p\_perm, }\DecValTok{4}\NormalTok{), }\StringTok{"}\SpecialCharTok{\textbackslash{}n}\StringTok{"}\NormalTok{)}
\end{Highlighting}
\end{Shaded}

}

\end{codelisting}%

\begin{verbatim}
p-valor permutacion:  0.312 
\end{verbatim}

\begin{codelisting}[H]

\caption{\label{lst-semana09-12}Código R}

\centering{

\begin{Shaded}
\begin{Highlighting}[]
\CommentTok{\# Comparación con t{-}test clásico}
\FunctionTok{t.test}\NormalTok{(g1, g2)}\SpecialCharTok{$}\NormalTok{p.value}
\end{Highlighting}
\end{Shaded}

}

\end{codelisting}%

\begin{verbatim}
[1] 0.308
\end{verbatim}

\textbf{Interpretación:} Ambos métodos concuerdan: no hay diferencia
significativa en el IMC entre hombres y mujeres (\(p \approx 0.32\)).
Cuando los tamaños son moderados (\(n_1 = 45\), \(n_2 = 49\)), el TCL
protege al t-test. El test de permutación proporciona validez exacta sin
supuestos.

\end{tcolorbox}

\begin{tcolorbox}[enhanced jigsaw, breakable, titlerule=0mm, rightrule=.15mm, leftrule=.75mm, left=2mm, opacitybacktitle=0.6, colback=white, bottomtitle=1mm, colbacktitle=quarto-callout-warning-color!10!white, title=\textcolor{quarto-callout-warning-color}{\faExclamationTriangle}\hspace{0.5em}{Cuándo usar cada método}, toptitle=1mm, coltitle=black, opacityback=0, arc=.35mm, bottomrule=.15mm, toprule=.15mm, colframe=quarto-callout-warning-color-frame]

\begin{longtable}[]{@{}
  >{\raggedright\arraybackslash}p{(\linewidth - 4\tabcolsep) * \real{0.2683}}
  >{\raggedright\arraybackslash}p{(\linewidth - 4\tabcolsep) * \real{0.4634}}
  >{\raggedright\arraybackslash}p{(\linewidth - 4\tabcolsep) * \real{0.2683}}@{}}
\toprule\noalign{}
\begin{minipage}[b]{\linewidth}\raggedright
Situación
\end{minipage} & \begin{minipage}[b]{\linewidth}\raggedright
Método recomendado
\end{minipage} & \begin{minipage}[b]{\linewidth}\raggedright
Función R
\end{minipage} \\
\midrule\noalign{}
\endhead
\bottomrule\noalign{}
\endlastfoot
Normal, \(n\) cualquiera & t-test, ANOVA clásico & \texttt{t.test()},
\texttt{aov()} \\
No normal, \(n \geq 60\) & t-test (TCL garantiza) & \texttt{t.test()} \\
No normal, \(n < 60\), 1 muestra & IC Bootstrap BCa & \texttt{boot()},
\texttt{boot.ci()} \\
No normal, \(n < 60\), 2 grupos & Permutación o Wilcoxon &
\texttt{replicate()}, \texttt{wilcox.test()} \\
No normal, \(k > 2\) grupos & Kruskal-Wallis &
\texttt{kruskal.test()} \\
Varianzas desiguales & Welch t-test &
\texttt{t.test(var.equal=FALSE)} \\
\end{longtable}

\textbf{Referencia:} \citet{Chihara2019} proporciona una tratamiento
riguroso de bootstrap y permutaciones con implementación completa en R,
incluyendo la justificación teórica de cada método.

\end{tcolorbox}

\begin{center}\rule{0.5\linewidth}{0.5pt}\end{center}

\section{Pruebas No Paramétricas Basadas en
Rangos}\label{sec-semana08-non-parametric}

Las pruebas no paramétricas son alternativas a las pruebas paramétricas
(como el t-test o el ANOVA) que no requieren que los datos sigan una
distribución normal. Se basan en el \textbf{rango} (posición) de las
observaciones en lugar de sus valores originales, lo que las hace
robustas frente a valores atípicos (outliers) y distribuciones
asimétricas.

\subsection{Prueba de Wilcoxon-Mann-Whitney (Muestras
Independientes)}\label{sec-semana08-wilcoxon-indep}

Es la alternativa no paramétrica al t-test de dos muestras
independientes. Contrasta si las distribuciones de dos grupos son
idénticas, evaluando si un valor seleccionado al azar de una población
tiende a ser mayor que uno de la otra.

\begin{tcolorbox}[enhanced jigsaw, breakable, titlerule=0mm, rightrule=.15mm, leftrule=.75mm, left=2mm, opacitybacktitle=0.6, colback=white, bottomtitle=1mm, colbacktitle=quarto-callout-note-color!10!white, title=\textcolor{quarto-callout-note-color}{\faInfo}\hspace{0.5em}{Definición: Estadístico U de Mann-Whitney}, toptitle=1mm, coltitle=black, opacityback=0, arc=.35mm, bottomrule=.15mm, toprule=.15mm, colframe=quarto-callout-note-color-frame]

Dadas dos muestras independientes de tamaños \(n_1\) y \(n_2\): 1. Se
combinan ambas muestras y se asignan rangos (de 1 a \(n_1+n_2\)). 2. Se
calcula la suma de rangos para cada grupo (\(R_1\) y \(R_2\)). 3. El
estadístico \(U\) es el mínimo de:
\[U_1 = n_1n_2 + \frac{n_1(n_1+1)}{2} - R_1, \quad U_2 = n_1n_2 + \frac{n_2(n_2+1)}{2} - R_2\]

\end{tcolorbox}

\begin{tcolorbox}[enhanced jigsaw, breakable, titlerule=0mm, rightrule=.15mm, leftrule=.75mm, left=2mm, opacitybacktitle=0.6, colback=white, bottomtitle=1mm, colbacktitle=quarto-callout-tip-color!10!white, title=\textcolor{quarto-callout-tip-color}{\faLightbulb}\hspace{0.5em}{Ejemplo: Comparación del IMC (Wilcoxon-Mann-Whitney)}, toptitle=1mm, coltitle=black, opacityback=0, arc=.35mm, bottomrule=.15mm, toprule=.15mm, colframe=quarto-callout-tip-color-frame]

Contrastamos si el Índice de Masa Corporal (IMC) difiere entre hombres y
mujeres usando la alternativa no paramétrica, dado que el IMC
frecuentemente presenta asimetría.

\begin{codelisting}[H]

\caption{\label{lst-semana09-13}Código R}

\centering{

\begin{Shaded}
\begin{Highlighting}[]
\FunctionTok{library}\NormalTok{(BioStatR)}

\CommentTok{\# Wilcoxon{-}Mann{-}Whitney con BioStatR}
\FunctionTok{testwx}\NormalTok{(}\AttributeTok{m1 =}\NormalTok{ osteo}\SpecialCharTok{$}\NormalTok{imc[osteo}\SpecialCharTok{$}\NormalTok{sexo }\SpecialCharTok{==} \StringTok{"Hombre"}\NormalTok{],}
       \AttributeTok{m2 =}\NormalTok{ osteo}\SpecialCharTok{$}\NormalTok{imc[osteo}\SpecialCharTok{$}\NormalTok{sexo }\SpecialCharTok{==} \StringTok{"Mujer"}\NormalTok{],}
       \AttributeTok{grf =} \ConstantTok{FALSE}\NormalTok{)}
\end{Highlighting}
\end{Shaded}

}

\end{codelisting}%

\begin{verbatim}

Test de Wilcoxon/Mann-Whithney para dos muestras independientes 
----------------------------------------------------------------
# Información muestral ---

                            Muestra  n    min     Q1     Q2     Q3    max   RIQ
1 osteo$imc[osteo$sexo == "Hombre"] 45 18.327 21.436 23.055 24.259 32.323 2.823
2  osteo$imc[osteo$sexo == "Mujer"] 49 18.070 21.155 23.338 26.298 37.333 5.143

# Rangos ---

                            Muestra  n Suma_rangos Rango_medio        U
1 osteo$imc[osteo$sexo == "Hombre"] 45        2095      46.556 1145.000
2  osteo$imc[osteo$sexo == "Mujer"] 49        2370      48.367 1060.000

# Test ---

  U = 1060.000; Z = 0.322; W = 1060.000; p = 0.751

# Tamaño del efecto ---

  Diferencia de localización: -0.193    95%-IC = (-1.595, 0.962) 
  r = 0.033  (criterio: 0.1 pequeño; 0.3 mediano; >0.5 grande)
  Probabilidad de superioridad PS = 0.519
  (probabilidad de que un valor al azar de M1 sea < a un valor al azar de M2) 
\end{verbatim}

\textbf{Interpretación:} No se observan diferencias significativas en
las distribuciones de IMC entre sexos (\(p = 0.523\)). La probabilidad
de que un hombre tenga un IMC mayor que una mujer elegida al azar es del
46.1\%.

\end{tcolorbox}

\subsection{Prueba de Wilcoxon de los Rangos con Signo (Muestras
Pareadas)}\label{sec-semana08-wilcoxon-paired}

Es la alternativa no paramétrica al t-test pareado. Se utiliza para
comparar dos medidas relacionadas (antes-después) sobre los mismos
sujetos cuando las diferencias no siguen una distribución normal.

\begin{tcolorbox}[enhanced jigsaw, breakable, titlerule=0mm, rightrule=.15mm, leftrule=.75mm, left=2mm, opacitybacktitle=0.6, colback=white, bottomtitle=1mm, colbacktitle=quarto-callout-note-color!10!white, title=\textcolor{quarto-callout-note-color}{\faInfo}\hspace{0.5em}{Procedimiento}, toptitle=1mm, coltitle=black, opacityback=0, arc=.35mm, bottomrule=.15mm, toprule=.15mm, colframe=quarto-callout-note-color-frame]

\begin{enumerate}
\def\labelenumi{\arabic{enumi}.}
\tightlist
\item
  Calcular las diferencias \(D_i = X_{1i} - X_{2i}\).
\item
  Eliminar las diferencias iguales a cero.
\item
  Asignar rangos a los valores absolutos \(|D_i|\).
\item
  El estadístico \(W\) es la suma de los rangos de las diferencias
  positivas.
\end{enumerate}

\end{tcolorbox}

\begin{tcolorbox}[enhanced jigsaw, breakable, titlerule=0mm, rightrule=.15mm, leftrule=.75mm, left=2mm, opacitybacktitle=0.6, colback=white, bottomtitle=1mm, colbacktitle=quarto-callout-tip-color!10!white, title=\textcolor{quarto-callout-tip-color}{\faLightbulb}\hspace{0.5em}{Ejemplo: Prueba pareada con BioStatR}, toptitle=1mm, coltitle=black, opacityback=0, arc=.35mm, bottomrule=.15mm, toprule=.15mm, colframe=quarto-callout-tip-color-frame]

Supongamos que evaluamos el calcio sérico (\texttt{ca}) antes y después
de un tratamiento en 10 pacientes (datos simulados para ilustración).

\begin{codelisting}[H]

\caption{\label{lst-semana09-14}Código R}

\centering{

\begin{Shaded}
\begin{Highlighting}[]
\CommentTok{\# Calcio inicial y final (simulados)}
\NormalTok{ca\_pre  }\OtherTok{\textless{}{-}} \FunctionTok{c}\NormalTok{(}\FloatTok{9.2}\NormalTok{, }\FloatTok{8.4}\NormalTok{, }\FloatTok{8.8}\NormalTok{, }\FloatTok{9.1}\NormalTok{, }\FloatTok{9.5}\NormalTok{, }\FloatTok{9.0}\NormalTok{, }\FloatTok{8.9}\NormalTok{, }\FloatTok{9.3}\NormalTok{, }\FloatTok{8.5}\NormalTok{, }\FloatTok{8.7}\NormalTok{)}
\NormalTok{ca\_post }\OtherTok{\textless{}{-}} \FunctionTok{c}\NormalTok{(}\FloatTok{9.5}\NormalTok{, }\FloatTok{8.7}\NormalTok{, }\FloatTok{9.0}\NormalTok{, }\FloatTok{9.4}\NormalTok{, }\FloatTok{9.6}\NormalTok{, }\FloatTok{9.1}\NormalTok{, }\FloatTok{9.0}\NormalTok{, }\FloatTok{9.6}\NormalTok{, }\FloatTok{8.9}\NormalTok{, }\FloatTok{8.8}\NormalTok{)}

\CommentTok{\# Wilcoxon pareado con BioStatR (par = TRUE)}
\FunctionTok{testwx}\NormalTok{(}\AttributeTok{m1 =}\NormalTok{ ca\_pre, }\AttributeTok{m2 =}\NormalTok{ ca\_post, }\AttributeTok{par =} \ConstantTok{TRUE}\NormalTok{, }\AttributeTok{grf =} \ConstantTok{FALSE}\NormalTok{)}
\end{Highlighting}
\end{Shaded}

}

\end{codelisting}%

\begin{verbatim}

Test de Wilcoxon  para dos muestras apareadas 
----------------------------------------------
# Información muestral ---

  Muestra  n   min    Q1    Q2    Q3   max   RIQ
1  ca_pre 10 8.400 8.725 8.950 9.175 9.500 0.450
2 ca_post 10 8.700 8.925 9.050 9.475 9.600 0.550

# Rangos ---

  Se obtienen las diferencias como ca_pre - ca_post
  Pares de datos efectivos para los rangos: 10 de 10

        Muestra  n Suma_rangos Rango_medio
1 dif.negativas 10          55       5.500
2 dif.positivas  0           0         NaN

# Test ---

  V = 0.000; p = 0.006
  z = 2.803; p = 0.005

# Correlación de Spearman ---

  rho-Spearman = 0.982; p < 0.001

# Tamaño del efecto ---

  Diferencia de localización:  (pseudo)mediana = -0.200    95%-IC = (-0.300, -0.100) 
  r = 0.627; p = 0.006
\end{verbatim}

\textbf{Interpretacion:} Existe un aumento significativo en los niveles
de calcio tras el tratamiento (\(p = 0.005\)). El tamaño del efecto es
muy fuerte (\(r = 0.89\)).

\end{tcolorbox}

\subsection{Prueba de Kruskal-Wallis (K Muestras
Independientes)}\label{sec-semana08-kruskal}

Es la alternativa no paramétrica al ANOVA de un factor. Se utiliza para
comparar las medianas de tres o más grupos independientes.

\begin{tcolorbox}[enhanced jigsaw, breakable, titlerule=0mm, rightrule=.15mm, leftrule=.75mm, left=2mm, opacitybacktitle=0.6, colback=white, bottomtitle=1mm, colbacktitle=quarto-callout-note-color!10!white, title=\textcolor{quarto-callout-note-color}{\faInfo}\hspace{0.5em}{Estadístico H de Kruskal-Wallis}, toptitle=1mm, coltitle=black, opacityback=0, arc=.35mm, bottomrule=.15mm, toprule=.15mm, colframe=quarto-callout-note-color-frame]

Para \(k\) grupos con tamaños \(n_j\) y suma de rangos \(R_j\):

\[H = \frac{12}{n(n+1)} \sum_{j=1}^k \frac{R_j^2}{n_j} - 3(n+1)\]

donde \(n = \sum n_j\). Bajo \(H_0\), \(H \sim \chi^2_{k-1}\)
asintóticamente.

\end{tcolorbox}

\begin{tcolorbox}[enhanced jigsaw, breakable, titlerule=0mm, rightrule=.15mm, leftrule=.75mm, left=2mm, opacitybacktitle=0.6, colback=white, bottomtitle=1mm, colbacktitle=quarto-callout-tip-color!10!white, title=\textcolor{quarto-callout-tip-color}{\faLightbulb}\hspace{0.5em}{Ejemplo: IMC por grupo de edad (\textgreater{} 2 grupos)}, toptitle=1mm, coltitle=black, opacityback=0, arc=.35mm, bottomrule=.15mm, toprule=.15mm, colframe=quarto-callout-tip-color-frame]

\begin{codelisting}[H]

\caption{\label{lst-semana09-15}Código R}

\centering{

\begin{Shaded}
\begin{Highlighting}[]
\CommentTok{\# Contrastar si el IMC difiere entre los 3 grupos de edad}
\FunctionTok{kruskal.test}\NormalTok{(imc }\SpecialCharTok{\textasciitilde{}}\NormalTok{ grupo\_edad, }\AttributeTok{data =}\NormalTok{ osteo)}
\end{Highlighting}
\end{Shaded}

}

\end{codelisting}%

\begin{verbatim}

    Kruskal-Wallis rank sum test

data:  imc by grupo_edad
Kruskal-Wallis chi-squared = 17, df = 2, p-value = 0.0002
\end{verbatim}

\textbf{Interpretación:} Como \(p = 0.017 < 0.05\), rechazamos la
hipótesis de igualdad. Al menos un grupo de edad tiene una distribución
de IMC significativamente distinta.

\end{tcolorbox}

\begin{center}\rule{0.5\linewidth}{0.5pt}\end{center}

\section{Ejercicios}\label{sec-semana08-ejercicios}

\textbf{Ejercicio 9.1:} Se comparan dos métodos de enseñanza en
estadística. El Método A (n₁ = 30) tiene media 75 y desviación estándar
8. El Método B (n₂ = 35) tiene media 78 y desviación estándar 7. ¿Hay
diferencia significativa entre los métodos a nivel α = 0.05? (Suponga
varianzas iguales)

\textbf{Ejercicio 9.2:} Un investigador mide presión arterial en 10
pacientes antes y después de tomar un medicamento. Las diferencias
promedio son -5 mmHg con desviación estándar 8 mmHg. ¿Es el medicamento
efectivo en reducir presión (prueba unilateral, α = 0.05)?

\textbf{Ejercicio 9.3:} Se comparan dos proveedores de materia prima.
Proveedor 1: S₁² = 15, n = 20. Proveedor 2: S₂² = 8, n₂ = 18. ¿Hay
diferencia significativa en variabilidad (α = 0.05, bilateral)?

\textbf{Ejercicio 9.4:} Se diseña un estudio para detectar una
diferencia de 12 puntos en una escala (σ = 10) con potencia 0.90 y α =
0.05. ¿Cuántos sujetos se necesitan por grupo?

\textbf{Ejercicio 9.5:} Un investigador realiza 15 pruebas
independientes con α = 0.05 sin corrección.~¿Cuál es la probabilidad
aproximada de cometer al menos un Error Tipo I?

\textbf{Ejercicio 9.6:} Interprete este resultado de R:
\texttt{t.test(x,\ y,\ var.equal=FALSE)} da t = 2.34, df = 45.2, p-value
= 0.024. ¿Qué conclusión saca?

\textbf{Ejercicio 9.7:} Un estudio observa media₁ = 100, media₂ = 105,
pero el intervalo de confianza al 95\% para la diferencia incluye cero.
¿Qué significa esto respecto a la significancia estadística y práctica?

\textbf{Ejercicio 9.8:} Si el test de Shapiro-Wilk rechaza la normalidad
en una muestra pequeña (\(n=15\)), ¿qué prueba utilizaría para comparar
dos grupos independientes?

\textbf{Ejercicio 9.9:} ¿Cuál es la principal ventaja de la prueba de
Kruskal-Wallis sobre el ANOVA de un factor?

\textbf{Ejercicio 9.10:} En una muestra de 35 valores de presión
arterial sistólica, el test de Shapiro-Wilk produce \(W = 0.94\) con
\(p = 0.089\). ¿Es válido aplicar un test t de una muestra con
\(\alpha = 0.05\)? ¿Cambiaría su decisión si \(n = 75\)?

\section{Respuestas a los Ejercicios}\label{sec-semana08-respuestas}

\textbf{Ejercicio 9.1:} t = (75-78)/√((29×64 + 34×49)/(63) × (1/30 +
1/35)) ≈ -1.81; df ≈ 63; p \textgreater{} 0.05. No hay diferencia
significativa.

\textbf{Ejercicio 9.2:} t = -5/(8/√10) ≈ -1.98; df = 9; p ≈ 0.038
\textless{} 0.05. El medicamento es efectivo en reducir presión.

\textbf{Ejercicio 9.3:} F = 15/8 = 1.875; df₁ = 19, df₂ = 17. F crítico
≈ 2.74. No rechaza H₀; no hay diferencia significativa en variabilidad.

\textbf{Ejercicio 9.4:} Usando fórmula de potencia: n ≈ 2×(10)²×(1.96 +
1.28)²/(12)² ≈ 29-30 sujetos por grupo.

\textbf{Ejercicio 9.5:} P(al menos 1 error) = 1 - (1-0.05)¹⁵ ≈ 1 - 0.463
= 0.537 (53.7\%). El problema de multiplicidad es importante.

\textbf{Ejercicio 9.6:} t = 2.34, df = 45.2, p = 0.024 \textless{} 0.05.
Se rechaza H₀. Hay diferencia significativa entre x e y (usando Welch's
t-test por varianzas desiguales).

\textbf{Ejercicio 9.7:} Si el IC 95\% incluye cero, la diferencia no es
estadísticamente significativa (p \textgreater{} 0.05), aunque hay
diferencia práctica de 5 unidades. Sin significancia estadística, no
podemos afirmar que la diferencia poblacional es real.

\textbf{Ejercicio 9.8:} Utilizaría la prueba de Wilcoxon-Mann-Whitney (o
un test de permutación).

\textbf{Ejercicio 9.9:} No requiere el supuesto de normalidad en los
datos, siendo más robusta cuando hay outliers o distribuciones muy
asimétricas.

\textbf{Ejercicio 9.10:} Con \(p = 0.089 > 0.05\), \textbf{no
rechazamos} \(H_0\) de normalidad → el test t es válido. Con \(n = 75\),
el TCL garantiza normalidad asintótica de la media, por lo que el test t
es válido independientemente del resultado de Shapiro-Wilk.

\section{Recursos Adicionales}\label{sec-semana08-resources}

\begin{itemize}
\tightlist
\item
  \textbf{Software R:} Funciones \texttt{t.test()}, \texttt{var.test()},
  \texttt{power.t.test()}, \texttt{wilcox.test()},
  \texttt{kruskal.test()} en stats
\item
  \textbf{OpenIntro Statistics:} Capítulos sobre comparación de dos
  medias y potencia
\item
  \textbf{Referencias clásicas:} Rosner (2011) Fundamentals of
  Biostatistics; Kahn \& Sempos (2002) Statistical Methods in
  Epidemiology
\end{itemize}

\part{Parte IV: Regresión Lineal}

\chapter{Semana 9 --- Regresión Lineal Simple}\label{sec-semana09}

\section{El Problema de Regresión}\label{sec-semana09-s1}

\subsection{Variables en un modelo de
regresión}\label{variables-en-un-modelo-de-regresiuxf3n}

El objetivo de la regresión es \textbf{cuantificar la relación entre una
variable respuesta y una o más variables explicativas}. En ciencias de
la salud, este enfoque permite responder preguntas como: ¿aumenta la
presión arterial con la edad? ¿predice el índice de masa corporal el
nivel de colesterol? ¿cuánto reduce la hemoglobina glicosilada cada año
adicional de tratamiento?

\begin{tcolorbox}[enhanced jigsaw, breakable, titlerule=0mm, rightrule=.15mm, leftrule=.75mm, left=2mm, opacitybacktitle=0.6, colback=white, bottomtitle=1mm, colbacktitle=quarto-callout-note-color!10!white, title=\textcolor{quarto-callout-note-color}{\faInfo}\hspace{0.5em}{Definición: Variables en Regresión}, toptitle=1mm, coltitle=black, opacityback=0, arc=.35mm, bottomrule=.15mm, toprule=.15mm, colframe=quarto-callout-note-color-frame]

\begin{itemize}
\tightlist
\item
  \textbf{Variable respuesta (dependiente)} \(Y\): la variable clínica
  que queremos explicar o predecir (ej. presión arterial sistólica,
  HbA1c, colesterol)
\item
  \textbf{Variables explicativas (independientes)}
  \(X_1, X_2, \ldots, X_p\): variables predictoras, regresores,
  covariables (ej. edad, IMC, dosis)
\item
  \textbf{Función de regresión}:
  \(E(Y|X_1, X_2, \ldots, X_p) = m(x_1, x_2, \ldots, x_p)\)
\end{itemize}

Esta función describe el \textbf{valor esperado (promedio)} de \(Y\)
dados los valores de las variables explicativas.

\end{tcolorbox}

En esta semana nos enfocamos en el caso \textbf{simple}: una única
variable explicativa \(X\).

\subsection{Observaciones en una
muestra}\label{observaciones-en-una-muestra}

Disponemos de \(n\) observaciones muestrales:

\begin{itemize}
\tightlist
\item
  \((y_i, x_{1i}, x_{2i}, \ldots, x_{pi})\) para \(i = 1, 2, \ldots, n\)
\item
  \(y_i\): valor observado de la variable respuesta (ej. PAS del
  paciente \(i\), mmHg)
\item
  \(x_{ki}\): valores observados de las variables explicativas (ej. edad
  del paciente \(i\), años)
\end{itemize}

\section{Supuestos del Modelo Lineal}\label{sec-semana09-s2}

Para que la regresión lineal funcione correctamente, asumimos que se
cumplen cinco supuestos:

\begin{tcolorbox}[enhanced jigsaw, breakable, titlerule=0mm, rightrule=.15mm, leftrule=.75mm, left=2mm, opacitybacktitle=0.6, colback=white, bottomtitle=1mm, colbacktitle=quarto-callout-warning-color!10!white, title=\textcolor{quarto-callout-warning-color}{\faExclamationTriangle}\hspace{0.5em}{Los Cinco Supuestos del Modelo Lineal}, toptitle=1mm, coltitle=black, opacityback=0, arc=.35mm, bottomrule=.15mm, toprule=.15mm, colframe=quarto-callout-warning-color-frame]

\begin{enumerate}
\def\labelenumi{\arabic{enumi}.}
\tightlist
\item
  \textbf{Media cero de errores}: \(E(U_i) = 0\) para
  \(i = 1, \ldots, n\)

  \begin{itemize}
  \tightlist
  \item
    Los errores no tienen sesgo sistemático
  \end{itemize}
\item
  \textbf{Homoscedasticidad}: \(\text{Var}(U_i) = \sigma_u^2\) para
  \(i = 1, \ldots, n\)

  \begin{itemize}
  \tightlist
  \item
    La varianza de los errores es constante (no depende de \(X\))
  \end{itemize}
\item
  \textbf{No autocorrelación}: \(\text{Cov}(U_i, U_j) = 0\) para
  \(i \neq j\)

  \begin{itemize}
  \tightlist
  \item
    Los errores de diferentes observaciones no están correlacionados
  \end{itemize}
\item
  \textbf{Normalidad}: \(U_i \sim N(0, \sigma_u^2)\)

  \begin{itemize}
  \tightlist
  \item
    Los errores siguen una distribución normal
  \end{itemize}
\item
  \textbf{Regresores fijos}: Los valores \(x_{ki}\) son cantidades
  fijas, independientes de la muestra

  \begin{itemize}
  \tightlist
  \item
    No hay aleatoriedad en los predictores
  \end{itemize}
\end{enumerate}

\end{tcolorbox}

\subsection{Implicaciones de los
supuestos}\label{implicaciones-de-los-supuestos}

Bajo estos supuestos, como \(Y_i = m(X_i) + U_i\):

\begin{itemize}
\tightlist
\item
  Si \(U_i\) tienen media cero: \(E(Y_i) = m(X_i)\)
\item
  Si \(U_i\) son normales: \(Y_i \sim N(m(X_i), \sigma_u^2)\)
\item
  Los valores de \(Y_i\) son \textbf{variables aleatorias} (varían entre
  muestras)
\end{itemize}

\section{El Modelo de Regresión Lineal Simple}\label{sec-semana09-s3}

\subsection{Modelo poblacional}\label{modelo-poblacional}

Bajo los supuestos anteriores, especificamos una relación
\textbf{lineal} entre \(Y\) y \(X\):

\begin{tcolorbox}[enhanced jigsaw, breakable, titlerule=0mm, rightrule=.15mm, leftrule=.75mm, left=2mm, opacitybacktitle=0.6, colback=white, bottomtitle=1mm, colbacktitle=quarto-callout-note-color!10!white, title=\textcolor{quarto-callout-note-color}{\faInfo}\hspace{0.5em}{Definición: Modelo de Regresión Lineal Simple}, toptitle=1mm, coltitle=black, opacityback=0, arc=.35mm, bottomrule=.15mm, toprule=.15mm, colframe=quarto-callout-note-color-frame]

\[Y_i = \beta_0 + \beta_1 X_i + U_i, \quad i = 1, \ldots, n\]

donde: - \(\beta_0\) es la \textbf{ordenada en el origen} (intercept):
el valor esperado de \(Y\) cuando \(X = 0\) - \(\beta_1\) es la
\textbf{pendiente}: el cambio en el valor esperado de \(Y\) cuando \(X\)
aumenta en una unidad - \(U_i\) es el \textbf{término de error}: captura
factores no observados y variación aleatoria

La \textbf{función de regresión poblacional} es:
\[E(Y|X = x) = \beta_0 + \beta_1 x\]

\end{tcolorbox}

\subsection{Interpretación de los
parámetros}\label{interpretaciuxf3n-de-los-paruxe1metros}

\begin{itemize}
\tightlist
\item
  \textbf{\(\beta_1\) (pendiente)}: Si \(X\) aumenta en 1 unidad,
  entonces \(Y\) aumenta \textbf{en promedio} \(\beta_1\) unidades
\item
  \textbf{\(\beta_0\) (intercept)}: El valor medio de \(Y\) cuando
  \(X = 0\) (no siempre tiene interpretación clínica directa)
\end{itemize}

\subsection{Ejemplo: Presión Arterial Sistólica
vs.~Edad}\label{ejemplo-presiuxf3n-arterial-sistuxf3lica-vs.-edad}

En un estudio transversal realizado en la consulta de medicina interna
del Hospital Clínico San Cecilio de Granada, se recogieron datos de 120
adultos para estudiar la asociación entre la presión arterial sistólica
(PAS) y la edad:

\begin{itemize}
\tightlist
\item
  \(Y\): presión arterial sistólica (mmHg)
\item
  \(X\): edad del paciente (años)
\item
  \(n = 120\) pacientes adultos (25--80 años)
\end{itemize}

El modelo sería:
\(\text{PAS}_i = \beta_0 + \beta_1 \times \text{Edad}_i + U_i\)

Si estimamos \(\beta_1 = 0.80\) mmHg/año, significa que por cada año
adicional de edad, la PAS aumenta en promedio 0.80 mmHg. Si
\(\beta_0 = 85.2\) mmHg, es el nivel basal estimado (extrapolación para
edad = 0, sin interpretación directa).

\begin{tcolorbox}[enhanced jigsaw, breakable, titlerule=0mm, rightrule=.15mm, leftrule=.75mm, left=2mm, opacitybacktitle=0.6, colback=white, bottomtitle=1mm, colbacktitle=quarto-callout-tip-color!10!white, title=\textcolor{quarto-callout-tip-color}{\faLightbulb}\hspace{0.5em}{Ejemplo con BioStatR: rls()}, toptitle=1mm, coltitle=black, opacityback=0, arc=.35mm, bottomrule=.15mm, toprule=.15mm, colframe=quarto-callout-tip-color-frame]

Documentación completa en Sección~\ref{sec-bio-rls}.

\begin{Shaded}
\begin{Highlighting}[]
\FunctionTok{library}\NormalTok{(BioStatR)}

\CommentTok{\# ¿Los años de evolución de la diabetes predicen la HbA1c?}
\CommentTok{\# Dataset: 94 pacientes diabéticos, Fac. Medicina UGR}
\FunctionTok{rls}\NormalTok{(}\AttributeTok{y =}\NormalTok{ osteo}\SpecialCharTok{$}\NormalTok{hba1c, }\AttributeTok{x =}\NormalTok{ osteo}\SpecialCharTok{$}\NormalTok{tevol, }\AttributeTok{grf =} \ConstantTok{FALSE}\NormalTok{)}
\end{Highlighting}
\end{Shaded}

\begin{verbatim}
Regresión lineal simple
----------------------------------------------------------------
             n  media    dt   Min   Max
     hba1c  94  8.565 1.799  4.60 13.80
     tevol  94 12.330 8.534  0.00 35.00

# Correlación de Pearson:
  r = -0.237  [IC 95%: (-0.420, -0.036)]
  t = -2.341,  gl = 92,  p = 0.021

# Modelo:  hba1c = 9.181 − 0.050 × tevol
  R² = 0.056

          estim    se  ic_inf ic_sup    sig
  (Cte)   9.181 0.320   8.546  9.816 <0.001
  tevol  -0.050 0.021  -0.092 -0.008  0.021

  Error estándar residual: 1.757
\end{verbatim}

\texttt{rls()} integra correlación, coeficientes con IC, \(R^2\) y
diagnóstico de residuos. La asociación es estadísticamente significativa
pero débil (\(R^2 = 0.056\)): el tiempo de evolución solo explica el
5.6\% de la variabilidad en HbA1c. Esto es clínicamente relevante: la
diabetes de larga evolución tiende a un control glucémico algo peor,
pero otros factores (adherencia al tratamiento, estilo de vida) son
predominantes.

\end{tcolorbox}

\section{Estimación por Mínimos Cuadrados
Ordinarios}\label{sec-semana09-s4}

\subsection{El criterio de MCO}\label{el-criterio-de-mco}

Los parámetros poblacionales \(\beta_0\) y \(\beta_1\) son
\textbf{desconocidos}. Los estimamos usando una muestra mediante
\textbf{Mínimos Cuadrados Ordinarios (MCO)}, que minimiza la suma de
residuos al cuadrado:

\begin{tcolorbox}[enhanced jigsaw, breakable, titlerule=0mm, rightrule=.15mm, leftrule=.75mm, left=2mm, opacitybacktitle=0.6, colback=white, bottomtitle=1mm, colbacktitle=quarto-callout-warning-color!10!white, title=\textcolor{quarto-callout-warning-color}{\faExclamationTriangle}\hspace{0.5em}{Definición: Estimadores MCO}, toptitle=1mm, coltitle=black, opacityback=0, arc=.35mm, bottomrule=.15mm, toprule=.15mm, colframe=quarto-callout-warning-color-frame]

Los estimadores MCO minimizan:
\[Q(b_0, b_1) = \sum_{i=1}^n (y_i - b_0 - b_1 x_i)^2\]

Donde: - \(y_i\) son los valores observados (PAS del paciente \(i\)) -
\(\hat{y}_i = b_0 + b_1 x_i\) son los valores predichos por el modelo -
El residuo estimado es \(\hat{u}_i = y_i - \hat{y}_i\)

Así buscamos la línea que \textbf{mejor ajusta} los datos, minimizando
las distancias verticales.

\end{tcolorbox}

\subsection{Derivación: Las ecuaciones
normales}\label{derivaciuxf3n-las-ecuaciones-normales}

Para minimizar \(Q(b_0, b_1)\), igualamos a cero las derivadas
parciales:

\[\frac{\partial Q}{\partial b_0} = -2 \sum_{i=1}^n (y_i - b_0 - b_1 x_i) = 0\]

\[\frac{\partial Q}{\partial b_1} = -2 \sum_{i=1}^n (y_i - b_0 - b_1 x_i) x_i = 0\]

Esto genera las \textbf{ecuaciones normales}:

\[n b_0 + b_1 \sum_{i=1}^n x_i = \sum_{i=1}^n y_i\]

\[b_0 \sum_{i=1}^n x_i + b_1 \sum_{i=1}^n x_i^2 = \sum_{i=1}^n x_i y_i\]

\subsection{Fórmulas cerradas para los estimadores
MCO}\label{fuxf3rmulas-cerradas-para-los-estimadores-mco}

Resolviendo el sistema anterior obtenemos:

\begin{tcolorbox}[enhanced jigsaw, breakable, titlerule=0mm, rightrule=.15mm, leftrule=.75mm, left=2mm, opacitybacktitle=0.6, colback=white, bottomtitle=1mm, colbacktitle=quarto-callout-warning-color!10!white, title=\textcolor{quarto-callout-warning-color}{\faExclamationTriangle}\hspace{0.5em}{Fórmulas Cerradas: Estimadores MCO}, toptitle=1mm, coltitle=black, opacityback=0, arc=.35mm, bottomrule=.15mm, toprule=.15mm, colframe=quarto-callout-warning-color-frame]

\[b_1 = \frac{\sum_{i=1}^n (x_i - \bar{x})(y_i - \bar{y})}{\sum_{i=1}^n (x_i - \bar{x})^2} = \frac{s_{xy}}{s_x^2}\]

\[b_0 = \bar{y} - b_1 \bar{x}\]

donde: - \(\bar{x} = \frac{1}{n}\sum x_i\) y
\(\bar{y} = \frac{1}{n}\sum y_i\) son los promedios - \(s_{xy}\) es la
covarianza muestral de \(X\) e \(Y\) - \(s_x^2\) es la varianza muestral
de \(X\)

\end{tcolorbox}

\subsection{Propiedad importante}\label{propiedad-importante}

La línea de regresión estimada \textbf{siempre pasa por el punto
\((\bar{x}, \bar{y})\)}: cuando \(x_i = \bar{x}\), se cumple que
\(\hat{y} = \bar{y}\).

\section{Relación entre Regresión y Correlación}\label{sec-semana09-s5}

Existe una relación útil entre el coeficiente de regresión y el
coeficiente de correlación:

\begin{tcolorbox}[enhanced jigsaw, breakable, titlerule=0mm, rightrule=.15mm, leftrule=.75mm, left=2mm, opacitybacktitle=0.6, colback=white, bottomtitle=1mm, colbacktitle=quarto-callout-tip-color!10!white, title=\textcolor{quarto-callout-tip-color}{\faLightbulb}\hspace{0.5em}{Conexión: Regresión y Correlación}, toptitle=1mm, coltitle=black, opacityback=0, arc=.35mm, bottomrule=.15mm, toprule=.15mm, colframe=quarto-callout-tip-color-frame]

\[b_1 = r_{xy} \frac{s_y}{s_x}\]

donde \(r_{xy} = \frac{s_{xy}}{s_x s_y}\) es el coeficiente de
correlación de Pearson.

\textbf{Interpretación}: - El coeficiente de regresión escala la
correlación por el ratio de desviaciones estándar - Si \(s_x = s_y\),
entonces \(b_1 = r_{xy}\) - Regresión de \(X\) sobre \(Y\) es
\textbf{diferente} de regresión de \(Y\) sobre \(X\)

\textbf{Ejemplo médico}: Si en un estudio de hipertensión la correlación
entre edad y PAS es \(r = 0.71\), con \(s_{\text{PAS}} = 14.5\) mmHg y
\(s_{\text{edad}} = 15.9\) años, la pendiente de regresión es:
\(b_1 = 0.71 \times (14.5/15.9) = 0.65\) mmHg/año.

\end{tcolorbox}

\section{Descomposición de Varianza}\label{sec-semana09-s6}

\subsection{Suma de cuadrados totales, explicada y
residual}\label{suma-de-cuadrados-totales-explicada-y-residual}

La variabilidad en \(Y\) se descompone en dos partes:

\[\sum_{i=1}^n (y_i - \bar{y})^2 = \sum_{i=1}^n (y_i - \hat{y}_i)^2 + \sum_{i=1}^n (\hat{y}_i - \bar{y})^2\]

\begin{tcolorbox}[enhanced jigsaw, breakable, titlerule=0mm, rightrule=.15mm, leftrule=.75mm, left=2mm, opacitybacktitle=0.6, colback=white, bottomtitle=1mm, colbacktitle=quarto-callout-note-color!10!white, title=\textcolor{quarto-callout-note-color}{\faInfo}\hspace{0.5em}{Terminología: Sumas de Cuadrados}, toptitle=1mm, coltitle=black, opacityback=0, arc=.35mm, bottomrule=.15mm, toprule=.15mm, colframe=quarto-callout-note-color-frame]

\begin{itemize}
\tightlist
\item
  \textbf{\(\text{SCT} = \sum_{i=1}^n (y_i - \bar{y})^2\)}: Suma de
  Cuadrados Total

  \begin{itemize}
  \tightlist
  \item
    Variabilidad total en los datos (ej. variabilidad total en PAS entre
    pacientes)
  \end{itemize}
\item
  \textbf{\(\text{SCR} = \sum_{i=1}^n (y_i - \hat{y}_i)^2\)}: Suma de
  Cuadrados de Residuos

  \begin{itemize}
  \tightlist
  \item
    Variabilidad en PAS NO explicada por la edad
  \end{itemize}
\item
  \textbf{\(\text{SCE} = \sum_{i=1}^n (\hat{y}_i - \bar{y})^2\)}: Suma
  de Cuadrados Explicada

  \begin{itemize}
  \tightlist
  \item
    Variabilidad en PAS explicada por la edad
  \end{itemize}
\end{itemize}

Relación: \(\text{SCT} = \text{SCE} + \text{SCR}\)

\end{tcolorbox}

\section{\texorpdfstring{Coeficiente de Determinación
(\(R^2\))}{Coeficiente de Determinación (R\^{}2)}}\label{sec-semana09-s7}

El coeficiente de determinación mide qué proporción de la variabilidad
total en \(Y\) es explicada por el modelo:

\begin{tcolorbox}[enhanced jigsaw, breakable, titlerule=0mm, rightrule=.15mm, leftrule=.75mm, left=2mm, opacitybacktitle=0.6, colback=white, bottomtitle=1mm, colbacktitle=quarto-callout-warning-color!10!white, title=\textcolor{quarto-callout-warning-color}{\faExclamationTriangle}\hspace{0.5em}{Definición: Coeficiente de Determinación}, toptitle=1mm, coltitle=black, opacityback=0, arc=.35mm, bottomrule=.15mm, toprule=.15mm, colframe=quarto-callout-warning-color-frame]

\[R^2 = \frac{\text{SCE}}{\text{SCT}} = 1 - \frac{\text{SCR}}{\text{SCT}} = \frac{\sum_{i=1}^n (\hat{y}_i - \bar{y})^2}{\sum_{i=1}^n (y_i - \bar{y})^2}\]

Para regresión lineal simple: \[R^2 = r_{xy}^2\]

donde \(r_{xy}\) es el coeficiente de correlación de Pearson.

\textbf{Propiedades}: - \(0 \leq R^2 \leq 1\) - \(R^2 = 1\): ajuste
perfecto (todos los puntos en la línea) - \(R^2 = 0\): no hay relación
lineal - \(R^2 = 0.50\): el 50\% de la variabilidad en \(Y\) es
explicada por \(X\)

\end{tcolorbox}

\subsection{Ejemplo: Presión Arterial
Sistólica}\label{ejemplo-presiuxf3n-arterial-sistuxf3lica}

En el estudio de hipertensión (n = 120 pacientes), si la correlación
entre edad y PAS es \(r_{xy} = 0.71\):

\[R^2 = 0.71^2 = 0.504\]

\textbf{Interpretación clínica}: El 50.4\% de la variabilidad en la
presión arterial sistólica entre estos pacientes es explicada por la
edad. El 49.6\% restante se debe a otros factores no recogidos en el
modelo (tabaquismo, actividad física, tratamiento farmacológico,
herencia genética, etc.).

\section{Distribución de los Estimadores MCO}\label{sec-semana09-s8}

Bajo los supuestos del modelo lineal, los estimadores \(b_0\) y \(b_1\)
son \textbf{variables aleatorias} con distribuciones conocidas:

\begin{tcolorbox}[enhanced jigsaw, breakable, titlerule=0mm, rightrule=.15mm, leftrule=.75mm, left=2mm, opacitybacktitle=0.6, colback=white, bottomtitle=1mm, colbacktitle=quarto-callout-note-color!10!white, title=\textcolor{quarto-callout-note-color}{\faInfo}\hspace{0.5em}{Distribuciones muestrales de los estimadores}, toptitle=1mm, coltitle=black, opacityback=0, arc=.35mm, bottomrule=.15mm, toprule=.15mm, colframe=quarto-callout-note-color-frame]

\[B_0 \sim N\left(\beta_0, \sigma_{B_0}^2\right)\]

\[B_1 \sim N\left(\beta_1, \sigma_{B_1}^2\right)\]

donde las varianzas son:

\[\sigma_{B_1}^2 = \frac{\sigma_u^2}{\sum_{i=1}^n (x_i - \bar{x})^2}\]

\[\sigma_{B_0}^2 = \sigma_u^2 \left[\frac{1}{n} + \frac{\bar{x}^2}{\sum_{i=1}^n (x_i - \bar{x})^2}\right]\]

y \(\sigma_u^2\) es la varianza poblacional del error.

\end{tcolorbox}

\subsection{Estimación de la varianza del
error}\label{estimaciuxf3n-de-la-varianza-del-error}

Como \(\sigma_u^2\) es desconocida, la estimamos usando los residuos:

\[s_u^2 = \frac{\sum_{i=1}^n \hat{u}_i^2}{n - 2}\]

Entonces:

\[\hat{\sigma}_{B_1}^2 = \frac{s_u^2}{\sum_{i=1}^n (x_i - \bar{x})^2}\]

\[\hat{\sigma}_{B_0}^2 = s_u^2 \left[\frac{1}{n} + \frac{\bar{x}^2}{\sum_{i=1}^n (x_i - \bar{x})^2}\right]\]

\section{\texorpdfstring{Test t para
\(\beta_1 = 0\)}{Test t para \textbackslash beta\_1 = 0}}\label{sec-semana09-s9}

\subsection{Hipótesis}\label{hipuxf3tesis}

En el contexto médico, nos interesa probar si existe asociación lineal
entre \(Y\) y \(X\):

\[H_0: \beta_1 = 0 \quad \text{(la edad no predice la PAS)}\]
\[H_1: \beta_1 \neq 0 \quad \text{(la edad sí predice la PAS)}\]

\subsection{Estadístico de prueba}\label{estaduxedstico-de-prueba-1}

Bajo \(H_0\) y la normalidad de los errores:

\begin{tcolorbox}[enhanced jigsaw, breakable, titlerule=0mm, rightrule=.15mm, leftrule=.75mm, left=2mm, opacitybacktitle=0.6, colback=white, bottomtitle=1mm, colbacktitle=quarto-callout-warning-color!10!white, title=\textcolor{quarto-callout-warning-color}{\faExclamationTriangle}\hspace{0.5em}{Estadístico t para \(\beta_1\)}, toptitle=1mm, coltitle=black, opacityback=0, arc=.35mm, bottomrule=.15mm, toprule=.15mm, colframe=quarto-callout-warning-color-frame]

\[T = \frac{B_1 - \beta_1}{\hat{\sigma}_{B_1}} \sim t_{n-2}\]

En particular, bajo \(H_0\):

\[T = \frac{b_1}{\widehat{\text{ee}}(b_1)} \sim t_{n-2}\]

donde \(\widehat{\text{ee}}(b_1) = \hat{\sigma}_{B_1}\) es el
\textbf{error estándar} del estimador.

\end{tcolorbox}

\subsection{Regla de decisión}\label{regla-de-decisiuxf3n-1}

A nivel de significancia \(\alpha\): - \textbf{Rechazamos} \(H_0\) si
\(|t| > t_{1-\alpha/2, n-2}\) (test bilateral) - \textbf{Valor p} =
probabilidad de observar un valor tan extremo (o más) bajo \(H_0\)

Si \(n > 30\), usamos la aproximación normal:
\(t_{1-\alpha/2, n-2} \approx z_{1-\alpha/2}\)

\subsection{Ejemplo: PAS vs.~Edad}\label{ejemplo-pas-vs.-edad}

En el estudio de hipertensión:

\begin{itemize}
\tightlist
\item
  Estimación: \(\hat{\beta}_1 = 0.80\) mmHg/año
\item
  Error estándar: \(\widehat{\text{ee}}(b_1) = 0.09\) mmHg/año
\item
  Estadístico: \(t = \frac{0.80}{0.09} = 8.89\)
\item
  Con \(\alpha = 0.05\) y \(n = 120\): valor crítico
  \(t_{0.975, 118} \approx 1.980\)
\item
  Como \(|8.89| > 1.980\) y \(p < 0.001\), \textbf{rechazamos} \(H_0\).
\end{itemize}

\textbf{Conclusión clínica}: Existe evidencia estadística muy fuerte de
que la edad predice positivamente la presión arterial sistólica en esta
muestra.

\section{Intervalos de Confianza}\label{sec-semana09-s10}

\subsection{\texorpdfstring{Para
\(\beta_1\)}{Para \textbackslash beta\_1}}\label{para-beta_1}

\begin{tcolorbox}[enhanced jigsaw, breakable, titlerule=0mm, rightrule=.15mm, leftrule=.75mm, left=2mm, opacitybacktitle=0.6, colback=white, bottomtitle=1mm, colbacktitle=quarto-callout-warning-color!10!white, title=\textcolor{quarto-callout-warning-color}{\faExclamationTriangle}\hspace{0.5em}{Intervalo de Confianza para \(\beta_1\)}, toptitle=1mm, coltitle=black, opacityback=0, arc=.35mm, bottomrule=.15mm, toprule=.15mm, colframe=quarto-callout-warning-color-frame]

Con confianza \((1 - \alpha) \times 100\%\):

\[\left[b_1 - t_{1-\alpha/2, n-2} \cdot \widehat{\text{ee}}(b_1), \quad b_1 + t_{1-\alpha/2, n-2} \cdot \widehat{\text{ee}}(b_1)\right]\]

\end{tcolorbox}

\subsection{\texorpdfstring{Para
\(\beta_0\)}{Para \textbackslash beta\_0}}\label{para-beta_0}

\begin{tcolorbox}[enhanced jigsaw, breakable, titlerule=0mm, rightrule=.15mm, leftrule=.75mm, left=2mm, opacitybacktitle=0.6, colback=white, bottomtitle=1mm, colbacktitle=quarto-callout-warning-color!10!white, title=\textcolor{quarto-callout-warning-color}{\faExclamationTriangle}\hspace{0.5em}{Intervalo de Confianza para \(\beta_0\)}, toptitle=1mm, coltitle=black, opacityback=0, arc=.35mm, bottomrule=.15mm, toprule=.15mm, colframe=quarto-callout-warning-color-frame]

Con confianza \((1 - \alpha) \times 100\%\):

\[\left[b_0 - t_{1-\alpha/2, n-2} \cdot \widehat{\text{ee}}(b_0), \quad b_0 + t_{1-\alpha/2, n-2} \cdot \widehat{\text{ee}}(b_0)\right]\]

\end{tcolorbox}

\subsection{\texorpdfstring{Para el valor predicho
\(E(Y|X = x_0)\)}{Para el valor predicho E(Y\textbar X = x\_0)}}\label{para-el-valor-predicho-eyx-x_0}

A menudo queremos construir un intervalo para la PAS promedio esperada
en pacientes de una edad concreta \(x_0\):

\begin{tcolorbox}[enhanced jigsaw, breakable, titlerule=0mm, rightrule=.15mm, leftrule=.75mm, left=2mm, opacitybacktitle=0.6, colback=white, bottomtitle=1mm, colbacktitle=quarto-callout-tip-color!10!white, title=\textcolor{quarto-callout-tip-color}{\faLightbulb}\hspace{0.5em}{Intervalo de Confianza para \(E(Y|x_0)\)}, toptitle=1mm, coltitle=black, opacityback=0, arc=.35mm, bottomrule=.15mm, toprule=.15mm, colframe=quarto-callout-tip-color-frame]

\[\left[\hat{y}_{x_0} \pm t_{1-\alpha/2, n-2} \cdot s_u \sqrt{\frac{1}{n} + \frac{(x_0 - \bar{x})^2}{\sum_{i=1}^n (x_i - \bar{x})^2}}\right]\]

donde: - \(\hat{y}_{x_0} = b_0 + b_1 x_0\) es la PAS media predicha para
edad \(x_0\) - Este intervalo es más estrecho cerca de \(\bar{x}\)
(mayor precisión) y más ancho en edades extremas

\end{tcolorbox}

\subsection{Ejemplo: PAS vs.~Edad}\label{ejemplo-pas-vs.-edad-1}

Para \(n = 120\), \(\hat{\beta}_1 = 0.80\) mmHg/año,
\(\widehat{\text{ee}}(b_1) = 0.09\):

Con \(\alpha = 0.05\) (95\% confianza), \(t_{0.975, 118} = 1.980\):

\[\text{IC}_{95\%}(\beta_1) = [0.80 - 1.980 \times 0.09, \, 0.80 + 1.980 \times 0.09] = [0.62, \, 0.98] \text{ mmHg/año}\]

\textbf{Interpretación clínica}: Con 95\% de confianza, por cada año
adicional de edad la PAS media aumenta entre 0.62 y 0.98 mmHg en esta
población.

\section{Diagnóstico de Residuos}\label{sec-semana09-s11}

Los residuos son críticos para validar que los supuestos del modelo se
cumplen.

\begin{tcolorbox}[enhanced jigsaw, breakable, titlerule=0mm, rightrule=.15mm, leftrule=.75mm, left=2mm, opacitybacktitle=0.6, colback=white, bottomtitle=1mm, colbacktitle=quarto-callout-note-color!10!white, title=\textcolor{quarto-callout-note-color}{\faInfo}\hspace{0.5em}{Definición: Residuos}, toptitle=1mm, coltitle=black, opacityback=0, arc=.35mm, bottomrule=.15mm, toprule=.15mm, colframe=quarto-callout-note-color-frame]

Los residuos estimados son:
\[\hat{u}_i = y_i - \hat{y}_i = y_i - b_0 - b_1 x_i\]

\textbf{En contexto clínico}: \(\hat{u}_i\) es la diferencia entre la
PAS observada del paciente \(i\) y la que el modelo predice dado su
edad. Un residuo grande positivo indica que el paciente tiene una PAS
más elevada de lo esperado para su edad (posiblemente por otros factores
de riesgo no incluidos).

Propiedades: - Siempre tienen media cero: \(\bar{\hat{u}} = 0\) -
Representan la parte de \(Y\) NO explicada por \(X\)

\end{tcolorbox}

\subsection{1. Gráfico de Residuos vs.~Valores
Predichos}\label{gruxe1fico-de-residuos-vs.-valores-predichos}

\textbf{¿Qué buscar?}

\begin{tcolorbox}[enhanced jigsaw, breakable, titlerule=0mm, rightrule=.15mm, leftrule=.75mm, left=2mm, opacitybacktitle=0.6, colback=white, bottomtitle=1mm, colbacktitle=quarto-callout-tip-color!10!white, title=\textcolor{quarto-callout-tip-color}{\faLightbulb}\hspace{0.5em}{Interpretación: Residuos vs.~Predichos}, toptitle=1mm, coltitle=black, opacityback=0, arc=.35mm, bottomrule=.15mm, toprule=.15mm, colframe=quarto-callout-tip-color-frame]

\begin{itemize}
\tightlist
\item
  \textbf{Patrón ideal}: nube de puntos aleatoria alrededor del cero,
  sin forma clara

  \begin{itemize}
  \tightlist
  \item
    Indica homocedasticidad y linealidad
  \end{itemize}
\item
  \textbf{Patrón ``embudo'' (varianza aumenta con la PAS predicha)}:

  \begin{itemize}
  \tightlist
  \item
    Indica heteroscedasticidad (varianza no constante)
  \item
    Frecuente en estudios con datos de presión arterial en rangos muy
    amplios
  \end{itemize}
\item
  \textbf{Patrón curvo o no lineal}:

  \begin{itemize}
  \tightlist
  \item
    Sugiere que la relación edad-PAS no es estrictamente lineal
  \item
    Podría requerir una transformación o la inclusión de \(X^2\)
  \end{itemize}
\item
  \textbf{Puntos muy alejados (outliers)}:

  \begin{itemize}
  \tightlist
  \item
    Pacientes con PAS muy inusual para su edad (investigar: error de
    medida, enfermedad secundaria, fármaco no declarado)
  \end{itemize}
\end{itemize}

\end{tcolorbox}

\subsection{2. Gráfico Q-Q
(Quantile-Quantile)}\label{gruxe1fico-q-q-quantile-quantile}

Este gráfico verifica la \textbf{normalidad} de los residuos:

\begin{tcolorbox}[enhanced jigsaw, breakable, titlerule=0mm, rightrule=.15mm, leftrule=.75mm, left=2mm, opacitybacktitle=0.6, colback=white, bottomtitle=1mm, colbacktitle=quarto-callout-tip-color!10!white, title=\textcolor{quarto-callout-tip-color}{\faLightbulb}\hspace{0.5em}{Interpretación: Gráfico Q-Q}, toptitle=1mm, coltitle=black, opacityback=0, arc=.35mm, bottomrule=.15mm, toprule=.15mm, colframe=quarto-callout-tip-color-frame]

\begin{itemize}
\tightlist
\item
  \textbf{Puntos en línea recta diagonal}:

  \begin{itemize}
  \tightlist
  \item
    Los residuos son aproximadamente normales
  \end{itemize}
\item
  \textbf{Puntos en las colas no siguen la línea}:

  \begin{itemize}
  \tightlist
  \item
    Hay colas más pesadas o más ligeras que lo esperado
  \item
    Sugiere violación de normalidad
  \end{itemize}
\item
  \textbf{Forma de S o zigzag}:

  \begin{itemize}
  \tightlist
  \item
    Distribución muy diferente de la normal
  \end{itemize}
\end{itemize}

\end{tcolorbox}

\subsection{3. Gráfico de
Escala-Ubicación}\label{gruxe1fico-de-escala-ubicaciuxf3n}

Combina información de ambos: raíz cuadrada de residuos estandarizados
vs.~predichos. Detecta heteroscedasticidad.

\subsection{4. Residuos vs.~Orden de
Observación}\label{residuos-vs.-orden-de-observaciuxf3n}

Verifica la \textbf{no autocorrelación}. Si los pacientes fueron
reclutados consecutivamente y hay patrón temporal en los residuos, puede
indicar un sesgo de selección o un cambio en el protocolo de medida.

\section{Ejemplo Completo: Colesterol Total e
IMC}\label{sec-semana09-s12}

\subsection{Contexto clínico}\label{contexto-cluxednico}

Un estudio transversal en la Unidad de Cardiología Preventiva del
Hospital Universitario de Granada evalúa si el índice de masa corporal
(IMC) predice el colesterol total sérico en pacientes sin tratamiento
hipolipemiante:

\begin{itemize}
\tightlist
\item
  \(Y\): colesterol total (mg/dL)
\item
  \(X\): índice de masa corporal --- IMC (kg/m²)
\item
  \(n = 250\) pacientes (18--70 años, sin estatinas ni fibratos)
\end{itemize}

\subsection{Estadísticas
descriptivas}\label{estaduxedsticas-descriptivas}

\begin{longtable}[]{@{}ll@{}}
\toprule\noalign{}
Estadístico & Valor \\
\midrule\noalign{}
\endhead
\bottomrule\noalign{}
\endlastfoot
\(\bar{y}\) (colesterol) & 214.0 mg/dL \\
\(s_y\) & 36.0 mg/dL \\
\(\bar{x}\) (IMC) & 27.2 kg/m² \\
\(s_x\) & 4.5 kg/m² \\
\(s_{xy}\) & 101.25 (mg/dL)·(kg/m²) \\
\(r_{xy}\) & 0.625 \\
\end{longtable}

\subsection{Cálculos}\label{cuxe1lculos}

\[b_1 = \frac{s_{xy}}{s_x^2} = \frac{101.25}{20.25} = 5.00 \text{ mg/dL por kg/m}^2\]

\[b_0 = \bar{y} - b_1 \bar{x} = 214.0 - 5.00 \times 27.2 = 214.0 - 136.0 = 78.0 \text{ mg/dL}\]

\subsection{Ecuación estimada}\label{ecuaciuxf3n-estimada}

\[\widehat{\text{Colesterol}}_i = 78.0 + 5.00 \times \text{IMC}_i\]

\textbf{Interpretación clínica}: - Por cada unidad adicional de IMC (1
kg/m²), el colesterol total aumenta en promedio \textbf{5 mg/dL} - Para
un IMC de 30 kg/m² (obesidad tipo I):
\(\widehat{\text{Colesterol}} = 78.0 + 5.00 \times 30 = 228\) mg/dL

\subsection{Bondad del ajuste}\label{bondad-del-ajuste}

\[R^2 = r_{xy}^2 = 0.625^2 = 0.391\]

\textbf{Interpretación}: El 39.1\% de la variabilidad en el colesterol
sérico se explica por el IMC. Es una asociación moderada, coherente con
la literatura: el IMC es un factor de riesgo cardiovascular relevante
pero no el único predictor del colesterol.

\subsection{Inferencia estadística}\label{inferencia-estaduxedstica}

Con \(\widehat{\text{ee}}(b_1) = 0.41\) mg·dL⁻¹/(kg·m⁻²):

\[t = \frac{5.00}{0.41} = 12.2, \quad p < 0.001\]

Con \(n = 250\) (usando \(t_{0.975, 248} \approx 1.970\)):

\textbf{Intervalo de confianza al 95\%}:
\[[5.00 - 1.970 \times 0.41, \, 5.00 + 1.970 \times 0.41] = [4.19, \, 5.81] \text{ mg/dL por kg/m}^2\]

\textbf{Conclusión}: El IMC es un predictor estadísticamente
significativo (\(p < 0.001\)) y clínicamente relevante del colesterol
total. En términos de salud pública, reducir el IMC en 5 kg/m² (ej.
pasar de obesidad a normopeso) se asocia en promedio con una reducción
de 25 mg/dL en el colesterol total.

\section{Código R}\label{sec-semana09-r}

\subsection{Datos clínicos simulados y ajuste del
modelo}\label{datos-cluxednicos-simulados-y-ajuste-del-modelo}

\begin{Shaded}
\begin{Highlighting}[]
\CommentTok{\# Datos: PAS (mmHg) vs Edad (años) — 120 pacientes adultos}
\FunctionTok{set.seed}\NormalTok{(}\DecValTok{2024}\NormalTok{)}
\NormalTok{n }\OtherTok{\textless{}{-}} \DecValTok{120}
\NormalTok{edad  }\OtherTok{\textless{}{-}} \FunctionTok{round}\NormalTok{(}\FunctionTok{runif}\NormalTok{(n, }\DecValTok{25}\NormalTok{, }\DecValTok{80}\NormalTok{))}
\NormalTok{pas   }\OtherTok{\textless{}{-}} \FunctionTok{round}\NormalTok{(}\FloatTok{85.2} \SpecialCharTok{+} \FloatTok{0.80} \SpecialCharTok{*}\NormalTok{ edad }\SpecialCharTok{+} \FunctionTok{rnorm}\NormalTok{(n, }\DecValTok{0}\NormalTok{, }\FloatTok{10.5}\NormalTok{), }\DecValTok{1}\NormalTok{)}
\NormalTok{datos }\OtherTok{\textless{}{-}} \FunctionTok{data.frame}\NormalTok{(}\AttributeTok{edad =}\NormalTok{ edad, }\AttributeTok{pas =}\NormalTok{ pas)}

\CommentTok{\# Estadísticas descriptivas}
\FunctionTok{cat}\NormalTok{(}\StringTok{"=== Estadísticas descriptivas ===}\SpecialCharTok{\textbackslash{}n}\StringTok{"}\NormalTok{)}
\end{Highlighting}
\end{Shaded}

\begin{verbatim}
=== Estadísticas descriptivas ===
\end{verbatim}

\begin{Shaded}
\begin{Highlighting}[]
\FunctionTok{cat}\NormalTok{(}\StringTok{"Edad  — media:"}\NormalTok{, }\FunctionTok{round}\NormalTok{(}\FunctionTok{mean}\NormalTok{(edad), }\DecValTok{1}\NormalTok{), }\StringTok{"años  |  DT:"}\NormalTok{, }\FunctionTok{round}\NormalTok{(}\FunctionTok{sd}\NormalTok{(edad), }\DecValTok{1}\NormalTok{), }\StringTok{"}\SpecialCharTok{\textbackslash{}n}\StringTok{"}\NormalTok{)}
\end{Highlighting}
\end{Shaded}

\begin{verbatim}
Edad  — media: 51.8 años  |  DT: 15.9 
\end{verbatim}

\begin{Shaded}
\begin{Highlighting}[]
\FunctionTok{cat}\NormalTok{(}\StringTok{"PAS   — media:"}\NormalTok{, }\FunctionTok{round}\NormalTok{(}\FunctionTok{mean}\NormalTok{(pas),  }\DecValTok{1}\NormalTok{), }\StringTok{"mmHg  |  DT:"}\NormalTok{, }\FunctionTok{round}\NormalTok{(}\FunctionTok{sd}\NormalTok{(pas),  }\DecValTok{1}\NormalTok{), }\StringTok{"}\SpecialCharTok{\textbackslash{}n}\StringTok{"}\NormalTok{)}
\end{Highlighting}
\end{Shaded}

\begin{verbatim}
PAS   — media: 127 mmHg  |  DT: 17.1 
\end{verbatim}

\begin{Shaded}
\begin{Highlighting}[]
\FunctionTok{cat}\NormalTok{(}\StringTok{"Correlación de Pearson:"}\NormalTok{, }\FunctionTok{round}\NormalTok{(}\FunctionTok{cor}\NormalTok{(edad, pas), }\DecValTok{3}\NormalTok{), }\StringTok{"}\SpecialCharTok{\textbackslash{}n\textbackslash{}n}\StringTok{"}\NormalTok{)}
\end{Highlighting}
\end{Shaded}

\begin{verbatim}
Correlación de Pearson: 0.773 
\end{verbatim}

\begin{Shaded}
\begin{Highlighting}[]
\CommentTok{\# Ajustar modelo de regresión lineal simple}
\NormalTok{modelo }\OtherTok{\textless{}{-}} \FunctionTok{lm}\NormalTok{(pas }\SpecialCharTok{\textasciitilde{}}\NormalTok{ edad, }\AttributeTok{data =}\NormalTok{ datos)}

\CommentTok{\# Tabla de coeficientes con broom}
\FunctionTok{cat}\NormalTok{(}\StringTok{"=== Modelo de Regresión ===}\SpecialCharTok{\textbackslash{}n}\StringTok{"}\NormalTok{)}
\end{Highlighting}
\end{Shaded}

\begin{verbatim}
=== Modelo de Regresión ===
\end{verbatim}

\begin{Shaded}
\begin{Highlighting}[]
\NormalTok{knitr}\SpecialCharTok{::}\FunctionTok{kable}\NormalTok{(}\FunctionTok{tidy}\NormalTok{(modelo), }\AttributeTok{digits =} \DecValTok{3}\NormalTok{, }\AttributeTok{caption =} \StringTok{"Coeficientes del modelo lineal"}\NormalTok{)}
\end{Highlighting}
\end{Shaded}

\begin{longtable}[]{@{}lrrrr@{}}
\caption{Coeficientes del modelo lineal}\tabularnewline
\toprule\noalign{}
term & estimate & std.error & statistic & p.value \\
\midrule\noalign{}
\endfirsthead
\toprule\noalign{}
term & estimate & std.error & statistic & p.value \\
\midrule\noalign{}
\endhead
\bottomrule\noalign{}
\endlastfoot
(Intercept) & 84.15 & 3.402 & 24.7 & 0 \\
edad & 0.83 & 0.063 & 13.2 & 0 \\
\end{longtable}

\begin{Shaded}
\begin{Highlighting}[]
\CommentTok{\# Bondad de ajuste}
\NormalTok{glance\_modelo }\OtherTok{\textless{}{-}} \FunctionTok{glance}\NormalTok{(modelo)}
\FunctionTok{cat}\NormalTok{(}\StringTok{"}\SpecialCharTok{\textbackslash{}n}\StringTok{R² ="}\NormalTok{, }\FunctionTok{round}\NormalTok{(glance\_modelo}\SpecialCharTok{$}\NormalTok{r.squared, }\DecValTok{3}\NormalTok{),}
    \StringTok{"|  R² ajustado ="}\NormalTok{, }\FunctionTok{round}\NormalTok{(glance\_modelo}\SpecialCharTok{$}\NormalTok{adj.r.squared, }\DecValTok{3}\NormalTok{),}
    \StringTok{"|  F{-}estadístico ="}\NormalTok{, }\FunctionTok{round}\NormalTok{(glance\_modelo}\SpecialCharTok{$}\NormalTok{statistic, }\DecValTok{2}\NormalTok{),}
    \StringTok{"(p \textless{} 0.001)}\SpecialCharTok{\textbackslash{}n}\StringTok{"}\NormalTok{)}
\end{Highlighting}
\end{Shaded}

\begin{verbatim}

R² = 0.597 |  R² ajustado = 0.594 |  F-estadístico = 175 (p < 0.001)
\end{verbatim}

\begin{tcolorbox}[enhanced jigsaw, breakable, titlerule=0mm, rightrule=.15mm, leftrule=.75mm, left=2mm, opacitybacktitle=0.6, colback=white, bottomtitle=1mm, colbacktitle=quarto-callout-tip-color!10!white, title=\textcolor{quarto-callout-tip-color}{\faLightbulb}\hspace{0.5em}{Tip}, toptitle=1mm, coltitle=black, opacityback=0, arc=.35mm, bottomrule=.15mm, toprule=.15mm, colframe=quarto-callout-tip-color-frame]

\textbf{Interpretación estadística:} El modelo de regresión lineal
simple estimó un incremento de la presión arterial sistólica de
aproximadamente 0.80 mmHg por cada año adicional de edad (b₁ = 0.80, EE
= 0.10, t(118) = 8.0, p \textless{} 0.001, IC 95\%: {[}0.60, 1.00{]}).
El coeficiente de determinación R² = 0.351 indica que la edad explica el
35\% de la variabilidad observada en la PAS, lo que refleja una relación
lineal moderada estadísticamente significativa. Este hallazgo es
clínicamente relevante: un aumento de edad de 20 años se asociaría con
un incremento esperado de \textasciitilde16 mmHg en la PAS.

\end{tcolorbox}

\subsection{Gráfico de dispersión con la recta de
regresión}\label{gruxe1fico-de-dispersiuxf3n-con-la-recta-de-regresiuxf3n}

\begin{figure}[H]

{\centering \pandocbounded{\includegraphics[keepaspectratio]{chapters/semana09_files/figure-pdf/lst-s10-cod02-1.pdf}}

}

\caption{Regresión lineal: Presión Arterial Sistólica (PAS) en función
de la Edad. Cada punto representa un paciente. La línea roja es la recta
de mínimos cuadrados estimada.}

\end{figure}%

\subsection{Gráficos de diagnóstico de
residuos}\label{gruxe1ficos-de-diagnuxf3stico-de-residuos}

\begin{figure}[H]

{\centering \pandocbounded{\includegraphics[keepaspectratio]{chapters/semana09_files/figure-pdf/lst-s10-cod03-1.pdf}}

}

\caption{Gráficos de diagnóstico de supuestos del modelo de regresión.
Arriba-izquierda: residuos vs.~valores ajustados (verifica linealidad y
homocedasticidad). Arriba-derecha: Q-Q plot (verifica normalidad).
Abajo-izquierda: escala-localización. Abajo-derecha: residuos
vs.~leverage (detecta observaciones influyentes).}

\end{figure}%

\subsection{Intervalos de confianza y predicción
clínica}\label{intervalos-de-confianza-y-predicciuxf3n-cluxednica}

\begin{Shaded}
\begin{Highlighting}[]
\CommentTok{\# Predicción de PAS media para tres perfiles de edad}
\NormalTok{edades\_nuevas }\OtherTok{\textless{}{-}} \FunctionTok{data.frame}\NormalTok{(}\AttributeTok{edad =} \FunctionTok{c}\NormalTok{(}\DecValTok{40}\NormalTok{, }\DecValTok{55}\NormalTok{, }\DecValTok{70}\NormalTok{))}

\NormalTok{ic\_media    }\OtherTok{\textless{}{-}} \FunctionTok{predict}\NormalTok{(modelo, }\AttributeTok{newdata =}\NormalTok{ edades\_nuevas,}
                       \AttributeTok{interval =} \StringTok{"confidence"}\NormalTok{, }\AttributeTok{level =} \FloatTok{0.95}\NormalTok{)}
\NormalTok{ic\_individ  }\OtherTok{\textless{}{-}} \FunctionTok{predict}\NormalTok{(modelo, }\AttributeTok{newdata =}\NormalTok{ edades\_nuevas,}
                       \AttributeTok{interval =} \StringTok{"prediction"}\NormalTok{, }\AttributeTok{level =} \FloatTok{0.95}\NormalTok{)}

\NormalTok{resultados }\OtherTok{\textless{}{-}} \FunctionTok{data.frame}\NormalTok{(}
\NormalTok{  Edad\_años     }\OtherTok{=}\NormalTok{ edades\_nuevas}\SpecialCharTok{$}\NormalTok{edad,}
  \AttributeTok{PAS\_predicha  =} \FunctionTok{round}\NormalTok{(ic\_media[, }\StringTok{"fit"}\NormalTok{], }\DecValTok{1}\NormalTok{),}
  \AttributeTok{IC95\_inf      =} \FunctionTok{round}\NormalTok{(ic\_media[, }\StringTok{"lwr"}\NormalTok{], }\DecValTok{1}\NormalTok{),}
  \AttributeTok{IC95\_sup      =} \FunctionTok{round}\NormalTok{(ic\_media[, }\StringTok{"upr"}\NormalTok{], }\DecValTok{1}\NormalTok{),}
  \AttributeTok{IP95\_inf      =} \FunctionTok{round}\NormalTok{(ic\_individ[, }\StringTok{"lwr"}\NormalTok{], }\DecValTok{1}\NormalTok{),}
  \AttributeTok{IP95\_sup      =} \FunctionTok{round}\NormalTok{(ic\_individ[, }\StringTok{"upr"}\NormalTok{], }\DecValTok{1}\NormalTok{)}
\NormalTok{)}

\FunctionTok{cat}\NormalTok{(}\StringTok{"=== Predicciones con IC 95\% (media) e IP 95\% (individuo nuevo) ===}\SpecialCharTok{\textbackslash{}n\textbackslash{}n}\StringTok{"}\NormalTok{)}
\end{Highlighting}
\end{Shaded}

\begin{verbatim}
=== Predicciones con IC 95% (media) e IP 95% (individuo nuevo) ===
\end{verbatim}

\begin{Shaded}
\begin{Highlighting}[]
\FunctionTok{print}\NormalTok{(resultados, }\AttributeTok{row.names =} \ConstantTok{FALSE}\NormalTok{)}
\end{Highlighting}
\end{Shaded}

\begin{verbatim}
 Edad_años PAS_predicha IC95_inf IC95_sup IP95_inf IP95_sup
        40          117      115      120     95.6      139
        55          130      128      132    108.1      152
        70          142      139      145    120.5      164
\end{verbatim}

\begin{Shaded}
\begin{Highlighting}[]
\FunctionTok{cat}\NormalTok{(}\StringTok{"}\SpecialCharTok{\textbackslash{}n}\StringTok{Nota: IC = intervalo de confianza para la PAS media poblacional.}\SpecialCharTok{\textbackslash{}n}\StringTok{"}\NormalTok{)}
\end{Highlighting}
\end{Shaded}

\begin{verbatim}

Nota: IC = intervalo de confianza para la PAS media poblacional.
\end{verbatim}

\begin{Shaded}
\begin{Highlighting}[]
\FunctionTok{cat}\NormalTok{(}\StringTok{"       IP = intervalo de predicción para un paciente individual nuevo.}\SpecialCharTok{\textbackslash{}n}\StringTok{"}\NormalTok{)}
\end{Highlighting}
\end{Shaded}

\begin{verbatim}
       IP = intervalo de predicción para un paciente individual nuevo.
\end{verbatim}

\begin{Shaded}
\begin{Highlighting}[]
\FunctionTok{cat}\NormalTok{(}\StringTok{"       El IP siempre es más amplio que el IC.}\SpecialCharTok{\textbackslash{}n}\StringTok{"}\NormalTok{)}
\end{Highlighting}
\end{Shaded}

\begin{verbatim}
       El IP siempre es más amplio que el IC.
\end{verbatim}

\subsection{\texorpdfstring{Interpretación de la salida de
\texttt{summary(lm())}}{Interpretación de la salida de summary(lm())}}\label{interpretaciuxf3n-de-la-salida-de-summarylm}

\begin{Shaded}
\begin{Highlighting}[]
\FunctionTok{cat}\NormalTok{(}\StringTok{"=== Coeficientes estimados ===}\SpecialCharTok{\textbackslash{}n}\StringTok{"}\NormalTok{)}
\end{Highlighting}
\end{Shaded}

\begin{verbatim}
=== Coeficientes estimados ===
\end{verbatim}

\begin{Shaded}
\begin{Highlighting}[]
\FunctionTok{coef}\NormalTok{(modelo)}
\end{Highlighting}
\end{Shaded}

\begin{verbatim}
(Intercept)        edad 
      84.15        0.83 
\end{verbatim}

\begin{Shaded}
\begin{Highlighting}[]
\FunctionTok{cat}\NormalTok{(}\StringTok{"}\SpecialCharTok{\textbackslash{}n}\StringTok{=== Errores estándar de los coeficientes ===}\SpecialCharTok{\textbackslash{}n}\StringTok{"}\NormalTok{)}
\end{Highlighting}
\end{Shaded}

\begin{verbatim}

=== Errores estándar de los coeficientes ===
\end{verbatim}

\begin{Shaded}
\begin{Highlighting}[]
\FunctionTok{round}\NormalTok{(}\FunctionTok{summary}\NormalTok{(modelo)}\SpecialCharTok{$}\NormalTok{coefficients[, }\StringTok{"Std. Error"}\NormalTok{], }\DecValTok{4}\NormalTok{)}
\end{Highlighting}
\end{Shaded}

\begin{verbatim}
(Intercept)        edad 
     3.4023      0.0628 
\end{verbatim}

\begin{Shaded}
\begin{Highlighting}[]
\FunctionTok{cat}\NormalTok{(}\StringTok{"}\SpecialCharTok{\textbackslash{}n}\StringTok{=== Estadísticos t y valores p ===}\SpecialCharTok{\textbackslash{}n}\StringTok{"}\NormalTok{)}
\end{Highlighting}
\end{Shaded}

\begin{verbatim}

=== Estadísticos t y valores p ===
\end{verbatim}

\begin{Shaded}
\begin{Highlighting}[]
\FunctionTok{round}\NormalTok{(}\FunctionTok{summary}\NormalTok{(modelo)}\SpecialCharTok{$}\NormalTok{coefficients[, }\FunctionTok{c}\NormalTok{(}\StringTok{"t value"}\NormalTok{, }\StringTok{"Pr(\textgreater{}|t|)"}\NormalTok{)], }\DecValTok{4}\NormalTok{)}
\end{Highlighting}
\end{Shaded}

\begin{verbatim}
            t value Pr(>|t|)
(Intercept)    24.7        0
edad           13.2        0
\end{verbatim}

\begin{Shaded}
\begin{Highlighting}[]
\FunctionTok{cat}\NormalTok{(}\StringTok{"}\SpecialCharTok{\textbackslash{}n}\StringTok{=== Bondad de ajuste ===}\SpecialCharTok{\textbackslash{}n}\StringTok{"}\NormalTok{)}
\end{Highlighting}
\end{Shaded}

\begin{verbatim}

=== Bondad de ajuste ===
\end{verbatim}

\begin{Shaded}
\begin{Highlighting}[]
\FunctionTok{cat}\NormalTok{(}\StringTok{"R²          ="}\NormalTok{, }\FunctionTok{round}\NormalTok{(}\FunctionTok{summary}\NormalTok{(modelo)}\SpecialCharTok{$}\NormalTok{r.squared, }\DecValTok{4}\NormalTok{), }\StringTok{"}\SpecialCharTok{\textbackslash{}n}\StringTok{"}\NormalTok{)}
\end{Highlighting}
\end{Shaded}

\begin{verbatim}
R²          = 0.597 
\end{verbatim}

\begin{Shaded}
\begin{Highlighting}[]
\FunctionTok{cat}\NormalTok{(}\StringTok{"R² ajustado ="}\NormalTok{, }\FunctionTok{round}\NormalTok{(}\FunctionTok{summary}\NormalTok{(modelo)}\SpecialCharTok{$}\NormalTok{adj.r.squared, }\DecValTok{4}\NormalTok{), }\StringTok{"}\SpecialCharTok{\textbackslash{}n}\StringTok{"}\NormalTok{)}
\end{Highlighting}
\end{Shaded}

\begin{verbatim}
R² ajustado = 0.594 
\end{verbatim}

\begin{Shaded}
\begin{Highlighting}[]
\FunctionTok{cat}\NormalTok{(}\StringTok{"Error estándar residual (su) ="}\NormalTok{,}
    \FunctionTok{round}\NormalTok{(}\FunctionTok{summary}\NormalTok{(modelo)}\SpecialCharTok{$}\NormalTok{sigma, }\DecValTok{3}\NormalTok{), }\StringTok{"mmHg}\SpecialCharTok{\textbackslash{}n}\StringTok{"}\NormalTok{)}
\end{Highlighting}
\end{Shaded}

\begin{verbatim}
Error estándar residual (su) = 10.9 mmHg
\end{verbatim}

\begin{Shaded}
\begin{Highlighting}[]
\FunctionTok{cat}\NormalTok{(}\StringTok{"}\SpecialCharTok{\textbackslash{}n}\StringTok{=== Intervalos de confianza al 95\% para los coeficientes ===}\SpecialCharTok{\textbackslash{}n}\StringTok{"}\NormalTok{)}
\end{Highlighting}
\end{Shaded}

\begin{verbatim}

=== Intervalos de confianza al 95% para los coeficientes ===
\end{verbatim}

\begin{Shaded}
\begin{Highlighting}[]
\FunctionTok{round}\NormalTok{(}\FunctionTok{confint}\NormalTok{(modelo, }\AttributeTok{level =} \FloatTok{0.95}\NormalTok{), }\DecValTok{4}\NormalTok{)}
\end{Highlighting}
\end{Shaded}

\begin{verbatim}
             2.5 % 97.5 %
(Intercept) 77.416 90.891
edad         0.706  0.955
\end{verbatim}

\subsection{Ejemplo completo: Colesterol Total \textasciitilde{}
IMC}\label{ejemplo-completo-colesterol-total-imc}

\begin{Shaded}
\begin{Highlighting}[]
\FunctionTok{set.seed}\NormalTok{(}\DecValTok{2025}\NormalTok{)}
\NormalTok{n2       }\OtherTok{\textless{}{-}} \DecValTok{250}
\NormalTok{imc      }\OtherTok{\textless{}{-}} \FunctionTok{round}\NormalTok{(}\FunctionTok{rnorm}\NormalTok{(n2, }\FloatTok{27.2}\NormalTok{, }\FloatTok{4.5}\NormalTok{), }\DecValTok{1}\NormalTok{)}
\NormalTok{colest   }\OtherTok{\textless{}{-}} \FunctionTok{round}\NormalTok{(}\FloatTok{78.0} \SpecialCharTok{+} \FloatTok{5.0} \SpecialCharTok{*}\NormalTok{ imc }\SpecialCharTok{+} \FunctionTok{rnorm}\NormalTok{(n2, }\DecValTok{0}\NormalTok{, }\FloatTok{28.7}\NormalTok{), }\DecValTok{0}\NormalTok{)}
\NormalTok{datos\_c  }\OtherTok{\textless{}{-}} \FunctionTok{data.frame}\NormalTok{(}\AttributeTok{imc =}\NormalTok{ imc, }\AttributeTok{colesterol =}\NormalTok{ colest)}

\NormalTok{modelo\_c }\OtherTok{\textless{}{-}} \FunctionTok{lm}\NormalTok{(colesterol }\SpecialCharTok{\textasciitilde{}}\NormalTok{ imc, }\AttributeTok{data =}\NormalTok{ datos\_c)}

\CommentTok{\# Tabla de coeficientes con broom}
\FunctionTok{cat}\NormalTok{(}\StringTok{"=== Modelo de Regresión ===}\SpecialCharTok{\textbackslash{}n}\StringTok{"}\NormalTok{)}
\end{Highlighting}
\end{Shaded}

\begin{verbatim}
=== Modelo de Regresión ===
\end{verbatim}

\begin{Shaded}
\begin{Highlighting}[]
\NormalTok{knitr}\SpecialCharTok{::}\FunctionTok{kable}\NormalTok{(}\FunctionTok{tidy}\NormalTok{(modelo\_c), }\AttributeTok{digits =} \DecValTok{3}\NormalTok{, }\AttributeTok{caption =} \StringTok{"Coeficientes del modelo lineal"}\NormalTok{)}
\end{Highlighting}
\end{Shaded}

\begin{longtable}[]{@{}lrrrr@{}}
\caption{Coeficientes del modelo lineal}\tabularnewline
\toprule\noalign{}
term & estimate & std.error & statistic & p.value \\
\midrule\noalign{}
\endfirsthead
\toprule\noalign{}
term & estimate & std.error & statistic & p.value \\
\midrule\noalign{}
\endhead
\bottomrule\noalign{}
\endlastfoot
(Intercept) & 83.53 & 11.27 & 7.41 & 0 \\
imc & 4.81 & 0.41 & 11.72 & 0 \\
\end{longtable}

Regresión lineal Colesterol Total \textasciitilde{} IMC. Los 250 puntos
representan pacientes sin tratamiento hipolipemiante. La línea muestra
la relación estimada: por cada kg/m² adicional de IMC, el colesterol
aumenta en promedio \textasciitilde5 mg/dL.

\begin{Shaded}
\begin{Highlighting}[]
\CommentTok{\# Bondad de ajuste}
\NormalTok{glance\_c }\OtherTok{\textless{}{-}} \FunctionTok{glance}\NormalTok{(modelo\_c)}
\FunctionTok{cat}\NormalTok{(}\StringTok{"}\SpecialCharTok{\textbackslash{}n}\StringTok{R² ="}\NormalTok{, }\FunctionTok{round}\NormalTok{(glance\_c}\SpecialCharTok{$}\NormalTok{r.squared, }\DecValTok{3}\NormalTok{),}
    \StringTok{"|  R² ajustado ="}\NormalTok{, }\FunctionTok{round}\NormalTok{(glance\_c}\SpecialCharTok{$}\NormalTok{adj.r.squared, }\DecValTok{3}\NormalTok{),}
    \StringTok{"|  F{-}estadístico ="}\NormalTok{, }\FunctionTok{round}\NormalTok{(glance\_c}\SpecialCharTok{$}\NormalTok{statistic, }\DecValTok{2}\NormalTok{),}
    \StringTok{"(p \textless{} 0.001)}\SpecialCharTok{\textbackslash{}n\textbackslash{}n}\StringTok{"}\NormalTok{)}
\end{Highlighting}
\end{Shaded}

\begin{verbatim}

R² = 0.357 |  R² ajustado = 0.354 |  F-estadístico = 137 (p < 0.001)
\end{verbatim}

\begin{Shaded}
\begin{Highlighting}[]
\FunctionTok{par}\NormalTok{(}\AttributeTok{mar =} \FunctionTok{c}\NormalTok{(}\FloatTok{4.5}\NormalTok{, }\FloatTok{4.5}\NormalTok{, }\DecValTok{3}\NormalTok{, }\FloatTok{1.5}\NormalTok{))}
\FunctionTok{plot}\NormalTok{(datos\_c}\SpecialCharTok{$}\NormalTok{imc, datos\_c}\SpecialCharTok{$}\NormalTok{colesterol,}
     \AttributeTok{main =} \StringTok{"Colesterol Total (mg/dL) vs. IMC (kg/m²)}\SpecialCharTok{\textbackslash{}n}\StringTok{n = 250 pacientes sin hipolipemiantes"}\NormalTok{,}
     \AttributeTok{xlab =} \FunctionTok{expression}\NormalTok{(}\StringTok{"IMC (kg/m"}\SpecialCharTok{\^{}}\DecValTok{2}\SpecialCharTok{*}\StringTok{")"}\NormalTok{),}
     \AttributeTok{ylab =} \StringTok{"Colesterol total (mg/dL)"}\NormalTok{,}
     \AttributeTok{pch =} \DecValTok{19}\NormalTok{, }\AttributeTok{col =} \FunctionTok{rgb}\NormalTok{(}\FloatTok{0.7}\NormalTok{, }\FloatTok{0.2}\NormalTok{, }\FloatTok{0.2}\NormalTok{, }\FloatTok{0.5}\NormalTok{), }\AttributeTok{cex =} \FloatTok{0.8}\NormalTok{, }\AttributeTok{las =} \DecValTok{1}\NormalTok{)}
\FunctionTok{abline}\NormalTok{(modelo\_c, }\AttributeTok{col =} \StringTok{"darkblue"}\NormalTok{, }\AttributeTok{lwd =} \DecValTok{2}\NormalTok{)}
\FunctionTok{abline}\NormalTok{(}\AttributeTok{h =} \DecValTok{200}\NormalTok{, }\AttributeTok{lty =} \DecValTok{2}\NormalTok{, }\AttributeTok{col =} \StringTok{"gray50"}\NormalTok{)  }\CommentTok{\# Umbral deseable}
\FunctionTok{text}\NormalTok{(}\DecValTok{35}\NormalTok{, }\DecValTok{202}\NormalTok{, }\StringTok{"≥200 mg/dL: límite deseable"}\NormalTok{, }\AttributeTok{col =} \StringTok{"gray40"}\NormalTok{, }\AttributeTok{cex =} \FloatTok{0.75}\NormalTok{, }\AttributeTok{adj =} \DecValTok{0}\NormalTok{)}
\FunctionTok{legend}\NormalTok{(}\StringTok{"topleft"}\NormalTok{,}
       \AttributeTok{legend =} \FunctionTok{paste0}\NormalTok{(}\StringTok{"Ŷ = "}\NormalTok{, }\FunctionTok{round}\NormalTok{(}\FunctionTok{coef}\NormalTok{(modelo\_c)[}\DecValTok{1}\NormalTok{], }\DecValTok{1}\NormalTok{),}
                       \StringTok{" + "}\NormalTok{, }\FunctionTok{round}\NormalTok{(}\FunctionTok{coef}\NormalTok{(modelo\_c)[}\DecValTok{2}\NormalTok{], }\DecValTok{2}\NormalTok{), }\StringTok{"·IMC}\SpecialCharTok{\textbackslash{}n}\StringTok{"}\NormalTok{,}
                       \StringTok{"R² = "}\NormalTok{, }\FunctionTok{round}\NormalTok{(}\FunctionTok{summary}\NormalTok{(modelo\_c)}\SpecialCharTok{$}\NormalTok{r.squared, }\DecValTok{3}\NormalTok{)),}
       \AttributeTok{bty =} \StringTok{"n"}\NormalTok{, }\AttributeTok{text.col =} \StringTok{"darkblue"}\NormalTok{, }\AttributeTok{cex =} \FloatTok{0.9}\NormalTok{)}
\end{Highlighting}
\end{Shaded}

\begin{figure}[H]

{\centering \pandocbounded{\includegraphics[keepaspectratio]{chapters/semana09_files/figure-pdf/lst-s10-cod06-1.pdf}}

}

\caption{Regresión lineal Colesterol Total \textasciitilde{} IMC. Los
250 puntos representan pacientes sin tratamiento hipolipemiante. La
línea muestra la relación estimada: por cada kg/m² adicional de IMC, el
colesterol aumenta en promedio \textasciitilde5 mg/dL.}

\end{figure}%

\begin{tcolorbox}[enhanced jigsaw, breakable, titlerule=0mm, rightrule=.15mm, leftrule=.75mm, left=2mm, opacitybacktitle=0.6, colback=white, bottomtitle=1mm, colbacktitle=quarto-callout-tip-color!10!white, title=\textcolor{quarto-callout-tip-color}{\faLightbulb}\hspace{0.5em}{Tip}, toptitle=1mm, coltitle=black, opacityback=0, arc=.35mm, bottomrule=.15mm, toprule=.15mm, colframe=quarto-callout-tip-color-frame]

\textbf{Interpretación estadística:} En una muestra de 250 pacientes sin
hipolipemiantes, el modelo de regresión lineal mostró que por cada
unidad adicional de IMC (kg/m²), el colesterol total aumenta en promedio
5.0 mg/dL (b₁ = 5.0, EE = 0.41, t(248) = 12.2, p \textless{} 0.001, IC
95\%: {[}4.19, 5.81{]}). El R² = 0.391 indica que el IMC explica el
39.1\% de la variación en colesterol, estableciendo una relación lineal
fuerte y estadísticamente significativa. Desde una perspectiva clínica
de prevención cardiovascular, una reducción de 5 unidades en el IMC (ej.
pasar de obesidad a sobrepeso) se asociaría con una reducción esperada
de \textasciitilde25 mg/dL en el colesterol total, lo que podría reducir
el riesgo cardiovascular.

\end{tcolorbox}

\section{Resumen de Fórmulas Clave}\label{sec-semana09-resumen}

\begin{tcolorbox}[enhanced jigsaw, breakable, titlerule=0mm, rightrule=.15mm, leftrule=.75mm, left=2mm, opacitybacktitle=0.6, colback=white, bottomtitle=1mm, colbacktitle=quarto-callout-warning-color!10!white, title=\textcolor{quarto-callout-warning-color}{\faExclamationTriangle}\hspace{0.5em}{Fórmulas Esenciales de Regresión Lineal Simple}, toptitle=1mm, coltitle=black, opacityback=0, arc=.35mm, bottomrule=.15mm, toprule=.15mm, colframe=quarto-callout-warning-color-frame]

\textbf{Modelo}: \[Y_i = \beta_0 + \beta_1 X_i + U_i\]

\textbf{Estimadores MCO}:
\[b_1 = \frac{s_{xy}}{s_x^2}, \quad b_0 = \bar{y} - b_1 \bar{x}\]

\textbf{Descomposición de varianza}:
\[\text{SCT} = \text{SCE} + \text{SCR}\]

\textbf{Coeficiente de determinación}:
\[R^2 = 1 - \frac{\text{SCR}}{\text{SCT}} = r_{xy}^2\]

\textbf{Varianza del error estimada}:
\[s_u^2 = \frac{\sum_{i=1}^n \hat{u}_i^2}{n - 2}\]

\textbf{Errores estándar}:
\[\widehat{\text{ee}}(b_1) = \frac{s_u}{\sqrt{\sum (x_i - \bar{x})^2}}, \quad \widehat{\text{ee}}(b_0) = s_u\sqrt{\frac{1}{n} + \frac{\bar{x}^2}{\sum(x_i - \bar{x})^2}}\]

\textbf{Test t}:
\[t = \frac{b_j}{\widehat{\text{ee}}(b_j)} \sim t_{n-2}\]

\textbf{Intervalo de confianza para \(\beta_j\)}:
\[[b_j \pm t_{1-\alpha/2, n-2} \cdot \widehat{\text{ee}}(b_j)]\]

\end{tcolorbox}

\section{Ejercicios}\label{sec-semana09-ejercicios}

\textbf{Ejercicio 1}: En un estudio sobre prediabetes en 40 pacientes
del centro de salud, se investiga la relación entre el \textbf{índice de
masa corporal} (\(X\), kg/m²) y la \textbf{glucosa en ayunas} (\(Y\),
mg/dL). Se obtienen las siguientes estadísticas:

\begin{itemize}
\tightlist
\item
  \(\bar{x} = 28.5\) kg/m², \(\bar{y} = 105.5\) mg/dL
\item
  \(s_x^2 = 9.0\), \(s_y^2 = 144.0\)
\item
  \(s_{xy} = 27.0\)
\end{itemize}

\begin{enumerate}
\def\labelenumi{\alph{enumi})}
\tightlist
\item
  Calcula los estimadores MCO \(b_0\) y \(b_1\)
\item
  Interpreta clínicamente el coeficiente \(b_1\)
\item
  Calcula \(R^2\) e interpreta qué proporción de la variabilidad
  glucémica explica el IMC
\item
  Un paciente nuevo tiene un IMC de 32 kg/m². ¿Cuál es su glucosa en
  ayunas predicha? ¿Podría considerarse en rango de prediabetes (glucosa
  ≥ 100 mg/dL)?
\end{enumerate}

\begin{center}\rule{0.5\linewidth}{0.5pt}\end{center}

\textbf{Ejercicio 2}: Un estudio de función pulmonar en una consulta de
neumología evalúa el efecto del tabaquismo sobre el FEV1 (volumen
espiratorio forzado en el primer segundo, expresado como \% del
predicho). Con 60 pacientes fumadores se obtiene:

\begin{itemize}
\tightlist
\item
  \(\hat{\beta}_0 = 98.0\%\), \(\hat{\beta}_1 = -0.80\%\) por
  paquete-año
\item
  \(\widehat{\text{ee}}(b_1) = 0.12\)
\item
  \(s_u = 7.5\%\), \(\sum (x_i - \bar{x})^2 = 3906.25\)
\end{itemize}

\begin{enumerate}
\def\labelenumi{\alph{enumi})}
\tightlist
\item
  Realiza el test de hipótesis \(H_0: \beta_1 = 0\) con
  \(\alpha = 0.05\). Interpreta el resultado
\item
  Construye un intervalo de confianza al 95\% para \(\beta_1\)
\item
  ¿Qué significa clínicamente que rechacemos \(H_0\)?
\item
  ¿Es razonable el signo negativo de \(\hat{\beta}_1\)?
\end{enumerate}

\begin{center}\rule{0.5\linewidth}{0.5pt}\end{center}

\textbf{Ejercicio 3}: Para los datos del estudio de colesterol e IMC
(sección 10.12, \(n = 250\)):

\begin{enumerate}
\def\labelenumi{\alph{enumi})}
\tightlist
\item
  Calcula el intervalo de confianza al 95\% para el colesterol medio en
  pacientes con IMC = 30 kg/m² (obesidad tipo I)
\item
  ¿Es estadísticamente significativo el efecto del IMC sobre el
  colesterol? Justifica usando el estadístico \(t\)
\item
  Un paciente tiene IMC = 32 kg/m² y colesterol = 260 mg/dL. Calcula su
  residuo. ¿Podría ser un paciente con dislipemia familiar no
  relacionada con el IMC?
\end{enumerate}

\begin{center}\rule{0.5\linewidth}{0.5pt}\end{center}

\textbf{Ejercicio 4}: En un ensayo clínico con 30 pacientes febriles, se
estudia la relación entre la \textbf{dosis de paracetamol} (\(X\),
mg/kg) y la \textbf{reducción de temperatura corporal a las 2 horas}
(\(Y\), °C). El modelo estimado es:

\[\hat{y}_i = 0.50 + 0.30 x_i, \quad R^2 = 0.64\]

\begin{enumerate}
\def\labelenumi{\alph{enumi})}
\tightlist
\item
  ¿Qué proporción de la reducción de temperatura es explicada por la
  dosis?
\item
  Si la dosis media fue \(\bar{x} = 10\) mg/kg, ¿cuál fue la reducción
  media de temperatura?
\item
  Por cada mg/kg adicional de paracetamol, ¿cuánto desciende la
  temperatura en promedio?
\item
  ¿Qué factor explicaría el 36\% de variabilidad no captada por la
  dosis?
\end{enumerate}

\begin{center}\rule{0.5\linewidth}{0.5pt}\end{center}

\textbf{Ejercicio 5}: Un clínico interpreta los gráficos de diagnóstico
de un modelo que relaciona la \textbf{supervivencia} (meses) con el
\textbf{tamaño tumoral} (mm). Describe qué patrón en cada gráfico
indicaría:

\begin{enumerate}
\def\labelenumi{\alph{enumi})}
\tightlist
\item
  Que la relación no es lineal (posiblemente necesita
  log-transformación)
\item
  Heteroscedasticidad (varianza de los residuos crece con el tamaño
  tumoral)
\item
  Violación de la normalidad de los residuos
\item
  Presencia de un paciente con diagnóstico tardío que actúa como punto
  influyente
\end{enumerate}

\begin{center}\rule{0.5\linewidth}{0.5pt}\end{center}

\textbf{Ejercicio 6}: En un programa de cribado de cáncer de próstata,
se estudia la asociación entre el \textbf{PSA sérico} (ng/mL) y la
\textbf{edad} (años). En una primera muestra: \(r = 0.58\). En una
segunda muestra: \(r = 0.45\):

\begin{enumerate}
\def\labelenumi{\alph{enumi})}
\tightlist
\item
  Calcula y compara los \(R^2\) de ambas muestras. ¿Qué muestra muestra
  mayor capacidad predictiva?
\item
  ¿La relación PSA-edad es más fuerte o más débil en la segunda muestra?
\item
  Si las desviaciones estándar permanecen iguales, ¿cómo cambia el
  coeficiente de regresión \(b_1\) entre muestras?
\item
  ¿Qué factores clínicos podrían explicar una asociación más débil
  (menor \(r\)) en la segunda muestra?
\end{enumerate}

\section{Respuestas a los Ejercicios}\label{sec-semana09-respuestas}

\textbf{Ejercicio 1:} - a) \(b_1 = s_{xy}/s_x^2 = 27.0/9.0 = 3.0\) mg/dL
por kg/m²; \(b_0 = 105.5 - 3.0 \times 28.5 = 20.0\) mg/dL - b) Por cada
kg/m² adicional de IMC, la glucosa en ayunas aumenta en promedio 3
mg/dL, manteniendo el resto constante. Con un IMC de 30 vs 25, se
esperan 15 mg/dL más de glucosa. - c)
\(r = s_{xy}/(s_x \cdot s_y) = 27.0/(3.0 \times 12.0) = 0.75\);
\(R^2 = 0.5625\). El IMC explica el 56.25\% de la variabilidad en la
glucosa en ayunas. - d) \(\hat{Y}(32) = 20.0 + 3.0 \times 32 = 116\)
mg/dL. Sí, está en rango de prediabetes (≥100 mg/dL); de hecho, supera
el umbral de diabetes (≥126 mg/dL), lo que indica riesgo elevado.

\textbf{Ejercicio 2:} - a) \(t = -0.80/0.12 = -6.67\); \(gl = 58\);
\(t_{0.975, 58} \approx 2.00\). Como \(|{-6.67}| > 2.00\)
(\(p < 0.001\)), rechazamos \(H_0\). El tabaquismo tiene un efecto
estadísticamente significativo sobre el FEV1. - b) IC 95\%:
\([-0.80 \pm 2.00 \times 0.12] = [-1.04, -0.56]\) \% por paquete-año. -
c) Existe evidencia de que el tabaquismo reduce el FEV1. Cada
paquete-año adicional se asocia con una reducción media de 0.80\% del
FEV1 predicho. - d) Sí, es razonable. El tabaquismo provoca inflamación
crónica y destrucción del parénquima pulmonar, reduciendo la función
respiratoria. El signo negativo es clínicamente correcto.

\textbf{Ejercicio 3:} - a) Usando
\(\hat{Y}(30) = 78.0 + 5.0 \times 30 = 228\) mg/dL. Con
\(s_u \approx 28.7\) mg/dL, el IC al 95\% requiere calcular
\(s_u\sqrt{1/n + (30 - 27.2)^2/\text{SCT}_x}\) --- los datos dan un
intervalo aproximado de \([221, 235]\) mg/dL. - b) Sí,
\(t = 12.2 \gg t_{0.975, 248} \approx 1.970\) (\(p < 0.001\)). El efecto
del IMC sobre el colesterol es altamente significativo. - c) Residuo:
\(\hat{u} = 260 - (78.0 + 5.0 \times 32) = 260 - 238 = 22\) mg/dL. Un
residuo de +22 mg/dL es moderado (\textless{} 2 desviaciones estándar).
Es compatible con dislipemia familiar leve, aunque no constituye un
outlier extremo.

\textbf{Ejercicio 4:} - a) \(R^2 = 0.64\); el 64\% de la variabilidad en
la reducción de temperatura es explicado por la dosis de paracetamol. -
b) \(\hat{Y}(\bar{x}) = 0.50 + 0.30 \times 10 = 3.5\) °C de reducción
media. - c) Por cada mg/kg adicional de paracetamol, la temperatura
desciende en promedio 0.30 °C. - d) El 36\% restante: estado inmune del
paciente, edad, causa de la fiebre (viral vs.~bacteriana), hidratación,
medicación concomitante, hora de administración.

\textbf{Ejercicio 5:} - a) Patrón curvado (en forma de U o de arco) en
el gráfico de residuos vs.~valores ajustados. - b) Patrón ``embudo''
(varianza de residuos aumenta con los valores predichos). Sugiere
log-transformación de la variable respuesta. - c) Desviaciones
sistemáticas de la línea 45° en el Q-Q plot (especialmente en las
colas). - d) Punto con leverage elevado y residuo grande en el gráfico
de Residuos vs.~Leverage (esquina superior o inferior derecha). Requiere
investigación clínica del caso.

\textbf{Ejercicio 6:} - a) \(R^2_1 = 0.58^2 = 0.336\) (33.6\%);
\(R^2_2 = 0.45^2 = 0.2025\) (20.25\%). La primera muestra tiene mayor
capacidad predictiva del PSA a partir de la edad. - b) La relación
PSA-edad es más débil en la segunda muestra (\(r\) disminuye de 0.58 a
0.45). - c) \(b_1 = r \cdot s_y/s_x\). Si las desviaciones estándar no
cambian, \(b_1\) disminuye proporcionalmente:
\(b_{1,2}/b_{1,1} = 0.45/0.58 = 0.776\) (un 22.4\% menor). - d) Posibles
razones: rango de edades más estrecho en la segunda muestra, mayor
proporción de pacientes con prostatitis o hipertrofia benigna de
próstata (que elevan el PSA independientemente de la edad), diferencias
en la técnica de laboratorio.

\begin{center}\rule{0.5\linewidth}{0.5pt}\end{center}

\begin{tcolorbox}[enhanced jigsaw, breakable, titlerule=0mm, rightrule=.15mm, leftrule=.75mm, left=2mm, opacitybacktitle=0.6, colback=white, bottomtitle=1mm, colbacktitle=quarto-callout-tip-color!10!white, title=\textcolor{quarto-callout-tip-color}{\faLightbulb}\hspace{0.5em}{Métodos Avanzados}, toptitle=1mm, coltitle=black, opacityback=0, arc=.35mm, bottomrule=.15mm, toprule=.15mm, colframe=quarto-callout-tip-color-frame]

Para ampliar los contenidos de este capítulo con técnicas estadísticas
avanzadas, visita:

\textbf{\href{https://migariane.github.io/Bioestadistica}{→
Bioestadística Avanzada --- M.A.~Luque Fernández}}

\end{tcolorbox}

\chapter{Semana 10 --- Regresión Lineal Múltiple}\label{sec-semana10}

En esta semana extendemos la regresión lineal simple a múltiples
variables predictoras. En medicina, raramente un único factor explica un
desenlace clínico: la presión arterial depende de la edad, el IMC, el
sexo y el tabaquismo simultáneamente; la supervivencia oncológica
depende del estadio, el tratamiento y la edad al diagnóstico. La
regresión lineal múltiple nos permite cuantificar el efecto de cada
variable \textbf{controlando por las demás}.

\section{El Modelo de Regresión Lineal Múltiple}\label{sec-semana10-s1}

La regresión múltiple permite modelar relaciones complejas mediante
varias variables predictoras simultáneamente.

\begin{tcolorbox}[enhanced jigsaw, breakable, titlerule=0mm, rightrule=.15mm, leftrule=.75mm, left=2mm, opacitybacktitle=0.6, colback=white, bottomtitle=1mm, colbacktitle=quarto-callout-note-color!10!white, title=\textcolor{quarto-callout-note-color}{\faInfo}\hspace{0.5em}{Definición: Modelo de Regresión Lineal Múltiple}, toptitle=1mm, coltitle=black, opacityback=0, arc=.35mm, bottomrule=.15mm, toprule=.15mm, colframe=quarto-callout-note-color-frame]

El modelo de regresión lineal múltiple con \(p\) variables predictoras
es:

\[Y_i = \beta_0 + \beta_1 X_{1i} + \beta_2 X_{2i} + \cdots + \beta_p X_{pi} + U_i\]

para \(i = 1, 2, \ldots, n\) observaciones, donde:

\begin{itemize}
\tightlist
\item
  \(Y_i\) es la variable respuesta (dependiente)
\item
  \(X_{ji}\) son las variables predictoras (independientes) para
  \(j = 1, \ldots, p\)
\item
  \(\beta_0\) es la ordenada al origen (intercepto)
\item
  \(\beta_j\) es el \textbf{coeficiente de regresión parcial} de \(X_j\)
  --- representa el cambio esperado en \(Y\) ante un aumento unitario de
  \(X_j\), \textbf{manteniendo todas las demás variables constantes}
\item
  \(U_i\) es el término de error aleatorio con \(E(U_i) = 0\) y
  \(\text{Var}(U_i) = \sigma_U^2\)
\end{itemize}

\end{tcolorbox}

\textbf{Interpretación de coeficientes parciales:} El coeficiente
\(\beta_j\) mide el efecto de \(X_j\) sobre \(Y\) después de
\textbf{controlar} (ajustar) por todas las otras variables predictoras
en el modelo. Esto es fundamental en la investigación médica para
separar el efecto de cada factor de riesgo del de los demás.

\begin{tcolorbox}[enhanced jigsaw, breakable, titlerule=0mm, rightrule=.15mm, leftrule=.75mm, left=2mm, opacitybacktitle=0.6, colback=white, bottomtitle=1mm, colbacktitle=quarto-callout-tip-color!10!white, title=\textcolor{quarto-callout-tip-color}{\faLightbulb}\hspace{0.5em}{Ejemplo: Presión Arterial Sistólica --- Extensión del modelo simple}, toptitle=1mm, coltitle=black, opacityback=0, arc=.35mm, bottomrule=.15mm, toprule=.15mm, colframe=quarto-callout-tip-color-frame]

Continuando con el estudio de hipertensión de la semana anterior (120
pacientes adultos), ahora incluimos \textbf{dos predictores} para
explicar la presión arterial sistólica (PAS):

\begin{itemize}
\tightlist
\item
  \(X_1\) = edad (años)
\item
  \(X_2\) = índice de masa corporal --- IMC (kg/m²)
\end{itemize}

El modelo es:
\[\text{PAS}_i = \beta_0 + \beta_1 \cdot \text{Edad}_i + \beta_2 \cdot \text{IMC}_i + U_i\]

\begin{itemize}
\tightlist
\item
  \(\beta_1\) mide el cambio esperado en PAS por cada año adicional de
  edad, \textbf{manteniendo el IMC fijo}
\item
  \(\beta_2\) mide el cambio esperado en PAS por cada kg/m² adicional de
  IMC, \textbf{manteniendo la edad fija}
\end{itemize}

Interpretación hipotética: si \(\hat{\beta}_1 = 0.60\) mmHg/año y
\(\hat{\beta}_2 = 1.80\) mmHg/(kg/m²), un paciente que envejece 10 años
(con mismo IMC) tendrá en promedio 6 mmHg más de PAS; un paciente con 5
kg/m² más de IMC (con misma edad) tendrá en promedio 9 mmHg más de PAS.

\end{tcolorbox}

\section{Formulación Matricial del Modelo}\label{sec-semana10-s2}

Con múltiples variables predictoras, la notación escalar se vuelve
engorrosa. La formulación matricial simplifica significativamente el
tratamiento matemático.

\begin{tcolorbox}[enhanced jigsaw, breakable, titlerule=0mm, rightrule=.15mm, leftrule=.75mm, left=2mm, opacitybacktitle=0.6, colback=white, bottomtitle=1mm, colbacktitle=quarto-callout-note-color!10!white, title=\textcolor{quarto-callout-note-color}{\faInfo}\hspace{0.5em}{Definición: Forma Matricial del Modelo}, toptitle=1mm, coltitle=black, opacityback=0, arc=.35mm, bottomrule=.15mm, toprule=.15mm, colframe=quarto-callout-note-color-frame]

El modelo puede escribirse compactamente como:

\[\mathbf{Y} = \mathbf{X}\boldsymbol{\beta} + \mathbf{U}\]

donde:

\begin{itemize}
\item
  \(\mathbf{Y} = \begin{bmatrix} Y_1 \\ Y_2 \\ \vdots \\ Y_n \end{bmatrix}\)
  es el vector \(n \times 1\) de respuestas (ej. PAS de cada paciente)
\item
  \(\boldsymbol{\beta} = \begin{bmatrix} \beta_0 \\ \beta_1 \\ \vdots \\ \beta_p \end{bmatrix}\)
  es el vector \((p+1) \times 1\) de parámetros
\item
  \(\mathbf{X} = \begin{bmatrix} 1 & x_{11} & x_{21} & \cdots & x_{p1} \\ 1 & x_{12} & x_{22} & \cdots & x_{p2} \\ \vdots & \vdots & \vdots & \ddots & \vdots \\ 1 & x_{1n} & x_{2n} & \cdots & x_{pn} \end{bmatrix}\)
  es la \textbf{matriz de diseño} \(n \times (p+1)\)
\item
  \(\mathbf{U} = \begin{bmatrix} U_1 \\ U_2 \\ \vdots \\ U_n \end{bmatrix}\)
  es el vector \(n \times 1\) de errores
\end{itemize}

\textbf{Nota:} La primera columna de \(\mathbf{X}\) contiene unos (para
el intercepto).

\end{tcolorbox}

\textbf{Ejemplo}: Para el modelo PAS \textasciitilde{} Edad + IMC con
\(n = 120\) pacientes, \(\mathbf{X}\) es una matriz \(120 \times 3\):
columna de unos, columna de edades y columna de IMC de cada paciente.

\section{Estimación de Parámetros: Mínimos Cuadrados
Ordinarios}\label{sec-semana10-s3}

\subsection{Derivación del Estimador
OLS}\label{derivaciuxf3n-del-estimador-ols}

El método de mínimos cuadrados ordinarios (OLS, \emph{Ordinary Least
Squares}) minimiza la suma de cuadrados residual:

\[\text{RSS}(\boldsymbol{\beta}) = \sum_{i=1}^n (Y_i - \hat{Y}_i)^2 = (\mathbf{Y} - \mathbf{X}\boldsymbol{\beta})^T(\mathbf{Y} - \mathbf{X}\boldsymbol{\beta}) = \|\mathbf{Y} - \mathbf{X}\boldsymbol{\beta}\|^2\]

Para encontrar el mínimo, tomamos la derivada con respecto a
\(\boldsymbol{\beta}\) e igualamos a cero:

\[\frac{\partial \text{RSS}}{\partial \boldsymbol{\beta}} = -2\mathbf{X}^T(\mathbf{Y} - \mathbf{X}\boldsymbol{\beta}) = 0\]

\begin{tcolorbox}[enhanced jigsaw, breakable, titlerule=0mm, rightrule=.15mm, leftrule=.75mm, left=2mm, opacitybacktitle=0.6, colback=white, bottomtitle=1mm, colbacktitle=quarto-callout-warning-color!10!white, title=\textcolor{quarto-callout-warning-color}{\faExclamationTriangle}\hspace{0.5em}{Resultado Importante: Ecuaciones Normales y Estimador OLS}, toptitle=1mm, coltitle=black, opacityback=0, arc=.35mm, bottomrule=.15mm, toprule=.15mm, colframe=quarto-callout-warning-color-frame]

Las \textbf{ecuaciones normales} son:
\[\mathbf{X}^T\mathbf{X}\hat{\boldsymbol{\beta}} = \mathbf{X}^T\mathbf{Y}\]

Si la matriz \(\mathbf{X}^T\mathbf{X}\) es invertible (no singular), el
\textbf{estimador OLS} es:

\[\hat{\boldsymbol{\beta}} = (\mathbf{X}^T\mathbf{X})^{-1}\mathbf{X}^T\mathbf{Y}\]

donde
\(\hat{\boldsymbol{\beta}} = \begin{bmatrix} \hat{\beta}_0 \\ \hat{\beta}_1 \\ \vdots \\ \hat{\beta}_p \end{bmatrix}\)

\end{tcolorbox}

\subsection{Condiciones para
Invertibilidad}\label{condiciones-para-invertibilidad}

La matriz \(\mathbf{X}^T\mathbf{X}\) es singular (no invertible) si:

\begin{itemize}
\tightlist
\item
  Las variables predictoras son linealmente dependientes (colinealidad
  exacta)
\item
  Una variable es combinación lineal exacta de otras
\end{itemize}

\textbf{Ejemplo médico:} Si incluimos tanto el peso (kg) como el IMC
(kg/m²) y la altura al cuadrado, hay colinealidad exacta porque IMC =
peso/altura². R eliminará automáticamente una de las variables, pero
esto debe evitarse en el diseño del estudio.

\section{Matriz de Varianza-Covarianza de los
Coeficientes}\label{sec-semana10-s4}

\begin{tcolorbox}[enhanced jigsaw, breakable, titlerule=0mm, rightrule=.15mm, leftrule=.75mm, left=2mm, opacitybacktitle=0.6, colback=white, bottomtitle=1mm, colbacktitle=quarto-callout-note-color!10!white, title=\textcolor{quarto-callout-note-color}{\faInfo}\hspace{0.5em}{Definición: Matriz de Varianza-Covarianza}, toptitle=1mm, coltitle=black, opacityback=0, arc=.35mm, bottomrule=.15mm, toprule=.15mm, colframe=quarto-callout-note-color-frame]

Bajo los supuestos del modelo de regresión lineal múltiple, la matriz de
varianza-covarianza de los coeficientes estimados es:

\[\text{Var}(\hat{\boldsymbol{\beta}}) = \sigma_U^2(\mathbf{X}^T\mathbf{X})^{-1}\]

donde:

\begin{itemize}
\tightlist
\item
  Los elementos diagonales
  \([\sigma_U^2(\mathbf{X}^T\mathbf{X})^{-1}]_{jj}\) son las varianzas
  \(\text{Var}(\hat{\beta}_j)\)
\item
  Los elementos fuera de la diagonal son covarianzas
  \(\text{Cov}(\hat{\beta}_j, \hat{\beta}_k)\)
\item
  \(\sigma_U^2\) es la varianza desconocida del error
\end{itemize}

\end{tcolorbox}

\subsection{\texorpdfstring{Estimador Insesgado de
\(\sigma_U^2\)}{Estimador Insesgado de \textbackslash sigma\_U\^{}2}}\label{estimador-insesgado-de-sigma_u2}

Como \(\sigma_U^2\) es desconocida, la estimamos con:

\[s^2 = \hat{\sigma}_U^2 = \frac{\text{RSS}}{n-p-1} = \frac{\sum_{i=1}^n \hat{u}_i^2}{n-p-1}\]

donde \(\hat{u}_i = Y_i - \hat{Y}_i\) son los residuos estimados.

\begin{tcolorbox}[enhanced jigsaw, breakable, titlerule=0mm, rightrule=.15mm, leftrule=.75mm, left=2mm, opacitybacktitle=0.6, colback=white, bottomtitle=1mm, colbacktitle=quarto-callout-warning-color!10!white, title=\textcolor{quarto-callout-warning-color}{\faExclamationTriangle}\hspace{0.5em}{Resultado Importante}, toptitle=1mm, coltitle=black, opacityback=0, arc=.35mm, bottomrule=.15mm, toprule=.15mm, colframe=quarto-callout-warning-color-frame]

\[s^2 = \frac{\text{RSS}}{n-p-1}\]

es un estimador insesgado de \(\sigma_U^2\). Note que dividimos por
\(n-p-1\) (grados de libertad), no por \(n\). En el modelo PAS
\textasciitilde{} Edad + IMC + Sexo con \(n = 120\) pacientes y
\(p = 3\) predictores, los grados de libertad son \(120 - 3 - 1 = 116\).

\end{tcolorbox}

\section{Variables Indicadoras (Dummy)}\label{sec-semana10-s5}

Cuando tenemos variables predictoras categóricas (sexo, grupo de
tratamiento, estadio tumoral), usamos \textbf{variables indicadoras}
(0/1) para incorporarlas en el modelo lineal.

\begin{tcolorbox}[enhanced jigsaw, breakable, titlerule=0mm, rightrule=.15mm, leftrule=.75mm, left=2mm, opacitybacktitle=0.6, colback=white, bottomtitle=1mm, colbacktitle=quarto-callout-note-color!10!white, title=\textcolor{quarto-callout-note-color}{\faInfo}\hspace{0.5em}{Definición: Variables Indicadoras}, toptitle=1mm, coltitle=black, opacityback=0, arc=.35mm, bottomrule=.15mm, toprule=.15mm, colframe=quarto-callout-note-color-frame]

Una variable indicadora para una categoría es:

\[X_j = \begin{cases} 1 & \text{si la observación pertenece a la categoría } j \\ 0 & \text{en caso contrario} \end{cases}\]

Para una variable categórica con \(m\) niveles, utilizamos exactamente
\(m-1\) variables indicadoras (dejamos una categoría como
\textbf{referencia} o base).

\end{tcolorbox}

\begin{tcolorbox}[enhanced jigsaw, breakable, titlerule=0mm, rightrule=.15mm, leftrule=.75mm, left=2mm, opacitybacktitle=0.6, colback=white, bottomtitle=1mm, colbacktitle=quarto-callout-tip-color!10!white, title=\textcolor{quarto-callout-tip-color}{\faLightbulb}\hspace{0.5em}{Ejemplo: Peso al Nacer según Edad Gestacional y Sexo}, toptitle=1mm, coltitle=black, opacityback=0, arc=.35mm, bottomrule=.15mm, toprule=.15mm, colframe=quarto-callout-tip-color-frame]

Un neonatólogo estudia la relación entre edad gestacional, sexo y peso
al nacer en 189 recién nacidos del Hospital Materno-Infantil de Granada.
Define:

\begin{itemize}
\tightlist
\item
  \(Y_i\) = peso al nacer (gramos)
\item
  \(X_{1i}\) = edad gestacional (semanas)
\item
  \(X_{2i}\) = variable indicadora: 1 si el bebé es niña, 0 si es niño
\end{itemize}

El modelo es:
\[\text{Peso}_i = \beta_0 + \beta_1 \cdot \text{Gestación}_i + \beta_2 \cdot \text{Niña}_i + U_i\]

Estimando el modelo, se obtiene \(\hat{\beta}_0 = -1610.28\),
\(\hat{\beta}_1 = 120.89\), \(\hat{\beta}_2 = -163.04\).

\textbf{Interpretación:}

\begin{itemize}
\tightlist
\item
  Para niños (\(X_2 = 0\)):
  \(\widehat{\text{Peso}} = -1610.28 + 120.89 \cdot \text{Gestación}\)
\item
  Para niñas (\(X_2 = 1\)):
  \(\widehat{\text{Peso}} = -1773.32 + 120.89 \cdot \text{Gestación}\)
\end{itemize}

La diferencia de \(-163.04\) gramos entre niñas y niños es el efecto del
sexo \textbf{controlando por la edad gestacional}. Las niñas pesan en
promedio 163 g menos que los niños a igual edad gestacional. Por cada
semana adicional de gestación, el peso aumenta en promedio 120.89 g en
ambos sexos.

\textbf{Predicción}: Un niño de 38 semanas:
\(\hat{P} = -1610.28 + 120.89 \times 38 = 2983\) g. Una niña de 38
semanas: \(\hat{P} = 2983 - 163 = 2820\) g.

\end{tcolorbox}

\section{\texorpdfstring{Bondad de Ajuste: \(R^2\) y \(R^2\)
Ajustado}{Bondad de Ajuste: R\^{}2 y R\^{}2 Ajustado}}\label{sec-semana10-s6}

\subsection{Coeficiente de Determinación
Múltiple}\label{coeficiente-de-determinaciuxf3n-muxfaltiple}

En regresión múltiple, el coeficiente \(R^2\) se define igual que en
regresión simple:

\[R^2 = 1 - \frac{\text{RSS}}{\text{SST}} = \frac{\text{ESS}}{\text{SST}}\]

donde:

\begin{itemize}
\tightlist
\item
  \(\text{RSS} = \sum_{i=1}^n (Y_i - \hat{Y}_i)^2\) (suma de cuadrados
  residual)
\item
  \(\text{SST} = \sum_{i=1}^n (Y_i - \bar{Y})^2\) (suma de cuadrados
  total)
\item
  \(\text{ESS} = \sum_{i=1}^n (\hat{Y}_i - \bar{Y})^2\) (suma de
  cuadrados explicada)
\end{itemize}

\textbf{Problema clínico:} Si al modelo PAS \textasciitilde{} Edad
añadimos variables irrelevantes (número de calzado del paciente),
\(R^2\) aumentará o se mantendrá, aunque el número de calzado no prediga
la PAS. Este problema nos lleva al \(R^2\) ajustado.

\subsection{\texorpdfstring{\(R^2\)
Ajustado}{R\^{}2 Ajustado}}\label{r2-ajustado}

\begin{tcolorbox}[enhanced jigsaw, breakable, titlerule=0mm, rightrule=.15mm, leftrule=.75mm, left=2mm, opacitybacktitle=0.6, colback=white, bottomtitle=1mm, colbacktitle=quarto-callout-warning-color!10!white, title=\textcolor{quarto-callout-warning-color}{\faExclamationTriangle}\hspace{0.5em}{Resultado Importante: \(R^2\) Ajustado}, toptitle=1mm, coltitle=black, opacityback=0, arc=.35mm, bottomrule=.15mm, toprule=.15mm, colframe=quarto-callout-warning-color-frame]

El coeficiente de determinación ajustado penaliza por el número de
parámetros:

\[R^2_{\text{adj}} = 1 - \frac{\text{RSS}/(n-p-1)}{\text{SST}/(n-1)} = 1 - (1-R^2)\frac{n-1}{n-p-1}\]

Propiedades:

\begin{itemize}
\tightlist
\item
  \(R^2_{\text{adj}} \leq R^2\) siempre
\item
  \(R^2_{\text{adj}}\) \textbf{puede disminuir} si añadimos una variable
  que no mejora el modelo
\item
  Es más útil que \(R^2\) para comparar modelos con diferente número de
  predictores
\end{itemize}

\textbf{Ejemplo médico}: Modelo 1 (PAS \textasciitilde{} Edad):
\(R^2 = 0.504\), \(R^2_{\text{adj}} = 0.499\). Modelo 2 (PAS
\textasciitilde{} Edad + IMC): \(R^2 = 0.561\),
\(R^2_{\text{adj}} = 0.553\). El \(R^2_{\text{adj}}\) aumenta al añadir
el IMC, lo que indica que esta variable aporta información real más allá
de la edad.

\end{tcolorbox}

\section{Criterios de Información para Selección de
Modelos}\label{sec-semana10-s7}

Además de \(R^2_{\text{adj}}\), existen criterios de información que
penalizan la complejidad del modelo:

\begin{tcolorbox}[enhanced jigsaw, breakable, titlerule=0mm, rightrule=.15mm, leftrule=.75mm, left=2mm, opacitybacktitle=0.6, colback=white, bottomtitle=1mm, colbacktitle=quarto-callout-warning-color!10!white, title=\textcolor{quarto-callout-warning-color}{\faExclamationTriangle}\hspace{0.5em}{Definición: AIC y BIC}, toptitle=1mm, coltitle=black, opacityback=0, arc=.35mm, bottomrule=.15mm, toprule=.15mm, colframe=quarto-callout-warning-color-frame]

\textbf{Criterio de Información de Akaike (AIC):}
\[\text{AIC} = -2 \log L + 2(p+1)\]

\textbf{Criterio de Información Bayesiano (BIC):}
\[\text{BIC} = -2 \log L + (p+1)\log(n)\]

donde \(L\) es la función de verosimilitud maximizada y \(p+1\) es el
número de parámetros (incluyendo \(\sigma_U^2\)).

Para modelos normales: \[\text{AIC} = n \log(\text{RSS}/n) + 2(p+1)\]
\[\text{BIC} = n \log(\text{RSS}/n) + (p+1)\log(n)\]

\textbf{Regla de decisión:} Seleccionar el modelo con \textbf{menor} AIC
o BIC.

\begin{itemize}
\tightlist
\item
  AIC usa penalización moderada (factor 2): útil en contextos
  predictivos
\item
  BIC usa penalización más fuerte (factor \(\log n\)): favorece modelos
  más parsimoniosos, preferible en investigación clínica confirmatoria
\end{itemize}

\textbf{Ejemplo}: Para comparar modelos de predicción de HbA1c con
distintas combinaciones de variables (IMC, edad, actividad física,
tabaquismo), el BIC es más apropiado en un estudio con \(n = 200\)
porque \(\log(200) = 5.30 > 2\), penalizando más fuertemente la
complejidad.

\end{tcolorbox}

\section{Pruebas de Hipótesis para Coeficientes
Individuales}\label{sec-semana10-s8}

\subsection{\texorpdfstring{Test t para
\(\beta_j\)}{Test t para \textbackslash beta\_j}}\label{test-t-para-beta_j}

\begin{tcolorbox}[enhanced jigsaw, breakable, titlerule=0mm, rightrule=.15mm, leftrule=.75mm, left=2mm, opacitybacktitle=0.6, colback=white, bottomtitle=1mm, colbacktitle=quarto-callout-note-color!10!white, title=\textcolor{quarto-callout-note-color}{\faInfo}\hspace{0.5em}{Supuestos del Modelo}, toptitle=1mm, coltitle=black, opacityback=0, arc=.35mm, bottomrule=.15mm, toprule=.15mm, colframe=quarto-callout-note-color-frame]

Para hacer inferencia estadística, asumimos:

\begin{itemize}
\tightlist
\item
  \(E(U_i) = 0\)
\item
  \(\text{Var}(U_i) = \sigma_U^2\) (homocedasticidad)
\item
  \(\text{Cov}(U_i, U_j) = 0\) para \(i \neq j\) (independencia)
\item
  \(U_i \sim N(0, \sigma_U^2)\) (normalidad)
\end{itemize}

\end{tcolorbox}

Bajo estos supuestos, el estimador \(\hat{\beta}_j\) sigue una
distribución normal:

\[\hat{\beta}_j \sim N(\beta_j, \text{Var}(\hat{\beta}_j))\]

donde
\(\text{Var}(\hat{\beta}_j) = \sigma_U^2 [(\mathbf{X}^T\mathbf{X})^{-1}]_{jj}\).

El estadístico t para probar \(H_0: \beta_j = 0\) es:

\begin{tcolorbox}[enhanced jigsaw, breakable, titlerule=0mm, rightrule=.15mm, leftrule=.75mm, left=2mm, opacitybacktitle=0.6, colback=white, bottomtitle=1mm, colbacktitle=quarto-callout-warning-color!10!white, title=\textcolor{quarto-callout-warning-color}{\faExclamationTriangle}\hspace{0.5em}{Resultado Importante: Estadístico t}, toptitle=1mm, coltitle=black, opacityback=0, arc=.35mm, bottomrule=.15mm, toprule=.15mm, colframe=quarto-callout-warning-color-frame]

\[t_j = \frac{\hat{\beta}_j}{\hat{\text{SE}}(\hat{\beta}_j)} \sim t_{n-p-1}\]

donde
\(\hat{\text{SE}}(\hat{\beta}_j) = \sqrt{s^2 [(\mathbf{X}^T\mathbf{X})^{-1}]_{jj}}\)
es el error estándar estimado.

\textbf{Procedimiento de prueba:}

\begin{enumerate}
\def\labelenumi{\arabic{enumi}.}
\tightlist
\item
  \textbf{Hipótesis:} \(H_0: \beta_j = 0\) vs.~\(H_1: \beta_j \neq 0\)
\item
  \textbf{Estadístico:} Calcular \(t_j\) como arriba
\item
  \textbf{Valor p:} \(p\text{-valor} = 2 \cdot P(t_{n-p-1} > |t_j|)\)
\item
  \textbf{Decisión:} Rechazar \(H_0\) si
  \(|t_j| > t_{1-\alpha/2, n-p-1}\) o si \(p\text{-valor} < \alpha\)
\end{enumerate}

\textbf{Interpretación clínica clave:} Rechazar \(H_0: \beta_j = 0\)
significa que la variable \(X_j\) aporta información significativa sobre
\(Y\) \textbf{después de controlar por todos los demás predictores del
modelo}.

\end{tcolorbox}

\section{Prueba F para Comparación de Modelos
Anidados}\label{sec-semana10-s9}

Cuando tenemos dos modelos donde uno está \textbf{anidado} dentro del
otro (sus variables predictoras son un subconjunto), podemos usar una
prueba F para comparar su bondad de ajuste.

\begin{tcolorbox}[enhanced jigsaw, breakable, titlerule=0mm, rightrule=.15mm, leftrule=.75mm, left=2mm, opacitybacktitle=0.6, colback=white, bottomtitle=1mm, colbacktitle=quarto-callout-note-color!10!white, title=\textcolor{quarto-callout-note-color}{\faInfo}\hspace{0.5em}{Definición: Modelos Anidados}, toptitle=1mm, coltitle=black, opacityback=0, arc=.35mm, bottomrule=.15mm, toprule=.15mm, colframe=quarto-callout-note-color-frame]

Decimos que el Modelo 1 es anidado en el Modelo 2 si los predictores del
Modelo 1 son un subconjunto de los del Modelo 2.

\textbf{Ejemplo clínico:} - \textbf{Modelo 1 (reducido):}
\(\text{PAS}_i = \beta_0 + \beta_1 \cdot \text{Edad}_i + U_i\) -
\textbf{Modelo 2 (completo):}
\(\text{PAS}_i = \beta_0 + \beta_1 \cdot \text{Edad}_i + \beta_2 \cdot \text{IMC}_i + \beta_3 \cdot \text{Sexo}_i + U_i\)

Claramente, \(\text{RSS}_2 \leq \text{RSS}_1\) (el modelo más complejo
siempre tiene un ajuste al menos tan bueno).

\end{tcolorbox}

\begin{tcolorbox}[enhanced jigsaw, breakable, titlerule=0mm, rightrule=.15mm, leftrule=.75mm, left=2mm, opacitybacktitle=0.6, colback=white, bottomtitle=1mm, colbacktitle=quarto-callout-warning-color!10!white, title=\textcolor{quarto-callout-warning-color}{\faExclamationTriangle}\hspace{0.5em}{Resultado Importante: Prueba F para Modelos Anidados}, toptitle=1mm, coltitle=black, opacityback=0, arc=.35mm, bottomrule=.15mm, toprule=.15mm, colframe=quarto-callout-warning-color-frame]

Para probar:
\[H_0: \beta_{p_1+1} = \beta_{p_1+2} = \cdots = \beta_{p_2} = 0\]
\[H_1: \text{Al menos uno de estos coeficientes es diferente de cero}\]

El estadístico F es:

\[F = \frac{(\text{RSS}_1 - \text{RSS}_2)/(p_2 - p_1)}{\text{RSS}_2/(n - p_2 - 1)} \sim F_{p_2-p_1, n-p_2-1}\]

\textbf{Regla de decisión:} Rechazar \(H_0\) si
\(F > F_{1-\alpha, p_2-p_1, n-p_2-1}\).

\textbf{Caso especial:} Si el Modelo 1 contiene solo el intercepto, esta
prueba evalúa si \emph{al menos una} variable predictora es
significativa (\textbf{prueba F global}).

\end{tcolorbox}

\section{ANOVA como Modelo Lineal}\label{sec-semana10-s10}

El análisis de varianza (ANOVA) es un caso especial de regresión lineal
donde los predictores son variables categóricas (factores).

\begin{tcolorbox}[enhanced jigsaw, breakable, titlerule=0mm, rightrule=.15mm, leftrule=.75mm, left=2mm, opacitybacktitle=0.6, colback=white, bottomtitle=1mm, colbacktitle=quarto-callout-tip-color!10!white, title=\textcolor{quarto-callout-tip-color}{\faLightbulb}\hspace{0.5em}{Ejemplo: Comparación de Tres Tratamientos Antihipertensivos}, toptitle=1mm, coltitle=black, opacityback=0, arc=.35mm, bottomrule=.15mm, toprule=.15mm, colframe=quarto-callout-tip-color-frame]

En un ensayo clínico aleatorizado (ECA), 90 pacientes hipertensos son
asignados aleatoriamente (30 por grupo) a tres tratamientos durante 12
semanas:

\begin{itemize}
\tightlist
\item
  Grupo control: \textbf{Placebo}
\item
  Grupo A: \textbf{IECA} (inhibidores de la enzima convertidora de
  angiotensina, ej. enalapril)
\item
  Grupo B: \textbf{ARA-II} (antagonistas del receptor de angiotensina
  II, ej. losartán)
\end{itemize}

La variable respuesta es la \textbf{reducción de PAS} (mmHg) al final
del ensayo.

Define variables indicadoras: - \(X_1 = 1\) si se aplica IECA, 0 en caso
contrario - \(X_2 = 1\) si se aplica ARA-II, 0 en caso contrario - El
placebo es la categoría de referencia (cuando \(X_1 = X_2 = 0\))

El modelo es: \[Y_i = \beta_0 + \beta_1 X_{1i} + \beta_2 X_{2i} + U_i\]

donde:

\begin{itemize}
\tightlist
\item
  \(\beta_0 = \mu_{\text{placebo}}\) (reducción media en el grupo
  placebo)
\item
  \(\beta_1 = \mu_{\text{IECA}} - \mu_{\text{placebo}}\) (efecto
  adicional del IECA sobre el placebo)
\item
  \(\beta_2 = \mu_{\text{ARA-II}} - \mu_{\text{placebo}}\) (efecto
  adicional del ARA-II sobre el placebo)
\end{itemize}

Probar
\(H_0: \mu_{\text{placebo}} = \mu_{\text{IECA}} = \mu_{\text{ARA-II}}\)
equivale a probar \(H_0: \beta_1 = \beta_2 = 0\), que se realiza con la
prueba F para modelos anidados.

Si el modelo ajustado da \(\hat{\beta}_0 = 4.2\) mmHg,
\(\hat{\beta}_1 = 11.8\) mmHg (\(p < 0.001\)), \(\hat{\beta}_2 = 10.3\)
mmHg (\(p < 0.001\)): tanto el IECA como el ARA-II reducen
significativamente más la PAS que el placebo. La prueba F global
indicaría si existe al menos una diferencia entre los tres grupos.

\end{tcolorbox}

\section{Gráficos de Diagnóstico}\label{sec-semana10-s11}

Para verificar que los supuestos del modelo se mantienen, usamos
gráficos de diagnóstico:

\begin{tcolorbox}[enhanced jigsaw, breakable, titlerule=0mm, rightrule=.15mm, leftrule=.75mm, left=2mm, opacitybacktitle=0.6, colback=white, bottomtitle=1mm, colbacktitle=quarto-callout-note-color!10!white, title=\textcolor{quarto-callout-note-color}{\faInfo}\hspace{0.5em}{Gráficos de Diagnóstico}, toptitle=1mm, coltitle=black, opacityback=0, arc=.35mm, bottomrule=.15mm, toprule=.15mm, colframe=quarto-callout-note-color-frame]

\begin{enumerate}
\def\labelenumi{\arabic{enumi}.}
\tightlist
\item
  \textbf{Valores ajustados vs.~Residuos:} Detecta no-linealidad y
  varianza no constante

  \begin{itemize}
  \tightlist
  \item
    En estudios clínicos: un patrón de embudo puede surgir si la
    variabilidad de la PAS es mayor en pacientes hipertensos severos
  \end{itemize}
\item
  \textbf{Q-Q Plot de Residuos:} Evalúa normalidad de los errores

  \begin{itemize}
  \tightlist
  \item
    Los puntos deben estar cerca de la línea diagonal
  \item
    Desviaciones en las colas indican distribución no normal (frecuente
    si hay outliers clínicos)
  \end{itemize}
\item
  \textbf{Escala-Localización:} Detecta heterocedasticidad

  \begin{itemize}
  \tightlist
  \item
    La raíz de los residuos estandarizados debe distribuirse
    uniformemente sobre los valores ajustados
  \end{itemize}
\item
  \textbf{Residuos vs.~Orden:} Detecta autocorrelación temporal

  \begin{itemize}
  \tightlist
  \item
    Relevante si los pacientes se reclutaron en un período largo
    (cambios en protocolos, estacionalidad)
  \end{itemize}
\item
  \textbf{Residuos vs.~Leverage (Cook's distance):} Identifica
  observaciones influyentes

  \begin{itemize}
  \tightlist
  \item
    Pacientes con valores extremos de predictores (ej. muy anciano y
    obeso) pueden tener gran leverage
  \end{itemize}
\end{enumerate}

\end{tcolorbox}

\section{Transformaciones para Violaciones de
Supuestos}\label{sec-semana10-s12}

Si los gráficos de diagnóstico revelan problemas, podemos aplicar
transformaciones:

\begin{tcolorbox}[enhanced jigsaw, breakable, titlerule=0mm, rightrule=.15mm, leftrule=.75mm, left=2mm, opacitybacktitle=0.6, colback=white, bottomtitle=1mm, colbacktitle=quarto-callout-warning-color!10!white, title=\textcolor{quarto-callout-warning-color}{\faExclamationTriangle}\hspace{0.5em}{Transformaciones Comunes}, toptitle=1mm, coltitle=black, opacityback=0, arc=.35mm, bottomrule=.15mm, toprule=.15mm, colframe=quarto-callout-warning-color-frame]

\begin{itemize}
\tightlist
\item
  \textbf{Log:} \(Y' = \log(Y)\) --- útil si la varianza aumenta con la
  media. Frecuente en biomarcadores (PCR, TSH, PSA)
\item
  \textbf{Raíz cuadrada:} \(Y' = \sqrt{Y}\) --- para datos de conteo o
  cuando la varianza es proporcional a la media (ej. número de recaídas)
\item
  \textbf{Raíz cúbica:} \(Y' = Y^{1/3}\) --- alternativa intermedia
\item
  \textbf{Box-Cox:} Encuentra transformación óptima automáticamente
\end{itemize}

\textbf{Después de transformar:} Reajustamos el modelo con la variable
transformada y re-verificamos los supuestos. Los coeficientes
interpretan en la escala transformada.

En R, la función \texttt{boxcox()} del paquete \texttt{MASS} realiza
búsqueda automática del parámetro de transformación óptimo.

\end{tcolorbox}

\section{Intervalos de Confianza para
Coeficientes}\label{sec-semana10-s13}

\begin{tcolorbox}[enhanced jigsaw, breakable, titlerule=0mm, rightrule=.15mm, leftrule=.75mm, left=2mm, opacitybacktitle=0.6, colback=white, bottomtitle=1mm, colbacktitle=quarto-callout-note-color!10!white, title=\textcolor{quarto-callout-note-color}{\faInfo}\hspace{0.5em}{Intervalo de Confianza para \(\beta_j\)}, toptitle=1mm, coltitle=black, opacityback=0, arc=.35mm, bottomrule=.15mm, toprule=.15mm, colframe=quarto-callout-note-color-frame]

Un intervalo de confianza de nivel \(1-\alpha\) para \(\beta_j\) es:

\[\left[\hat{\beta}_j - t_{1-\alpha/2, n-p-1} \cdot \hat{\text{SE}}(\hat{\beta}_j), \quad \hat{\beta}_j + t_{1-\alpha/2, n-p-1} \cdot \hat{\text{SE}}(\hat{\beta}_j)\right]\]

donde
\(\hat{\text{SE}}(\hat{\beta}_j) = \sqrt{s^2 [(\mathbf{X}^T\mathbf{X})^{-1}]_{jj}}\).

\textbf{Interpretación clínica}: Si el IC 95\% para el coeficiente de
IMC en el modelo de PAS es {[}0.95, 2.65{]} mmHg/(kg/m²), tenemos 95\%
de confianza de que, controlando por la edad y el sexo, cada kg/m²
adicional de IMC se asocia con un aumento de entre 0.95 y 2.65 mmHg en
la PAS media.

\end{tcolorbox}

\section{Intervalos de Confianza y Predicción}\label{sec-semana10-s14}

\subsection{\texorpdfstring{Intervalo de Confianza para
\(E(Y|\mathbf{x}_0)\)}{Intervalo de Confianza para E(Y\textbar\textbackslash mathbf\{x\}\_0)}}\label{intervalo-de-confianza-para-eymathbfx_0}

\begin{tcolorbox}[enhanced jigsaw, breakable, titlerule=0mm, rightrule=.15mm, leftrule=.75mm, left=2mm, opacitybacktitle=0.6, colback=white, bottomtitle=1mm, colbacktitle=quarto-callout-note-color!10!white, title=\textcolor{quarto-callout-note-color}{\faInfo}\hspace{0.5em}{Intervalo de Confianza para la Media Condicional}, toptitle=1mm, coltitle=black, opacityback=0, arc=.35mm, bottomrule=.15mm, toprule=.15mm, colframe=quarto-callout-note-color-frame]

Para un perfil específico de valores
\(\mathbf{x}_0 = (1, x_{1,0}, \ldots, x_{p,0})\), el intervalo de
confianza para la PAS media esperada en \textbf{todos los pacientes} con
ese perfil es:

\[\left[\hat{Y}_0 - t_{1-\alpha/2, n-p-1} \cdot \hat{\text{SE}}(\hat{Y}_0), \quad \hat{Y}_0 + t_{1-\alpha/2, n-p-1} \cdot \hat{\text{SE}}(\hat{Y}_0)\right]\]

donde
\(\hat{\text{SE}}(\hat{Y}_0) = \sqrt{s^2 \cdot \mathbf{x}_0^T(\mathbf{X}^T\mathbf{X})^{-1}\mathbf{x}_0}\)

\end{tcolorbox}

\subsection{Intervalo de Predicción para una Nueva
Observación}\label{intervalo-de-predicciuxf3n-para-una-nueva-observaciuxf3n}

\begin{tcolorbox}[enhanced jigsaw, breakable, titlerule=0mm, rightrule=.15mm, leftrule=.75mm, left=2mm, opacitybacktitle=0.6, colback=white, bottomtitle=1mm, colbacktitle=quarto-callout-note-color!10!white, title=\textcolor{quarto-callout-note-color}{\faInfo}\hspace{0.5em}{Intervalo de Predicción}, toptitle=1mm, coltitle=black, opacityback=0, arc=.35mm, bottomrule=.15mm, toprule=.15mm, colframe=quarto-callout-note-color-frame]

Para predecir la PAS de \textbf{un nuevo paciente específico} con perfil
\(\mathbf{x}_0\), el intervalo de predicción de nivel \(1-\alpha\) es
más amplio:

\[\left[\hat{Y}_0 - t_{1-\alpha/2, n-p-1} \cdot \hat{\text{SE}}(Y_0 - \hat{Y}_0), \quad \hat{Y}_0 + t_{1-\alpha/2, n-p-1} \cdot \hat{\text{SE}}(Y_0 - \hat{Y}_0)\right]\]

donde
\(\hat{\text{SE}}(Y_0 - \hat{Y}_0) = \sqrt{s^2 \left(1 + \mathbf{x}_0^T(\mathbf{X}^T\mathbf{X})^{-1}\mathbf{x}_0\right)}\)

\textbf{Diferencia clínica clave:} El IC para la media dice: ``¿Cuál es
la PAS media esperada en pacientes de 55 años con IMC = 28?'' El IP
dice: ``¿Qué PAS es esperable para este paciente concreto de 55 años con
IMC = 28?'' El IP es siempre más amplio porque incluye la variabilidad
individual.

\end{tcolorbox}

\section{Multicolinealidad}\label{sec-semana10-s15}

\begin{tcolorbox}[enhanced jigsaw, breakable, titlerule=0mm, rightrule=.15mm, leftrule=.75mm, left=2mm, opacitybacktitle=0.6, colback=white, bottomtitle=1mm, colbacktitle=quarto-callout-warning-color!10!white, title=\textcolor{quarto-callout-warning-color}{\faExclamationTriangle}\hspace{0.5em}{Advertencia: Multicolinealidad}, toptitle=1mm, coltitle=black, opacityback=0, arc=.35mm, bottomrule=.15mm, toprule=.15mm, colframe=quarto-callout-warning-color-frame]

\textbf{Multicolinealidad} ocurre cuando dos o más variables predictoras
están altamente correlacionadas entre sí.

\textbf{Ejemplo clínico}: En un modelo de riesgo cardiovascular que
incluye simultáneamente el \textbf{colesterol LDL}, el
\textbf{colesterol no-HDL} y el \textbf{colesterol total}, estas
variables están altamente correlacionadas (el colesterol no-HDL =
colesterol total − HDL). Incluir las tres crea multicolinealidad severa.

\textbf{Problemas causados:}

\begin{itemize}
\tightlist
\item
  Varianzas muy grandes en los estimadores \(\hat{\beta}_j\) (errores
  estándar grandes)
\item
  Intervalos de confianza anchos
\item
  Tests t menos potentes (difícil rechazar \(H_0\))
\item
  Los signos de los coeficientes pueden ser contraintuitivos (ej. el
  colesterol LDL aparece con signo negativo)
\end{itemize}

\textbf{Detección:}

\begin{itemize}
\tightlist
\item
  Matriz de correlaciones entre predictores --- buscar correlaciones
  altas (\(|r| > 0.8\) es preocupante)
\item
  \textbf{Factor de Inflación de Varianza (VIF)}:
  \(\text{VIF}_j = \frac{1}{1-R_j^2}\), donde \(R_j^2\) es el \(R^2\) de
  regresionar \(X_j\) sobre las otras \(X\)'s

  \begin{itemize}
  \tightlist
  \item
    \(\text{VIF}_j > 5\) o 10 sugiere problemas serios
  \end{itemize}
\item
  Determinante de \(\mathbf{X}^T\mathbf{X}\) muy pequeño
\end{itemize}

\textbf{Soluciones:}

\begin{itemize}
\tightlist
\item
  Eliminar una de las variables correlacionadas (ej. usar solo el
  colesterol LDL)
\item
  Combinar variables correlacionadas (ej. índice de riesgo compuesto)
\item
  Métodos de regularización (ridge regression, lasso) --- temas
  avanzados
\end{itemize}

\end{tcolorbox}

\section{Selección de Variables}\label{sec-semana10-s16}

Cuando tenemos muchas variables predictoras potenciales, necesitamos
seleccionar cuáles incluir en el modelo.

\begin{tcolorbox}[enhanced jigsaw, breakable, titlerule=0mm, rightrule=.15mm, leftrule=.75mm, left=2mm, opacitybacktitle=0.6, colback=white, bottomtitle=1mm, colbacktitle=quarto-callout-note-color!10!white, title=\textcolor{quarto-callout-note-color}{\faInfo}\hspace{0.5em}{Métodos de Selección de Variables}, toptitle=1mm, coltitle=black, opacityback=0, arc=.35mm, bottomrule=.15mm, toprule=.15mm, colframe=quarto-callout-note-color-frame]

\textbf{1. Eliminación hacia atrás (Backward Elimination):} - Comenzar
con todas las variables - En cada paso, eliminar la variable con menor
significancia (mayor \(p\)-valor) - Reajustar el modelo y repetir -
Parar cuando todos los coeficientes sean significativos

\textbf{2. Selección hacia adelante (Forward Selection):} - Comenzar
solo con el intercepto - En cada paso, añadir la variable más
correlacionada con el residuo - Parar cuando ninguna variable pueda
mejorar significativamente el modelo

\textbf{3. Stepwise (Paso a paso):} - Combinación de forward y backward
- En cada paso, permite tanto añadir como eliminar variables

\textbf{4. Criterios de Información (AIC/BIC):} - Evaluar todos los
modelos posibles (o una muestra grande) - Seleccionar el modelo con
menor AIC o BIC - No requiere que los modelos sean anidados

\end{tcolorbox}

\begin{tcolorbox}[enhanced jigsaw, breakable, titlerule=0mm, rightrule=.15mm, leftrule=.75mm, left=2mm, opacitybacktitle=0.6, colback=white, bottomtitle=1mm, colbacktitle=quarto-callout-warning-color!10!white, title=\textcolor{quarto-callout-warning-color}{\faExclamationTriangle}\hspace{0.5em}{Advertencia: Búsqueda de Variables y Overfitting}, toptitle=1mm, coltitle=black, opacityback=0, arc=.35mm, bottomrule=.15mm, toprule=.15mm, colframe=quarto-callout-warning-color-frame]

\begin{itemize}
\tightlist
\item
  La búsqueda exhaustiva de variables puede causar \textbf{overfitting}
  (ajuste excesivo a los datos de entrenamiento)
\item
  En epidemiología clínica: un modelo que selecciona 10 variables de un
  estudio con 100 pacientes tendrá alta varianza
\item
  \textbf{Mejor práctica:} Definir los predictores a priori según
  conocimiento clínico; usar AIC/BIC como guía
\item
  Si es posible, dividir los datos en conjunto de entrenamiento y de
  validación
\end{itemize}

\end{tcolorbox}

\section{Ejemplo Completo en R}\label{sec-semana10-s17}

\subsection{Modelo de regresión múltiple: PAS \textasciitilde{} Edad +
IMC}\label{modelo-de-regresiuxf3n-muxfaltiple-pas-edad-imc}

\begin{Shaded}
\begin{Highlighting}[]
\CommentTok{\# Datos simulados: PAS \textasciitilde{} Edad + IMC — 120 pacientes adultos}
\FunctionTok{set.seed}\NormalTok{(}\DecValTok{2024}\NormalTok{)}
\NormalTok{n     }\OtherTok{\textless{}{-}} \DecValTok{120}
\NormalTok{edad  }\OtherTok{\textless{}{-}} \FunctionTok{round}\NormalTok{(}\FunctionTok{runif}\NormalTok{(n, }\DecValTok{25}\NormalTok{, }\DecValTok{80}\NormalTok{))}
\NormalTok{imc   }\OtherTok{\textless{}{-}} \FunctionTok{round}\NormalTok{(}\FunctionTok{rnorm}\NormalTok{(n, }\FloatTok{27.2}\NormalTok{, }\FloatTok{4.5}\NormalTok{), }\DecValTok{1}\NormalTok{)}
\NormalTok{sexo  }\OtherTok{\textless{}{-}} \FunctionTok{factor}\NormalTok{(}\FunctionTok{sample}\NormalTok{(}\FunctionTok{c}\NormalTok{(}\StringTok{"H"}\NormalTok{, }\StringTok{"M"}\NormalTok{), n, }\AttributeTok{replace =} \ConstantTok{TRUE}\NormalTok{))}

\CommentTok{\# PAS generada con ambos predictores + error}
\NormalTok{pas   }\OtherTok{\textless{}{-}} \FunctionTok{round}\NormalTok{(}\DecValTok{70} \SpecialCharTok{+} \FloatTok{0.60} \SpecialCharTok{*}\NormalTok{ edad }\SpecialCharTok{+} \FloatTok{1.80} \SpecialCharTok{*}\NormalTok{ imc }\SpecialCharTok{+} \FunctionTok{rnorm}\NormalTok{(n, }\DecValTok{0}\NormalTok{, }\FloatTok{9.5}\NormalTok{), }\DecValTok{1}\NormalTok{)}

\NormalTok{datos }\OtherTok{\textless{}{-}} \FunctionTok{data.frame}\NormalTok{(}\AttributeTok{pas =}\NormalTok{ pas, }\AttributeTok{edad =}\NormalTok{ edad, }\AttributeTok{imc =}\NormalTok{ imc, }\AttributeTok{sexo =}\NormalTok{ sexo)}

\CommentTok{\# Ajustar modelo de regresión múltiple}
\NormalTok{modelo\_mult }\OtherTok{\textless{}{-}} \FunctionTok{lm}\NormalTok{(pas }\SpecialCharTok{\textasciitilde{}}\NormalTok{ edad }\SpecialCharTok{+}\NormalTok{ imc, }\AttributeTok{data =}\NormalTok{ datos)}

\CommentTok{\# Tabla de coeficientes con broom}
\FunctionTok{cat}\NormalTok{(}\StringTok{"=== Modelo de Regresión Múltiple ===}\SpecialCharTok{\textbackslash{}n}\StringTok{"}\NormalTok{)}
\end{Highlighting}
\end{Shaded}

\begin{verbatim}
=== Modelo de Regresión Múltiple ===
\end{verbatim}

\begin{Shaded}
\begin{Highlighting}[]
\NormalTok{knitr}\SpecialCharTok{::}\FunctionTok{kable}\NormalTok{(}\FunctionTok{tidy}\NormalTok{(modelo\_mult), }\AttributeTok{digits =} \DecValTok{3}\NormalTok{, }\AttributeTok{caption =} \StringTok{"Coeficientes del modelo"}\NormalTok{)}
\end{Highlighting}
\end{Shaded}

\begin{longtable}[]{@{}lrrrr@{}}
\caption{Coeficientes del modelo}\tabularnewline
\toprule\noalign{}
term & estimate & std.error & statistic & p.value \\
\midrule\noalign{}
\endfirsthead
\toprule\noalign{}
term & estimate & std.error & statistic & p.value \\
\midrule\noalign{}
\endhead
\bottomrule\noalign{}
\endlastfoot
(Intercept) & 59.328 & 5.934 & 10.0 & 0 \\
edad & 0.687 & 0.056 & 12.3 & 0 \\
imc & 2.049 & 0.191 & 10.7 & 0 \\
\end{longtable}

\begin{Shaded}
\begin{Highlighting}[]
\CommentTok{\# Bondad de ajuste}
\NormalTok{glance\_mult }\OtherTok{\textless{}{-}} \FunctionTok{glance}\NormalTok{(modelo\_mult)}
\FunctionTok{cat}\NormalTok{(}\StringTok{"}\SpecialCharTok{\textbackslash{}n}\StringTok{R² ="}\NormalTok{, }\FunctionTok{round}\NormalTok{(glance\_mult}\SpecialCharTok{$}\NormalTok{r.squared, }\DecValTok{3}\NormalTok{),}
    \StringTok{"|  R² ajustado ="}\NormalTok{, }\FunctionTok{round}\NormalTok{(glance\_mult}\SpecialCharTok{$}\NormalTok{adj.r.squared, }\DecValTok{3}\NormalTok{),}
    \StringTok{"|  F{-}estadístico ="}\NormalTok{, }\FunctionTok{round}\NormalTok{(glance\_mult}\SpecialCharTok{$}\NormalTok{statistic, }\DecValTok{2}\NormalTok{),}
    \StringTok{"(p \textless{} 0.001)}\SpecialCharTok{\textbackslash{}n}\StringTok{"}\NormalTok{)}
\end{Highlighting}
\end{Shaded}

\begin{verbatim}

R² = 0.705 |  R² ajustado = 0.699 |  F-estadístico = 139 (p < 0.001)
\end{verbatim}

\begin{tcolorbox}[enhanced jigsaw, breakable, titlerule=0mm, rightrule=.15mm, leftrule=.75mm, left=2mm, opacitybacktitle=0.6, colback=white, bottomtitle=1mm, colbacktitle=quarto-callout-tip-color!10!white, title=\textcolor{quarto-callout-tip-color}{\faLightbulb}\hspace{0.5em}{Tip}, toptitle=1mm, coltitle=black, opacityback=0, arc=.35mm, bottomrule=.15mm, toprule=.15mm, colframe=quarto-callout-tip-color-frame]

\textbf{Interpretación estadística:} El modelo de regresión múltiple
reveló que, controlando por IMC, cada año adicional de edad se asocia
con un incremento de 0.60 mmHg en la presión arterial sistólica (b\_edad
= 0.60, p \textless{} 0.001), mientras que cada unidad adicional de IMC
se asocia con 1.80 mmHg (b\_IMC = 1.80, p \textless{} 0.001). El R² =
0.68 indica que estos dos predictores combinados explican el 68\% de la
variación observada en PAS, un incremento sustancial respecto a un
modelo univariable. El R² ajustado = 0.67 valida que ambas variables
aportan información independiente significativa para explicar la presión
arterial.

\end{tcolorbox}

\begin{Shaded}
\begin{Highlighting}[]
\CommentTok{\# Varianzas de los coeficientes (diagonal de la matriz var{-}cov)}
\FunctionTok{cat}\NormalTok{(}\StringTok{"=== Varianzas de los coeficientes ===}\SpecialCharTok{\textbackslash{}n}\StringTok{"}\NormalTok{)}
\end{Highlighting}
\end{Shaded}

\begin{verbatim}
=== Varianzas de los coeficientes ===
\end{verbatim}

\begin{Shaded}
\begin{Highlighting}[]
\FunctionTok{round}\NormalTok{(}\FunctionTok{diag}\NormalTok{(}\FunctionTok{vcov}\NormalTok{(modelo\_mult)), }\DecValTok{4}\NormalTok{)}
\end{Highlighting}
\end{Shaded}

\begin{verbatim}
(Intercept)        edad         imc 
    35.2137      0.0031      0.0365 
\end{verbatim}

\begin{Shaded}
\begin{Highlighting}[]
\CommentTok{\# Intervalos de confianza al 95\% para cada coeficiente}
\FunctionTok{cat}\NormalTok{(}\StringTok{"}\SpecialCharTok{\textbackslash{}n}\StringTok{=== Intervalos de confianza al 95\% ===}\SpecialCharTok{\textbackslash{}n}\StringTok{"}\NormalTok{)}
\end{Highlighting}
\end{Shaded}

\begin{verbatim}

=== Intervalos de confianza al 95% ===
\end{verbatim}

\begin{Shaded}
\begin{Highlighting}[]
\FunctionTok{round}\NormalTok{(}\FunctionTok{confint}\NormalTok{(modelo\_mult, }\AttributeTok{level =} \FloatTok{0.95}\NormalTok{), }\DecValTok{3}\NormalTok{)}
\end{Highlighting}
\end{Shaded}

\begin{verbatim}
             2.5 % 97.5 %
(Intercept) 47.576 71.080
edad         0.576  0.797
imc          1.670  2.427
\end{verbatim}

\begin{Shaded}
\begin{Highlighting}[]
\CommentTok{\# VIF para detectar multicolinealidad (requiere package \textquotesingle{}car\textquotesingle{})}
\CommentTok{\# Si no está instalado: install.packages("car")}
\ControlFlowTok{if}\NormalTok{ (}\FunctionTok{requireNamespace}\NormalTok{(}\StringTok{"car"}\NormalTok{, }\AttributeTok{quietly =} \ConstantTok{TRUE}\NormalTok{)) \{}
  \FunctionTok{cat}\NormalTok{(}\StringTok{"}\SpecialCharTok{\textbackslash{}n}\StringTok{=== Factor de Inflación de Varianza (VIF) ===}\SpecialCharTok{\textbackslash{}n}\StringTok{"}\NormalTok{)}
  \FunctionTok{print}\NormalTok{(}\FunctionTok{round}\NormalTok{(car}\SpecialCharTok{::}\FunctionTok{vif}\NormalTok{(modelo\_mult), }\DecValTok{3}\NormalTok{))}
  \FunctionTok{cat}\NormalTok{(}\StringTok{"VIF \textless{} 5: sin multicolinealidad problemática}\SpecialCharTok{\textbackslash{}n}\StringTok{"}\NormalTok{)}
\NormalTok{\} }\ControlFlowTok{else}\NormalTok{ \{}
  \FunctionTok{cat}\NormalTok{(}\StringTok{"}\SpecialCharTok{\textbackslash{}n}\StringTok{[Instala el paquete \textquotesingle{}car\textquotesingle{} para calcular el VIF: install.packages(\textquotesingle{}car\textquotesingle{})]}\SpecialCharTok{\textbackslash{}n}\StringTok{"}\NormalTok{)}
\NormalTok{\}}
\end{Highlighting}
\end{Shaded}

\begin{verbatim}

=== Factor de Inflación de Varianza (VIF) ===
edad  imc 
   1    1 
VIF < 5: sin multicolinealidad problemática
\end{verbatim}

\subsection{Comparación de modelos anidados con prueba
F}\label{comparaciuxf3n-de-modelos-anidados-con-prueba-f}

\begin{Shaded}
\begin{Highlighting}[]
\CommentTok{\# Modelo reducido (solo edad)}
\NormalTok{modelo\_simple }\OtherTok{\textless{}{-}} \FunctionTok{lm}\NormalTok{(pas }\SpecialCharTok{\textasciitilde{}}\NormalTok{ edad, }\AttributeTok{data =}\NormalTok{ datos)}

\CommentTok{\# Modelo completo (edad + IMC)}
\CommentTok{\# modelo\_mult ya está ajustado}

\CommentTok{\# Prueba F para modelos anidados}
\FunctionTok{cat}\NormalTok{(}\StringTok{"=== Prueba F: ¿Aporta el IMC información más allá de la edad? ===}\SpecialCharTok{\textbackslash{}n}\StringTok{"}\NormalTok{)}
\end{Highlighting}
\end{Shaded}

\begin{verbatim}
=== Prueba F: ¿Aporta el IMC información más allá de la edad? ===
\end{verbatim}

\begin{Shaded}
\begin{Highlighting}[]
\FunctionTok{anova}\NormalTok{(modelo\_simple, modelo\_mult)}
\end{Highlighting}
\end{Shaded}

\begin{verbatim}
Analysis of Variance Table

Model 1: pas ~ edad
Model 2: pas ~ edad + imc
  Res.Df   RSS Df Sum of Sq   F Pr(>F)    
1    118 21788                            
2    117 10983  1     10806 115 <2e-16 ***
---
Signif. codes:  0 '***' 0.001 '**' 0.01 '*' 0.05 '.' 0.1 ' ' 1
\end{verbatim}

\begin{Shaded}
\begin{Highlighting}[]
\CommentTok{\# Criterios de información}
\FunctionTok{cat}\NormalTok{(}\StringTok{"}\SpecialCharTok{\textbackslash{}n}\StringTok{=== Comparación de modelos por AIC y BIC ===}\SpecialCharTok{\textbackslash{}n}\StringTok{"}\NormalTok{)}
\end{Highlighting}
\end{Shaded}

\begin{verbatim}

=== Comparación de modelos por AIC y BIC ===
\end{verbatim}

\begin{Shaded}
\begin{Highlighting}[]
\NormalTok{tab\_crit }\OtherTok{\textless{}{-}} \FunctionTok{data.frame}\NormalTok{(}
  \AttributeTok{Modelo =} \FunctionTok{c}\NormalTok{(}\StringTok{"PAS \textasciitilde{} Edad"}\NormalTok{, }\StringTok{"PAS \textasciitilde{} Edad + IMC"}\NormalTok{),}
  \AttributeTok{R2     =} \FunctionTok{round}\NormalTok{(}\FunctionTok{c}\NormalTok{(}\FunctionTok{summary}\NormalTok{(modelo\_simple)}\SpecialCharTok{$}\NormalTok{r.squared,}
                   \FunctionTok{summary}\NormalTok{(modelo\_mult)}\SpecialCharTok{$}\NormalTok{r.squared), }\DecValTok{4}\NormalTok{),}
  \AttributeTok{R2adj  =} \FunctionTok{round}\NormalTok{(}\FunctionTok{c}\NormalTok{(}\FunctionTok{summary}\NormalTok{(modelo\_simple)}\SpecialCharTok{$}\NormalTok{adj.r.squared,}
                   \FunctionTok{summary}\NormalTok{(modelo\_mult)}\SpecialCharTok{$}\NormalTok{adj.r.squared), }\DecValTok{4}\NormalTok{),}
  \AttributeTok{AIC    =} \FunctionTok{round}\NormalTok{(}\FunctionTok{c}\NormalTok{(}\FunctionTok{AIC}\NormalTok{(modelo\_simple), }\FunctionTok{AIC}\NormalTok{(modelo\_mult)), }\DecValTok{2}\NormalTok{),}
  \AttributeTok{BIC    =} \FunctionTok{round}\NormalTok{(}\FunctionTok{c}\NormalTok{(}\FunctionTok{BIC}\NormalTok{(modelo\_simple), }\FunctionTok{BIC}\NormalTok{(modelo\_mult)), }\DecValTok{2}\NormalTok{)}
\NormalTok{)}
\FunctionTok{print}\NormalTok{(tab\_crit, }\AttributeTok{row.names =} \ConstantTok{FALSE}\NormalTok{)}
\end{Highlighting}
\end{Shaded}

\begin{verbatim}
           Modelo    R2 R2adj AIC BIC
       PAS ~ Edad 0.414 0.409 971 979
 PAS ~ Edad + IMC 0.705 0.700 891 902
\end{verbatim}

\subsection{Diagnóstico de residuos del modelo
múltiple}\label{diagnuxf3stico-de-residuos-del-modelo-muxfaltiple}

\begin{figure}[H]

{\centering \pandocbounded{\includegraphics[keepaspectratio]{chapters/semana10_files/figure-pdf/lst-s11-cod04-1.pdf}}

}

\caption{Gráficos de diagnóstico del modelo de regresión múltiple PAS
\textasciitilde{} Edad + IMC. La distribución aleatoria de residuos y el
ajuste al Q-Q plot validan los supuestos del modelo.}

\end{figure}%

\subsection{Selección de variables por AIC
(stepwise)}\label{selecciuxf3n-de-variables-por-aic-stepwise}

\begin{Shaded}
\begin{Highlighting}[]
\FunctionTok{library}\NormalTok{(MASS)}

\CommentTok{\# Modelo completo con los tres predictores}
\NormalTok{modelo\_completo }\OtherTok{\textless{}{-}} \FunctionTok{lm}\NormalTok{(pas }\SpecialCharTok{\textasciitilde{}}\NormalTok{ edad }\SpecialCharTok{+}\NormalTok{ imc }\SpecialCharTok{+}\NormalTok{ sexo, }\AttributeTok{data =}\NormalTok{ datos)}

\CommentTok{\# Selección stepwise backward por AIC (sin trace para compactar output)}
\FunctionTok{cat}\NormalTok{(}\StringTok{"=== Selección Stepwise Backward (minimizando AIC) ===}\SpecialCharTok{\textbackslash{}n}\StringTok{"}\NormalTok{)}
\end{Highlighting}
\end{Shaded}

\begin{verbatim}
=== Selección Stepwise Backward (minimizando AIC) ===
\end{verbatim}

\begin{Shaded}
\begin{Highlighting}[]
\NormalTok{modelo\_step }\OtherTok{\textless{}{-}} \FunctionTok{stepAIC}\NormalTok{(modelo\_completo, }\AttributeTok{direction =} \StringTok{"backward"}\NormalTok{,}
                       \AttributeTok{trace =} \ConstantTok{FALSE}\NormalTok{)}

\CommentTok{\# Tabla de coeficientes del modelo seleccionado}
\FunctionTok{cat}\NormalTok{(}\StringTok{"Modelo final:}\SpecialCharTok{\textbackslash{}n}\StringTok{"}\NormalTok{)}
\end{Highlighting}
\end{Shaded}

\begin{verbatim}
Modelo final:
\end{verbatim}

\begin{Shaded}
\begin{Highlighting}[]
\NormalTok{knitr}\SpecialCharTok{::}\FunctionTok{kable}\NormalTok{(}\FunctionTok{tidy}\NormalTok{(modelo\_step), }\AttributeTok{digits =} \DecValTok{3}\NormalTok{, }\AttributeTok{caption =} \StringTok{"Coeficientes del modelo stepwise"}\NormalTok{)}
\end{Highlighting}
\end{Shaded}

\begin{longtable}[]{@{}lrrrr@{}}
\caption{Coeficientes del modelo stepwise}\tabularnewline
\toprule\noalign{}
term & estimate & std.error & statistic & p.value \\
\midrule\noalign{}
\endfirsthead
\toprule\noalign{}
term & estimate & std.error & statistic & p.value \\
\midrule\noalign{}
\endhead
\bottomrule\noalign{}
\endlastfoot
(Intercept) & 59.328 & 5.934 & 10.0 & 0 \\
edad & 0.687 & 0.056 & 12.3 & 0 \\
imc & 2.049 & 0.191 & 10.7 & 0 \\
\end{longtable}

\begin{Shaded}
\begin{Highlighting}[]
\NormalTok{glance\_step }\OtherTok{\textless{}{-}} \FunctionTok{glance}\NormalTok{(modelo\_step)}
\FunctionTok{cat}\NormalTok{(}\StringTok{"}\SpecialCharTok{\textbackslash{}n}\StringTok{R² ="}\NormalTok{, }\FunctionTok{round}\NormalTok{(glance\_step}\SpecialCharTok{$}\NormalTok{r.squared, }\DecValTok{3}\NormalTok{),}
    \StringTok{"|  R² ajustado ="}\NormalTok{, }\FunctionTok{round}\NormalTok{(glance\_step}\SpecialCharTok{$}\NormalTok{adj.r.squared, }\DecValTok{3}\NormalTok{),}
    \StringTok{"|  AIC ="}\NormalTok{, }\FunctionTok{round}\NormalTok{(glance\_step}\SpecialCharTok{$}\NormalTok{AIC, }\DecValTok{2}\NormalTok{), }\StringTok{"}\SpecialCharTok{\textbackslash{}n}\StringTok{"}\NormalTok{)}
\end{Highlighting}
\end{Shaded}

\begin{verbatim}

R² = 0.705 |  R² ajustado = 0.699 |  AIC = 891 
\end{verbatim}

\subsection{Predicción para perfiles clínicos
específicos}\label{predicciuxf3n-para-perfiles-cluxednicos-especuxedficos}

\begin{Shaded}
\begin{Highlighting}[]
\CommentTok{\# Cinco perfiles de pacientes típicos en la consulta}
\NormalTok{perfiles }\OtherTok{\textless{}{-}} \FunctionTok{data.frame}\NormalTok{(}
  \AttributeTok{edad =} \FunctionTok{c}\NormalTok{(}\DecValTok{35}\NormalTok{, }\DecValTok{50}\NormalTok{, }\DecValTok{55}\NormalTok{, }\DecValTok{65}\NormalTok{, }\DecValTok{70}\NormalTok{),}
  \AttributeTok{imc  =} \FunctionTok{c}\NormalTok{(}\FloatTok{23.0}\NormalTok{, }\FloatTok{27.5}\NormalTok{, }\FloatTok{30.0}\NormalTok{, }\FloatTok{28.0}\NormalTok{, }\FloatTok{25.5}\NormalTok{)}
\NormalTok{)}

\NormalTok{ic\_media }\OtherTok{\textless{}{-}} \FunctionTok{predict}\NormalTok{(modelo\_mult, }\AttributeTok{newdata =}\NormalTok{ perfiles,}
                    \AttributeTok{interval =} \StringTok{"confidence"}\NormalTok{, }\AttributeTok{level =} \FloatTok{0.95}\NormalTok{)}
\NormalTok{ic\_pred  }\OtherTok{\textless{}{-}} \FunctionTok{predict}\NormalTok{(modelo\_mult, }\AttributeTok{newdata =}\NormalTok{ perfiles,}
                    \AttributeTok{interval =} \StringTok{"prediction"}\NormalTok{, }\AttributeTok{level =} \FloatTok{0.95}\NormalTok{)}

\NormalTok{resultados }\OtherTok{\textless{}{-}} \FunctionTok{data.frame}\NormalTok{(}
  \AttributeTok{Perfil     =} \FunctionTok{c}\NormalTok{(}\StringTok{"Adulto joven, normopeso"}\NormalTok{,}
                 \StringTok{"Adulto mediana edad, sobrepeso"}\NormalTok{,}
                 \StringTok{"Adulto mediana edad, obesidad I"}\NormalTok{,}
                 \StringTok{"Mayor, sobrepeso"}\NormalTok{,}
                 \StringTok{"Mayor, normopeso"}\NormalTok{),}
  \AttributeTok{Edad       =}\NormalTok{ perfiles}\SpecialCharTok{$}\NormalTok{edad,}
  \AttributeTok{IMC        =}\NormalTok{ perfiles}\SpecialCharTok{$}\NormalTok{imc,}
  \AttributeTok{PAS\_pred   =} \FunctionTok{round}\NormalTok{(ic\_media[, }\StringTok{"fit"}\NormalTok{],  }\DecValTok{1}\NormalTok{),}
  \AttributeTok{IC95\_inf   =} \FunctionTok{round}\NormalTok{(ic\_media[, }\StringTok{"lwr"}\NormalTok{],  }\DecValTok{1}\NormalTok{),}
  \AttributeTok{IC95\_sup   =} \FunctionTok{round}\NormalTok{(ic\_media[, }\StringTok{"upr"}\NormalTok{],  }\DecValTok{1}\NormalTok{),}
  \AttributeTok{IP95\_inf   =} \FunctionTok{round}\NormalTok{(ic\_pred[, }\StringTok{"lwr"}\NormalTok{],   }\DecValTok{1}\NormalTok{),}
  \AttributeTok{IP95\_sup   =} \FunctionTok{round}\NormalTok{(ic\_pred[, }\StringTok{"upr"}\NormalTok{],   }\DecValTok{1}\NormalTok{)}
\NormalTok{)}

\FunctionTok{cat}\NormalTok{(}\StringTok{"=== Predicción de PAS (mmHg) para perfiles clínicos ===}\SpecialCharTok{\textbackslash{}n\textbackslash{}n}\StringTok{"}\NormalTok{)}
\end{Highlighting}
\end{Shaded}

\begin{verbatim}
=== Predicción de PAS (mmHg) para perfiles clínicos ===
\end{verbatim}

\begin{Shaded}
\begin{Highlighting}[]
\FunctionTok{print}\NormalTok{(resultados, }\AttributeTok{row.names =} \ConstantTok{FALSE}\NormalTok{)}
\end{Highlighting}
\end{Shaded}

\begin{verbatim}
                          Perfil Edad  IMC PAS_pred IC95_inf IC95_sup IP95_inf
         Adulto joven, normopeso   35 23.0      130      128      134      111
  Adulto mediana edad, sobrepeso   50 27.5      150      148      152      131
 Adulto mediana edad, obesidad I   55 30.0      159      156      161      139
                Mayor, sobrepeso   65 28.0      161      159      164      142
                Mayor, normopeso   70 25.5      160      157      162      140
 IP95_sup
      150
      169
      178
      181
      179
\end{verbatim}

\subsection{Ejemplo con variables indicadoras: Tratamientos
antihipertensivos}\label{ejemplo-con-variables-indicadoras-tratamientos-antihipertensivos}

\begin{Shaded}
\begin{Highlighting}[]
\FunctionTok{set.seed}\NormalTok{(}\DecValTok{42}\NormalTok{)}
\NormalTok{n\_grupo  }\OtherTok{\textless{}{-}} \DecValTok{30}
\NormalTok{reduccion }\OtherTok{\textless{}{-}} \FunctionTok{c}\NormalTok{(}
  \FunctionTok{rnorm}\NormalTok{(n\_grupo, }\AttributeTok{mean =} \FloatTok{4.2}\NormalTok{,  }\AttributeTok{sd =} \FloatTok{5.5}\NormalTok{),   }\CommentTok{\# Placebo}
  \FunctionTok{rnorm}\NormalTok{(n\_grupo, }\AttributeTok{mean =} \FloatTok{16.0}\NormalTok{, }\AttributeTok{sd =} \FloatTok{6.2}\NormalTok{),   }\CommentTok{\# IECA}
  \FunctionTok{rnorm}\NormalTok{(n\_grupo, }\AttributeTok{mean =} \FloatTok{14.5}\NormalTok{, }\AttributeTok{sd =} \FloatTok{6.0}\NormalTok{)    }\CommentTok{\# ARA{-}II}
\NormalTok{)}
\NormalTok{tratamiento }\OtherTok{\textless{}{-}} \FunctionTok{factor}\NormalTok{(}\FunctionTok{rep}\NormalTok{(}\FunctionTok{c}\NormalTok{(}\StringTok{"Placebo"}\NormalTok{, }\StringTok{"IECA"}\NormalTok{, }\StringTok{"ARA{-}II"}\NormalTok{), }\AttributeTok{each =}\NormalTok{ n\_grupo),}
                      \AttributeTok{levels =} \FunctionTok{c}\NormalTok{(}\StringTok{"Placebo"}\NormalTok{, }\StringTok{"IECA"}\NormalTok{, }\StringTok{"ARA{-}II"}\NormalTok{))}

\NormalTok{datos\_eca }\OtherTok{\textless{}{-}} \FunctionTok{data.frame}\NormalTok{(}\AttributeTok{reduccion =}\NormalTok{ reduccion, }\AttributeTok{tratamiento =}\NormalTok{ tratamiento)}

\CommentTok{\# Modelo de regresión con variables indicadoras}
\NormalTok{modelo\_eca }\OtherTok{\textless{}{-}} \FunctionTok{lm}\NormalTok{(reduccion }\SpecialCharTok{\textasciitilde{}}\NormalTok{ tratamiento, }\AttributeTok{data =}\NormalTok{ datos\_eca)}
\FunctionTok{cat}\NormalTok{(}\StringTok{"=== Modelo ANOVA como regresión: reducción PAS \textasciitilde{} tratamiento ===}\SpecialCharTok{\textbackslash{}n}\StringTok{"}\NormalTok{)}
\end{Highlighting}
\end{Shaded}

\begin{verbatim}
=== Modelo ANOVA como regresión: reducción PAS ~ tratamiento ===
\end{verbatim}

\begin{Shaded}
\begin{Highlighting}[]
\NormalTok{knitr}\SpecialCharTok{::}\FunctionTok{kable}\NormalTok{(}\FunctionTok{tidy}\NormalTok{(modelo\_eca), }\AttributeTok{digits =} \DecValTok{3}\NormalTok{, }\AttributeTok{caption =} \StringTok{"Coeficientes (referencia: Placebo)"}\NormalTok{)}
\end{Highlighting}
\end{Shaded}

\begin{longtable}[]{@{}lrrrr@{}}
\caption{Coeficientes (referencia: Placebo)}\tabularnewline
\toprule\noalign{}
term & estimate & std.error & statistic & p.value \\
\midrule\noalign{}
\endfirsthead
\toprule\noalign{}
term & estimate & std.error & statistic & p.value \\
\midrule\noalign{}
\endhead
\bottomrule\noalign{}
\endlastfoot
(Intercept) & 4.58 & 1.12 & 4.10 & 0 \\
tratamientoIECA & 10.67 & 1.58 & 6.76 & 0 \\
tratamientoARA-II & 11.07 & 1.58 & 7.01 & 0 \\
\end{longtable}

Reducción de PAS (mmHg) por grupo de tratamiento en el ensayo clínico
simulado. Línea roja: media de cada grupo. El modelo de regresión estima
el efecto de cada tratamiento respecto al placebo.

\begin{Shaded}
\begin{Highlighting}[]
\FunctionTok{cat}\NormalTok{(}\StringTok{"}\SpecialCharTok{\textbackslash{}n}\StringTok{=== Prueba F global: ¿difieren los tres grupos? ===}\SpecialCharTok{\textbackslash{}n}\StringTok{"}\NormalTok{)}
\end{Highlighting}
\end{Shaded}

\begin{verbatim}

=== Prueba F global: ¿difieren los tres grupos? ===
\end{verbatim}

\begin{Shaded}
\begin{Highlighting}[]
\NormalTok{glance\_eca }\OtherTok{\textless{}{-}} \FunctionTok{glance}\NormalTok{(modelo\_eca)}
\FunctionTok{cat}\NormalTok{(}\StringTok{"F{-}estadístico ="}\NormalTok{, }\FunctionTok{round}\NormalTok{(glance\_eca}\SpecialCharTok{$}\NormalTok{statistic, }\DecValTok{2}\NormalTok{),}
    \StringTok{"|  gl ="}\NormalTok{, glance\_eca}\SpecialCharTok{$}\NormalTok{df, }\StringTok{","}\NormalTok{, glance\_eca}\SpecialCharTok{$}\NormalTok{nobs }\SpecialCharTok{{-}}\NormalTok{ glance\_eca}\SpecialCharTok{$}\NormalTok{df }\SpecialCharTok{{-}} \DecValTok{1}\NormalTok{,}
    \StringTok{"|  p \textless{} 0.001}\SpecialCharTok{\textbackslash{}n\textbackslash{}n}\StringTok{"}\NormalTok{)}
\end{Highlighting}
\end{Shaded}

\begin{verbatim}
F-estadístico = 31.6 |  gl = 2 , 87 |  p < 0.001
\end{verbatim}

\begin{Shaded}
\begin{Highlighting}[]
\CommentTok{\# Gráfico}
\FunctionTok{par}\NormalTok{(}\AttributeTok{mar =} \FunctionTok{c}\NormalTok{(}\FloatTok{4.5}\NormalTok{, }\FloatTok{4.5}\NormalTok{, }\DecValTok{3}\NormalTok{, }\FloatTok{1.5}\NormalTok{))}
\FunctionTok{boxplot}\NormalTok{(reduccion }\SpecialCharTok{\textasciitilde{}}\NormalTok{ tratamiento, }\AttributeTok{data =}\NormalTok{ datos\_eca,}
        \AttributeTok{main =} \StringTok{"Reducción de PAS por Tratamiento (ECA simulado)}\SpecialCharTok{\textbackslash{}n}\StringTok{n = 30 por grupo"}\NormalTok{,}
        \AttributeTok{ylab =} \StringTok{"Reducción de PAS (mmHg)"}\NormalTok{,}
        \AttributeTok{xlab =} \StringTok{"Tratamiento"}\NormalTok{,}
        \AttributeTok{col  =} \FunctionTok{c}\NormalTok{(}\StringTok{"lightgray"}\NormalTok{, }\StringTok{"\#5B9BD5"}\NormalTok{, }\StringTok{"\#ED7D31"}\NormalTok{),}
        \AttributeTok{las  =} \DecValTok{1}\NormalTok{, }\AttributeTok{notch =} \ConstantTok{FALSE}\NormalTok{)}
\FunctionTok{abline}\NormalTok{(}\AttributeTok{h =} \DecValTok{0}\NormalTok{, }\AttributeTok{lty =} \DecValTok{2}\NormalTok{, }\AttributeTok{col =} \StringTok{"red"}\NormalTok{)}

\CommentTok{\# Medias por grupo}
\NormalTok{medias }\OtherTok{\textless{}{-}} \FunctionTok{tapply}\NormalTok{(datos\_eca}\SpecialCharTok{$}\NormalTok{reduccion, datos\_eca}\SpecialCharTok{$}\NormalTok{tratamiento, mean)}
\FunctionTok{points}\NormalTok{(}\DecValTok{1}\SpecialCharTok{:}\DecValTok{3}\NormalTok{, medias, }\AttributeTok{pch =} \DecValTok{18}\NormalTok{, }\AttributeTok{cex =} \FloatTok{1.5}\NormalTok{, }\AttributeTok{col =} \StringTok{"red"}\NormalTok{)}
\end{Highlighting}
\end{Shaded}

\begin{figure}[H]

{\centering \pandocbounded{\includegraphics[keepaspectratio]{chapters/semana10_files/figure-pdf/lst-s11-cod07-1.pdf}}

}

\caption{Reducción de PAS (mmHg) por grupo de tratamiento en el ensayo
clínico simulado. Línea roja: media de cada grupo. El modelo de
regresión estima el efecto de cada tratamiento respecto al placebo.}

\end{figure}%

\begin{tcolorbox}[enhanced jigsaw, breakable, titlerule=0mm, rightrule=.15mm, leftrule=.75mm, left=2mm, opacitybacktitle=0.6, colback=white, bottomtitle=1mm, colbacktitle=quarto-callout-tip-color!10!white, title=\textcolor{quarto-callout-tip-color}{\faLightbulb}\hspace{0.5em}{Tip}, toptitle=1mm, coltitle=black, opacityback=0, arc=.35mm, bottomrule=.15mm, toprule=.15mm, colframe=quarto-callout-tip-color-frame]

\textbf{Interpretación estadística:} El análisis de varianza reveló
diferencias estadísticamente significativas en la reducción de presión
arterial entre los tres grupos de tratamiento (F(2,87) = 23.4, p
\textless{} 0.001). Los coeficientes de regresión muestran que, respecto
al placebo (reducción media \textasciitilde4 mmHg), tanto el tratamiento
IECA (reducción adicional \textasciitilde12 mmHg) como el ARA-II
(reducción adicional \textasciitilde11 mmHg) produjeron mejoras
clínicamente relevantes y estadísticamente significativas. La magnitud
del efecto es importante para la práctica clínica: ambos
antihipertensivos superan al placebo en aproximadamente 10-12 mmHg, una
diferencia clínicamente significativa en el manejo de la hipertensión
arterial.

\end{tcolorbox}

\subsection{\texorpdfstring{Ejemplo con \texttt{BioStatR}: regresión
lineal
múltiple}{Ejemplo con BioStatR: regresión lineal múltiple}}\label{ejemplo-con-biostatr-regresiuxf3n-lineal-muxfaltiple}

La función \texttt{rlm()} del paquete \texttt{BioStatR} (ver
Apéndice~\ref{sec-apendiceB}) extiende \texttt{rls()} al caso
multivariante. Integra en una sola llamada la estimación de los
coeficientes con intervalos de confianza, el coeficiente de
determinación \(R^2\) y \(R^2\) ajustado, el test de normalidad de los
residuos (Shapiro--Wilk) y los gráficos de diagnóstico (residuos
\emph{vs.} ajustados, Q--Q plot e histograma de residuos
estandarizados).

\begin{tcolorbox}[enhanced jigsaw, breakable, titlerule=0mm, rightrule=.15mm, leftrule=.75mm, left=2mm, opacitybacktitle=0.6, colback=white, bottomtitle=1mm, colbacktitle=quarto-callout-tip-color!10!white, title=\textcolor{quarto-callout-tip-color}{\faLightbulb}\hspace{0.5em}{Modelo: HbA1c \textasciitilde{} tiempo de evolución + edad + IMC}, toptitle=1mm, coltitle=black, opacityback=0, arc=.35mm, bottomrule=.15mm, toprule=.15mm, colframe=quarto-callout-tip-color-frame]

Aplicamos \texttt{rlm()} al conjunto de datos \texttt{osteo} para
evaluar cómo el tiempo de evolución de la diabetes, la edad y el índice
de masa corporal (IMC) explican conjuntamente la hemoglobina glicosilada
(HbA1c).

\begin{Shaded}
\begin{Highlighting}[]
\FunctionTok{library}\NormalTok{(BioStatR)}
\FunctionTok{data}\NormalTok{(osteo)}

\FunctionTok{rlm}\NormalTok{(hba1c }\SpecialCharTok{\textasciitilde{}}\NormalTok{ tevol }\SpecialCharTok{+}\NormalTok{ edad }\SpecialCharTok{+}\NormalTok{ imc, }\AttributeTok{data =}\NormalTok{ osteo, }\AttributeTok{grf =} \ConstantTok{FALSE}\NormalTok{)}
\end{Highlighting}
\end{Shaded}

\begin{verbatim}

Regresión lineal múltiple 
----------------------------------------------------------------
# Información muestral --- 

      Variable  n Media   DT  Min  Max
hba1c    hba1c 94  8.56 1.80  4.6 13.8
tevol    tevol 94 12.33 8.53  0.0 35.0
edad      edad 94 30.19 9.37 18.0 56.0
imc        imc 94 23.92 3.75 18.1 37.3

# Modelo lineal --- 

   Modelo :  hba1c ~ tevol + edad + imc 
  R² =  0.063  (R²  ajustado  =  0.031 )
  S²residual =  3.13 

   Coeficientes del modelo :

      Termino Estimacion Error_Std IC_inf IC_sup  t_exp      sig
1 (Intercept)      8.601     1.200  6.217 10.984  7.169  < 0.001
2       tevol     -0.046     0.025 -0.095  0.003 -1.883  = 0.063
3        edad     -0.012     0.024 -0.059  0.035 -0.521  = 0.604
4         imc      0.038     0.052 -0.066  0.142  0.722  = 0.472

# Distribución residual --- 
  Error estándar residual:  1.77 
    Residuos Res_Est
min   -4.519  -2.589
Q1    -0.945  -0.545
Q2    -0.010  -0.006
Q3     0.995   0.572
max    4.938   2.923

  Test de normalidad residual (Shapiro-Wilk): 
   w = 0.991 ,   = 0.744 
\end{verbatim}

\textbf{Interpretación:} Tras ajustar simultáneamente por la edad y el
IMC, el coeficiente del tiempo de evolución (\texttt{tevol}) es negativo
y cercano a la significación estadística (\(p \approx 0.06\)). Esto
sugiere que, controlando por la edad y la composición corporal, cada año
adicional de evolución de la diabetes se asocia con una ligera reducción
de la HbA1c, posiblemente reflejo de una mejor adherencia terapéutica en
pacientes con mayor experiencia clínica. El coeficiente de determinación
ajustado (\(R^2_{\text{adj}} \approx 0.03\)) indica que estas tres
variables explican apenas una pequeña fracción de la variabilidad de la
HbA1c, lo que es coherente con la naturaleza multifactorial del control
glucémico.

El test de Shapiro--Wilk sobre los residuos no rechaza la normalidad,
validando los intervalos de confianza y los contrastes individuales.

\end{tcolorbox}

\begin{tcolorbox}[enhanced jigsaw, breakable, titlerule=0mm, rightrule=.15mm, leftrule=.75mm, left=2mm, opacitybacktitle=0.6, colback=white, bottomtitle=1mm, colbacktitle=quarto-callout-note-color!10!white, title=\textcolor{quarto-callout-note-color}{\faInfo}\hspace{0.5em}{Diagnóstico gráfico con \texttt{grf\ =\ TRUE}}, toptitle=1mm, coltitle=black, opacityback=0, arc=.35mm, bottomrule=.15mm, toprule=.15mm, colframe=quarto-callout-note-color-frame]

Llamando a \texttt{rlm()} con \texttt{grf\ =\ TRUE} se obtienen los
cuatro gráficos clásicos de diagnóstico (residuos \emph{vs.} ajustados,
Q--Q plot de residuos estandarizados, escala-localización y residuos
\emph{vs.} leverage), junto con un histograma de residuos
estandarizados:

\begin{Shaded}
\begin{Highlighting}[]
\FunctionTok{rlm}\NormalTok{(hba1c }\SpecialCharTok{\textasciitilde{}}\NormalTok{ tevol }\SpecialCharTok{+}\NormalTok{ edad }\SpecialCharTok{+}\NormalTok{ imc, }\AttributeTok{data =}\NormalTok{ osteo, }\AttributeTok{grf =} \ConstantTok{TRUE}\NormalTok{)}
\end{Highlighting}
\end{Shaded}

\begin{verbatim}

Regresión lineal múltiple 
----------------------------------------------------------------
# Información muestral --- 

      Variable  n Media   DT  Min  Max
hba1c    hba1c 94  8.56 1.80  4.6 13.8
tevol    tevol 94 12.33 8.53  0.0 35.0
edad      edad 94 30.19 9.37 18.0 56.0
imc        imc 94 23.92 3.75 18.1 37.3

# Modelo lineal --- 

   Modelo :  hba1c ~ tevol + edad + imc 
  R² =  0.063  (R²  ajustado  =  0.031 )
  S²residual =  3.13 

   Coeficientes del modelo :

      Termino Estimacion Error_Std IC_inf IC_sup  t_exp      sig
1 (Intercept)      8.601     1.200  6.217 10.984  7.169  < 0.001
2       tevol     -0.046     0.025 -0.095  0.003 -1.883  = 0.063
3        edad     -0.012     0.024 -0.059  0.035 -0.521  = 0.604
4         imc      0.038     0.052 -0.066  0.142  0.722  = 0.472

# Distribución residual --- 
  Error estándar residual:  1.77 
    Residuos Res_Est
min   -4.519  -2.589
Q1    -0.945  -0.545
Q2    -0.010  -0.006
Q3     0.995   0.572
max    4.938   2.923

  Test de normalidad residual (Shapiro-Wilk): 
   w = 0.991 ,   = 0.744 
\end{verbatim}

\pandocbounded{\includegraphics[keepaspectratio]{chapters/semana10_files/figure-pdf/lst-s10-bio-rlm-diag-1.pdf}}

\pandocbounded{\includegraphics[keepaspectratio]{chapters/semana10_files/figure-pdf/lst-s10-bio-rlm-diag-2.pdf}}

\pandocbounded{\includegraphics[keepaspectratio]{chapters/semana10_files/figure-pdf/lst-s10-bio-rlm-diag-3.pdf}}

\pandocbounded{\includegraphics[keepaspectratio]{chapters/semana10_files/figure-pdf/lst-s10-bio-rlm-diag-4.pdf}}

\end{tcolorbox}

\section{Resumen}\label{sec-semana10-resumen}

\begin{tcolorbox}[enhanced jigsaw, breakable, titlerule=0mm, rightrule=.15mm, leftrule=.75mm, left=2mm, opacitybacktitle=0.6, colback=white, bottomtitle=1mm, colbacktitle=quarto-callout-note-color!10!white, title=\textcolor{quarto-callout-note-color}{\faInfo}\hspace{0.5em}{Fórmulas Clave}, toptitle=1mm, coltitle=black, opacityback=0, arc=.35mm, bottomrule=.15mm, toprule=.15mm, colframe=quarto-callout-note-color-frame]

\textbf{Modelo y Estimación:}
\[Y_i = \beta_0 + \beta_1 X_{1i} + \cdots + \beta_p X_{pi} + U_i\]
\[\hat{\boldsymbol{\beta}} = (\mathbf{X}^T\mathbf{X})^{-1}\mathbf{X}^T\mathbf{Y}\]

\textbf{Varianza de Coeficientes:}
\[\text{Var}(\hat{\boldsymbol{\beta}}) = \sigma_U^2(\mathbf{X}^T\mathbf{X})^{-1}\]
\[s^2 = \frac{\text{RSS}}{n-p-1}\]

\textbf{Bondad de Ajuste:} \[R^2 = 1 - \frac{\text{RSS}}{\text{SST}}\]
\[R^2_{\text{adj}} = 1 - (1-R^2)\frac{n-1}{n-p-1}\]

\textbf{Tests e Intervalos:}
\[t_j = \frac{\hat{\beta}_j}{\hat{\text{SE}}(\hat{\beta}_j)} \sim t_{n-p-1}\]
\[F = \frac{(\text{RSS}_1 - \text{RSS}_2)/(p_2 - p_1)}{\text{RSS}_2/(n-p_2-1)} \sim F_{p_2-p_1, n-p_2-1}\]

\textbf{Criterios de Información:}
\[\text{AIC} = n \log(\text{RSS}/n) + 2(p+1)\]
\[\text{BIC} = n \log(\text{RSS}/n) + (p+1)\log(n)\]

Seleccionar el modelo con \textbf{menor} AIC o BIC.

\end{tcolorbox}

\section{Ejercicios}\label{sec-semana10-ejercicios}

\begin{enumerate}
\def\labelenumi{\arabic{enumi}.}
\item
  \textbf{Interpretación de Coeficientes Parciales:} En un modelo que
  predice la \textbf{hemoglobina glicosilada} (HbA1c, \%) usando años de
  evolución de la diabetes, edad del paciente e IMC, el coeficiente de
  los años de evolución es 0.05. ¿Qué significa este número? ¿Cómo
  interpretaría el coeficiente del IMC?
\item
  \textbf{Matriz de Diseño:} Un estudio epidemiológico analiza la
  \textbf{presión arterial diastólica} con 3 variables predictoras
  (edad, IMC y tabaquismo en paquetes-año) y 80 observaciones. ¿Cuál es
  la dimensión de la matriz \(\mathbf{X}\)? ¿Qué contiene la primera
  columna? ¿Cuántos grados de libertad tienen los residuos?
\item
  \textbf{Problema de Multicolinealidad:} En un modelo de riesgo
  cardiovascular se incluyen simultáneamente el colesterol LDL (mg/dL) y
  el colesterol no-HDL (mg/dL). ¿Qué problema clínico y estadístico
  espera encontrar? ¿Cómo lo detectaría con el VIF y cómo lo resolvería?
\item
  \textbf{\(R^2\) vs.~\(R^2_{\text{adj}}\):} Un investigador ajusta un
  modelo de predicción de supervivencia en cáncer de mama e incluye
  secuencialmente 15 variables (estadio, ganglios, Ki67, HER2, RE, RP,
  tamaño tumoral, edad, tratamiento, comorbilidades, etc.). ¿Por qué
  \(R^2\) siempre aumenta al añadir cada variable? ¿Cuál criterio (AIC,
  BIC o \(R^2_{\text{adj}}\)) es más apropiado para seleccionar el
  modelo final?
\item
  \textbf{Prueba F para Modelos Anidados:} En un estudio de HbA1c se
  ajustan dos modelos:

  \begin{itemize}
  \tightlist
  \item
    Modelo A:
    \(\text{HbA1c} = \beta_0 + \beta_1 \cdot \text{tevol} + U\)
  \item
    Modelo B:
    \(\text{HbA1c} = \beta_0 + \beta_1 \cdot \text{tevol} + \beta_2 \cdot \text{Edad} + \beta_3 \cdot \text{IMC} + U\)
  \end{itemize}

  Se obtiene \(\text{RSS}_A = 188.4\), \(\text{RSS}_B = 172.6\),
  \(n = 94\).

  \begin{enumerate}
  \def\labelenumii{\alph{enumii})}
  \tightlist
  \item
    ¿Qué hipótesis nula se está probando con la prueba F?
  \item
    Calcula el estadístico F
  \item
    Con \(F_{0.05, 2, 90} \approx 3.10\), ¿es el modelo B
    significativamente mejor? Interpreta clínicamente
  \end{enumerate}
\item
  \textbf{Variables Indicadoras en Ensayos Clínicos:} En un ensayo
  clínico de 4 brazos se comparan tres tratamientos antidiabéticos
  (metformina, sitagliptina, empagliflozina) frente a placebo. ¿Cuántas
  variables indicadoras necesitas? Si el coeficiente estimado de
  empagliflozina es \(-0.72\%\) de HbA1c con \(p = 0.003\), interpreta
  este resultado en el contexto del ensayo. ¿Cuál es la categoría de
  referencia?
\end{enumerate}

\section{Respuestas a los Ejercicios}\label{sec-semana10-respuestas}

\textbf{Ejercicio 1:} El coeficiente de años de evolución (0.05)
significa que, manteniendo la edad del paciente y el IMC constantes,
cada año adicional de evolución de la diabetes se asocia con un aumento
promedio de 0.05 puntos porcentuales en la HbA1c. El coeficiente del IMC
representaría el cambio esperado en HbA1c por cada kg/m² adicional de
IMC, ajustando por los años de evolución y la edad: si es positivo,
mayor obesidad se asocia con peor control glucémico.

\textbf{Ejercicio 2:} Dimensión de \(\mathbf{X}\): \(80 \times 4\) (80
observaciones, 1 columna de unos para el intercepto + 3 columnas de
predictores). La primera columna contiene todos los 1s (para estimar
\(\beta_0\)). Los grados de libertad de los residuos son
\(n - p - 1 = 80 - 3 - 1 = 76\).

\textbf{Ejercicio 3:} El colesterol no-HDL = Colesterol Total − HDL, por
lo que está determinísticamente relacionado con el LDL más el VLDL. Si
ambos entran en el modelo, hay multicolinealidad severa (potencialmente
casi perfecta). Consecuencias: errores estándar muy grandes,
coeficientes inestables, posibles signos contraintuitivos. Detección:
VIF tendería a infinito o sería muy alto (\textgreater10). Solución:
utilizar solo uno de los dos indicadores lipídicos, preferiblemente el
LDL por su mayor validación clínica, o usar el colesterol no-HDL como
alternativa cuando el LDL no es medible.

\textbf{Ejercicio 4:} \(R^2\) siempre aumenta porque es el ratio
SCE/SCT, y añadir variables (aunque sean ruido puro) solo puede
disminuir o mantener el RSS, nunca aumentarlo. El modelo siempre
``aprovecha'' cualquier correlación espuria con los datos de muestra.
Para 15 variables y \(n\) no muy grande, el \(R^2\) puede estar muy
inflado. El BIC es el criterio más apropiado en investigación
confirmatoria: penaliza más fuertemente la complejidad (\(\log n\) vs.~2
en AIC), favoreciendo modelos parsimoniosos. El \(R^2_{\text{adj}}\)
también es útil como referencia, pero no tiene la misma base
probabilística.

\textbf{Ejercicio 5:} - a) \(H_0: \beta_2 = \beta_3 = 0\) (la edad y el
IMC no aportan información sobre la HbA1c más allá del tiempo de
evolución) - b)
\(F = \frac{(\text{RSS}_A - \text{RSS}_B)/(p_2 - p_1)}{\text{RSS}_B/(n - p_2 - 1)} = \frac{(188.4 - 172.6)/2}{172.6/90} = \frac{15.8/2}{1.918} = \frac{7.9}{1.918} = 4.12\)
- c) \(F = 4.12 > F_{0.05, 2, 90} = 3.10\), por lo que rechazamos
\(H_0\) (\(p < 0.05\)). El modelo B es significativamente mejor.
Clínicamente: la edad y el IMC aportan información independiente sobre
el control glucémico, más allá de los años de evolución de la diabetes.
Esto sugiere que el tratamiento de la hiperglucemia debe considerar
también la edad del paciente y su estado ponderal.

\textbf{Ejercicio 6:} Se necesitan \textbf{3 variables indicadoras}
(\(m - 1 = 4 - 1 = 3\)): una para metformina, otra para sitagliptina y
otra para empagliflozina. El placebo es la categoría de referencia. El
coeficiente de empagliflozina (\(-0.72\%\), \(p = 0.003\)) significa
que, en promedio, los pacientes tratados con empagliflozina tienen una
HbA1c 0.72 puntos porcentuales menor que los pacientes con placebo, y
esta diferencia es estadísticamente significativa. En términos clínicos,
una reducción de \textasciitilde0.7\% en HbA1c es clínicamente relevante
(el umbral suele ser ≥0.5\%).

\begin{center}\rule{0.5\linewidth}{0.5pt}\end{center}

\begin{tcolorbox}[enhanced jigsaw, breakable, titlerule=0mm, rightrule=.15mm, leftrule=.75mm, left=2mm, opacitybacktitle=0.6, colback=white, bottomtitle=1mm, colbacktitle=quarto-callout-tip-color!10!white, title=\textcolor{quarto-callout-tip-color}{\faLightbulb}\hspace{0.5em}{Métodos Avanzados}, toptitle=1mm, coltitle=black, opacityback=0, arc=.35mm, bottomrule=.15mm, toprule=.15mm, colframe=quarto-callout-tip-color-frame]

Para ampliar los contenidos de este capítulo con técnicas estadísticas
avanzadas, visita:

\textbf{\href{https://migariane.github.io/Bioestadistica}{→
Bioestadística Avanzada --- M.A.~Luque Fernández}}

\end{tcolorbox}

\part{Parte V: Análisis de Datos Categóricos}

\chapter{Semana 11 --- Análisis de Datos
Categóricos}\label{sec-semana11}

En estadística médica y epidemiológica, muchas variables de interés son
categóricas: enfermedad (sí/no), exposición a un factor de riesgo
(presente/ausente), grupo de tratamiento, estadio clínico. Este capítulo
introduce las herramientas para analizar este tipo de datos: las pruebas
Chi-cuadrado de independencia y homogeneidad, la prueba de McNemar para
datos apareados, las medidas de asociación en tablas de contingencia
según el diseño del estudio (cohorte, transversal, caso-control), y los
métodos exactos cuando los supuestos paramétricos no se cumplen.

\begin{center}\rule{0.5\linewidth}{0.5pt}\end{center}

\section{Variables Categóricas y Tablas de
Contingencia}\label{sec-semana11-tablas}

Una \textbf{tabla de contingencia} es una representación matricial de
las frecuencias conjuntas de dos o más variables categóricas. Es la base
del análisis de datos categóricos.

\subsection{Tabla 2×2: Estructura y Notación
Epidemiológica}\label{tabla-22-estructura-y-notaciuxf3n-epidemioluxf3gica}

La tabla 2×2 es el caso más frecuente en epidemiología y ensayos
clínicos:

\begin{longtable}[]{@{}lccc@{}}
\toprule\noalign{}
& \textbf{Resultado: Sí} & \textbf{Resultado: No} & \textbf{Total} \\
\midrule\noalign{}
\endhead
\bottomrule\noalign{}
\endlastfoot
\textbf{Exposición: Sí} & \(a\) & \(b\) & \(a+b\) \\
\textbf{Exposición: No} & \(c\) & \(d\) & \(c+d\) \\
\textbf{Total} & \(a+c\) & \(b+d\) & \(n\) \\
\end{longtable}

Donde: - \(a\) = expuestos con resultado positivo (casos expuestos) -
\(b\) = expuestos sin resultado - \(c\) = no expuestos con resultado
positivo - \(d\) = no expuestos sin resultado - \(n = a + b + c + d\) =
total de observaciones

\subsection{Tablas r×c Generalizadas}\label{tablas-rc-generalizadas}

Para \(r\) categorías de fila y \(c\) categorías de columna, la tabla es
de dimensión \(r \times c\) con frecuencias observadas \(O_{ij}\),
marginales de fila \(O_{i\cdot} = \sum_j O_{ij}\) y de columna
\(O_{\cdot j} = \sum_i O_{ij}\).

\begin{center}\rule{0.5\linewidth}{0.5pt}\end{center}

\section{Prueba χ² de Independencia}\label{sec-semana11-independencia}

La prueba Chi-cuadrado de independencia contrasta si dos variables
categóricas son estadísticamente independientes en la población.

\subsection{Hipótesis}\label{hipuxf3tesis-1}

\[H_0: \text{las variables son independientes} \quad \Leftrightarrow \quad P(X = i,\, Y = j) = P(X = i)\,P(Y = j)\]
\[H_1: \text{las variables no son independientes (están asociadas)}\]

\subsection{\texorpdfstring{Frecuencias Esperadas Bajo
\(H_0\)}{Frecuencias Esperadas Bajo H\_0}}\label{frecuencias-esperadas-bajo-h_0}

Si \(H_0\) es verdadera, las frecuencias esperadas son:

\[E_{ij} = \frac{O_{i\cdot}\, \cdot\, O_{\cdot j}}{n}\]

\begin{tcolorbox}[enhanced jigsaw, breakable, titlerule=0mm, rightrule=.15mm, leftrule=.75mm, left=2mm, opacitybacktitle=0.6, colback=white, bottomtitle=1mm, colbacktitle=quarto-callout-warning-color!10!white, title=\textcolor{quarto-callout-warning-color}{\faExclamationTriangle}\hspace{0.5em}{Estadístico: Pearson χ²}, toptitle=1mm, coltitle=black, opacityback=0, arc=.35mm, bottomrule=.15mm, toprule=.15mm, colframe=quarto-callout-warning-color-frame]

\[\chi^2 = \sum_{i=1}^r \sum_{j=1}^c \frac{(O_{ij} - E_{ij})^2}{E_{ij}}\]

Bajo \(H_0\): \(\chi^2 \sim \chi^2_{(r-1)(c-1)}\) (asintóticamente).

Para tabla 2×2: \((r-1)(c-1) = 1\) grado de libertad.

\textbf{Condición de validez:} Todas las frecuencias esperadas deben ser
\(E_{ij} \geq 5\). Si alguna es menor, usar la corrección de Yates o la
prueba exacta de Fisher.

\end{tcolorbox}

\subsection{Cálculo Manual para Tabla
2×2}\label{cuxe1lculo-manual-para-tabla-22}

Para tabla 2×2 con notación \((a, b, c, d)\):

\[\chi^2 = \frac{n(ad - bc)^2}{(a+b)(c+d)(a+c)(b+d)}\]

\subsection{Ejemplo Médico: Tabaquismo y
Osteoporosis}\label{ejemplo-muxe9dico-tabaquismo-y-osteoporosis}

\begin{tcolorbox}[enhanced jigsaw, breakable, titlerule=0mm, rightrule=.15mm, leftrule=.75mm, left=2mm, opacitybacktitle=0.6, colback=white, bottomtitle=1mm, colbacktitle=quarto-callout-tip-color!10!white, title=\textcolor{quarto-callout-tip-color}{\faLightbulb}\hspace{0.5em}{Ejemplo: tabla2x2() --- datos osteo (Fac. Medicina UGR)}, toptitle=1mm, coltitle=black, opacityback=0, arc=.35mm, bottomrule=.15mm, toprule=.15mm, colframe=quarto-callout-tip-color-frame]

Analizamos la asociación entre tabaquismo (\texttt{tabaco}) y
osteoporosis de cúbito (\texttt{osteo\_cue}) en 94 pacientes diabéticos.
Documentación completa en Sección~\ref{sec-bio-tablas}.

\begin{verbatim}

# Análisis de tablas 2x2
# ----------------------

# Frecuencias observadas
               Sin osteoporosis   Con osteoporosis Total
  No fumador                 44                  9    53
  Fumador                    26                 15    41
  Total                      70                 24    94


# Test Chi-cuadrado para un estudio indefinido

  χ² = 4.651,   gl = 1,  p = 0.031, (cpc = 1) 
  Validez: Frecuencia mínima esperada = 10.47  > 7.7  

  Test exacto de Fisher (bilateral): p = 0.035 

  --- Otros criterios χ²:  
  χ² = 4.673,   gl = 1,  p = 0.054, (sin cpc) 
  χ² = 3.699,   gl = 1,  p = 0.054, (cpc de Yates = 47.00) 

# Medidas de asociación para un estudio indefinido

  Razón del producto cruzado (odds ratio):
  OR=2.821; 95%-IC(OR)= (1.070, 7.013) 
\end{verbatim}

\textbf{Comprobación manual de \(E_{ij}\):}
\[E_{11} = \frac{53 \times 70}{94} = 39.47, \quad E_{12} = \frac{53 \times 24}{94} = 13.53\]
\[E_{21} = \frac{41 \times 70}{94} = 30.53, \quad E_{22} = \frac{41 \times 24}{94} = 10.47\]

Mínima frecuencia esperada: \(E_{22} = 10.47 > 5\) ✓ (condición de
validez cumplida).

\begin{quote}
\textbf{Interpretación}

El estadístico χ² = 4.65 con 1 grado de libertad produce p = 0.031
\textless{} 0.05, rechazando la hipótesis nula de independencia. Existe
asociación estadísticamente significativa entre tabaquismo y
osteoporosis en esta muestra de 94 pacientes diabéticos. Los fumadores
presentan una odds ratio de 2.82 (IC 95\%: 1.07--7.01), indicando
aproximadamente 2.8 veces mayor probabilidad de osteoporosis de cúbito
comparado con no fumadores. En el contexto clínico, esta asociación es
relevante para estratificar riesgo en pacientes diabéticos, sugiriendo
que el cese del tabaquismo podría ser una intervención preventiva
importante en esta población de alto riesgo metabólico.
\end{quote}

\end{tcolorbox}

\begin{center}\rule{0.5\linewidth}{0.5pt}\end{center}

\section{Prueba χ² de Homogeneidad}\label{sec-semana11-homogeneidad}

La prueba de homogeneidad contrasta si la distribución de una variable
categórica es \textbf{la misma en \(k\) poblaciones o grupos}
predefinidos. Matemáticamente usa el mismo estadístico \(\chi^2\), pero
el diseño del estudio es diferente.

\subsection{Independencia
vs.~Homogeneidad}\label{independencia-vs.-homogeneidad}

\begin{longtable}[]{@{}
  >{\raggedright\arraybackslash}p{(\linewidth - 4\tabcolsep) * \real{0.2432}}
  >{\raggedright\arraybackslash}p{(\linewidth - 4\tabcolsep) * \real{0.3784}}
  >{\raggedright\arraybackslash}p{(\linewidth - 4\tabcolsep) * \real{0.3784}}@{}}
\toprule\noalign{}
\begin{minipage}[b]{\linewidth}\raggedright
Aspecto
\end{minipage} & \begin{minipage}[b]{\linewidth}\raggedright
Independencia
\end{minipage} & \begin{minipage}[b]{\linewidth}\raggedright
Homogeneidad
\end{minipage} \\
\midrule\noalign{}
\endhead
\bottomrule\noalign{}
\endlastfoot
\textbf{Muestreo} & Muestra única & \(k\) muestras independientes \\
\textbf{\(H_0\)} & Independencia entre variables & Igual distribución en
\(k\) grupos \\
\textbf{Ejemplo} & ¿Tabaquismo y cáncer relacionados? & ¿Mismo \% de
curación en 3 tratamientos? \\
\end{longtable}

\subsection{Procedimiento}\label{procedimiento-1}

Se aplica el mismo estadístico
\(\chi^2 = \sum \frac{(O_{ij}-E_{ij})^2}{E_{ij}}\) con \((r-1)(c-1)\)
grados de libertad.

\begin{tcolorbox}[enhanced jigsaw, breakable, titlerule=0mm, rightrule=.15mm, leftrule=.75mm, left=2mm, opacitybacktitle=0.6, colback=white, bottomtitle=1mm, colbacktitle=quarto-callout-tip-color!10!white, title=\textcolor{quarto-callout-tip-color}{\faLightbulb}\hspace{0.5em}{Ejemplo: tablarxc() --- eficacia de tratamiento en 3 grupos}, toptitle=1mm, coltitle=black, opacityback=0, arc=.35mm, bottomrule=.15mm, toprule=.15mm, colframe=quarto-callout-tip-color-frame]

Ensayo clínico con 3 grupos de tratamiento (Control, Dosis baja, Dosis
alta) y resultado binario (Curado/No curado):

\begin{verbatim}

# Test Chi-cuadrado para tablas RxC
# ---------------------------------

# Frecuencias observadas
               Curado   No_curado Total
  Control          20          30    50
  Dosis_baja       35          25    60
  Dosis_alta       45          15    60
  Total           100          70   170

# Test chi-cuadrado 
  Validez: Frecuencia mínima esperada = 20.59 
           0 frecuencias esperadas son menores a 1
           0 son menores a 5 (el 0% de la tabla) 
  χ²(2 gl) = 13.802, p = 0.001 
\end{verbatim}

\begin{quote}
\textbf{Interpretación}

El estadístico χ² = 13.38 con 2 grados de libertad produce p \textless{}
0.001, rechazando la hipótesis nula de homogeneidad entre grupos. La
distribución de curación no es homogénea entre los 3 grupos de
tratamiento (Control: 40\%, Dosis baja: 58\%, Dosis alta: 75\%). Los
resultados demuestran que la dosis alta de tratamiento se asocia
significativamente con tasas de curación más altas comparado con control
y dosis baja. Este patrón dosis-respuesta es clínicamente relevante y
sugiere que la escalada de dosis podría optimizar eficacia, aunque el
balance entre beneficio y tolerabilidad debe evaluarse conjuntamente en
la toma de decisiones terapéuticas.
\end{quote}

\end{tcolorbox}

\begin{center}\rule{0.5\linewidth}{0.5pt}\end{center}

\section{Condiciones de Aplicación y
Alternativas}\label{sec-semana11-condiciones}

\subsection{Condición de Frecuencias
Esperadas}\label{condiciuxf3n-de-frecuencias-esperadas}

La regla clásica de Cochran (1954) exige, para que la aproximación
\(\chi^2 \sim \chi^2_{(r-1)(c-1)}\) sea válida, que \textbf{todas} las
frecuencias esperadas \(E_{ij}\) sean \(\geq 5\) en una tabla 2×2 y que,
en tablas \(r\times c\), al menos el 80\% de las \(E_{ij}\) sean
\(\geq 5\) y ninguna \(< 1\).

Sin embargo, este umbral único es excesivamente conservador. La escuela
granadina de bioestadística (\citet{MartinAndresLuna2004};
\citet{MartinAndres1994}) ha demostrado, mediante estudios de simulación
extensivos, que el valor mínimo de \(E_{ij}\) que mantiene un error de
tipo I cercano al nominal \textbf{depende del tamaño muestral y del
diseño del estudio}:

\begin{tcolorbox}[enhanced jigsaw, breakable, titlerule=0mm, rightrule=.15mm, leftrule=.75mm, left=2mm, opacitybacktitle=0.6, colback=white, bottomtitle=1mm, colbacktitle=quarto-callout-important-color!10!white, title=\textcolor{quarto-callout-important-color}{\faExclamation}\hspace{0.5em}{Valores mínimos de \(E_{ij}\) recomendados --- Tablas 2×2}, toptitle=1mm, coltitle=black, opacityback=0, arc=.35mm, bottomrule=.15mm, toprule=.15mm, colframe=quarto-callout-important-color-frame]

\begin{longtable}[]{@{}
  >{\raggedright\arraybackslash}p{(\linewidth - 6\tabcolsep) * \real{0.2500}}
  >{\raggedright\arraybackslash}p{(\linewidth - 6\tabcolsep) * \real{0.2500}}
  >{\raggedright\arraybackslash}p{(\linewidth - 6\tabcolsep) * \real{0.2500}}
  >{\raggedright\arraybackslash}p{(\linewidth - 6\tabcolsep) * \real{0.2500}}@{}}
\caption{Umbrales para la validez del \(\chi^2\) en tablas 2×2 según
diseño y tamaño muestral. Adaptado de \citet{MartinAndresLuna2004} y
\citet{MartinAndres1994}.}\tabularnewline
\toprule\noalign{}
\begin{minipage}[b]{\linewidth}\raggedright
Tamaño muestral \(n\)
\end{minipage} & \begin{minipage}[b]{\linewidth}\raggedright
Diseño con un margen fijo (cohortes, caso-control)
\end{minipage} & \begin{minipage}[b]{\linewidth}\raggedright
Diseño con ambos márgenes libres (transversal, prevalencia)
\end{minipage} & \begin{minipage}[b]{\linewidth}\raggedright
Prueba recomendada
\end{minipage} \\
\midrule\noalign{}
\endfirsthead
\toprule\noalign{}
\begin{minipage}[b]{\linewidth}\raggedright
Tamaño muestral \(n\)
\end{minipage} & \begin{minipage}[b]{\linewidth}\raggedright
Diseño con un margen fijo (cohortes, caso-control)
\end{minipage} & \begin{minipage}[b]{\linewidth}\raggedright
Diseño con ambos márgenes libres (transversal, prevalencia)
\end{minipage} & \begin{minipage}[b]{\linewidth}\raggedright
Prueba recomendada
\end{minipage} \\
\midrule\noalign{}
\endhead
\bottomrule\noalign{}
\endlastfoot
\(n \leq 20\) & --- & --- & \textbf{Fisher exacta} (siempre) \\
\(20 < n \leq 40\) & \(E_{\min} \geq 5\) & \(E_{\min} \geq 5\) &
\(\chi^2\) con corrección de Yates \\
\(40 < n \leq 100\) & \(E_{\min} \geq 3\) & \(E_{\min} \geq 4\) &
\(\chi^2\) de Pearson (sin corrección) \\
\(100 < n \leq 200\) & \(E_{\min} \geq 2\) & \(E_{\min} \geq 3\) &
\(\chi^2\) de Pearson \\
\(n > 200\) & \(E_{\min} \geq 1\) & \(E_{\min} \geq 2\) & \(\chi^2\) de
Pearson \\
Ambos márgenes fijos (raro) & --- & --- & \textbf{Fisher exacta} (test
condicional óptimo) \\
\end{longtable}

\end{tcolorbox}

\begin{tcolorbox}[enhanced jigsaw, breakable, titlerule=0mm, rightrule=.15mm, leftrule=.75mm, left=2mm, opacitybacktitle=0.6, colback=white, bottomtitle=1mm, colbacktitle=quarto-callout-note-color!10!white, title=\textcolor{quarto-callout-note-color}{\faInfo}\hspace{0.5em}{Cuándo el diseño es de ``margen fijo'' vs.~``márgenes libres''}, toptitle=1mm, coltitle=black, opacityback=0, arc=.35mm, bottomrule=.15mm, toprule=.15mm, colframe=quarto-callout-note-color-frame]

\begin{itemize}
\tightlist
\item
  \textbf{Un margen fijo} (lo más habitual en epidemiología): el
  investigador fija a priori el tamaño de los grupos de comparación.

  \begin{itemize}
  \tightlist
  \item
    \emph{Cohortes}: se fija el número de expuestos y no expuestos; se
    observa la incidencia.
  \item
    \emph{Caso-control}: se fija el número de casos y controles; se
    observa la exposición.
  \end{itemize}
\item
  \textbf{Ambos márgenes libres}: el muestreo es transversal/de
  prevalencia y solo se fija \(n\) total; los dos márgenes (exposición y
  enfermedad) son aleatorios.
\item
  \textbf{Ambos márgenes fijos}: situación poco frecuente en
  investigación clínica (p.ej., experimento del té de Fisher). Es la
  única en la que la \textbf{prueba exacta de Fisher} es el test
  condicional óptimo desde un punto de vista frecuentista.
\end{itemize}

Para tablas \(r\times c\) con \(r\) o \(c > 2\), la regla práctica
recomendada por estos autores es: \(\geq 80\%\) de las \(E_{ij}\) por
encima del umbral correspondiente al \(n\) total (columna apropiada), y
ninguna \(E_{ij} < 1\). Para detalles teóricos y simulaciones que
sustentan estos umbrales, véase \citet{MartinAndresLuna2004}.

\end{tcolorbox}

\subsection{Corrección de Continuidad de Yates (Tabla
2×2)}\label{sec-semana11-yates}

El estadístico \(\chi^2\) de Pearson es una suma de términos calculados
sobre frecuencias enteras que se aproxima por la distribución continua
\(\chi^2_{(r-1)(c-1)}\). La \textbf{corrección de continuidad de Yates}
reduce el estadístico para mejorar la aproximación:

\begin{tcolorbox}[enhanced jigsaw, breakable, titlerule=0mm, rightrule=.15mm, leftrule=.75mm, left=2mm, opacitybacktitle=0.6, colback=white, bottomtitle=1mm, colbacktitle=quarto-callout-note-color!10!white, title=\textcolor{quarto-callout-note-color}{\faInfo}\hspace{0.5em}{Estadístico: \(\chi^2\) de Yates (con corrección de continuidad)}, toptitle=1mm, coltitle=black, opacityback=0, arc=.35mm, bottomrule=.15mm, toprule=.15mm, colframe=quarto-callout-note-color-frame]

Para una tabla 2×2 con frecuencias \(a, b, c, d\):

\[\chi^2_{\text{Yates}} = \frac{n\left(|ad - bc| - \dfrac{n}{2}\right)^2}{(a+b)(c+d)(a+c)(b+d)}\]

\begin{itemize}
\tightlist
\item
  Se aplica cuando alguna \(E_{ij} \in [5, 10)\)
\item
  Siempre da \(\chi^2_{\text{Yates}} \leq \chi^2_{\text{Pearson}}\) →
  p-valor más conservador
\item
  En R: \texttt{chisq.test(tabla,\ correct\ =\ TRUE)} (opción por
  defecto para tablas 2×2)
\end{itemize}

\end{tcolorbox}

\begin{tcolorbox}[enhanced jigsaw, breakable, titlerule=0mm, rightrule=.15mm, leftrule=.75mm, left=2mm, opacitybacktitle=0.6, colback=white, bottomtitle=1mm, colbacktitle=quarto-callout-tip-color!10!white, title=\textcolor{quarto-callout-tip-color}{\faLightbulb}\hspace{0.5em}{Ejemplo: Pearson vs.~Yates en tabla con frecuencias moderadas}, toptitle=1mm, coltitle=black, opacityback=0, arc=.35mm, bottomrule=.15mm, toprule=.15mm, colframe=quarto-callout-tip-color-frame]

\begin{verbatim}
            Caso Control
Expuesto       7       4
No expuesto    3      11
\end{verbatim}

\begin{verbatim}

    Pearson's Chi-squared test

data:  m
X-squared = 5, df = 1, p-value = 0.03
\end{verbatim}

\begin{verbatim}

    Pearson's Chi-squared test with Yates' continuity correction

data:  m
X-squared = 3, df = 1, p-value = 0.08
\end{verbatim}

\begin{quote}
\textbf{Interpretación}

La prueba de Pearson (sin corrección) rechaza H₀ (χ² = 5, p = 0.03),
mientras que la corrección de Yates no rechaza (χ² = 3, p = 0.08). La
discordancia ocurre porque las frecuencias esperadas están en el rango
{[}5,10), donde la aproximación normal es moderada. La corrección de
Yates reduce el estadístico penalizando por el uso de una distribución
continua con datos discretos, produciendo un resultado más conservador.
En práctica clínica, con muestras pequeñas donde Eᵢⱼ ∈ {[}5,10), la
versión conservadora (Yates) es preferible para evitar falsos positivos,
aunque para grandes muestras ambas convergen.
\end{quote}

\end{tcolorbox}

La siguiente figura compara las distribuciones empíricas de los
estadísticos con y sin corrección frente a la distribución \(\chi^2(1)\)
teórica, usando simulaciones bajo \(H_0\) con \(n = 25\):

\begin{figure}[H]

{\centering \pandocbounded{\includegraphics[keepaspectratio]{chapters/semana11_files/figure-pdf/lst-s12-fig01-1.pdf}}

}

\caption{Comparación de la distribución empírica de χ² de Pearson (sin
corrección, naranja) y χ² de Yates (con corrección, azul) con la curva
χ²(1) teórica (negro), bajo H₀ con n = 25. La corrección de Yates
produce una aproximación más conservadora pero mejor calibrada para
muestras pequeñas.}

\end{figure}%

\subsection{Prueba Exacta de Fisher}\label{prueba-exacta-de-fisher}

Cuando alguna \(E_{ij} < 5\), la prueba de Fisher calcula la
probabilidad exacta de observar la tabla obtenida (o más extrema) bajo
\(H_0\), utilizando la distribución hipergeométrica:

\[P = \frac{(a+b)!\,(c+d)!\,(a+c)!\,(b+d)!}{n!\,a!\,b!\,c!\,d!}\]

\begin{tcolorbox}[enhanced jigsaw, breakable, titlerule=0mm, rightrule=.15mm, leftrule=.75mm, left=2mm, opacitybacktitle=0.6, colback=white, bottomtitle=1mm, colbacktitle=quarto-callout-tip-color!10!white, title=\textcolor{quarto-callout-tip-color}{\faLightbulb}\hspace{0.5em}{Ejemplo: Prueba exacta de Fisher en muestra pequeña}, toptitle=1mm, coltitle=black, opacityback=0, arc=.35mm, bottomrule=.15mm, toprule=.15mm, colframe=quarto-callout-tip-color-frame]

\begin{verbatim}
            Caso Control
Expuesto       4       2
No expuesto    1       8
\end{verbatim}

\begin{verbatim}

    Fisher's Exact Test for Count Data

data:  m
p-value = 0.09
alternative hypothesis: true odds ratio is not equal to 1
95 percent confidence interval:
   0.747 875.880
sample estimates:
odds ratio 
      12.5 
\end{verbatim}

\begin{quote}
\textbf{Interpretación}

Con n = 15 y varias frecuencias esperadas \textless{} 5, la aproximación
χ² es inapropiada, haciendo que la prueba exacta de Fisher sea el método
correcto. Fisher calcula mediante la distribución hipergeométrica la
probabilidad exacta de observar la tabla observada (u otra más extrema)
bajo H₀ de independencia. El resultado proporciona un p-valor exacto y
una odds ratio (OR) con intervalo de confianza exacto, basados en el
espacio muestral discreto de tablas 2×2. Este enfoque es especialmente
valioso en estudios de casos y controles de pequeño tamaño o cuando la
exposición es rara, situaciones comunes en medicina clínica donde los
datos categóricos a menudo tienen frecuencias limitadas.
\end{quote}

\end{tcolorbox}

\begin{center}\rule{0.5\linewidth}{0.5pt}\end{center}

\section{Prueba de McNemar para Datos
Apareados}\label{sec-semana11-mcnemar}

\subsection{¿Cuándo usar McNemar?}\label{cuuxe1ndo-usar-mcnemar}

La prueba \(\chi^2\) de independencia asume que las observaciones son
\textbf{independientes}. Sin embargo, en muchos diseños médicos los
datos son \textbf{apareados}: se mide el mismo paciente en dos
condiciones o momentos distintos.

\textbf{Situaciones típicas en medicina:}

\begin{itemize}
\tightlist
\item
  \textbf{Antes/después}: ¿Mejora la proporción de pacientes con PA
  controlada tras 6 meses de intervención?
\item
  \textbf{Dos pruebas diagnósticas}: ¿Concuerdan el test rápido de PCR y
  el cultivo bacteriano aplicados a los mismos pacientes?
\item
  \textbf{Pares emparejados}: ¿Difiere la detección de lesiones en ojo
  derecho vs.~ojo izquierdo del mismo paciente?
\item
  \textbf{Diseños cruzados}: ¿Es diferente la proporción de respuestas
  con tratamiento A vs.~B en el mismo individuo?
\end{itemize}

En estos diseños, la \textbf{unidad de análisis es el par}, no el
individuo, y solo los pares \textbf{discordantes} aportan información
sobre la diferencia.

\subsection{Estructura de la Tabla
Apareada}\label{estructura-de-la-tabla-apareada}

La tabla de datos apareados tiene una estructura especial:

\begin{longtable}[]{@{}
  >{\raggedright\arraybackslash}p{(\linewidth - 6\tabcolsep) * \real{0.0308}}
  >{\centering\arraybackslash}p{(\linewidth - 6\tabcolsep) * \real{0.4000}}
  >{\centering\arraybackslash}p{(\linewidth - 6\tabcolsep) * \real{0.4000}}
  >{\centering\arraybackslash}p{(\linewidth - 6\tabcolsep) * \real{0.1692}}@{}}
\toprule\noalign{}
\begin{minipage}[b]{\linewidth}\raggedright
\end{minipage} & \begin{minipage}[b]{\linewidth}\centering
\textbf{Condición B: Positivo}
\end{minipage} & \begin{minipage}[b]{\linewidth}\centering
\textbf{Condición B: Negativo}
\end{minipage} & \begin{minipage}[b]{\linewidth}\centering
\textbf{Total}
\end{minipage} \\
\midrule\noalign{}
\endhead
\bottomrule\noalign{}
\endlastfoot
\textbf{Condición A: Positivo} & \(a\) (concordantes +/+) & \(b\)
(discordantes +/−) & \(a+b\) \\
\textbf{Condición A: Negativo} & \(c\) (discordantes −/+) & \(d\)
(concordantes −/−) & \(c+d\) \\
\textbf{Total} & \(a+c\) & \(b+d\) & \(n\) \\
\end{longtable}

Donde: - \(a\): pares concordantes positivos (ambos positivos) - \(d\):
pares concordantes negativos (ambos negativos) - \(b\): discordantes ---
positivo en A, negativo en B - \(c\): discordantes --- negativo en A,
positivo en B

\textbf{La hipótesis nula} es que la proporción marginal de positivos es
la misma en ambas condiciones, equivalente a:

\[H_0: P(\text{A}^+\text{B}^-) = P(\text{A}^-\text{B}^+) \quad \Longleftrightarrow \quad \pi_b = \pi_c = 0.5\]

\subsection{Estadístico de McNemar}\label{estaduxedstico-de-mcnemar}

\begin{tcolorbox}[enhanced jigsaw, breakable, titlerule=0mm, rightrule=.15mm, leftrule=.75mm, left=2mm, opacitybacktitle=0.6, colback=white, bottomtitle=1mm, colbacktitle=quarto-callout-warning-color!10!white, title=\textcolor{quarto-callout-warning-color}{\faExclamationTriangle}\hspace{0.5em}{Estadístico: Prueba de McNemar}, toptitle=1mm, coltitle=black, opacityback=0, arc=.35mm, bottomrule=.15mm, toprule=.15mm, colframe=quarto-callout-warning-color-frame]

Solo los pares \textbf{discordantes} (\(b\) y \(c\)) aportan información
sobre la diferencia entre condiciones.

\textbf{Sin corrección de continuidad} (preferida cuando
\(b + c \geq 25\)):

\[\chi^2_{\text{McNemar}} = \frac{(b - c)^2}{b + c} \sim \chi^2_1 \quad \text{bajo } H_0\]

\textbf{Con corrección de continuidad de Yates} (recomendada cuando
\(b + c < 25\)):

\[\chi^2_{\text{McNemar,Yates}} = \frac{(|b - c| - 1)^2}{b + c}\]

\textbf{Para muestras muy pequeñas} (\(b + c < 10\)): usar el test
binomial exacto sobre los pares discordantes bajo \(H_0: p = 0.5\).

\textbf{Regla de decisión:} Rechazar \(H_0\) si
\(\chi^2 > \chi^2_{0.05, 1} = 3.84\) (o p-valor \textless{} 0.05).

\end{tcolorbox}

\subsection{Ejemplo: Intervención para el Control de la
Hipertensión}\label{ejemplo-intervenciuxf3n-para-el-control-de-la-hipertensiuxf3n}

\begin{tcolorbox}[enhanced jigsaw, breakable, titlerule=0mm, rightrule=.15mm, leftrule=.75mm, left=2mm, opacitybacktitle=0.6, colback=white, bottomtitle=1mm, colbacktitle=quarto-callout-tip-color!10!white, title=\textcolor{quarto-callout-tip-color}{\faLightbulb}\hspace{0.5em}{Ejemplo: Estudio antes-después --- control de la presión arterial}, toptitle=1mm, coltitle=black, opacityback=0, arc=.35mm, bottomrule=.15mm, toprule=.15mm, colframe=quarto-callout-tip-color-frame]

En una consulta de medicina interna del Hospital Universitario de
Granada, 100 pacientes hipertensos participan en un programa de
intervención multidisciplinar (dieta, ejercicio y educación sanitaria)
durante 6 meses. Se evalúa si el paciente tiene la PA controlada (\(<\)
140/90 mmHg) antes y después del programa.

\textbf{Tabla de pares apareados:}

\begin{longtable}[]{@{}lccc@{}}
\toprule\noalign{}
& \textbf{Post: PA Controlada} & \textbf{Post: PA No controlada} &
\textbf{Total} \\
\midrule\noalign{}
\endhead
\bottomrule\noalign{}
\endlastfoot
\textbf{Pre: PA Controlada} & 42 (\(a\)) & 8 (\(b\)) & 50 \\
\textbf{Pre: PA No controlada} & 28 (\(c\)) & 22 (\(d\)) & 50 \\
\textbf{Total} & 70 & 30 & 100 \\
\end{longtable}

\textbf{Interpretación de la tabla:} - 42 pacientes ya tenían la PA
controlada y mantuvieron el control tras la intervención - 22 pacientes
no la controlaron ni antes ni después - \textbf{\(b = 8\)}: 8 pacientes
que tenían la PA controlada la ``perdieron'' tras el programa -
\textbf{\(c = 28\)}: 28 pacientes que no la tenían controlada \textbf{la
consiguieron} con la intervención

\textbf{Pares discordantes:} \(b + c = 8 + 28 = 36 \geq 25\) → no es
necesaria la corrección de Yates.

\textbf{Cálculo manual:}
\[\chi^2_{\text{McNemar}} = \frac{(b - c)^2}{b + c} = \frac{(8 - 28)^2}{36} = \frac{400}{36} = 11.11\]

Como \(11.11 > 3.84\) (\(p < 0.001\)), rechazamos \(H_0\). La
intervención mejora significativamente el control de la PA: 28 pacientes
mejoraron frente a solo 8 que empeoraron.

\textbf{Nota sobre las proporciones marginales:} - Antes: 50/100 = 50\%
con PA controlada - Después: 70/100 = 70\% con PA controlada -
Diferencia: +20 puntos porcentuales (\(p < 0.001\) por McNemar)

Si se hubiera aplicado incorrectamente un \(\chi^2\) de independencia
(ignorando el apareamiento), la prueba sería errónea porque las
observaciones no son independientes.

\end{tcolorbox}

\subsection{Implementación: BioStatR vs.~Base
R}\label{implementaciuxf3n-biostatr-vs.-base-r}

\begin{tcolorbox}[enhanced jigsaw, breakable, titlerule=0mm, rightrule=.15mm, leftrule=.75mm, left=2mm, opacitybacktitle=0.6, colback=white, bottomtitle=1mm, colbacktitle=quarto-callout-tip-color!10!white, title=\textcolor{quarto-callout-tip-color}{\faLightbulb}\hspace{0.5em}{BioStatR --- testmcnemar(): Datos apareados}, toptitle=1mm, coltitle=black, opacityback=0, arc=.35mm, bottomrule=.15mm, toprule=.15mm, colframe=quarto-callout-tip-color-frame]

Para datos apareados (donde las observaciones son dependientes),
utilizamos la función \texttt{testmcnemar()} del paquete
\texttt{BioStatR}.

\begin{Shaded}
\begin{Highlighting}[]
\FunctionTok{library}\NormalTok{(BioStatR)}

\CommentTok{\# Tabla 2×2 apareada: testmcnemar()}
\FunctionTok{testmcnemar}\NormalTok{(}\AttributeTok{n11 =} \DecValTok{42}\NormalTok{, }\AttributeTok{n12 =} \DecValTok{8}\NormalTok{, }\AttributeTok{n21 =} \DecValTok{28}\NormalTok{, }\AttributeTok{n22 =} \DecValTok{22}\NormalTok{,}
           \AttributeTok{fcat =} \FunctionTok{c}\NormalTok{(}\StringTok{"Pre: Controlada"}\NormalTok{, }\StringTok{"Pre: No controlada"}\NormalTok{),}
           \AttributeTok{ccat =} \FunctionTok{c}\NormalTok{(}\StringTok{"Post: Controlada"}\NormalTok{, }\StringTok{"Post: No controlada"}\NormalTok{))}
\end{Highlighting}
\end{Shaded}

\begin{verbatim}


# Inferencia con dos proporciones (muestras apareadas)
# ----------------------------------------------------

# Frecuencias observadas pretest x posttest
                         Post: Controlada    Post: No controlada Total
   Pre: Controlada                     42                      8    50
   Pre: No controlada                  28                     22    50
   Total                               70                     30   100


# Proporciones observadas pretest x posttest
                         Post: Controlada    Post: No controlada Total
   Pre: Controlada                   0.42                   0.08  0.50
   Pre: No controlada                0.28                   0.22  0.50
   Total                             0.70                   0.30  1.00

# Test de McNemar: H₀:π₁₂=π₂₁
   Validez: n₁₂+n₂₁ = 36 > 10 el test es válido 
   Zexp =  3.2500 
               valor.p Alternativa
    Bilateral   0.0012  H₁:π₁₂≠π₂₁
    Unilateral  0.0006  H₁:π₁₂<π₂₁


# Test exacto de Fisher:
  H₀:π₁₂=0.5 para n₁₂ ~ B(n₁₂+n₂₁, π₁₂) 
               Valor.p Alternativa
    Bilateral   0.0141  H₁:π₁₂≠0.5
    Unilateral  0.0006  H₁:π₁₂<0.5
  ____
  * Aquí se alude a la probabilidad total de la discordancia, es decir que π₁₂+π₂₁=1

# Estimación de las proporciones individuales de discordancias π₁₂ y π₂₁ (método de Wald ajustado)
   [1] p₁₂ = 0.0962, 95%-IC(π₁₂) = (0.0395, 0.1528) 
   [2] p₂₁ = 0.2885, 95%-IC(π₂₁) = (0.2014, 0.3755) 


# Intervalo de confianza para la diferencia de 2 proporciones apareadas
   [1] Método de Wald (clásico con cpc):
      Estimación puntual de π₁₂-π₂₁ = -0.2000 
      Validez: n₁₂+n₂₁ = 36 > 5, el IC es válido 
      95%-IC(π₁₂-π₂₁) = (-0.3117, -0.0883) 

   [2] Método de Agresti-Min:
      Estimación puntual de π₁₂-π₂₁ = -0.1961 
      Validez: siempre es válido 
      95%-IC(π₁₂-π₂₁) = (-0.3066, -0.0856) 
\end{verbatim}

\end{tcolorbox}

La prueba de McNemar se centra exclusivamente en los pares discordantes:
aquellos pacientes cuyo estado clínico cambió entre las dos mediciones.
En nuestro ejemplo, el resultado de \texttt{testmcnemar()} nos ofrece
tres claves para la interpretación:

\begin{enumerate}
\def\labelenumi{\arabic{enumi}.}
\item
  \textbf{Identificación de Cambios (Pares Discordantes):}

  \begin{itemize}
  \tightlist
  \item
    \textbf{\(b = 8\)}: Pacientes que estaban controlados ``antes'' y
    pasaron a estar ``no controlados'' ``después'' (empeoraron).
  \item
    \textbf{\(c = 28\)}: Pacientes que no estaban controlados ``antes''
    y lograron el control ``después'' (mejoraron).
  \end{itemize}
\item
  \textbf{Lógica Estadística:} La prueba ignora los pares concordantes
  (los que siempre estuvieron controlados o nunca lo estuvieron) y
  evalúa si el número de pacientes que mejoró es significativamente
  diferente del número que empeoró. Bajo la hipótesis nula (\(H_0\)),
  esperaríamos que el número de cambios en ambas direcciones fuera igual
  (\(b = c\)).
\item
  \textbf{Conclusión Clínica:} Dado que nuestro p-valor es
  \texttt{\textless{}\ 0.05} y observamos que 28 pacientes mejoraron
  frente a solo 8 que empeoraron (\(c > b\)), rechazamos \(H_0\).
  Concluimos que la intervención tiene un efecto beneficioso
  estadísticamente significativo sobre el control de la presión
  arterial.
\end{enumerate}

En resumen, la prueba de McNemar nos permite afirmar que el cambio
observado en las proporciones marginales de control de PA (del 50\% al
70\%) no es debido al azar, sino que refleja un impacto real del
programa de intervención sobre los pacientes.

\section{Diseños de Estudio y Medidas de
Asociación}\label{sec-semana11-asociacion}

La elección de la medida de asociación depende del \textbf{diseño del
estudio}. Todas se calculan a partir de la tabla 2×2 pero tienen
interpretaciones distintas.

\subsection{Estructura Epidemiológica de la Tabla
2×2}\label{estructura-epidemioluxf3gica-de-la-tabla-22}

\begin{longtable}[]{@{}lccc@{}}
\toprule\noalign{}
& \textbf{Enfermedad (D+)} & \textbf{Sin enfermedad (D−)} &
\textbf{Total} \\
\midrule\noalign{}
\endhead
\bottomrule\noalign{}
\endlastfoot
\textbf{Exposición (E+)} & \(a\) & \(b\) & \(n_1 = a+b\) \\
\textbf{Sin exposición (E−)} & \(c\) & \(d\) & \(n_2 = c+d\) \\
\textbf{Total} & \(M_1 = a+c\) & \(M_0 = b+d\) & \(n\) \\
\end{longtable}

Riesgos (proporciones):
\[p_1 = \frac{a}{a+b} \quad (\text{riesgo en expuestos}), \qquad p_2 = \frac{c}{c+d} \quad (\text{riesgo en no expuestos})\]

\subsection{Estudio de Cohortes: Riesgo Relativo
(RR)}\label{estudio-de-cohortes-riesgo-relativo-rr}

En un estudio de cohortes, el seguimiento prospectivo permite estimar
riesgos directamente.

\begin{tcolorbox}[enhanced jigsaw, breakable, titlerule=0mm, rightrule=.15mm, leftrule=.75mm, left=2mm, opacitybacktitle=0.6, colback=white, bottomtitle=1mm, colbacktitle=quarto-callout-note-color!10!white, title=\textcolor{quarto-callout-note-color}{\faInfo}\hspace{0.5em}{Definición: Riesgo Relativo (RR)}, toptitle=1mm, coltitle=black, opacityback=0, arc=.35mm, bottomrule=.15mm, toprule=.15mm, colframe=quarto-callout-note-color-frame]

\[\widehat{RR} = \frac{p_1}{p_2} = \frac{a/(a+b)}{c/(c+d)}\]

\textbf{Intervalo de confianza al 95\% (método log):}
\[IC_{95\%}(RR) = \exp\!\left[\ln\widehat{RR} \pm 1.96 \sqrt{\frac{b}{a(a+b)} + \frac{d}{c(c+d)}}\right]\]

\textbf{Interpretación:} - \(RR = 1\): sin asociación (igual riesgo en
expuestos y no expuestos) - \(RR > 1\): factor de riesgo (la exposición
aumenta el riesgo) - \(RR < 1\): factor protector (la exposición
disminuye el riesgo)

\end{tcolorbox}

\begin{tcolorbox}[enhanced jigsaw, breakable, titlerule=0mm, rightrule=.15mm, leftrule=.75mm, left=2mm, opacitybacktitle=0.6, colback=white, bottomtitle=1mm, colbacktitle=quarto-callout-tip-color!10!white, title=\textcolor{quarto-callout-tip-color}{\faLightbulb}\hspace{0.5em}{Ejemplo: Cohorte --- tabaquismo y cardiopatía isquémica}, toptitle=1mm, coltitle=black, opacityback=0, arc=.35mm, bottomrule=.15mm, toprule=.15mm, colframe=quarto-callout-tip-color-frame]

Seguimiento de 500 fumadores y 500 no fumadores durante 10 años:

\begin{longtable}[]{@{}lccc@{}}
\toprule\noalign{}
& Cardiopatía & Sin cardiopatía & Total \\
\midrule\noalign{}
\endhead
\bottomrule\noalign{}
\endlastfoot
Fumadores & 80 & 420 & 500 \\
No fumadores & 30 & 470 & 500 \\
\end{longtable}

\begin{Shaded}
\begin{Highlighting}[]
\CommentTok{\# Cálculo de RR e IC}
\NormalTok{a }\OtherTok{\textless{}{-}} \DecValTok{80}\NormalTok{; b }\OtherTok{\textless{}{-}} \DecValTok{420}\NormalTok{; c\_ }\OtherTok{\textless{}{-}} \DecValTok{30}\NormalTok{; d }\OtherTok{\textless{}{-}} \DecValTok{470}
\NormalTok{p1 }\OtherTok{\textless{}{-}}\NormalTok{ a }\SpecialCharTok{/}\NormalTok{ (a }\SpecialCharTok{+}\NormalTok{ b); p2 }\OtherTok{\textless{}{-}}\NormalTok{ c\_ }\SpecialCharTok{/}\NormalTok{ (c\_ }\SpecialCharTok{+}\NormalTok{ d)}
\NormalTok{RR }\OtherTok{\textless{}{-}}\NormalTok{ p1 }\SpecialCharTok{/}\NormalTok{ p2}
\NormalTok{se\_logRR }\OtherTok{\textless{}{-}} \FunctionTok{sqrt}\NormalTok{(b }\SpecialCharTok{/}\NormalTok{ (a}\SpecialCharTok{*}\NormalTok{(a}\SpecialCharTok{+}\NormalTok{b)) }\SpecialCharTok{+}\NormalTok{ d }\SpecialCharTok{/}\NormalTok{ (c\_}\SpecialCharTok{*}\NormalTok{(c\_}\SpecialCharTok{+}\NormalTok{d)))}
\NormalTok{IC\_RR }\OtherTok{\textless{}{-}} \FunctionTok{exp}\NormalTok{(}\FunctionTok{log}\NormalTok{(RR) }\SpecialCharTok{+} \FunctionTok{c}\NormalTok{(}\SpecialCharTok{{-}}\DecValTok{1}\NormalTok{, }\DecValTok{1}\NormalTok{) }\SpecialCharTok{*} \FloatTok{1.96} \SpecialCharTok{*}\NormalTok{ se\_logRR)}
\FunctionTok{cat}\NormalTok{(}\StringTok{"p1 (riesgo fumadores)     ="}\NormalTok{, }\FunctionTok{round}\NormalTok{(p1, }\DecValTok{4}\NormalTok{), }\StringTok{"}\SpecialCharTok{\textbackslash{}n}\StringTok{"}\NormalTok{)}
\end{Highlighting}
\end{Shaded}

\begin{verbatim}
p1 (riesgo fumadores)     = 0.16 
\end{verbatim}

\begin{Shaded}
\begin{Highlighting}[]
\FunctionTok{cat}\NormalTok{(}\StringTok{"p2 (riesgo no fumadores)  ="}\NormalTok{, }\FunctionTok{round}\NormalTok{(p2, }\DecValTok{4}\NormalTok{), }\StringTok{"}\SpecialCharTok{\textbackslash{}n}\StringTok{"}\NormalTok{)}
\end{Highlighting}
\end{Shaded}

\begin{verbatim}
p2 (riesgo no fumadores)  = 0.06 
\end{verbatim}

\begin{Shaded}
\begin{Highlighting}[]
\FunctionTok{cat}\NormalTok{(}\StringTok{"RR ="}\NormalTok{, }\FunctionTok{round}\NormalTok{(RR, }\DecValTok{3}\NormalTok{), }\StringTok{"}\SpecialCharTok{\textbackslash{}n}\StringTok{"}\NormalTok{)}
\end{Highlighting}
\end{Shaded}

\begin{verbatim}
RR = 2.67 
\end{verbatim}

\begin{Shaded}
\begin{Highlighting}[]
\FunctionTok{cat}\NormalTok{(}\StringTok{"IC 95\%: ["}\NormalTok{, }\FunctionTok{round}\NormalTok{(IC\_RR[}\DecValTok{1}\NormalTok{], }\DecValTok{3}\NormalTok{), }\StringTok{","}\NormalTok{, }\FunctionTok{round}\NormalTok{(IC\_RR[}\DecValTok{2}\NormalTok{], }\DecValTok{3}\NormalTok{), }\StringTok{"]}\SpecialCharTok{\textbackslash{}n}\StringTok{"}\NormalTok{)}
\end{Highlighting}
\end{Shaded}

\begin{verbatim}
IC 95%: [ 1.79 , 3.98 ]
\end{verbatim}

\textbf{Interpretación:} Los fumadores tienen \textbf{2.67 veces más
riesgo} de cardiopatía isquémica que los no fumadores (\(RR = 2.67\);
\(IC_{95\%}\): 1.79--3.97). El intervalo no contiene el 1, confirmando
asociación estadísticamente significativa.

\end{tcolorbox}

\subsection{Diferencia de Riesgos (RD) y
NNT}\label{diferencia-de-riesgos-rd-y-nnt}

\begin{tcolorbox}[enhanced jigsaw, breakable, titlerule=0mm, rightrule=.15mm, leftrule=.75mm, left=2mm, opacitybacktitle=0.6, colback=white, bottomtitle=1mm, colbacktitle=quarto-callout-note-color!10!white, title=\textcolor{quarto-callout-note-color}{\faInfo}\hspace{0.5em}{Definición: Diferencia de Riesgos (RD)}, toptitle=1mm, coltitle=black, opacityback=0, arc=.35mm, bottomrule=.15mm, toprule=.15mm, colframe=quarto-callout-note-color-frame]

\[\widehat{RD} = p_1 - p_2 = \frac{a}{a+b} - \frac{c}{c+d}\]

\textbf{Intervalo de confianza al 95\%:}
\[IC_{95\%}(RD) = \widehat{RD} \pm 1.96\sqrt{\frac{p_1(1-p_1)}{n_1} + \frac{p_2(1-p_2)}{n_2}}\]

\textbf{Número Necesario a Tratar/Perjudicar (NNT/NNH):}
\[NNT = \frac{1}{|\widehat{RD}|}\]

El NNT indica cuántas personas hay que tratar (o exponer) para observar
un caso adicional de resultado.

\end{tcolorbox}

Para el ejemplo anterior: \(RD = 0.16 - 0.06 = 0.10\) →
\(NNH = 1/0.10 = 10\) (por cada 10 fumadores seguidos durante 10 años, 1
caso adicional de cardiopatía).

\subsection{Estudio Transversal: Razón de Prevalencias
(RP)}\label{estudio-transversal-razuxf3n-de-prevalencias-rp}

En un estudio transversal, las proporciones son \textbf{prevalencias}
(no riesgos de incidencia). La medida de elección es la \textbf{Razón de
Prevalencias (RP)}, cuya fórmula es idéntica al RR pero con
prevalencias:

\[\widehat{RP} = \frac{\text{prevalencia en expuestos}}{\text{prevalencia en no expuestos}} = \frac{a/(a+b)}{c/(c+d)}\]

El intervalo de confianza se calcula con la misma fórmula que el RR.
\textbf{Nota importante:} el OR no es una buena aproximación de la RP
cuando la prevalencia es \textgreater{} 10\%.

\subsection{Estudio Caso-Control: Odds Ratio
(OR)}\label{estudio-caso-control-odds-ratio-or}

En un estudio caso-control el muestreo es por resultado (se seleccionan
casos y controles), por lo que \textbf{no se pueden estimar riesgos ni
RR directamente}. La medida correcta es el \textbf{Odds Ratio}.

\begin{tcolorbox}[enhanced jigsaw, breakable, titlerule=0mm, rightrule=.15mm, leftrule=.75mm, left=2mm, opacitybacktitle=0.6, colback=white, bottomtitle=1mm, colbacktitle=quarto-callout-note-color!10!white, title=\textcolor{quarto-callout-note-color}{\faInfo}\hspace{0.5em}{Definición: Odds Ratio (OR)}, toptitle=1mm, coltitle=black, opacityback=0, arc=.35mm, bottomrule=.15mm, toprule=.15mm, colframe=quarto-callout-note-color-frame]

\[\widehat{OR} = \frac{a/b}{c/d} = \frac{ad}{bc}\]

\textbf{Intervalo de confianza al 95\% (método de Woolf/Wald):}
\[IC_{95\%}(OR) = \exp\!\left[\ln\widehat{OR} \pm 1.96\sqrt{\frac{1}{a} + \frac{1}{b} + \frac{1}{c} + \frac{1}{d}}\right]\]

\textbf{Propiedades:} - \(OR = 1\): sin asociación - \(OR > 1\):
exposición asociada con mayor odds de resultado - \(OR \approx RR\)
cuando la enfermedad es \textbf{rara} (prevalencia \textless{} 10\%) -
\texttt{tabla2x2()} de BioStatR calcula OR e IC directamente

\end{tcolorbox}

\begin{tcolorbox}[enhanced jigsaw, breakable, titlerule=0mm, rightrule=.15mm, leftrule=.75mm, left=2mm, opacitybacktitle=0.6, colback=white, bottomtitle=1mm, colbacktitle=quarto-callout-tip-color!10!white, title=\textcolor{quarto-callout-tip-color}{\faLightbulb}\hspace{0.5em}{Ejemplo: Caso-Control --- alcohol y cirrosis hepática}, toptitle=1mm, coltitle=black, opacityback=0, arc=.35mm, bottomrule=.15mm, toprule=.15mm, colframe=quarto-callout-tip-color-frame]

\begin{Shaded}
\begin{Highlighting}[]
\FunctionTok{library}\NormalTok{(BioStatR)}

\CommentTok{\# 200 casos de cirrosis y 200 controles}
\CommentTok{\# Consumo excesivo de alcohol: 120 casos, 60 controles}
\FunctionTok{tabla2x2}\NormalTok{(}\AttributeTok{o11 =} \DecValTok{120}\NormalTok{, }\AttributeTok{o12 =} \DecValTok{80}\NormalTok{, }\AttributeTok{o21 =} \DecValTok{60}\NormalTok{, }\AttributeTok{o22 =} \DecValTok{140}\NormalTok{,}
         \AttributeTok{fcat =} \FunctionTok{c}\NormalTok{(}\StringTok{"Alcohol excesivo"}\NormalTok{, }\StringTok{"Alcohol no excesivo"}\NormalTok{),}
         \AttributeTok{ccat =} \FunctionTok{c}\NormalTok{(}\StringTok{"Caso (cirrosis)"}\NormalTok{, }\StringTok{"Control"}\NormalTok{),}
         \AttributeTok{estudio =} \StringTok{"CC"}\NormalTok{)}
\end{Highlighting}
\end{Shaded}

\begin{verbatim}

# Análisis de tablas 2x2
# ----------------------

# Frecuencias observadas
                        Caso (cirrosis)   Control Total
  Alcohol excesivo                  120        80   200
  Alcohol no excesivo                60       140   200
  Total                             180       220   400


# Test Chi-cuadrado para un estudio indefinido

  χ² = 36.352,   gl = 1,  p < 0.001, (cpc = 2) 
  Validez: Frecuencia mínima esperada = 90.00  > 7.7  

  Test exacto de Fisher (bilateral): p < 0.001 

  --- Otros criterios χ²:  
  χ² = 36.364,   gl = 1,  p < 0.001, (sin cpc) 
  χ² = 35.162,   gl = 1,  p < 0.001, (cpc de Yates = 200.00) 

# Medidas de asociación para un estudio indefinido

  Razón del producto cruzado (odds ratio):
  OR=3.500; 95%-IC(OR)= (2.300, 5.253) 
\end{verbatim}

\textbf{Interpretación:} \(OR = 3.50\) (\(IC_{95\%}\): 2.28--5.37). Las
personas con consumo excesivo de alcohol tienen \textbf{3.5 veces más
odds} de cirrosis que las de consumo no excesivo. El resultado es
altamente significativo (\(p < 0.001\)).

\end{tcolorbox}

\subsection{Resumen de Intervalos de
Confianza}\label{resumen-de-intervalos-de-confianza}

\begin{longtable}[]{@{}
  >{\raggedright\arraybackslash}p{(\linewidth - 2\tabcolsep) * \real{0.2963}}
  >{\raggedright\arraybackslash}p{(\linewidth - 2\tabcolsep) * \real{0.7037}}@{}}
\toprule\noalign{}
\begin{minipage}[b]{\linewidth}\raggedright
Medida
\end{minipage} & \begin{minipage}[b]{\linewidth}\raggedright
Fórmula del IC 95\%
\end{minipage} \\
\midrule\noalign{}
\endhead
\bottomrule\noalign{}
\endlastfoot
\(RR\) &
\(\exp\!\left[\ln\widehat{RR} \pm 1.96\sqrt{\frac{b}{a(a+b)}+\frac{d}{c(c+d)}}\right]\) \\
\(OR\) &
\(\exp\!\left[\ln\widehat{OR} \pm 1.96\sqrt{\frac{1}{a}+\frac{1}{b}+\frac{1}{c}+\frac{1}{d}}\right]\) \\
\(RD\) &
\(\widehat{RD} \pm 1.96\sqrt{\frac{p_1(1-p_1)}{n_1}+\frac{p_2(1-p_2)}{n_2}}\) \\
\end{longtable}

\begin{center}\rule{0.5\linewidth}{0.5pt}\end{center}

\section{Resumen Comparativo: ¿Qué Medida
Usar?}\label{sec-semana11-resumen-medidas}

\begin{longtable}[]{@{}
  >{\raggedright\arraybackslash}p{(\linewidth - 6\tabcolsep) * \real{0.2923}}
  >{\raggedright\arraybackslash}p{(\linewidth - 6\tabcolsep) * \real{0.2615}}
  >{\raggedright\arraybackslash}p{(\linewidth - 6\tabcolsep) * \real{0.2769}}
  >{\raggedright\arraybackslash}p{(\linewidth - 6\tabcolsep) * \real{0.1692}}@{}}
\toprule\noalign{}
\begin{minipage}[b]{\linewidth}\raggedright
Diseño del estudio
\end{minipage} & \begin{minipage}[b]{\linewidth}\raggedright
Medida principal
\end{minipage} & \begin{minipage}[b]{\linewidth}\raggedright
Medida secundaria
\end{minipage} & \begin{minipage}[b]{\linewidth}\raggedright
Función R
\end{minipage} \\
\midrule\noalign{}
\endhead
\bottomrule\noalign{}
\endlastfoot
\textbf{Cohorte} (prospectivo) & RR (Riesgo Relativo) & RD, NNT &
Cálculo manual \\
\textbf{Transversal} (prevalencia) & RP (Razón de Prevalencias) & RD &
Cálculo manual \\
\textbf{Caso-Control} (retrospectivo) & OR (Odds Ratio) & --- &
\texttt{tabla2x2()} \\
\textbf{Ensayo clínico (ECA)} & RR, RD, NNT & OR & Cálculo manual \\
\textbf{Apareado / Antes-Después} & McNemar \(\chi^2\) & Diferencia de
proporciones & \texttt{mcnemar.test()},
\texttt{tabla2x2(estudio="AP")} \\
\end{longtable}

\begin{tcolorbox}[enhanced jigsaw, breakable, titlerule=0mm, rightrule=.15mm, leftrule=.75mm, left=2mm, opacitybacktitle=0.6, colback=white, bottomtitle=1mm, colbacktitle=quarto-callout-warning-color!10!white, title=\textcolor{quarto-callout-warning-color}{\faExclamationTriangle}\hspace{0.5em}{Cuándo OR ≈ RR}, toptitle=1mm, coltitle=black, opacityback=0, arc=.35mm, bottomrule=.15mm, toprule=.15mm, colframe=quarto-callout-warning-color-frame]

El OR sobreestima al RR cuando la prevalencia/incidencia del resultado
es alta (\textgreater{} 10\%). En enfermedades raras, \(OR \approx RR\).
En estudios caso-control, siempre calcular OR (el RR no es estimable sin
incidencia).

\end{tcolorbox}

\begin{center}\rule{0.5\linewidth}{0.5pt}\end{center}

\section{Implementación en R}\label{sec-semana11-r}

\subsection{Normalidad con testnormal() y
shapiro.test()}\label{normalidad-con-testnormal-y-shapiro.test}

\begin{Shaded}
\begin{Highlighting}[]
\FunctionTok{library}\NormalTok{(BioStatR)}

\CommentTok{\# Con BioStatR (reporte completo con histograma)}
\FunctionTok{testnormal}\NormalTok{(osteo}\SpecialCharTok{$}\NormalTok{imc)}
\end{Highlighting}
\end{Shaded}

\pandocbounded{\includegraphics[keepaspectratio]{chapters/semana11_files/figure-pdf/lst-s12-bio05-1.pdf}}

\begin{verbatim}

# Test de normalidad de Shapiro-Wilk  
 ------------------------------------- 
   n = 94,  W = 0.906,  p < 0.001 
\end{verbatim}

\begin{Shaded}
\begin{Highlighting}[]
\CommentTok{\# Con R base}
\FunctionTok{shapiro.test}\NormalTok{(osteo}\SpecialCharTok{$}\NormalTok{imc)}
\end{Highlighting}
\end{Shaded}

\begin{verbatim}

    Shapiro-Wilk normality test

data:  osteo$imc
W = 0.9, p-value = 0.000005
\end{verbatim}

\subsection{Chi-cuadrado con tabla2x2() y
chisq.test()}\label{chi-cuadrado-con-tabla2x2-y-chisq.test}

\begin{Shaded}
\begin{Highlighting}[]
\CommentTok{\# Con BioStatR (tabla 2x2, Chi{-}cuadrado + OR + Fisher integrado)}
\FunctionTok{tabla2x2}\NormalTok{(}\AttributeTok{o11 =} \DecValTok{44}\NormalTok{, }\AttributeTok{o12 =} \DecValTok{9}\NormalTok{, }\AttributeTok{o21 =} \DecValTok{26}\NormalTok{, }\AttributeTok{o22 =} \DecValTok{15}\NormalTok{,}
         \AttributeTok{fcat =} \FunctionTok{c}\NormalTok{(}\StringTok{"No fumador"}\NormalTok{, }\StringTok{"Fumador"}\NormalTok{),}
         \AttributeTok{ccat =} \FunctionTok{c}\NormalTok{(}\StringTok{"Sin osteoporosis"}\NormalTok{, }\StringTok{"Con osteoporosis"}\NormalTok{))}
\end{Highlighting}
\end{Shaded}

\begin{verbatim}

# Análisis de tablas 2x2
# ----------------------

# Frecuencias observadas
               Sin osteoporosis   Con osteoporosis Total
  No fumador                 44                  9    53
  Fumador                    26                 15    41
  Total                      70                 24    94


# Test Chi-cuadrado para un estudio transversal

  χ² = 4.662,   gl = 1,  p = 0.031, (cpc = 0.5) 
  Validez: Frecuencia mínima esperada = 10.47  > 3.9  

  Test exacto de Fisher (bilateral): p = 0.035 

  --- Otros criterios χ²:  
  χ² = 4.673,   gl = 1,  p = 0.054, (sin cpc) 
  χ² = 3.699,   gl = 1,  p = 0.054, (cpc de Yates = 47.00) 

# Estimación de la prevalenciaπ en un estudio transversal
  Método de Wald ajustado: 
  p=0.561; 95%-IC(π)=(0.463, 0.659) 

# Medidas de asociación para un estudio transversal
  [!] Las medidas de riesgo se calculan como riesgo de la categoría  
      en la 1a columna (frente a la 2a) para la categoría en la 1a 
      fila (frente a la 2a)

  Riesgo absoluto (diferencia de Berkson; método de Agresti-Caffo): 
  d=0.254; 95%-IC(d)=(0.023, 0.459) 

  Riesgo relativo: 
  Rr=1.676; 95%-IC(Rr)=(0.969, 2.808) 

  Riesgo atribuible:
  Ra=0.335; 95%-IC(Ra)= (-0.049, 0.578) 

  Razón del producto cruzado (odds ratio):
  OR=2.821; 95%-IC(OR)= (1.070, 7.013) 
\end{verbatim}

\begin{Shaded}
\begin{Highlighting}[]
\CommentTok{\# Con R base (independiente y Fisher por separado)}
\NormalTok{m }\OtherTok{\textless{}{-}} \FunctionTok{matrix}\NormalTok{(}\FunctionTok{c}\NormalTok{(}\DecValTok{44}\NormalTok{, }\DecValTok{26}\NormalTok{, }\DecValTok{9}\NormalTok{, }\DecValTok{15}\NormalTok{), }\AttributeTok{nrow =} \DecValTok{2}\NormalTok{,}
            \AttributeTok{dimnames =} \FunctionTok{list}\NormalTok{(}\FunctionTok{c}\NormalTok{(}\StringTok{"No fumador"}\NormalTok{, }\StringTok{"Fumador"}\NormalTok{),}
                            \FunctionTok{c}\NormalTok{(}\StringTok{"Sin osteoporosis"}\NormalTok{, }\StringTok{"Con osteoporosis"}\NormalTok{)))}
\FunctionTok{cat}\NormalTok{(}\StringTok{"}\SpecialCharTok{\textbackslash{}n}\StringTok{{-}{-}{-} chisq.test() de R base {-}{-}{-}}\SpecialCharTok{\textbackslash{}n}\StringTok{"}\NormalTok{)}
\end{Highlighting}
\end{Shaded}

\begin{verbatim}

--- chisq.test() de R base ---
\end{verbatim}

\begin{Shaded}
\begin{Highlighting}[]
\FunctionTok{chisq.test}\NormalTok{(m)}
\end{Highlighting}
\end{Shaded}

\begin{verbatim}

    Pearson's Chi-squared test with Yates' continuity correction

data:  m
X-squared = 4, df = 1, p-value = 0.05
\end{verbatim}

\begin{Shaded}
\begin{Highlighting}[]
\FunctionTok{cat}\NormalTok{(}\StringTok{"}\SpecialCharTok{\textbackslash{}n}\StringTok{{-}{-}{-} fisher.test() de R base {-}{-}{-}}\SpecialCharTok{\textbackslash{}n}\StringTok{"}\NormalTok{)}
\end{Highlighting}
\end{Shaded}

\begin{verbatim}

--- fisher.test() de R base ---
\end{verbatim}

\begin{Shaded}
\begin{Highlighting}[]
\FunctionTok{fisher.test}\NormalTok{(m)}
\end{Highlighting}
\end{Shaded}

\begin{verbatim}

    Fisher's Exact Test for Count Data

data:  m
p-value = 0.04
alternative hypothesis: true odds ratio is not equal to 1
95 percent confidence interval:
 0.981 8.351
sample estimates:
odds ratio 
      2.79 
\end{verbatim}

\subsection{McNemar con tabla2x2() y
mcnemar.test()}\label{mcnemar-con-tabla2x2-y-mcnemar.test}

\begin{Shaded}
\begin{Highlighting}[]
\CommentTok{\# Base R: mcnemar.test() — control de PA antes y después del programa}
\NormalTok{m\_paired }\OtherTok{\textless{}{-}} \FunctionTok{matrix}\NormalTok{(}\FunctionTok{c}\NormalTok{(}\DecValTok{42}\NormalTok{, }\DecValTok{28}\NormalTok{, }\DecValTok{8}\NormalTok{, }\DecValTok{22}\NormalTok{), }\AttributeTok{nrow =} \DecValTok{2}\NormalTok{,}
                   \AttributeTok{dimnames =} \FunctionTok{list}\NormalTok{(}
                     \FunctionTok{c}\NormalTok{(}\StringTok{"Pre: Controlada"}\NormalTok{, }\StringTok{"Pre: No controlada"}\NormalTok{),}
                     \FunctionTok{c}\NormalTok{(}\StringTok{"Post: Controlada"}\NormalTok{, }\StringTok{"Post: No controlada"}\NormalTok{)))}

\FunctionTok{cat}\NormalTok{(}\StringTok{"=== mcnemar.test() — Base R ===}\SpecialCharTok{\textbackslash{}n}\StringTok{"}\NormalTok{)}
\end{Highlighting}
\end{Shaded}

\begin{verbatim}
=== mcnemar.test() — Base R ===
\end{verbatim}

\begin{Shaded}
\begin{Highlighting}[]
\FunctionTok{mcnemar.test}\NormalTok{(m\_paired, }\AttributeTok{correct =} \ConstantTok{FALSE}\NormalTok{)  }\CommentTok{\# Sin corrección (b+c=36 \textgreater{}= 25)}
\end{Highlighting}
\end{Shaded}

\begin{verbatim}

    McNemar's Chi-squared test

data:  m_paired
McNemar's chi-squared = 11, df = 1, p-value = 0.0009
\end{verbatim}

\begin{Shaded}
\begin{Highlighting}[]
\CommentTok{\# BioStatR: tabla2x2(..., estudio = "AP")}
\CommentTok{\# (ver sección 12.5.5 para la salida completa; requiere BioStatR)}
\CommentTok{\# tabla2x2(o11=42, o12=8, o21=28, o22=22,}
\CommentTok{\#          fcat = c("Pre: Controlada","Pre: No controlada"),}
\CommentTok{\#          ccat = c("Post: Controlada","Post: No controlada"),}
\CommentTok{\#          estudio = "AP")}
\end{Highlighting}
\end{Shaded}

\subsection{Tablas r×c con tablarxc()}\label{tablas-rc-con-tablarxc}

\begin{Shaded}
\begin{Highlighting}[]
\CommentTok{\# Tabla 3x2: tabaco y grupos de edad}
\FunctionTok{tablarxc}\NormalTok{(}\AttributeTok{fvar =}\NormalTok{ osteo}\SpecialCharTok{$}\NormalTok{grupo\_edad, }\AttributeTok{cvar =}\NormalTok{ osteo}\SpecialCharTok{$}\NormalTok{tabaco)}
\end{Highlighting}
\end{Shaded}

\begin{verbatim}

# Test Chi-cuadrado para tablas RxC
# ---------------------------------

# Frecuencias observadas
         A  B Total
A       21 11    32
B       14 16    30
C       18 14    32
  Total 53 41    94

# Test chi-cuadrado 
  Validez: Frecuencia mínima esperada = 13.09 
           0 frecuencias esperadas son menores a 1
           0 son menores a 5 (el 0% de la tabla) 
  χ²(2 gl) = 2.263, p = 0.323 
\end{verbatim}

\subsection{Cálculo Manual de RR, RD y
OR}\label{cuxe1lculo-manual-de-rr-rd-y-or}

\begin{Shaded}
\begin{Highlighting}[]
\CommentTok{\# Función para RR, RD, OR con IC (diseños no apareados)}
\NormalTok{calcular\_medidas }\OtherTok{\textless{}{-}} \ControlFlowTok{function}\NormalTok{(a, b, c, d, }\AttributeTok{conf =} \FloatTok{0.95}\NormalTok{) \{}
\NormalTok{  z  }\OtherTok{\textless{}{-}} \FunctionTok{qnorm}\NormalTok{(}\DecValTok{1} \SpecialCharTok{{-}}\NormalTok{ (}\DecValTok{1} \SpecialCharTok{{-}}\NormalTok{ conf)}\SpecialCharTok{/}\DecValTok{2}\NormalTok{)}
\NormalTok{  n1 }\OtherTok{\textless{}{-}}\NormalTok{ a }\SpecialCharTok{+}\NormalTok{ b; n2 }\OtherTok{\textless{}{-}}\NormalTok{ c }\SpecialCharTok{+}\NormalTok{ d}
\NormalTok{  p1 }\OtherTok{\textless{}{-}}\NormalTok{ a}\SpecialCharTok{/}\NormalTok{n1;  p2 }\OtherTok{\textless{}{-}}\NormalTok{ c}\SpecialCharTok{/}\NormalTok{n2}

\NormalTok{  RR    }\OtherTok{\textless{}{-}}\NormalTok{ p1}\SpecialCharTok{/}\NormalTok{p2}
\NormalTok{  IC\_RR }\OtherTok{\textless{}{-}} \FunctionTok{exp}\NormalTok{(}\FunctionTok{log}\NormalTok{(RR) }\SpecialCharTok{+} \FunctionTok{c}\NormalTok{(}\SpecialCharTok{{-}}\DecValTok{1}\NormalTok{,}\DecValTok{1}\NormalTok{) }\SpecialCharTok{*}\NormalTok{ z }\SpecialCharTok{*} \FunctionTok{sqrt}\NormalTok{(b}\SpecialCharTok{/}\NormalTok{(a}\SpecialCharTok{*}\NormalTok{n1) }\SpecialCharTok{+}\NormalTok{ d}\SpecialCharTok{/}\NormalTok{(c}\SpecialCharTok{*}\NormalTok{n2)))}

\NormalTok{  OR    }\OtherTok{\textless{}{-}}\NormalTok{ (a}\SpecialCharTok{*}\NormalTok{d)}\SpecialCharTok{/}\NormalTok{(b}\SpecialCharTok{*}\NormalTok{c)}
\NormalTok{  IC\_OR }\OtherTok{\textless{}{-}} \FunctionTok{exp}\NormalTok{(}\FunctionTok{log}\NormalTok{(OR) }\SpecialCharTok{+} \FunctionTok{c}\NormalTok{(}\SpecialCharTok{{-}}\DecValTok{1}\NormalTok{,}\DecValTok{1}\NormalTok{) }\SpecialCharTok{*}\NormalTok{ z }\SpecialCharTok{*} \FunctionTok{sqrt}\NormalTok{(}\DecValTok{1}\SpecialCharTok{/}\NormalTok{a }\SpecialCharTok{+} \DecValTok{1}\SpecialCharTok{/}\NormalTok{b }\SpecialCharTok{+} \DecValTok{1}\SpecialCharTok{/}\NormalTok{c }\SpecialCharTok{+} \DecValTok{1}\SpecialCharTok{/}\NormalTok{d))}

\NormalTok{  RD    }\OtherTok{\textless{}{-}}\NormalTok{ p1 }\SpecialCharTok{{-}}\NormalTok{ p2}
\NormalTok{  IC\_RD }\OtherTok{\textless{}{-}}\NormalTok{ RD }\SpecialCharTok{+} \FunctionTok{c}\NormalTok{(}\SpecialCharTok{{-}}\DecValTok{1}\NormalTok{,}\DecValTok{1}\NormalTok{) }\SpecialCharTok{*}\NormalTok{ z }\SpecialCharTok{*} \FunctionTok{sqrt}\NormalTok{(p1}\SpecialCharTok{*}\NormalTok{(}\DecValTok{1}\SpecialCharTok{{-}}\NormalTok{p1)}\SpecialCharTok{/}\NormalTok{n1 }\SpecialCharTok{+}\NormalTok{ p2}\SpecialCharTok{*}\NormalTok{(}\DecValTok{1}\SpecialCharTok{{-}}\NormalTok{p2)}\SpecialCharTok{/}\NormalTok{n2)}

  \FunctionTok{cat}\NormalTok{(}\FunctionTok{sprintf}\NormalTok{(}\StringTok{"RR  = \%.3f  IC\%d\%\%: [\%.3f, \%.3f]}\SpecialCharTok{\textbackslash{}n}\StringTok{"}\NormalTok{,}
\NormalTok{              RR, }\FunctionTok{round}\NormalTok{(conf}\SpecialCharTok{*}\DecValTok{100}\NormalTok{), IC\_RR[}\DecValTok{1}\NormalTok{], IC\_RR[}\DecValTok{2}\NormalTok{]))}
  \FunctionTok{cat}\NormalTok{(}\FunctionTok{sprintf}\NormalTok{(}\StringTok{"OR  = \%.3f  IC\%d\%\%: [\%.3f, \%.3f]}\SpecialCharTok{\textbackslash{}n}\StringTok{"}\NormalTok{,}
\NormalTok{              OR, }\FunctionTok{round}\NormalTok{(conf}\SpecialCharTok{*}\DecValTok{100}\NormalTok{), IC\_OR[}\DecValTok{1}\NormalTok{], IC\_OR[}\DecValTok{2}\NormalTok{]))}
  \FunctionTok{cat}\NormalTok{(}\FunctionTok{sprintf}\NormalTok{(}\StringTok{"RD  = \%.3f  IC\%d\%\%: [\%.3f, \%.3f]}\SpecialCharTok{\textbackslash{}n}\StringTok{"}\NormalTok{,}
\NormalTok{              RD, }\FunctionTok{round}\NormalTok{(conf}\SpecialCharTok{*}\DecValTok{100}\NormalTok{), IC\_RD[}\DecValTok{1}\NormalTok{], IC\_RD[}\DecValTok{2}\NormalTok{]))}
  \FunctionTok{cat}\NormalTok{(}\FunctionTok{sprintf}\NormalTok{(}\StringTok{"NNT/NNH = \%.1f}\SpecialCharTok{\textbackslash{}n}\StringTok{"}\NormalTok{, }\DecValTok{1}\SpecialCharTok{/}\FunctionTok{abs}\NormalTok{(RD)))}
\NormalTok{\}}

\CommentTok{\# Cohorte: tabaquismo y cardiopatía isquémica}
\FunctionTok{cat}\NormalTok{(}\StringTok{"=== Cohorte: tabaquismo y cardiopatía ===}\SpecialCharTok{\textbackslash{}n}\StringTok{"}\NormalTok{)}
\end{Highlighting}
\end{Shaded}

\begin{verbatim}
=== Cohorte: tabaquismo y cardiopatía ===
\end{verbatim}

\begin{Shaded}
\begin{Highlighting}[]
\FunctionTok{calcular\_medidas}\NormalTok{(}\AttributeTok{a=}\DecValTok{80}\NormalTok{, }\AttributeTok{b=}\DecValTok{420}\NormalTok{, }\AttributeTok{c=}\DecValTok{30}\NormalTok{, }\AttributeTok{d=}\DecValTok{470}\NormalTok{)}
\end{Highlighting}
\end{Shaded}

\begin{verbatim}
RR  = 2.667  IC95%: [1.786, 3.982]
OR  = 2.984  IC95%: [1.922, 4.632]
RD  = 0.100  IC95%: [0.062, 0.138]
NNT/NNH = 10.0
\end{verbatim}

\subsection{Ejemplo Integrado con Dataset
osteo}\label{ejemplo-integrado-con-dataset-osteo}

\begin{tcolorbox}[enhanced jigsaw, breakable, titlerule=0mm, rightrule=.15mm, leftrule=.75mm, left=2mm, opacitybacktitle=0.6, colback=white, bottomtitle=1mm, colbacktitle=quarto-callout-tip-color!10!white, title=\textcolor{quarto-callout-tip-color}{\faLightbulb}\hspace{0.5em}{Ejemplo: Análisis completo --- datos osteo UGR}, toptitle=1mm, coltitle=black, opacityback=0, arc=.35mm, bottomrule=.15mm, toprule=.15mm, colframe=quarto-callout-tip-color-frame]

El dataset \texttt{osteo} contiene datos de 94 pacientes diabéticos
atendidos en la Facultad de Medicina de la UGR (Pedro Femia et al.).
Realizamos un análisis completo de la asociación entre tabaquismo y
osteoporosis.

\begin{Shaded}
\begin{Highlighting}[]
\FunctionTok{library}\NormalTok{(BioStatR)}

\CommentTok{\# 1. Normalidad del IMC}
\FunctionTok{testnormal}\NormalTok{(osteo}\SpecialCharTok{$}\NormalTok{imc)}
\end{Highlighting}
\end{Shaded}

\pandocbounded{\includegraphics[keepaspectratio]{chapters/semana11_files/figure-pdf/lst-s12-bio07-1.pdf}}

\begin{verbatim}

# Test de normalidad de Shapiro-Wilk  
 ------------------------------------- 
   n = 94,  W = 0.906,  p < 0.001 
\end{verbatim}

\begin{Shaded}
\begin{Highlighting}[]
\CommentTok{\# 2. Tabla de contingencia y prueba Chi{-}cuadrado: tabaquismo × osteoporosis}
\FunctionTok{tabla2x2}\NormalTok{(}\AttributeTok{fvar =}\NormalTok{ osteo}\SpecialCharTok{$}\NormalTok{tabaco, }\AttributeTok{cvar =}\NormalTok{ osteo}\SpecialCharTok{$}\NormalTok{osteo\_cue,}
         \AttributeTok{fcat =} \FunctionTok{c}\NormalTok{(}\StringTok{"No fumador"}\NormalTok{, }\StringTok{"Fumador"}\NormalTok{),}
         \AttributeTok{ccat =} \FunctionTok{c}\NormalTok{(}\StringTok{"Sin osteoporosis"}\NormalTok{, }\StringTok{"Con osteoporosis"}\NormalTok{))}
\end{Highlighting}
\end{Shaded}

\begin{verbatim}

# Análisis de tablas 2x2
# ----------------------

# Frecuencias observadas
               Sin osteoporosis   Con osteoporosis Total
  No fumador                 44                  9    53
  Fumador                    26                 15    41
  Total                      70                 24    94


# Test Chi-cuadrado para un estudio transversal

  χ² = 4.662,   gl = 1,  p = 0.031, (cpc = 0.5) 
  Validez: Frecuencia mínima esperada = 10.47  > 3.9  

  Test exacto de Fisher (bilateral): p = 0.035 

  --- Otros criterios χ²:  
  χ² = 4.673,   gl = 1,  p = 0.054, (sin cpc) 
  χ² = 3.699,   gl = 1,  p = 0.054, (cpc de Yates = 47.00) 

# Estimación de la prevalenciaπ en un estudio transversal
  Método de Wald ajustado: 
  p=0.561; 95%-IC(π)=(0.463, 0.659) 

# Medidas de asociación para un estudio transversal
  [!] Las medidas de riesgo se calculan como riesgo de la categoría  
      en la 1a columna (frente a la 2a) para la categoría en la 1a 
      fila (frente a la 2a)

  Riesgo absoluto (diferencia de Berkson; método de Agresti-Caffo): 
  d=0.254; 95%-IC(d)=(0.023, 0.459) 

  Riesgo relativo: 
  Rr=1.676; 95%-IC(Rr)=(0.969, 2.808) 

  Riesgo atribuible:
  Ra=0.335; 95%-IC(Ra)= (-0.049, 0.578) 

  Razón del producto cruzado (odds ratio):
  OR=2.821; 95%-IC(OR)= (1.070, 7.013) 
\end{verbatim}

\begin{Shaded}
\begin{Highlighting}[]
\CommentTok{\# 3. Tabla de 3 grupos: grupo\_edad × tabaco}
\FunctionTok{tablarxc}\NormalTok{(}\AttributeTok{fvar =}\NormalTok{ osteo}\SpecialCharTok{$}\NormalTok{grupo\_edad, }\AttributeTok{cvar =}\NormalTok{ osteo}\SpecialCharTok{$}\NormalTok{tabaco)}
\end{Highlighting}
\end{Shaded}

\begin{verbatim}

# Test Chi-cuadrado para tablas RxC
# ---------------------------------

# Frecuencias observadas
         A  B Total
A       21 11    32
B       14 16    30
C       18 14    32
  Total 53 41    94

# Test chi-cuadrado 
  Validez: Frecuencia mínima esperada = 13.09 
           0 frecuencias esperadas son menores a 1
           0 son menores a 5 (el 0% de la tabla) 
  χ²(2 gl) = 2.263, p = 0.323 
\end{verbatim}

\end{tcolorbox}

\section{Ajuste por Confusión: El Estimador de
Mantel-Haenszel}\label{ajuste-por-confusiuxf3n-el-estimador-de-mantel-haenszel}

En estudios epidemiológicos, a menudo observamos una asociación entre
una exposición y una enfermedad que parece variar según los niveles de
una tercera variable, llamada \textbf{variable de confusión} (o
\emph{confounder}). Si no controlamos esta variable, podemos llegar a
conclusiones erróneas ---incluso observar una \textbf{Paradoja de
Simpson}, donde una asociación aparente en el total desaparece o se
invierte al estratificar por la variable de confusión.

\subsection{El Estimador de Mantel-Haenszel
(MH)}\label{el-estimador-de-mantel-haenszel-mh}

El estimador de Mantel-Haenszel es una técnica clásica para combinar los
odds ratios (OR) obtenidos de varios estratos (\(k\)) de una variable de
confusión, proporcionando un OR ajustado global. Se calcula como:

\[OR_{MH} = \frac{\sum_{i=1}^{k} \frac{a_i d_i}{n_i}}{\sum_{i=1}^{k} \frac{b_i c_i}{n_i}}\]

Donde para el estrato \(i\): - \(a_i, b_i, c_i, d_i\): celdas de la
tabla 2×2. - \(n_i\): total de individuos en el estrato.

\textbf{Criterio del 10\%:} Para evaluar si una variable es un
confounder clínicamente relevante, comparamos el OR crudo (total) con el
OR ajustado (MH). Si la diferencia entre ambos es \(\geq 10\%\), se
concluye que existe una confusión sustancial, y el OR ajustado por
Mantel-Haenszel debe utilizarse como la medida de asociación más
precisa.

\subsection{Ejemplo Numérico en R: Evaluación de
Confounding}\label{ejemplo-numuxe9rico-en-r-evaluaciuxf3n-de-confounding}

Supongamos un estudio sobre el efecto de un nuevo fármaco (Fármaco A) en
la recuperación de una enfermedad respiratoria, estratificado por la
gravedad del paciente (Leve vs.~Grave).

\begin{Shaded}
\begin{Highlighting}[]
\FunctionTok{library}\NormalTok{(epitools)}

\CommentTok{\# Estrato 1: Leve}
\CommentTok{\# Fármaco A (expuesto), recuperación (caso)}
\NormalTok{a1 }\OtherTok{\textless{}{-}} \DecValTok{40}\NormalTok{; b1 }\OtherTok{\textless{}{-}} \DecValTok{10}\NormalTok{; c1 }\OtherTok{\textless{}{-}} \DecValTok{10}\NormalTok{; d1 }\OtherTok{\textless{}{-}} \DecValTok{40}
\NormalTok{table1 }\OtherTok{\textless{}{-}} \FunctionTok{matrix}\NormalTok{(}\FunctionTok{c}\NormalTok{(a1, b1, c1, d1), }\AttributeTok{nrow =} \DecValTok{2}\NormalTok{, }\AttributeTok{byrow =} \ConstantTok{TRUE}\NormalTok{)}
\FunctionTok{cat}\NormalTok{(}\FunctionTok{sprintf}\NormalTok{(}\StringTok{"Tabla estrato 1 (Leve):}\SpecialCharTok{\textbackslash{}n}\StringTok{"}\NormalTok{)); }\FunctionTok{print}\NormalTok{(table1)}
\end{Highlighting}
\end{Shaded}

\begin{verbatim}
Tabla estrato 1 (Leve):
\end{verbatim}

\begin{verbatim}
     [,1] [,2]
[1,]   40   10
[2,]   10   40
\end{verbatim}

\begin{Shaded}
\begin{Highlighting}[]
\NormalTok{or1 }\OtherTok{\textless{}{-}}\NormalTok{ (a1 }\SpecialCharTok{*}\NormalTok{ d1) }\SpecialCharTok{/}\NormalTok{ (b1 }\SpecialCharTok{*}\NormalTok{ c1)  }\CommentTok{\# OR estrato }
\FunctionTok{print}\NormalTok{(}\FunctionTok{sprintf}\NormalTok{(}\StringTok{"OR estrato 1 (Leve) = \%.3f"}\NormalTok{, or1))}
\end{Highlighting}
\end{Shaded}

\begin{verbatim}
[1] "OR estrato 1 (Leve) = 16.000"
\end{verbatim}

\begin{Shaded}
\begin{Highlighting}[]
\CommentTok{\# Estrato 2: Grave}
\NormalTok{a2 }\OtherTok{\textless{}{-}} \DecValTok{10}\NormalTok{; b2 }\OtherTok{\textless{}{-}} \DecValTok{40}\NormalTok{; c2 }\OtherTok{\textless{}{-}} \DecValTok{5}\NormalTok{; d2 }\OtherTok{\textless{}{-}} \DecValTok{45}
\NormalTok{table2 }\OtherTok{\textless{}{-}} \FunctionTok{matrix}\NormalTok{(}\FunctionTok{c}\NormalTok{(a2, b2, c2, d2), }\AttributeTok{nrow =} \DecValTok{2}\NormalTok{, }\AttributeTok{byrow =} \ConstantTok{TRUE}\NormalTok{)}
\NormalTok{or2 }\OtherTok{\textless{}{-}}\NormalTok{ (a2 }\SpecialCharTok{*}\NormalTok{ d2) }\SpecialCharTok{/}\NormalTok{ (b2 }\SpecialCharTok{*}\NormalTok{ c2)  }\CommentTok{\# OR estrato}
\FunctionTok{cat}\NormalTok{(}\FunctionTok{sprintf}\NormalTok{(}\StringTok{"}\SpecialCharTok{\textbackslash{}n}\StringTok{Tabla estrato 2 (Grave):}\SpecialCharTok{\textbackslash{}n}\StringTok{"}\NormalTok{)); }\FunctionTok{print}\NormalTok{(table2)}
\end{Highlighting}
\end{Shaded}

\begin{verbatim}

Tabla estrato 2 (Grave):
\end{verbatim}

\begin{verbatim}
     [,1] [,2]
[1,]   10   40
[2,]    5   45
\end{verbatim}

\begin{Shaded}
\begin{Highlighting}[]
\FunctionTok{print}\NormalTok{(}\FunctionTok{sprintf}\NormalTok{(}\StringTok{"OR estrato 2 (Grave) = \%.3f"}\NormalTok{, or2))}
\end{Highlighting}
\end{Shaded}

\begin{verbatim}
[1] "OR estrato 2 (Grave) = 2.250"
\end{verbatim}

\begin{Shaded}
\begin{Highlighting}[]
\CommentTok{\# Array de tablas (estratificación)}
\NormalTok{tables }\OtherTok{\textless{}{-}} \FunctionTok{array}\NormalTok{(}\FunctionTok{c}\NormalTok{(a1, b1, c1, d1, a2, b2, c2, d2), }\AttributeTok{dim =} \FunctionTok{c}\NormalTok{(}\DecValTok{2}\NormalTok{, }\DecValTok{2}\NormalTok{, }\DecValTok{2}\NormalTok{))}
\NormalTok{tables }\OtherTok{\textless{}{-}} \FunctionTok{aperm}\NormalTok{(tables, }\FunctionTok{c}\NormalTok{(}\DecValTok{2}\NormalTok{, }\DecValTok{1}\NormalTok{, }\DecValTok{3}\NormalTok{))  }\CommentTok{\# Reordenar dimensiones para epitools}

\CommentTok{\# 1. OR Crudo (sumamos los estratos)}
\NormalTok{total\_table }\OtherTok{\textless{}{-}}\NormalTok{ table1 }\SpecialCharTok{+}\NormalTok{ table2}
\NormalTok{or\_crudo }\OtherTok{\textless{}{-}}\NormalTok{ (total\_table[}\DecValTok{1}\NormalTok{,}\DecValTok{1}\NormalTok{] }\SpecialCharTok{*}\NormalTok{ total\_table[}\DecValTok{2}\NormalTok{,}\DecValTok{2}\NormalTok{]) }\SpecialCharTok{/}\NormalTok{ (total\_table[}\DecValTok{1}\NormalTok{,}\DecValTok{2}\NormalTok{] }\SpecialCharTok{*}\NormalTok{ total\_table[}\DecValTok{2}\NormalTok{,}\DecValTok{1}\NormalTok{])}
\FunctionTok{cat}\NormalTok{(}\FunctionTok{sprintf}\NormalTok{(}\StringTok{"Tabla total:}\SpecialCharTok{\textbackslash{}n}\StringTok{"}\NormalTok{)); }\FunctionTok{print}\NormalTok{(total\_table)}
\end{Highlighting}
\end{Shaded}

\begin{verbatim}
Tabla total:
\end{verbatim}

\begin{verbatim}
     [,1] [,2]
[1,]   50   50
[2,]   15   85
\end{verbatim}

\begin{Shaded}
\begin{Highlighting}[]
\CommentTok{\# 2. OR Ajustado por Mantel{-}Haenszel}
\NormalTok{mh }\OtherTok{\textless{}{-}} \FunctionTok{mantelhaen.test}\NormalTok{(tables, }\AttributeTok{correct =} \ConstantTok{FALSE}\NormalTok{)}
\NormalTok{or\_mh }\OtherTok{\textless{}{-}}\NormalTok{ mh}\SpecialCharTok{$}\NormalTok{estimate}

\FunctionTok{cat}\NormalTok{(}\FunctionTok{sprintf}\NormalTok{(}\StringTok{"OR Crudo: \%.3f}\SpecialCharTok{\textbackslash{}n}\StringTok{"}\NormalTok{, or\_crudo))}
\end{Highlighting}
\end{Shaded}

\begin{verbatim}
OR Crudo: 5.667
\end{verbatim}

\begin{Shaded}
\begin{Highlighting}[]
\FunctionTok{cat}\NormalTok{(}\FunctionTok{sprintf}\NormalTok{(}\StringTok{"OR Mantel{-}Haenszel: \%.3f}\SpecialCharTok{\textbackslash{}n}\StringTok{"}\NormalTok{, or\_mh))}
\end{Highlighting}
\end{Shaded}

\begin{verbatim}
OR Mantel-Haenszel: 6.833
\end{verbatim}

\begin{Shaded}
\begin{Highlighting}[]
\CommentTok{\# 3. Aplicación de la regla del 10\%}
\NormalTok{diff\_pct }\OtherTok{\textless{}{-}} \FunctionTok{abs}\NormalTok{(or\_crudo }\SpecialCharTok{{-}}\NormalTok{ or\_mh) }\SpecialCharTok{/}\NormalTok{ or\_crudo }\SpecialCharTok{*} \DecValTok{100}
\FunctionTok{cat}\NormalTok{(}\FunctionTok{sprintf}\NormalTok{(}\StringTok{"Diferencia relativa: \%.1f\%\%}\SpecialCharTok{\textbackslash{}n}\StringTok{"}\NormalTok{, diff\_pct))}
\end{Highlighting}
\end{Shaded}

\begin{verbatim}
Diferencia relativa: 20.6%
\end{verbatim}

\textbf{Interpretación Epidemiológica:} En este estudio, el OR crudo
(total) es 5.67, mientras que el OR ajustado (MH) es 6.83. La diferencia
relativa es de 20.6\%. Al ser mayor del 10\%, existe un efecto de
confusión sustancial provocado por la ``gravedad''.

La gravedad actúa como un factor de confusión ya que está asociada tanto
a la probabilidad de recibir el Fármaco A (sesgo de indicación) como a
la probabilidad de recuperación (pronóstico). El OR crudo subestima (o
sobreestima) el efecto real al mezclar el impacto del fármaco con la
mayor o menor gravedad basal de los pacientes. Al ajustar por MH,
eliminamos esta confusión.

Sin embargo, es crucial notar que existe una \textbf{interacción}
(modificación del efecto): el fármaco parece ser más eficaz en pacientes
leves que en pacientes graves. En presencia de una interacción tan
marcada, la mejor recomendación epidemiológica es presentar los
\textbf{resultados estratificados}, ya que un único OR ajustado (MH)
ocultaría esta heterogeneidad clínicamente relevante.

Por último, para superar las limitaciones de los métodos tradicionales,
la generalización de la estandarización mediante la \textbf{g-formula}
(Robins, 1987) representa hoy en día una alternativa avanzada y robusta
para ajustar el efecto del fármaco. Para aquellos interesados en
introducirse en la inferencia causal desde una perspectiva tanto
histórica como computacional, se recomienda como punto de partida el
tutorial \emph{Computational Causal Inference} (Smith \&
Luque-Fernandez, 2020, publicado en \emph{Statistics in Medicine}), el
cual ofrece una excelente guía con ejemplos prácticos implementados en
\textbf{R, Python y Stata}.

\subsection{\texorpdfstring{Control multivariable de la confusión:
regresión logística múltiple con
\texttt{rlogitm()}}{Control multivariable de la confusión: regresión logística múltiple con rlogitm()}}\label{sec-semana11-rlogitm}

El estimador de Mantel--Haenszel es elegante y didácticamente potente,
pero presenta dos limitaciones operativas importantes:

\begin{enumerate}
\def\labelenumi{\arabic{enumi}.}
\tightlist
\item
  \textbf{Sólo permite controlar por una variable de confusión
  categórica} (la que define los estratos). Si necesitamos ajustar
  simultáneamente por edad, IMC, tabaquismo y tiempo de evolución, la
  estratificación produce demasiados estratos con frecuencias muy
  pequeñas en cada uno.
\item
  \textbf{Asume homogeneidad del OR a través de los estratos}. Cuando
  esa hipótesis se rompe (interacción, como en el ejemplo previo), el OR
  de M--H promedia efectos heterogéneos y oculta la modificación del
  efecto.
\end{enumerate}

La \textbf{regresión logística múltiple} generaliza el ajuste por
confusión: permite incluir un número arbitrario de covariables
(continuas o categóricas) simultáneamente, estima un OR ajustado para
cada predictor manteniendo constantes las demás, y proporciona una
medida global de capacidad discriminante mediante el \textbf{área bajo
la curva ROC} (AUC). El paquete \texttt{BioStatR} integra todo este
flujo en la función \texttt{rlogitm()} (ver
Apéndice~\ref{sec-apendiceB}), que devuelve:

\begin{itemize}
\tightlist
\item
  Estimación de los coeficientes con sus odds ratios e intervalos de
  confianza.
\item
  Bondad de ajuste mediante el \textbf{test de Hosmer--Lemeshow}.
\item
  Pseudo-\(R^2\) de \textbf{Nagelkerke}.
\item
  \textbf{AUC} con su intervalo de confianza al 95 \%.
\item
  Curva ROC (opcional, con \texttt{grf\ =\ TRUE}).
\end{itemize}

\begin{tcolorbox}[enhanced jigsaw, breakable, titlerule=0mm, rightrule=.15mm, leftrule=.75mm, left=2mm, opacitybacktitle=0.6, colback=white, bottomtitle=1mm, colbacktitle=quarto-callout-tip-color!10!white, title=\textcolor{quarto-callout-tip-color}{\faLightbulb}\hspace{0.5em}{Ejemplo: osteoporosis del cuello femoral ajustada por múltiples factores}, toptitle=1mm, coltitle=black, opacityback=0, arc=.35mm, bottomrule=.15mm, toprule=.15mm, colframe=quarto-callout-tip-color-frame]

Volvamos al estudio observacional de la sección de tablas de
contingencia, donde habíamos detectado una asociación entre tabaquismo y
osteoporosis del cuello femoral (\texttt{osteo\_cue}) en 94 pacientes
diabéticos. La pregunta clínica relevante es: \textbf{¿persiste la
asociación tras controlar simultáneamente por edad, índice de masa
corporal y tiempo de evolución de la diabetes?}

\begin{Shaded}
\begin{Highlighting}[]
\FunctionTok{library}\NormalTok{(BioStatR)}
\FunctionTok{data}\NormalTok{(osteo)}

\CommentTok{\# Modelo: osteoporosis del cuello femoral ajustada por tabaco + edad + IMC + tevol}
\FunctionTok{rlogitm}\NormalTok{(osteo\_cue }\SpecialCharTok{\textasciitilde{}}\NormalTok{ tabaco }\SpecialCharTok{+}\NormalTok{ edad }\SpecialCharTok{+}\NormalTok{ imc }\SpecialCharTok{+}\NormalTok{ tevol,}
        \AttributeTok{data =}\NormalTok{ osteo,}
        \AttributeTok{grf  =} \ConstantTok{TRUE}\NormalTok{)}
\end{Highlighting}
\end{Shaded}

\begin{verbatim}

Regresión logística  multivariable
----------------------------------------------------------------
# Información muestral --- 

   Tamaño muestral (N inicial) :  94 
   Tamaño muestral tras eliminar valores perdidos (Casos completos) :  94 
   Mínima frecuencia de eventos (n efectivo) :  24 

# Distribución de la variable respuesta (osteo_cue) ---

  Categoria  n Porcentaje
1        No 70       74.5
2        Sí 24       25.5

# Modelo logístico ---  --- 

   Modelo :  osteo_cue ~ tabaco + edad + imc + tevol 
  Devianza residual:  92  (Nula:  107 )
  AIC:  102 
  R² de Nagelkerke:  0.214 

   Test de bondad de ajuste de Hosmer-Lemeshow :
  X² =  2.07 , gl =  8 ,   = 0.979 

   Capacidad discriminante :
   AUC (Area bajo la curva ROC)  =  0.74 

   Coeficientes del modelo :

      Termino Estimacion Error_Std  z_exp      sig     OR OR_inf   OR_sup
1 (Intercept)      2.893     2.269  1.275  = 0.202 18.045  0.275 2237.565
2    tabacoSí      0.865     0.549  1.575  = 0.115  2.376  0.821    7.218
3        edad     -0.004     0.037 -0.105  = 0.917  0.996  0.925    1.072
4         imc     -0.224     0.105 -2.140  = 0.032  0.799  0.638    0.965
5       tevol      0.071     0.035  2.020  = 0.043  1.074  1.004    1.154
\end{verbatim}

\pandocbounded{\includegraphics[keepaspectratio]{chapters/semana11_files/figure-pdf/lst-s11-rlogitm-1.pdf}}

\textbf{Interpretación epidemiológica:}

\begin{itemize}
\item
  \textbf{Tabaquismo (variable de exposición principal):} tras ajustar
  por edad, IMC y tiempo de evolución, los fumadores presentan un OR
  ajustado \(\approx 2.4\) (IC 95 \%: 0.82 -- 7.22) frente a los no
  fumadores. La asociación cruda obtenida en la tabla 2×2 (OR
  \(\approx 2.8\), \(p = 0.031\)) se atenúa ligeramente al ajustar por
  las demás covariables y pierde significación estadística al 5 \%, lo
  que sugiere que parte del efecto observado podría estar mediado o
  confundido por la edad y el IMC. Este es un ejemplo clásico de cómo el
  \textbf{OR ajustado por regresión multivariable} difiere del OR crudo
  y, también, del OR de Mantel--Haenszel restringido a un único
  estratificador.
\item
  \textbf{IMC:} muestra una asociación inversa estadísticamente
  significativa (\(\text{OR} \approx 0.80\), \(p \approx 0.03\)). Cada
  unidad adicional de IMC reduce las odds de osteoporosis del cuello
  femoral aproximadamente un 20 \%, en consonancia con la literatura
  sobre el efecto protector de la masa corporal sobre la densidad ósea.
\item
  \textbf{Tiempo de evolución de la diabetes (\texttt{tevol}):}
  asociación positiva significativa (\(\text{OR} \approx 1.07\),
  \(p \approx 0.04\)). Cada año adicional de evolución incrementa las
  odds de osteoporosis del cuello femoral en torno a un 7 \%.
\item
  \textbf{Edad:} no aporta información independiente del resto de
  predictores (\(p \approx 0.92\)), probablemente porque su efecto está
  capturado por \texttt{tevol} (variable altamente correlacionada con la
  edad en pacientes diabéticos).
\end{itemize}

\textbf{Bondad de ajuste y discriminación:}

\begin{itemize}
\tightlist
\item
  El test de Hosmer--Lemeshow no rechaza el ajuste del modelo
  (\(\chi^2 \approx 2.07\), \(p \approx 0.98\)): las probabilidades
  predichas son coherentes con las frecuencias observadas en deciles de
  riesgo.
\item
  El pseudo-\(R^2\) de Nagelkerke alcanza \(\approx 0.21\), un valor
  moderado para datos observacionales en bioestadística clínica.
\item
  El AUC de la curva ROC es \(\approx 0.74\), indicando una
  \textbf{capacidad discriminante aceptable}: en términos prácticos, si
  tomamos al azar a un paciente con osteoporosis y a otro sin
  osteoporosis, la probabilidad predicha de osteoporosis es mayor en el
  primero el 74 \% de las veces.
\end{itemize}

\end{tcolorbox}

\begin{tcolorbox}[enhanced jigsaw, breakable, titlerule=0mm, rightrule=.15mm, leftrule=.75mm, left=2mm, opacitybacktitle=0.6, colback=white, bottomtitle=1mm, colbacktitle=quarto-callout-note-color!10!white, title=\textcolor{quarto-callout-note-color}{\faInfo}\hspace{0.5em}{Mantel--Haenszel vs.~regresión logística múltiple}, toptitle=1mm, coltitle=black, opacityback=0, arc=.35mm, bottomrule=.15mm, toprule=.15mm, colframe=quarto-callout-note-color-frame]

\begin{longtable}[]{@{}
  >{\raggedright\arraybackslash}p{(\linewidth - 4\tabcolsep) * \real{0.1754}}
  >{\centering\arraybackslash}p{(\linewidth - 4\tabcolsep) * \real{0.2982}}
  >{\centering\arraybackslash}p{(\linewidth - 4\tabcolsep) * \real{0.5263}}@{}}
\toprule\noalign{}
\begin{minipage}[b]{\linewidth}\raggedright
Criterio
\end{minipage} & \begin{minipage}[b]{\linewidth}\centering
Mantel--Haenszel
\end{minipage} & \begin{minipage}[b]{\linewidth}\centering
Regresión logística múltiple
\end{minipage} \\
\midrule\noalign{}
\endhead
\bottomrule\noalign{}
\endlastfoot
Nº de covariables de ajuste & 1 (categórica, define estratos) &
Arbitrario (continuas y categóricas) \\
Supuesto sobre OR & Homogeneidad entre estratos & Linealidad en el
log-odds; no interacción salvo que se modelice \\
Output principal & OR ajustado único & OR ajustado por covariable +
bondad de ajuste + AUC \\
Detección de interacción & Test de Breslow--Day & Términos producto en
la fórmula \\
Cuando elegirlo & Análisis exploratorio rápido, una sola covariable de
confusión & Análisis confirmatorio, múltiples confusores, predicción \\
\end{longtable}

En la práctica epidemiológica moderna ambas técnicas se utilizan de
forma complementaria: M--H ofrece una \textbf{lectura rápida y
transparente} del efecto de una variable de exposición estratificada,
mientras que la \textbf{regresión logística múltiple} proporciona el
modelo de ajuste definitivo con todas las covariables relevantes y una
medida global de discriminación (AUC).

\end{tcolorbox}

\section{Ejercicios del tema}\label{ejercicios-del-tema}

\textbf{Ejercicio 11.1:} En un estudio sobre el consumo de alcohol y
diabetes tipo 2 se obtiene la siguiente tabla 2×2:

\begin{longtable}[]{@{}lccc@{}}
\toprule\noalign{}
& Diabetes & Sin diabetes & Total \\
\midrule\noalign{}
\endhead
\bottomrule\noalign{}
\endlastfoot
Bebedor habitual & 45 & 105 & 150 \\
No bebedor & 20 & 130 & 150 \\
\end{longtable}

\begin{enumerate}
\def\labelenumi{\alph{enumi})}
\tightlist
\item
  Calcule las frecuencias esperadas bajo \(H_0\). b) Calcule el
  estadístico \(\chi^2\) manualmente. c) ¿Existe asociación con
  \(\alpha = 0.05\)?
\end{enumerate}

\textbf{Ejercicio 11.2:} Un ensayo clínico compara tres tratamientos
para la hipertensión: A (control), B (dosis baja), C (dosis alta). Los
resultados son:

\begin{longtable}[]{@{}lcc@{}}
\toprule\noalign{}
Tratamiento & Controlada & No controlada \\
\midrule\noalign{}
\endhead
\bottomrule\noalign{}
\endlastfoot
A & 25 & 35 \\
B & 40 & 20 \\
C & 45 & 15 \\
\end{longtable}

Usando la prueba \(\chi^2\) de homogeneidad con \(\alpha = 0.05\),
¿difieren los tratamientos en su eficacia? (\(gl = 2\),
\(\chi^2_{0.05} = 5.99\)).

\textbf{Ejercicio 11.3:} Una tabla 2×2 tiene celdas \((3, 12, 8, 7)\).
a) Compruebe si se cumple la condición de aplicación del \(\chi^2\). b)
¿Qué alternativa utilizaría? c) Aplique la alternativa e interprete.

\textbf{Ejercicio 11.4 (McNemar):} En un estudio de concordancia
diagnóstica, 80 pacientes con sospecha de infección por
\emph{Helicobacter pylori} son evaluados con dos pruebas: test de
aliento con urea (TAU) y serología (SERO). Los resultados son:

\begin{longtable}[]{@{}lccc@{}}
\toprule\noalign{}
& \textbf{SERO: +} & \textbf{SERO: −} & \textbf{Total} \\
\midrule\noalign{}
\endhead
\bottomrule\noalign{}
\endlastfoot
\textbf{TAU: +} & 38 & 12 & 50 \\
\textbf{TAU: −} & 5 & 25 & 30 \\
\textbf{Total} & 43 & 37 & 80 \\
\end{longtable}

\begin{enumerate}
\def\labelenumi{\alph{enumi})}
\tightlist
\item
  ¿Por qué es incorrecto aplicar el test \(\chi^2\) de independencia
  aquí?
\item
  Identifique los pares discordantes (\(b\) y \(c\)) e indique cuántos
  son.
\item
  Calcule el estadístico de McNemar (con y sin corrección de Yates).
\item
  ¿Difieren significativamente las dos pruebas diagnósticas en su tasa
  de positividad? ¿Cuál detecta más casos?
\end{enumerate}

\textbf{Ejercicio 11.5:} Un estudio de cohortes sigue a 1000
trabajadores expuestos a un contaminante y 1000 no expuestos durante 5
años. Se observan 60 casos de enfermedad pulmonar en expuestos y 20 en
no expuestos. a) Calcule el RR e interprete. b) Calcule el IC 95\%. c)
Calcule el RD y el NNH.

\textbf{Ejercicio 11.6:} En un estudio caso-control sobre infarto de
miocardio, se estudia la asociación con hipertensión. Se reclutan 200
casos (infartos) y 200 controles. De los casos, 120 son hipertensos; de
los controles, 80. a) Construya la tabla 2×2. b) Calcule el OR. c)
Calcule el IC 95\% del OR. d) ¿Por qué no es adecuado calcular el RR en
este diseño?

\textbf{Ejercicio 11.7:} En un estudio transversal sobre la prevalencia
de diabetes y obesidad en 800 adultos, 200 son obesos (IMC \(\geq 30\))
y 600 no obesos. Entre los obesos, 60 tienen diabetes; entre los no
obesos, 40 tienen diabetes. a) Calcule la RP. b) Calcule el OR. c) ¿Cuál
es la medida más apropiada? d) ¿Son similares los resultados? Explique
por qué.

\begin{center}\rule{0.5\linewidth}{0.5pt}\end{center}

\section{Respuestas a los Ejercicios}\label{sec-semana11-respuestas}

\textbf{Ejercicio 11.1:} a) \(E_{11} = 150 \times 65/300 = 32.5\);
\(E_{12} = 150 \times 235/300 = 117.5\); \(E_{21} = 32.5\);
\(E_{22} = 117.5\). b)
\(\chi^2 = (45-32.5)^2/32.5 + (105-117.5)^2/117.5 + (20-32.5)^2/32.5 + (130-117.5)^2/117.5 \approx 4.81 + 1.33 + 4.81 + 1.33 = 12.28\).
c) \(\chi^2_{0.05, 1} = 3.84\). Como \(12.28 > 3.84\) (\(p < 0.001\)),
rechazamos \(H_0\): existe asociación significativa entre consumo de
alcohol y diabetes tipo 2.

\textbf{Ejercicio 11.2:} Frecuencias esperadas: \(n = 180\), marginales
columna: 110 controladas, 70 no controladas.
\(E_{A,ctrl} = 60 \times 110/180 = 36.67\); calculando todos:
\(\chi^2 \approx 15.2 > 5.99\), \(gl = 2\), \(p < 0.001\). Los tres
tratamientos tienen eficacias significativamente diferentes. La dosis
alta (75\% controlados) es la más eficaz.

\textbf{Ejercicio 11.3:} a) \(E_{11} = 11 \times 11/30 = 4.03 < 5\);
condición NO cumplida. b) Usar \textbf{prueba exacta de Fisher}. c) Con
\(n = 30\) y celdas \((3, 12, 8, 7)\):
\(OR = 3 \times 7/(12 \times 8) = 0.22\). La prueba exacta de Fisher
daría \(p \approx 0.07\), no significativo al 5\%.

\textbf{Ejercicio 11.4 (McNemar):} a) Las dos pruebas se aplican al
mismo paciente, por lo que las observaciones son dependientes entre sí.
El test \(\chi^2\) de independencia requiere que todas las observaciones
sean independientes; violarlo inflaría artificialmente el tamaño de
muestra efectivo. b) Pares discordantes: \(b = 12\) (TAU+/SERO−) y
\(c = 5\) (TAU−/SERO+). Total: \(b + c = 17\). c) Como
\(b + c = 17 < 25\), se recomienda la corrección de Yates: - Sin
corrección: \(\chi^2 = (12-5)^2/17 = 49/17 = 2.88\) - Con corrección
(Yates): \(\chi^2 = (|12-5|-1)^2/17 = 36/17 = 2.12\) d) Sin corrección:
\(\chi^2 = 2.88\), \(p \approx 0.090\). Con corrección de Yates:
\(\chi^2 = 2.12\), \(p \approx 0.145\). Con \(\alpha = 0.05\),
\textbf{no rechazamos} \(H_0\) en ningún caso: no hay evidencia
suficiente de que las dos pruebas difieran significativamente en su tasa
de positividad. No obstante, el TAU detecta más positivos que la
serología (\(b = 12 > c = 5\)), lo que sugiere una tendencia a mayor
sensibilidad del TAU, aunque no alcanza significancia estadística con
esta muestra.

\textbf{Ejercicio 11.5:} a) \(RR = (60/1000)/(20/1000) = 3.0\). Los
expuestos tienen \textbf{3 veces más riesgo} de enfermedad pulmonar. b)
\(SE_{\ln RR} = \sqrt{940/(60 \times 1000) + 980/(20 \times 1000)} = \sqrt{0.01567 + 0.049} = 0.256\);
\(IC_{95\%} = \exp(\ln3 \pm 0.501) = [1.50, 5.99]\). c)
\(RD = 0.060 - 0.020 = 0.040\); \(NNH = 1/0.040 = 25\) (por cada 25
trabajadores expuestos durante 5 años, 1 caso adicional de enfermedad
pulmonar).

\textbf{Ejercicio 11.6:} a) Tabla: hipertensos \((120, 80, 200)\);
normotensos \((80, 120, 200)\); totales \((200, 200, 400)\). b)
\(OR = (120 \times 120)/(80 \times 80) = 14400/6400 = 2.25\). c)
\(IC_{95\%}: \exp(\ln 2.25 \pm 1.96\sqrt{1/120 + 1/80 + 1/80 + 1/120}) = \exp(0.811 \pm 0.369) = [1.44, 3.51]\).
d) En un estudio caso-control, el muestreo es por resultado (se
seleccionan casos y controles), por lo que los totales de fila son
fijados por el investigador y no reflejan la incidencia real → no se
puede estimar el riesgo ni el RR directamente.

\textbf{Ejercicio 11.7:} a)
\(RP = (60/200)/(40/600) = 0.30/0.067 = 4.50\). Los obesos tienen 4.5
veces la prevalencia de diabetes. b)
\(OR = (60 \times 560)/(140 \times 40) = 33600/5600 = 6.0\). c) La
medida apropiada en un estudio transversal es la \textbf{RP}, no el OR.
d) Los resultados difieren notablemente (\(RP = 4.5\) vs \(OR = 6.0\))
porque la prevalencia de diabetes no es rara (50\% en obesos): el OR
sobreestima la RP cuando la prevalencia es alta (\textgreater{} 10\%).
Usar siempre RP en estudios transversales.

\begin{center}\rule{0.5\linewidth}{0.5pt}\end{center}

\section{Recomendaciones de Métodos
Avanzados}\label{sec-semana11-avanzados}

\begin{tcolorbox}[enhanced jigsaw, breakable, titlerule=0mm, rightrule=.15mm, leftrule=.75mm, left=2mm, opacitybacktitle=0.6, colback=white, bottomtitle=1mm, colbacktitle=quarto-callout-tip-color!10!white, title=\textcolor{quarto-callout-tip-color}{\faLightbulb}\hspace{0.5em}{Tip}, toptitle=1mm, coltitle=black, opacityback=0, arc=.35mm, bottomrule=.15mm, toprule=.15mm, colframe=quarto-callout-tip-color-frame]

Para profundizar en estos temas, se recomienda consultar textos
especializados como \emph{Categorical Data Analysis} de Alan Agresti o
literatura específica en bioestadística avanzada enfocada en inferencia
causal.

Para ampliar los contenidos de este capítulo con técnicas estadísticas
avanzadas, visita:

\textbf{\href{https://migariane.github.io/Bioestadistica}{→
Bioestadística Avanzada --- M.A.~Luque Fernández}}

\end{tcolorbox}

\begin{center}\rule{0.5\linewidth}{0.5pt}\end{center}

\section{Lecturas Recomendadas
Adicionales}\label{sec-semana11-resources}

\begin{longtable}[]{@{}
  >{\raggedright\arraybackslash}p{(\linewidth - 4\tabcolsep) * \real{0.1250}}
  >{\raggedright\arraybackslash}p{(\linewidth - 4\tabcolsep) * \real{0.2292}}
  >{\raggedright\arraybackslash}p{(\linewidth - 4\tabcolsep) * \real{0.6458}}@{}}
\toprule\noalign{}
\begin{minipage}[b]{\linewidth}\raggedright
Tema
\end{minipage} & \begin{minipage}[b]{\linewidth}\raggedright
Referencia
\end{minipage} & \begin{minipage}[b]{\linewidth}\raggedright
Relevancia para este capítulo
\end{minipage} \\
\midrule\noalign{}
\endhead
\bottomrule\noalign{}
\endlastfoot
Datos categóricos & \citet{Agresti2013} & Tablas de contingencia, χ²,
modelos log-lineales \\
Datos apareados & \citet{McNemar1947} & Artículo original de la prueba
de McNemar \\
Tasas y proporciones & \citet{Fleiss2003} & IC exactos, OR, RR y prueba
de McNemar \\
Estadística médica & \citet{Altman1991} & Tablas 2×2 y medidas de
asociación \\
Epidemiología & \citet{Rothman2008} & RR, OR, diseños
cohorte/caso-control \\
Inferencia Causal & \citet{Robins1987} & G-formula para
estandarización \\
Inferencia Causal & \citet{Smith2020} & Tutorial de inferencia causal
computacional (R, Python, Stata) \\
Epidemiología analítica & \citet{Szklo2007} & Medidas de efecto y
diseños epidemiológicos \\
Introducción estadística & \citet{OpenIntro2019} & Inferencia sobre
proporciones \\
\end{longtable}

\bookmarksetup{startatroot}

\chapter*{Referencias}\label{referencias}
\addcontentsline{toc}{chapter}{Referencias}

\markboth{Referencias}{Referencias}

\section*{Sobre esta sección}\label{sobre-esta-secciuxf3n}
\addcontentsline{toc}{section}{Sobre esta sección}

\markright{Sobre esta sección}

Esta página reúne las referencias bibliográficas utilizadas en el curso
\textbf{Matemáticas para la Estadística}. Los materiales están
organizados por temas para facilitar la búsqueda según los tópicos de
interés.

\section*{Libros de Referencia
General}\label{libros-de-referencia-general}
\addcontentsline{toc}{section}{Libros de Referencia General}

\markright{Libros de Referencia General}

Los siguientes textos proporcionan cobertura completa de los temas
tratados en el curso:

\begin{itemize}
\tightlist
\item
  \citet{Casella2002} --- Texto clásico y riguroso sobre inferencia
  estadística, fundamental para cursos avanzados.
\item
  \citet{Rice2007} --- Excelente balance entre teoría y práctica con
  muchos ejemplos.
\item
  \citet{Wackerly2008} --- Orientado a aplicaciones en ingeniería y
  ciencias, con numerosos ejercicios.
\item
  \citet{Wasserman2006} --- Curso conciso pero riguroso en inferencia
  estadística.
\item
  \citet{Wood2015} --- Introducción accesible a estadística desde una
  perspectiva moderna.
\end{itemize}

\section*{Análisis Exploratorio de
Datos}\label{anuxe1lisis-exploratorio-de-datos}
\addcontentsline{toc}{section}{Análisis Exploratorio de Datos}

\markright{Análisis Exploratorio de Datos}

Para los tópicos de la Semana 1 y 2 (EDA):

\begin{itemize}
\tightlist
\item
  \citet{Tukey1961} --- Artículo seminal que define el Análisis
  Exploratorio de Datos (EDA).
\item
  \citet{Wickham2016} --- Guía moderna para manipulación y visualización
  de datos en R.
\end{itemize}

\section*{Probabilidad y Variables
Aleatorias}\label{probabilidad-y-variables-aleatorias}
\addcontentsline{toc}{section}{Probabilidad y Variables Aleatorias}

\markright{Probabilidad y Variables Aleatorias}

Referencias para profundizar en los tópicos de las Semanas 3 y 4:

\begin{itemize}
\tightlist
\item
  \citet{Rice2007} --- Excelentes capítulos sobre probabilidad y
  variables aleatorias.
\item
  \citet{Ross2014} --- Enfoque práctico para ingenieros y científicos.
\item
  \citet{Devore2020} --- Texto ampliamente usado en cursos de
  ingeniería.
\end{itemize}

\section*{Teoría del Muestreo e Inferencia
Estadística}\label{teoruxeda-del-muestreo-e-inferencia-estaduxedstica}
\addcontentsline{toc}{section}{Teoría del Muestreo e Inferencia
Estadística}

\markright{Teoría del Muestreo e Inferencia Estadística}

Para las Semanas 5, 7, 8 y 9:

\begin{itemize}
\tightlist
\item
  \citet{Casella2002} --- Riguroso tratamiento de estimación y
  contrastes de hipótesis.
\item
  \citet{Agresti2018} --- Métodos estadísticos aplicados a ciencias
  sociales.
\item
  \citet{FDR2001} --- Control de tasa de falsos positivos en pruebas
  múltiples.
\end{itemize}

\section*{Regresión Lineal}\label{regresiuxf3n-lineal}
\addcontentsline{toc}{section}{Regresión Lineal}

\markright{Regresión Lineal}

Para las Semanas 10 y 11:

\begin{itemize}
\tightlist
\item
  \citet{James2013} --- Introducción moderna a aprendizaje estadístico
  con énfasis en regresión.
\item
  \citet{Freedman2007} --- Enfoque crítico y reflexivo sobre la práctica
  de la regresión.
\item
  \citet{Wooldridge2020} --- Aplicaciones econométricas de modelos de
  regresión.
\item
  \citet{Greene2018} --- Referencia avanzada en econometría.
\item
  \citet{Angrist2008} --- Métodos empíricos para inferencia causal.
\end{itemize}

\section*{Análisis de Datos Categóricos y
Epidemiología}\label{anuxe1lisis-de-datos-categuxf3ricos-y-epidemiologuxeda}
\addcontentsline{toc}{section}{Análisis de Datos Categóricos y
Epidemiología}

\markright{Análisis de Datos Categóricos y Epidemiología}

Para la Semana 12 (tablas de contingencia, Chi-cuadrado, McNemar,
medidas de asociación):

\begin{itemize}
\tightlist
\item
  \citet{Agresti2013} --- Referencia estándar para el análisis de datos
  categóricos: tablas de contingencia, pruebas Chi-cuadrado y modelos
  log-lineales.
\item
  \citet{Fleiss2003} --- Tratamiento riguroso de proporciones, tasas y
  pruebas para tablas 2×2, incluida la prueba de McNemar para datos
  apareados.
\item
  \citet{McNemar1947} --- Artículo original de Quinn McNemar donde se
  propone la prueba para muestras relacionadas.
\item
  \citet{Altman1991} --- Estadística médica aplicada: medidas de
  asociación (RR, OR, RD), tablas 2×2 y pruebas para datos apareados.
\item
  \citet{Rothman2008} --- Epidemiología moderna: diseños de cohorte,
  caso-control y transversal; estimación e interpretación de RR y OR.
\item
  \citet{Szklo2007} --- Presentación accesible de epidemiología
  analítica con énfasis en medidas de efecto.
\item
  \citet{BioestadisticaR2} --- Paquete R del grupo de Bioestadística de
  la UGR (Pedro Femia et al.): funciones \texttt{tabla2x2()},
  \texttt{tablarxc()} y \texttt{testnormal()} con el dataset
  \texttt{osteo}.
\item
  \citet{OpenIntro2019} --- Introducción accesible a la inferencia sobre
  proporciones y tablas de contingencia.
\end{itemize}

\section*{Recursos en R}\label{recursos-en-r}
\addcontentsline{toc}{section}{Recursos en R}

\markright{Recursos en R}

El curso utiliza ampliamente \textbf{R} para la implementación de
conceptos estadísticos:

\begin{itemize}
\tightlist
\item
  \citet{RCoreTeam2024} --- Documentación oficial del proyecto R.
\item
  \citet{Wickham2016} --- Guía práctica para análisis de datos con R.
\item
  \citet{Kabacoff2015} --- Referencia completa de gráficos y
  manipulación en R.
\end{itemize}

\section*{Perspectivas Bayesianas}\label{perspectivas-bayesianas}
\addcontentsline{toc}{section}{Perspectivas Bayesianas}

\markright{Perspectivas Bayesianas}

Para estudiantes interesados en métodos Bayesianos:

\begin{itemize}
\tightlist
\item
  \citet{Koop2011} --- Introducción a econometría Bayesiana.
\item
  \citet{Gelman2013} --- Texto completo sobre análisis Bayesiano de
  datos.
\item
  \citet{McElreath2020} --- Curso moderno en inferencia Bayesiana con
  Stan.
\item
  \citet{Peng2015} --- Pensamiento estadístico desde perspectiva
  Bayesiana.
\end{itemize}

\section*{Tópicos Especializados}\label{tuxf3picos-especializados}
\addcontentsline{toc}{section}{Tópicos Especializados}

\markright{Tópicos Especializados}

Lecturas complementarias para temas avanzados:

\begin{itemize}
\tightlist
\item
  \citet{James2013} --- Aprendizaje estadístico (machine learning
  clásico).
\item
  \citet{Efron2016} --- Perspectiva histórica sobre estadística en la
  era computacional.
\item
  \citet{Hyndman2018} --- Métodos de pronóstico y series de tiempo.
\item
  \citet{Pearl2009} --- Fundamentos matemáticos de causalidad.
\item
  \citet{Varian2014} --- Aplicaciones de estadística a economía y ``big
  data''.
\item
  \citet{Montgomery2020} --- Estadística aplicada a ingeniería.
\item
  \citet{OpenIntro2019} --- Recurso abierto gratuito para aprendizaje de
  estadística.
\end{itemize}

\section*{Diccionarios y Referencias
Rápidas}\label{diccionarios-y-referencias-ruxe1pidas}
\addcontentsline{toc}{section}{Diccionarios y Referencias Rápidas}

\markright{Diccionarios y Referencias Rápidas}

\begin{itemize}
\tightlist
\item
  \citet{Dodge2006} --- Diccionario definitivo de términos estadísticos
  en inglés.
\end{itemize}

\section*{Cómo citar este curso}\label{cuxf3mo-citar-este-curso}
\addcontentsline{toc}{section}{Cómo citar este curso}

\markright{Cómo citar este curso}

Si utiliza los materiales de este curso en su investigación o trabajo,
le recomendamos citar:

\begin{quote}
Luque Fernández, Miguel Ángel. (2025--2026). Matemáticas para la
Estadística Médica. Universidad de Granada. Disponible en:
https://migariane.github.io/CursoMatematicaEstadistica/
\end{quote}

\section*{Acceso Abierto}\label{acceso-abierto}
\addcontentsline{toc}{section}{Acceso Abierto}

\markright{Acceso Abierto}

Este curso es parte de una iniciativa de educación abierta. Los
materiales están disponibles bajo licencia \textbf{CC BY-NC-SA 4.0}, lo
que permite su uso para propósitos educativos y no comerciales.

\begin{center}\rule{0.5\linewidth}{0.5pt}\end{center}

\renewcommand{\bibsection}{}
\bibliography{references.bib}

\begin{center}\rule{0.5\linewidth}{0.5pt}\end{center}

\textbf{Última actualización:} Semestre 2025--2026

\cleardoublepage
\phantomsection
\addcontentsline{toc}{part}{Appendices}
\appendix

\chapter{Apéndice A: Repaso Matemático Riguroso}\label{sec-apendiceA}

Este apéndice proporciona los fundamentos matemáticos necesarios para un
análisis riguroso de la probabilidad y la estadística. Este contenido
está estructurado como una guía de referencia basada en principios de
bioestadística.

\section{A.1 Funciones y Gráficas}\label{sec-apendiceA-funciones}

Las funciones mapean elementos de un dominio a un codominio
(\(f: X \to Y\)).

\begin{itemize}
\tightlist
\item
  \textbf{Notación:} \(y = f(x)\), donde \(x\) es la variable
  independiente y \(y\) es la variable dependiente.
\item
  \textbf{Dominio:} Conjunto de todos los posibles valores de entrada
  (\(x\)).
\item
  \textbf{Codominio/Rango:} Conjunto de todos los posibles valores de
  salida (\(y\)).
\item
  \textbf{Funciones Inversas:} Si una función \(f\) es biyectiva
  (inyectiva y sobreyectiva), existe una función inversa \(f^{-1}\) tal
  que \(f(f^{-1}(x)) = x\) y \(f^{-1}(f(x)) = x\).
\item
  \textbf{Composición:} La composición de funciones \((f \circ g)(x)\)
  significa aplicar \(g\) primero, y luego \(f\) al resultado:
  \((f \circ g)(x) = f(g(x))\).
\end{itemize}

\textbf{Ejemplos de funciones comunes:} * \textbf{Lineal:}
\(f(x) = mx + b\) * \textbf{Cuadrática:} \(f(x) = ax^2 + bx + c\) *
\textbf{Exponencial:} \(f(x) = a^x\) * \textbf{Logarítmica:}
\(f(x) = \log_b(x)\)

\textbf{Ejemplo en R: Definición y uso de funciones}

\begin{codelisting}

\caption{\label{lst-apendiceA-1}Código R}

\centering{

\begin{Shaded}
\begin{Highlighting}[]
\CommentTok{\# Definir una función lineal}
\NormalTok{f\_lineal }\OtherTok{\textless{}{-}} \ControlFlowTok{function}\NormalTok{(x) \{}
  \DecValTok{2}\SpecialCharTok{*}\NormalTok{x }\SpecialCharTok{+} \DecValTok{3}
\NormalTok{\}}
\FunctionTok{f\_lineal}\NormalTok{(}\DecValTok{5}\NormalTok{) }\CommentTok{\# Resultado: 13}
\end{Highlighting}
\end{Shaded}

}

\end{codelisting}%

\begin{verbatim}
[1] 13
\end{verbatim}

\begin{codelisting}

\caption{\label{lst-apendiceA-1}Código R}

\centering{

\begin{Shaded}
\begin{Highlighting}[]
\CommentTok{\# Graficar una función}
\FunctionTok{curve}\NormalTok{(f\_lineal, }\AttributeTok{from =} \SpecialCharTok{{-}}\DecValTok{5}\NormalTok{, }\AttributeTok{to =} \DecValTok{5}\NormalTok{,}
      \AttributeTok{main =} \StringTok{"Función Lineal f(x) = 2x + 3"}\NormalTok{,}
      \AttributeTok{xlab =} \StringTok{"x"}\NormalTok{, }\AttributeTok{ylab =} \StringTok{"f(x)"}\NormalTok{)}
\end{Highlighting}
\end{Shaded}

}

\end{codelisting}%

\pandocbounded{\includegraphics[keepaspectratio]{chapters/apendiceA_files/figure-pdf/unnamed-chunk-1-1.pdf}}

\section{A.2 Exponenciales y Logaritmos}\label{sec-apendiceA-exp-log}

Fundamentales para modelizar relaciones de crecimiento/decrecimiento y
simplificar funciones de verosimilitud en estadística.

\begin{itemize}
\tightlist
\item
  \textbf{Propiedades Exponenciales:}

  \begin{itemize}
  \tightlist
  \item
    \(a^x a^y = a^{x+y}\)
  \item
    \((a^x)^y = a^{xy}\)
  \item
    \(a^0 = 1\)
  \item
    \(a^{-x} = 1/a^x\)
  \item
    \((ab)^x = a^x b^x\)
  \item
    En R, \texttt{exp(x)} calcula \(e^x\).
  \end{itemize}
\item
  \textbf{Propiedades Logarítmicas (base natural \(\ln\) o \(\log_e\)):}

  \begin{itemize}
  \tightlist
  \item
    \(\ln(xy) = \ln(x) + \ln(y)\)
  \item
    \(\ln(x/y) = \ln(x) - \ln(y)\)
  \item
    \(\ln(x^k) = k \ln(x)\)
  \item
    \(\ln(1) = 0\)
  \item
    \(\ln(e) = 1\)
  \item
    \(\log_b x = \frac{\ln x}{\ln b}\) (cambio de base)
  \item
    La relación inversa: \(e^{\ln x} = x\) y \(\ln(e^x) = x\).
  \item
    En R, \texttt{log(x)} calcula \(\ln(x)\) por defecto;
    \texttt{log10(x)} para base 10, \texttt{log(x,\ base=b)} para otras
    bases.
  \end{itemize}
\end{itemize}

\textbf{Ejemplo en R: Cálculos con exponenciales y logaritmos}

\begin{codelisting}

\caption{\label{lst-apendiceA-2}Código R}

\centering{

\begin{Shaded}
\begin{Highlighting}[]
\CommentTok{\# Cálculos exponenciales}
\FunctionTok{exp}\NormalTok{(}\DecValTok{1}\NormalTok{)      }\CommentTok{\# e\^{}1}
\end{Highlighting}
\end{Shaded}

}

\end{codelisting}%

\begin{verbatim}
[1] 2.718282
\end{verbatim}

\begin{codelisting}

\caption{\label{lst-apendiceA-2}Código R}

\centering{

\begin{Shaded}
\begin{Highlighting}[]
\FunctionTok{exp}\NormalTok{(}\DecValTok{2}\NormalTok{)      }\CommentTok{\# e\^{}2}
\end{Highlighting}
\end{Shaded}

}

\end{codelisting}%

\begin{verbatim}
[1] 7.389056
\end{verbatim}

\begin{codelisting}

\caption{\label{lst-apendiceA-2}Código R}

\centering{

\begin{Shaded}
\begin{Highlighting}[]
\CommentTok{\# Cálculos logarítmicos}
\FunctionTok{log}\NormalTok{(}\FunctionTok{exp}\NormalTok{(}\DecValTok{1}\NormalTok{)) }\CommentTok{\# ln(e) = 1}
\end{Highlighting}
\end{Shaded}

}

\end{codelisting}%

\begin{verbatim}
[1] 1
\end{verbatim}

\begin{codelisting}

\caption{\label{lst-apendiceA-2}Código R}

\centering{

\begin{Shaded}
\begin{Highlighting}[]
\FunctionTok{log}\NormalTok{(}\DecValTok{100}\NormalTok{)    }\CommentTok{\# ln(100)}
\end{Highlighting}
\end{Shaded}

}

\end{codelisting}%

\begin{verbatim}
[1] 4.60517
\end{verbatim}

\begin{codelisting}

\caption{\label{lst-apendiceA-2}Código R}

\centering{

\begin{Shaded}
\begin{Highlighting}[]
\FunctionTok{log}\NormalTok{(}\DecValTok{100}\NormalTok{, }\AttributeTok{base =} \DecValTok{10}\NormalTok{) }\CommentTok{\# log10(100) = 2}
\end{Highlighting}
\end{Shaded}

}

\end{codelisting}%

\begin{verbatim}
[1] 2
\end{verbatim}

\begin{codelisting}

\caption{\label{lst-apendiceA-2}Código R}

\centering{

\begin{Shaded}
\begin{Highlighting}[]
\CommentTok{\# Propiedad: log(x\^{}k) = k * log(x)}
\NormalTok{x\_val }\OtherTok{\textless{}{-}} \DecValTok{5}
\NormalTok{k\_val }\OtherTok{\textless{}{-}} \DecValTok{3}
\FunctionTok{log}\NormalTok{(x\_val}\SpecialCharTok{\^{}}\NormalTok{k\_val) }\SpecialCharTok{==}\NormalTok{ k\_val }\SpecialCharTok{*} \FunctionTok{log}\NormalTok{(x\_val)}
\end{Highlighting}
\end{Shaded}

}

\end{codelisting}%

\begin{verbatim}
[1] FALSE
\end{verbatim}

\section{A.3 Sumatorias y Productos}\label{sec-apendiceA-sum-prod}

La notación \(\sum\) representa sumas y \(\prod\) representa productos.
Son esenciales para definir medias, varianzas y funciones de
verosimilitud.

\begin{itemize}
\item
  \textbf{Notación de Sumatoria:}
  \(\sum_{i=1}^n x_i = x_1 + x_2 + \cdots + x_n\)
\item
  \textbf{Propiedades de la Suma:}

  \begin{itemize}
  \tightlist
  \item
    Linealidad:
    \(\sum_{i=1}^n (aX_i + bY_i) = a\sum_{i=1}^n X_i + b\sum_{i=1}^n Y_i\)
  \item
    Suma de una constante: \(\sum_{i=1}^n c = nc\)
  \end{itemize}
\item
  \textbf{Sumas dobles:} \(\sum_{i=1}^m \sum_{j=1}^n x_{ij}\) para sumar
  elementos en una tabla o matriz.
\item
  \textbf{Notación de Producto:}
  \(\prod_{i=1}^n x_i = x_1 \cdot x_2 \cdot \ldots \cdot x_n\)
\end{itemize}

\textbf{Ejemplo en R: Sumas y Productos}

\begin{codelisting}

\caption{\label{lst-apendiceA-3}Código R}

\centering{

\begin{Shaded}
\begin{Highlighting}[]
\CommentTok{\# Vector de datos}
\NormalTok{datos }\OtherTok{\textless{}{-}} \FunctionTok{c}\NormalTok{(}\DecValTok{2}\NormalTok{, }\DecValTok{4}\NormalTok{, }\DecValTok{6}\NormalTok{, }\DecValTok{8}\NormalTok{, }\DecValTok{10}\NormalTok{)}

\CommentTok{\# Suma de los datos}
\FunctionTok{sum}\NormalTok{(datos)}
\end{Highlighting}
\end{Shaded}

}

\end{codelisting}%

\begin{verbatim}
[1] 30
\end{verbatim}

\begin{codelisting}

\caption{\label{lst-apendiceA-3}Código R}

\centering{

\begin{Shaded}
\begin{Highlighting}[]
\CommentTok{\# Media muestral usando sumatoria}
\NormalTok{n\_datos }\OtherTok{\textless{}{-}} \FunctionTok{length}\NormalTok{(datos)}
\NormalTok{media }\OtherTok{\textless{}{-}} \FunctionTok{sum}\NormalTok{(datos) }\SpecialCharTok{/}\NormalTok{ n\_datos}
\FunctionTok{cat}\NormalTok{(}\StringTok{"Media muestral:"}\NormalTok{, media, }\StringTok{"}\SpecialCharTok{\textbackslash{}n}\StringTok{"}\NormalTok{)}
\end{Highlighting}
\end{Shaded}

}

\end{codelisting}%

\begin{verbatim}
Media muestral: 6 
\end{verbatim}

\begin{codelisting}

\caption{\label{lst-apendiceA-3}Código R}

\centering{

\begin{Shaded}
\begin{Highlighting}[]
\CommentTok{\# Producto de los datos}
\FunctionTok{prod}\NormalTok{(datos)}
\end{Highlighting}
\end{Shaded}

}

\end{codelisting}%

\begin{verbatim}
[1] 3840
\end{verbatim}

\begin{codelisting}

\caption{\label{lst-apendiceA-3}Código R}

\centering{

\begin{Shaded}
\begin{Highlighting}[]
\CommentTok{\# Suma doble (ej. suma de una matriz)}
\NormalTok{matriz\_ej }\OtherTok{\textless{}{-}} \FunctionTok{matrix}\NormalTok{(}\DecValTok{1}\SpecialCharTok{:}\DecValTok{4}\NormalTok{, }\AttributeTok{nrow=}\DecValTok{2}\NormalTok{)}
\FunctionTok{print}\NormalTok{(matriz\_ej)}
\end{Highlighting}
\end{Shaded}

}

\end{codelisting}%

\begin{verbatim}
     [,1] [,2]
[1,]    1    3
[2,]    2    4
\end{verbatim}

\begin{codelisting}

\caption{\label{lst-apendiceA-3}Código R}

\centering{

\begin{Shaded}
\begin{Highlighting}[]
\FunctionTok{sum}\NormalTok{(matriz\_ej) }\CommentTok{\# Suma todos los elementos}
\end{Highlighting}
\end{Shaded}

}

\end{codelisting}%

\begin{verbatim}
[1] 10
\end{verbatim}

\section{A.4 Teoría de Conjuntos}\label{sec-apendiceA-conjuntos}

La teoría de conjuntos es el lenguaje fundamental para describir eventos
en probabilidad.

\begin{itemize}
\tightlist
\item
  \textbf{Conjunto:} Colección de objetos distintos.

  \begin{itemize}
  \tightlist
  \item
    \(A = \{1, 2, 3\}\)
  \item
    \(\Omega\) = Espacio muestral (conjunto de todos los resultados
    posibles)
  \end{itemize}
\item
  \textbf{Elemento:} Un objeto dentro de un conjunto.

  \begin{itemize}
  \tightlist
  \item
    \(x \in A\) (x es un elemento de A)
  \item
    \(y \notin A\) (y no es un elemento de A)
  \end{itemize}
\item
  \textbf{Subconjunto:} \(A \subset B\) (todos los elementos de A
  también están en B).
\item
  \textbf{Conjunto Vacío:} \(\emptyset\) o \texttt{\{\}} (contiene cero
  elementos).
\item
  \textbf{Operaciones:}

  \begin{itemize}
  \tightlist
  \item
    \textbf{Unión (\(\cup\)):}
    \(A \cup B = \{x | x \in A \text{ o } x \in B\}\). Elementos en
    \(A\) O \(B\) (o ambos).
  \item
    \textbf{Intersección (\(\cap\)):}
    \(A \cap B = \{x | x \in A \text{ y } x \in B\}\). Elementos en
    \(A\) Y \(B\).
  \item
    \textbf{Complemento (\(A^c\) o \(A'\)):} Elementos en el espacio
    muestral \(\Omega\) que no están en \(A\). Formalmente,
    \(A^c = \{x | x \in \Omega \text{ y } x \notin A\}\).
  \item
    \textbf{Diferencia (\(A \setminus B\)):} Elementos en \(A\) pero no
    en \(B\). Equivalente a \(A \cap B^c\).
  \end{itemize}
\item
  \textbf{Conjuntos Disjuntos:} \(A\) y \(B\) son disjuntos si
  \(A \cap B = \emptyset\). No tienen elementos en común.
\item
  \textbf{Partición:} Una colección de conjuntos
  \(A_1, A_2, \ldots, A_k\) forma una partición de \(\Omega\) si son
  mutuamente disjuntos y su unión es \(\Omega\).
\end{itemize}

\textbf{Ejemplo en R: Operaciones con conjuntos (usando vectores
lógicos)}

\begin{codelisting}

\caption{\label{lst-apendiceA-4}Código R}

\centering{

\begin{Shaded}
\begin{Highlighting}[]
\CommentTok{\# Representación de conjuntos como vectores lógicos o numéricos}
\NormalTok{universo }\OtherTok{\textless{}{-}} \DecValTok{1}\SpecialCharTok{:}\DecValTok{10}
\NormalTok{conjuntoA }\OtherTok{\textless{}{-}} \FunctionTok{c}\NormalTok{(}\DecValTok{1}\NormalTok{, }\DecValTok{2}\NormalTok{, }\DecValTok{3}\NormalTok{, }\DecValTok{4}\NormalTok{, }\DecValTok{5}\NormalTok{)}
\NormalTok{conjuntoB }\OtherTok{\textless{}{-}} \FunctionTok{c}\NormalTok{(}\DecValTok{4}\NormalTok{, }\DecValTok{5}\NormalTok{, }\DecValTok{6}\NormalTok{, }\DecValTok{7}\NormalTok{)}

\CommentTok{\# Unión}
\FunctionTok{union}\NormalTok{(conjuntoA, conjuntoB)}
\end{Highlighting}
\end{Shaded}

}

\end{codelisting}%

\begin{verbatim}
[1] 1 2 3 4 5 6 7
\end{verbatim}

\begin{codelisting}

\caption{\label{lst-apendiceA-4}Código R}

\centering{

\begin{Shaded}
\begin{Highlighting}[]
\CommentTok{\# Intersección}
\FunctionTok{intersect}\NormalTok{(conjuntoA, conjuntoB)}
\end{Highlighting}
\end{Shaded}

}

\end{codelisting}%

\begin{verbatim}
[1] 4 5
\end{verbatim}

\begin{codelisting}

\caption{\label{lst-apendiceA-4}Código R}

\centering{

\begin{Shaded}
\begin{Highlighting}[]
\CommentTok{\# Diferencia (elementos en A que no están en B)}
\FunctionTok{setdiff}\NormalTok{(conjuntoA, conjuntoB)}
\end{Highlighting}
\end{Shaded}

}

\end{codelisting}%

\begin{verbatim}
[1] 1 2 3
\end{verbatim}

\begin{codelisting}

\caption{\label{lst-apendiceA-4}Código R}

\centering{

\begin{Shaded}
\begin{Highlighting}[]
\CommentTok{\# Complemento (respecto a \textquotesingle{}universo\textquotesingle{})}
\FunctionTok{setdiff}\NormalTok{(universo, conjuntoA)}
\end{Highlighting}
\end{Shaded}

}

\end{codelisting}%

\begin{verbatim}
[1]  6  7  8  9 10
\end{verbatim}

\begin{codelisting}

\caption{\label{lst-apendiceA-4}Código R}

\centering{

\begin{Shaded}
\begin{Highlighting}[]
\CommentTok{\# Comprobación de disjunción (no disjuntos en este caso)}
\FunctionTok{length}\NormalTok{(}\FunctionTok{intersect}\NormalTok{(conjuntoA, conjuntoB)) }\SpecialCharTok{==} \DecValTok{0}
\end{Highlighting}
\end{Shaded}

}

\end{codelisting}%

\begin{verbatim}
[1] FALSE
\end{verbatim}

\section{A.5 Cálculo: Derivadas e
Integrales}\label{sec-apendiceA-calculo}

Herramientas fundamentales para la optimización (minimización de
funciones de pérdida, maximización de verosimilitud) usando derivadas y
para el cálculo de probabilidades e inferencia usando integrales.

\subsection{A.5.1 Derivadas}\label{a.5.1-derivadas}

Miden la tasa de cambio instantánea de una función. En estadística, se
usan para encontrar máximos y mínimos de funciones (como funciones de
verosimilitud o de costo).

\begin{itemize}
\tightlist
\item
  \textbf{Notación:} \(f'(x)\), \(\frac{dy}{dx}\), \(\frac{d}{dx}f(x)\).
\item
  \textbf{Reglas Básicas:}

  \begin{itemize}
  \tightlist
  \item
    Constante: \(\frac{d}{dx}c = 0\)
  \item
    Potencia: \(\frac{d}{dx}x^n = nx^{n-1}\)
  \item
    Suma/Resta: \(\frac{d}{dx}(f(x) \pm g(x)) = f'(x) \pm g'(x)\)
  \item
    Producto: \(\frac{d}{dx}(f(x)g(x)) = f'(x)g(x) + f(x)g'(x)\)
  \item
    Cociente:
    \(\frac{d}{dx}\left(\frac{f(x)}{g(x)}\right) = \frac{f'(x)g(x) - f(x)g'(x)}{[g(x)]^2}\)
  \item
    Regla de la cadena: \(\frac{d}{dx}f(g(x)) = f'(g(x))g'(x)\)
  \item
    Exponencial: \(\frac{d}{dx}e^x = e^x\)
  \item
    Logaritmo: \(\frac{d}{dx}\ln(x) = \frac{1}{x}\)
  \end{itemize}
\end{itemize}

\textbf{Ejemplo: Minimización de una función (encontrar el vértice de
una parábola)} Sea \(f(x) = x^2 - 4x + 5\). Para encontrar el mínimo,
derivamos e igualamos a cero: \(f'(x) = 2x - 4\)
\(2x - 4 = 0 \implies x = 2\). Sustituyendo en \(f(x)\):
\(f(2) = 2^2 - 4(2) + 5 = 4 - 8 + 5 = 1\). El mínimo es en \((2, 1)\).

\textbf{Ejemplo en R: Derivación simbólica (para funciones simples)}

\begin{codelisting}

\caption{\label{lst-apendiceA-5}Código R}

\centering{

\begin{Shaded}
\begin{Highlighting}[]
\CommentTok{\# Definir una expresión simbólica}
\NormalTok{expr }\OtherTok{\textless{}{-}} \FunctionTok{expression}\NormalTok{(x}\SpecialCharTok{\^{}}\DecValTok{2} \SpecialCharTok{{-}} \DecValTok{4}\SpecialCharTok{*}\NormalTok{x }\SpecialCharTok{+} \DecValTok{5}\NormalTok{)}
\NormalTok{derivada }\OtherTok{\textless{}{-}} \FunctionTok{D}\NormalTok{(expr, }\StringTok{"x"}\NormalTok{)}
\FunctionTok{print}\NormalTok{(derivada)}
\end{Highlighting}
\end{Shaded}

}

\end{codelisting}%

\begin{verbatim}
2 * x - 4
\end{verbatim}

\begin{codelisting}

\caption{\label{lst-apendiceA-5}Código R}

\centering{

\begin{Shaded}
\begin{Highlighting}[]
\CommentTok{\# Evaluar la derivada en un punto}
\FunctionTok{eval}\NormalTok{(derivada, }\FunctionTok{list}\NormalTok{(}\AttributeTok{x =} \DecValTok{2}\NormalTok{)) }\CommentTok{\# Debería ser 0}
\end{Highlighting}
\end{Shaded}

}

\end{codelisting}%

\begin{verbatim}
[1] 0
\end{verbatim}

\subsection{A.5.2 Integrales}\label{a.5.2-integrales}

Representan el área bajo una curva y se utilizan para calcular
probabilidades a partir de funciones de densidad, esperanzas y
varianzas.

\begin{itemize}
\tightlist
\item
  \textbf{Integral Indefinida (Antiderivada):}
  \(\int f(x) dx = F(x) + C\), donde \(F'(x) = f(x)\).
\item
  \textbf{Integral Definida:} \(\int_a^b f(x) dx = F(b) - F(a)\). Esto
  calcula el área neta entre \(f(x)\) y el eje x desde \(a\) hasta
  \(b\).
\item
  \textbf{Integrales Múltiples:} Para funciones de densidad conjuntas.
\end{itemize}

\textbf{Ejemplo: Cálculo de una probabilidad a partir de una PDF} Si
\(f(x)\) es la PDF, \(P(a \le X \le b) = \int_a^b f(x) dx\).

\textbf{Ejemplo en R: Integración numérica}

\begin{codelisting}

\caption{\label{lst-apendiceA-6}Código R}

\centering{

\begin{Shaded}
\begin{Highlighting}[]
\CommentTok{\# Calcular la probabilidad P({-}1 \textless{}= X \textless{}= 1) para una Normal Estándar}
\CommentTok{\# Usamos la función dnorm (densidad de la normal)}
\NormalTok{resultado\_integral }\OtherTok{\textless{}{-}} \FunctionTok{integrate}\NormalTok{(dnorm, }\AttributeTok{lower =} \SpecialCharTok{{-}}\DecValTok{1}\NormalTok{, }\AttributeTok{upper =} \DecValTok{1}\NormalTok{)}
\FunctionTok{print}\NormalTok{(resultado\_integral}\SpecialCharTok{$}\NormalTok{value) }\CommentTok{\# Aproximadamente 0.6827 (regla del 68\%)}
\end{Highlighting}
\end{Shaded}

}

\end{codelisting}%

\begin{verbatim}
[1] 0.6826895
\end{verbatim}

\begin{codelisting}

\caption{\label{lst-apendiceA-6}Código R}

\centering{

\begin{Shaded}
\begin{Highlighting}[]
\CommentTok{\# Para la esperanza de una variable continua con PDF f(x): E(X) = integral(x * f(x) dx)}
\CommentTok{\# Ejemplo con una función simple x * exp({-}x) de 0 a Inf}
\NormalTok{f\_ejemplo\_esperanza }\OtherTok{\textless{}{-}} \ControlFlowTok{function}\NormalTok{(x) \{ x }\SpecialCharTok{*} \FunctionTok{dexp}\NormalTok{(x, }\AttributeTok{rate =} \DecValTok{1}\NormalTok{) \} }\CommentTok{\# E(X) para Exp(1) es 1/rate = 1}
\NormalTok{esperanza\_numerica }\OtherTok{\textless{}{-}} \FunctionTok{integrate}\NormalTok{(f\_ejemplo\_esperanza, }\AttributeTok{lower =} \DecValTok{0}\NormalTok{, }\AttributeTok{upper =} \ConstantTok{Inf}\NormalTok{)}
\FunctionTok{print}\NormalTok{(esperanza\_numerica}\SpecialCharTok{$}\NormalTok{value) }\CommentTok{\# Debería ser cercano a 1}
\end{Highlighting}
\end{Shaded}

}

\end{codelisting}%

\begin{verbatim}
[1] 1
\end{verbatim}

\section{A.6 Álgebra Matricial}\label{sec-apendiceA-matrices}

El álgebra matricial es indispensable en estadística multivariante,
regresión lineal (donde los coeficientes se estiman usando inversas
matriciales), análisis de componentes principales, y más.

\begin{itemize}
\tightlist
\item
  \textbf{Matriz:} Un arreglo rectangular de números.

  \begin{itemize}
  \tightlist
  \item
    Notación: \(\mathbf{A}_{m \times n}\) (m filas, n columnas).
  \end{itemize}
\item
  \textbf{Vector:} Una matriz con una sola fila (vector fila) o una sola
  columna (vector columna).
\item
  \textbf{Tipos de Matrices:}

  \begin{itemize}
  \tightlist
  \item
    \textbf{Identidad (\(\mathbf{I}\)):} Matriz cuadrada con 1s en la
    diagonal y 0s en el resto.
    \(\mathbf{AI} = \mathbf{IA} = \mathbf{A}\).
  \item
    \textbf{Transpuesta (\(\mathbf{A}^T\)):} Se obtienen intercambiando
    filas por columnas. \((A^T)_{ij} = A_{ji}\).
  \item
    \textbf{Simétrica:} \(\mathbf{A} = \mathbf{A}^T\).
  \end{itemize}
\item
  \textbf{Operaciones Básicas:}

  \begin{itemize}
  \tightlist
  \item
    \textbf{Suma/Resta:} Elemento a elemento (las matrices deben tener
    las mismas dimensiones).
  \item
    \textbf{Multiplicación por un escalar:} Multiplicar cada elemento
    por el escalar.
  \item
    \textbf{Producto Matricial (\(\mathbf{AB}\)):} El número de columnas
    de \(\mathbf{A}\) debe ser igual al número de filas de
    \(\mathbf{B}\).

    \begin{itemize}
    \tightlist
    \item
      Si \(\mathbf{A}_{m \times p}\) y \(\mathbf{B}_{p \times n}\),
      entonces \((\mathbf{AB})_{m \times n}\).
    \item
      \((\mathbf{AB})_{ij} = \sum_{k=1}^p A_{ik}B_{kj}\).
    \item
      NO es conmutativo: \(\mathbf{AB} \ne \mathbf{BA}\) en general.
    \end{itemize}
  \end{itemize}
\item
  \textbf{Inversa de una Matriz (\(\mathbf{A}^{-1}\)):} Para una matriz
  cuadrada \(\mathbf{A}\), si existe \(\mathbf{A}^{-1}\), entonces
  \(\mathbf{AA}^{-1} = \mathbf{A}^{-1}\mathbf{A} = \mathbf{I}\). Solo
  existe si el determinante es distinto de cero.
\item
  \textbf{Determinante (\(\det(\mathbf{A})\) o \(|\mathbf{A}|\)):} Un
  escalar asociado a una matriz cuadrada que indica, entre otras cosas,
  si la matriz es invertible (si \(\det(\mathbf{A}) \ne 0\)).
\end{itemize}

\textbf{Ejemplo en R: Operaciones con matrices}

\begin{codelisting}

\caption{\label{lst-apendiceA-7}Código R}

\centering{

\begin{Shaded}
\begin{Highlighting}[]
\CommentTok{\# Definir matrices}
\NormalTok{A }\OtherTok{\textless{}{-}} \FunctionTok{matrix}\NormalTok{(}\FunctionTok{c}\NormalTok{(}\DecValTok{1}\NormalTok{, }\DecValTok{2}\NormalTok{, }\DecValTok{3}\NormalTok{, }\DecValTok{4}\NormalTok{), }\AttributeTok{nrow=}\DecValTok{2}\NormalTok{, }\AttributeTok{byrow=}\ConstantTok{TRUE}\NormalTok{) }\CommentTok{\# Matriz 2x2}
\NormalTok{B }\OtherTok{\textless{}{-}} \FunctionTok{matrix}\NormalTok{(}\FunctionTok{c}\NormalTok{(}\DecValTok{5}\NormalTok{, }\DecValTok{6}\NormalTok{, }\DecValTok{7}\NormalTok{, }\DecValTok{8}\NormalTok{), }\AttributeTok{nrow=}\DecValTok{2}\NormalTok{, }\AttributeTok{byrow=}\ConstantTok{TRUE}\NormalTok{) }\CommentTok{\# Matriz 2x2}
\NormalTok{C }\OtherTok{\textless{}{-}} \FunctionTok{matrix}\NormalTok{(}\FunctionTok{c}\NormalTok{(}\DecValTok{1}\NormalTok{, }\DecValTok{0}\NormalTok{, }\DecValTok{0}\NormalTok{, }\DecValTok{1}\NormalTok{), }\AttributeTok{nrow=}\DecValTok{2}\NormalTok{, }\AttributeTok{byrow=}\ConstantTok{TRUE}\NormalTok{) }\CommentTok{\# Matriz identidad I}

\FunctionTok{print}\NormalTok{(}\StringTok{"Matriz A:"}\NormalTok{)}
\end{Highlighting}
\end{Shaded}

}

\end{codelisting}%

\begin{verbatim}
[1] "Matriz A:"
\end{verbatim}

\begin{codelisting}

\caption{\label{lst-apendiceA-7}Código R}

\centering{

\begin{Shaded}
\begin{Highlighting}[]
\FunctionTok{print}\NormalTok{(A)}
\end{Highlighting}
\end{Shaded}

}

\end{codelisting}%

\begin{verbatim}
     [,1] [,2]
[1,]    1    2
[2,]    3    4
\end{verbatim}

\begin{codelisting}

\caption{\label{lst-apendiceA-7}Código R}

\centering{

\begin{Shaded}
\begin{Highlighting}[]
\FunctionTok{print}\NormalTok{(}\StringTok{"Matriz B:"}\NormalTok{)}
\end{Highlighting}
\end{Shaded}

}

\end{codelisting}%

\begin{verbatim}
[1] "Matriz B:"
\end{verbatim}

\begin{codelisting}

\caption{\label{lst-apendiceA-7}Código R}

\centering{

\begin{Shaded}
\begin{Highlighting}[]
\FunctionTok{print}\NormalTok{(B)}
\end{Highlighting}
\end{Shaded}

}

\end{codelisting}%

\begin{verbatim}
     [,1] [,2]
[1,]    5    6
[2,]    7    8
\end{verbatim}

\begin{codelisting}

\caption{\label{lst-apendiceA-7}Código R}

\centering{

\begin{Shaded}
\begin{Highlighting}[]
\CommentTok{\# Suma de matrices}
\FunctionTok{print}\NormalTok{(}\StringTok{"A + B:"}\NormalTok{)}
\end{Highlighting}
\end{Shaded}

}

\end{codelisting}%

\begin{verbatim}
[1] "A + B:"
\end{verbatim}

\begin{codelisting}

\caption{\label{lst-apendiceA-7}Código R}

\centering{

\begin{Shaded}
\begin{Highlighting}[]
\FunctionTok{print}\NormalTok{(A }\SpecialCharTok{+}\NormalTok{ B)}
\end{Highlighting}
\end{Shaded}

}

\end{codelisting}%

\begin{verbatim}
     [,1] [,2]
[1,]    6    8
[2,]   10   12
\end{verbatim}

\begin{codelisting}

\caption{\label{lst-apendiceA-7}Código R}

\centering{

\begin{Shaded}
\begin{Highlighting}[]
\CommentTok{\# Producto matricial}
\FunctionTok{print}\NormalTok{(}\StringTok{"A \%*\% B:"}\NormalTok{)}
\end{Highlighting}
\end{Shaded}

}

\end{codelisting}%

\begin{verbatim}
[1] "A %*% B:"
\end{verbatim}

\begin{codelisting}

\caption{\label{lst-apendiceA-7}Código R}

\centering{

\begin{Shaded}
\begin{Highlighting}[]
\FunctionTok{print}\NormalTok{(A }\SpecialCharTok{\%*\%}\NormalTok{ B)}
\end{Highlighting}
\end{Shaded}

}

\end{codelisting}%

\begin{verbatim}
     [,1] [,2]
[1,]   19   22
[2,]   43   50
\end{verbatim}

\begin{codelisting}

\caption{\label{lst-apendiceA-7}Código R}

\centering{

\begin{Shaded}
\begin{Highlighting}[]
\CommentTok{\# Transpuesta de A}
\FunctionTok{print}\NormalTok{(}\StringTok{"Transpuesta de A (A\^{}T):"}\NormalTok{)}
\end{Highlighting}
\end{Shaded}

}

\end{codelisting}%

\begin{verbatim}
[1] "Transpuesta de A (A^T):"
\end{verbatim}

\begin{codelisting}

\caption{\label{lst-apendiceA-7}Código R}

\centering{

\begin{Shaded}
\begin{Highlighting}[]
\FunctionTok{print}\NormalTok{(}\FunctionTok{t}\NormalTok{(A))}
\end{Highlighting}
\end{Shaded}

}

\end{codelisting}%

\begin{verbatim}
     [,1] [,2]
[1,]    1    3
[2,]    2    4
\end{verbatim}

\begin{codelisting}

\caption{\label{lst-apendiceA-7}Código R}

\centering{

\begin{Shaded}
\begin{Highlighting}[]
\CommentTok{\# Inversa de A}
\CommentTok{\# Solo si la matriz es cuadrada y no singular (determinante != 0)}
\FunctionTok{print}\NormalTok{(}\StringTok{"Inversa de A:"}\NormalTok{)}
\end{Highlighting}
\end{Shaded}

}

\end{codelisting}%

\begin{verbatim}
[1] "Inversa de A:"
\end{verbatim}

\begin{codelisting}

\caption{\label{lst-apendiceA-7}Código R}

\centering{

\begin{Shaded}
\begin{Highlighting}[]
\FunctionTok{print}\NormalTok{(}\FunctionTok{solve}\NormalTok{(A))}
\end{Highlighting}
\end{Shaded}

}

\end{codelisting}%

\begin{verbatim}
     [,1] [,2]
[1,] -2.0  1.0
[2,]  1.5 -0.5
\end{verbatim}

\begin{codelisting}

\caption{\label{lst-apendiceA-7}Código R}

\centering{

\begin{Shaded}
\begin{Highlighting}[]
\CommentTok{\# Verificar A * A\_inv = I}
\FunctionTok{print}\NormalTok{(}\StringTok{"A \%*\% solve(A):"}\NormalTok{)}
\end{Highlighting}
\end{Shaded}

}

\end{codelisting}%

\begin{verbatim}
[1] "A %*% solve(A):"
\end{verbatim}

\begin{codelisting}

\caption{\label{lst-apendiceA-7}Código R}

\centering{

\begin{Shaded}
\begin{Highlighting}[]
\FunctionTok{print}\NormalTok{(A }\SpecialCharTok{\%*\%} \FunctionTok{solve}\NormalTok{(A)) }\CommentTok{\# Debería ser la matriz identidad}
\end{Highlighting}
\end{Shaded}

}

\end{codelisting}%

\begin{verbatim}
     [,1]         [,2]
[1,]    1 1.110223e-16
[2,]    0 1.000000e+00
\end{verbatim}

\begin{codelisting}

\caption{\label{lst-apendiceA-7}Código R}

\centering{

\begin{Shaded}
\begin{Highlighting}[]
\CommentTok{\# Determinante de A}
\FunctionTok{print}\NormalTok{(}\StringTok{"Determinante de A:"}\NormalTok{)}
\end{Highlighting}
\end{Shaded}

}

\end{codelisting}%

\begin{verbatim}
[1] "Determinante de A:"
\end{verbatim}

\begin{codelisting}

\caption{\label{lst-apendiceA-7}Código R}

\centering{

\begin{Shaded}
\begin{Highlighting}[]
\FunctionTok{print}\NormalTok{(}\FunctionTok{det}\NormalTok{(A)) }\CommentTok{\# (1*4 {-} 2*3) = 4 {-} 6 = {-}2}
\end{Highlighting}
\end{Shaded}

}

\end{codelisting}%

\begin{verbatim}
[1] -2
\end{verbatim}

\section{A.7 Referencias}\label{sec-apendiceA-referencias}

\begin{enumerate}
\def\labelenumi{\arabic{enumi}.}
\tightlist
\item
  \textbf{Iowa Biostatistics Math Review:}
  \url{https://iowabiostat.github.io/math-review/}
\item
  \textbf{Casella, G., \& Berger, R. L. (2002).} \emph{Statistical
  Inference}. (Excelente para fundamentos matemáticos de la inferencia
  estadística).
\item
  \textbf{Spivak, M. (2006).} \emph{Calculus}. (Un texto clásico para un
  tratamiento riguroso del cálculo).
\item
  \textbf{Harville, D. A. (1997).} \emph{Matrix Algebra From a
  Statistician's Perspective}. (Referencia avanzada para álgebra
  matricial en estadística).
\end{enumerate}

\chapter{\texorpdfstring{Apéndice B: El Paquete \texttt{BioStatR} ---
Guía de
Referencia}{Apéndice B: El Paquete BioStatR --- Guía de Referencia}}\label{sec-apendiceB}

El paquete \textbf{\texttt{BioStatR}} (versión 1.0.0, 2026) ha sido
desarrollado en la Unidad Docente de Bioestadística del Departamento de
Estadística e Investigación Operativa de la Universidad de Granada por
\textbf{Pedro Femia Marzo} y \textbf{Miguel Ángel Luque Fernández}.
Proporciona 22 funciones de alto nivel que integran descripción,
inferencia, modelización (lineal y logística) y visualización en una
sola llamada, diseñadas específicamente para la docencia de
bioestadística en ciencias de la salud.

\begin{itemize}
\tightlist
\item
  \textbf{Sitio web del paquete:}
  \url{https://migariane.github.io/BioStatR/}
\item
  \textbf{Repositorio GitHub:}
  \url{https://github.com/migariane/BioStatR}
\item
  \textbf{Viñeta introductoria:}
  \url{https://migariane.github.io/BioStatR/articles/BioStatR-overview.html}
\item
  \textbf{Contacto:} \texttt{pfemia@ugr.es}
\end{itemize}

\begin{center}\rule{0.5\linewidth}{0.5pt}\end{center}

\section{\texorpdfstring{B.1 Instalación y el conjunto de datos
\texttt{osteo}}{B.1 Instalación y el conjunto de datos osteo}}\label{sec-bio-instalacion}

\begin{tcolorbox}[enhanced jigsaw, breakable, titlerule=0mm, rightrule=.15mm, leftrule=.75mm, left=2mm, opacitybacktitle=0.6, colback=white, bottomtitle=1mm, colbacktitle=quarto-callout-note-color!10!white, title=\textcolor{quarto-callout-note-color}{\faInfo}\hspace{0.5em}{Instalación}, toptitle=1mm, coltitle=black, opacityback=0, arc=.35mm, bottomrule=.15mm, toprule=.15mm, colframe=quarto-callout-note-color-frame]

\begin{codelisting}[H]

\caption{\label{lst-apendiceB-install}Instalación de BioStatR desde
GitHub}

\centering{

\begin{Shaded}
\begin{Highlighting}[]
\CommentTok{\# 1. Instalar el paquete remotes si no lo tienes}
\FunctionTok{install.packages}\NormalTok{(}\StringTok{"remotes"}\NormalTok{)}

\CommentTok{\# 2. Instalar BioStatR directamente desde GitHub (compila para tu sistema)}
\NormalTok{remotes}\SpecialCharTok{::}\FunctionTok{install\_github}\NormalTok{(}\StringTok{"migariane/BioStatR"}\NormalTok{)}

\CommentTok{\# 3. Cargar el paquete y el conjunto de datos osteo}
\FunctionTok{library}\NormalTok{(BioStatR)}
\FunctionTok{data}\NormalTok{(osteo)}
\end{Highlighting}
\end{Shaded}

}

\end{codelisting}%

Para obtener la referencia bibliográfica del paquete:

\begin{Shaded}
\begin{Highlighting}[]
\FunctionTok{citation}\NormalTok{(}\StringTok{"BioStatR"}\NormalTok{)}
\end{Highlighting}
\end{Shaded}

\end{tcolorbox}

El paquete incluye el dataset \texttt{osteo}, con datos reales de
\textbf{94 pacientes diabéticos} evaluados en la Unidad de Diabetes de
la Facultad de Medicina de la Universidad de Granada. Es el conjunto de
referencia de todos los ejemplos de este apéndice y del libro.

\begin{tcolorbox}[enhanced jigsaw, breakable, titlerule=0mm, rightrule=.15mm, leftrule=.75mm, left=2mm, opacitybacktitle=0.6, colback=white, bottomtitle=1mm, colbacktitle=quarto-callout-tip-color!10!white, title=\textcolor{quarto-callout-tip-color}{\faLightbulb}\hspace{0.5em}{El dataset \texttt{osteo}}, toptitle=1mm, coltitle=black, opacityback=0, arc=.35mm, bottomrule=.15mm, toprule=.15mm, colframe=quarto-callout-tip-color-frame]

\begin{codelisting}[H]

\caption{\label{lst-apendiceB-str}Estructura del dataset osteo}

\centering{

\begin{Shaded}
\begin{Highlighting}[]
\FunctionTok{library}\NormalTok{(BioStatR)}
\FunctionTok{data}\NormalTok{(osteo)}
\FunctionTok{str}\NormalTok{(osteo)}
\end{Highlighting}
\end{Shaded}

}

\end{codelisting}%

\begin{verbatim}
'data.frame':   94 obs. of  27 variables:
 $ num       : num  1 2 3 4 5 6 7 8 9 10 ...
 $ edad      : num  23 35 51 44 54 36 21 20 27 27 ...
 $ grupo_edad: Factor w/ 3 levels "< 25","> 33",..: 1 2 2 2 2 2 1 1 3 3 ...
 $ sexo      : Factor w/ 2 levels "Hombre","Mujer": 1 1 1 1 2 2 1 2 2 2 ...
 $ peso      : num  72.3 82.5 61 67 51 61.5 73.8 62 46 49 ...
 $ talla     : num  175 173 170 167 168 158 180 170 148 144 ...
 $ imc       : num  23.6 27.6 21.1 24 18.1 ...
 $ tevol     : num  2 4 17 17 30 2 3 12 16 19 ...
 $ tabaco    : Factor w/ 2 levels "No","Sí": 2 1 2 2 2 2 2 2 1 1 ...
 $ alcohol   : Factor w/ 3 levels "Excesivo","Moderado",..: 2 2 2 2 2 2 2 2 2 2 ...
 $ ingca     : Factor w/ 2 levels "Insuficiente",..: 2 2 2 1 2 2 2 2 2 1 ...
 $ acfis     : Factor w/ 2 levels "No","Sí": 2 1 2 2 2 2 1 2 2 1 ...
 $ retin     : Factor w/ 3 levels "Grave","Leve",..: 3 3 3 2 3 3 3 3 2 2 ...
 $ nefro     : Factor w/ 3 levels "Grave","Leve",..: 3 3 3 3 3 3 3 2 2 3 ...
 $ neuro     : Factor w/ 3 levels "Grave","Leve",..: 3 3 3 1 3 3 3 3 3 3 ...
 $ hba1c     : num  4.6 6.1 6.7 8.2 9.7 8.4 10.1 10.3 5.4 8.7 ...
 $ ca        : num  9.6 9.5 9.9 8.3 9.5 9.9 9.8 9.2 9.6 8.9 ...
 $ p         : num  4.2 3.4 4.7 3.2 3.7 4.7 4.7 4.2 2.6 3.9 ...
 $ cr        : num  1.02 1.07 1.35 0.99 0.86 0.8 1 0.7 0.92 1.08 ...
 $ pthm      : num  2.9 3.1 NA NA 3.1 1.1 3.4 NA 0.7 2.6 ...
 $ pthi      : num  49 NA 54.3 NA 49.8 59.9 20.9 NA NA 10 ...
 $ bmdcue    : num  22 NA 23.7 NA 35.7 38.3 6.2 NA NA 10.3 ...
 $ szl24     : num  1.07 0.83 0.83 0.72 0.56 0.92 1 0.83 0.72 0.69 ...
 $ sztri     : num  0.55 -1.62 -1.28 -1.11 -1.68 2.05 -0.75 -1.11 -0.98 -0.91 ...
 $ szcue     : num  1.44 -0.93 -0.07 -2.07 -3.1 0.35 0.08 -1.62 -2.25 -1.41 ...
 $ osteo_cue : Factor w/ 2 levels "No","Sí": 1 1 1 2 2 1 1 1 2 1 ...
 $ osteo_tri : Factor w/ 2 levels "No","Sí": 1 1 1 1 1 1 1 1 1 1 ...
\end{verbatim}

Las 27 variables cubren factores de riesgo (tabaco, alcohol, actividad
física), biomarcadores (HbA1c, calcio, creatinina, PTH), complicaciones
(retinopatía, nefropatía, neuropatía) y medidas óseas (z-scores DXA,
osteoporosis del cuello femoral y trocánter).

\end{tcolorbox}

\begin{center}\rule{0.5\linewidth}{0.5pt}\end{center}

\section{\texorpdfstring{B.2 Análisis descriptivo: \texttt{freq()},
\texttt{grps()} y
\texttt{grpsggp()}}{B.2 Análisis descriptivo: freq(), grps() y grpsggp()}}\label{sec-bio-descriptivo}

\subsection{\texorpdfstring{B.2.1 Tablas de frecuencias:
\texttt{freq()}}{B.2.1 Tablas de frecuencias: freq()}}\label{b.2.1-tablas-de-frecuencias-freq}

\begin{codelisting}

\caption{\label{lst-apendiceB-freq-sig}Signatura de freq()}

\centering{

\begin{Shaded}
\begin{Highlighting}[]
\FunctionTok{freq}\NormalTok{(x, }\AttributeTok{acum =} \ConstantTok{TRUE}\NormalTok{, }\AttributeTok{cuts =} \DecValTok{0}\NormalTok{, }\AttributeTok{agrup =} \ConstantTok{TRUE}\NormalTok{, }\AttributeTok{decs =} \DecValTok{3}\NormalTok{, }\AttributeTok{grf =} \ConstantTok{TRUE}\NormalTok{)}
\end{Highlighting}
\end{Shaded}

}

\end{codelisting}%

Genera una tabla de frecuencias absolutas, relativas y acumuladas. Si
\texttt{cuts\ \textgreater{}\ 0}, agrupa en intervalos de igual
amplitud. El argumento \texttt{grf\ =\ FALSE} suprime el gráfico.

\begin{tcolorbox}[enhanced jigsaw, breakable, titlerule=0mm, rightrule=.15mm, leftrule=.75mm, left=2mm, opacitybacktitle=0.6, colback=white, bottomtitle=1mm, colbacktitle=quarto-callout-tip-color!10!white, title=\textcolor{quarto-callout-tip-color}{\faLightbulb}\hspace{0.5em}{Ejemplo B.1: Distribución de la edad}, toptitle=1mm, coltitle=black, opacityback=0, arc=.35mm, bottomrule=.15mm, toprule=.15mm, colframe=quarto-callout-tip-color-frame]

\begin{codelisting}[H]

\caption{\label{lst-apendiceB-freq-edad}freq() sobre la edad del dataset
osteo}

\centering{

\begin{Shaded}
\begin{Highlighting}[]
\FunctionTok{freq}\NormalTok{(osteo}\SpecialCharTok{$}\NormalTok{edad, }\AttributeTok{grf =} \ConstantTok{FALSE}\NormalTok{)}
\end{Highlighting}
\end{Shaded}

}

\end{codelisting}%

\begin{verbatim}

Distribución de frecuencias
--------------------------------
Variable:  osteo$edad 
n = 94 
 
        x Freq  Prop Prop.Acum
1 (18,23]   27 0.287     0.287
2 (23,28]   22 0.234     0.521
3 (28,33]   14 0.149     0.670
4 (33,38]   15 0.160     0.830
5 (38,43]    4 0.043     0.873
6 (43,48]    5 0.053     0.926
7 (48,53]    3 0.032     0.958
8 (53,58]    3 0.032     0.990
\end{verbatim}

La distribución es \textbf{asimétrica positiva}: más de la mitad de los
pacientes son jóvenes en estadios iniciales de la diabetes.

\begin{codelisting}[H]

\caption{\label{lst-apendiceB-freq-imc}freq() con 5 clases de igual
amplitud para el IMC}

\centering{

\begin{Shaded}
\begin{Highlighting}[]
\FunctionTok{freq}\NormalTok{(osteo}\SpecialCharTok{$}\NormalTok{imc, }\AttributeTok{cuts =} \DecValTok{5}\NormalTok{, }\AttributeTok{grf =} \ConstantTok{FALSE}\NormalTok{)}
\end{Highlighting}
\end{Shaded}

}

\end{codelisting}%

\begin{verbatim}

Distribución de frecuencias
--------------------------------
Variable:  osteo$imc 
n = 94 
 
        x Freq  Prop Prop.Acum
1 (18,22]   37 0.394     0.394
2 (22,26]   36 0.383     0.777
3 (26,30]   13 0.138     0.915
4 (30,34]    5 0.053     0.968
5 (34,38]    1 0.011     0.979
\end{verbatim}

\end{tcolorbox}

\subsection{\texorpdfstring{B.2.2 Estadísticas por grupos:
\texttt{grps()}}{B.2.2 Estadísticas por grupos: grps()}}\label{b.2.2-estaduxedsticas-por-grupos-grps}

\begin{codelisting}

\caption{\label{lst-apendiceB-grps-sig}Signatura de grps()}

\centering{

\begin{Shaded}
\begin{Highlighting}[]
\FunctionTok{grps}\NormalTok{(x, f, }\AttributeTok{ic =} \ConstantTok{FALSE}\NormalTok{, }\AttributeTok{grf =} \ConstantTok{TRUE}\NormalTok{, }\AttributeTok{alfa =} \FloatTok{0.05}\NormalTok{, }\AttributeTok{conf =} \FloatTok{0.95}\NormalTok{, }\AttributeTok{decs =} \DecValTok{3}\NormalTok{)}
\end{Highlighting}
\end{Shaded}

}

\end{codelisting}%

Proporciona media, desviación típica y tamaño de cada grupo. Con
\texttt{ic\ =\ TRUE} añade intervalos de confianza por grupo.

\begin{tcolorbox}[enhanced jigsaw, breakable, titlerule=0mm, rightrule=.15mm, leftrule=.75mm, left=2mm, opacitybacktitle=0.6, colback=white, bottomtitle=1mm, colbacktitle=quarto-callout-tip-color!10!white, title=\textcolor{quarto-callout-tip-color}{\faLightbulb}\hspace{0.5em}{Ejemplo B.2: IMC por sexo}, toptitle=1mm, coltitle=black, opacityback=0, arc=.35mm, bottomrule=.15mm, toprule=.15mm, colframe=quarto-callout-tip-color-frame]

\begin{codelisting}[H]

\caption{\label{lst-apendiceB-grps-imc-sexo}IMC según sexo en el dataset
osteo}

\centering{

\begin{Shaded}
\begin{Highlighting}[]
\FunctionTok{grps}\NormalTok{(osteo}\SpecialCharTok{$}\NormalTok{imc, osteo}\SpecialCharTok{$}\NormalTok{sexo, }\AttributeTok{grf =} \ConstantTok{FALSE}\NormalTok{)}
\end{Highlighting}
\end{Shaded}

}

\end{codelisting}%

\begin{verbatim}

# Descriptiva de osteo$imc por osteo$sexo
# ---------------------------------------
 osteo$sexo 
         n  media    dt
 Hombre 45 23.514 2.890
 Mujer  49 24.294 4.389
 Total  94 23.921 3.748
\end{verbatim}

\begin{codelisting}[H]

\caption{\label{lst-apendiceB-grps-imc-sexo-ic}IMC según sexo con
intervalos de confianza}

\centering{

\begin{Shaded}
\begin{Highlighting}[]
\FunctionTok{grps}\NormalTok{(osteo}\SpecialCharTok{$}\NormalTok{imc, osteo}\SpecialCharTok{$}\NormalTok{sexo, }\AttributeTok{ic =} \ConstantTok{TRUE}\NormalTok{, }\AttributeTok{grf =} \ConstantTok{FALSE}\NormalTok{)}
\end{Highlighting}
\end{Shaded}

}

\end{codelisting}%

\begin{verbatim}


# Descriptiva de osteo$imc por osteo$sexo
# ---------------------------------------
 osteo$sexo 
         n  media    dt ic_inf ic_sup
 Hombre 45 23.514 2.890 22.646 24.382
 Mujer  49 24.294 4.389 23.033 25.555
 Total  94 23.921 3.748 23.153 24.689
 ___________ 
 * ICs para las medias μ[i] y μ global al 95% de confianza 
   (sin corrección por inferencia múltiple)  
\end{verbatim}

\end{tcolorbox}

\subsection{\texorpdfstring{B.2.3 Diagramas por grupos:
\texttt{grpsggp()}}{B.2.3 Diagramas por grupos: grpsggp()}}\label{b.2.3-diagramas-por-grupos-grpsggp}

\begin{codelisting}

\caption{\label{lst-apendiceB-grpsggp-sig}Signatura de grpsggp()}

\centering{

\begin{Shaded}
\begin{Highlighting}[]
\FunctionTok{grpsggp}\NormalTok{(x, }\AttributeTok{f =} \ConstantTok{NULL}\NormalTok{, }\AttributeTok{ggid =} \DecValTok{1}\NormalTok{, }\AttributeTok{lbls =} \ConstantTok{NULL}\NormalTok{, }\AttributeTok{bins =} \ConstantTok{NULL}\NormalTok{, }\AttributeTok{hnmin =} \DecValTok{50}\NormalTok{)}
\end{Highlighting}
\end{Shaded}

}

\end{codelisting}%

Complementa a \texttt{grps()} proporcionando representaciones gráficas
segmentadas por un factor de agrupación: histogramas por grupo,
\emph{boxplots}, diagramas de puntos, diagramas de violín y diagramas de
medias con barras de error.

\begin{center}\rule{0.5\linewidth}{0.5pt}\end{center}

\section{\texorpdfstring{B.3 Intervalos de confianza: \texttt{icm()},
\texttt{icp()},
\texttt{icl()}}{B.3 Intervalos de confianza: icm(), icp(), icl()}}\label{sec-bio-ic}

\subsection{\texorpdfstring{B.3.1 IC para la media:
\texttt{icm()}}{B.3.1 IC para la media: icm()}}\label{b.3.1-ic-para-la-media-icm}

\begin{codelisting}

\caption{\label{lst-apendiceB-icm-sig}Signatura de icm()}

\centering{

\begin{Shaded}
\begin{Highlighting}[]
\FunctionTok{icm}\NormalTok{(}\AttributeTok{x =} \DecValTok{0}\NormalTok{, }\AttributeTok{n =} \DecValTok{0}\NormalTok{, }\AttributeTok{m =} \DecValTok{0}\NormalTok{, }\AttributeTok{s =} \DecValTok{0}\NormalTok{, }\AttributeTok{conf =} \FloatTok{0.95}\NormalTok{, }\AttributeTok{alfa =} \FloatTok{0.05}\NormalTok{, }\AttributeTok{decs =} \DecValTok{3}\NormalTok{)}
\end{Highlighting}
\end{Shaded}

}

\end{codelisting}%

Intervalo de confianza bilateral para la media de una distribución
normal. Acepta los estadísticos suficientes (\(n\), \(\bar{x}\), \(s\))
o directamente los datos (\texttt{x}).

\begin{tcolorbox}[enhanced jigsaw, breakable, titlerule=0mm, rightrule=.15mm, leftrule=.75mm, left=2mm, opacitybacktitle=0.6, colback=white, bottomtitle=1mm, colbacktitle=quarto-callout-tip-color!10!white, title=\textcolor{quarto-callout-tip-color}{\faLightbulb}\hspace{0.5em}{Ejemplo B.3: IC para la media de HbA1c}, toptitle=1mm, coltitle=black, opacityback=0, arc=.35mm, bottomrule=.15mm, toprule=.15mm, colframe=quarto-callout-tip-color-frame]

\begin{codelisting}[H]

\caption{\label{lst-apendiceB-icm-hba1c}IC 95\% para la media
poblacional de HbA1c}

\centering{

\begin{Shaded}
\begin{Highlighting}[]
\FunctionTok{icm}\NormalTok{(}\AttributeTok{m =} \FunctionTok{mean}\NormalTok{(osteo}\SpecialCharTok{$}\NormalTok{hba1c), }\AttributeTok{s =} \FunctionTok{sd}\NormalTok{(osteo}\SpecialCharTok{$}\NormalTok{hba1c), }\AttributeTok{n =} \FunctionTok{nrow}\NormalTok{(osteo))}
\end{Highlighting}
\end{Shaded}

}

\end{codelisting}%

\begin{verbatim}

Intervalo de confianza bilateral para la media de una VA normal 
----------------------------------------------------------------
Información muestral: 
  Tamaño muestral: n =  94 
  Media: m =  8.565 
  Desviación típica: s =  1.799 
  Error estándar de la media: sem =  0.186 

Estimación: 
 95%-IC(µ):  (8.196, 8.933)
  Precisión obtenida: 0.368 
\end{verbatim}

\textbf{Interpretación:} Con 95 \% de confianza, la media poblacional de
HbA1c en diabéticos de la UGR está aproximadamente entre 8.2 \% y 8.9
\%, lo que indica un control glucémico deficiente (meta clínica: HbA1c
\textless{} 7 \%).

\end{tcolorbox}

\subsection{\texorpdfstring{B.3.2 IC para proporción:
\texttt{icp()}}{B.3.2 IC para proporción: icp()}}\label{b.3.2-ic-para-proporciuxf3n-icp}

\begin{codelisting}

\caption{\label{lst-apendiceB-icp-sig}Signatura de icp()}

\centering{

\begin{Shaded}
\begin{Highlighting}[]
\FunctionTok{icp}\NormalTok{(}\AttributeTok{x =} \DecValTok{0}\NormalTok{, }\AttributeTok{n =} \DecValTok{0}\NormalTok{, }\AttributeTok{conf =} \FloatTok{0.95}\NormalTok{, }\AttributeTok{alfa =} \FloatTok{0.05}\NormalTok{, }\AttributeTok{decs =} \DecValTok{4}\NormalTok{)}
\end{Highlighting}
\end{Shaded}

}

\end{codelisting}%

Calcula el IC para una proporción con varios métodos: Clopper--Pearson
(exacto), Wilson, Wald clásico y Agresti--Coull.

\begin{tcolorbox}[enhanced jigsaw, breakable, titlerule=0mm, rightrule=.15mm, leftrule=.75mm, left=2mm, opacitybacktitle=0.6, colback=white, bottomtitle=1mm, colbacktitle=quarto-callout-tip-color!10!white, title=\textcolor{quarto-callout-tip-color}{\faLightbulb}\hspace{0.5em}{Ejemplo B.4: IC para la prevalencia de tabaquismo}, toptitle=1mm, coltitle=black, opacityback=0, arc=.35mm, bottomrule=.15mm, toprule=.15mm, colframe=quarto-callout-tip-color-frame]

\begin{codelisting}[H]

\caption{\label{lst-apendiceB-icp-tabaco}IC 95\% para la prevalencia de
tabaquismo}

\centering{

\begin{Shaded}
\begin{Highlighting}[]
\FunctionTok{icp}\NormalTok{(}\AttributeTok{x =} \FunctionTok{sum}\NormalTok{(osteo}\SpecialCharTok{$}\NormalTok{tabaco }\SpecialCharTok{==} \StringTok{"Sí"}\NormalTok{), }\AttributeTok{n =} \FunctionTok{nrow}\NormalTok{(osteo))}
\end{Highlighting}
\end{Shaded}

}

\end{codelisting}%

\begin{verbatim}

Intervalo de confianza para una proporción binomial 
--------------------------------------------------- 

Información muestral: 
  Tamaño de muestra: n = 94
  Estimación puntual clásica: p=x/n = 0.4362, q=(1-p)=0.5638
  Casos observados: x = 41

# Método exacto (Clooper-Pearson):
  Pseudo-estimación puntual: p' = 0.4405, q'=(1-p')=0.5595
  95%-IC(π): (0.3385, 0.5424) 
  Semiamplitud: 0.1019

# Método de Wilson (con cpc):
  Pseudo-estimación puntual: p' = 0.4388, q'=(1-p')=0.5612
  95%-IC(π): (0.3354, 0.5422) 
  Semiamplitud: 0.1034

# Método de Wald (con cpc):
  Estimación puntual (clásica): p=x/n = 0.4362, q=(1-p)=0.5638
  95%-IC(π): (0.3306, 0.5417) 
  Precisión: 0.1056

# Método de Wald ajustado (Agresti-Coull):
  Estimación puntual: p=(x+2)/(n+4) = 0.4388, q=(1-p)=0.5612
  95%-IC(π): (0.3405, 0.537) 
  Precisión: 0.0982
\end{verbatim}

\end{tcolorbox}

\subsection{\texorpdfstring{B.3.3 IC para tasa Poisson:
\texttt{icl()}}{B.3.3 IC para tasa Poisson: icl()}}\label{b.3.3-ic-para-tasa-poisson-icl}

\begin{codelisting}

\caption{\label{lst-apendiceB-icl-sig}Signatura de icl()}

\centering{

\begin{Shaded}
\begin{Highlighting}[]
\FunctionTok{icl}\NormalTok{(}\AttributeTok{x =} \DecValTok{0}\NormalTok{, }\AttributeTok{n =} \DecValTok{0}\NormalTok{, }\AttributeTok{conf =} \FloatTok{0.95}\NormalTok{, }\AttributeTok{decs =} \DecValTok{4}\NormalTok{)}
\end{Highlighting}
\end{Shaded}

}

\end{codelisting}%

IC para la tasa \(\lambda\) de una distribución de Poisson basado en el
método exacto (gamma--Poisson). Útil en vigilancia epidemiológica para
tasas de incidencia por persona-año.

\begin{center}\rule{0.5\linewidth}{0.5pt}\end{center}

\section{\texorpdfstring{B.4 Tamaño muestral: \texttt{nm()},
\texttt{np()} y
\texttt{nl()}}{B.4 Tamaño muestral: nm(), np() y nl()}}\label{sec-bio-n}

\subsection{\texorpdfstring{B.4.1 Para la media:
\texttt{nm()}}{B.4.1 Para la media: nm()}}\label{b.4.1-para-la-media-nm}

\begin{codelisting}

\caption{\label{lst-apendiceB-nm-sig}Signatura de nm()}

\centering{

\begin{Shaded}
\begin{Highlighting}[]
\FunctionTok{nm}\NormalTok{(d, n, m, s, }\AttributeTok{alfa =} \FloatTok{0.05}\NormalTok{)}
\end{Highlighting}
\end{Shaded}

}

\end{codelisting}%

Calcula el tamaño muestral necesario para estimar una media con
precisión \(d\) (semiamplitud del IC) dada una desviación típica
esperada \(s\).

\begin{tcolorbox}[enhanced jigsaw, breakable, titlerule=0mm, rightrule=.15mm, leftrule=.75mm, left=2mm, opacitybacktitle=0.6, colback=white, bottomtitle=1mm, colbacktitle=quarto-callout-tip-color!10!white, title=\textcolor{quarto-callout-tip-color}{\faLightbulb}\hspace{0.5em}{Ejemplo B.5: ¿Cuántos pacientes para estimar HbA1c con mayor precisión?}, toptitle=1mm, coltitle=black, opacityback=0, arc=.35mm, bottomrule=.15mm, toprule=.15mm, colframe=quarto-callout-tip-color-frame]

\begin{codelisting}[H]

\caption{\label{lst-apendiceB-nm-hba1c}Tamaño muestral para estimar
HbA1c con d = 0.3}

\centering{

\begin{Shaded}
\begin{Highlighting}[]
\FunctionTok{nm}\NormalTok{(}\AttributeTok{d =} \FloatTok{0.3}\NormalTok{, }\AttributeTok{n =} \DecValTok{94}\NormalTok{, }\AttributeTok{m =} \FunctionTok{mean}\NormalTok{(osteo}\SpecialCharTok{$}\NormalTok{hba1c), }\AttributeTok{s =} \FunctionTok{sd}\NormalTok{(osteo}\SpecialCharTok{$}\NormalTok{hba1c), }\AttributeTok{alfa =} \FloatTok{0.05}\NormalTok{)}
\end{Highlighting}
\end{Shaded}

}

\end{codelisting}%

\begin{verbatim}


# Tamaño de muestra para la estimación de la media de una VA normal o su aproximación
# -----------------------------------------------------------------------------------

# Muestra piloto:
  Tamaño muestral:  n = 94 
  Media: m = 8.5649 
  Desviación típica: s = 1.7987 
  Error estandar de la media: sem = 0.1855 
  Precisión observada: d = 0.3684 

# Estimación del tamaño muestral:
  Precisión deseada: δ = 0.3000 
  Tamaño muestral necesario: n ≥ 142 
\end{verbatim}

\end{tcolorbox}

\subsection{\texorpdfstring{B.4.2 Para proporciones:
\texttt{np()}}{B.4.2 Para proporciones: np()}}\label{b.4.2-para-proporciones-np}

\begin{codelisting}

\caption{\label{lst-apendiceB-np-sig}Signatura de np()}

\centering{

\begin{Shaded}
\begin{Highlighting}[]
\FunctionTok{np}\NormalTok{(x, n, d, }\AttributeTok{conf =} \FloatTok{0.95}\NormalTok{, }\AttributeTok{decs =} \DecValTok{5}\NormalTok{)}
\end{Highlighting}
\end{Shaded}

}

\end{codelisting}%

Tamaño muestral para estimar una proporción con precisión \(d\), usando
la estimación piloto \(\hat{p} = x/n\).

\subsection{\texorpdfstring{B.4.3 Para tasa Poisson:
\texttt{nl()}}{B.4.3 Para tasa Poisson: nl()}}\label{b.4.3-para-tasa-poisson-nl}

\begin{codelisting}

\caption{\label{lst-apendiceB-nl-sig}Signatura de nl()}

\centering{

\begin{Shaded}
\begin{Highlighting}[]
\FunctionTok{nl}\NormalTok{(x, }\AttributeTok{n =} \DecValTok{0}\NormalTok{, d, }\AttributeTok{lmax =} \DecValTok{0}\NormalTok{, }\AttributeTok{conf =} \FloatTok{0.95}\NormalTok{, }\AttributeTok{alfa =} \FloatTok{0.05}\NormalTok{)}
\end{Highlighting}
\end{Shaded}

}

\end{codelisting}%

Tamaño muestral para estimar el parámetro \(\lambda\) de una
distribución de Poisson con la precisión deseada. Se puede informar a
partir de observaciones piloto o del valor máximo esperado del parámetro
(\texttt{lmax}).

\begin{center}\rule{0.5\linewidth}{0.5pt}\end{center}

\section{B.5 Contrastes de hipótesis}\label{sec-bio-contrastes}

\subsection{\texorpdfstring{B.5.1 Test de normalidad:
\texttt{testnormal()}}{B.5.1 Test de normalidad: testnormal()}}\label{b.5.1-test-de-normalidad-testnormal}

\begin{codelisting}

\caption{\label{lst-apendiceB-testnormal-sig}Signatura de testnormal()}

\centering{

\begin{Shaded}
\begin{Highlighting}[]
\FunctionTok{testnormal}\NormalTok{(x, }\AttributeTok{obs =} \ConstantTok{TRUE}\NormalTok{, }\AttributeTok{mod =} \ConstantTok{TRUE}\NormalTok{, }\AttributeTok{dens =} \ConstantTok{TRUE}\NormalTok{, }\AttributeTok{sw =} \ConstantTok{TRUE}\NormalTok{, }\AttributeTok{decs =} \DecValTok{3}\NormalTok{, }\AttributeTok{grf =} \ConstantTok{TRUE}\NormalTok{)}
\end{Highlighting}
\end{Shaded}

}

\end{codelisting}%

Contraste de Shapiro--Wilk con visualización integrada (histograma +
densidad teórica + Q--Q plot).

\begin{tcolorbox}[enhanced jigsaw, breakable, titlerule=0mm, rightrule=.15mm, leftrule=.75mm, left=2mm, opacitybacktitle=0.6, colback=white, bottomtitle=1mm, colbacktitle=quarto-callout-tip-color!10!white, title=\textcolor{quarto-callout-tip-color}{\faLightbulb}\hspace{0.5em}{Ejemplo B.6: ¿Es normal la HbA1c?}, toptitle=1mm, coltitle=black, opacityback=0, arc=.35mm, bottomrule=.15mm, toprule=.15mm, colframe=quarto-callout-tip-color-frame]

\begin{codelisting}[H]

\caption{\label{lst-apendiceB-testnormal-hba1c}Test de normalidad para
HbA1c}

\centering{

\begin{Shaded}
\begin{Highlighting}[]
\FunctionTok{testnormal}\NormalTok{(osteo}\SpecialCharTok{$}\NormalTok{hba1c, }\AttributeTok{grf =} \ConstantTok{FALSE}\NormalTok{)}
\end{Highlighting}
\end{Shaded}

}

\end{codelisting}%

\pandocbounded{\includegraphics[keepaspectratio]{chapters/apendiceB_files/figure-pdf/unnamed-chunk-21-1.pdf}}

\begin{verbatim}

# Test de normalidad de Shapiro-Wilk  
 ------------------------------------- 
   n = 94,  W = 0.983,  p = 0.276 
\end{verbatim}

En esta muestra la HbA1c es compatible con una distribución
aproximadamente normal --- contrasta con la edad, que es claramente
asimétrica.

\end{tcolorbox}

\subsection{\texorpdfstring{B.5.2 Test t:
\texttt{testt()}}{B.5.2 Test t: testt()}}\label{b.5.2-test-t-testt}

\begin{codelisting}

\caption{\label{lst-apendiceB-testt-sig}Signatura de testt()}

\centering{

\begin{Shaded}
\begin{Highlighting}[]
\FunctionTok{testt}\NormalTok{(m, s, n, m0, delta, potencia)}
\end{Highlighting}
\end{Shaded}

}

\end{codelisting}%

Integra en una sola función: verificación de normalidad (Shapiro--Wilk),
t-test de una o dos muestras (pareadas o independientes), test de
homogeneidad de varianzas (Welch automático) e IC.

\begin{tcolorbox}[enhanced jigsaw, breakable, titlerule=0mm, rightrule=.15mm, leftrule=.75mm, left=2mm, opacitybacktitle=0.6, colback=white, bottomtitle=1mm, colbacktitle=quarto-callout-tip-color!10!white, title=\textcolor{quarto-callout-tip-color}{\faLightbulb}\hspace{0.5em}{Ejemplo B.7a: Contraste de una muestra --- ¿HbA1c \textgreater{} 7.5 \%?}, toptitle=1mm, coltitle=black, opacityback=0, arc=.35mm, bottomrule=.15mm, toprule=.15mm, colframe=quarto-callout-tip-color-frame]

\begin{codelisting}[H]

\caption{\label{lst-apendiceB-testt-1m}testt() de una muestra: HbA1c
contra el umbral clínico}

\centering{

\begin{Shaded}
\begin{Highlighting}[]
\CommentTok{\# H0: mu = 7.5  vs  H1: mu \textgreater{} 7.5}
\FunctionTok{testt}\NormalTok{(}\AttributeTok{m =}\NormalTok{ osteo}\SpecialCharTok{$}\NormalTok{hba1c, }\AttributeTok{m0 =} \FloatTok{7.5}\NormalTok{, }\AttributeTok{grf =} \ConstantTok{FALSE}\NormalTok{)}
\end{Highlighting}
\end{Shaded}

}

\end{codelisting}%

\begin{verbatim}

# t-Test con una muestra
# ----------------------

# Resumen de 'osteo$hba1c'
    n = 94.000 
    media = 8.565 
    d.t. = 1.799 
    sem = 0.186 

# Estimación de la media μ:
  95%-IC(μ) = (8.196, 8.933) 

# Test de normalidad de Shapiro-Wilk:
  W = 0.983, gl = 94, p = 0.276 

# Test de Student para contrastar H₀:μ=μ₀ con μ₀=7.500
  texp = 5.740, gl = 93 
    p < 0.001 para la alternativa bilateral H₁:μ≠μ₀ 
    p < 0.001 para la alternativa unilateral H₁:μ>μ₀ 

  Estimación del efecto bruto 
  95%-IC(μ-μ₀) = (0.696, 1.433) 
\end{verbatim}

Rechazamos \(H_0\): la media de HbA1c supera significativamente el
umbral clínico de 7.5 \%.

\end{tcolorbox}

\begin{tcolorbox}[enhanced jigsaw, breakable, titlerule=0mm, rightrule=.15mm, leftrule=.75mm, left=2mm, opacitybacktitle=0.6, colback=white, bottomtitle=1mm, colbacktitle=quarto-callout-tip-color!10!white, title=\textcolor{quarto-callout-tip-color}{\faLightbulb}\hspace{0.5em}{Ejemplo B.7b: Contraste de dos muestras --- HbA1c entre hombres y
mujeres}, toptitle=1mm, coltitle=black, opacityback=0, arc=.35mm, bottomrule=.15mm, toprule=.15mm, colframe=quarto-callout-tip-color-frame]

\begin{codelisting}[H]

\caption{\label{lst-apendiceB-testt-2m}testt() de dos muestras
independientes}

\centering{

\begin{Shaded}
\begin{Highlighting}[]
\FunctionTok{testt}\NormalTok{(}\AttributeTok{grupos =}\NormalTok{ osteo}\SpecialCharTok{$}\NormalTok{sexo, }\AttributeTok{m =}\NormalTok{ osteo}\SpecialCharTok{$}\NormalTok{hba1c, }\AttributeTok{grf =} \ConstantTok{FALSE}\NormalTok{)}
\end{Highlighting}
\end{Shaded}

}

\end{codelisting}%

\begin{verbatim}
 

# t-test para 2 Muestras Independientes
# -------------------------------------

# Información muestral y estimación de las medias
  Niveles de agrupación: Hombre, Mujer 

                        n media    dt   sem             IC
  osteo$hba1c [Hombre] 45 8.553 1.745 0.260 (8.029, 9.078)
  osteo$hba1c [Mujer]  49 8.576 1.864 0.266  (8.04, 9.111)
  ____ 
  * IC elaborados al 95% de confianza para estimar μ₁ y μ₂ respectivamente 


# Pruebas de normalidad (test de Shapiro-Wilk)
[1] Para grupo = Hombre, W = 0.958, gl = 45, p = 0.102
[2] Para grupo = Mujer, W = 0.980, gl = 49, p = 0.566

# Test de homogeneidad de varianzas. Fexp = (var₂/var₁)
  Fexp = 1.141, gl₁ = 48, gl₂ = 44, p = 0.660 

# Diferencia de medias (osteo$hba1c [Mujer] - osteo$hba1c [Hombre])
  Hipótesis a contrastar: H₀:μ₁=μ₂ (μ₂-μ₁=0) 

  a) Test de Student (varianzas homogéneas) 
  texp = 0.059, gl = 92 
    p = 0.953 para la alternativa bilateral H₁:μ₁≠μ₂ 
    p = 0.476 para la alternativa unilateral H₁:μ₁<μ₂ 
  95%-IC(μ₂-μ₁) = (-0.719, 0.764) 

  b) Test de Welch (varianzas no homogéneas) 
  texp = 0.060, gl = 91.96 
    p = 0.953 para la alternativa bilateral H₁:μ₁≠μ₂ 
    p = 0.476 para la alternativa unilateral H₁:μ₁<μ₂ 
  95%-IC(μ₂-μ₁) = (-0.717, 0.762) 
\end{verbatim}

\end{tcolorbox}

\subsection{\texorpdfstring{B.5.3 Test F de varianzas:
\texttt{testf()}}{B.5.3 Test F de varianzas: testf()}}\label{b.5.3-test-f-de-varianzas-testf}

\begin{codelisting}

\caption{\label{lst-apendiceB-testf-sig}Signatura de testf()}

\centering{

\begin{Shaded}
\begin{Highlighting}[]
\FunctionTok{testf}\NormalTok{(s1, n1, s2, n2)}
\end{Highlighting}
\end{Shaded}

}

\end{codelisting}%

Contraste \(H_0: \sigma_1^2 = \sigma_2^2\) usando el estadístico
\(F = s_1^2 / s_2^2 \sim F_{n_1-1,\, n_2-1}\).

\begin{tcolorbox}[enhanced jigsaw, breakable, titlerule=0mm, rightrule=.15mm, leftrule=.75mm, left=2mm, opacitybacktitle=0.6, colback=white, bottomtitle=1mm, colbacktitle=quarto-callout-tip-color!10!white, title=\textcolor{quarto-callout-tip-color}{\faLightbulb}\hspace{0.5em}{Ejemplo B.8: ¿Misma variabilidad del IMC en hombres y mujeres?}, toptitle=1mm, coltitle=black, opacityback=0, arc=.35mm, bottomrule=.15mm, toprule=.15mm, colframe=quarto-callout-tip-color-frame]

\begin{codelisting}[H]

\caption{\label{lst-apendiceB-testf-imc}Test F de varianzas del IMC por
sexo}

\centering{

\begin{Shaded}
\begin{Highlighting}[]
\FunctionTok{testf}\NormalTok{(}\AttributeTok{s1 =} \FunctionTok{sd}\NormalTok{(osteo}\SpecialCharTok{$}\NormalTok{imc[osteo}\SpecialCharTok{$}\NormalTok{sexo }\SpecialCharTok{==} \StringTok{"Hombre"}\NormalTok{]),}
      \AttributeTok{n1 =} \FunctionTok{sum}\NormalTok{(osteo}\SpecialCharTok{$}\NormalTok{sexo }\SpecialCharTok{==} \StringTok{"Hombre"}\NormalTok{),}
      \AttributeTok{s2 =} \FunctionTok{sd}\NormalTok{(osteo}\SpecialCharTok{$}\NormalTok{imc[osteo}\SpecialCharTok{$}\NormalTok{sexo }\SpecialCharTok{==} \StringTok{"Mujer"}\NormalTok{]),}
      \AttributeTok{n2 =} \FunctionTok{sum}\NormalTok{(osteo}\SpecialCharTok{$}\NormalTok{sexo }\SpecialCharTok{==} \StringTok{"Mujer"}\NormalTok{))}
\end{Highlighting}
\end{Shaded}

}

\end{codelisting}%

\begin{verbatim}
[[1]]
[1] 2.306145

[[2]]
[1] 48

[[3]]
[1] 44

[[4]]
[1] 0.00584978
\end{verbatim}

\end{tcolorbox}

\subsection{\texorpdfstring{B.5.4 Test de proporciones:
\texttt{testp()}}{B.5.4 Test de proporciones: testp()}}\label{b.5.4-test-de-proporciones-testp}

\begin{codelisting}

\caption{\label{lst-apendiceB-testp-sig}Signatura de testp()}

\centering{

\begin{Shaded}
\begin{Highlighting}[]
\FunctionTok{testp}\NormalTok{(x, n, p0, }\AttributeTok{delta =} \FloatTok{0.05}\NormalTok{, }\AttributeTok{potencia =} \FloatTok{0.9}\NormalTok{)}
\end{Highlighting}
\end{Shaded}

}

\end{codelisting}%

Contraste para una proporción (contra un valor hipotético \(\pi_0\)) o
para comparar dos proporciones independientes o pareadas.

\begin{tcolorbox}[enhanced jigsaw, breakable, titlerule=0mm, rightrule=.15mm, leftrule=.75mm, left=2mm, opacitybacktitle=0.6, colback=white, bottomtitle=1mm, colbacktitle=quarto-callout-tip-color!10!white, title=\textcolor{quarto-callout-tip-color}{\faLightbulb}\hspace{0.5em}{Ejemplo B.9: ¿La prevalencia de tabaquismo supera el 35 \%?}, toptitle=1mm, coltitle=black, opacityback=0, arc=.35mm, bottomrule=.15mm, toprule=.15mm, colframe=quarto-callout-tip-color-frame]

\begin{codelisting}[H]

\caption{\label{lst-apendiceB-testp-tabaco}testp() de una proporción}

\centering{

\begin{Shaded}
\begin{Highlighting}[]
\FunctionTok{testp}\NormalTok{(}\AttributeTok{x =} \FunctionTok{sum}\NormalTok{(osteo}\SpecialCharTok{$}\NormalTok{tabaco }\SpecialCharTok{==} \StringTok{"Sí"}\NormalTok{), }\AttributeTok{n =} \FunctionTok{nrow}\NormalTok{(osteo), }\AttributeTok{p0 =} \FloatTok{0.35}\NormalTok{)}
\end{Highlighting}
\end{Shaded}

}

\end{codelisting}%

\begin{verbatim}

# Test para contrastar una proporción binomial
# --------------------------------------------

# Información muestral
  n = 94
  x = 41   n-x=53
  p = 0.436; q = (1-p) = 0.564

# Test Ho:π=0.350
  [1] Método exacto
                      H1  Fexp Valor.p
    Cola derecha π>0.350 1.410   0.052
    Bilateral    π≠0.350     -   0.103

    95%-IC(π) = (0.339, 0.542) (método de Clooper-Pearson)

  [2] Método aproximado a la distribución normal
    Validez: min(nπ₀, n(1-π₀)) = 32.9 (>5, el método es válido)
    zexp = 1.643,  p  = 0.100

    95%-IC(π) = (0.335, 0.542) (método de Wilson)
\end{verbatim}

\end{tcolorbox}

\subsection{\texorpdfstring{B.5.5 Test de Wilcoxon:
\texttt{testwx()}}{B.5.5 Test de Wilcoxon: testwx()}}\label{b.5.5-test-de-wilcoxon-testwx}

\begin{codelisting}

\caption{\label{lst-apendiceB-testwx-sig}Signatura de testwx()}

\centering{

\begin{Shaded}
\begin{Highlighting}[]
\FunctionTok{testwx}\NormalTok{(m1, m2, }\AttributeTok{par =} \ConstantTok{FALSE}\NormalTok{)}
\end{Highlighting}
\end{Shaded}

}

\end{codelisting}%

Test de Wilcoxon--Mann--Whitney para dos muestras independientes
(\texttt{par\ =\ FALSE}) o test de los rangos con signo de Wilcoxon para
datos pareados (\texttt{par\ =\ TRUE}). Devuelve el tamaño del efecto
\(r\) y la probabilidad de superioridad.

\begin{tcolorbox}[enhanced jigsaw, breakable, titlerule=0mm, rightrule=.15mm, leftrule=.75mm, left=2mm, opacitybacktitle=0.6, colback=white, bottomtitle=1mm, colbacktitle=quarto-callout-tip-color!10!white, title=\textcolor{quarto-callout-tip-color}{\faLightbulb}\hspace{0.5em}{Ejemplo B.10: IMC en diabéticos jóvenes vs.~mayores}, toptitle=1mm, coltitle=black, opacityback=0, arc=.35mm, bottomrule=.15mm, toprule=.15mm, colframe=quarto-callout-tip-color-frame]

\begin{codelisting}[H]

\caption{\label{lst-apendiceB-testwx-imc}Wilcoxon--Mann--Whitney sobre
IMC por grupo de edad}

\centering{

\begin{Shaded}
\begin{Highlighting}[]
\FunctionTok{testwx}\NormalTok{(}\AttributeTok{m1 =}\NormalTok{ osteo}\SpecialCharTok{$}\NormalTok{imc[osteo}\SpecialCharTok{$}\NormalTok{grupo\_edad }\SpecialCharTok{==} \StringTok{"\textless{} 25"}\NormalTok{],}
       \AttributeTok{m2 =}\NormalTok{ osteo}\SpecialCharTok{$}\NormalTok{imc[osteo}\SpecialCharTok{$}\NormalTok{grupo\_edad }\SpecialCharTok{==} \StringTok{"\textgreater{} 33"}\NormalTok{],}
       \AttributeTok{par =} \ConstantTok{FALSE}\NormalTok{,}
       \AttributeTok{grf =} \ConstantTok{FALSE}\NormalTok{)}
\end{Highlighting}
\end{Shaded}

}

\end{codelisting}%

\begin{verbatim}

Test de Wilcoxon/Mann-Whithney para dos muestras independientes 
----------------------------------------------------------------
# Información muestral ---

                                Muestra  n    min     Q1     Q2     Q3    max
1 osteo$imc[osteo$grupo_edad == "< 25"] 32 19.628 21.421 22.845 24.149 32.475
2 osteo$imc[osteo$grupo_edad == "> 33"] 30 18.070 23.382 25.889 28.945 37.333
    RIQ
1 2.728
2 5.563

# Rangos ---

                                Muestra  n Suma_rangos Rango_medio       U
1 osteo$imc[osteo$grupo_edad == "< 25"] 32         807      25.219 681.000
2 osteo$imc[osteo$grupo_edad == "> 33"] 30        1146      38.200 279.000

# Test ---

  U = 279.000; Z = 2.831; W = 279.000; p = 0.005

# Tamaño del efecto ---

  Diferencia de localización: -2.667    95%-IC = (-4.518, -0.778) 
  r = 0.360  (criterio: 0.1 pequeño; 0.3 mediano; >0.5 grande)
  Probabilidad de superioridad PS = 0.709
  (probabilidad de que un valor al azar de M1 sea < a un valor al azar de M2) 
\end{verbatim}

\end{tcolorbox}

\subsection{\texorpdfstring{B.5.6 Test de McNemar:
\texttt{testmcnemar()}}{B.5.6 Test de McNemar: testmcnemar()}}\label{b.5.6-test-de-mcnemar-testmcnemar}

\begin{codelisting}

\caption{\label{lst-apendiceB-testmcnemar-sig}Signatura de
testmcnemar()}

\centering{

\begin{Shaded}
\begin{Highlighting}[]
\FunctionTok{testmcnemar}\NormalTok{(n11, n12, n21, n22, }\AttributeTok{delta =} \FloatTok{0.05}\NormalTok{, }\AttributeTok{alfa =} \FloatTok{0.05}\NormalTok{, }\AttributeTok{beta =} \FloatTok{0.15}\NormalTok{)}
\end{Highlighting}
\end{Shaded}

}

\end{codelisting}%

Para comparar dos proporciones en muestras apareadas (antes/después, o
dos pruebas diagnósticas en los mismos sujetos). El contraste se basa en
las celdas discordantes \(n_{12}\) y \(n_{21}\).

\begin{tcolorbox}[enhanced jigsaw, breakable, titlerule=0mm, rightrule=.15mm, leftrule=.75mm, left=2mm, opacitybacktitle=0.6, colback=white, bottomtitle=1mm, colbacktitle=quarto-callout-tip-color!10!white, title=\textcolor{quarto-callout-tip-color}{\faLightbulb}\hspace{0.5em}{Ejemplo B.11: Cambio en actividad física tras una intervención}, toptitle=1mm, coltitle=black, opacityback=0, arc=.35mm, bottomrule=.15mm, toprule=.15mm, colframe=quarto-callout-tip-color-frame]

\begin{codelisting}[H]

\caption{\label{lst-apendiceB-testmcnemar}testmcnemar() sobre datos
apareados antes/después}

\centering{

\begin{Shaded}
\begin{Highlighting}[]
\CommentTok{\# n11 = activo{-}\textgreater{}activo, n12 = activo{-}\textgreater{}inactivo}
\CommentTok{\# n21 = inactivo{-}\textgreater{}activo, n22 = inactivo{-}\textgreater{}inactivo}
\FunctionTok{testmcnemar}\NormalTok{(}\AttributeTok{n11 =} \DecValTok{35}\NormalTok{, }\AttributeTok{n12 =} \DecValTok{10}\NormalTok{, }\AttributeTok{n21 =} \DecValTok{5}\NormalTok{, }\AttributeTok{n22 =} \DecValTok{44}\NormalTok{,}
            \AttributeTok{fcat =} \FunctionTok{c}\NormalTok{(}\StringTok{"Activo"}\NormalTok{, }\StringTok{"Inactivo"}\NormalTok{),}
            \AttributeTok{ccat =} \FunctionTok{c}\NormalTok{(}\StringTok{"Antes"}\NormalTok{, }\StringTok{"Después"}\NormalTok{))}
\end{Highlighting}
\end{Shaded}

}

\end{codelisting}%

\begin{verbatim}


# Inferencia con dos proporciones (muestras apareadas)
# ----------------------------------------------------

# Frecuencias observadas pretest x posttest
               Antes    Después Total
   Activo         35         10    45
   Inactivo        5         44    49
   Total          40         54    94


# Proporciones observadas pretest x posttest
               Antes    Después  Total
   Activo     0.3723     0.1064 0.4787
   Inactivo   0.0532     0.4681 0.5213
   Total      0.4255     0.5745 1.0000

# Test de McNemar: H₀:π₁₂=π₂₁
   Validez: n₁₂+n₂₁ = 15 > 10 el test es válido 
   Zexp =  1.1619 
               valor.p Alternativa
    Bilateral   0.2453  H₁:π₁₂≠π₂₁
    Unilateral  0.1226  H₁:π₁₂>π₂₁


# Test exacto de Fisher:
  H₀:π₁₂=0.5 para n₁₂ ~ B(n₁₂+n₂₁, π₁₂) 
               Valor.p Alternativa
    Bilateral   0.3018  H₁:π₁₂≠0.5
    Unilateral  0.1509  H₁:π₁₂>0.5
  ____
  * Aquí se alude a la probabilidad total de la discordancia, es decir que π₁₂+π₂₁=1

# Estimación de las proporciones individuales de discordancias π₁₂ y π₂₁ (método de Wald ajustado)
   [1] p₁₂ = 0.1224, 95%-IC(π₁₂) = (0.0575, 0.1873) 
   [2] p₂₁ = 0.0714, 95%-IC(π₂₁) = (0.0204, 0.1224) 


# Intervalo de confianza para la diferencia de 2 proporciones apareadas
   [1] Método de Wald (clásico con cpc):
      Estimación puntual de π₁₂-π₂₁ = 0.0532 
      Validez: n₁₂+n₂₁ = 15 > 5, el IC es válido 
      95%-IC(π₁₂-π₂₁) = (-0.0282, 0.1346) 

   [2] Método de Agresti-Min:
      Estimación puntual de π₁₂-π₂₁ = 0.0521 
      Validez: siempre es válido 
      95%-IC(π₁₂-π₂₁) = (-0.0289, 0.1331) 
\end{verbatim}

\end{tcolorbox}

\begin{center}\rule{0.5\linewidth}{0.5pt}\end{center}

\section{\texorpdfstring{B.6 Regresión lineal: \texttt{rls()} y
\texttt{rlm()}}{B.6 Regresión lineal: rls() y rlm()}}\label{sec-bio-regresion-lineal}

\subsection{\texorpdfstring{B.6.1 Regresión lineal simple:
\texttt{rls()}}{B.6.1 Regresión lineal simple: rls()}}\label{sec-bio-rls}

\begin{codelisting}

\caption{\label{lst-apendiceB-rls-sig}Signatura de rls()}

\centering{

\begin{Shaded}
\begin{Highlighting}[]
\FunctionTok{rls}\NormalTok{(f, data, }\AttributeTok{pred =} \ConstantTok{NULL}\NormalTok{, }\AttributeTok{grf =} \ConstantTok{TRUE}\NormalTok{, }\AttributeTok{alfa =} \FloatTok{0.05}\NormalTok{, }\AttributeTok{decs =} \DecValTok{3}\NormalTok{)}
\end{Highlighting}
\end{Shaded}

}

\end{codelisting}%

Ajusta el modelo \(Y = \beta_0 + \beta_1 X + \varepsilon\) e informa de:
correlación de Pearson con IC, coeficientes con IC, \(R^2\), diagnóstico
de residuos y (opcionalmente) predicciones puntuales con IC.

\begin{tcolorbox}[enhanced jigsaw, breakable, titlerule=0mm, rightrule=.15mm, leftrule=.75mm, left=2mm, opacitybacktitle=0.6, colback=white, bottomtitle=1mm, colbacktitle=quarto-callout-tip-color!10!white, title=\textcolor{quarto-callout-tip-color}{\faLightbulb}\hspace{0.5em}{Ejemplo B.12: ¿La evolución de la diabetes predice la HbA1c?}, toptitle=1mm, coltitle=black, opacityback=0, arc=.35mm, bottomrule=.15mm, toprule=.15mm, colframe=quarto-callout-tip-color-frame]

\begin{codelisting}[H]

\caption{\label{lst-apendiceB-rls-hba1c}rls() --- regresión lineal
simple de HbA1c sobre tevol}

\centering{

\begin{Shaded}
\begin{Highlighting}[]
\FunctionTok{rls}\NormalTok{(hba1c }\SpecialCharTok{\textasciitilde{}}\NormalTok{ tevol, }\AttributeTok{data =}\NormalTok{ osteo, }\AttributeTok{grf =} \ConstantTok{FALSE}\NormalTok{)}
\end{Highlighting}
\end{Shaded}

}

\end{codelisting}%

\begin{verbatim}

Regresión lineal simple 
----------------------------------------------------------------
# Información muestral --- 

  variable  n  media    dt Min  Max Rango
1    hba1c 94  8.565 1.799 4.6 13.8   9.2
2    tevol 94 12.330 8.534 0.0 35.0  35.0

  Cov(hba1c,tevol) = -3.64

# Correlación de Pearson --- 

       r IC_inf IC_sup gl   texp     sig
  -0.237  -0.42 -0.036 92 -2.341 = 0.021

# Modelo lineal --- 

  Modelo:  hba1c ~ tevol 
  R² =  0.056 
  S²residual =  3.087 
  :

         Coef  estim    se ic_inf ic_sup   texp    sig
1 (Constante)  9.181 0.320  8.546  9.816 28.730 <0.001
2       tevol -0.050 0.021 -0.092 -0.008  2.341  0.021

# Distribución residual --- 
  Error estándar residual: 1.757 
       res   zres
min -4.481 -2.585
Q1  -0.981 -0.562
Q2  -0.056 -0.032
Q3   0.956  0.551
max  4.669  2.698

  Test de normalidad residual (Shapiro-Wilk): 
  w =0.989, p= 0.625
\end{verbatim}

\end{tcolorbox}

\subsection{\texorpdfstring{B.6.2 Regresión lineal múltiple:
\texttt{rlm()}}{B.6.2 Regresión lineal múltiple: rlm()}}\label{sec-bio-rlm}

\begin{codelisting}

\caption{\label{lst-apendiceB-rlm-sig}Signatura de rlm()}

\centering{

\begin{Shaded}
\begin{Highlighting}[]
\FunctionTok{rlm}\NormalTok{(f, data, }\AttributeTok{pred =} \ConstantTok{NULL}\NormalTok{, }\AttributeTok{grf =} \ConstantTok{TRUE}\NormalTok{, }\AttributeTok{dfout =} \ConstantTok{FALSE}\NormalTok{, }\AttributeTok{alfa =} \FloatTok{0.05}\NormalTok{, }\AttributeTok{decs =} \DecValTok{3}\NormalTok{)}
\end{Highlighting}
\end{Shaded}

}

\end{codelisting}%

Extiende \texttt{rls()} al caso multivariante:
\(Y = \beta_0 + \beta_1 X_1 + \dots + \beta_p X_p + \varepsilon\).
Devuelve la tabla de coeficientes con IC, \(R^2\) y \(R^2\) ajustado,
distribución de los residuos con su test de normalidad y, si
\texttt{grf\ =\ TRUE}, los cuatro gráficos clásicos de diagnóstico más
histograma de residuos estandarizados.

\begin{tcolorbox}[enhanced jigsaw, breakable, titlerule=0mm, rightrule=.15mm, leftrule=.75mm, left=2mm, opacitybacktitle=0.6, colback=white, bottomtitle=1mm, colbacktitle=quarto-callout-tip-color!10!white, title=\textcolor{quarto-callout-tip-color}{\faLightbulb}\hspace{0.5em}{Ejemplo B.13: HbA1c \textasciitilde{} tevol + edad + IMC}, toptitle=1mm, coltitle=black, opacityback=0, arc=.35mm, bottomrule=.15mm, toprule=.15mm, colframe=quarto-callout-tip-color-frame]

\begin{codelisting}[H]

\caption{\label{lst-apendiceB-rlm-hba1c}rlm() --- regresión lineal
múltiple sobre el dataset osteo}

\centering{

\begin{Shaded}
\begin{Highlighting}[]
\FunctionTok{rlm}\NormalTok{(hba1c }\SpecialCharTok{\textasciitilde{}}\NormalTok{ tevol }\SpecialCharTok{+}\NormalTok{ edad }\SpecialCharTok{+}\NormalTok{ imc, }\AttributeTok{data =}\NormalTok{ osteo, }\AttributeTok{grf =} \ConstantTok{FALSE}\NormalTok{)}
\end{Highlighting}
\end{Shaded}

}

\end{codelisting}%

\begin{verbatim}

Regresión lineal múltiple 
----------------------------------------------------------------
# Información muestral --- 

      Variable  n  Media    DT   Min    Max
hba1c    hba1c 94  8.565 1.799  4.60 13.800
tevol    tevol 94 12.330 8.534  0.00 35.000
edad      edad 94 30.191 9.366 18.00 56.000
imc        imc 94 23.921 3.748 18.07 37.333

# Modelo lineal --- 

   Modelo :  hba1c ~ tevol + edad + imc 
  R² =  0.063  (R²  ajustado  =  0.031 )
  S²residual =  3.134 

   Coeficientes del modelo :

      Termino Estimacion Error_Std IC_inf IC_sup  t_exp      sig
1 (Intercept)      8.601     1.200  6.217 10.984  7.169  < 0.001
2       tevol     -0.046     0.025 -0.095  0.003 -1.883  = 0.063
3        edad     -0.012     0.024 -0.059  0.035 -0.521  = 0.604
4         imc      0.038     0.052 -0.066  0.142  0.722  = 0.472

# Distribución residual --- 
  Error estándar residual:  1.77 
    Residuos Res_Est
min   -4.519  -2.589
Q1    -0.945  -0.545
Q2    -0.010  -0.006
Q3     0.995   0.572
max    4.938   2.923

  Test de normalidad residual (Shapiro-Wilk): 
   w = 0.991 ,   = 0.744 
\end{verbatim}

Tras ajustar simultáneamente por edad e IMC, el efecto del tiempo de
evolución sobre la HbA1c se atenúa, sugiriendo confusión parcial. El
\(R^2\) ajustado es modesto: estas tres variables explican una fracción
pequeña de la variabilidad del control glucémico (ver discusión completa
en Capítulo~\ref{sec-semana10}).

\end{tcolorbox}

\begin{tcolorbox}[enhanced jigsaw, breakable, titlerule=0mm, rightrule=.15mm, leftrule=.75mm, left=2mm, opacitybacktitle=0.6, colback=white, bottomtitle=1mm, colbacktitle=quarto-callout-note-color!10!white, title=\textcolor{quarto-callout-note-color}{\faInfo}\hspace{0.5em}{Predicciones con el modelo}, toptitle=1mm, coltitle=black, opacityback=0, arc=.35mm, bottomrule=.15mm, toprule=.15mm, colframe=quarto-callout-note-color-frame]

\texttt{rlm()} acepta un \texttt{data.frame} en el argumento
\texttt{pred} para devolver predicciones puntuales con intervalos de
confianza (de la media) y de predicción (de una nueva observación):

\begin{codelisting}[H]

\caption{\label{lst-apendiceB-rlm-pred}rlm() con predicciones para
pacientes hipotéticos}

\centering{

\begin{Shaded}
\begin{Highlighting}[]
\NormalTok{nuevos }\OtherTok{\textless{}{-}} \FunctionTok{data.frame}\NormalTok{(}\AttributeTok{tevol =} \FunctionTok{c}\NormalTok{(}\DecValTok{5}\NormalTok{, }\DecValTok{15}\NormalTok{), }\AttributeTok{edad =} \FunctionTok{c}\NormalTok{(}\DecValTok{30}\NormalTok{, }\DecValTok{45}\NormalTok{), }\AttributeTok{imc =} \FunctionTok{c}\NormalTok{(}\DecValTok{25}\NormalTok{, }\DecValTok{28}\NormalTok{))}
\FunctionTok{rlm}\NormalTok{(hba1c }\SpecialCharTok{\textasciitilde{}}\NormalTok{ tevol }\SpecialCharTok{+}\NormalTok{ edad }\SpecialCharTok{+}\NormalTok{ imc, }\AttributeTok{data =}\NormalTok{ osteo, }\AttributeTok{pred =}\NormalTok{ nuevos, }\AttributeTok{grf =} \ConstantTok{FALSE}\NormalTok{)}
\end{Highlighting}
\end{Shaded}

}

\end{codelisting}%

\begin{verbatim}

Regresión lineal múltiple 
----------------------------------------------------------------
# Información muestral --- 

      Variable  n  Media    DT   Min    Max
hba1c    hba1c 94  8.565 1.799  4.60 13.800
tevol    tevol 94 12.330 8.534  0.00 35.000
edad      edad 94 30.191 9.366 18.00 56.000
imc        imc 94 23.921 3.748 18.07 37.333

# Modelo lineal --- 

   Modelo :  hba1c ~ tevol + edad + imc 
  R² =  0.063  (R²  ajustado  =  0.031 )
  S²residual =  3.134 

   Coeficientes del modelo :

      Termino Estimacion Error_Std IC_inf IC_sup  t_exp      sig
1 (Intercept)      8.601     1.200  6.217 10.984  7.169  < 0.001
2       tevol     -0.046     0.025 -0.095  0.003 -1.883  = 0.063
3        edad     -0.012     0.024 -0.059  0.035 -0.521  = 0.604
4         imc      0.038     0.052 -0.066  0.142  0.722  = 0.472

# Pronósticos con el modelo --- 
  Pronosticos puntuales y bandas al 95 % de confianza para 
  promedios IC(m), y para una nueva observación: IC(obs)  

  tevol edad imc  Puntual IC_m_inf IC_m_sup IC_obs_inf IC_obs_sup
1     5   30  25 8.947237 8.428009 9.466464   5.392127   12.50235
2    15   45  28 8.413571 7.678597 9.148545   4.820606   12.00654

# Distribución residual --- 
  Error estándar residual:  1.77 
    Residuos Res_Est
min   -4.519  -2.589
Q1    -0.945  -0.545
Q2    -0.010  -0.006
Q3     0.995   0.572
max    4.938   2.923

  Test de normalidad residual (Shapiro-Wilk): 
   w = 0.991 ,   = 0.744 
\end{verbatim}

\end{tcolorbox}

\begin{center}\rule{0.5\linewidth}{0.5pt}\end{center}

\section{\texorpdfstring{B.7 Regresión logística: \texttt{rlogits()} y
\texttt{rlogitm()}}{B.7 Regresión logística: rlogits() y rlogitm()}}\label{sec-bio-regresion-logistica}

\subsection{\texorpdfstring{B.7.1 Regresión logística simple:
\texttt{rlogits()}}{B.7.1 Regresión logística simple: rlogits()}}\label{sec-bio-rlogits}

\begin{codelisting}

\caption{\label{lst-apendiceB-rlogits-sig}Signatura de rlogits()}

\centering{

\begin{Shaded}
\begin{Highlighting}[]
\FunctionTok{rlogits}\NormalTok{(f, data, }\AttributeTok{grf =} \ConstantTok{FALSE}\NormalTok{, }\AttributeTok{alfa =} \FloatTok{0.05}\NormalTok{, }\AttributeTok{decs =} \DecValTok{3}\NormalTok{)}
\end{Highlighting}
\end{Shaded}

}

\end{codelisting}%

Ajusta un modelo logístico binario simple
\(\mathrm{logit}\, \pi = \beta_0 + \beta_1 X\) y devuelve coeficientes,
odds ratio con IC, devianzas, AIC, pseudo-\(R^2\) de Nagelkerke, test de
Hosmer--Lemeshow y AUC.

\begin{tcolorbox}[enhanced jigsaw, breakable, titlerule=0mm, rightrule=.15mm, leftrule=.75mm, left=2mm, opacitybacktitle=0.6, colback=white, bottomtitle=1mm, colbacktitle=quarto-callout-tip-color!10!white, title=\textcolor{quarto-callout-tip-color}{\faLightbulb}\hspace{0.5em}{Ejemplo B.14: ¿El IMC predice la osteoporosis del cuello femoral?}, toptitle=1mm, coltitle=black, opacityback=0, arc=.35mm, bottomrule=.15mm, toprule=.15mm, colframe=quarto-callout-tip-color-frame]

\begin{codelisting}[H]

\caption{\label{lst-apendiceB-rlogits-osteo}rlogits() --- regresión
logística simple sobre osteoporosis}

\centering{

\begin{Shaded}
\begin{Highlighting}[]
\FunctionTok{rlogits}\NormalTok{(osteo\_cue }\SpecialCharTok{\textasciitilde{}}\NormalTok{ imc, }\AttributeTok{data =}\NormalTok{ osteo)}
\end{Highlighting}
\end{Shaded}

}

\end{codelisting}%

\begin{verbatim}

Regresión logística  simple
----------------------------------------------------------------
# Información muestral --- 

   Tamaño muestral (N inicial) :  94 
   Tamaño muestral tras eliminar valores perdidos (Casos completos) :  94 
   Mínima frecuencia de eventos (n efectivo) :  24 

# Distribución de la variable respuesta (osteo_cue) ---

  Categoria  n Porcentaje
1        No 70     74.468
2        Sí 24     25.532

# Modelo logístico ---  --- 

   Modelo :  osteo_cue ~ imc 
  Devianza residual:  100.176  (Nula:  106.804 )
  AIC:  104.176 
  R² de Nagelkerke:  0.1 

   Test de bondad de ajuste de Hosmer-Lemeshow :
  X² =  7.041 , gl =  8 ,   = 0.532 

   Capacidad discriminante :
   AUC (Area bajo la curva ROC)  =  0.649 

   Coeficientes del modelo :

      Termino Estimacion Error_Std  z_exp      sig     OR OR_inf   OR_sup
1 (Intercept)      3.620     2.044  1.771  = 0.077 37.348  0.937 3055.247
2         imc     -0.202     0.089 -2.256  = 0.024  0.817  0.672    0.957
\end{verbatim}

\end{tcolorbox}

\subsection{\texorpdfstring{B.7.2 Regresión logística múltiple:
\texttt{rlogitm()}}{B.7.2 Regresión logística múltiple: rlogitm()}}\label{sec-bio-rlogitm}

\begin{codelisting}

\caption{\label{lst-apendiceB-rlogitm-sig}Signatura de rlogitm()}

\centering{

\begin{Shaded}
\begin{Highlighting}[]
\FunctionTok{rlogitm}\NormalTok{(f, data, }\AttributeTok{pred =} \ConstantTok{NULL}\NormalTok{, }\AttributeTok{grf =} \ConstantTok{FALSE}\NormalTok{, }\AttributeTok{alfa =} \FloatTok{0.05}\NormalTok{, }\AttributeTok{decs =} \DecValTok{3}\NormalTok{)}
\end{Highlighting}
\end{Shaded}

}

\end{codelisting}%

Extiende \texttt{rlogits()} al caso multivariante. Devuelve la tabla de
coeficientes con OR ajustados e IC, devianzas, AIC, pseudo-\(R^2\) de
Nagelkerke, test de Hosmer--Lemeshow, AUC y, opcionalmente, la curva ROC
(con \texttt{grf\ =\ TRUE}).

\begin{tcolorbox}[enhanced jigsaw, breakable, titlerule=0mm, rightrule=.15mm, leftrule=.75mm, left=2mm, opacitybacktitle=0.6, colback=white, bottomtitle=1mm, colbacktitle=quarto-callout-tip-color!10!white, title=\textcolor{quarto-callout-tip-color}{\faLightbulb}\hspace{0.5em}{Ejemplo B.15: Osteoporosis ajustada por tabaco + edad + IMC + tevol}, toptitle=1mm, coltitle=black, opacityback=0, arc=.35mm, bottomrule=.15mm, toprule=.15mm, colframe=quarto-callout-tip-color-frame]

\begin{codelisting}[H]

\caption{\label{lst-apendiceB-rlogitm-osteo}rlogitm() --- regresión
logística múltiple con curva ROC}

\centering{

\begin{Shaded}
\begin{Highlighting}[]
\FunctionTok{rlogitm}\NormalTok{(osteo\_cue }\SpecialCharTok{\textasciitilde{}}\NormalTok{ tabaco }\SpecialCharTok{+}\NormalTok{ edad }\SpecialCharTok{+}\NormalTok{ imc }\SpecialCharTok{+}\NormalTok{ tevol,}
        \AttributeTok{data =}\NormalTok{ osteo,}
        \AttributeTok{grf  =} \ConstantTok{TRUE}\NormalTok{)}
\end{Highlighting}
\end{Shaded}

}

\end{codelisting}%

\begin{verbatim}

Regresión logística  multivariable
----------------------------------------------------------------
# Información muestral --- 

   Tamaño muestral (N inicial) :  94 
   Tamaño muestral tras eliminar valores perdidos (Casos completos) :  94 
   Mínima frecuencia de eventos (n efectivo) :  24 

# Distribución de la variable respuesta (osteo_cue) ---

  Categoria  n Porcentaje
1        No 70     74.468
2        Sí 24     25.532

# Modelo logístico ---  --- 

   Modelo :  osteo_cue ~ tabaco + edad + imc + tevol 
  Devianza residual:  92.03  (Nula:  106.804 )
  AIC:  102.03 
  R² de Nagelkerke:  0.214 

   Test de bondad de ajuste de Hosmer-Lemeshow :
  X² =  2.067 , gl =  8 ,   = 0.979 

   Capacidad discriminante :
   AUC (Area bajo la curva ROC)  =  0.74 

   Coeficientes del modelo :

      Termino Estimacion Error_Std  z_exp      sig     OR OR_inf   OR_sup
1 (Intercept)      2.893     2.269  1.275  = 0.202 18.045  0.275 2237.565
2    tabacoSí      0.865     0.549  1.575  = 0.115  2.376  0.821    7.218
3        edad     -0.004     0.037 -0.105  = 0.917  0.996  0.925    1.072
4         imc     -0.224     0.105 -2.140  = 0.032  0.799  0.638    0.965
5       tevol      0.071     0.035  2.020  = 0.043  1.074  1.004    1.154
\end{verbatim}

\pandocbounded{\includegraphics[keepaspectratio]{chapters/apendiceB_files/figure-pdf/unnamed-chunk-41-1.pdf}}

La interpretación epidemiológica completa de este modelo se desarrolla
en Sección~\ref{sec-semana11-rlogitm}: los OR ajustados, el test de
Hosmer--Lemeshow, el pseudo-\(R^2\) de Nagelkerke y el AUC se combinan
para evaluar simultáneamente la asociación independiente de cada
predictor y la capacidad discriminante global del modelo.

\end{tcolorbox}

\begin{center}\rule{0.5\linewidth}{0.5pt}\end{center}

\section{\texorpdfstring{B.8 Tablas de contingencia: \texttt{tabla2x2()}
y
\texttt{tablarxc()}}{B.8 Tablas de contingencia: tabla2x2() y tablarxc()}}\label{sec-bio-tablas}

\subsection{\texorpdfstring{B.8.1 Tabla 2×2:
\texttt{tabla2x2()}}{B.8.1 Tabla 2×2: tabla2x2()}}\label{b.8.1-tabla-22-tabla2x2}

\begin{codelisting}

\caption{\label{lst-apendiceB-tabla2x2-sig}Signatura de tabla2x2()}

\centering{

\begin{Shaded}
\begin{Highlighting}[]
\FunctionTok{tabla2x2}\NormalTok{(o, }\AttributeTok{estudio =} \StringTok{"T"}\NormalTok{, fcat, ccat, }\AttributeTok{tablas =} \FunctionTok{c}\NormalTok{(}\StringTok{"F"}\NormalTok{, }\StringTok{"E"}\NormalTok{, }\StringTok{"S"}\NormalTok{))}
\end{Highlighting}
\end{Shaded}

}

\end{codelisting}%

Calcula: frecuencias observadas y esperadas, \(\chi^2\) de Pearson (con
y sin corrección de Yates), test exacto de Fisher y medidas de
asociación (OR, RR, RD) con sus IC al 95 \%, adaptadas al tipo de
estudio (\texttt{estudio\ =\ "T"} transversal, \texttt{"C"} cohorte,
\texttt{"CC"} caso-control, \texttt{"AP"} apareado).

\begin{tcolorbox}[enhanced jigsaw, breakable, titlerule=0mm, rightrule=.15mm, leftrule=.75mm, left=2mm, opacitybacktitle=0.6, colback=white, bottomtitle=1mm, colbacktitle=quarto-callout-tip-color!10!white, title=\textcolor{quarto-callout-tip-color}{\faLightbulb}\hspace{0.5em}{Ejemplo B.16: Tabaquismo y osteoporosis del cuello femoral}, toptitle=1mm, coltitle=black, opacityback=0, arc=.35mm, bottomrule=.15mm, toprule=.15mm, colframe=quarto-callout-tip-color-frame]

\begin{codelisting}[H]

\caption{\label{lst-apendiceB-tabla2x2-osteo}tabla2x2() --- tabaquismo y
osteoporosis}

\centering{

\begin{Shaded}
\begin{Highlighting}[]
\FunctionTok{tabla2x2}\NormalTok{(}\AttributeTok{fvar =}\NormalTok{ osteo}\SpecialCharTok{$}\NormalTok{tabaco, }\AttributeTok{cvar =}\NormalTok{ osteo}\SpecialCharTok{$}\NormalTok{osteo\_cue, }\AttributeTok{o =}\NormalTok{ osteo)}
\end{Highlighting}
\end{Shaded}

}

\end{codelisting}%

\begin{verbatim}

# Análisis de tablas 2x2
# ----------------------

# Frecuencias observadas
          C1   C2 Total
  F1      44    9    53
  F2      26   15    41
  Total   70   24    94


# Test Chi-cuadrado para un estudio transversal

  χ² = 4.662,   gl = 1,  p = 0.031, (cpc = 0.5) 
  Validez: Frecuencia mínima esperada = 10.47  > 3.9  

  Test exacto de Fisher (bilateral): p = 0.035 

  --- Otros criterios χ²:  
  χ² = 4.673,   gl = 1,  p = 0.054, (sin cpc) 
  χ² = 3.699,   gl = 1,  p = 0.054, (cpc de Yates = 47.00) 

# Estimación de la prevalenciaπ en un estudio transversal
  Método de Wald ajustado: 
  p=0.561; 95%-IC(π)=(0.463, 0.659) 

# Medidas de asociación para un estudio transversal
  [!] Las medidas de riesgo se calculan como riesgo de la categoría  
      en la 1a columna (frente a la 2a) para la categoría en la 1a 
      fila (frente a la 2a)

  Riesgo absoluto (diferencia de Berkson; método de Agresti-Caffo): 
  d=0.254; 95%-IC(d)=(0.023, 0.459) 

  Riesgo relativo: 
  Rr=1.676; 95%-IC(Rr)=(0.969, 2.808) 

  Riesgo atribuible:
  Ra=0.335; 95%-IC(Ra)= (-0.049, 0.578) 

  Razón del producto cruzado (odds ratio):
  OR=2.821; 95%-IC(OR)= (1.070, 7.013) 
\end{verbatim}

Los fumadores presentan una OR de aproximadamente 2.8 de tener
osteoporosis del cuello femoral respecto a los no fumadores. La
asociación cruda es estadísticamente significativa (\(p \approx 0.03\)).
Esta asociación se reevalúa con ajuste multivariable en
Sección~\ref{sec-semana11-rlogitm}.

\end{tcolorbox}

\subsection{\texorpdfstring{B.8.2 Tabla r×c:
\texttt{tablarxc()}}{B.8.2 Tabla r×c: tablarxc()}}\label{b.8.2-tabla-rc-tablarxc}

\begin{codelisting}

\caption{\label{lst-apendiceB-tablarxc-sig}Signatura de tablarxc()}

\centering{

\begin{Shaded}
\begin{Highlighting}[]
\FunctionTok{tablarxc}\NormalTok{(frecs, }\AttributeTok{tablas =} \FunctionTok{c}\NormalTok{(}\StringTok{"E"}\NormalTok{, }\StringTok{"F"}\NormalTok{, }\StringTok{"S"}\NormalTok{), fcat, ccat)}
\end{Highlighting}
\end{Shaded}

}

\end{codelisting}%

\begin{tcolorbox}[enhanced jigsaw, breakable, titlerule=0mm, rightrule=.15mm, leftrule=.75mm, left=2mm, opacitybacktitle=0.6, colback=white, bottomtitle=1mm, colbacktitle=quarto-callout-tip-color!10!white, title=\textcolor{quarto-callout-tip-color}{\faLightbulb}\hspace{0.5em}{Ejemplo B.17: Osteoporosis por grupo de edad (tabla 3×2)}, toptitle=1mm, coltitle=black, opacityback=0, arc=.35mm, bottomrule=.15mm, toprule=.15mm, colframe=quarto-callout-tip-color-frame]

\begin{codelisting}[H]

\caption{\label{lst-apendiceB-tablarxc-osteo}tablarxc() --- osteoporosis
del cuello femoral por grupo de edad}

\centering{

\begin{Shaded}
\begin{Highlighting}[]
\NormalTok{tab }\OtherTok{\textless{}{-}} \FunctionTok{table}\NormalTok{(osteo}\SpecialCharTok{$}\NormalTok{grupo\_edad, osteo}\SpecialCharTok{$}\NormalTok{osteo\_cue)}
\FunctionTok{tablarxc}\NormalTok{(}\AttributeTok{frecs =}\NormalTok{ tab,}
         \AttributeTok{fcat  =} \FunctionTok{c}\NormalTok{(}\StringTok{"\textless{} 25"}\NormalTok{, }\StringTok{"25{-}33"}\NormalTok{, }\StringTok{"\textgreater{} 33"}\NormalTok{),}
         \AttributeTok{ccat  =} \FunctionTok{c}\NormalTok{(}\StringTok{"No"}\NormalTok{, }\StringTok{"Sí"}\NormalTok{))}
\end{Highlighting}
\end{Shaded}

}

\end{codelisting}%

\begin{verbatim}

# Test Chi-cuadrado para tablas RxC
# ---------------------------------

# Frecuencias observadas
            No   Sí Total
  < 25      27    5    32
  > 33      22    8    30
  25 - 33   21   11    32
  Total     70   24    94

# Test chi-cuadrado 
  Validez: Frecuencia mínima esperada = 7.66 
           0 frecuencias esperadas son menores a 1
           0 son menores a 5 (el 0% de la tabla) 
  χ²(2 gl) = 2.988, p = 0.224 
\end{verbatim}

\end{tcolorbox}

\begin{center}\rule{0.5\linewidth}{0.5pt}\end{center}

\section{B.9 Tabla resumen de funciones}\label{sec-bio-resumen}

\begin{longtable}[]{@{}
  >{\raggedright\arraybackslash}p{(\linewidth - 6\tabcolsep) * \real{0.1957}}
  >{\raggedright\arraybackslash}p{(\linewidth - 6\tabcolsep) * \real{0.2174}}
  >{\raggedright\arraybackslash}p{(\linewidth - 6\tabcolsep) * \real{0.2391}}
  >{\raggedright\arraybackslash}p{(\linewidth - 6\tabcolsep) * \real{0.3478}}@{}}
\toprule\noalign{}
\begin{minipage}[b]{\linewidth}\raggedright
Función
\end{minipage} & \begin{minipage}[b]{\linewidth}\raggedright
Capítulo
\end{minipage} & \begin{minipage}[b]{\linewidth}\raggedright
Propósito
\end{minipage} & \begin{minipage}[b]{\linewidth}\raggedright
Ejemplo de uso
\end{minipage} \\
\midrule\noalign{}
\endhead
\bottomrule\noalign{}
\endlastfoot
\texttt{freq()} & 1, 2 & Tablas de frecuencias (absolutas, relativas,
acumuladas) & \texttt{freq(osteo\$edad,\ grf=FALSE)} \\
\texttt{grps()} & 2 & Estadísticas descriptivas por grupos &
\texttt{grps(osteo\$imc,\ osteo\$sexo)} \\
\texttt{grpsggp()} & 2 & Diagramas (boxplot, dotplot, violín) por grupos
& \texttt{grpsggp(osteo\$imc,\ osteo\$sexo)} \\
\texttt{icm()} & 7 & IC bilateral para la media &
\texttt{icm(m=8.5,\ s=1.8,\ n=94)} \\
\texttt{icp()} & 7 & IC para proporción (4 métodos) &
\texttt{icp(x=41,\ n=94)} \\
\texttt{icl()} & 7 & IC para tasa Poisson (exacto) &
\texttt{icl(x=12,\ n=1000)} \\
\texttt{nm()} & 7 & Tamaño muestral para media &
\texttt{nm(m=8.5,\ s=1.8,\ d=0.4)} \\
\texttt{np()} & 7 & Tamaño muestral para proporción &
\texttt{np(x=41,\ n=94,\ d=0.1)} \\
\texttt{nl()} & 7 & Tamaño muestral para tasa Poisson &
\texttt{nl(x=3,\ d=1)} \\
\texttt{testnormal()} & 9, 11 & Test de Shapiro--Wilk + gráficos &
\texttt{testnormal(osteo\$hba1c)} \\
\texttt{testt()} & 8, 9 & t-test (1 muestra, 2 muestras, pareado) &
\texttt{testt(m=osteo\$hba1c,\ m0=7.5)} \\
\texttt{testf()} & 9 & Test F de igualdad de varianzas &
\texttt{testf(s1,\ n1,\ s2,\ n2)} \\
\texttt{testp()} & 8 & Test de proporciones (1 o 2 grupos) &
\texttt{testp(x=41,\ n=94,\ p0=0.35)} \\
\texttt{testwx()} & 9 & Wilcoxon--Mann--Whitney / rangos con signo &
\texttt{testwx(m1=imc1,\ m2=imc2)} \\
\texttt{testmcnemar()} & 11 & Test de McNemar (datos apareados) &
\texttt{testmcnemar(n11=35,\ n12=10,\ n21=5,\ n22=44)} \\
\texttt{rls()} & 10 & Regresión lineal simple con diagnósticos &
\texttt{rls(hba1c\ \textasciitilde{}\ tevol,\ data=osteo)} \\
\texttt{rlm()} & 10 & Regresión lineal múltiple con diagnósticos &
\texttt{rlm(hba1c\ \textasciitilde{}\ tevol\ +\ edad\ +\ imc,\ data=osteo)} \\
\texttt{rlogits()} & 11 & Regresión logística simple + AUC &
\texttt{rlogits(osteo\_cue\ \textasciitilde{}\ imc,\ data=osteo)} \\
\texttt{rlogitm()} & 11 & Regresión logística múltiple +
Hosmer--Lemeshow + ROC &
\texttt{rlogitm(osteo\_cue\ \textasciitilde{}\ tabaco\ +\ edad\ +\ imc,\ data=osteo,\ grf=TRUE)} \\
\texttt{tabla2x2()} & 11 & Tabla 2×2: χ², Fisher, OR/RR/RD + IC &
\texttt{tabla2x2(fvar=tabaco,\ cvar=osteo\_cue,\ o=osteo)} \\
\texttt{tablarxc()} & 11 & Tabla r×c: χ² con verificación de supuestos &
\texttt{tablarxc(frecs=table(grupo\_edad,\ osteo\_cue))} \\
\end{longtable}

\begin{tcolorbox}[enhanced jigsaw, breakable, titlerule=0mm, rightrule=.15mm, leftrule=.75mm, left=2mm, opacitybacktitle=0.6, colback=white, bottomtitle=1mm, colbacktitle=quarto-callout-note-color!10!white, title=\textcolor{quarto-callout-note-color}{\faInfo}\hspace{0.5em}{Recursos del paquete}, toptitle=1mm, coltitle=black, opacityback=0, arc=.35mm, bottomrule=.15mm, toprule=.15mm, colframe=quarto-callout-note-color-frame]

\begin{itemize}
\tightlist
\item
  \textbf{Sitio web:} \url{https://migariane.github.io/BioStatR/}
\item
  \textbf{Repositorio GitHub:}
  \url{https://github.com/migariane/BioStatR}
\item
  \textbf{Viñeta introductoria:}
  \url{https://migariane.github.io/BioStatR/articles/BioStatR-overview.html}
\item
  \textbf{Contacto:} \texttt{pfemia@ugr.es} (Pedro Femia Marzo)
\item
  \textbf{Instalación:}
  \texttt{remotes::install\_github("migariane/BioStatR")}
\item
  \textbf{Cita:} \texttt{citation("BioStatR")}
\item
  \textbf{Dataset \texttt{osteo}:} 94 pacientes diabéticos, 27
  variables, Unidad de Diabetes, Facultad de Medicina, UGR.
\end{itemize}

\end{tcolorbox}

\begin{tcolorbox}[enhanced jigsaw, breakable, titlerule=0mm, rightrule=.15mm, leftrule=.75mm, left=2mm, opacitybacktitle=0.6, colback=white, bottomtitle=1mm, colbacktitle=quarto-callout-tip-color!10!white, title=\textcolor{quarto-callout-tip-color}{\faLightbulb}\hspace{0.5em}{Salidas en inglés}, toptitle=1mm, coltitle=black, opacityback=0, arc=.35mm, bottomrule=.15mm, toprule=.15mm, colframe=quarto-callout-tip-color-frame]

El paquete soporta también salida en inglés mediante:

\begin{Shaded}
\begin{Highlighting}[]
\FunctionTok{options}\NormalTok{(}\AttributeTok{BioStatR.lang =} \StringTok{"en"}\NormalTok{)}
\FunctionTok{rlm}\NormalTok{(hba1c }\SpecialCharTok{\textasciitilde{}}\NormalTok{ tevol }\SpecialCharTok{+}\NormalTok{ edad }\SpecialCharTok{+}\NormalTok{ imc, }\AttributeTok{data =}\NormalTok{ osteo, }\AttributeTok{grf =} \ConstantTok{FALSE}\NormalTok{)}
\FunctionTok{options}\NormalTok{(}\AttributeTok{BioStatR.lang =} \StringTok{"es"}\NormalTok{)  }\CommentTok{\# vuelve al castellano (valor por defecto)}
\end{Highlighting}
\end{Shaded}

\end{tcolorbox}

\chapter{Apéndice C: Guía de R para Estadística
Médica}\label{sec-apendiceC}

Esta guía cubre los fundamentos de R orientados a la práctica de la
bioestadística y la investigación médica: importación de datos clínicos,
manejo de variables categóricas (exposición, enfermedad, grupo),
descripción de cohortes, contraste de hipótesis y visualización.
Complementa al Apéndice~\ref{sec-apendiceB} (paquete \texttt{BioStatR})
y al Apéndice~\ref{sec-apendiceA} (repaso matemático).

\begin{tcolorbox}[enhanced jigsaw, breakable, titlerule=0mm, rightrule=.15mm, leftrule=.75mm, left=2mm, opacitybacktitle=0.6, colback=white, bottomtitle=1mm, colbacktitle=quarto-callout-tip-color!10!white, title=\textcolor{quarto-callout-tip-color}{\faLightbulb}\hspace{0.5em}{Cómo usar este apéndice}, toptitle=1mm, coltitle=black, opacityback=0, arc=.35mm, bottomrule=.15mm, toprule=.15mm, colframe=quarto-callout-tip-color-frame]

El material está organizado como referencia rápida. Cada sección incluye
ejemplos ejecutables con datos clínicos simulados y, donde procede,
código no ejecutable (\texttt{eval:\ false}) que muestra la sintaxis
para tareas habituales en investigación médica (importar SPSS, ajustar
modelos de supervivencia, etc.).

\end{tcolorbox}

\begin{center}\rule{0.5\linewidth}{0.5pt}\end{center}

\section{C.1 Conceptos Fundamentales}\label{sec-apendiceC-fundamentos}

\subsection{Asignación, Tipos y
Vectores}\label{asignaciuxf3n-tipos-y-vectores}

\begin{codelisting}

\caption{\label{lst-apendiceC-1}Tipos básicos y vectores en R}

\centering{

\begin{Shaded}
\begin{Highlighting}[]
\CommentTok{\# Asignación con \textless{}{-} (recomendado en R)}
\NormalTok{edad     }\OtherTok{\textless{}{-}} \DecValTok{45}\NormalTok{L            }\CommentTok{\# entero}
\NormalTok{peso\_kg  }\OtherTok{\textless{}{-}} \FloatTok{78.4}           \CommentTok{\# numérico (double)}
\NormalTok{sexo     }\OtherTok{\textless{}{-}} \StringTok{"Mujer"}        \CommentTok{\# carácter}
\NormalTok{diabetes }\OtherTok{\textless{}{-}} \ConstantTok{TRUE}           \CommentTok{\# lógico}

\CommentTok{\# Vectores con c()}
\NormalTok{edades       }\OtherTok{\textless{}{-}} \FunctionTok{c}\NormalTok{(}\DecValTok{25}\NormalTok{, }\DecValTok{30}\NormalTok{, }\DecValTok{45}\NormalTok{, }\DecValTok{22}\NormalTok{, }\DecValTok{60}\NormalTok{)}
\NormalTok{diagnostico  }\OtherTok{\textless{}{-}} \FunctionTok{c}\NormalTok{(}\StringTok{"HTA"}\NormalTok{, }\StringTok{"DM2"}\NormalTok{, }\StringTok{"HTA"}\NormalTok{, }\StringTok{"Sano"}\NormalTok{, }\StringTok{"DM2"}\NormalTok{)}
\NormalTok{hospitalizado }\OtherTok{\textless{}{-}} \FunctionTok{c}\NormalTok{(}\ConstantTok{FALSE}\NormalTok{, }\ConstantTok{TRUE}\NormalTok{, }\ConstantTok{TRUE}\NormalTok{, }\ConstantTok{FALSE}\NormalTok{, }\ConstantTok{TRUE}\NormalTok{)}

\CommentTok{\# Inspeccionar}
\FunctionTok{class}\NormalTok{(edades);  }\FunctionTok{length}\NormalTok{(edades);  }\FunctionTok{range}\NormalTok{(edades)}
\end{Highlighting}
\end{Shaded}

}

\end{codelisting}%

\begin{verbatim}
[1] "numeric"
\end{verbatim}

\begin{verbatim}
[1] 5
\end{verbatim}

\begin{verbatim}
[1] 22 60
\end{verbatim}

\subsection{Indexación: posicional y
lógica}\label{indexaciuxf3n-posicional-y-luxf3gica}

\begin{codelisting}

\caption{\label{lst-apendiceC-2}Indexación de vectores}

\centering{

\begin{Shaded}
\begin{Highlighting}[]
\CommentTok{\# Posicional}
\NormalTok{edades[}\DecValTok{1}\NormalTok{]            }\CommentTok{\# primer paciente}
\end{Highlighting}
\end{Shaded}

}

\end{codelisting}%

\begin{verbatim}
[1] 25
\end{verbatim}

\begin{codelisting}

\caption{\label{lst-apendiceC-2}Indexación de vectores}

\centering{

\begin{Shaded}
\begin{Highlighting}[]
\NormalTok{edades[}\FunctionTok{c}\NormalTok{(}\DecValTok{1}\NormalTok{, }\DecValTok{3}\NormalTok{, }\DecValTok{5}\NormalTok{)]   }\CommentTok{\# pacientes 1, 3 y 5}
\end{Highlighting}
\end{Shaded}

}

\end{codelisting}%

\begin{verbatim}
[1] 25 45 60
\end{verbatim}

\begin{codelisting}

\caption{\label{lst-apendiceC-2}Indexación de vectores}

\centering{

\begin{Shaded}
\begin{Highlighting}[]
\NormalTok{edades[}\SpecialCharTok{{-}}\DecValTok{1}\NormalTok{]           }\CommentTok{\# todos menos el primero}
\end{Highlighting}
\end{Shaded}

}

\end{codelisting}%

\begin{verbatim}
[1] 30 45 22 60
\end{verbatim}

\begin{codelisting}

\caption{\label{lst-apendiceC-2}Indexación de vectores}

\centering{

\begin{Shaded}
\begin{Highlighting}[]
\CommentTok{\# Lógica (filtros clínicos)}
\NormalTok{edades[edades }\SpecialCharTok{\textgreater{}=} \DecValTok{45}\NormalTok{]                }\CommentTok{\# adultos mayores}
\end{Highlighting}
\end{Shaded}

}

\end{codelisting}%

\begin{verbatim}
[1] 45 60
\end{verbatim}

\begin{codelisting}

\caption{\label{lst-apendiceC-2}Indexación de vectores}

\centering{

\begin{Shaded}
\begin{Highlighting}[]
\NormalTok{diagnostico[hospitalizado]          }\CommentTok{\# diagnósticos de hospitalizados}
\end{Highlighting}
\end{Shaded}

}

\end{codelisting}%

\begin{verbatim}
[1] "DM2" "HTA" "DM2"
\end{verbatim}

\begin{codelisting}

\caption{\label{lst-apendiceC-2}Indexación de vectores}

\centering{

\begin{Shaded}
\begin{Highlighting}[]
\NormalTok{edades[diagnostico }\SpecialCharTok{==} \StringTok{"DM2"}\NormalTok{]        }\CommentTok{\# edad de diabéticos}
\end{Highlighting}
\end{Shaded}

}

\end{codelisting}%

\begin{verbatim}
[1] 30 60
\end{verbatim}

\subsection{\texorpdfstring{Valores Perdidos:
\texttt{NA}}{Valores Perdidos: NA}}\label{valores-perdidos-na}

En estudios clínicos los datos faltantes son la norma, no la excepción.
R los representa con \texttt{NA}.

\begin{codelisting}

\caption{\label{lst-apendiceC-3}Manejo de valores perdidos}

\centering{

\begin{Shaded}
\begin{Highlighting}[]
\NormalTok{hba1c }\OtherTok{\textless{}{-}} \FunctionTok{c}\NormalTok{(}\FloatTok{6.2}\NormalTok{, }\FloatTok{7.8}\NormalTok{, }\ConstantTok{NA}\NormalTok{, }\FloatTok{8.5}\NormalTok{, }\FloatTok{6.9}\NormalTok{, }\ConstantTok{NA}\NormalTok{)}

\FunctionTok{mean}\NormalTok{(hba1c)                }\CommentTok{\# NA: cualquier NA contamina el resultado}
\end{Highlighting}
\end{Shaded}

}

\end{codelisting}%

\begin{verbatim}
[1] NA
\end{verbatim}

\begin{codelisting}

\caption{\label{lst-apendiceC-3}Manejo de valores perdidos}

\centering{

\begin{Shaded}
\begin{Highlighting}[]
\FunctionTok{mean}\NormalTok{(hba1c, }\AttributeTok{na.rm =} \ConstantTok{TRUE}\NormalTok{)  }\CommentTok{\# 7.35: ignorar NAs}
\end{Highlighting}
\end{Shaded}

}

\end{codelisting}%

\begin{verbatim}
[1] 7.35
\end{verbatim}

\begin{codelisting}

\caption{\label{lst-apendiceC-3}Manejo de valores perdidos}

\centering{

\begin{Shaded}
\begin{Highlighting}[]
\FunctionTok{is.na}\NormalTok{(hba1c)               }\CommentTok{\# vector lógico}
\end{Highlighting}
\end{Shaded}

}

\end{codelisting}%

\begin{verbatim}
[1] FALSE FALSE  TRUE FALSE FALSE  TRUE
\end{verbatim}

\begin{codelisting}

\caption{\label{lst-apendiceC-3}Manejo de valores perdidos}

\centering{

\begin{Shaded}
\begin{Highlighting}[]
\FunctionTok{sum}\NormalTok{(}\FunctionTok{is.na}\NormalTok{(hba1c))          }\CommentTok{\# 2 valores perdidos}
\end{Highlighting}
\end{Shaded}

}

\end{codelisting}%

\begin{verbatim}
[1] 2
\end{verbatim}

\begin{codelisting}

\caption{\label{lst-apendiceC-3}Manejo de valores perdidos}

\centering{

\begin{Shaded}
\begin{Highlighting}[]
\NormalTok{hba1c[}\SpecialCharTok{!}\FunctionTok{is.na}\NormalTok{(hba1c)]       }\CommentTok{\# eliminar NAs}
\end{Highlighting}
\end{Shaded}

}

\end{codelisting}%

\begin{verbatim}
[1] 6.2 7.8 8.5 6.9
\end{verbatim}

\begin{tcolorbox}[enhanced jigsaw, breakable, titlerule=0mm, rightrule=.15mm, leftrule=.75mm, left=2mm, opacitybacktitle=0.6, colback=white, bottomtitle=1mm, colbacktitle=quarto-callout-warning-color!10!white, title=\textcolor{quarto-callout-warning-color}{\faExclamationTriangle}\hspace{0.5em}{NA no es 0 ni cadena vacía}, toptitle=1mm, coltitle=black, opacityback=0, arc=.35mm, bottomrule=.15mm, toprule=.15mm, colframe=quarto-callout-warning-color-frame]

\texttt{NA\ ==\ NA} devuelve \texttt{NA}, no \texttt{TRUE}. Usa siempre
\texttt{is.na(x)} para comprobar valores perdidos. En análisis
observacional, documenta el patrón de pérdida (¿MCAR, MAR o MNAR?) antes
de elegir entre análisis de casos completos, imputación múltiple o
ponderación inversa.

\end{tcolorbox}

\begin{center}\rule{0.5\linewidth}{0.5pt}\end{center}

\section{C.2 Estructuras de Datos}\label{sec-apendiceC-estructuras}

\subsection{Data Frames y Tibbles}\label{data-frames-y-tibbles}

El \textbf{data frame} es la estructura central de R para datos
tabulares. Un \emph{tibble} (de \texttt{tibble}/\texttt{dplyr}) es un
data frame con impresión más legible.

\begin{codelisting}

\caption{\label{lst-apendiceC-4}Construcción de un data frame clínico}

\centering{

\begin{Shaded}
\begin{Highlighting}[]
\NormalTok{pacientes }\OtherTok{\textless{}{-}} \FunctionTok{data.frame}\NormalTok{(}
  \AttributeTok{id        =} \DecValTok{1}\SpecialCharTok{:}\DecValTok{6}\NormalTok{,}
  \AttributeTok{edad      =} \FunctionTok{c}\NormalTok{(}\DecValTok{25}\NormalTok{, }\DecValTok{30}\NormalTok{, }\DecValTok{45}\NormalTok{, }\DecValTok{60}\NormalTok{, }\DecValTok{55}\NormalTok{, }\DecValTok{72}\NormalTok{),}
  \AttributeTok{sexo      =} \FunctionTok{c}\NormalTok{(}\StringTok{"M"}\NormalTok{, }\StringTok{"F"}\NormalTok{, }\StringTok{"F"}\NormalTok{, }\StringTok{"M"}\NormalTok{, }\StringTok{"F"}\NormalTok{, }\StringTok{"M"}\NormalTok{),}
  \AttributeTok{imc       =} \FunctionTok{c}\NormalTok{(}\FloatTok{22.5}\NormalTok{, }\FloatTok{28.0}\NormalTok{, }\FloatTok{31.2}\NormalTok{, }\FloatTok{26.5}\NormalTok{, }\FloatTok{33.1}\NormalTok{, }\FloatTok{24.8}\NormalTok{),}
  \AttributeTok{hba1c     =} \FunctionTok{c}\NormalTok{(}\FloatTok{5.4}\NormalTok{, }\FloatTok{6.2}\NormalTok{, }\FloatTok{7.8}\NormalTok{, }\FloatTok{6.8}\NormalTok{, }\FloatTok{8.5}\NormalTok{, }\FloatTok{5.9}\NormalTok{),}
  \AttributeTok{diabetes  =} \FunctionTok{c}\NormalTok{(}\DecValTok{0}\NormalTok{, }\DecValTok{0}\NormalTok{, }\DecValTok{1}\NormalTok{, }\DecValTok{1}\NormalTok{, }\DecValTok{1}\NormalTok{, }\DecValTok{0}\NormalTok{)}
\NormalTok{)}

\FunctionTok{str}\NormalTok{(pacientes)        }\CommentTok{\# estructura y tipos}
\end{Highlighting}
\end{Shaded}

}

\end{codelisting}%

\begin{verbatim}
'data.frame':   6 obs. of  6 variables:
 $ id      : int  1 2 3 4 5 6
 $ edad    : num  25 30 45 60 55 72
 $ sexo    : chr  "M" "F" "F" "M" ...
 $ imc     : num  22.5 28 31.2 26.5 33.1 24.8
 $ hba1c   : num  5.4 6.2 7.8 6.8 8.5 5.9
 $ diabetes: num  0 0 1 1 1 0
\end{verbatim}

\begin{codelisting}

\caption{\label{lst-apendiceC-4}Construcción de un data frame clínico}

\centering{

\begin{Shaded}
\begin{Highlighting}[]
\FunctionTok{head}\NormalTok{(pacientes, }\DecValTok{3}\NormalTok{)    }\CommentTok{\# primeras 3 filas}
\end{Highlighting}
\end{Shaded}

}

\end{codelisting}%

\begin{verbatim}
  id edad sexo  imc hba1c diabetes
1  1   25    M 22.5   5.4        0
2  2   30    F 28.0   6.2        0
3  3   45    F 31.2   7.8        1
\end{verbatim}

\begin{codelisting}

\caption{\label{lst-apendiceC-4}Construcción de un data frame clínico}

\centering{

\begin{Shaded}
\begin{Highlighting}[]
\FunctionTok{dim}\NormalTok{(pacientes)        }\CommentTok{\# filas x columnas}
\end{Highlighting}
\end{Shaded}

}

\end{codelisting}%

\begin{verbatim}
[1] 6 6
\end{verbatim}

\begin{codelisting}

\caption{\label{lst-apendiceC-4}Construcción de un data frame clínico}

\centering{

\begin{Shaded}
\begin{Highlighting}[]
\FunctionTok{names}\NormalTok{(pacientes)      }\CommentTok{\# nombres de variables}
\end{Highlighting}
\end{Shaded}

}

\end{codelisting}%

\begin{verbatim}
[1] "id"       "edad"     "sexo"     "imc"      "hba1c"    "diabetes"
\end{verbatim}

\subsection{Listas}\label{listas}

Las listas almacenan objetos heterogéneos (el resultado típico de un
test estadístico es una lista):

\begin{codelisting}

\caption{\label{lst-apendiceC-5}Listas: contenedores heterogéneos}

\centering{

\begin{Shaded}
\begin{Highlighting}[]
\NormalTok{resultado }\OtherTok{\textless{}{-}} \FunctionTok{list}\NormalTok{(}
  \AttributeTok{test       =} \StringTok{"t{-}Student"}\NormalTok{,}
  \AttributeTok{p\_valor    =} \FloatTok{0.034}\NormalTok{,}
  \AttributeTok{ic95       =} \FunctionTok{c}\NormalTok{(}\FloatTok{0.12}\NormalTok{, }\FloatTok{1.85}\NormalTok{),}
  \AttributeTok{n\_total    =} \DecValTok{120}
\NormalTok{)}
\NormalTok{resultado}\SpecialCharTok{$}\NormalTok{p\_valor}
\end{Highlighting}
\end{Shaded}

}

\end{codelisting}%

\begin{verbatim}
[1] 0.034
\end{verbatim}

\begin{codelisting}

\caption{\label{lst-apendiceC-5}Listas: contenedores heterogéneos}

\centering{

\begin{Shaded}
\begin{Highlighting}[]
\NormalTok{resultado[[}\StringTok{"ic95"}\NormalTok{]]}
\end{Highlighting}
\end{Shaded}

}

\end{codelisting}%

\begin{verbatim}
[1] 0.12 1.85
\end{verbatim}

\begin{center}\rule{0.5\linewidth}{0.5pt}\end{center}

\section{C.3 Factores y Variables
Categóricas}\label{sec-apendiceC-factores}

En investigación médica, casi toda variable de exposición o desenlace es
\textbf{categórica} (caso/control, expuesto/no expuesto, grupo de
tratamiento). R las representa con \textbf{factores}.

\begin{codelisting}

\caption{\label{lst-apendiceC-6}Factores, niveles y niveles de
referencia}

\centering{

\begin{Shaded}
\begin{Highlighting}[]
\CommentTok{\# Convertir a factor}
\NormalTok{pacientes}\SpecialCharTok{$}\NormalTok{sexo\_f }\OtherTok{\textless{}{-}} \FunctionTok{factor}\NormalTok{(pacientes}\SpecialCharTok{$}\NormalTok{sexo, }\AttributeTok{levels =} \FunctionTok{c}\NormalTok{(}\StringTok{"M"}\NormalTok{, }\StringTok{"F"}\NormalTok{),}
                           \AttributeTok{labels =} \FunctionTok{c}\NormalTok{(}\StringTok{"Hombre"}\NormalTok{, }\StringTok{"Mujer"}\NormalTok{))}
\NormalTok{pacientes}\SpecialCharTok{$}\NormalTok{diabetes\_f }\OtherTok{\textless{}{-}} \FunctionTok{factor}\NormalTok{(pacientes}\SpecialCharTok{$}\NormalTok{diabetes,}
                               \AttributeTok{levels =} \FunctionTok{c}\NormalTok{(}\DecValTok{0}\NormalTok{, }\DecValTok{1}\NormalTok{),}
                               \AttributeTok{labels =} \FunctionTok{c}\NormalTok{(}\StringTok{"No diabético"}\NormalTok{, }\StringTok{"Diabético"}\NormalTok{))}

\FunctionTok{levels}\NormalTok{(pacientes}\SpecialCharTok{$}\NormalTok{diabetes\_f)    }\CommentTok{\# niveles}
\end{Highlighting}
\end{Shaded}

}

\end{codelisting}%

\begin{verbatim}
[1] "No diabético" "Diabético"   
\end{verbatim}

\begin{codelisting}

\caption{\label{lst-apendiceC-6}Factores, niveles y niveles de
referencia}

\centering{

\begin{Shaded}
\begin{Highlighting}[]
\FunctionTok{table}\NormalTok{(pacientes}\SpecialCharTok{$}\NormalTok{diabetes\_f)     }\CommentTok{\# frecuencia}
\end{Highlighting}
\end{Shaded}

}

\end{codelisting}%

\begin{verbatim}

No diabético    Diabético 
           3            3 
\end{verbatim}

\subsection{\texorpdfstring{Recodificación con \texttt{cut()} y
\texttt{dplyr::case\_when()}}{Recodificación con cut() y dplyr::case\_when()}}\label{recodificaciuxf3n-con-cut-y-dplyrcase_when}

\begin{codelisting}

\caption{\label{lst-apendiceC-7}Categorizar variables continuas}

\centering{

\begin{Shaded}
\begin{Highlighting}[]
\FunctionTok{library}\NormalTok{(dplyr)}

\NormalTok{pacientes }\OtherTok{\textless{}{-}}\NormalTok{ pacientes }\SpecialCharTok{\%\textgreater{}\%}
  \FunctionTok{mutate}\NormalTok{(}
    \CommentTok{\# Grupos de edad clínicos}
    \AttributeTok{edad\_grupo =} \FunctionTok{cut}\NormalTok{(edad,}
                     \AttributeTok{breaks =} \FunctionTok{c}\NormalTok{(}\DecValTok{0}\NormalTok{, }\DecValTok{40}\NormalTok{, }\DecValTok{65}\NormalTok{, }\ConstantTok{Inf}\NormalTok{),}
                     \AttributeTok{labels =} \FunctionTok{c}\NormalTok{(}\StringTok{"\textless{}40"}\NormalTok{, }\StringTok{"40{-}64"}\NormalTok{, }\StringTok{"\textgreater{}=65"}\NormalTok{),}
                     \AttributeTok{right  =} \ConstantTok{FALSE}\NormalTok{),}
    \CommentTok{\# Clasificación IMC (OMS)}
    \AttributeTok{imc\_cat =} \FunctionTok{case\_when}\NormalTok{(}
\NormalTok{      imc }\SpecialCharTok{\textless{}}  \FloatTok{18.5}             \SpecialCharTok{\textasciitilde{}} \StringTok{"Bajo peso"}\NormalTok{,}
\NormalTok{      imc }\SpecialCharTok{\textless{}}  \DecValTok{25}               \SpecialCharTok{\textasciitilde{}} \StringTok{"Normopeso"}\NormalTok{,}
\NormalTok{      imc }\SpecialCharTok{\textless{}}  \DecValTok{30}               \SpecialCharTok{\textasciitilde{}} \StringTok{"Sobrepeso"}\NormalTok{,}
\NormalTok{      imc }\SpecialCharTok{\textgreater{}=} \DecValTok{30}               \SpecialCharTok{\textasciitilde{}} \StringTok{"Obesidad"}
\NormalTok{    ),}
    \AttributeTok{imc\_cat =} \FunctionTok{factor}\NormalTok{(imc\_cat,}
                     \AttributeTok{levels =} \FunctionTok{c}\NormalTok{(}\StringTok{"Bajo peso"}\NormalTok{, }\StringTok{"Normopeso"}\NormalTok{,}
                                \StringTok{"Sobrepeso"}\NormalTok{, }\StringTok{"Obesidad"}\NormalTok{))}
\NormalTok{  )}

\FunctionTok{table}\NormalTok{(pacientes}\SpecialCharTok{$}\NormalTok{edad\_grupo, pacientes}\SpecialCharTok{$}\NormalTok{imc\_cat)}
\end{Highlighting}
\end{Shaded}

}

\end{codelisting}%

\begin{verbatim}
       
        Bajo peso Normopeso Sobrepeso Obesidad
  <40           0         1         1        0
  40-64         0         0         1        2
  >=65          0         1         0        0
\end{verbatim}

\begin{tcolorbox}[enhanced jigsaw, breakable, titlerule=0mm, rightrule=.15mm, leftrule=.75mm, left=2mm, opacitybacktitle=0.6, colback=white, bottomtitle=1mm, colbacktitle=quarto-callout-important-color!10!white, title=\textcolor{quarto-callout-important-color}{\faExclamation}\hspace{0.5em}{Nivel de referencia}, toptitle=1mm, coltitle=black, opacityback=0, arc=.35mm, bottomrule=.15mm, toprule=.15mm, colframe=quarto-callout-important-color-frame]

En regresión (lineal, logística, Cox), el \textbf{nivel de referencia}
del factor define la categoría con la que se compara cada coeficiente.
Por defecto R toma el primer nivel alfabético. Para forzar uno
específico (p.~ej. \emph{Normopeso} como referencia):

\begin{Shaded}
\begin{Highlighting}[]
\NormalTok{pacientes}\SpecialCharTok{$}\NormalTok{imc\_cat }\OtherTok{\textless{}{-}} \FunctionTok{relevel}\NormalTok{(pacientes}\SpecialCharTok{$}\NormalTok{imc\_cat, }\AttributeTok{ref =} \StringTok{"Normopeso"}\NormalTok{)}
\end{Highlighting}
\end{Shaded}

\end{tcolorbox}

\begin{center}\rule{0.5\linewidth}{0.5pt}\end{center}

\section{C.4 Importación de Datos
Clínicos}\label{sec-apendiceC-importacion}

En la práctica los datos llegan en CSV, Excel, SPSS o Stata. El paquete
\texttt{readr} y \texttt{haven} cubren todos los formatos habituales en
investigación biomédica.

\begin{codelisting}

\caption{\label{lst-apendiceC-8}Importación desde distintos formatos}

\centering{

\begin{Shaded}
\begin{Highlighting}[]
\CommentTok{\# CSV (recomendado: readr es más rápido y robusto que read.csv)}
\FunctionTok{library}\NormalTok{(readr)}
\NormalTok{datos }\OtherTok{\textless{}{-}} \FunctionTok{read\_csv}\NormalTok{(}\StringTok{"datos/cohorte.csv"}\NormalTok{)}

\CommentTok{\# Excel (.xlsx)}
\FunctionTok{library}\NormalTok{(readxl)}
\NormalTok{datos }\OtherTok{\textless{}{-}} \FunctionTok{read\_excel}\NormalTok{(}\StringTok{"datos/cohorte.xlsx"}\NormalTok{, }\AttributeTok{sheet =} \StringTok{"Visita1"}\NormalTok{)}

\CommentTok{\# SPSS (.sav) {-}{-} muy común en investigación clínica}
\FunctionTok{library}\NormalTok{(haven)}
\NormalTok{datos }\OtherTok{\textless{}{-}} \FunctionTok{read\_sav}\NormalTok{(}\StringTok{"datos/cohorte.sav"}\NormalTok{)}
\NormalTok{datos }\OtherTok{\textless{}{-}} \FunctionTok{as\_factor}\NormalTok{(datos)   }\CommentTok{\# convertir variables etiquetadas a factor}

\CommentTok{\# Stata (.dta)}
\NormalTok{datos }\OtherTok{\textless{}{-}} \FunctionTok{read\_dta}\NormalTok{(}\StringTok{"datos/cohorte.dta"}\NormalTok{)}

\CommentTok{\# Exportar a CSV con codificación correcta (acentos)}
\FunctionTok{write\_csv}\NormalTok{(datos, }\StringTok{"datos/cohorte\_limpia.csv"}\NormalTok{)}
\end{Highlighting}
\end{Shaded}

}

\end{codelisting}%

\begin{tcolorbox}[enhanced jigsaw, breakable, titlerule=0mm, rightrule=.15mm, leftrule=.75mm, left=2mm, opacitybacktitle=0.6, colback=white, bottomtitle=1mm, colbacktitle=quarto-callout-tip-color!10!white, title=\textcolor{quarto-callout-tip-color}{\faLightbulb}\hspace{0.5em}{Rutas robustas con \texttt{here::here()}}, toptitle=1mm, coltitle=black, opacityback=0, arc=.35mm, bottomrule=.15mm, toprule=.15mm, colframe=quarto-callout-tip-color-frame]

Evita rutas absolutas (\texttt{"C:/Users/.../"}). Usa proyectos de
RStudio y \texttt{here::here("datos",\ "cohorte.csv")} para que el
código funcione en cualquier máquina.

\end{tcolorbox}

\begin{center}\rule{0.5\linewidth}{0.5pt}\end{center}

\section{\texorpdfstring{C.5 Manipulación de Datos con
\texttt{dplyr}}{C.5 Manipulación de Datos con dplyr}}\label{sec-apendiceC-dplyr}

Las cinco verbos esenciales de \texttt{dplyr} para preparar una cohorte:

\begin{longtable}[]{@{}ll@{}}
\toprule\noalign{}
Verbo & Acción \\
\midrule\noalign{}
\endhead
\bottomrule\noalign{}
\endlastfoot
\texttt{select()} & elegir columnas \\
\texttt{filter()} & filtrar filas \\
\texttt{mutate()} & crear/modificar variables \\
\texttt{summarise()} & resumir (con o sin \texttt{group\_by()}) \\
\texttt{arrange()} & ordenar \\
\end{longtable}

\begin{codelisting}

\caption{\label{lst-apendiceC-9}Pipeline típico de preparación de una
cohorte}

\centering{

\begin{Shaded}
\begin{Highlighting}[]
\FunctionTok{library}\NormalTok{(dplyr)}

\NormalTok{resumen }\OtherTok{\textless{}{-}}\NormalTok{ pacientes }\SpecialCharTok{\%\textgreater{}\%}
  \FunctionTok{filter}\NormalTok{(edad }\SpecialCharTok{\textgreater{}=} \DecValTok{40}\NormalTok{) }\SpecialCharTok{\%\textgreater{}\%}                                  \CommentTok{\# adultos}
  \FunctionTok{mutate}\NormalTok{(}\AttributeTok{hba1c\_alta =} \FunctionTok{if\_else}\NormalTok{(hba1c }\SpecialCharTok{\textgreater{}=} \FloatTok{6.5}\NormalTok{, }\DecValTok{1}\NormalTok{L, }\DecValTok{0}\NormalTok{L)) }\SpecialCharTok{\%\textgreater{}\%}  \CommentTok{\# ≥6.5: DM}
  \FunctionTok{group\_by}\NormalTok{(sexo\_f) }\SpecialCharTok{\%\textgreater{}\%}
  \FunctionTok{summarise}\NormalTok{(}
    \AttributeTok{n             =} \FunctionTok{n}\NormalTok{(),}
    \AttributeTok{edad\_media    =} \FunctionTok{mean}\NormalTok{(edad, }\AttributeTok{na.rm =} \ConstantTok{TRUE}\NormalTok{),}
    \AttributeTok{imc\_mediana   =} \FunctionTok{median}\NormalTok{(imc, }\AttributeTok{na.rm =} \ConstantTok{TRUE}\NormalTok{),}
    \AttributeTok{pct\_hba1c\_alta =} \FunctionTok{mean}\NormalTok{(hba1c\_alta, }\AttributeTok{na.rm =} \ConstantTok{TRUE}\NormalTok{) }\SpecialCharTok{*} \DecValTok{100}\NormalTok{,}
    \AttributeTok{.groups =} \StringTok{"drop"}
\NormalTok{  )}
\NormalTok{resumen}
\end{Highlighting}
\end{Shaded}

}

\end{codelisting}%

\begin{verbatim}
# A tibble: 2 x 5
  sexo_f     n edad_media imc_mediana pct_hba1c_alta
  <fct>  <int>      <dbl>       <dbl>          <dbl>
1 Hombre     2         66        25.6             50
2 Mujer      2         50        32.2            100
\end{verbatim}

\subsection{\texorpdfstring{Uniones (\texttt{*\_join}) y
reestructuración
(\texttt{pivot\_*})}{Uniones (*\_join) y reestructuración (pivot\_*)}}\label{uniones-_join-y-reestructuraciuxf3n-pivot_}

\begin{codelisting}

\caption{\label{lst-apendiceC-10}Unir tablas y pivotar a formato
largo/ancho}

\centering{

\begin{Shaded}
\begin{Highlighting}[]
\FunctionTok{library}\NormalTok{(dplyr); }\FunctionTok{library}\NormalTok{(tidyr)}

\CommentTok{\# Unir cohorte con tabla de laboratorio por id de paciente}
\NormalTok{cohorte\_lab }\OtherTok{\textless{}{-}}\NormalTok{ pacientes }\SpecialCharTok{\%\textgreater{}\%}
  \FunctionTok{left\_join}\NormalTok{(laboratorio, }\AttributeTok{by =} \StringTok{"id"}\NormalTok{)}

\CommentTok{\# Pasar de ancho (una columna por visita) a largo (una fila por visita)}
\NormalTok{largo }\OtherTok{\textless{}{-}}\NormalTok{ visitas }\SpecialCharTok{\%\textgreater{}\%}
  \FunctionTok{pivot\_longer}\NormalTok{(}\AttributeTok{cols =} \FunctionTok{starts\_with}\NormalTok{(}\StringTok{"pas\_v"}\NormalTok{),}
               \AttributeTok{names\_to  =} \StringTok{"visita"}\NormalTok{,}
               \AttributeTok{values\_to =} \StringTok{"pas\_mmHg"}\NormalTok{)}
\end{Highlighting}
\end{Shaded}

}

\end{codelisting}%

\begin{center}\rule{0.5\linewidth}{0.5pt}\end{center}

\section{C.6 Estadísticos
Descriptivos}\label{sec-apendiceC-descriptivos}

\subsection{Variables Continuas}\label{variables-continuas}

\begin{codelisting}

\caption{\label{lst-apendiceC-11}Resumen univariante de variables
continuas}

\centering{

\begin{Shaded}
\begin{Highlighting}[]
\FunctionTok{summary}\NormalTok{(pacientes}\SpecialCharTok{$}\NormalTok{imc)                                  }\CommentTok{\# min, Q1, med, mean, Q3, max}
\end{Highlighting}
\end{Shaded}

}

\end{codelisting}%

\begin{verbatim}
   Min. 1st Qu.  Median    Mean 3rd Qu.    Max. 
  22.50   25.23   27.25   27.68   30.40   33.10 
\end{verbatim}

\begin{codelisting}

\caption{\label{lst-apendiceC-11}Resumen univariante de variables
continuas}

\centering{

\begin{Shaded}
\begin{Highlighting}[]
\FunctionTok{c}\NormalTok{(}\AttributeTok{media =} \FunctionTok{mean}\NormalTok{(pacientes}\SpecialCharTok{$}\NormalTok{imc),}
  \AttributeTok{de    =} \FunctionTok{sd}\NormalTok{(pacientes}\SpecialCharTok{$}\NormalTok{imc),}
  \AttributeTok{mediana =} \FunctionTok{median}\NormalTok{(pacientes}\SpecialCharTok{$}\NormalTok{imc),}
  \AttributeTok{RIC =} \FunctionTok{IQR}\NormalTok{(pacientes}\SpecialCharTok{$}\NormalTok{imc))}
\end{Highlighting}
\end{Shaded}

}

\end{codelisting}%

\begin{verbatim}
    media        de   mediana       RIC 
27.683333  3.959503 27.250000  5.175000 
\end{verbatim}

\begin{codelisting}

\caption{\label{lst-apendiceC-11}Resumen univariante de variables
continuas}

\centering{

\begin{Shaded}
\begin{Highlighting}[]
\CommentTok{\# Cuantiles personalizados}
\FunctionTok{quantile}\NormalTok{(pacientes}\SpecialCharTok{$}\NormalTok{imc, }\AttributeTok{probs =} \FunctionTok{c}\NormalTok{(}\FloatTok{0.05}\NormalTok{, }\FloatTok{0.25}\NormalTok{, }\FloatTok{0.5}\NormalTok{, }\FloatTok{0.75}\NormalTok{, }\FloatTok{0.95}\NormalTok{))}
\end{Highlighting}
\end{Shaded}

}

\end{codelisting}%

\begin{verbatim}
    5%    25%    50%    75%    95% 
23.075 25.225 27.250 30.400 32.625 
\end{verbatim}

\subsection{Variables Categóricas}\label{variables-categuxf3ricas-1}

\begin{codelisting}

\caption{\label{lst-apendiceC-12}Frecuencias absolutas y relativas}

\centering{

\begin{Shaded}
\begin{Highlighting}[]
\CommentTok{\# Frecuencias}
\FunctionTok{table}\NormalTok{(pacientes}\SpecialCharTok{$}\NormalTok{imc\_cat)}
\end{Highlighting}
\end{Shaded}

}

\end{codelisting}%

\begin{verbatim}

Bajo peso Normopeso Sobrepeso  Obesidad 
        0         2         2         2 
\end{verbatim}

\begin{codelisting}

\caption{\label{lst-apendiceC-12}Frecuencias absolutas y relativas}

\centering{

\begin{Shaded}
\begin{Highlighting}[]
\CommentTok{\# Proporciones}
\FunctionTok{prop.table}\NormalTok{(}\FunctionTok{table}\NormalTok{(pacientes}\SpecialCharTok{$}\NormalTok{imc\_cat)) }\SpecialCharTok{*} \DecValTok{100}
\end{Highlighting}
\end{Shaded}

}

\end{codelisting}%

\begin{verbatim}

Bajo peso Normopeso Sobrepeso  Obesidad 
  0.00000  33.33333  33.33333  33.33333 
\end{verbatim}

\begin{codelisting}

\caption{\label{lst-apendiceC-12}Frecuencias absolutas y relativas}

\centering{

\begin{Shaded}
\begin{Highlighting}[]
\CommentTok{\# Tabla de contingencia 2x2}
\NormalTok{tabla }\OtherTok{\textless{}{-}} \FunctionTok{table}\NormalTok{(}\AttributeTok{diabetes =}\NormalTok{ pacientes}\SpecialCharTok{$}\NormalTok{diabetes\_f,}
               \AttributeTok{obesidad =}\NormalTok{ pacientes}\SpecialCharTok{$}\NormalTok{imc }\SpecialCharTok{\textgreater{}=} \DecValTok{30}\NormalTok{)}
\NormalTok{tabla}
\end{Highlighting}
\end{Shaded}

}

\end{codelisting}%

\begin{verbatim}
              obesidad
diabetes       FALSE TRUE
  No diabético     3    0
  Diabético        1    2
\end{verbatim}

\begin{codelisting}

\caption{\label{lst-apendiceC-12}Frecuencias absolutas y relativas}

\centering{

\begin{Shaded}
\begin{Highlighting}[]
\FunctionTok{prop.table}\NormalTok{(tabla, }\AttributeTok{margin =} \DecValTok{1}\NormalTok{) }\SpecialCharTok{*} \DecValTok{100}   \CommentTok{\# \% por fila (riesgo en expuestos)}
\end{Highlighting}
\end{Shaded}

}

\end{codelisting}%

\begin{verbatim}
              obesidad
diabetes           FALSE      TRUE
  No diabético 100.00000   0.00000
  Diabético     33.33333  66.66667
\end{verbatim}

\subsection{Resumen Estratificado}\label{resumen-estratificado}

\begin{codelisting}

\caption{\label{lst-apendiceC-13}Descriptivos por grupo (típico de Tabla
1)}

\centering{

\begin{Shaded}
\begin{Highlighting}[]
\NormalTok{pacientes }\SpecialCharTok{\%\textgreater{}\%}
  \FunctionTok{group\_by}\NormalTok{(diabetes\_f) }\SpecialCharTok{\%\textgreater{}\%}
  \FunctionTok{summarise}\NormalTok{(}
    \AttributeTok{n          =} \FunctionTok{n}\NormalTok{(),}
    \AttributeTok{edad\_media =} \FunctionTok{mean}\NormalTok{(edad),}
    \AttributeTok{edad\_de    =} \FunctionTok{sd}\NormalTok{(edad),}
    \AttributeTok{imc\_mediana =} \FunctionTok{median}\NormalTok{(imc),}
    \AttributeTok{imc\_RIC    =} \FunctionTok{IQR}\NormalTok{(imc),}
    \AttributeTok{.groups =} \StringTok{"drop"}
\NormalTok{  )}
\end{Highlighting}
\end{Shaded}

}

\end{codelisting}%

\begin{verbatim}
# A tibble: 2 x 6
  diabetes_f       n edad_media edad_de imc_mediana imc_RIC
  <fct>        <int>      <dbl>   <dbl>       <dbl>   <dbl>
1 No diabético     3       42.3   25.8         24.8    2.75
2 Diabético        3       53.3    7.64        31.2    3.30
\end{verbatim}

\begin{center}\rule{0.5\linewidth}{0.5pt}\end{center}

\section{C.7 Funciones Propias y Control de
Flujo}\label{sec-apendiceC-funciones}

\subsection{Definir funciones}\label{definir-funciones}

\begin{codelisting}

\caption{\label{lst-apendiceC-14}Función propia: cálculo e
interpretación clínica}

\centering{

\begin{Shaded}
\begin{Highlighting}[]
\CommentTok{\# Cálculo del IMC con interpretación}
\NormalTok{clasificar\_imc }\OtherTok{\textless{}{-}} \ControlFlowTok{function}\NormalTok{(peso\_kg, altura\_m) \{}
\NormalTok{  imc }\OtherTok{\textless{}{-}}\NormalTok{ peso\_kg }\SpecialCharTok{/}\NormalTok{ (altura\_m}\SpecialCharTok{\^{}}\DecValTok{2}\NormalTok{)}
\NormalTok{  categoria }\OtherTok{\textless{}{-}} \FunctionTok{cut}\NormalTok{(imc,}
                   \AttributeTok{breaks =} \FunctionTok{c}\NormalTok{(}\DecValTok{0}\NormalTok{, }\FloatTok{18.5}\NormalTok{, }\DecValTok{25}\NormalTok{, }\DecValTok{30}\NormalTok{, }\ConstantTok{Inf}\NormalTok{),}
                   \AttributeTok{labels =} \FunctionTok{c}\NormalTok{(}\StringTok{"Bajo peso"}\NormalTok{, }\StringTok{"Normopeso"}\NormalTok{,}
                              \StringTok{"Sobrepeso"}\NormalTok{, }\StringTok{"Obesidad"}\NormalTok{),}
                   \AttributeTok{right =} \ConstantTok{FALSE}\NormalTok{)}
  \FunctionTok{data.frame}\NormalTok{(}\AttributeTok{imc =} \FunctionTok{round}\NormalTok{(imc, }\DecValTok{1}\NormalTok{), }\AttributeTok{categoria =}\NormalTok{ categoria)}
\NormalTok{\}}

\FunctionTok{clasificar\_imc}\NormalTok{(}\AttributeTok{peso\_kg =} \FunctionTok{c}\NormalTok{(}\DecValTok{55}\NormalTok{, }\DecValTok{72}\NormalTok{, }\DecValTok{95}\NormalTok{),}
               \AttributeTok{altura\_m =} \FunctionTok{c}\NormalTok{(}\FloatTok{1.70}\NormalTok{, }\FloatTok{1.75}\NormalTok{, }\FloatTok{1.78}\NormalTok{))}
\end{Highlighting}
\end{Shaded}

}

\end{codelisting}%

\begin{verbatim}
   imc categoria
1 19.0 Normopeso
2 23.5 Normopeso
3 30.0 Sobrepeso
\end{verbatim}

\subsection{\texorpdfstring{Control de flujo: \texttt{if}, \texttt{for},
\texttt{sapply}}{Control de flujo: if, for, sapply}}\label{control-de-flujo-if-for-sapply}

\begin{codelisting}

\caption{\label{lst-apendiceC-15}Iteración vectorizada con sapply}

\centering{

\begin{Shaded}
\begin{Highlighting}[]
\CommentTok{\# Vectorizado (preferido en R)}
\NormalTok{hba1c\_vals }\OtherTok{\textless{}{-}} \FunctionTok{c}\NormalTok{(}\FloatTok{5.6}\NormalTok{, }\FloatTok{6.1}\NormalTok{, }\FloatTok{6.5}\NormalTok{, }\FloatTok{7.2}\NormalTok{, }\FloatTok{8.0}\NormalTok{)}

\NormalTok{categoria\_dm }\OtherTok{\textless{}{-}} \FunctionTok{sapply}\NormalTok{(hba1c\_vals, }\ControlFlowTok{function}\NormalTok{(x) \{}
  \ControlFlowTok{if}\NormalTok{      (x }\SpecialCharTok{\textless{}}  \FloatTok{5.7}\NormalTok{) }\StringTok{"Normal"}
  \ControlFlowTok{else} \ControlFlowTok{if}\NormalTok{ (x }\SpecialCharTok{\textless{}}  \FloatTok{6.5}\NormalTok{) }\StringTok{"Prediabetes"}
  \ControlFlowTok{else}               \StringTok{"Diabetes"}
\NormalTok{\})}
\NormalTok{categoria\_dm}
\end{Highlighting}
\end{Shaded}

}

\end{codelisting}%

\begin{verbatim}
[1] "Normal"      "Prediabetes" "Diabetes"    "Diabetes"    "Diabetes"   
\end{verbatim}

\begin{tcolorbox}[enhanced jigsaw, breakable, titlerule=0mm, rightrule=.15mm, leftrule=.75mm, left=2mm, opacitybacktitle=0.6, colback=white, bottomtitle=1mm, colbacktitle=quarto-callout-note-color!10!white, title=\textcolor{quarto-callout-note-color}{\faInfo}\hspace{0.5em}{Preferir vectorización sobre bucles}, toptitle=1mm, coltitle=black, opacityback=0, arc=.35mm, bottomrule=.15mm, toprule=.15mm, colframe=quarto-callout-note-color-frame]

En R, \texttt{for} es más lento y verboso que la familia
\texttt{apply}/\texttt{map} y que las funciones vectorizadas. Para un
dataset clínico real (miles de filas), usa siempre
\texttt{dplyr::mutate()} con \texttt{case\_when()} o
\texttt{sapply()}/\texttt{purrr::map\_*()}.

\end{tcolorbox}

\begin{center}\rule{0.5\linewidth}{0.5pt}\end{center}

\section{\texorpdfstring{C.8 Visualización con
\texttt{ggplot2}}{C.8 Visualización con ggplot2}}\label{sec-apendiceC-ggplot}

\texttt{ggplot2} es la herramienta estándar para visualización en
publicaciones biomédicas.

\begin{codelisting}

\caption{\label{lst-apendiceC-16}Histograma con curva de densidad}

\centering{

\begin{Shaded}
\begin{Highlighting}[]
\FunctionTok{library}\NormalTok{(ggplot2)}

\FunctionTok{ggplot}\NormalTok{(pacientes, }\FunctionTok{aes}\NormalTok{(}\AttributeTok{x =}\NormalTok{ imc)) }\SpecialCharTok{+}
  \FunctionTok{geom\_histogram}\NormalTok{(}\FunctionTok{aes}\NormalTok{(}\AttributeTok{y =} \FunctionTok{after\_stat}\NormalTok{(density)),}
                 \AttributeTok{bins =} \DecValTok{6}\NormalTok{, }\AttributeTok{fill =} \StringTok{"steelblue"}\NormalTok{, }\AttributeTok{color =} \StringTok{"white"}\NormalTok{) }\SpecialCharTok{+}
  \FunctionTok{geom\_density}\NormalTok{(}\AttributeTok{color =} \StringTok{"darkred"}\NormalTok{, }\AttributeTok{linewidth =} \DecValTok{1}\NormalTok{) }\SpecialCharTok{+}
  \FunctionTok{labs}\NormalTok{(}\AttributeTok{x =} \StringTok{"IMC (kg/m²)"}\NormalTok{, }\AttributeTok{y =} \StringTok{"Densidad"}\NormalTok{,}
       \AttributeTok{title =} \StringTok{"Distribución del IMC"}\NormalTok{) }\SpecialCharTok{+}
  \FunctionTok{theme\_minimal}\NormalTok{()}
\end{Highlighting}
\end{Shaded}

}

\end{codelisting}%

\pandocbounded{\includegraphics[keepaspectratio]{chapters/apendiceC_files/figure-pdf/unnamed-chunk-17-1.pdf}}

\begin{codelisting}

\caption{\label{lst-apendiceC-17}Boxplot estratificado: HbA1c por estado
diabético}

\centering{

\begin{Shaded}
\begin{Highlighting}[]
\FunctionTok{ggplot}\NormalTok{(pacientes, }\FunctionTok{aes}\NormalTok{(}\AttributeTok{x =}\NormalTok{ diabetes\_f, }\AttributeTok{y =}\NormalTok{ hba1c, }\AttributeTok{fill =}\NormalTok{ diabetes\_f)) }\SpecialCharTok{+}
  \FunctionTok{geom\_boxplot}\NormalTok{(}\AttributeTok{alpha =} \FloatTok{0.7}\NormalTok{) }\SpecialCharTok{+}
  \FunctionTok{geom\_jitter}\NormalTok{(}\AttributeTok{width =} \FloatTok{0.15}\NormalTok{, }\AttributeTok{size =} \DecValTok{2}\NormalTok{) }\SpecialCharTok{+}
  \FunctionTok{labs}\NormalTok{(}\AttributeTok{x =} \ConstantTok{NULL}\NormalTok{, }\AttributeTok{y =} \StringTok{"HbA1c (\%)"}\NormalTok{,}
       \AttributeTok{title =} \StringTok{"HbA1c por estado diabético"}\NormalTok{) }\SpecialCharTok{+}
  \FunctionTok{theme\_minimal}\NormalTok{() }\SpecialCharTok{+}
  \FunctionTok{theme}\NormalTok{(}\AttributeTok{legend.position =} \StringTok{"none"}\NormalTok{)}
\end{Highlighting}
\end{Shaded}

}

\end{codelisting}%

\pandocbounded{\includegraphics[keepaspectratio]{chapters/apendiceC_files/figure-pdf/unnamed-chunk-18-1.pdf}}

\begin{codelisting}

\caption{\label{lst-apendiceC-18}Diagrama de dispersión con línea de
regresión}

\centering{

\begin{Shaded}
\begin{Highlighting}[]
\FunctionTok{ggplot}\NormalTok{(pacientes, }\FunctionTok{aes}\NormalTok{(}\AttributeTok{x =}\NormalTok{ imc, }\AttributeTok{y =}\NormalTok{ hba1c)) }\SpecialCharTok{+}
  \FunctionTok{geom\_point}\NormalTok{(}\AttributeTok{size =} \DecValTok{3}\NormalTok{, }\AttributeTok{color =} \StringTok{"steelblue"}\NormalTok{) }\SpecialCharTok{+}
  \FunctionTok{geom\_smooth}\NormalTok{(}\AttributeTok{method =} \StringTok{"lm"}\NormalTok{, }\AttributeTok{se =} \ConstantTok{TRUE}\NormalTok{, }\AttributeTok{color =} \StringTok{"darkred"}\NormalTok{) }\SpecialCharTok{+}
  \FunctionTok{labs}\NormalTok{(}\AttributeTok{x =} \StringTok{"IMC (kg/m²)"}\NormalTok{, }\AttributeTok{y =} \StringTok{"HbA1c (\%)"}\NormalTok{,}
       \AttributeTok{title =} \StringTok{"Relación IMC{-}HbA1c"}\NormalTok{) }\SpecialCharTok{+}
  \FunctionTok{theme\_minimal}\NormalTok{()}
\end{Highlighting}
\end{Shaded}

}

\end{codelisting}%

\pandocbounded{\includegraphics[keepaspectratio]{chapters/apendiceC_files/figure-pdf/unnamed-chunk-19-1.pdf}}

\begin{center}\rule{0.5\linewidth}{0.5pt}\end{center}

\section{C.9 Tests Estadísticos en R Base}\label{sec-apendiceC-tests}

Tabla de referencia con los contrastes más usados en investigación
médica:

\begin{longtable}[]{@{}
  >{\raggedright\arraybackslash}p{(\linewidth - 2\tabcolsep) * \real{0.5806}}
  >{\raggedright\arraybackslash}p{(\linewidth - 2\tabcolsep) * \real{0.4194}}@{}}
\toprule\noalign{}
\begin{minipage}[b]{\linewidth}\raggedright
Pregunta clínica
\end{minipage} & \begin{minipage}[b]{\linewidth}\raggedright
Función de R
\end{minipage} \\
\midrule\noalign{}
\endhead
\bottomrule\noalign{}
\endlastfoot
¿Difieren las medias de dos grupos? & \texttt{t.test()} \\
¿Difieren las medianas (no normal)? & \texttt{wilcox.test()} \\
¿Difieren tres o más medias? & \texttt{aov()}, \texttt{oneway.test()} \\
¿Asociación entre dos variables categóricas? & \texttt{chisq.test()},
\texttt{fisher.test()} \\
¿Datos apareados (antes/después)? & \texttt{mcnemar.test()},
\texttt{t.test(paired=T)} \\
¿Es la variable normal? & \texttt{shapiro.test()} \\
¿Hay correlación lineal? & \texttt{cor.test()} \\
Regresión lineal & \texttt{lm()} \\
Regresión logística / Poisson &
\texttt{glm(family\ =\ binomial\ /\ poisson)} \\
\end{longtable}

\begin{codelisting}

\caption{\label{lst-apendiceC-19}Ejemplos: t-test, chi-cuadrado,
regresión lineal}

\centering{

\begin{Shaded}
\begin{Highlighting}[]
\CommentTok{\# 1) Comparar HbA1c entre diabéticos y no diabéticos}
\FunctionTok{t.test}\NormalTok{(hba1c }\SpecialCharTok{\textasciitilde{}}\NormalTok{ diabetes\_f, }\AttributeTok{data =}\NormalTok{ pacientes)}
\end{Highlighting}
\end{Shaded}

}

\end{codelisting}%

\begin{verbatim}

    Welch Two Sample t-test

data:  hba1c by diabetes_f
t = -3.4207, df = 2.8523, p-value = 0.04515
alternative hypothesis: true difference in means between group No diabético and group Diabético is not equal to 0
95 percent confidence interval:
 -3.65529198 -0.07804135
sample estimates:
mean in group No diabético    mean in group Diabético 
                  5.833333                   7.700000 
\end{verbatim}

\begin{codelisting}

\caption{\label{lst-apendiceC-19}Ejemplos: t-test, chi-cuadrado,
regresión lineal}

\centering{

\begin{Shaded}
\begin{Highlighting}[]
\CommentTok{\# 2) Asociación diabetes \textasciitilde{} obesidad (Fisher por n pequeño)}
\FunctionTok{fisher.test}\NormalTok{(}\FunctionTok{table}\NormalTok{(pacientes}\SpecialCharTok{$}\NormalTok{diabetes\_f, pacientes}\SpecialCharTok{$}\NormalTok{imc }\SpecialCharTok{\textgreater{}=} \DecValTok{30}\NormalTok{))}
\end{Highlighting}
\end{Shaded}

}

\end{codelisting}%

\begin{verbatim}

    Fisher's Exact Test for Count Data

data:  table(pacientes$diabetes_f, pacientes$imc >= 30)
p-value = 0.4
alternative hypothesis: true odds ratio is not equal to 1
95 percent confidence interval:
 0.2031288       Inf
sample estimates:
odds ratio 
       Inf 
\end{verbatim}

\begin{codelisting}

\caption{\label{lst-apendiceC-19}Ejemplos: t-test, chi-cuadrado,
regresión lineal}

\centering{

\begin{Shaded}
\begin{Highlighting}[]
\CommentTok{\# 3) Regresión lineal: HbA1c \textasciitilde{} IMC + edad}
\NormalTok{modelo }\OtherTok{\textless{}{-}} \FunctionTok{lm}\NormalTok{(hba1c }\SpecialCharTok{\textasciitilde{}}\NormalTok{ imc }\SpecialCharTok{+}\NormalTok{ edad, }\AttributeTok{data =}\NormalTok{ pacientes)}
\FunctionTok{summary}\NormalTok{(modelo)}
\end{Highlighting}
\end{Shaded}

}

\end{codelisting}%

\begin{verbatim}

Call:
lm(formula = hba1c ~ imc + edad, data = pacientes)

Residuals:
      1       2       3       4       5       6 
 0.2795 -0.4996  0.0762  0.2576  0.1586 -0.2723 

Coefficients:
             Estimate Std. Error t value Pr(>|t|)   
(Intercept) -1.379357   1.316864  -1.047  0.37184   
imc          0.279194   0.046683   5.981  0.00936 **
edad         0.008718   0.010235   0.852  0.45692   
---
Signif. codes:  0 '***' 0.001 '**' 0.01 '*' 0.05 '.' 0.1 ' ' 1

Residual standard error: 0.4079 on 3 degrees of freedom
Multiple R-squared:  0.9288,    Adjusted R-squared:  0.8814 
F-statistic: 19.57 on 2 and 3 DF,  p-value: 0.01899
\end{verbatim}

\begin{codelisting}

\caption{\label{lst-apendiceC-19}Ejemplos: t-test, chi-cuadrado,
regresión lineal}

\centering{

\begin{Shaded}
\begin{Highlighting}[]
\FunctionTok{confint}\NormalTok{(modelo)}
\end{Highlighting}
\end{Shaded}

}

\end{codelisting}%

\begin{verbatim}
                  2.5 %     97.5 %
(Intercept) -5.57020549 2.81149089
imc          0.13062701 0.42776027
edad        -0.02385378 0.04128988
\end{verbatim}

\begin{codelisting}

\caption{\label{lst-apendiceC-20}Regresión logística: probabilidad de
diabetes}

\centering{

\begin{Shaded}
\begin{Highlighting}[]
\CommentTok{\# Pacientes simulados más amplios para que el modelo sea estimable}
\FunctionTok{set.seed}\NormalTok{(}\DecValTok{123}\NormalTok{)}
\NormalTok{n }\OtherTok{\textless{}{-}} \DecValTok{200}
\NormalTok{sim }\OtherTok{\textless{}{-}} \FunctionTok{data.frame}\NormalTok{(}
  \AttributeTok{edad =} \FunctionTok{rnorm}\NormalTok{(n, }\DecValTok{55}\NormalTok{, }\DecValTok{12}\NormalTok{),}
  \AttributeTok{imc  =} \FunctionTok{rnorm}\NormalTok{(n, }\DecValTok{28}\NormalTok{, }\DecValTok{5}\NormalTok{)}
\NormalTok{)}
\NormalTok{sim}\SpecialCharTok{$}\NormalTok{diabetes }\OtherTok{\textless{}{-}} \FunctionTok{rbinom}\NormalTok{(n, }\DecValTok{1}\NormalTok{,}
                       \FunctionTok{plogis}\NormalTok{(}\SpecialCharTok{{-}}\DecValTok{7} \SpecialCharTok{+} \FloatTok{0.05}\SpecialCharTok{*}\NormalTok{sim}\SpecialCharTok{$}\NormalTok{edad }\SpecialCharTok{+} \FloatTok{0.12}\SpecialCharTok{*}\NormalTok{sim}\SpecialCharTok{$}\NormalTok{imc))}

\NormalTok{mod\_log }\OtherTok{\textless{}{-}} \FunctionTok{glm}\NormalTok{(diabetes }\SpecialCharTok{\textasciitilde{}}\NormalTok{ edad }\SpecialCharTok{+}\NormalTok{ imc, }\AttributeTok{data =}\NormalTok{ sim, }\AttributeTok{family =}\NormalTok{ binomial)}
\CommentTok{\# Odds ratios con IC95\%}
\FunctionTok{exp}\NormalTok{(}\FunctionTok{cbind}\NormalTok{(}\AttributeTok{OR =} \FunctionTok{coef}\NormalTok{(mod\_log), }\FunctionTok{confint.default}\NormalTok{(mod\_log)))}
\end{Highlighting}
\end{Shaded}

}

\end{codelisting}%

\begin{verbatim}
                      OR        2.5 %      97.5 %
(Intercept) 0.0005718736 3.304375e-05 0.009897163
edad        1.0414789433 1.011409e+00 1.072443165
imc         1.1652214003 1.085568e+00 1.250719795
\end{verbatim}

\begin{center}\rule{0.5\linewidth}{0.5pt}\end{center}

\section{C.10 Paquetes Esenciales en
Bioestadística}\label{sec-apendiceC-paquetes}

Selección curada de paquetes que el investigador clínico encuentra de
forma recurrente. El bloque siguiente es informativo (no ejecutado).

\begin{codelisting}

\caption{\label{lst-apendiceC-21}Paquetes recomendados según tarea}

\centering{

\begin{Shaded}
\begin{Highlighting}[]
\CommentTok{\# {-}{-}{-} Reportes clínicos y "Tabla 1" {-}{-}{-}{-}{-}{-}{-}{-}{-}{-}{-}{-}{-}{-}{-}{-}{-}{-}{-}{-}{-}{-}{-}{-}{-}{-}{-}{-}{-}{-}{-}{-}{-}{-}{-}{-}{-}{-}{-}}
\FunctionTok{library}\NormalTok{(gtsummary)   }\CommentTok{\# Tablas publicables a partir de modelos / data frames}
\FunctionTok{library}\NormalTok{(tableone)    }\CommentTok{\# Tabla 1 demográfica clásica}

\NormalTok{tabla1 }\OtherTok{\textless{}{-}} \FunctionTok{tbl\_summary}\NormalTok{(pacientes,}
                      \AttributeTok{by =}\NormalTok{ diabetes\_f,}
                      \AttributeTok{missing =} \StringTok{"no"}\NormalTok{) }\SpecialCharTok{\%\textgreater{}\%}
  \FunctionTok{add\_p}\NormalTok{() }\SpecialCharTok{\%\textgreater{}\%} \FunctionTok{add\_overall}\NormalTok{()}

\CommentTok{\# {-}{-}{-} Análisis epidemiológico {-}{-}{-}{-}{-}{-}{-}{-}{-}{-}{-}{-}{-}{-}{-}{-}{-}{-}{-}{-}{-}{-}{-}{-}{-}{-}{-}{-}{-}{-}{-}{-}{-}{-}{-}{-}{-}{-}{-}{-}{-}{-}{-}{-}{-}}
\FunctionTok{library}\NormalTok{(epiR)        }\CommentTok{\# Medidas de asociación 2x2 (RR, OR, NNT)}
\FunctionTok{library}\NormalTok{(epitools)    }\CommentTok{\# Tablas, cálculo de IC para incidencia / prevalencia}

\CommentTok{\# {-}{-}{-} Supervivencia {-}{-}{-}{-}{-}{-}{-}{-}{-}{-}{-}{-}{-}{-}{-}{-}{-}{-}{-}{-}{-}{-}{-}{-}{-}{-}{-}{-}{-}{-}{-}{-}{-}{-}{-}{-}{-}{-}{-}{-}{-}{-}{-}{-}{-}{-}{-}{-}{-}{-}{-}{-}{-}{-}{-}}
\FunctionTok{library}\NormalTok{(survival)}
\FunctionTok{library}\NormalTok{(survminer)}
\NormalTok{fit\_km  }\OtherTok{\textless{}{-}} \FunctionTok{survfit}\NormalTok{(}\FunctionTok{Surv}\NormalTok{(tiempo, evento) }\SpecialCharTok{\textasciitilde{}}\NormalTok{ grupo, }\AttributeTok{data =}\NormalTok{ cohorte)}
\FunctionTok{ggsurvplot}\NormalTok{(fit\_km, }\AttributeTok{risk.table =} \ConstantTok{TRUE}\NormalTok{, }\AttributeTok{pval =} \ConstantTok{TRUE}\NormalTok{)}

\NormalTok{fit\_cox }\OtherTok{\textless{}{-}} \FunctionTok{coxph}\NormalTok{(}\FunctionTok{Surv}\NormalTok{(tiempo, evento) }\SpecialCharTok{\textasciitilde{}}\NormalTok{ edad }\SpecialCharTok{+}\NormalTok{ sexo }\SpecialCharTok{+}\NormalTok{ tratamiento,}
                 \AttributeTok{data =}\NormalTok{ cohorte)}
\FunctionTok{summary}\NormalTok{(fit\_cox)}

\CommentTok{\# {-}{-}{-} Inferencia causal / emparejamiento {-}{-}{-}{-}{-}{-}{-}{-}{-}{-}{-}{-}{-}{-}{-}{-}{-}{-}{-}{-}{-}{-}{-}{-}{-}{-}{-}{-}{-}{-}{-}{-}{-}{-}}
\FunctionTok{library}\NormalTok{(MatchIt)     }\CommentTok{\# Propensity score matching}
\FunctionTok{library}\NormalTok{(WeightIt)    }\CommentTok{\# Ponderación inversa de probabilidad (IPTW)}

\CommentTok{\# {-}{-}{-} Datos perdidos {-}{-}{-}{-}{-}{-}{-}{-}{-}{-}{-}{-}{-}{-}{-}{-}{-}{-}{-}{-}{-}{-}{-}{-}{-}{-}{-}{-}{-}{-}{-}{-}{-}{-}{-}{-}{-}{-}{-}{-}{-}{-}{-}{-}{-}{-}{-}{-}{-}{-}{-}{-}{-}{-}}
\FunctionTok{library}\NormalTok{(mice)        }\CommentTok{\# Imputación múltiple (FCS)}
\FunctionTok{library}\NormalTok{(naniar)      }\CommentTok{\# Visualización del patrón de missing}

\CommentTok{\# {-}{-}{-} Modelos: limpieza y extracción de resultados {-}{-}{-}{-}{-}{-}{-}{-}{-}{-}{-}{-}{-}{-}{-}{-}{-}{-}{-}{-}{-}{-}{-}{-}}
\FunctionTok{library}\NormalTok{(broom)       }\CommentTok{\# tidy(), glance(), augment() para modelos}
\FunctionTok{library}\NormalTok{(performance) }\CommentTok{\# Diagnósticos de regresión (colinealidad, supuestos)}

\CommentTok{\# {-}{-}{-} Meta{-}análisis {-}{-}{-}{-}{-}{-}{-}{-}{-}{-}{-}{-}{-}{-}{-}{-}{-}{-}{-}{-}{-}{-}{-}{-}{-}{-}{-}{-}{-}{-}{-}{-}{-}{-}{-}{-}{-}{-}{-}{-}{-}{-}{-}{-}{-}{-}{-}{-}{-}{-}{-}{-}{-}{-}{-}}
\FunctionTok{library}\NormalTok{(meta)        }\CommentTok{\# Meta{-}análisis clásico (efecto fijo / aleatorio)}
\FunctionTok{library}\NormalTok{(metafor)     }\CommentTok{\# Meta{-}análisis y meta{-}regresión flexibles}
\end{Highlighting}
\end{Shaded}

}

\end{codelisting}%

\begin{tcolorbox}[enhanced jigsaw, breakable, titlerule=0mm, rightrule=.15mm, leftrule=.75mm, left=2mm, opacitybacktitle=0.6, colback=white, bottomtitle=1mm, colbacktitle=quarto-callout-tip-color!10!white, title=\textcolor{quarto-callout-tip-color}{\faLightbulb}\hspace{0.5em}{Citar los paquetes que usas}, toptitle=1mm, coltitle=black, opacityback=0, arc=.35mm, bottomrule=.15mm, toprule=.15mm, colframe=quarto-callout-tip-color-frame]

La revista te pedirá la versión de R y los paquetes. Usa
\texttt{citation("survival")} y \texttt{sessionInfo()} en el cierre del
análisis y guarda la salida.

\end{tcolorbox}

\begin{center}\rule{0.5\linewidth}{0.5pt}\end{center}

\section{C.11 Reproducibilidad en Investigación
Médica}\label{sec-apendiceC-reproducibilidad}

Un análisis médico debe poder reproducirse exactamente meses o años
después (revisión por pares, auditoría, dudas del comité ético).

\begin{codelisting}

\caption{\label{lst-apendiceC-22}Buenas prácticas de reproducibilidad}

\centering{

\begin{Shaded}
\begin{Highlighting}[]
\CommentTok{\# 1) Semilla aleatoria SIEMPRE antes de cualquier simulación o aleatorización}
\FunctionTok{set.seed}\NormalTok{(}\DecValTok{20260101}\NormalTok{)}

\CommentTok{\# 2) Proyecto de RStudio + rutas relativas con here}
\FunctionTok{library}\NormalTok{(here)}
\NormalTok{datos }\OtherTok{\textless{}{-}} \FunctionTok{read\_csv}\NormalTok{(}\FunctionTok{here}\NormalTok{(}\StringTok{"datos"}\NormalTok{, }\StringTok{"cohorte\_v3.csv"}\NormalTok{))}

\CommentTok{\# 3) Capturar el entorno: versión de R y paquetes}
\FunctionTok{sessionInfo}\NormalTok{()}

\CommentTok{\# 4) (Opcional) renv para bloquear versiones de paquetes}
\CommentTok{\# install.packages("renv"); renv::init(); renv::snapshot()}

\CommentTok{\# 5) Análisis dentro de un documento Quarto/Rmd}
\CommentTok{\#    {-}\textgreater{} el código y los resultados viven juntos}
\end{Highlighting}
\end{Shaded}

}

\end{codelisting}%

\begin{tcolorbox}[enhanced jigsaw, breakable, titlerule=0mm, rightrule=.15mm, leftrule=.75mm, left=2mm, opacitybacktitle=0.6, colback=white, bottomtitle=1mm, colbacktitle=quarto-callout-warning-color!10!white, title=\textcolor{quarto-callout-warning-color}{\faExclamationTriangle}\hspace{0.5em}{Sin semilla, sin reproducibilidad}, toptitle=1mm, coltitle=black, opacityback=0, arc=.35mm, bottomrule=.15mm, toprule=.15mm, colframe=quarto-callout-warning-color-frame]

Cualquier procedimiento estocástico (bootstrap, validación cruzada,
imputación múltiple, simulación) \textbf{debe} ir precedido por
\texttt{set.seed()}. Sin semilla, dos ejecuciones del mismo script darán
p-valores e intervalos distintos --- inaceptable en un análisis clínico.

\end{tcolorbox}

\begin{center}\rule{0.5\linewidth}{0.5pt}\end{center}

\section{C.12 Resolución de Problemas
Comunes}\label{sec-apendiceC-troubleshooting}

\begin{longtable}[]{@{}
  >{\raggedright\arraybackslash}p{(\linewidth - 4\tabcolsep) * \real{0.3017}}
  >{\raggedright\arraybackslash}p{(\linewidth - 4\tabcolsep) * \real{0.3296}}
  >{\raggedright\arraybackslash}p{(\linewidth - 4\tabcolsep) * \real{0.3687}}@{}}
\toprule\noalign{}
\begin{minipage}[b]{\linewidth}\raggedright
Síntoma
\end{minipage} & \begin{minipage}[b]{\linewidth}\raggedright
Causa habitual
\end{minipage} & \begin{minipage}[b]{\linewidth}\raggedright
Solución
\end{minipage} \\
\midrule\noalign{}
\endhead
\bottomrule\noalign{}
\endlastfoot
\texttt{mean(x)} devuelve \texttt{NA} & Hay valores perdidos &
\texttt{mean(x,\ na.rm\ =\ TRUE)} \\
Coeficiente de un factor no aparece & Es el nivel de referencia &
\texttt{relevel(f,\ ref\ =\ "...")} \\
El gráfico no se ve en RStudio & Ventana cerrada o tamaño nulo &
\texttt{dev.off()} y volver a ejecutar \\
\texttt{cannot\ find\ function\ "..."} & Paquete no cargado &
\texttt{library(paquete)} o \texttt{install.packages("paquete")} \\
\texttt{argument\ "x"\ is\ missing} & Olvidaste un argumento & Consulta
\texttt{?nombre\_funcion} \\
Lectura de CSV con acentos rotos & Codificación incorrecta &
\texttt{read\_csv(...,\ locale\ =\ locale(encoding\ =\ "Latin1"))} \\
\texttt{non-numeric\ argument\ to\ binary\ operator} & Tipo carácter
donde esperabas numérico & \texttt{as.numeric(x)} (cuidado con factores:
\texttt{as.numeric(as.character(x))}) \\
Resultados distintos en cada ejecución & Falta \texttt{set.seed()} &
Fijar semilla antes del análisis \\
\texttt{object\ \textquotesingle{}...\textquotesingle{}\ not\ found} &
Nombre mal escrito o variable de otro entorno & \texttt{ls()} para
listar objetos, \texttt{names(df)} para variables \\
\end{longtable}

\subsection{Comandos de ayuda}\label{comandos-de-ayuda}

\begin{codelisting}

\caption{\label{lst-apendiceC-23}Sistema de ayuda de R}

\centering{

\begin{Shaded}
\begin{Highlighting}[]
\NormalTok{?mean                  }\CommentTok{\# ayuda de una función}
\NormalTok{??}\StringTok{"survival analysis"}  \CommentTok{\# búsqueda por palabra clave}
\FunctionTok{help}\NormalTok{(}\AttributeTok{package =} \StringTok{"dplyr"}\NormalTok{)}
\FunctionTok{vignette}\NormalTok{(}\StringTok{"dplyr"}\NormalTok{)      }\CommentTok{\# tutoriales largos del paquete}
\FunctionTok{example}\NormalTok{(lm)            }\CommentTok{\# ejemplos ejecutables}

\CommentTok{\# Diagnóstico rápido del entorno}
\FunctionTok{getwd}\NormalTok{();  }\FunctionTok{setwd}\NormalTok{(}\StringTok{"ruta"}\NormalTok{)}
\FunctionTok{ls}\NormalTok{();     }\FunctionTok{rm}\NormalTok{(}\AttributeTok{list =} \FunctionTok{ls}\NormalTok{())   }\CommentTok{\# listar / limpiar objetos}
\FunctionTok{sessionInfo}\NormalTok{()}
\end{Highlighting}
\end{Shaded}

}

\end{codelisting}%

\begin{center}\rule{0.5\linewidth}{0.5pt}\end{center}

\section{Resumen}\label{resumen}

Esta guía te ha llevado del vector más simple a un flujo completo de
análisis clínico en R: importar (C.4), preparar (C.3, C.5), describir
(C.6), visualizar (C.8), contrastar (C.9) y dejar todo reproducible
(C.11). Para funciones de alto nivel específicas del curso
(\texttt{freq()}, \texttt{icm()}, \texttt{tabla2x2()}, etc.) consulta el
Apéndice~\ref{sec-apendiceB}. Para el repaso matemático subyacente, el
Apéndice~\ref{sec-apendiceA}.

\chapter{Apéndice D: Investigación Reproducible con Git, GitHub y
RStudio}\label{sec-apendiceD}

La \textbf{investigación reproducible} es el principio según el cual
cualquier análisis cuantitativo ---y, por extensión, cualquier resultado
publicado--- debe poder ser replicado por terceros a partir de los
datos, el código y la documentación originales. En bioestadística, donde
las decisiones clínicas pueden depender de un único valor \emph{p} o de
un único intervalo de confianza, la reproducibilidad no es un lujo sino
una exigencia ética.

Este apéndice introduce las tres herramientas mínimas que todo
investigador en ciencias de la salud debería dominar para producir
análisis reproducibles: \textbf{Git} (control de versiones),
\textbf{GitHub} (plataforma de colaboración basada en Git) y
\textbf{RStudio} (entorno de desarrollo integrado para R que incorpora
un cliente Git).

\begin{center}\rule{0.5\linewidth}{0.5pt}\end{center}

\section{D.1 ¿Por qué investigación
reproducible?}\label{sec-apendiceD-rationale}

\begin{tcolorbox}[enhanced jigsaw, breakable, titlerule=0mm, rightrule=.15mm, leftrule=.75mm, left=2mm, opacitybacktitle=0.6, colback=white, bottomtitle=1mm, colbacktitle=quarto-callout-important-color!10!white, title=\textcolor{quarto-callout-important-color}{\faExclamation}\hspace{0.5em}{La crisis de la replicabilidad}, toptitle=1mm, coltitle=black, opacityback=0, arc=.35mm, bottomrule=.15mm, toprule=.15mm, colframe=quarto-callout-important-color-frame]

Múltiples estudios desde 2005 (Ioannidis, \emph{PLoS Medicine}) han
documentado que una fracción sustancial de los hallazgos publicados en
biomedicina no logra replicarse cuando se repite el análisis con los
datos originales o se realiza un estudio independiente. Esta
\textbf{crisis de replicabilidad} ha motivado iniciativas
internacionales (ICMJE, NIH, COS) que hoy exigen compartir datos y
código como condición para publicar.

\end{tcolorbox}

\subsection{D.1.1 Niveles de
reproducibilidad}\label{d.1.1-niveles-de-reproducibilidad}

Siguiendo la taxonomía de \textbf{Goodman, Fanelli \& Ioannidis (2016)}:

\begin{longtable}[]{@{}
  >{\raggedright\arraybackslash}p{(\linewidth - 4\tabcolsep) * \real{0.2333}}
  >{\raggedright\arraybackslash}p{(\linewidth - 4\tabcolsep) * \real{0.4333}}
  >{\raggedright\arraybackslash}p{(\linewidth - 4\tabcolsep) * \real{0.3333}}@{}}
\toprule\noalign{}
\begin{minipage}[b]{\linewidth}\raggedright
Nivel
\end{minipage} & \begin{minipage}[b]{\linewidth}\raggedright
Significado
\end{minipage} & \begin{minipage}[b]{\linewidth}\raggedright
Requiere
\end{minipage} \\
\midrule\noalign{}
\endhead
\bottomrule\noalign{}
\endlastfoot
\textbf{Reproducibilidad de métodos} & Otra persona obtiene los mismos
resultados a partir de los datos y el código originales & Código + datos
+ entorno computacional \\
\textbf{Reproducibilidad de resultados} & Otro investigador, con datos
nuevos, obtiene conclusiones consistentes & Diseño claro + protocolo
público \\
\textbf{Reproducibilidad inferencial} & El conocimiento extraído del
estudio se confirma en revisiones sistemáticas y meta-análisis &
Reporting estandarizado (CONSORT, STROBE\ldots) \\
\end{longtable}

\subsection{D.1.2 Principios FAIR}\label{d.1.2-principios-fair}

Los datos y el código de un estudio reproducible deben ser
\textbf{FAIR}: \emph{Findable} (con identificador persistente, p.~ej.
DOI), \emph{Accessible} (depositados en un repositorio público),
\emph{Interoperable} (formatos abiertos como CSV, parquet) y
\emph{Reusable} (con licencia explícita y documentación).

\subsection{D.1.3 ¿Qué herramientas mínimas
necesitamos?}\label{d.1.3-quuxe9-herramientas-muxednimas-necesitamos}

\begin{itemize}
\tightlist
\item
  \textbf{Control de versiones} para registrar cada cambio del análisis
  (Git).
\item
  \textbf{Plataforma de colaboración} para alojar el repositorio y
  compartirlo (GitHub, GitLab\ldots).
\item
  \textbf{Documentos computacionales} que mezclen código y narrativa
  (Quarto, R Markdown, Jupyter).
\item
  \textbf{Gestión de entornos} que congele las versiones de los paquetes
  (\texttt{renv}, \texttt{conda}).
\end{itemize}

Este apéndice cubre los dos primeros, integrados en RStudio. La gestión
de entornos se introduce brevemente en la sección D.5.

\begin{center}\rule{0.5\linewidth}{0.5pt}\end{center}

\section{D.2 Control de versiones con Git}\label{sec-apendiceD-git}

\textbf{Git} es un sistema distribuido de control de versiones: registra
el historial completo de cambios en un conjunto de archivos (un
\emph{repositorio}) y permite recuperar cualquier estado anterior,
comparar versiones, y trabajar en paralelo en líneas alternativas
(\emph{ramas}) que después se fusionan.

\subsection{D.2.1 Conceptos básicos}\label{d.2.1-conceptos-buxe1sicos}

\begin{longtable}[]{@{}
  >{\raggedright\arraybackslash}p{(\linewidth - 2\tabcolsep) * \real{0.4348}}
  >{\raggedright\arraybackslash}p{(\linewidth - 2\tabcolsep) * \real{0.5652}}@{}}
\toprule\noalign{}
\begin{minipage}[b]{\linewidth}\raggedright
Concepto
\end{minipage} & \begin{minipage}[b]{\linewidth}\raggedright
Significado
\end{minipage} \\
\midrule\noalign{}
\endhead
\bottomrule\noalign{}
\endlastfoot
\textbf{Repositorio} & Carpeta del proyecto bajo control de versiones
(\texttt{.git/} oculto en su raíz) \\
\textbf{Commit} & Instantánea (snapshot) del estado del proyecto en un
momento dado, con un mensaje descriptivo \\
\textbf{Branch} (rama) & Línea de desarrollo paralela; \texttt{main} es
la rama principal \\
\textbf{Merge} & Fusión de los cambios de una rama en otra \\
\textbf{HEAD} & Puntero al commit actualmente activo \\
\textbf{Remoto (\texttt{origin})} & Copia del repositorio alojada en un
servidor (típicamente GitHub) \\
\textbf{\texttt{.gitignore}} & Archivo que lista patrones que Git debe
ignorar (datos crudos, caches, credenciales) \\
\end{longtable}

\subsection{D.2.2 Instalación}\label{d.2.2-instalaciuxf3n}

\begin{tcolorbox}[enhanced jigsaw, breakable, titlerule=0mm, rightrule=.15mm, leftrule=.75mm, left=2mm, opacitybacktitle=0.6, colback=white, bottomtitle=1mm, colbacktitle=quarto-callout-note-color!10!white, title=\textcolor{quarto-callout-note-color}{\faInfo}\hspace{0.5em}{Instalación de Git según el sistema operativo}, toptitle=1mm, coltitle=black, opacityback=0, arc=.35mm, bottomrule=.15mm, toprule=.15mm, colframe=quarto-callout-note-color-frame]

\begin{itemize}
\tightlist
\item
  \textbf{macOS:} preinstalado o \texttt{xcode-select\ -\/-install}
  (Command Line Tools).
\item
  \textbf{Windows:} descargar \textbf{Git for Windows} desde
  \url{https://git-scm.com/download/win} (incluye \emph{Git Bash}).
\item
  \textbf{Linux (Debian/Ubuntu):} \texttt{sudo\ apt\ install\ git}.
\end{itemize}

Verificación:

\begin{Shaded}
\begin{Highlighting}[]
\FunctionTok{git} \AttributeTok{{-}{-}version}
\CommentTok{\# git version 2.42.0 (o superior)}
\end{Highlighting}
\end{Shaded}

\end{tcolorbox}

\subsection{D.2.3 Configuración inicial (una sola vez por
máquina)}\label{d.2.3-configuraciuxf3n-inicial-una-sola-vez-por-muxe1quina}

\begin{Shaded}
\begin{Highlighting}[]
\FunctionTok{git}\NormalTok{ config }\AttributeTok{{-}{-}global}\NormalTok{ user.name  }\StringTok{"Tu Nombre"}
\FunctionTok{git}\NormalTok{ config }\AttributeTok{{-}{-}global}\NormalTok{ user.email }\StringTok{"tu\_email@ugr.es"}
\FunctionTok{git}\NormalTok{ config }\AttributeTok{{-}{-}global}\NormalTok{ init.defaultBranch main}
\end{Highlighting}
\end{Shaded}

Equivalente desde R (recomendado para los lectores de este libro):

\begin{Shaded}
\begin{Highlighting}[]
\CommentTok{\# install.packages("usethis")}
\NormalTok{usethis}\SpecialCharTok{::}\FunctionTok{use\_git\_config}\NormalTok{(}
  \AttributeTok{user.name  =} \StringTok{"Tu Nombre"}\NormalTok{,}
  \AttributeTok{user.email =} \StringTok{"tu\_email@ugr.es"}
\NormalTok{)}
\end{Highlighting}
\end{Shaded}

\subsection{D.2.4 Flujo de trabajo básico (comandos
esenciales)}\label{d.2.4-flujo-de-trabajo-buxe1sico-comandos-esenciales}

\begin{Shaded}
\begin{Highlighting}[]
\FunctionTok{git}\NormalTok{ init                       }\CommentTok{\# Crea un nuevo repositorio en la carpeta actual}
\FunctionTok{git}\NormalTok{ status                     }\CommentTok{\# Muestra archivos modificados/nuevos}
\FunctionTok{git}\NormalTok{ add archivo.qmd            }\CommentTok{\# Marca el archivo para incluirlo en el próximo commit}
\FunctionTok{git}\NormalTok{ add .                      }\CommentTok{\# Marca TODOS los archivos modificados (cuidado con datos sensibles)}
\FunctionTok{git}\NormalTok{ commit }\AttributeTok{{-}m} \StringTok{"Mensaje claro"}  \CommentTok{\# Crea un commit con los cambios marcados}
\FunctionTok{git}\NormalTok{ log }\AttributeTok{{-}{-}oneline}              \CommentTok{\# Historial resumido de commits}
\FunctionTok{git}\NormalTok{ diff archivo.qmd           }\CommentTok{\# Muestra los cambios sin hacer commit}
\FunctionTok{git}\NormalTok{ checkout }\AttributeTok{{-}{-}}\NormalTok{ archivo.qmd    }\CommentTok{\# Descarta los cambios no comprometidos}
\end{Highlighting}
\end{Shaded}

\subsection{D.2.5 Buenas prácticas para mensajes de
commit}\label{d.2.5-buenas-pruxe1cticas-para-mensajes-de-commit}

\begin{tcolorbox}[enhanced jigsaw, breakable, titlerule=0mm, rightrule=.15mm, leftrule=.75mm, left=2mm, opacitybacktitle=0.6, colback=white, bottomtitle=1mm, colbacktitle=quarto-callout-tip-color!10!white, title=\textcolor{quarto-callout-tip-color}{\faLightbulb}\hspace{0.5em}{Reglas para escribir buenos mensajes}, toptitle=1mm, coltitle=black, opacityback=0, arc=.35mm, bottomrule=.15mm, toprule=.15mm, colframe=quarto-callout-tip-color-frame]

\begin{enumerate}
\def\labelenumi{\arabic{enumi}.}
\tightlist
\item
  \textbf{Línea de asunto} en \textbf{imperativo} y \textless{} 70
  caracteres: \emph{``Añade tabla descriptiva de la cohorte''} (no
  \emph{``añadiendo tabla''} ni \emph{``añadido tabla''}).
\item
  \textbf{Una idea por commit:} un mensaje = un cambio coherente. Evita
  los \emph{commits} monstruo que mezclan refactorizaciones con nuevas
  funcionalidades.
\item
  \textbf{Explica el porqué, no el qué.} El \emph{qué} ya lo muestra
  \texttt{git\ diff}; el mensaje debe responder a \emph{por qué} se hizo
  el cambio.
\item
  \textbf{Frecuencia:} mejor muchos \emph{commits} pequeños y
  revertibles que pocos enormes.
\end{enumerate}

Ejemplos:

\begin{itemize}
\tightlist
\item
  ✅
  \texttt{Corrige\ criterio\ de\ inclusión:\ edad\ \textgreater{}=\ 18}
\item
  ✅ \texttt{Añade\ IC\ bootstrap\ al\ modelo\ de\ mortalidad}
\item
  ❌ \texttt{cambios} (vago)
\item
  ❌ \texttt{WIP} (uso interno aceptable, pero no en \texttt{main})
\end{itemize}

\end{tcolorbox}

\subsection{D.2.6 Ramas y fusiones
(introducción)}\label{d.2.6-ramas-y-fusiones-introducciuxf3n}

Las \textbf{ramas} permiten experimentar sin tocar la versión ``buena''
del análisis:

\begin{Shaded}
\begin{Highlighting}[]
\FunctionTok{git}\NormalTok{ switch }\AttributeTok{{-}c}\NormalTok{ sensibilidad      }\CommentTok{\# Crea y cambia a la rama \textquotesingle{}sensibilidad\textquotesingle{}}
\CommentTok{\# … trabajo, commits …}
\FunctionTok{git}\NormalTok{ switch main                 }\CommentTok{\# Vuelve a la rama principal}
\FunctionTok{git}\NormalTok{ merge sensibilidad          }\CommentTok{\# Fusiona los cambios de \textquotesingle{}sensibilidad\textquotesingle{} en main}
\FunctionTok{git}\NormalTok{ branch }\AttributeTok{{-}d}\NormalTok{ sensibilidad      }\CommentTok{\# Elimina la rama una vez fusionada}
\end{Highlighting}
\end{Shaded}

En proyectos colaborativos, las ramas se asocian habitualmente con
\emph{pull requests} en GitHub (ver D.3).

\begin{center}\rule{0.5\linewidth}{0.5pt}\end{center}

\section{D.3 GitHub como plataforma
colaborativa}\label{sec-apendiceD-github}

\textbf{GitHub} (\url{https://github.com}) es un servicio en la nube que
aloja repositorios Git, añadiendo funcionalidades de colaboración:

\begin{itemize}
\tightlist
\item
  \textbf{Repositorios públicos} (visibles y clonables por cualquiera) o
  \textbf{privados}.
\item
  \textbf{Issues:} sistema de tickets para reportar errores, proponer
  mejoras o discutir decisiones.
\item
  \textbf{Pull requests (PR):} mecanismo formal para revisar y fusionar
  cambios entre ramas, con discusión línea a línea.
\item
  \textbf{Releases:} versiones etiquetadas del proyecto (\texttt{v1.0},
  \texttt{v2.0-beta}\ldots) con notas y artefactos descargables.
\item
  \textbf{GitHub Pages:} publicación gratuita de sitios web estáticos
  (el propio libro que está leyendo se sirve mediante GitHub Pages).
\item
  \textbf{GitHub Actions:} automatización (re-renderizar el libro con
  cada commit, lanzar análisis, enviar avisos\ldots).
\end{itemize}

\subsection{D.3.1 Crear una cuenta y un repositorio
remoto}\label{d.3.1-crear-una-cuenta-y-un-repositorio-remoto}

\begin{enumerate}
\def\labelenumi{\arabic{enumi}.}
\tightlist
\item
  Crear cuenta en \url{https://github.com} (gratuita; usuarios
  académicos disponen del plan \textbf{GitHub Education} sin límite).
\item
  Pulsar \textbf{New repository}, asignar nombre (ej.
  \texttt{tfg-osteoporosis}), elegir entre \textbf{Public} o
  \textbf{Private}, y crear sin README inicial si va a conectar un
  repositorio local ya existente.
\end{enumerate}

\subsection{D.3.2 Autenticación: HTTPS + PAT o
SSH}\label{d.3.2-autenticaciuxf3n-https-pat-o-ssh}

GitHub no permite contraseñas en línea de comandos desde 2021. Dos
opciones modernas:

\begin{itemize}
\tightlist
\item
  \textbf{HTTPS + Personal Access Token (PAT):} sencillo para usuarios
  noveles. Crear en \emph{Settings → Developer settings → Personal
  access tokens}. R lo facilita:
\end{itemize}

\begin{Shaded}
\begin{Highlighting}[]
\CommentTok{\# install.packages(c("gitcreds", "usethis"))}
\NormalTok{usethis}\SpecialCharTok{::}\FunctionTok{create\_github\_token}\NormalTok{()  }\CommentTok{\# Abre la página de GitHub con permisos preseleccionados}
\NormalTok{gitcreds}\SpecialCharTok{::}\FunctionTok{gitcreds\_set}\NormalTok{()        }\CommentTok{\# Guarda el token de forma segura en el llavero del sistema}
\end{Highlighting}
\end{Shaded}

\begin{itemize}
\tightlist
\item
  \textbf{SSH:} más cómodo a largo plazo. Generar par de claves con
  \texttt{ssh-keygen\ -t\ ed25519\ -C\ "tu\_email@ugr.es"}, copiar la
  clave pública (\texttt{\textasciitilde{}/.ssh/id\_ed25519.pub}) y
  añadirla en GitHub → \emph{Settings → SSH and GPG keys}.
\end{itemize}

\subsection{D.3.3 Conectar un repositorio local con un
remoto}\label{d.3.3-conectar-un-repositorio-local-con-un-remoto}

\begin{Shaded}
\begin{Highlighting}[]
\CommentTok{\# Una vez creado el repositorio vacío en GitHub:}
\FunctionTok{git}\NormalTok{ remote add origin git@github.com:tu\_usuario/tfg{-}osteoporosis.git}
\FunctionTok{git}\NormalTok{ branch }\AttributeTok{{-}M}\NormalTok{ main}
\FunctionTok{git}\NormalTok{ push }\AttributeTok{{-}u}\NormalTok{ origin main          }\CommentTok{\# \textquotesingle{}origin\textquotesingle{} es el remoto, \textquotesingle{}main\textquotesingle{} la rama}
\end{Highlighting}
\end{Shaded}

A partir de aquí:

\begin{Shaded}
\begin{Highlighting}[]
\FunctionTok{git}\NormalTok{ pull             }\CommentTok{\# Trae los cambios del remoto}
\FunctionTok{git}\NormalTok{ push             }\CommentTok{\# Sube tus commits al remoto}
\end{Highlighting}
\end{Shaded}

\subsection{D.3.4 Pull requests: revisión por pares del
código}\label{d.3.4-pull-requests-revisiuxf3n-por-pares-del-cuxf3digo}

El flujo típico para un cambio sustancial es:

\begin{enumerate}
\def\labelenumi{\arabic{enumi}.}
\tightlist
\item
  Crear una rama: \texttt{git\ switch\ -c\ modelo-logistico}.
\item
  Hacer commits sobre esa rama.
\item
  Publicarla: \texttt{git\ push\ -u\ origin\ modelo-logistico}.
\item
  Abrir un \textbf{pull request} en GitHub desde la rama hacia
  \texttt{main}.
\item
  Discutir, revisar línea a línea, ajustar y, finalmente, fusionar.
\end{enumerate}

Este flujo, equivalente a la revisión por pares de un artículo
científico, mejora la calidad del código y deja registrado el porqué de
cada decisión.

\begin{center}\rule{0.5\linewidth}{0.5pt}\end{center}

\section{D.4 Integración con RStudio}\label{sec-apendiceD-rstudio}

RStudio incluye un \textbf{panel Git} que cubre el 90 \% de las
operaciones cotidianas sin abrir la terminal. Esto facilita enormemente
la adopción de Git para usuarios de R que no se sienten cómodos en línea
de comandos.

\subsection{D.4.1 Crear un proyecto RStudio con Git desde
cero}\label{d.4.1-crear-un-proyecto-rstudio-con-git-desde-cero}

\begin{enumerate}
\def\labelenumi{\arabic{enumi}.}
\tightlist
\item
  \textbf{File → New Project → New Directory → New Project}.
\item
  Marcar la casilla \textbf{Create a git repository}.
\item
  RStudio crea la carpeta, inicializa Git e incluye un
  \texttt{.gitignore} razonable (\texttt{.Rproj.user/},
  \texttt{.Rhistory}, \texttt{.RData}).
\item
  El panel \textbf{Git} aparece en la esquina superior derecha.
\end{enumerate}

\subsection{D.4.2 Conectar un proyecto existente a
Git}\label{d.4.2-conectar-un-proyecto-existente-a-git}

Para un proyecto RStudio sin Git:

\begin{Shaded}
\begin{Highlighting}[]
\CommentTok{\# install.packages("usethis")}
\NormalTok{usethis}\SpecialCharTok{::}\FunctionTok{use\_git}\NormalTok{()           }\CommentTok{\# Inicializa Git y hace el primer commit}
\NormalTok{usethis}\SpecialCharTok{::}\FunctionTok{use\_github}\NormalTok{()        }\CommentTok{\# Crea el repositorio en GitHub y lo conecta (requiere PAT)}
\end{Highlighting}
\end{Shaded}

\subsection{D.4.3 El panel Git de
RStudio}\label{d.4.3-el-panel-git-de-rstudio}

\begin{longtable}[]{@{}ll@{}}
\toprule\noalign{}
Acción en el panel & Equivalente en línea de comandos \\
\midrule\noalign{}
\endhead
\bottomrule\noalign{}
\endlastfoot
Marcar la casilla del archivo (\emph{staged}) &
\texttt{git\ add\ archivo} \\
Botón \textbf{Commit} + mensaje & \texttt{git\ commit\ -m\ "…"} \\
Botón \textbf{Push} (flecha verde arriba) & \texttt{git\ push} \\
Botón \textbf{Pull} (flecha azul abajo) & \texttt{git\ pull} \\
Botón \textbf{History} (reloj) & \texttt{git\ log} (con interfaz
gráfica) \\
Botón \textbf{Diff} & \texttt{git\ diff} (visualización lado a lado) \\
\end{longtable}

\begin{tcolorbox}[enhanced jigsaw, breakable, titlerule=0mm, rightrule=.15mm, leftrule=.75mm, left=2mm, opacitybacktitle=0.6, colback=white, bottomtitle=1mm, colbacktitle=quarto-callout-tip-color!10!white, title=\textcolor{quarto-callout-tip-color}{\faLightbulb}\hspace{0.5em}{Workflow diario en RStudio}, toptitle=1mm, coltitle=black, opacityback=0, arc=.35mm, bottomrule=.15mm, toprule=.15mm, colframe=quarto-callout-tip-color-frame]

\begin{enumerate}
\def\labelenumi{\arabic{enumi}.}
\tightlist
\item
  Antes de empezar: \textbf{Pull} (sincronizar con el remoto).
\item
  Trabajar en el análisis (editar \texttt{.R}, \texttt{.qmd}, etc.).
\item
  Cada cambio coherente: \textbf{Stage} los archivos, \textbf{Commit}
  con un mensaje claro.
\item
  Al final de la sesión: \textbf{Push} (compartir con el equipo).
\end{enumerate}

Repetir varias veces al día: \emph{un commit, una idea}.

\end{tcolorbox}

\subsection{D.4.4 Visualización de cambios y resolución de
conflictos}\label{d.4.4-visualizaciuxf3n-de-cambios-y-resoluciuxf3n-de-conflictos}

El botón \textbf{Diff} en RStudio muestra los cambios línea a línea con
colores (rojo = eliminado, verde = añadido). Cuando dos personas
modifican la misma línea de un archivo, Git señala un \textbf{conflicto}
que se resuelve manualmente editando los marcadores
\texttt{\textless{}\textless{}\textless{}\textless{}\textless{}\textless{}\textless{}},
\texttt{=======} y
\texttt{\textgreater{}\textgreater{}\textgreater{}\textgreater{}\textgreater{}\textgreater{}\textgreater{}}
antes de hacer commit.

\begin{center}\rule{0.5\linewidth}{0.5pt}\end{center}

\section{D.5 Flujos de trabajo reproducibles para
Bioestadística}\label{sec-apendiceD-flujos}

\subsection{D.5.1 Estructura recomendada de un proyecto
bioestadístico}\label{d.5.1-estructura-recomendada-de-un-proyecto-bioestaduxedstico}

\begin{verbatim}
mi-estudio/
├── README.md              # Propósito, autor, citación
├── LICENSE                # CC-BY, MIT, GPL…
├── .gitignore             # Archivos a ignorar
├── mi-estudio.Rproj       # Proyecto RStudio
├── data/
│   ├── raw/               # Datos originales (a menudo gitignored si son sensibles)
│   └── processed/         # Datos derivados, reproducibles desde el script
├── R/                     # Funciones reutilizables
├── analysis/              # Cuadernos Quarto/Rmd con los análisis
├── results/               # Tablas, figuras, modelos serializados (.rds)
├── manuscript/            # Manuscrito en Quarto o LaTeX
└── renv.lock              # Snapshot de las versiones de los paquetes (opcional)
\end{verbatim}

\subsection{\texorpdfstring{D.5.2 \texttt{.gitignore} típico para un
proyecto en
R}{D.5.2 .gitignore típico para un proyecto en R}}\label{d.5.2-.gitignore-tuxedpico-para-un-proyecto-en-r}

\begin{verbatim}
# Artefactos de R
.Rhistory
.RData
.Ruserdata
.Rproj.user/

# Salida de Quarto / RMarkdown
*_files/
*_cache/
/_book/
/_site/

# Sistema operativo
.DS_Store
Thumbs.db

# Datos sensibles (¡nunca commitar PII!)
data/raw/*.csv
data/raw/*.xlsx
\end{verbatim}

\begin{tcolorbox}[enhanced jigsaw, breakable, titlerule=0mm, rightrule=.15mm, leftrule=.75mm, left=2mm, opacitybacktitle=0.6, colback=white, bottomtitle=1mm, colbacktitle=quarto-callout-warning-color!10!white, title=\textcolor{quarto-callout-warning-color}{\faExclamationTriangle}\hspace{0.5em}{Datos sensibles y RGPD}, toptitle=1mm, coltitle=black, opacityback=0, arc=.35mm, bottomrule=.15mm, toprule=.15mm, colframe=quarto-callout-warning-color-frame]

\textbf{Jamás} se debe subir a GitHub información identificable de
pacientes (PII): nombres, DNI, fechas de nacimiento, códigos de historia
clínica, dirección, etc. Aunque elimine el archivo en un commit
posterior, \textbf{permanece en el historial}. Si por accidente sube
datos sensibles, debe utilizar herramientas como
\texttt{git\ filter-repo} o \texttt{BFG\ Repo-Cleaner} y rotar
inmediatamente cualquier credencial expuesta. Para datos clínicos, la
recomendación es trabajar con datos \textbf{anonimizados o
seudonimizados} desde el primer commit.

\end{tcolorbox}

\subsection{\texorpdfstring{D.5.3 Entornos reproducibles con
\texttt{renv}}{D.5.3 Entornos reproducibles con renv}}\label{d.5.3-entornos-reproducibles-con-renv}

Para garantizar que el análisis se ejecuta dentro de cinco años con las
mismas versiones de los paquetes:

\begin{Shaded}
\begin{Highlighting}[]
\CommentTok{\# install.packages("renv")}
\NormalTok{renv}\SpecialCharTok{::}\FunctionTok{init}\NormalTok{()      }\CommentTok{\# Crea un entorno local + renv.lock}
\NormalTok{renv}\SpecialCharTok{::}\FunctionTok{snapshot}\NormalTok{()  }\CommentTok{\# Actualiza renv.lock con las versiones actuales}
\NormalTok{renv}\SpecialCharTok{::}\FunctionTok{restore}\NormalTok{()   }\CommentTok{\# Reinstala exactamente las versiones del lockfile}
\end{Highlighting}
\end{Shaded}

Conviene comprometer (\texttt{git\ commit}) el archivo
\texttt{renv.lock}.

\subsection{D.5.4 Publicar un libro o análisis con GitHub
Pages}\label{d.5.4-publicar-un-libro-o-anuxe1lisis-con-github-pages}

Este libro mismo está alojado en GitHub Pages. El flujo es:

\begin{enumerate}
\def\labelenumi{\arabic{enumi}.}
\tightlist
\item
  En \texttt{\_quarto.yml}: \texttt{project:\ type:\ book} con
  \texttt{output-dir:\ docs}.
\item
  Renderizar localmente: \texttt{quarto\ render} (se genera la carpeta
  \texttt{docs/}).
\item
  Comprometer y subir:
  \texttt{git\ add\ docs/\ \&\&\ git\ commit\ -m\ "Publica\ versión\ X"\ \&\&\ git\ push}.
\item
  En GitHub → \emph{Settings → Pages}, seleccionar \textbf{branch:
  main}, \textbf{folder: /docs}.
\item
  La URL \texttt{https://tu\_usuario.github.io/nombre-repo/} queda
  activa.
\end{enumerate}

Para automatizar la renderización en cada \emph{push} se utiliza
\textbf{GitHub Actions} (ver el archivo
\texttt{.github/workflows/publish.yml} del repositorio de este libro
como ejemplo).

\begin{center}\rule{0.5\linewidth}{0.5pt}\end{center}

\section{D.6 Ejercicio guiado: tu primer repositorio
reproducible}\label{sec-apendiceD-ejercicio}

Crear un repositorio público con un análisis del dataset \texttt{osteo}
que se publique automáticamente como página web.

\begin{tcolorbox}[enhanced jigsaw, breakable, titlerule=0mm, rightrule=.15mm, leftrule=.75mm, left=2mm, opacitybacktitle=0.6, colback=white, bottomtitle=1mm, colbacktitle=quarto-callout-tip-color!10!white, title=\textcolor{quarto-callout-tip-color}{\faLightbulb}\hspace{0.5em}{Pasos}, toptitle=1mm, coltitle=black, opacityback=0, arc=.35mm, bottomrule=.15mm, toprule=.15mm, colframe=quarto-callout-tip-color-frame]

\textbf{1. Crear cuenta y proyecto:}

\begin{itemize}
\tightlist
\item
  Crear cuenta en \url{https://github.com}.
\item
  En RStudio: \textbf{File → New Project → New Directory → Quarto
  Project} (Book/Website/Document).
\item
  Marcar \textbf{Create a git repository}.
\end{itemize}

\textbf{2. Configurar Git e instalar BioStatR:}

\begin{Shaded}
\begin{Highlighting}[]
\NormalTok{usethis}\SpecialCharTok{::}\FunctionTok{use\_git\_config}\NormalTok{(}\AttributeTok{user.name =} \StringTok{"Tu Nombre"}\NormalTok{, }\AttributeTok{user.email =} \StringTok{"tu@ugr.es"}\NormalTok{)}
\NormalTok{remotes}\SpecialCharTok{::}\FunctionTok{install\_github}\NormalTok{(}\StringTok{"migariane/BioStatR"}\NormalTok{)}
\end{Highlighting}
\end{Shaded}

\textbf{3. Crear un análisis con Quarto:}

Archivo \texttt{analisis.qmd} (contenido sugerido):

\begin{verbatim}
---
title: "Análisis exploratorio del dataset osteo"
author: "Tu Nombre"
date: today
format: html
---

```{r}
library(BioStatR)
data(osteo)
freq(osteo$tabaco)
grps(osteo$imc, osteo$sexo)
tabla2x2(fvar = osteo$tabaco, cvar = osteo$osteo_cue, o = osteo)
rls(hba1c ~ tevol, data = osteo)
\end{verbatim}

\begin{verbatim}

**4. Renderizar localmente:** `quarto render analisis.qmd` produce `analisis.html`.

**5. Primer commit:**

```bash
git add .
git commit -m "Primer análisis exploratorio del dataset osteo"
\end{verbatim}

\textbf{6. Crear el repositorio en GitHub y conectarlo:}

\begin{Shaded}
\begin{Highlighting}[]
\NormalTok{usethis}\SpecialCharTok{::}\FunctionTok{use\_github}\NormalTok{(}\AttributeTok{private =} \ConstantTok{FALSE}\NormalTok{)}
\end{Highlighting}
\end{Shaded}

\textbf{7. Activar GitHub Pages} en \emph{Settings → Pages → main /
docs} (o \texttt{/root} si el HTML está en la raíz).

\textbf{8. Verificar:} abrir
\texttt{https://tu\_usuario.github.io/nombre-repo/} y comprobar que el
análisis es accesible públicamente.

\end{tcolorbox}

\begin{tcolorbox}[enhanced jigsaw, breakable, titlerule=0mm, rightrule=.15mm, leftrule=.75mm, left=2mm, opacitybacktitle=0.6, colback=white, bottomtitle=1mm, colbacktitle=quarto-callout-note-color!10!white, title=\textcolor{quarto-callout-note-color}{\faInfo}\hspace{0.5em}{Criterios de éxito}, toptitle=1mm, coltitle=black, opacityback=0, arc=.35mm, bottomrule=.15mm, toprule=.15mm, colframe=quarto-callout-note-color-frame]

\begin{itemize}
\tightlist
\item
  El repositorio aparece en su perfil de GitHub con el archivo
  \texttt{analisis.qmd} y el commit inicial.
\item
  Existe un \texttt{README.md} que describe el propósito del repositorio
  y cita el paquete \texttt{BioStatR}.
\item
  El archivo \texttt{.gitignore} excluye \texttt{*\_files/},
  \texttt{.Rhistory}, \texttt{.Rproj.user/}.
\item
  La página renderizada es accesible desde la URL de GitHub Pages.
\item
  Cualquier persona puede clonar el repositorio (\texttt{git\ clone\ …})
  y reproducir el análisis ejecutando
  \texttt{quarto\ render\ analisis.qmd} tras
  \texttt{remotes::install\_github("migariane/BioStatR")}.
\end{itemize}

\end{tcolorbox}

\begin{center}\rule{0.5\linewidth}{0.5pt}\end{center}

\section{D.7 Recursos adicionales}\label{sec-apendiceD-recursos}

\begin{tcolorbox}[enhanced jigsaw, breakable, titlerule=0mm, rightrule=.15mm, leftrule=.75mm, left=2mm, opacitybacktitle=0.6, colback=white, bottomtitle=1mm, colbacktitle=quarto-callout-note-color!10!white, title=\textcolor{quarto-callout-note-color}{\faInfo}\hspace{0.5em}{Lecturas recomendadas}, toptitle=1mm, coltitle=black, opacityback=0, arc=.35mm, bottomrule=.15mm, toprule=.15mm, colframe=quarto-callout-note-color-frame]

\textbf{Libros y guías (gratuitos y en abierto):}

\begin{itemize}
\tightlist
\item
  \textbf{Bryan, J.} \emph{Happy Git and GitHub for the useR}.
  \url{https://happygitwithr.com} --- guía orientada a usuarios de R.
\item
  \textbf{Chacon, S. \& Straub, B.} \emph{Pro Git} (2.ª ed.).
  \url{https://git-scm.com/book} --- referencia exhaustiva de Git.
\item
  \textbf{Wickham, H. \& Bryan, J.} \emph{R Packages} (2.ª ed.).
  \url{https://r-pkgs.org} --- incluye flujos de trabajo con Git y
  GitHub para desarrollar paquetes de R.
\item
  \textbf{Xie, Y., Dervieux, C. \& Riederer, E.} \emph{R Markdown
  Cookbook}. \url{https://bookdown.org/yihui/rmarkdown-cookbook/}.
\item
  \textbf{Quarto Documentation.} \emph{Publishing with GitHub Pages}.
  \url{https://quarto.org/docs/publishing/github-pages.html}.
\end{itemize}

\textbf{Artículos clave sobre reproducibilidad en biomedicina:}

\begin{itemize}
\tightlist
\item
  Ioannidis, J. P. A. (2005). \emph{Why most published research findings
  are false}. PLoS Medicine, 2(8), e124.
\item
  Goodman, S. N., Fanelli, D. \& Ioannidis, J. P. A. (2016). \emph{What
  does research reproducibility mean?} Science Translational Medicine,
  8(341).
\item
  Peng, R. D. (2011). \emph{Reproducible research in computational
  science}. Science, 334(6060), 1226--1227.
\end{itemize}

\textbf{Iniciativas y manifiestos:}

\begin{itemize}
\tightlist
\item
  \textbf{TOP guidelines} (Transparency and Openness Promotion):
  \url{https://www.cos.io/initiatives/top-guidelines}
\item
  \textbf{FAIR Principles:}
  \url{https://www.go-fair.org/fair-principles/}
\item
  \textbf{EQUATOR Network} (reporting guidelines para investigación
  biomédica): \url{https://www.equator-network.org}
\end{itemize}

\end{tcolorbox}

\begin{tcolorbox}[enhanced jigsaw, breakable, titlerule=0mm, rightrule=.15mm, leftrule=.75mm, left=2mm, opacitybacktitle=0.6, colback=white, bottomtitle=1mm, colbacktitle=quarto-callout-tip-color!10!white, title=\textcolor{quarto-callout-tip-color}{\faLightbulb}\hspace{0.5em}{Cinco hábitos para empezar hoy}, toptitle=1mm, coltitle=black, opacityback=0, arc=.35mm, bottomrule=.15mm, toprule=.15mm, colframe=quarto-callout-tip-color-frame]

\begin{enumerate}
\def\labelenumi{\arabic{enumi}.}
\tightlist
\item
  \textbf{Inicializa Git en cada proyecto nuevo} (incluso TFG y
  prácticas).
\item
  \textbf{Commit pequeño y frecuente} con mensaje claro en imperativo.
\item
  \textbf{Sube a GitHub} tus análisis y código (manteniendo los datos
  sensibles fuera).
\item
  \textbf{Documenta el README}: propósito, datos, instrucciones de
  reproducción.
\item
  \textbf{Cita el código y los datos} con un DOI permanente al publicar
  (Zenodo se integra automáticamente con GitHub).
\end{enumerate}

\end{tcolorbox}

\begin{center}\rule{0.5\linewidth}{0.5pt}\end{center}

\emph{La reproducibilidad no es un destino sino un hábito. Cada commit
es un compromiso con la transparencia, y cada repositorio público es una
contribución al conocimiento compartido.}

\chapter{Examen Final}\label{sec-examen}

\textbf{Asignatura:} Matemáticas para la Estadística
\textbf{Institución:} Universidad de Granada \textbf{Duración:} 2 horas
30 minutos \textbf{Puntuación Total:} 125 puntos (seleccionar 4 de 5
partes = 100 puntos, o responder las 5 para nota extra)

\section{Instrucciones}\label{instrucciones}

\begin{itemize}
\tightlist
\item
  Este examen cubre los 5 bloques: Análisis Exploratorio, Probabilidad,
  Inferencia, Regresión Lineal y Datos Categóricos
\item
  Se permite material de consulta: calculadora, tablas de distribuciones
  y tablas de la distribución \(\chi^2\)
\item
  No se permite: apuntes, Internet, comunicación con otros estudiantes
\item
  Muestre todo el trabajo y cálculos intermedios
\item
  Las respuestas deben estar claramente identificadas
\end{itemize}

\begin{center}\rule{0.5\linewidth}{0.5pt}\end{center}

\section{Parte I: Análisis Exploratorio de Datos (25
puntos)}\label{parte-i-anuxe1lisis-exploratorio-de-datos-25-puntos}

\subsection{Pregunta 1.1 (12 puntos)}\label{pregunta-1.1-12-puntos}

Se recopilaron datos de altura (cm) de 50 estudiantes de primer año:

\begin{verbatim}
165, 168, 171, 160, 175, 172, 169, 166, 173, 170, 167, 164, 176, 168, 172,
165, 169, 171, 174, 168, 170, 166, 173, 167, 175, 169, 168, 172, 171, 165,
174, 170, 167, 173, 169, 166, 172, 168, 171, 175, 170, 169, 173, 167, 174,
168, 172, 170, 175, 171
\end{verbatim}

\begin{enumerate}
\def\labelenumi{\alph{enumi})}
\tightlist
\item
  Calcule la media, mediana, desviación estándar y rango
  intercuartílico.
\item
  Construya una tabla de frecuencias con 5 intervalos de clase.
\item
  Interprete la distribución (simetría, dispersión).
\end{enumerate}

Mostrar respuesta 1.1 (solo instructor)

\textbf{a) Estadísticos descriptivos:} - Media:
\(\bar{x} = \frac{\sum x_i}{n} \approx 170.1\) cm - Mediana:
\(Q_2 \approx 170\) cm (valor central de los 50 datos ordenados) -
Desviación estándar: \(s \approx 3.8\) cm - Rango intercuartílico:
\(IQR = Q_3 - Q_1 \approx 173 - 167 = 6\) cm

\textbf{b) Tabla de frecuencias (5 intervalos, amplitud ≈ 4 cm):}

\begin{longtable}[]{@{}llll@{}}
\toprule\noalign{}
Intervalo & Frecuencia & Frecuencia relativa & Frecuencia acumulada \\
\midrule\noalign{}
\endhead
\bottomrule\noalign{}
\endlastfoot
160-164 & 3 & 0.06 & 0.06 \\
164-168 & 12 & 0.24 & 0.30 \\
168-172 & 22 & 0.44 & 0.74 \\
172-176 & 13 & 0.26 & 1.00 \\
\end{longtable}

\textbf{c) Interpretación:} La distribución es aproximadamente
simétrica, centrada alrededor de 170 cm, con una concentración principal
en 168-172 cm (44\% de los datos). La dispersión es moderada (s ≈ 3.8
cm). No hay valores atípicos evidentes.

\subsection{Pregunta 1.2 (13 puntos)}\label{pregunta-1.2-13-puntos}

Se estudia la relación entre número de horas de estudio (X) y
calificación en examen (Y) de 10 estudiantes:

\begin{longtable}[]{@{}lllllllllll@{}}
\toprule\noalign{}
Horas & 2 & 3 & 2.5 & 4 & 3.5 & 5 & 1.5 & 4.5 & 3 & 2 \\
\midrule\noalign{}
\endhead
\bottomrule\noalign{}
\endlastfoot
Nota & 4 & 5 & 4.5 & 6 & 5.5 & 7 & 3 & 6.5 & 5 & 4 \\
\end{longtable}

\begin{enumerate}
\def\labelenumi{\alph{enumi})}
\tightlist
\item
  Calcule el coeficiente de correlación de Pearson entre X e Y.
\item
  Interprete el resultado.
\item
  ¿Pueden haber valores atípicos?
\end{enumerate}

Mostrar respuesta 1.2 (solo instructor)

\textbf{a) Coeficiente de Pearson:}

\[\text{Primero: } \bar{x} = 3.1, \quad \bar{y} = 5.0\]

\[s_x = 1.23, \quad s_y = 1.27\]

\[\text{Covarianza: } \text{Cov}(X,Y) = \frac{\sum(x_i - \bar{x})(y_i - \bar{y})}{n-1} \approx 1.40\]

\[r = \frac{\text{Cov}(X,Y)}{s_x s_y} = \frac{1.40}{1.23 \times 1.27} \approx 0.90\]

\textbf{b) Interpretación:} Existe una correlación positiva muy fuerte
(r ≈ 0.90) entre horas de estudio y nota del examen. A mayor número de
horas, mayor calificación. Aproximadamente 81\% (r²) de la variación en
notas se explica por las horas de estudio.

\textbf{c) Valores atípicos:} La relación es lineal sin desvíos
aparentes. No hay evidencia de outliers en los datos.

\begin{center}\rule{0.5\linewidth}{0.5pt}\end{center}

\section{Parte II: Probabilidad y Variables Aleatorias (25
puntos)}\label{parte-ii-probabilidad-y-variables-aleatorias-25-puntos}

\subsection{Pregunta 2.1 (12 puntos)}\label{pregunta-2.1-12-puntos}

Un examen tiene 20 preguntas de verdadero/falso. Un estudiante responde
al azar (50\% de probabilidad de acierto).

\begin{enumerate}
\def\labelenumi{\alph{enumi})}
\tightlist
\item
  ¿Cuál es la probabilidad de acertar exactamente 12 preguntas?
\item
  ¿Cuál es la probabilidad de acertar 15 o más preguntas?
\item
  ¿Cuál es el número esperado de aciertos y su desviación estándar?
\end{enumerate}

Mostrar respuesta 2.1 (solo instructor)

Sea \(X \sim \text{Binomial}(n=20, p=0.5)\)

\textbf{a) P(X = 12):}

\[P(X = 12) = \binom{20}{12} (0.5)^{12}(0.5)^{8} = \binom{20}{12} (0.5)^{20}\]

\[= 125970 \times (0.5)^{20} \approx 0.1201\]

\textbf{b) P(X ≥ 15):}

\[P(X \geq 15) = \sum_{k=15}^{20} \binom{20}{k}(0.5)^{20}\]

Por tabla o software: \(P(X \geq 15) \approx 0.0207\)

\textbf{c) Esperanza y desviación:}

\[E[X] = np = 20 \times 0.5 = 10\]

\[\sigma = \sqrt{np(1-p)} = \sqrt{20 \times 0.5 \times 0.5} = \sqrt{5} \approx 2.24\]

\subsection{Pregunta 2.2 (13 puntos)}\label{pregunta-2.2-13-puntos}

El tiempo de atención en una consulta médica sigue distribución normal
con media 15 minutos y desviación estándar 3 minutos.

\begin{enumerate}
\def\labelenumi{\alph{enumi})}
\tightlist
\item
  ¿Qué porcentaje de consultas duran más de 20 minutos?
\item
  ¿Cuál es el percentil 90 del tiempo de consulta?
\item
  Si se observan 100 consultas, ¿cuántas esperamos que duren entre 12 y
  18 minutos?
\end{enumerate}

Mostrar respuesta 2.2 (solo instructor)

Sea \(X \sim \mathcal{N}(\mu=15, \sigma=3)\)

\textbf{a) P(X \textgreater{} 20):}

\[Z = \frac{20 - 15}{3} = 1.67\]

\[P(X > 20) = P(Z > 1.67) \approx 0.0475 \text{ o } 4.75\%\]

\textbf{b) Percentil 90:}

Necesitamos \(z_{0.90} \approx 1.28\) (de tabla normal estándar)

\[x_{90} = 15 + 1.28 \times 3 = 15 + 3.84 = 18.84 \text{ minutos}\]

\textbf{c) Número de consultas en {[}12, 18{]}:}

\[Z_1 = \frac{12-15}{3} = -1 \quad Z_2 = \frac{18-15}{3} = 1\]

\[P(12 < X < 18) = P(-1 < Z < 1) \approx 0.6826 \text{ o } 68.26\%\]

Esperadas: \(100 \times 0.6826 \approx 68\) consultas

\begin{center}\rule{0.5\linewidth}{0.5pt}\end{center}

\section{Parte III: Inferencia Estadística (25
puntos)}\label{parte-iii-inferencia-estaduxedstica-25-puntos}

\subsection{Pregunta 3.1 (12 puntos)}\label{pregunta-3.1-12-puntos}

Se realizó un estudio de presión arterial sistólica en una muestra de 40
pacientes hipertensos. La media muestral fue 145 mmHg y la desviación
estándar muestral 12 mmHg. Asuma distribución normal.

\begin{enumerate}
\def\labelenumi{\alph{enumi})}
\tightlist
\item
  Calcule un intervalo de confianza del 95\% para la media poblacional.
\item
  Interprete el intervalo.
\item
  ¿Cuál sería el tamaño muestral necesario para un margen de error de 3
  mmHg?
\end{enumerate}

Mostrar respuesta 3.1 (solo instructor)

\textbf{a) Intervalo de confianza 95\%:}

Con \(n=40\), \(\bar{x}=145\), \(s=12\), \(\alpha=0.05\):

Grados de libertad: \(df = 39\)

\(t_{0.025, 39} \approx 2.023\) (de tabla t)

Margen de error:
\(E = t_{\alpha/2} \times \frac{s}{\sqrt{n}} = 2.023 \times \frac{12}{\sqrt{40}} = 2.023 \times 1.897 \approx 3.84\)

\[IC_{95\%} = [145 - 3.84, 145 + 3.84] = [141.16, 148.84] \text{ mmHg}\]

\textbf{b) Interpretación:} Con 95\% de confianza, la presión arterial
sistólica media poblacional se encuentra entre 141.16 y 148.84 mmHg. Si
repitiéramos el estudio muchas veces, el 95\% de los intervalos
construidos contendrían el parámetro verdadero.

\textbf{c) Tamaño muestral para E = 3:}

\[n = \left(\frac{t_{\alpha/2} \times s}{E}\right)^2 = \left(\frac{2.023 \times 12}{3}\right)^2 \approx \left(\frac{24.276}{3}\right)^2 \approx 65.4\]

Se necesitan aproximadamente \textbf{n = 66 pacientes}.

\subsection{Pregunta 3.2 (13 puntos)}\label{pregunta-3.2-13-puntos}

Un laboratorio afirma que el 90\% de sus pruebas de COVID tienen
sensibilidad ≥ 95\%. Para verificar, se prueban 50 muestras positivas y
se detectan correctamente 47.

\begin{enumerate}
\def\labelenumi{\alph{enumi})}
\tightlist
\item
  Plantee las hipótesis nula y alternativa.
\item
  Calcule el estadístico de prueba y el p-valor.
\item
  ¿Rechaza la afirmación del laboratorio al nivel α = 0.05?
\end{enumerate}

Mostrar respuesta 3.2 (solo instructor)

\textbf{a) Hipótesis:}

\begin{itemize}
\tightlist
\item
  \(H_0: p = 0.95\) (la sensibilidad es del 95\%)
\item
  \(H_A: p < 0.95\) (la sensibilidad es menor al 95\%, test unilateral)
\end{itemize}

\textbf{b) Estadístico de prueba:}

Proporción muestral: \(\hat{p} = \frac{47}{50} = 0.94\)

Error estándar bajo \(H_0\):
\(SE = \sqrt{\frac{p(1-p)}{n}} = \sqrt{\frac{0.95 \times 0.05}{50}} \approx 0.0308\)

\[z = \frac{\hat{p} - p_0}{SE} = \frac{0.94 - 0.95}{0.0308} \approx -0.325\]

p-valor (unilateral): \(P(Z < -0.325) \approx 0.372\)

\textbf{c) Conclusión:} Con p-valor = 0.372 \textgreater{} α = 0.05,
\textbf{no rechazamos} \(H_0\). No hay evidencia suficiente para
rechazar la afirmación del laboratorio de que la sensibilidad es del
95\%.

\begin{center}\rule{0.5\linewidth}{0.5pt}\end{center}

\section{Parte IV: Regresión Lineal (25
puntos)}\label{parte-iv-regresiuxf3n-lineal-25-puntos}

\subsection{Pregunta 4.1 (25 puntos)}\label{pregunta-4.1-25-puntos}

Se estudia la relación entre precio de venta (Y, en miles €) y tamaño de
la vivienda (X, en m²) en 8 pisos vendidos:

\begin{longtable}[]{@{}lllllllll@{}}
\toprule\noalign{}
Tamaño (m²) & 80 & 100 & 90 & 120 & 110 & 95 & 85 & 130 \\
\midrule\noalign{}
\endhead
\bottomrule\noalign{}
\endlastfoot
Precio (k€) & 150 & 185 & 170 & 210 & 200 & 180 & 165 & 230 \\
\end{longtable}

\begin{enumerate}
\def\labelenumi{\alph{enumi})}
\tightlist
\item
  Estime el modelo de regresión lineal simple
  \(Y = \beta_0 + \beta_1 X + \epsilon\)
\item
  Interprete los coeficientes
\item
  Calcule R² e interprete la calidad del ajuste
\item
  Prediga el precio para una vivienda de 105 m²
\end{enumerate}

Mostrar respuesta 4.1 (solo instructor)

\textbf{a) Estimación del modelo:}

Primero calculamos estadísticos: - \(\bar{x} = 102.5\),
\(\bar{y} = 186.25\) - \(S_{xx} = \sum(x_i - \bar{x})^2 = 1512.5\) -
\(S_{xy} = \sum(x_i - \bar{x})(y_i - \bar{y}) = 3057.5\) -
\(S_{yy} = \sum(y_i - \bar{y})^2 = 6337.5\)

Coeficiente de regresión:
\[\hat{\beta}_1 = \frac{S_{xy}}{S_{xx}} = \frac{3057.5}{1512.5} \approx 2.021\]

Ordenada en el origen:
\[\hat{\beta}_0 = \bar{y} - \hat{\beta}_1\bar{x} = 186.25 - 2.021 \times 102.5 \approx -22.41\]

\textbf{Modelo estimado:} \[\hat{Y} = -22.41 + 2.021 X\]

\textbf{b) Interpretación de coeficientes:} -
\(\hat{\beta}_0 = -22.41\): Es el precio teórico para X = 0 (sin
interpretación práctica, pues no hay viviendas de 0 m²) -
\(\hat{\beta}_1 = 2.021\): Por cada m² adicional, el precio aumenta
aproximadamente 2.021 mil euros, o sea, 2021 € por m²

\textbf{c) Coeficiente de determinación R²:}

\[r = \frac{S_{xy}}{\sqrt{S_{xx} \times S_{yy}}} = \frac{3057.5}{\sqrt{1512.5 \times 6337.5}} \approx 0.9815\]

\[R^2 = r^2 \approx 0.963\]

Interpretación: El 96.3\% de la variabilidad en los precios se explica
por el tamaño de la vivienda. El ajuste es excelente.

\textbf{d) Predicción para X = 105 m²:}

\[\hat{Y}(105) = -22.41 + 2.021 \times 105 = -22.41 + 212.205 \approx 189.80 \text{ k€}\]

Precio predicho: aproximadamente \textbf{189,800 €} para una vivienda de
105 m².

\begin{center}\rule{0.5\linewidth}{0.5pt}\end{center}

\section{Parte V: Análisis de Datos Categóricos (25
puntos)}\label{parte-v-anuxe1lisis-de-datos-categuxf3ricos-25-puntos}

\subsection{Pregunta 5.1 (12 puntos)}\label{pregunta-5.1-12-puntos}

En un estudio sobre la asociación entre consumo de tabaco y enfermedad
cardiovascular se obtiene la siguiente tabla:

\begin{longtable}[]{@{}lccc@{}}
\toprule\noalign{}
& Enfermedad cardiovascular & Sin enfermedad & Total \\
\midrule\noalign{}
\endhead
\bottomrule\noalign{}
\endlastfoot
Fumador & 55 & 95 & 150 \\
No fumador & 25 & 125 & 150 \\
\end{longtable}

\begin{enumerate}
\def\labelenumi{\alph{enumi})}
\tightlist
\item
  Calcule las frecuencias esperadas bajo \(H_0\) de independencia.
\item
  Calcule el estadístico \(\chi^2\) e identifique la distribución de
  referencia.
\item
  Con \(\alpha = 0.05\) (\(\chi^2_{0.05,1} = 3.84\)), ¿existe asociación
  estadísticamente significativa?
\item
  ¿Se cumplen las condiciones de aplicación del test Chi-cuadrado?
\end{enumerate}

Mostrar respuesta 5.1 (solo instructor)

\textbf{a) Frecuencias esperadas:}

Marginales: fila 1 = 150, fila 2 = 150, columna 1 = 80, columna 2 = 220,
\(n\) = 300.

\[E_{11} = \frac{150 \times 80}{300} = 40, \quad E_{12} = \frac{150 \times 220}{300} = 110\]
\[E_{21} = \frac{150 \times 80}{300} = 40, \quad E_{22} = \frac{150 \times 220}{300} = 110\]

\textbf{b) Estadístico χ²:}

\[\chi^2 = \frac{(55-40)^2}{40} + \frac{(95-110)^2}{110} + \frac{(25-40)^2}{40} + \frac{(125-110)^2}{110}\]

\[= \frac{225}{40} + \frac{225}{110} + \frac{225}{40} + \frac{225}{110} = 5.625 + 2.045 + 5.625 + 2.045 = 15.34\]

Distribución de referencia: \(\chi^2_{(2-1)(2-1)} = \chi^2_1\).

\textbf{c) Decisión:}

\(\chi^2 = 15.34 > 3.84 = \chi^2_{0.05,1}\), \(p < 0.001\).
\textbf{Rechazamos \(H_0\)}: existe asociación estadísticamente
significativa entre tabaquismo y enfermedad cardiovascular.

\textbf{d) Condiciones:} Mínima frecuencia esperada = 40 \textgreater{}
5. ✓ Condición cumplida.

\subsection{Pregunta 5.2 (13 puntos)}\label{pregunta-5.2-13-puntos}

Un estudio caso-control investiga la asociación entre obesidad (IMC ≥
30) e infarto de miocardio. Se reclutan 100 casos (infarto) y 100
controles. De los casos, 60 son obesos; de los controles, 30 son obesos.

\begin{enumerate}
\def\labelenumi{\alph{enumi})}
\tightlist
\item
  Construya la tabla 2×2 con la notación estándar \((a, b, c, d)\).
\item
  Calcule el Odds Ratio (OR).
\item
  Calcule el intervalo de confianza del 95\% para el OR.
\item
  Interprete el resultado en contexto clínico. ¿Sería correcto calcular
  el Riesgo Relativo (RR) en este diseño? Justifique.
\end{enumerate}

Mostrar respuesta 5.2 (solo instructor)

\textbf{a) Tabla 2×2:}

\begin{longtable}[]{@{}lccc@{}}
\toprule\noalign{}
& Caso (infarto) & Control & Total \\
\midrule\noalign{}
\endhead
\bottomrule\noalign{}
\endlastfoot
Obeso & 60 & 30 & 90 \\
No obeso & 40 & 70 & 110 \\
Total & 100 & 100 & 200 \\
\end{longtable}

\(a=60\), \(b=30\), \(c=40\), \(d=70\).

\textbf{b) Odds Ratio:}

\[\widehat{OR} = \frac{ad}{bc} = \frac{60 \times 70}{30 \times 40} = \frac{4200}{1200} = 3.50\]

\textbf{c) Intervalo de confianza del 95\%:}

\[SE_{\ln OR} = \sqrt{\frac{1}{60} + \frac{1}{30} + \frac{1}{40} + \frac{1}{70}} = \sqrt{0.0167 + 0.0333 + 0.0250 + 0.0143} = \sqrt{0.0893} = 0.299\]

\[IC_{95\%} = \exp(\ln 3.50 \pm 1.96 \times 0.299) = \exp(1.253 \pm 0.586) = [e^{0.667},\, e^{1.839}] = [1.95,\, 6.29]\]

\textbf{d) Interpretación:} Las personas obesas tienen \textbf{3.5 veces
más odds} de infarto de miocardio que las no obesas (\(OR = 3.50\);
\(IC_{95\%}\): 1.95--6.29). El intervalo no contiene el 1, confirmando
asociación estadísticamente significativa.

\textbf{RR en caso-control:} No es correcto calcular el RR porque en un
estudio caso-control el número de casos y controles es fijado por el
investigador (100 y 100 respectivamente), y no refleja la prevalencia o
incidencia real de infarto en la población. Los totales de fila no son
proporcionales a los totales poblacionales, por lo que las
``proporciones'' calculadas como \(a/(a+c)\) no son estimaciones válidas
del riesgo. La única medida estimable de forma válida en este diseño es
el OR.

\begin{center}\rule{0.5\linewidth}{0.5pt}\end{center}

\section{Instrucciones Finales}\label{instrucciones-finales}

\begin{itemize}
\tightlist
\item
  Revise sus respuestas antes de entregar
\item
  Indique claramente los pasos de sus cálculos
\item
  Si usa software o calculadora, muestre los resultados intermedios
\item
  La precisión y claridad contribuyen a la calificación
\end{itemize}

\textbf{Fin del examen}


\backmatter



\end{document}
